% Juge de la troisième passe : document minimal écrit de zéro pour la cause A (thmtools sous LaTeX 2026-06-01).
% Variantes, par une macro posée sur la ligne de commande (deux passes chacune, depuis memoire/essai/juge3/minimal/) :
%   lualatex -jobname=nu       "% Juge de la troisième passe : document minimal écrit de zéro pour la cause A (thmtools sous LaTeX 2026-06-01).
% Variantes, par une macro posée sur la ligne de commande (deux passes chacune, depuis memoire/essai/juge3/minimal/) :
%   lualatex -jobname=nu       "% Juge de la troisième passe : document minimal écrit de zéro pour la cause A (thmtools sous LaTeX 2026-06-01).
% Variantes, par une macro posée sur la ligne de commande (deux passes chacune, depuis memoire/essai/juge3/minimal/) :
%   lualatex -jobname=nu       "% Juge de la troisième passe : document minimal écrit de zéro pour la cause A (thmtools sous LaTeX 2026-06-01).
% Variantes, par une macro posée sur la ligne de commande (deux passes chacune, depuis memoire/essai/juge3/minimal/) :
%   lualatex -jobname=nu       "\input{minimal-juge3}"                                   sans correctif
%   lualatex -jobname=corr     "\def\CORR{1}\input{minimal-juge3}"                       correctif A du préambule
%   lualatex -jobname=crochet  "\def\SANSCROCHET{1}\input{minimal-juge3}"                sans correctif, crochet retiré
%   lualatex -jobname=amsseul  "\def\AMSSEUL{1}\input{minimal-juge3}"                    amsthm seul (témoin)
%   lualatex -jobname=faux     "\def\CORR{1}\def\FAUX{1}\input{minimal-juge3}"           deux vrais défauts : doivent tomber
%   lualatex -jobname=trace    "\def\TRACE{1}\input{minimal-juge3}"                      trace de l'appel
% Les lignes du correctif sont tirées du préambule commité par extraire_correctifs.py, pas retapées.
\documentclass{scrbook}
\ifdefined\SANSCROCHET \RemoveFromHook{package/amsthm/after}[firstaid/aliascounter]\fi
\usepackage{amsthm}
\usepackage{hyperref}
\usepackage[french, nameinlink, noabbrev]{cleveref}
\ifdefined\AMSSEUL
  \newtheorem{theoreme}{Théorème}[chapter]
  \newtheorem{proposition}[theoreme]{Proposition}
  \newtheorem{corollaire}[theoreme]{Corollaire}
  \newtheorem{remarque}{Remarque}[chapter]
\else
  \usepackage{thmtools}
  \ifdefined\CORR \input{correctif-a-pose}\fi
  \declaretheorem[numberwithin=chapter, name=Théorème]{theoreme}
  \ifdefined\TRACE \tracingmacros=2 \tracingonline=0 \fi
  \declaretheorem[sibling=theoreme, name=Proposition]{proposition}
  \ifdefined\TRACE \tracingmacros=0 \fi
  \declaretheorem[sibling=theoreme, name=Corollaire]{corollaire}
  \declaretheorem[numberwithin=chapter, name=Remarque]{remarque}
  \ifdefined\FAUX \newcounteralias{corollaire}{section}\newcounteralias{orphelin}{inexistant}\fi
  \ifdefined\CORR \input{correctif-a-remise}\fi
\fi
\crefname{theoreme}{théorème}{théorèmes}
\crefname{proposition}{proposition}{propositions}
\crefname{corollaire}{corollaire}{corollaires}
\crefname{remarque}{remarque}{remarques}
\begin{document}
\chapter{Premier}
\begin{proposition}\label{p1}a\end{proposition}
\begin{theoreme}\label{t1}b\end{theoreme}
\begin{remarque}\label{r1}c\end{remarque}
\begin{corollaire}\label{c1}d\end{corollaire}
\chapter{Second}
\begin{corollaire}\label{c2}e\end{corollaire}
\begin{theoreme}\label{t2}f\end{theoreme}
RENVOIS \ref{p1} \ref{t1} \ref{r1} \ref{c1} \ref{c2} \ref{t2} ; \cref{p1} ; \cref{t1} ; \cref{r1} ; \cref{c1} ; \cref{c2} ; \cref{t2} FIN
\end{document}
"                                   sans correctif
%   lualatex -jobname=corr     "\def\CORR{1}% Juge de la troisième passe : document minimal écrit de zéro pour la cause A (thmtools sous LaTeX 2026-06-01).
% Variantes, par une macro posée sur la ligne de commande (deux passes chacune, depuis memoire/essai/juge3/minimal/) :
%   lualatex -jobname=nu       "\input{minimal-juge3}"                                   sans correctif
%   lualatex -jobname=corr     "\def\CORR{1}\input{minimal-juge3}"                       correctif A du préambule
%   lualatex -jobname=crochet  "\def\SANSCROCHET{1}\input{minimal-juge3}"                sans correctif, crochet retiré
%   lualatex -jobname=amsseul  "\def\AMSSEUL{1}\input{minimal-juge3}"                    amsthm seul (témoin)
%   lualatex -jobname=faux     "\def\CORR{1}\def\FAUX{1}\input{minimal-juge3}"           deux vrais défauts : doivent tomber
%   lualatex -jobname=trace    "\def\TRACE{1}\input{minimal-juge3}"                      trace de l'appel
% Les lignes du correctif sont tirées du préambule commité par extraire_correctifs.py, pas retapées.
\documentclass{scrbook}
\ifdefined\SANSCROCHET \RemoveFromHook{package/amsthm/after}[firstaid/aliascounter]\fi
\usepackage{amsthm}
\usepackage{hyperref}
\usepackage[french, nameinlink, noabbrev]{cleveref}
\ifdefined\AMSSEUL
  \newtheorem{theoreme}{Théorème}[chapter]
  \newtheorem{proposition}[theoreme]{Proposition}
  \newtheorem{corollaire}[theoreme]{Corollaire}
  \newtheorem{remarque}{Remarque}[chapter]
\else
  \usepackage{thmtools}
  \ifdefined\CORR \input{correctif-a-pose}\fi
  \declaretheorem[numberwithin=chapter, name=Théorème]{theoreme}
  \ifdefined\TRACE \tracingmacros=2 \tracingonline=0 \fi
  \declaretheorem[sibling=theoreme, name=Proposition]{proposition}
  \ifdefined\TRACE \tracingmacros=0 \fi
  \declaretheorem[sibling=theoreme, name=Corollaire]{corollaire}
  \declaretheorem[numberwithin=chapter, name=Remarque]{remarque}
  \ifdefined\FAUX \newcounteralias{corollaire}{section}\newcounteralias{orphelin}{inexistant}\fi
  \ifdefined\CORR \input{correctif-a-remise}\fi
\fi
\crefname{theoreme}{théorème}{théorèmes}
\crefname{proposition}{proposition}{propositions}
\crefname{corollaire}{corollaire}{corollaires}
\crefname{remarque}{remarque}{remarques}
\begin{document}
\chapter{Premier}
\begin{proposition}\label{p1}a\end{proposition}
\begin{theoreme}\label{t1}b\end{theoreme}
\begin{remarque}\label{r1}c\end{remarque}
\begin{corollaire}\label{c1}d\end{corollaire}
\chapter{Second}
\begin{corollaire}\label{c2}e\end{corollaire}
\begin{theoreme}\label{t2}f\end{theoreme}
RENVOIS \ref{p1} \ref{t1} \ref{r1} \ref{c1} \ref{c2} \ref{t2} ; \cref{p1} ; \cref{t1} ; \cref{r1} ; \cref{c1} ; \cref{c2} ; \cref{t2} FIN
\end{document}
"                       correctif A du préambule
%   lualatex -jobname=crochet  "\def\SANSCROCHET{1}% Juge de la troisième passe : document minimal écrit de zéro pour la cause A (thmtools sous LaTeX 2026-06-01).
% Variantes, par une macro posée sur la ligne de commande (deux passes chacune, depuis memoire/essai/juge3/minimal/) :
%   lualatex -jobname=nu       "\input{minimal-juge3}"                                   sans correctif
%   lualatex -jobname=corr     "\def\CORR{1}\input{minimal-juge3}"                       correctif A du préambule
%   lualatex -jobname=crochet  "\def\SANSCROCHET{1}\input{minimal-juge3}"                sans correctif, crochet retiré
%   lualatex -jobname=amsseul  "\def\AMSSEUL{1}\input{minimal-juge3}"                    amsthm seul (témoin)
%   lualatex -jobname=faux     "\def\CORR{1}\def\FAUX{1}\input{minimal-juge3}"           deux vrais défauts : doivent tomber
%   lualatex -jobname=trace    "\def\TRACE{1}\input{minimal-juge3}"                      trace de l'appel
% Les lignes du correctif sont tirées du préambule commité par extraire_correctifs.py, pas retapées.
\documentclass{scrbook}
\ifdefined\SANSCROCHET \RemoveFromHook{package/amsthm/after}[firstaid/aliascounter]\fi
\usepackage{amsthm}
\usepackage{hyperref}
\usepackage[french, nameinlink, noabbrev]{cleveref}
\ifdefined\AMSSEUL
  \newtheorem{theoreme}{Théorème}[chapter]
  \newtheorem{proposition}[theoreme]{Proposition}
  \newtheorem{corollaire}[theoreme]{Corollaire}
  \newtheorem{remarque}{Remarque}[chapter]
\else
  \usepackage{thmtools}
  \ifdefined\CORR \input{correctif-a-pose}\fi
  \declaretheorem[numberwithin=chapter, name=Théorème]{theoreme}
  \ifdefined\TRACE \tracingmacros=2 \tracingonline=0 \fi
  \declaretheorem[sibling=theoreme, name=Proposition]{proposition}
  \ifdefined\TRACE \tracingmacros=0 \fi
  \declaretheorem[sibling=theoreme, name=Corollaire]{corollaire}
  \declaretheorem[numberwithin=chapter, name=Remarque]{remarque}
  \ifdefined\FAUX \newcounteralias{corollaire}{section}\newcounteralias{orphelin}{inexistant}\fi
  \ifdefined\CORR \input{correctif-a-remise}\fi
\fi
\crefname{theoreme}{théorème}{théorèmes}
\crefname{proposition}{proposition}{propositions}
\crefname{corollaire}{corollaire}{corollaires}
\crefname{remarque}{remarque}{remarques}
\begin{document}
\chapter{Premier}
\begin{proposition}\label{p1}a\end{proposition}
\begin{theoreme}\label{t1}b\end{theoreme}
\begin{remarque}\label{r1}c\end{remarque}
\begin{corollaire}\label{c1}d\end{corollaire}
\chapter{Second}
\begin{corollaire}\label{c2}e\end{corollaire}
\begin{theoreme}\label{t2}f\end{theoreme}
RENVOIS \ref{p1} \ref{t1} \ref{r1} \ref{c1} \ref{c2} \ref{t2} ; \cref{p1} ; \cref{t1} ; \cref{r1} ; \cref{c1} ; \cref{c2} ; \cref{t2} FIN
\end{document}
"                sans correctif, crochet retiré
%   lualatex -jobname=amsseul  "\def\AMSSEUL{1}% Juge de la troisième passe : document minimal écrit de zéro pour la cause A (thmtools sous LaTeX 2026-06-01).
% Variantes, par une macro posée sur la ligne de commande (deux passes chacune, depuis memoire/essai/juge3/minimal/) :
%   lualatex -jobname=nu       "\input{minimal-juge3}"                                   sans correctif
%   lualatex -jobname=corr     "\def\CORR{1}\input{minimal-juge3}"                       correctif A du préambule
%   lualatex -jobname=crochet  "\def\SANSCROCHET{1}\input{minimal-juge3}"                sans correctif, crochet retiré
%   lualatex -jobname=amsseul  "\def\AMSSEUL{1}\input{minimal-juge3}"                    amsthm seul (témoin)
%   lualatex -jobname=faux     "\def\CORR{1}\def\FAUX{1}\input{minimal-juge3}"           deux vrais défauts : doivent tomber
%   lualatex -jobname=trace    "\def\TRACE{1}\input{minimal-juge3}"                      trace de l'appel
% Les lignes du correctif sont tirées du préambule commité par extraire_correctifs.py, pas retapées.
\documentclass{scrbook}
\ifdefined\SANSCROCHET \RemoveFromHook{package/amsthm/after}[firstaid/aliascounter]\fi
\usepackage{amsthm}
\usepackage{hyperref}
\usepackage[french, nameinlink, noabbrev]{cleveref}
\ifdefined\AMSSEUL
  \newtheorem{theoreme}{Théorème}[chapter]
  \newtheorem{proposition}[theoreme]{Proposition}
  \newtheorem{corollaire}[theoreme]{Corollaire}
  \newtheorem{remarque}{Remarque}[chapter]
\else
  \usepackage{thmtools}
  \ifdefined\CORR \input{correctif-a-pose}\fi
  \declaretheorem[numberwithin=chapter, name=Théorème]{theoreme}
  \ifdefined\TRACE \tracingmacros=2 \tracingonline=0 \fi
  \declaretheorem[sibling=theoreme, name=Proposition]{proposition}
  \ifdefined\TRACE \tracingmacros=0 \fi
  \declaretheorem[sibling=theoreme, name=Corollaire]{corollaire}
  \declaretheorem[numberwithin=chapter, name=Remarque]{remarque}
  \ifdefined\FAUX \newcounteralias{corollaire}{section}\newcounteralias{orphelin}{inexistant}\fi
  \ifdefined\CORR \input{correctif-a-remise}\fi
\fi
\crefname{theoreme}{théorème}{théorèmes}
\crefname{proposition}{proposition}{propositions}
\crefname{corollaire}{corollaire}{corollaires}
\crefname{remarque}{remarque}{remarques}
\begin{document}
\chapter{Premier}
\begin{proposition}\label{p1}a\end{proposition}
\begin{theoreme}\label{t1}b\end{theoreme}
\begin{remarque}\label{r1}c\end{remarque}
\begin{corollaire}\label{c1}d\end{corollaire}
\chapter{Second}
\begin{corollaire}\label{c2}e\end{corollaire}
\begin{theoreme}\label{t2}f\end{theoreme}
RENVOIS \ref{p1} \ref{t1} \ref{r1} \ref{c1} \ref{c2} \ref{t2} ; \cref{p1} ; \cref{t1} ; \cref{r1} ; \cref{c1} ; \cref{c2} ; \cref{t2} FIN
\end{document}
"                    amsthm seul (témoin)
%   lualatex -jobname=faux     "\def\CORR{1}\def\FAUX{1}% Juge de la troisième passe : document minimal écrit de zéro pour la cause A (thmtools sous LaTeX 2026-06-01).
% Variantes, par une macro posée sur la ligne de commande (deux passes chacune, depuis memoire/essai/juge3/minimal/) :
%   lualatex -jobname=nu       "\input{minimal-juge3}"                                   sans correctif
%   lualatex -jobname=corr     "\def\CORR{1}\input{minimal-juge3}"                       correctif A du préambule
%   lualatex -jobname=crochet  "\def\SANSCROCHET{1}\input{minimal-juge3}"                sans correctif, crochet retiré
%   lualatex -jobname=amsseul  "\def\AMSSEUL{1}\input{minimal-juge3}"                    amsthm seul (témoin)
%   lualatex -jobname=faux     "\def\CORR{1}\def\FAUX{1}\input{minimal-juge3}"           deux vrais défauts : doivent tomber
%   lualatex -jobname=trace    "\def\TRACE{1}\input{minimal-juge3}"                      trace de l'appel
% Les lignes du correctif sont tirées du préambule commité par extraire_correctifs.py, pas retapées.
\documentclass{scrbook}
\ifdefined\SANSCROCHET \RemoveFromHook{package/amsthm/after}[firstaid/aliascounter]\fi
\usepackage{amsthm}
\usepackage{hyperref}
\usepackage[french, nameinlink, noabbrev]{cleveref}
\ifdefined\AMSSEUL
  \newtheorem{theoreme}{Théorème}[chapter]
  \newtheorem{proposition}[theoreme]{Proposition}
  \newtheorem{corollaire}[theoreme]{Corollaire}
  \newtheorem{remarque}{Remarque}[chapter]
\else
  \usepackage{thmtools}
  \ifdefined\CORR \input{correctif-a-pose}\fi
  \declaretheorem[numberwithin=chapter, name=Théorème]{theoreme}
  \ifdefined\TRACE \tracingmacros=2 \tracingonline=0 \fi
  \declaretheorem[sibling=theoreme, name=Proposition]{proposition}
  \ifdefined\TRACE \tracingmacros=0 \fi
  \declaretheorem[sibling=theoreme, name=Corollaire]{corollaire}
  \declaretheorem[numberwithin=chapter, name=Remarque]{remarque}
  \ifdefined\FAUX \newcounteralias{corollaire}{section}\newcounteralias{orphelin}{inexistant}\fi
  \ifdefined\CORR \input{correctif-a-remise}\fi
\fi
\crefname{theoreme}{théorème}{théorèmes}
\crefname{proposition}{proposition}{propositions}
\crefname{corollaire}{corollaire}{corollaires}
\crefname{remarque}{remarque}{remarques}
\begin{document}
\chapter{Premier}
\begin{proposition}\label{p1}a\end{proposition}
\begin{theoreme}\label{t1}b\end{theoreme}
\begin{remarque}\label{r1}c\end{remarque}
\begin{corollaire}\label{c1}d\end{corollaire}
\chapter{Second}
\begin{corollaire}\label{c2}e\end{corollaire}
\begin{theoreme}\label{t2}f\end{theoreme}
RENVOIS \ref{p1} \ref{t1} \ref{r1} \ref{c1} \ref{c2} \ref{t2} ; \cref{p1} ; \cref{t1} ; \cref{r1} ; \cref{c1} ; \cref{c2} ; \cref{t2} FIN
\end{document}
"           deux vrais défauts : doivent tomber
%   lualatex -jobname=trace    "\def\TRACE{1}% Juge de la troisième passe : document minimal écrit de zéro pour la cause A (thmtools sous LaTeX 2026-06-01).
% Variantes, par une macro posée sur la ligne de commande (deux passes chacune, depuis memoire/essai/juge3/minimal/) :
%   lualatex -jobname=nu       "\input{minimal-juge3}"                                   sans correctif
%   lualatex -jobname=corr     "\def\CORR{1}\input{minimal-juge3}"                       correctif A du préambule
%   lualatex -jobname=crochet  "\def\SANSCROCHET{1}\input{minimal-juge3}"                sans correctif, crochet retiré
%   lualatex -jobname=amsseul  "\def\AMSSEUL{1}\input{minimal-juge3}"                    amsthm seul (témoin)
%   lualatex -jobname=faux     "\def\CORR{1}\def\FAUX{1}\input{minimal-juge3}"           deux vrais défauts : doivent tomber
%   lualatex -jobname=trace    "\def\TRACE{1}\input{minimal-juge3}"                      trace de l'appel
% Les lignes du correctif sont tirées du préambule commité par extraire_correctifs.py, pas retapées.
\documentclass{scrbook}
\ifdefined\SANSCROCHET \RemoveFromHook{package/amsthm/after}[firstaid/aliascounter]\fi
\usepackage{amsthm}
\usepackage{hyperref}
\usepackage[french, nameinlink, noabbrev]{cleveref}
\ifdefined\AMSSEUL
  \newtheorem{theoreme}{Théorème}[chapter]
  \newtheorem{proposition}[theoreme]{Proposition}
  \newtheorem{corollaire}[theoreme]{Corollaire}
  \newtheorem{remarque}{Remarque}[chapter]
\else
  \usepackage{thmtools}
  \ifdefined\CORR \input{correctif-a-pose}\fi
  \declaretheorem[numberwithin=chapter, name=Théorème]{theoreme}
  \ifdefined\TRACE \tracingmacros=2 \tracingonline=0 \fi
  \declaretheorem[sibling=theoreme, name=Proposition]{proposition}
  \ifdefined\TRACE \tracingmacros=0 \fi
  \declaretheorem[sibling=theoreme, name=Corollaire]{corollaire}
  \declaretheorem[numberwithin=chapter, name=Remarque]{remarque}
  \ifdefined\FAUX \newcounteralias{corollaire}{section}\newcounteralias{orphelin}{inexistant}\fi
  \ifdefined\CORR \input{correctif-a-remise}\fi
\fi
\crefname{theoreme}{théorème}{théorèmes}
\crefname{proposition}{proposition}{propositions}
\crefname{corollaire}{corollaire}{corollaires}
\crefname{remarque}{remarque}{remarques}
\begin{document}
\chapter{Premier}
\begin{proposition}\label{p1}a\end{proposition}
\begin{theoreme}\label{t1}b\end{theoreme}
\begin{remarque}\label{r1}c\end{remarque}
\begin{corollaire}\label{c1}d\end{corollaire}
\chapter{Second}
\begin{corollaire}\label{c2}e\end{corollaire}
\begin{theoreme}\label{t2}f\end{theoreme}
RENVOIS \ref{p1} \ref{t1} \ref{r1} \ref{c1} \ref{c2} \ref{t2} ; \cref{p1} ; \cref{t1} ; \cref{r1} ; \cref{c1} ; \cref{c2} ; \cref{t2} FIN
\end{document}
"                      trace de l'appel
% Les lignes du correctif sont tirées du préambule commité par extraire_correctifs.py, pas retapées.
\documentclass{scrbook}
\ifdefined\SANSCROCHET \RemoveFromHook{package/amsthm/after}[firstaid/aliascounter]\fi
\usepackage{amsthm}
\usepackage{hyperref}
\usepackage[french, nameinlink, noabbrev]{cleveref}
\ifdefined\AMSSEUL
  \newtheorem{theoreme}{Théorème}[chapter]
  \newtheorem{proposition}[theoreme]{Proposition}
  \newtheorem{corollaire}[theoreme]{Corollaire}
  \newtheorem{remarque}{Remarque}[chapter]
\else
  \usepackage{thmtools}
  \ifdefined\CORR \input{correctif-a-pose}\fi
  \declaretheorem[numberwithin=chapter, name=Théorème]{theoreme}
  \ifdefined\TRACE \tracingmacros=2 \tracingonline=0 \fi
  \declaretheorem[sibling=theoreme, name=Proposition]{proposition}
  \ifdefined\TRACE \tracingmacros=0 \fi
  \declaretheorem[sibling=theoreme, name=Corollaire]{corollaire}
  \declaretheorem[numberwithin=chapter, name=Remarque]{remarque}
  \ifdefined\FAUX \newcounteralias{corollaire}{section}\newcounteralias{orphelin}{inexistant}\fi
  \ifdefined\CORR \input{correctif-a-remise}\fi
\fi
\crefname{theoreme}{théorème}{théorèmes}
\crefname{proposition}{proposition}{propositions}
\crefname{corollaire}{corollaire}{corollaires}
\crefname{remarque}{remarque}{remarques}
\begin{document}
\chapter{Premier}
\begin{proposition}\label{p1}a\end{proposition}
\begin{theoreme}\label{t1}b\end{theoreme}
\begin{remarque}\label{r1}c\end{remarque}
\begin{corollaire}\label{c1}d\end{corollaire}
\chapter{Second}
\begin{corollaire}\label{c2}e\end{corollaire}
\begin{theoreme}\label{t2}f\end{theoreme}
RENVOIS \ref{p1} \ref{t1} \ref{r1} \ref{c1} \ref{c2} \ref{t2} ; \cref{p1} ; \cref{t1} ; \cref{r1} ; \cref{c1} ; \cref{c2} ; \cref{t2} FIN
\end{document}
"                                   sans correctif
%   lualatex -jobname=corr     "\def\CORR{1}% Juge de la troisième passe : document minimal écrit de zéro pour la cause A (thmtools sous LaTeX 2026-06-01).
% Variantes, par une macro posée sur la ligne de commande (deux passes chacune, depuis memoire/essai/juge3/minimal/) :
%   lualatex -jobname=nu       "% Juge de la troisième passe : document minimal écrit de zéro pour la cause A (thmtools sous LaTeX 2026-06-01).
% Variantes, par une macro posée sur la ligne de commande (deux passes chacune, depuis memoire/essai/juge3/minimal/) :
%   lualatex -jobname=nu       "\input{minimal-juge3}"                                   sans correctif
%   lualatex -jobname=corr     "\def\CORR{1}\input{minimal-juge3}"                       correctif A du préambule
%   lualatex -jobname=crochet  "\def\SANSCROCHET{1}\input{minimal-juge3}"                sans correctif, crochet retiré
%   lualatex -jobname=amsseul  "\def\AMSSEUL{1}\input{minimal-juge3}"                    amsthm seul (témoin)
%   lualatex -jobname=faux     "\def\CORR{1}\def\FAUX{1}\input{minimal-juge3}"           deux vrais défauts : doivent tomber
%   lualatex -jobname=trace    "\def\TRACE{1}\input{minimal-juge3}"                      trace de l'appel
% Les lignes du correctif sont tirées du préambule commité par extraire_correctifs.py, pas retapées.
\documentclass{scrbook}
\ifdefined\SANSCROCHET \RemoveFromHook{package/amsthm/after}[firstaid/aliascounter]\fi
\usepackage{amsthm}
\usepackage{hyperref}
\usepackage[french, nameinlink, noabbrev]{cleveref}
\ifdefined\AMSSEUL
  \newtheorem{theoreme}{Théorème}[chapter]
  \newtheorem{proposition}[theoreme]{Proposition}
  \newtheorem{corollaire}[theoreme]{Corollaire}
  \newtheorem{remarque}{Remarque}[chapter]
\else
  \usepackage{thmtools}
  \ifdefined\CORR \input{correctif-a-pose}\fi
  \declaretheorem[numberwithin=chapter, name=Théorème]{theoreme}
  \ifdefined\TRACE \tracingmacros=2 \tracingonline=0 \fi
  \declaretheorem[sibling=theoreme, name=Proposition]{proposition}
  \ifdefined\TRACE \tracingmacros=0 \fi
  \declaretheorem[sibling=theoreme, name=Corollaire]{corollaire}
  \declaretheorem[numberwithin=chapter, name=Remarque]{remarque}
  \ifdefined\FAUX \newcounteralias{corollaire}{section}\newcounteralias{orphelin}{inexistant}\fi
  \ifdefined\CORR \input{correctif-a-remise}\fi
\fi
\crefname{theoreme}{théorème}{théorèmes}
\crefname{proposition}{proposition}{propositions}
\crefname{corollaire}{corollaire}{corollaires}
\crefname{remarque}{remarque}{remarques}
\begin{document}
\chapter{Premier}
\begin{proposition}\label{p1}a\end{proposition}
\begin{theoreme}\label{t1}b\end{theoreme}
\begin{remarque}\label{r1}c\end{remarque}
\begin{corollaire}\label{c1}d\end{corollaire}
\chapter{Second}
\begin{corollaire}\label{c2}e\end{corollaire}
\begin{theoreme}\label{t2}f\end{theoreme}
RENVOIS \ref{p1} \ref{t1} \ref{r1} \ref{c1} \ref{c2} \ref{t2} ; \cref{p1} ; \cref{t1} ; \cref{r1} ; \cref{c1} ; \cref{c2} ; \cref{t2} FIN
\end{document}
"                                   sans correctif
%   lualatex -jobname=corr     "\def\CORR{1}% Juge de la troisième passe : document minimal écrit de zéro pour la cause A (thmtools sous LaTeX 2026-06-01).
% Variantes, par une macro posée sur la ligne de commande (deux passes chacune, depuis memoire/essai/juge3/minimal/) :
%   lualatex -jobname=nu       "\input{minimal-juge3}"                                   sans correctif
%   lualatex -jobname=corr     "\def\CORR{1}\input{minimal-juge3}"                       correctif A du préambule
%   lualatex -jobname=crochet  "\def\SANSCROCHET{1}\input{minimal-juge3}"                sans correctif, crochet retiré
%   lualatex -jobname=amsseul  "\def\AMSSEUL{1}\input{minimal-juge3}"                    amsthm seul (témoin)
%   lualatex -jobname=faux     "\def\CORR{1}\def\FAUX{1}\input{minimal-juge3}"           deux vrais défauts : doivent tomber
%   lualatex -jobname=trace    "\def\TRACE{1}\input{minimal-juge3}"                      trace de l'appel
% Les lignes du correctif sont tirées du préambule commité par extraire_correctifs.py, pas retapées.
\documentclass{scrbook}
\ifdefined\SANSCROCHET \RemoveFromHook{package/amsthm/after}[firstaid/aliascounter]\fi
\usepackage{amsthm}
\usepackage{hyperref}
\usepackage[french, nameinlink, noabbrev]{cleveref}
\ifdefined\AMSSEUL
  \newtheorem{theoreme}{Théorème}[chapter]
  \newtheorem{proposition}[theoreme]{Proposition}
  \newtheorem{corollaire}[theoreme]{Corollaire}
  \newtheorem{remarque}{Remarque}[chapter]
\else
  \usepackage{thmtools}
  \ifdefined\CORR \input{correctif-a-pose}\fi
  \declaretheorem[numberwithin=chapter, name=Théorème]{theoreme}
  \ifdefined\TRACE \tracingmacros=2 \tracingonline=0 \fi
  \declaretheorem[sibling=theoreme, name=Proposition]{proposition}
  \ifdefined\TRACE \tracingmacros=0 \fi
  \declaretheorem[sibling=theoreme, name=Corollaire]{corollaire}
  \declaretheorem[numberwithin=chapter, name=Remarque]{remarque}
  \ifdefined\FAUX \newcounteralias{corollaire}{section}\newcounteralias{orphelin}{inexistant}\fi
  \ifdefined\CORR \input{correctif-a-remise}\fi
\fi
\crefname{theoreme}{théorème}{théorèmes}
\crefname{proposition}{proposition}{propositions}
\crefname{corollaire}{corollaire}{corollaires}
\crefname{remarque}{remarque}{remarques}
\begin{document}
\chapter{Premier}
\begin{proposition}\label{p1}a\end{proposition}
\begin{theoreme}\label{t1}b\end{theoreme}
\begin{remarque}\label{r1}c\end{remarque}
\begin{corollaire}\label{c1}d\end{corollaire}
\chapter{Second}
\begin{corollaire}\label{c2}e\end{corollaire}
\begin{theoreme}\label{t2}f\end{theoreme}
RENVOIS \ref{p1} \ref{t1} \ref{r1} \ref{c1} \ref{c2} \ref{t2} ; \cref{p1} ; \cref{t1} ; \cref{r1} ; \cref{c1} ; \cref{c2} ; \cref{t2} FIN
\end{document}
"                       correctif A du préambule
%   lualatex -jobname=crochet  "\def\SANSCROCHET{1}% Juge de la troisième passe : document minimal écrit de zéro pour la cause A (thmtools sous LaTeX 2026-06-01).
% Variantes, par une macro posée sur la ligne de commande (deux passes chacune, depuis memoire/essai/juge3/minimal/) :
%   lualatex -jobname=nu       "\input{minimal-juge3}"                                   sans correctif
%   lualatex -jobname=corr     "\def\CORR{1}\input{minimal-juge3}"                       correctif A du préambule
%   lualatex -jobname=crochet  "\def\SANSCROCHET{1}\input{minimal-juge3}"                sans correctif, crochet retiré
%   lualatex -jobname=amsseul  "\def\AMSSEUL{1}\input{minimal-juge3}"                    amsthm seul (témoin)
%   lualatex -jobname=faux     "\def\CORR{1}\def\FAUX{1}\input{minimal-juge3}"           deux vrais défauts : doivent tomber
%   lualatex -jobname=trace    "\def\TRACE{1}\input{minimal-juge3}"                      trace de l'appel
% Les lignes du correctif sont tirées du préambule commité par extraire_correctifs.py, pas retapées.
\documentclass{scrbook}
\ifdefined\SANSCROCHET \RemoveFromHook{package/amsthm/after}[firstaid/aliascounter]\fi
\usepackage{amsthm}
\usepackage{hyperref}
\usepackage[french, nameinlink, noabbrev]{cleveref}
\ifdefined\AMSSEUL
  \newtheorem{theoreme}{Théorème}[chapter]
  \newtheorem{proposition}[theoreme]{Proposition}
  \newtheorem{corollaire}[theoreme]{Corollaire}
  \newtheorem{remarque}{Remarque}[chapter]
\else
  \usepackage{thmtools}
  \ifdefined\CORR \input{correctif-a-pose}\fi
  \declaretheorem[numberwithin=chapter, name=Théorème]{theoreme}
  \ifdefined\TRACE \tracingmacros=2 \tracingonline=0 \fi
  \declaretheorem[sibling=theoreme, name=Proposition]{proposition}
  \ifdefined\TRACE \tracingmacros=0 \fi
  \declaretheorem[sibling=theoreme, name=Corollaire]{corollaire}
  \declaretheorem[numberwithin=chapter, name=Remarque]{remarque}
  \ifdefined\FAUX \newcounteralias{corollaire}{section}\newcounteralias{orphelin}{inexistant}\fi
  \ifdefined\CORR \input{correctif-a-remise}\fi
\fi
\crefname{theoreme}{théorème}{théorèmes}
\crefname{proposition}{proposition}{propositions}
\crefname{corollaire}{corollaire}{corollaires}
\crefname{remarque}{remarque}{remarques}
\begin{document}
\chapter{Premier}
\begin{proposition}\label{p1}a\end{proposition}
\begin{theoreme}\label{t1}b\end{theoreme}
\begin{remarque}\label{r1}c\end{remarque}
\begin{corollaire}\label{c1}d\end{corollaire}
\chapter{Second}
\begin{corollaire}\label{c2}e\end{corollaire}
\begin{theoreme}\label{t2}f\end{theoreme}
RENVOIS \ref{p1} \ref{t1} \ref{r1} \ref{c1} \ref{c2} \ref{t2} ; \cref{p1} ; \cref{t1} ; \cref{r1} ; \cref{c1} ; \cref{c2} ; \cref{t2} FIN
\end{document}
"                sans correctif, crochet retiré
%   lualatex -jobname=amsseul  "\def\AMSSEUL{1}% Juge de la troisième passe : document minimal écrit de zéro pour la cause A (thmtools sous LaTeX 2026-06-01).
% Variantes, par une macro posée sur la ligne de commande (deux passes chacune, depuis memoire/essai/juge3/minimal/) :
%   lualatex -jobname=nu       "\input{minimal-juge3}"                                   sans correctif
%   lualatex -jobname=corr     "\def\CORR{1}\input{minimal-juge3}"                       correctif A du préambule
%   lualatex -jobname=crochet  "\def\SANSCROCHET{1}\input{minimal-juge3}"                sans correctif, crochet retiré
%   lualatex -jobname=amsseul  "\def\AMSSEUL{1}\input{minimal-juge3}"                    amsthm seul (témoin)
%   lualatex -jobname=faux     "\def\CORR{1}\def\FAUX{1}\input{minimal-juge3}"           deux vrais défauts : doivent tomber
%   lualatex -jobname=trace    "\def\TRACE{1}\input{minimal-juge3}"                      trace de l'appel
% Les lignes du correctif sont tirées du préambule commité par extraire_correctifs.py, pas retapées.
\documentclass{scrbook}
\ifdefined\SANSCROCHET \RemoveFromHook{package/amsthm/after}[firstaid/aliascounter]\fi
\usepackage{amsthm}
\usepackage{hyperref}
\usepackage[french, nameinlink, noabbrev]{cleveref}
\ifdefined\AMSSEUL
  \newtheorem{theoreme}{Théorème}[chapter]
  \newtheorem{proposition}[theoreme]{Proposition}
  \newtheorem{corollaire}[theoreme]{Corollaire}
  \newtheorem{remarque}{Remarque}[chapter]
\else
  \usepackage{thmtools}
  \ifdefined\CORR \input{correctif-a-pose}\fi
  \declaretheorem[numberwithin=chapter, name=Théorème]{theoreme}
  \ifdefined\TRACE \tracingmacros=2 \tracingonline=0 \fi
  \declaretheorem[sibling=theoreme, name=Proposition]{proposition}
  \ifdefined\TRACE \tracingmacros=0 \fi
  \declaretheorem[sibling=theoreme, name=Corollaire]{corollaire}
  \declaretheorem[numberwithin=chapter, name=Remarque]{remarque}
  \ifdefined\FAUX \newcounteralias{corollaire}{section}\newcounteralias{orphelin}{inexistant}\fi
  \ifdefined\CORR \input{correctif-a-remise}\fi
\fi
\crefname{theoreme}{théorème}{théorèmes}
\crefname{proposition}{proposition}{propositions}
\crefname{corollaire}{corollaire}{corollaires}
\crefname{remarque}{remarque}{remarques}
\begin{document}
\chapter{Premier}
\begin{proposition}\label{p1}a\end{proposition}
\begin{theoreme}\label{t1}b\end{theoreme}
\begin{remarque}\label{r1}c\end{remarque}
\begin{corollaire}\label{c1}d\end{corollaire}
\chapter{Second}
\begin{corollaire}\label{c2}e\end{corollaire}
\begin{theoreme}\label{t2}f\end{theoreme}
RENVOIS \ref{p1} \ref{t1} \ref{r1} \ref{c1} \ref{c2} \ref{t2} ; \cref{p1} ; \cref{t1} ; \cref{r1} ; \cref{c1} ; \cref{c2} ; \cref{t2} FIN
\end{document}
"                    amsthm seul (témoin)
%   lualatex -jobname=faux     "\def\CORR{1}\def\FAUX{1}% Juge de la troisième passe : document minimal écrit de zéro pour la cause A (thmtools sous LaTeX 2026-06-01).
% Variantes, par une macro posée sur la ligne de commande (deux passes chacune, depuis memoire/essai/juge3/minimal/) :
%   lualatex -jobname=nu       "\input{minimal-juge3}"                                   sans correctif
%   lualatex -jobname=corr     "\def\CORR{1}\input{minimal-juge3}"                       correctif A du préambule
%   lualatex -jobname=crochet  "\def\SANSCROCHET{1}\input{minimal-juge3}"                sans correctif, crochet retiré
%   lualatex -jobname=amsseul  "\def\AMSSEUL{1}\input{minimal-juge3}"                    amsthm seul (témoin)
%   lualatex -jobname=faux     "\def\CORR{1}\def\FAUX{1}\input{minimal-juge3}"           deux vrais défauts : doivent tomber
%   lualatex -jobname=trace    "\def\TRACE{1}\input{minimal-juge3}"                      trace de l'appel
% Les lignes du correctif sont tirées du préambule commité par extraire_correctifs.py, pas retapées.
\documentclass{scrbook}
\ifdefined\SANSCROCHET \RemoveFromHook{package/amsthm/after}[firstaid/aliascounter]\fi
\usepackage{amsthm}
\usepackage{hyperref}
\usepackage[french, nameinlink, noabbrev]{cleveref}
\ifdefined\AMSSEUL
  \newtheorem{theoreme}{Théorème}[chapter]
  \newtheorem{proposition}[theoreme]{Proposition}
  \newtheorem{corollaire}[theoreme]{Corollaire}
  \newtheorem{remarque}{Remarque}[chapter]
\else
  \usepackage{thmtools}
  \ifdefined\CORR \input{correctif-a-pose}\fi
  \declaretheorem[numberwithin=chapter, name=Théorème]{theoreme}
  \ifdefined\TRACE \tracingmacros=2 \tracingonline=0 \fi
  \declaretheorem[sibling=theoreme, name=Proposition]{proposition}
  \ifdefined\TRACE \tracingmacros=0 \fi
  \declaretheorem[sibling=theoreme, name=Corollaire]{corollaire}
  \declaretheorem[numberwithin=chapter, name=Remarque]{remarque}
  \ifdefined\FAUX \newcounteralias{corollaire}{section}\newcounteralias{orphelin}{inexistant}\fi
  \ifdefined\CORR \input{correctif-a-remise}\fi
\fi
\crefname{theoreme}{théorème}{théorèmes}
\crefname{proposition}{proposition}{propositions}
\crefname{corollaire}{corollaire}{corollaires}
\crefname{remarque}{remarque}{remarques}
\begin{document}
\chapter{Premier}
\begin{proposition}\label{p1}a\end{proposition}
\begin{theoreme}\label{t1}b\end{theoreme}
\begin{remarque}\label{r1}c\end{remarque}
\begin{corollaire}\label{c1}d\end{corollaire}
\chapter{Second}
\begin{corollaire}\label{c2}e\end{corollaire}
\begin{theoreme}\label{t2}f\end{theoreme}
RENVOIS \ref{p1} \ref{t1} \ref{r1} \ref{c1} \ref{c2} \ref{t2} ; \cref{p1} ; \cref{t1} ; \cref{r1} ; \cref{c1} ; \cref{c2} ; \cref{t2} FIN
\end{document}
"           deux vrais défauts : doivent tomber
%   lualatex -jobname=trace    "\def\TRACE{1}% Juge de la troisième passe : document minimal écrit de zéro pour la cause A (thmtools sous LaTeX 2026-06-01).
% Variantes, par une macro posée sur la ligne de commande (deux passes chacune, depuis memoire/essai/juge3/minimal/) :
%   lualatex -jobname=nu       "\input{minimal-juge3}"                                   sans correctif
%   lualatex -jobname=corr     "\def\CORR{1}\input{minimal-juge3}"                       correctif A du préambule
%   lualatex -jobname=crochet  "\def\SANSCROCHET{1}\input{minimal-juge3}"                sans correctif, crochet retiré
%   lualatex -jobname=amsseul  "\def\AMSSEUL{1}\input{minimal-juge3}"                    amsthm seul (témoin)
%   lualatex -jobname=faux     "\def\CORR{1}\def\FAUX{1}\input{minimal-juge3}"           deux vrais défauts : doivent tomber
%   lualatex -jobname=trace    "\def\TRACE{1}\input{minimal-juge3}"                      trace de l'appel
% Les lignes du correctif sont tirées du préambule commité par extraire_correctifs.py, pas retapées.
\documentclass{scrbook}
\ifdefined\SANSCROCHET \RemoveFromHook{package/amsthm/after}[firstaid/aliascounter]\fi
\usepackage{amsthm}
\usepackage{hyperref}
\usepackage[french, nameinlink, noabbrev]{cleveref}
\ifdefined\AMSSEUL
  \newtheorem{theoreme}{Théorème}[chapter]
  \newtheorem{proposition}[theoreme]{Proposition}
  \newtheorem{corollaire}[theoreme]{Corollaire}
  \newtheorem{remarque}{Remarque}[chapter]
\else
  \usepackage{thmtools}
  \ifdefined\CORR \input{correctif-a-pose}\fi
  \declaretheorem[numberwithin=chapter, name=Théorème]{theoreme}
  \ifdefined\TRACE \tracingmacros=2 \tracingonline=0 \fi
  \declaretheorem[sibling=theoreme, name=Proposition]{proposition}
  \ifdefined\TRACE \tracingmacros=0 \fi
  \declaretheorem[sibling=theoreme, name=Corollaire]{corollaire}
  \declaretheorem[numberwithin=chapter, name=Remarque]{remarque}
  \ifdefined\FAUX \newcounteralias{corollaire}{section}\newcounteralias{orphelin}{inexistant}\fi
  \ifdefined\CORR \input{correctif-a-remise}\fi
\fi
\crefname{theoreme}{théorème}{théorèmes}
\crefname{proposition}{proposition}{propositions}
\crefname{corollaire}{corollaire}{corollaires}
\crefname{remarque}{remarque}{remarques}
\begin{document}
\chapter{Premier}
\begin{proposition}\label{p1}a\end{proposition}
\begin{theoreme}\label{t1}b\end{theoreme}
\begin{remarque}\label{r1}c\end{remarque}
\begin{corollaire}\label{c1}d\end{corollaire}
\chapter{Second}
\begin{corollaire}\label{c2}e\end{corollaire}
\begin{theoreme}\label{t2}f\end{theoreme}
RENVOIS \ref{p1} \ref{t1} \ref{r1} \ref{c1} \ref{c2} \ref{t2} ; \cref{p1} ; \cref{t1} ; \cref{r1} ; \cref{c1} ; \cref{c2} ; \cref{t2} FIN
\end{document}
"                      trace de l'appel
% Les lignes du correctif sont tirées du préambule commité par extraire_correctifs.py, pas retapées.
\documentclass{scrbook}
\ifdefined\SANSCROCHET \RemoveFromHook{package/amsthm/after}[firstaid/aliascounter]\fi
\usepackage{amsthm}
\usepackage{hyperref}
\usepackage[french, nameinlink, noabbrev]{cleveref}
\ifdefined\AMSSEUL
  \newtheorem{theoreme}{Théorème}[chapter]
  \newtheorem{proposition}[theoreme]{Proposition}
  \newtheorem{corollaire}[theoreme]{Corollaire}
  \newtheorem{remarque}{Remarque}[chapter]
\else
  \usepackage{thmtools}
  \ifdefined\CORR \input{correctif-a-pose}\fi
  \declaretheorem[numberwithin=chapter, name=Théorème]{theoreme}
  \ifdefined\TRACE \tracingmacros=2 \tracingonline=0 \fi
  \declaretheorem[sibling=theoreme, name=Proposition]{proposition}
  \ifdefined\TRACE \tracingmacros=0 \fi
  \declaretheorem[sibling=theoreme, name=Corollaire]{corollaire}
  \declaretheorem[numberwithin=chapter, name=Remarque]{remarque}
  \ifdefined\FAUX \newcounteralias{corollaire}{section}\newcounteralias{orphelin}{inexistant}\fi
  \ifdefined\CORR \input{correctif-a-remise}\fi
\fi
\crefname{theoreme}{théorème}{théorèmes}
\crefname{proposition}{proposition}{propositions}
\crefname{corollaire}{corollaire}{corollaires}
\crefname{remarque}{remarque}{remarques}
\begin{document}
\chapter{Premier}
\begin{proposition}\label{p1}a\end{proposition}
\begin{theoreme}\label{t1}b\end{theoreme}
\begin{remarque}\label{r1}c\end{remarque}
\begin{corollaire}\label{c1}d\end{corollaire}
\chapter{Second}
\begin{corollaire}\label{c2}e\end{corollaire}
\begin{theoreme}\label{t2}f\end{theoreme}
RENVOIS \ref{p1} \ref{t1} \ref{r1} \ref{c1} \ref{c2} \ref{t2} ; \cref{p1} ; \cref{t1} ; \cref{r1} ; \cref{c1} ; \cref{c2} ; \cref{t2} FIN
\end{document}
"                       correctif A du préambule
%   lualatex -jobname=crochet  "\def\SANSCROCHET{1}% Juge de la troisième passe : document minimal écrit de zéro pour la cause A (thmtools sous LaTeX 2026-06-01).
% Variantes, par une macro posée sur la ligne de commande (deux passes chacune, depuis memoire/essai/juge3/minimal/) :
%   lualatex -jobname=nu       "% Juge de la troisième passe : document minimal écrit de zéro pour la cause A (thmtools sous LaTeX 2026-06-01).
% Variantes, par une macro posée sur la ligne de commande (deux passes chacune, depuis memoire/essai/juge3/minimal/) :
%   lualatex -jobname=nu       "\input{minimal-juge3}"                                   sans correctif
%   lualatex -jobname=corr     "\def\CORR{1}\input{minimal-juge3}"                       correctif A du préambule
%   lualatex -jobname=crochet  "\def\SANSCROCHET{1}\input{minimal-juge3}"                sans correctif, crochet retiré
%   lualatex -jobname=amsseul  "\def\AMSSEUL{1}\input{minimal-juge3}"                    amsthm seul (témoin)
%   lualatex -jobname=faux     "\def\CORR{1}\def\FAUX{1}\input{minimal-juge3}"           deux vrais défauts : doivent tomber
%   lualatex -jobname=trace    "\def\TRACE{1}\input{minimal-juge3}"                      trace de l'appel
% Les lignes du correctif sont tirées du préambule commité par extraire_correctifs.py, pas retapées.
\documentclass{scrbook}
\ifdefined\SANSCROCHET \RemoveFromHook{package/amsthm/after}[firstaid/aliascounter]\fi
\usepackage{amsthm}
\usepackage{hyperref}
\usepackage[french, nameinlink, noabbrev]{cleveref}
\ifdefined\AMSSEUL
  \newtheorem{theoreme}{Théorème}[chapter]
  \newtheorem{proposition}[theoreme]{Proposition}
  \newtheorem{corollaire}[theoreme]{Corollaire}
  \newtheorem{remarque}{Remarque}[chapter]
\else
  \usepackage{thmtools}
  \ifdefined\CORR \input{correctif-a-pose}\fi
  \declaretheorem[numberwithin=chapter, name=Théorème]{theoreme}
  \ifdefined\TRACE \tracingmacros=2 \tracingonline=0 \fi
  \declaretheorem[sibling=theoreme, name=Proposition]{proposition}
  \ifdefined\TRACE \tracingmacros=0 \fi
  \declaretheorem[sibling=theoreme, name=Corollaire]{corollaire}
  \declaretheorem[numberwithin=chapter, name=Remarque]{remarque}
  \ifdefined\FAUX \newcounteralias{corollaire}{section}\newcounteralias{orphelin}{inexistant}\fi
  \ifdefined\CORR \input{correctif-a-remise}\fi
\fi
\crefname{theoreme}{théorème}{théorèmes}
\crefname{proposition}{proposition}{propositions}
\crefname{corollaire}{corollaire}{corollaires}
\crefname{remarque}{remarque}{remarques}
\begin{document}
\chapter{Premier}
\begin{proposition}\label{p1}a\end{proposition}
\begin{theoreme}\label{t1}b\end{theoreme}
\begin{remarque}\label{r1}c\end{remarque}
\begin{corollaire}\label{c1}d\end{corollaire}
\chapter{Second}
\begin{corollaire}\label{c2}e\end{corollaire}
\begin{theoreme}\label{t2}f\end{theoreme}
RENVOIS \ref{p1} \ref{t1} \ref{r1} \ref{c1} \ref{c2} \ref{t2} ; \cref{p1} ; \cref{t1} ; \cref{r1} ; \cref{c1} ; \cref{c2} ; \cref{t2} FIN
\end{document}
"                                   sans correctif
%   lualatex -jobname=corr     "\def\CORR{1}% Juge de la troisième passe : document minimal écrit de zéro pour la cause A (thmtools sous LaTeX 2026-06-01).
% Variantes, par une macro posée sur la ligne de commande (deux passes chacune, depuis memoire/essai/juge3/minimal/) :
%   lualatex -jobname=nu       "\input{minimal-juge3}"                                   sans correctif
%   lualatex -jobname=corr     "\def\CORR{1}\input{minimal-juge3}"                       correctif A du préambule
%   lualatex -jobname=crochet  "\def\SANSCROCHET{1}\input{minimal-juge3}"                sans correctif, crochet retiré
%   lualatex -jobname=amsseul  "\def\AMSSEUL{1}\input{minimal-juge3}"                    amsthm seul (témoin)
%   lualatex -jobname=faux     "\def\CORR{1}\def\FAUX{1}\input{minimal-juge3}"           deux vrais défauts : doivent tomber
%   lualatex -jobname=trace    "\def\TRACE{1}\input{minimal-juge3}"                      trace de l'appel
% Les lignes du correctif sont tirées du préambule commité par extraire_correctifs.py, pas retapées.
\documentclass{scrbook}
\ifdefined\SANSCROCHET \RemoveFromHook{package/amsthm/after}[firstaid/aliascounter]\fi
\usepackage{amsthm}
\usepackage{hyperref}
\usepackage[french, nameinlink, noabbrev]{cleveref}
\ifdefined\AMSSEUL
  \newtheorem{theoreme}{Théorème}[chapter]
  \newtheorem{proposition}[theoreme]{Proposition}
  \newtheorem{corollaire}[theoreme]{Corollaire}
  \newtheorem{remarque}{Remarque}[chapter]
\else
  \usepackage{thmtools}
  \ifdefined\CORR \input{correctif-a-pose}\fi
  \declaretheorem[numberwithin=chapter, name=Théorème]{theoreme}
  \ifdefined\TRACE \tracingmacros=2 \tracingonline=0 \fi
  \declaretheorem[sibling=theoreme, name=Proposition]{proposition}
  \ifdefined\TRACE \tracingmacros=0 \fi
  \declaretheorem[sibling=theoreme, name=Corollaire]{corollaire}
  \declaretheorem[numberwithin=chapter, name=Remarque]{remarque}
  \ifdefined\FAUX \newcounteralias{corollaire}{section}\newcounteralias{orphelin}{inexistant}\fi
  \ifdefined\CORR \input{correctif-a-remise}\fi
\fi
\crefname{theoreme}{théorème}{théorèmes}
\crefname{proposition}{proposition}{propositions}
\crefname{corollaire}{corollaire}{corollaires}
\crefname{remarque}{remarque}{remarques}
\begin{document}
\chapter{Premier}
\begin{proposition}\label{p1}a\end{proposition}
\begin{theoreme}\label{t1}b\end{theoreme}
\begin{remarque}\label{r1}c\end{remarque}
\begin{corollaire}\label{c1}d\end{corollaire}
\chapter{Second}
\begin{corollaire}\label{c2}e\end{corollaire}
\begin{theoreme}\label{t2}f\end{theoreme}
RENVOIS \ref{p1} \ref{t1} \ref{r1} \ref{c1} \ref{c2} \ref{t2} ; \cref{p1} ; \cref{t1} ; \cref{r1} ; \cref{c1} ; \cref{c2} ; \cref{t2} FIN
\end{document}
"                       correctif A du préambule
%   lualatex -jobname=crochet  "\def\SANSCROCHET{1}% Juge de la troisième passe : document minimal écrit de zéro pour la cause A (thmtools sous LaTeX 2026-06-01).
% Variantes, par une macro posée sur la ligne de commande (deux passes chacune, depuis memoire/essai/juge3/minimal/) :
%   lualatex -jobname=nu       "\input{minimal-juge3}"                                   sans correctif
%   lualatex -jobname=corr     "\def\CORR{1}\input{minimal-juge3}"                       correctif A du préambule
%   lualatex -jobname=crochet  "\def\SANSCROCHET{1}\input{minimal-juge3}"                sans correctif, crochet retiré
%   lualatex -jobname=amsseul  "\def\AMSSEUL{1}\input{minimal-juge3}"                    amsthm seul (témoin)
%   lualatex -jobname=faux     "\def\CORR{1}\def\FAUX{1}\input{minimal-juge3}"           deux vrais défauts : doivent tomber
%   lualatex -jobname=trace    "\def\TRACE{1}\input{minimal-juge3}"                      trace de l'appel
% Les lignes du correctif sont tirées du préambule commité par extraire_correctifs.py, pas retapées.
\documentclass{scrbook}
\ifdefined\SANSCROCHET \RemoveFromHook{package/amsthm/after}[firstaid/aliascounter]\fi
\usepackage{amsthm}
\usepackage{hyperref}
\usepackage[french, nameinlink, noabbrev]{cleveref}
\ifdefined\AMSSEUL
  \newtheorem{theoreme}{Théorème}[chapter]
  \newtheorem{proposition}[theoreme]{Proposition}
  \newtheorem{corollaire}[theoreme]{Corollaire}
  \newtheorem{remarque}{Remarque}[chapter]
\else
  \usepackage{thmtools}
  \ifdefined\CORR \input{correctif-a-pose}\fi
  \declaretheorem[numberwithin=chapter, name=Théorème]{theoreme}
  \ifdefined\TRACE \tracingmacros=2 \tracingonline=0 \fi
  \declaretheorem[sibling=theoreme, name=Proposition]{proposition}
  \ifdefined\TRACE \tracingmacros=0 \fi
  \declaretheorem[sibling=theoreme, name=Corollaire]{corollaire}
  \declaretheorem[numberwithin=chapter, name=Remarque]{remarque}
  \ifdefined\FAUX \newcounteralias{corollaire}{section}\newcounteralias{orphelin}{inexistant}\fi
  \ifdefined\CORR \input{correctif-a-remise}\fi
\fi
\crefname{theoreme}{théorème}{théorèmes}
\crefname{proposition}{proposition}{propositions}
\crefname{corollaire}{corollaire}{corollaires}
\crefname{remarque}{remarque}{remarques}
\begin{document}
\chapter{Premier}
\begin{proposition}\label{p1}a\end{proposition}
\begin{theoreme}\label{t1}b\end{theoreme}
\begin{remarque}\label{r1}c\end{remarque}
\begin{corollaire}\label{c1}d\end{corollaire}
\chapter{Second}
\begin{corollaire}\label{c2}e\end{corollaire}
\begin{theoreme}\label{t2}f\end{theoreme}
RENVOIS \ref{p1} \ref{t1} \ref{r1} \ref{c1} \ref{c2} \ref{t2} ; \cref{p1} ; \cref{t1} ; \cref{r1} ; \cref{c1} ; \cref{c2} ; \cref{t2} FIN
\end{document}
"                sans correctif, crochet retiré
%   lualatex -jobname=amsseul  "\def\AMSSEUL{1}% Juge de la troisième passe : document minimal écrit de zéro pour la cause A (thmtools sous LaTeX 2026-06-01).
% Variantes, par une macro posée sur la ligne de commande (deux passes chacune, depuis memoire/essai/juge3/minimal/) :
%   lualatex -jobname=nu       "\input{minimal-juge3}"                                   sans correctif
%   lualatex -jobname=corr     "\def\CORR{1}\input{minimal-juge3}"                       correctif A du préambule
%   lualatex -jobname=crochet  "\def\SANSCROCHET{1}\input{minimal-juge3}"                sans correctif, crochet retiré
%   lualatex -jobname=amsseul  "\def\AMSSEUL{1}\input{minimal-juge3}"                    amsthm seul (témoin)
%   lualatex -jobname=faux     "\def\CORR{1}\def\FAUX{1}\input{minimal-juge3}"           deux vrais défauts : doivent tomber
%   lualatex -jobname=trace    "\def\TRACE{1}\input{minimal-juge3}"                      trace de l'appel
% Les lignes du correctif sont tirées du préambule commité par extraire_correctifs.py, pas retapées.
\documentclass{scrbook}
\ifdefined\SANSCROCHET \RemoveFromHook{package/amsthm/after}[firstaid/aliascounter]\fi
\usepackage{amsthm}
\usepackage{hyperref}
\usepackage[french, nameinlink, noabbrev]{cleveref}
\ifdefined\AMSSEUL
  \newtheorem{theoreme}{Théorème}[chapter]
  \newtheorem{proposition}[theoreme]{Proposition}
  \newtheorem{corollaire}[theoreme]{Corollaire}
  \newtheorem{remarque}{Remarque}[chapter]
\else
  \usepackage{thmtools}
  \ifdefined\CORR \input{correctif-a-pose}\fi
  \declaretheorem[numberwithin=chapter, name=Théorème]{theoreme}
  \ifdefined\TRACE \tracingmacros=2 \tracingonline=0 \fi
  \declaretheorem[sibling=theoreme, name=Proposition]{proposition}
  \ifdefined\TRACE \tracingmacros=0 \fi
  \declaretheorem[sibling=theoreme, name=Corollaire]{corollaire}
  \declaretheorem[numberwithin=chapter, name=Remarque]{remarque}
  \ifdefined\FAUX \newcounteralias{corollaire}{section}\newcounteralias{orphelin}{inexistant}\fi
  \ifdefined\CORR \input{correctif-a-remise}\fi
\fi
\crefname{theoreme}{théorème}{théorèmes}
\crefname{proposition}{proposition}{propositions}
\crefname{corollaire}{corollaire}{corollaires}
\crefname{remarque}{remarque}{remarques}
\begin{document}
\chapter{Premier}
\begin{proposition}\label{p1}a\end{proposition}
\begin{theoreme}\label{t1}b\end{theoreme}
\begin{remarque}\label{r1}c\end{remarque}
\begin{corollaire}\label{c1}d\end{corollaire}
\chapter{Second}
\begin{corollaire}\label{c2}e\end{corollaire}
\begin{theoreme}\label{t2}f\end{theoreme}
RENVOIS \ref{p1} \ref{t1} \ref{r1} \ref{c1} \ref{c2} \ref{t2} ; \cref{p1} ; \cref{t1} ; \cref{r1} ; \cref{c1} ; \cref{c2} ; \cref{t2} FIN
\end{document}
"                    amsthm seul (témoin)
%   lualatex -jobname=faux     "\def\CORR{1}\def\FAUX{1}% Juge de la troisième passe : document minimal écrit de zéro pour la cause A (thmtools sous LaTeX 2026-06-01).
% Variantes, par une macro posée sur la ligne de commande (deux passes chacune, depuis memoire/essai/juge3/minimal/) :
%   lualatex -jobname=nu       "\input{minimal-juge3}"                                   sans correctif
%   lualatex -jobname=corr     "\def\CORR{1}\input{minimal-juge3}"                       correctif A du préambule
%   lualatex -jobname=crochet  "\def\SANSCROCHET{1}\input{minimal-juge3}"                sans correctif, crochet retiré
%   lualatex -jobname=amsseul  "\def\AMSSEUL{1}\input{minimal-juge3}"                    amsthm seul (témoin)
%   lualatex -jobname=faux     "\def\CORR{1}\def\FAUX{1}\input{minimal-juge3}"           deux vrais défauts : doivent tomber
%   lualatex -jobname=trace    "\def\TRACE{1}\input{minimal-juge3}"                      trace de l'appel
% Les lignes du correctif sont tirées du préambule commité par extraire_correctifs.py, pas retapées.
\documentclass{scrbook}
\ifdefined\SANSCROCHET \RemoveFromHook{package/amsthm/after}[firstaid/aliascounter]\fi
\usepackage{amsthm}
\usepackage{hyperref}
\usepackage[french, nameinlink, noabbrev]{cleveref}
\ifdefined\AMSSEUL
  \newtheorem{theoreme}{Théorème}[chapter]
  \newtheorem{proposition}[theoreme]{Proposition}
  \newtheorem{corollaire}[theoreme]{Corollaire}
  \newtheorem{remarque}{Remarque}[chapter]
\else
  \usepackage{thmtools}
  \ifdefined\CORR \input{correctif-a-pose}\fi
  \declaretheorem[numberwithin=chapter, name=Théorème]{theoreme}
  \ifdefined\TRACE \tracingmacros=2 \tracingonline=0 \fi
  \declaretheorem[sibling=theoreme, name=Proposition]{proposition}
  \ifdefined\TRACE \tracingmacros=0 \fi
  \declaretheorem[sibling=theoreme, name=Corollaire]{corollaire}
  \declaretheorem[numberwithin=chapter, name=Remarque]{remarque}
  \ifdefined\FAUX \newcounteralias{corollaire}{section}\newcounteralias{orphelin}{inexistant}\fi
  \ifdefined\CORR \input{correctif-a-remise}\fi
\fi
\crefname{theoreme}{théorème}{théorèmes}
\crefname{proposition}{proposition}{propositions}
\crefname{corollaire}{corollaire}{corollaires}
\crefname{remarque}{remarque}{remarques}
\begin{document}
\chapter{Premier}
\begin{proposition}\label{p1}a\end{proposition}
\begin{theoreme}\label{t1}b\end{theoreme}
\begin{remarque}\label{r1}c\end{remarque}
\begin{corollaire}\label{c1}d\end{corollaire}
\chapter{Second}
\begin{corollaire}\label{c2}e\end{corollaire}
\begin{theoreme}\label{t2}f\end{theoreme}
RENVOIS \ref{p1} \ref{t1} \ref{r1} \ref{c1} \ref{c2} \ref{t2} ; \cref{p1} ; \cref{t1} ; \cref{r1} ; \cref{c1} ; \cref{c2} ; \cref{t2} FIN
\end{document}
"           deux vrais défauts : doivent tomber
%   lualatex -jobname=trace    "\def\TRACE{1}% Juge de la troisième passe : document minimal écrit de zéro pour la cause A (thmtools sous LaTeX 2026-06-01).
% Variantes, par une macro posée sur la ligne de commande (deux passes chacune, depuis memoire/essai/juge3/minimal/) :
%   lualatex -jobname=nu       "\input{minimal-juge3}"                                   sans correctif
%   lualatex -jobname=corr     "\def\CORR{1}\input{minimal-juge3}"                       correctif A du préambule
%   lualatex -jobname=crochet  "\def\SANSCROCHET{1}\input{minimal-juge3}"                sans correctif, crochet retiré
%   lualatex -jobname=amsseul  "\def\AMSSEUL{1}\input{minimal-juge3}"                    amsthm seul (témoin)
%   lualatex -jobname=faux     "\def\CORR{1}\def\FAUX{1}\input{minimal-juge3}"           deux vrais défauts : doivent tomber
%   lualatex -jobname=trace    "\def\TRACE{1}\input{minimal-juge3}"                      trace de l'appel
% Les lignes du correctif sont tirées du préambule commité par extraire_correctifs.py, pas retapées.
\documentclass{scrbook}
\ifdefined\SANSCROCHET \RemoveFromHook{package/amsthm/after}[firstaid/aliascounter]\fi
\usepackage{amsthm}
\usepackage{hyperref}
\usepackage[french, nameinlink, noabbrev]{cleveref}
\ifdefined\AMSSEUL
  \newtheorem{theoreme}{Théorème}[chapter]
  \newtheorem{proposition}[theoreme]{Proposition}
  \newtheorem{corollaire}[theoreme]{Corollaire}
  \newtheorem{remarque}{Remarque}[chapter]
\else
  \usepackage{thmtools}
  \ifdefined\CORR \input{correctif-a-pose}\fi
  \declaretheorem[numberwithin=chapter, name=Théorème]{theoreme}
  \ifdefined\TRACE \tracingmacros=2 \tracingonline=0 \fi
  \declaretheorem[sibling=theoreme, name=Proposition]{proposition}
  \ifdefined\TRACE \tracingmacros=0 \fi
  \declaretheorem[sibling=theoreme, name=Corollaire]{corollaire}
  \declaretheorem[numberwithin=chapter, name=Remarque]{remarque}
  \ifdefined\FAUX \newcounteralias{corollaire}{section}\newcounteralias{orphelin}{inexistant}\fi
  \ifdefined\CORR \input{correctif-a-remise}\fi
\fi
\crefname{theoreme}{théorème}{théorèmes}
\crefname{proposition}{proposition}{propositions}
\crefname{corollaire}{corollaire}{corollaires}
\crefname{remarque}{remarque}{remarques}
\begin{document}
\chapter{Premier}
\begin{proposition}\label{p1}a\end{proposition}
\begin{theoreme}\label{t1}b\end{theoreme}
\begin{remarque}\label{r1}c\end{remarque}
\begin{corollaire}\label{c1}d\end{corollaire}
\chapter{Second}
\begin{corollaire}\label{c2}e\end{corollaire}
\begin{theoreme}\label{t2}f\end{theoreme}
RENVOIS \ref{p1} \ref{t1} \ref{r1} \ref{c1} \ref{c2} \ref{t2} ; \cref{p1} ; \cref{t1} ; \cref{r1} ; \cref{c1} ; \cref{c2} ; \cref{t2} FIN
\end{document}
"                      trace de l'appel
% Les lignes du correctif sont tirées du préambule commité par extraire_correctifs.py, pas retapées.
\documentclass{scrbook}
\ifdefined\SANSCROCHET \RemoveFromHook{package/amsthm/after}[firstaid/aliascounter]\fi
\usepackage{amsthm}
\usepackage{hyperref}
\usepackage[french, nameinlink, noabbrev]{cleveref}
\ifdefined\AMSSEUL
  \newtheorem{theoreme}{Théorème}[chapter]
  \newtheorem{proposition}[theoreme]{Proposition}
  \newtheorem{corollaire}[theoreme]{Corollaire}
  \newtheorem{remarque}{Remarque}[chapter]
\else
  \usepackage{thmtools}
  \ifdefined\CORR \input{correctif-a-pose}\fi
  \declaretheorem[numberwithin=chapter, name=Théorème]{theoreme}
  \ifdefined\TRACE \tracingmacros=2 \tracingonline=0 \fi
  \declaretheorem[sibling=theoreme, name=Proposition]{proposition}
  \ifdefined\TRACE \tracingmacros=0 \fi
  \declaretheorem[sibling=theoreme, name=Corollaire]{corollaire}
  \declaretheorem[numberwithin=chapter, name=Remarque]{remarque}
  \ifdefined\FAUX \newcounteralias{corollaire}{section}\newcounteralias{orphelin}{inexistant}\fi
  \ifdefined\CORR \input{correctif-a-remise}\fi
\fi
\crefname{theoreme}{théorème}{théorèmes}
\crefname{proposition}{proposition}{propositions}
\crefname{corollaire}{corollaire}{corollaires}
\crefname{remarque}{remarque}{remarques}
\begin{document}
\chapter{Premier}
\begin{proposition}\label{p1}a\end{proposition}
\begin{theoreme}\label{t1}b\end{theoreme}
\begin{remarque}\label{r1}c\end{remarque}
\begin{corollaire}\label{c1}d\end{corollaire}
\chapter{Second}
\begin{corollaire}\label{c2}e\end{corollaire}
\begin{theoreme}\label{t2}f\end{theoreme}
RENVOIS \ref{p1} \ref{t1} \ref{r1} \ref{c1} \ref{c2} \ref{t2} ; \cref{p1} ; \cref{t1} ; \cref{r1} ; \cref{c1} ; \cref{c2} ; \cref{t2} FIN
\end{document}
"                sans correctif, crochet retiré
%   lualatex -jobname=amsseul  "\def\AMSSEUL{1}% Juge de la troisième passe : document minimal écrit de zéro pour la cause A (thmtools sous LaTeX 2026-06-01).
% Variantes, par une macro posée sur la ligne de commande (deux passes chacune, depuis memoire/essai/juge3/minimal/) :
%   lualatex -jobname=nu       "% Juge de la troisième passe : document minimal écrit de zéro pour la cause A (thmtools sous LaTeX 2026-06-01).
% Variantes, par une macro posée sur la ligne de commande (deux passes chacune, depuis memoire/essai/juge3/minimal/) :
%   lualatex -jobname=nu       "\input{minimal-juge3}"                                   sans correctif
%   lualatex -jobname=corr     "\def\CORR{1}\input{minimal-juge3}"                       correctif A du préambule
%   lualatex -jobname=crochet  "\def\SANSCROCHET{1}\input{minimal-juge3}"                sans correctif, crochet retiré
%   lualatex -jobname=amsseul  "\def\AMSSEUL{1}\input{minimal-juge3}"                    amsthm seul (témoin)
%   lualatex -jobname=faux     "\def\CORR{1}\def\FAUX{1}\input{minimal-juge3}"           deux vrais défauts : doivent tomber
%   lualatex -jobname=trace    "\def\TRACE{1}\input{minimal-juge3}"                      trace de l'appel
% Les lignes du correctif sont tirées du préambule commité par extraire_correctifs.py, pas retapées.
\documentclass{scrbook}
\ifdefined\SANSCROCHET \RemoveFromHook{package/amsthm/after}[firstaid/aliascounter]\fi
\usepackage{amsthm}
\usepackage{hyperref}
\usepackage[french, nameinlink, noabbrev]{cleveref}
\ifdefined\AMSSEUL
  \newtheorem{theoreme}{Théorème}[chapter]
  \newtheorem{proposition}[theoreme]{Proposition}
  \newtheorem{corollaire}[theoreme]{Corollaire}
  \newtheorem{remarque}{Remarque}[chapter]
\else
  \usepackage{thmtools}
  \ifdefined\CORR \input{correctif-a-pose}\fi
  \declaretheorem[numberwithin=chapter, name=Théorème]{theoreme}
  \ifdefined\TRACE \tracingmacros=2 \tracingonline=0 \fi
  \declaretheorem[sibling=theoreme, name=Proposition]{proposition}
  \ifdefined\TRACE \tracingmacros=0 \fi
  \declaretheorem[sibling=theoreme, name=Corollaire]{corollaire}
  \declaretheorem[numberwithin=chapter, name=Remarque]{remarque}
  \ifdefined\FAUX \newcounteralias{corollaire}{section}\newcounteralias{orphelin}{inexistant}\fi
  \ifdefined\CORR \input{correctif-a-remise}\fi
\fi
\crefname{theoreme}{théorème}{théorèmes}
\crefname{proposition}{proposition}{propositions}
\crefname{corollaire}{corollaire}{corollaires}
\crefname{remarque}{remarque}{remarques}
\begin{document}
\chapter{Premier}
\begin{proposition}\label{p1}a\end{proposition}
\begin{theoreme}\label{t1}b\end{theoreme}
\begin{remarque}\label{r1}c\end{remarque}
\begin{corollaire}\label{c1}d\end{corollaire}
\chapter{Second}
\begin{corollaire}\label{c2}e\end{corollaire}
\begin{theoreme}\label{t2}f\end{theoreme}
RENVOIS \ref{p1} \ref{t1} \ref{r1} \ref{c1} \ref{c2} \ref{t2} ; \cref{p1} ; \cref{t1} ; \cref{r1} ; \cref{c1} ; \cref{c2} ; \cref{t2} FIN
\end{document}
"                                   sans correctif
%   lualatex -jobname=corr     "\def\CORR{1}% Juge de la troisième passe : document minimal écrit de zéro pour la cause A (thmtools sous LaTeX 2026-06-01).
% Variantes, par une macro posée sur la ligne de commande (deux passes chacune, depuis memoire/essai/juge3/minimal/) :
%   lualatex -jobname=nu       "\input{minimal-juge3}"                                   sans correctif
%   lualatex -jobname=corr     "\def\CORR{1}\input{minimal-juge3}"                       correctif A du préambule
%   lualatex -jobname=crochet  "\def\SANSCROCHET{1}\input{minimal-juge3}"                sans correctif, crochet retiré
%   lualatex -jobname=amsseul  "\def\AMSSEUL{1}\input{minimal-juge3}"                    amsthm seul (témoin)
%   lualatex -jobname=faux     "\def\CORR{1}\def\FAUX{1}\input{minimal-juge3}"           deux vrais défauts : doivent tomber
%   lualatex -jobname=trace    "\def\TRACE{1}\input{minimal-juge3}"                      trace de l'appel
% Les lignes du correctif sont tirées du préambule commité par extraire_correctifs.py, pas retapées.
\documentclass{scrbook}
\ifdefined\SANSCROCHET \RemoveFromHook{package/amsthm/after}[firstaid/aliascounter]\fi
\usepackage{amsthm}
\usepackage{hyperref}
\usepackage[french, nameinlink, noabbrev]{cleveref}
\ifdefined\AMSSEUL
  \newtheorem{theoreme}{Théorème}[chapter]
  \newtheorem{proposition}[theoreme]{Proposition}
  \newtheorem{corollaire}[theoreme]{Corollaire}
  \newtheorem{remarque}{Remarque}[chapter]
\else
  \usepackage{thmtools}
  \ifdefined\CORR \input{correctif-a-pose}\fi
  \declaretheorem[numberwithin=chapter, name=Théorème]{theoreme}
  \ifdefined\TRACE \tracingmacros=2 \tracingonline=0 \fi
  \declaretheorem[sibling=theoreme, name=Proposition]{proposition}
  \ifdefined\TRACE \tracingmacros=0 \fi
  \declaretheorem[sibling=theoreme, name=Corollaire]{corollaire}
  \declaretheorem[numberwithin=chapter, name=Remarque]{remarque}
  \ifdefined\FAUX \newcounteralias{corollaire}{section}\newcounteralias{orphelin}{inexistant}\fi
  \ifdefined\CORR \input{correctif-a-remise}\fi
\fi
\crefname{theoreme}{théorème}{théorèmes}
\crefname{proposition}{proposition}{propositions}
\crefname{corollaire}{corollaire}{corollaires}
\crefname{remarque}{remarque}{remarques}
\begin{document}
\chapter{Premier}
\begin{proposition}\label{p1}a\end{proposition}
\begin{theoreme}\label{t1}b\end{theoreme}
\begin{remarque}\label{r1}c\end{remarque}
\begin{corollaire}\label{c1}d\end{corollaire}
\chapter{Second}
\begin{corollaire}\label{c2}e\end{corollaire}
\begin{theoreme}\label{t2}f\end{theoreme}
RENVOIS \ref{p1} \ref{t1} \ref{r1} \ref{c1} \ref{c2} \ref{t2} ; \cref{p1} ; \cref{t1} ; \cref{r1} ; \cref{c1} ; \cref{c2} ; \cref{t2} FIN
\end{document}
"                       correctif A du préambule
%   lualatex -jobname=crochet  "\def\SANSCROCHET{1}% Juge de la troisième passe : document minimal écrit de zéro pour la cause A (thmtools sous LaTeX 2026-06-01).
% Variantes, par une macro posée sur la ligne de commande (deux passes chacune, depuis memoire/essai/juge3/minimal/) :
%   lualatex -jobname=nu       "\input{minimal-juge3}"                                   sans correctif
%   lualatex -jobname=corr     "\def\CORR{1}\input{minimal-juge3}"                       correctif A du préambule
%   lualatex -jobname=crochet  "\def\SANSCROCHET{1}\input{minimal-juge3}"                sans correctif, crochet retiré
%   lualatex -jobname=amsseul  "\def\AMSSEUL{1}\input{minimal-juge3}"                    amsthm seul (témoin)
%   lualatex -jobname=faux     "\def\CORR{1}\def\FAUX{1}\input{minimal-juge3}"           deux vrais défauts : doivent tomber
%   lualatex -jobname=trace    "\def\TRACE{1}\input{minimal-juge3}"                      trace de l'appel
% Les lignes du correctif sont tirées du préambule commité par extraire_correctifs.py, pas retapées.
\documentclass{scrbook}
\ifdefined\SANSCROCHET \RemoveFromHook{package/amsthm/after}[firstaid/aliascounter]\fi
\usepackage{amsthm}
\usepackage{hyperref}
\usepackage[french, nameinlink, noabbrev]{cleveref}
\ifdefined\AMSSEUL
  \newtheorem{theoreme}{Théorème}[chapter]
  \newtheorem{proposition}[theoreme]{Proposition}
  \newtheorem{corollaire}[theoreme]{Corollaire}
  \newtheorem{remarque}{Remarque}[chapter]
\else
  \usepackage{thmtools}
  \ifdefined\CORR \input{correctif-a-pose}\fi
  \declaretheorem[numberwithin=chapter, name=Théorème]{theoreme}
  \ifdefined\TRACE \tracingmacros=2 \tracingonline=0 \fi
  \declaretheorem[sibling=theoreme, name=Proposition]{proposition}
  \ifdefined\TRACE \tracingmacros=0 \fi
  \declaretheorem[sibling=theoreme, name=Corollaire]{corollaire}
  \declaretheorem[numberwithin=chapter, name=Remarque]{remarque}
  \ifdefined\FAUX \newcounteralias{corollaire}{section}\newcounteralias{orphelin}{inexistant}\fi
  \ifdefined\CORR \input{correctif-a-remise}\fi
\fi
\crefname{theoreme}{théorème}{théorèmes}
\crefname{proposition}{proposition}{propositions}
\crefname{corollaire}{corollaire}{corollaires}
\crefname{remarque}{remarque}{remarques}
\begin{document}
\chapter{Premier}
\begin{proposition}\label{p1}a\end{proposition}
\begin{theoreme}\label{t1}b\end{theoreme}
\begin{remarque}\label{r1}c\end{remarque}
\begin{corollaire}\label{c1}d\end{corollaire}
\chapter{Second}
\begin{corollaire}\label{c2}e\end{corollaire}
\begin{theoreme}\label{t2}f\end{theoreme}
RENVOIS \ref{p1} \ref{t1} \ref{r1} \ref{c1} \ref{c2} \ref{t2} ; \cref{p1} ; \cref{t1} ; \cref{r1} ; \cref{c1} ; \cref{c2} ; \cref{t2} FIN
\end{document}
"                sans correctif, crochet retiré
%   lualatex -jobname=amsseul  "\def\AMSSEUL{1}% Juge de la troisième passe : document minimal écrit de zéro pour la cause A (thmtools sous LaTeX 2026-06-01).
% Variantes, par une macro posée sur la ligne de commande (deux passes chacune, depuis memoire/essai/juge3/minimal/) :
%   lualatex -jobname=nu       "\input{minimal-juge3}"                                   sans correctif
%   lualatex -jobname=corr     "\def\CORR{1}\input{minimal-juge3}"                       correctif A du préambule
%   lualatex -jobname=crochet  "\def\SANSCROCHET{1}\input{minimal-juge3}"                sans correctif, crochet retiré
%   lualatex -jobname=amsseul  "\def\AMSSEUL{1}\input{minimal-juge3}"                    amsthm seul (témoin)
%   lualatex -jobname=faux     "\def\CORR{1}\def\FAUX{1}\input{minimal-juge3}"           deux vrais défauts : doivent tomber
%   lualatex -jobname=trace    "\def\TRACE{1}\input{minimal-juge3}"                      trace de l'appel
% Les lignes du correctif sont tirées du préambule commité par extraire_correctifs.py, pas retapées.
\documentclass{scrbook}
\ifdefined\SANSCROCHET \RemoveFromHook{package/amsthm/after}[firstaid/aliascounter]\fi
\usepackage{amsthm}
\usepackage{hyperref}
\usepackage[french, nameinlink, noabbrev]{cleveref}
\ifdefined\AMSSEUL
  \newtheorem{theoreme}{Théorème}[chapter]
  \newtheorem{proposition}[theoreme]{Proposition}
  \newtheorem{corollaire}[theoreme]{Corollaire}
  \newtheorem{remarque}{Remarque}[chapter]
\else
  \usepackage{thmtools}
  \ifdefined\CORR \input{correctif-a-pose}\fi
  \declaretheorem[numberwithin=chapter, name=Théorème]{theoreme}
  \ifdefined\TRACE \tracingmacros=2 \tracingonline=0 \fi
  \declaretheorem[sibling=theoreme, name=Proposition]{proposition}
  \ifdefined\TRACE \tracingmacros=0 \fi
  \declaretheorem[sibling=theoreme, name=Corollaire]{corollaire}
  \declaretheorem[numberwithin=chapter, name=Remarque]{remarque}
  \ifdefined\FAUX \newcounteralias{corollaire}{section}\newcounteralias{orphelin}{inexistant}\fi
  \ifdefined\CORR \input{correctif-a-remise}\fi
\fi
\crefname{theoreme}{théorème}{théorèmes}
\crefname{proposition}{proposition}{propositions}
\crefname{corollaire}{corollaire}{corollaires}
\crefname{remarque}{remarque}{remarques}
\begin{document}
\chapter{Premier}
\begin{proposition}\label{p1}a\end{proposition}
\begin{theoreme}\label{t1}b\end{theoreme}
\begin{remarque}\label{r1}c\end{remarque}
\begin{corollaire}\label{c1}d\end{corollaire}
\chapter{Second}
\begin{corollaire}\label{c2}e\end{corollaire}
\begin{theoreme}\label{t2}f\end{theoreme}
RENVOIS \ref{p1} \ref{t1} \ref{r1} \ref{c1} \ref{c2} \ref{t2} ; \cref{p1} ; \cref{t1} ; \cref{r1} ; \cref{c1} ; \cref{c2} ; \cref{t2} FIN
\end{document}
"                    amsthm seul (témoin)
%   lualatex -jobname=faux     "\def\CORR{1}\def\FAUX{1}% Juge de la troisième passe : document minimal écrit de zéro pour la cause A (thmtools sous LaTeX 2026-06-01).
% Variantes, par une macro posée sur la ligne de commande (deux passes chacune, depuis memoire/essai/juge3/minimal/) :
%   lualatex -jobname=nu       "\input{minimal-juge3}"                                   sans correctif
%   lualatex -jobname=corr     "\def\CORR{1}\input{minimal-juge3}"                       correctif A du préambule
%   lualatex -jobname=crochet  "\def\SANSCROCHET{1}\input{minimal-juge3}"                sans correctif, crochet retiré
%   lualatex -jobname=amsseul  "\def\AMSSEUL{1}\input{minimal-juge3}"                    amsthm seul (témoin)
%   lualatex -jobname=faux     "\def\CORR{1}\def\FAUX{1}\input{minimal-juge3}"           deux vrais défauts : doivent tomber
%   lualatex -jobname=trace    "\def\TRACE{1}\input{minimal-juge3}"                      trace de l'appel
% Les lignes du correctif sont tirées du préambule commité par extraire_correctifs.py, pas retapées.
\documentclass{scrbook}
\ifdefined\SANSCROCHET \RemoveFromHook{package/amsthm/after}[firstaid/aliascounter]\fi
\usepackage{amsthm}
\usepackage{hyperref}
\usepackage[french, nameinlink, noabbrev]{cleveref}
\ifdefined\AMSSEUL
  \newtheorem{theoreme}{Théorème}[chapter]
  \newtheorem{proposition}[theoreme]{Proposition}
  \newtheorem{corollaire}[theoreme]{Corollaire}
  \newtheorem{remarque}{Remarque}[chapter]
\else
  \usepackage{thmtools}
  \ifdefined\CORR \input{correctif-a-pose}\fi
  \declaretheorem[numberwithin=chapter, name=Théorème]{theoreme}
  \ifdefined\TRACE \tracingmacros=2 \tracingonline=0 \fi
  \declaretheorem[sibling=theoreme, name=Proposition]{proposition}
  \ifdefined\TRACE \tracingmacros=0 \fi
  \declaretheorem[sibling=theoreme, name=Corollaire]{corollaire}
  \declaretheorem[numberwithin=chapter, name=Remarque]{remarque}
  \ifdefined\FAUX \newcounteralias{corollaire}{section}\newcounteralias{orphelin}{inexistant}\fi
  \ifdefined\CORR \input{correctif-a-remise}\fi
\fi
\crefname{theoreme}{théorème}{théorèmes}
\crefname{proposition}{proposition}{propositions}
\crefname{corollaire}{corollaire}{corollaires}
\crefname{remarque}{remarque}{remarques}
\begin{document}
\chapter{Premier}
\begin{proposition}\label{p1}a\end{proposition}
\begin{theoreme}\label{t1}b\end{theoreme}
\begin{remarque}\label{r1}c\end{remarque}
\begin{corollaire}\label{c1}d\end{corollaire}
\chapter{Second}
\begin{corollaire}\label{c2}e\end{corollaire}
\begin{theoreme}\label{t2}f\end{theoreme}
RENVOIS \ref{p1} \ref{t1} \ref{r1} \ref{c1} \ref{c2} \ref{t2} ; \cref{p1} ; \cref{t1} ; \cref{r1} ; \cref{c1} ; \cref{c2} ; \cref{t2} FIN
\end{document}
"           deux vrais défauts : doivent tomber
%   lualatex -jobname=trace    "\def\TRACE{1}% Juge de la troisième passe : document minimal écrit de zéro pour la cause A (thmtools sous LaTeX 2026-06-01).
% Variantes, par une macro posée sur la ligne de commande (deux passes chacune, depuis memoire/essai/juge3/minimal/) :
%   lualatex -jobname=nu       "\input{minimal-juge3}"                                   sans correctif
%   lualatex -jobname=corr     "\def\CORR{1}\input{minimal-juge3}"                       correctif A du préambule
%   lualatex -jobname=crochet  "\def\SANSCROCHET{1}\input{minimal-juge3}"                sans correctif, crochet retiré
%   lualatex -jobname=amsseul  "\def\AMSSEUL{1}\input{minimal-juge3}"                    amsthm seul (témoin)
%   lualatex -jobname=faux     "\def\CORR{1}\def\FAUX{1}\input{minimal-juge3}"           deux vrais défauts : doivent tomber
%   lualatex -jobname=trace    "\def\TRACE{1}\input{minimal-juge3}"                      trace de l'appel
% Les lignes du correctif sont tirées du préambule commité par extraire_correctifs.py, pas retapées.
\documentclass{scrbook}
\ifdefined\SANSCROCHET \RemoveFromHook{package/amsthm/after}[firstaid/aliascounter]\fi
\usepackage{amsthm}
\usepackage{hyperref}
\usepackage[french, nameinlink, noabbrev]{cleveref}
\ifdefined\AMSSEUL
  \newtheorem{theoreme}{Théorème}[chapter]
  \newtheorem{proposition}[theoreme]{Proposition}
  \newtheorem{corollaire}[theoreme]{Corollaire}
  \newtheorem{remarque}{Remarque}[chapter]
\else
  \usepackage{thmtools}
  \ifdefined\CORR \input{correctif-a-pose}\fi
  \declaretheorem[numberwithin=chapter, name=Théorème]{theoreme}
  \ifdefined\TRACE \tracingmacros=2 \tracingonline=0 \fi
  \declaretheorem[sibling=theoreme, name=Proposition]{proposition}
  \ifdefined\TRACE \tracingmacros=0 \fi
  \declaretheorem[sibling=theoreme, name=Corollaire]{corollaire}
  \declaretheorem[numberwithin=chapter, name=Remarque]{remarque}
  \ifdefined\FAUX \newcounteralias{corollaire}{section}\newcounteralias{orphelin}{inexistant}\fi
  \ifdefined\CORR \input{correctif-a-remise}\fi
\fi
\crefname{theoreme}{théorème}{théorèmes}
\crefname{proposition}{proposition}{propositions}
\crefname{corollaire}{corollaire}{corollaires}
\crefname{remarque}{remarque}{remarques}
\begin{document}
\chapter{Premier}
\begin{proposition}\label{p1}a\end{proposition}
\begin{theoreme}\label{t1}b\end{theoreme}
\begin{remarque}\label{r1}c\end{remarque}
\begin{corollaire}\label{c1}d\end{corollaire}
\chapter{Second}
\begin{corollaire}\label{c2}e\end{corollaire}
\begin{theoreme}\label{t2}f\end{theoreme}
RENVOIS \ref{p1} \ref{t1} \ref{r1} \ref{c1} \ref{c2} \ref{t2} ; \cref{p1} ; \cref{t1} ; \cref{r1} ; \cref{c1} ; \cref{c2} ; \cref{t2} FIN
\end{document}
"                      trace de l'appel
% Les lignes du correctif sont tirées du préambule commité par extraire_correctifs.py, pas retapées.
\documentclass{scrbook}
\ifdefined\SANSCROCHET \RemoveFromHook{package/amsthm/after}[firstaid/aliascounter]\fi
\usepackage{amsthm}
\usepackage{hyperref}
\usepackage[french, nameinlink, noabbrev]{cleveref}
\ifdefined\AMSSEUL
  \newtheorem{theoreme}{Théorème}[chapter]
  \newtheorem{proposition}[theoreme]{Proposition}
  \newtheorem{corollaire}[theoreme]{Corollaire}
  \newtheorem{remarque}{Remarque}[chapter]
\else
  \usepackage{thmtools}
  \ifdefined\CORR \input{correctif-a-pose}\fi
  \declaretheorem[numberwithin=chapter, name=Théorème]{theoreme}
  \ifdefined\TRACE \tracingmacros=2 \tracingonline=0 \fi
  \declaretheorem[sibling=theoreme, name=Proposition]{proposition}
  \ifdefined\TRACE \tracingmacros=0 \fi
  \declaretheorem[sibling=theoreme, name=Corollaire]{corollaire}
  \declaretheorem[numberwithin=chapter, name=Remarque]{remarque}
  \ifdefined\FAUX \newcounteralias{corollaire}{section}\newcounteralias{orphelin}{inexistant}\fi
  \ifdefined\CORR \input{correctif-a-remise}\fi
\fi
\crefname{theoreme}{théorème}{théorèmes}
\crefname{proposition}{proposition}{propositions}
\crefname{corollaire}{corollaire}{corollaires}
\crefname{remarque}{remarque}{remarques}
\begin{document}
\chapter{Premier}
\begin{proposition}\label{p1}a\end{proposition}
\begin{theoreme}\label{t1}b\end{theoreme}
\begin{remarque}\label{r1}c\end{remarque}
\begin{corollaire}\label{c1}d\end{corollaire}
\chapter{Second}
\begin{corollaire}\label{c2}e\end{corollaire}
\begin{theoreme}\label{t2}f\end{theoreme}
RENVOIS \ref{p1} \ref{t1} \ref{r1} \ref{c1} \ref{c2} \ref{t2} ; \cref{p1} ; \cref{t1} ; \cref{r1} ; \cref{c1} ; \cref{c2} ; \cref{t2} FIN
\end{document}
"                    amsthm seul (témoin)
%   lualatex -jobname=faux     "\def\CORR{1}\def\FAUX{1}% Juge de la troisième passe : document minimal écrit de zéro pour la cause A (thmtools sous LaTeX 2026-06-01).
% Variantes, par une macro posée sur la ligne de commande (deux passes chacune, depuis memoire/essai/juge3/minimal/) :
%   lualatex -jobname=nu       "% Juge de la troisième passe : document minimal écrit de zéro pour la cause A (thmtools sous LaTeX 2026-06-01).
% Variantes, par une macro posée sur la ligne de commande (deux passes chacune, depuis memoire/essai/juge3/minimal/) :
%   lualatex -jobname=nu       "\input{minimal-juge3}"                                   sans correctif
%   lualatex -jobname=corr     "\def\CORR{1}\input{minimal-juge3}"                       correctif A du préambule
%   lualatex -jobname=crochet  "\def\SANSCROCHET{1}\input{minimal-juge3}"                sans correctif, crochet retiré
%   lualatex -jobname=amsseul  "\def\AMSSEUL{1}\input{minimal-juge3}"                    amsthm seul (témoin)
%   lualatex -jobname=faux     "\def\CORR{1}\def\FAUX{1}\input{minimal-juge3}"           deux vrais défauts : doivent tomber
%   lualatex -jobname=trace    "\def\TRACE{1}\input{minimal-juge3}"                      trace de l'appel
% Les lignes du correctif sont tirées du préambule commité par extraire_correctifs.py, pas retapées.
\documentclass{scrbook}
\ifdefined\SANSCROCHET \RemoveFromHook{package/amsthm/after}[firstaid/aliascounter]\fi
\usepackage{amsthm}
\usepackage{hyperref}
\usepackage[french, nameinlink, noabbrev]{cleveref}
\ifdefined\AMSSEUL
  \newtheorem{theoreme}{Théorème}[chapter]
  \newtheorem{proposition}[theoreme]{Proposition}
  \newtheorem{corollaire}[theoreme]{Corollaire}
  \newtheorem{remarque}{Remarque}[chapter]
\else
  \usepackage{thmtools}
  \ifdefined\CORR \input{correctif-a-pose}\fi
  \declaretheorem[numberwithin=chapter, name=Théorème]{theoreme}
  \ifdefined\TRACE \tracingmacros=2 \tracingonline=0 \fi
  \declaretheorem[sibling=theoreme, name=Proposition]{proposition}
  \ifdefined\TRACE \tracingmacros=0 \fi
  \declaretheorem[sibling=theoreme, name=Corollaire]{corollaire}
  \declaretheorem[numberwithin=chapter, name=Remarque]{remarque}
  \ifdefined\FAUX \newcounteralias{corollaire}{section}\newcounteralias{orphelin}{inexistant}\fi
  \ifdefined\CORR \input{correctif-a-remise}\fi
\fi
\crefname{theoreme}{théorème}{théorèmes}
\crefname{proposition}{proposition}{propositions}
\crefname{corollaire}{corollaire}{corollaires}
\crefname{remarque}{remarque}{remarques}
\begin{document}
\chapter{Premier}
\begin{proposition}\label{p1}a\end{proposition}
\begin{theoreme}\label{t1}b\end{theoreme}
\begin{remarque}\label{r1}c\end{remarque}
\begin{corollaire}\label{c1}d\end{corollaire}
\chapter{Second}
\begin{corollaire}\label{c2}e\end{corollaire}
\begin{theoreme}\label{t2}f\end{theoreme}
RENVOIS \ref{p1} \ref{t1} \ref{r1} \ref{c1} \ref{c2} \ref{t2} ; \cref{p1} ; \cref{t1} ; \cref{r1} ; \cref{c1} ; \cref{c2} ; \cref{t2} FIN
\end{document}
"                                   sans correctif
%   lualatex -jobname=corr     "\def\CORR{1}% Juge de la troisième passe : document minimal écrit de zéro pour la cause A (thmtools sous LaTeX 2026-06-01).
% Variantes, par une macro posée sur la ligne de commande (deux passes chacune, depuis memoire/essai/juge3/minimal/) :
%   lualatex -jobname=nu       "\input{minimal-juge3}"                                   sans correctif
%   lualatex -jobname=corr     "\def\CORR{1}\input{minimal-juge3}"                       correctif A du préambule
%   lualatex -jobname=crochet  "\def\SANSCROCHET{1}\input{minimal-juge3}"                sans correctif, crochet retiré
%   lualatex -jobname=amsseul  "\def\AMSSEUL{1}\input{minimal-juge3}"                    amsthm seul (témoin)
%   lualatex -jobname=faux     "\def\CORR{1}\def\FAUX{1}\input{minimal-juge3}"           deux vrais défauts : doivent tomber
%   lualatex -jobname=trace    "\def\TRACE{1}\input{minimal-juge3}"                      trace de l'appel
% Les lignes du correctif sont tirées du préambule commité par extraire_correctifs.py, pas retapées.
\documentclass{scrbook}
\ifdefined\SANSCROCHET \RemoveFromHook{package/amsthm/after}[firstaid/aliascounter]\fi
\usepackage{amsthm}
\usepackage{hyperref}
\usepackage[french, nameinlink, noabbrev]{cleveref}
\ifdefined\AMSSEUL
  \newtheorem{theoreme}{Théorème}[chapter]
  \newtheorem{proposition}[theoreme]{Proposition}
  \newtheorem{corollaire}[theoreme]{Corollaire}
  \newtheorem{remarque}{Remarque}[chapter]
\else
  \usepackage{thmtools}
  \ifdefined\CORR \input{correctif-a-pose}\fi
  \declaretheorem[numberwithin=chapter, name=Théorème]{theoreme}
  \ifdefined\TRACE \tracingmacros=2 \tracingonline=0 \fi
  \declaretheorem[sibling=theoreme, name=Proposition]{proposition}
  \ifdefined\TRACE \tracingmacros=0 \fi
  \declaretheorem[sibling=theoreme, name=Corollaire]{corollaire}
  \declaretheorem[numberwithin=chapter, name=Remarque]{remarque}
  \ifdefined\FAUX \newcounteralias{corollaire}{section}\newcounteralias{orphelin}{inexistant}\fi
  \ifdefined\CORR \input{correctif-a-remise}\fi
\fi
\crefname{theoreme}{théorème}{théorèmes}
\crefname{proposition}{proposition}{propositions}
\crefname{corollaire}{corollaire}{corollaires}
\crefname{remarque}{remarque}{remarques}
\begin{document}
\chapter{Premier}
\begin{proposition}\label{p1}a\end{proposition}
\begin{theoreme}\label{t1}b\end{theoreme}
\begin{remarque}\label{r1}c\end{remarque}
\begin{corollaire}\label{c1}d\end{corollaire}
\chapter{Second}
\begin{corollaire}\label{c2}e\end{corollaire}
\begin{theoreme}\label{t2}f\end{theoreme}
RENVOIS \ref{p1} \ref{t1} \ref{r1} \ref{c1} \ref{c2} \ref{t2} ; \cref{p1} ; \cref{t1} ; \cref{r1} ; \cref{c1} ; \cref{c2} ; \cref{t2} FIN
\end{document}
"                       correctif A du préambule
%   lualatex -jobname=crochet  "\def\SANSCROCHET{1}% Juge de la troisième passe : document minimal écrit de zéro pour la cause A (thmtools sous LaTeX 2026-06-01).
% Variantes, par une macro posée sur la ligne de commande (deux passes chacune, depuis memoire/essai/juge3/minimal/) :
%   lualatex -jobname=nu       "\input{minimal-juge3}"                                   sans correctif
%   lualatex -jobname=corr     "\def\CORR{1}\input{minimal-juge3}"                       correctif A du préambule
%   lualatex -jobname=crochet  "\def\SANSCROCHET{1}\input{minimal-juge3}"                sans correctif, crochet retiré
%   lualatex -jobname=amsseul  "\def\AMSSEUL{1}\input{minimal-juge3}"                    amsthm seul (témoin)
%   lualatex -jobname=faux     "\def\CORR{1}\def\FAUX{1}\input{minimal-juge3}"           deux vrais défauts : doivent tomber
%   lualatex -jobname=trace    "\def\TRACE{1}\input{minimal-juge3}"                      trace de l'appel
% Les lignes du correctif sont tirées du préambule commité par extraire_correctifs.py, pas retapées.
\documentclass{scrbook}
\ifdefined\SANSCROCHET \RemoveFromHook{package/amsthm/after}[firstaid/aliascounter]\fi
\usepackage{amsthm}
\usepackage{hyperref}
\usepackage[french, nameinlink, noabbrev]{cleveref}
\ifdefined\AMSSEUL
  \newtheorem{theoreme}{Théorème}[chapter]
  \newtheorem{proposition}[theoreme]{Proposition}
  \newtheorem{corollaire}[theoreme]{Corollaire}
  \newtheorem{remarque}{Remarque}[chapter]
\else
  \usepackage{thmtools}
  \ifdefined\CORR \input{correctif-a-pose}\fi
  \declaretheorem[numberwithin=chapter, name=Théorème]{theoreme}
  \ifdefined\TRACE \tracingmacros=2 \tracingonline=0 \fi
  \declaretheorem[sibling=theoreme, name=Proposition]{proposition}
  \ifdefined\TRACE \tracingmacros=0 \fi
  \declaretheorem[sibling=theoreme, name=Corollaire]{corollaire}
  \declaretheorem[numberwithin=chapter, name=Remarque]{remarque}
  \ifdefined\FAUX \newcounteralias{corollaire}{section}\newcounteralias{orphelin}{inexistant}\fi
  \ifdefined\CORR \input{correctif-a-remise}\fi
\fi
\crefname{theoreme}{théorème}{théorèmes}
\crefname{proposition}{proposition}{propositions}
\crefname{corollaire}{corollaire}{corollaires}
\crefname{remarque}{remarque}{remarques}
\begin{document}
\chapter{Premier}
\begin{proposition}\label{p1}a\end{proposition}
\begin{theoreme}\label{t1}b\end{theoreme}
\begin{remarque}\label{r1}c\end{remarque}
\begin{corollaire}\label{c1}d\end{corollaire}
\chapter{Second}
\begin{corollaire}\label{c2}e\end{corollaire}
\begin{theoreme}\label{t2}f\end{theoreme}
RENVOIS \ref{p1} \ref{t1} \ref{r1} \ref{c1} \ref{c2} \ref{t2} ; \cref{p1} ; \cref{t1} ; \cref{r1} ; \cref{c1} ; \cref{c2} ; \cref{t2} FIN
\end{document}
"                sans correctif, crochet retiré
%   lualatex -jobname=amsseul  "\def\AMSSEUL{1}% Juge de la troisième passe : document minimal écrit de zéro pour la cause A (thmtools sous LaTeX 2026-06-01).
% Variantes, par une macro posée sur la ligne de commande (deux passes chacune, depuis memoire/essai/juge3/minimal/) :
%   lualatex -jobname=nu       "\input{minimal-juge3}"                                   sans correctif
%   lualatex -jobname=corr     "\def\CORR{1}\input{minimal-juge3}"                       correctif A du préambule
%   lualatex -jobname=crochet  "\def\SANSCROCHET{1}\input{minimal-juge3}"                sans correctif, crochet retiré
%   lualatex -jobname=amsseul  "\def\AMSSEUL{1}\input{minimal-juge3}"                    amsthm seul (témoin)
%   lualatex -jobname=faux     "\def\CORR{1}\def\FAUX{1}\input{minimal-juge3}"           deux vrais défauts : doivent tomber
%   lualatex -jobname=trace    "\def\TRACE{1}\input{minimal-juge3}"                      trace de l'appel
% Les lignes du correctif sont tirées du préambule commité par extraire_correctifs.py, pas retapées.
\documentclass{scrbook}
\ifdefined\SANSCROCHET \RemoveFromHook{package/amsthm/after}[firstaid/aliascounter]\fi
\usepackage{amsthm}
\usepackage{hyperref}
\usepackage[french, nameinlink, noabbrev]{cleveref}
\ifdefined\AMSSEUL
  \newtheorem{theoreme}{Théorème}[chapter]
  \newtheorem{proposition}[theoreme]{Proposition}
  \newtheorem{corollaire}[theoreme]{Corollaire}
  \newtheorem{remarque}{Remarque}[chapter]
\else
  \usepackage{thmtools}
  \ifdefined\CORR \input{correctif-a-pose}\fi
  \declaretheorem[numberwithin=chapter, name=Théorème]{theoreme}
  \ifdefined\TRACE \tracingmacros=2 \tracingonline=0 \fi
  \declaretheorem[sibling=theoreme, name=Proposition]{proposition}
  \ifdefined\TRACE \tracingmacros=0 \fi
  \declaretheorem[sibling=theoreme, name=Corollaire]{corollaire}
  \declaretheorem[numberwithin=chapter, name=Remarque]{remarque}
  \ifdefined\FAUX \newcounteralias{corollaire}{section}\newcounteralias{orphelin}{inexistant}\fi
  \ifdefined\CORR \input{correctif-a-remise}\fi
\fi
\crefname{theoreme}{théorème}{théorèmes}
\crefname{proposition}{proposition}{propositions}
\crefname{corollaire}{corollaire}{corollaires}
\crefname{remarque}{remarque}{remarques}
\begin{document}
\chapter{Premier}
\begin{proposition}\label{p1}a\end{proposition}
\begin{theoreme}\label{t1}b\end{theoreme}
\begin{remarque}\label{r1}c\end{remarque}
\begin{corollaire}\label{c1}d\end{corollaire}
\chapter{Second}
\begin{corollaire}\label{c2}e\end{corollaire}
\begin{theoreme}\label{t2}f\end{theoreme}
RENVOIS \ref{p1} \ref{t1} \ref{r1} \ref{c1} \ref{c2} \ref{t2} ; \cref{p1} ; \cref{t1} ; \cref{r1} ; \cref{c1} ; \cref{c2} ; \cref{t2} FIN
\end{document}
"                    amsthm seul (témoin)
%   lualatex -jobname=faux     "\def\CORR{1}\def\FAUX{1}% Juge de la troisième passe : document minimal écrit de zéro pour la cause A (thmtools sous LaTeX 2026-06-01).
% Variantes, par une macro posée sur la ligne de commande (deux passes chacune, depuis memoire/essai/juge3/minimal/) :
%   lualatex -jobname=nu       "\input{minimal-juge3}"                                   sans correctif
%   lualatex -jobname=corr     "\def\CORR{1}\input{minimal-juge3}"                       correctif A du préambule
%   lualatex -jobname=crochet  "\def\SANSCROCHET{1}\input{minimal-juge3}"                sans correctif, crochet retiré
%   lualatex -jobname=amsseul  "\def\AMSSEUL{1}\input{minimal-juge3}"                    amsthm seul (témoin)
%   lualatex -jobname=faux     "\def\CORR{1}\def\FAUX{1}\input{minimal-juge3}"           deux vrais défauts : doivent tomber
%   lualatex -jobname=trace    "\def\TRACE{1}\input{minimal-juge3}"                      trace de l'appel
% Les lignes du correctif sont tirées du préambule commité par extraire_correctifs.py, pas retapées.
\documentclass{scrbook}
\ifdefined\SANSCROCHET \RemoveFromHook{package/amsthm/after}[firstaid/aliascounter]\fi
\usepackage{amsthm}
\usepackage{hyperref}
\usepackage[french, nameinlink, noabbrev]{cleveref}
\ifdefined\AMSSEUL
  \newtheorem{theoreme}{Théorème}[chapter]
  \newtheorem{proposition}[theoreme]{Proposition}
  \newtheorem{corollaire}[theoreme]{Corollaire}
  \newtheorem{remarque}{Remarque}[chapter]
\else
  \usepackage{thmtools}
  \ifdefined\CORR \input{correctif-a-pose}\fi
  \declaretheorem[numberwithin=chapter, name=Théorème]{theoreme}
  \ifdefined\TRACE \tracingmacros=2 \tracingonline=0 \fi
  \declaretheorem[sibling=theoreme, name=Proposition]{proposition}
  \ifdefined\TRACE \tracingmacros=0 \fi
  \declaretheorem[sibling=theoreme, name=Corollaire]{corollaire}
  \declaretheorem[numberwithin=chapter, name=Remarque]{remarque}
  \ifdefined\FAUX \newcounteralias{corollaire}{section}\newcounteralias{orphelin}{inexistant}\fi
  \ifdefined\CORR \input{correctif-a-remise}\fi
\fi
\crefname{theoreme}{théorème}{théorèmes}
\crefname{proposition}{proposition}{propositions}
\crefname{corollaire}{corollaire}{corollaires}
\crefname{remarque}{remarque}{remarques}
\begin{document}
\chapter{Premier}
\begin{proposition}\label{p1}a\end{proposition}
\begin{theoreme}\label{t1}b\end{theoreme}
\begin{remarque}\label{r1}c\end{remarque}
\begin{corollaire}\label{c1}d\end{corollaire}
\chapter{Second}
\begin{corollaire}\label{c2}e\end{corollaire}
\begin{theoreme}\label{t2}f\end{theoreme}
RENVOIS \ref{p1} \ref{t1} \ref{r1} \ref{c1} \ref{c2} \ref{t2} ; \cref{p1} ; \cref{t1} ; \cref{r1} ; \cref{c1} ; \cref{c2} ; \cref{t2} FIN
\end{document}
"           deux vrais défauts : doivent tomber
%   lualatex -jobname=trace    "\def\TRACE{1}% Juge de la troisième passe : document minimal écrit de zéro pour la cause A (thmtools sous LaTeX 2026-06-01).
% Variantes, par une macro posée sur la ligne de commande (deux passes chacune, depuis memoire/essai/juge3/minimal/) :
%   lualatex -jobname=nu       "\input{minimal-juge3}"                                   sans correctif
%   lualatex -jobname=corr     "\def\CORR{1}\input{minimal-juge3}"                       correctif A du préambule
%   lualatex -jobname=crochet  "\def\SANSCROCHET{1}\input{minimal-juge3}"                sans correctif, crochet retiré
%   lualatex -jobname=amsseul  "\def\AMSSEUL{1}\input{minimal-juge3}"                    amsthm seul (témoin)
%   lualatex -jobname=faux     "\def\CORR{1}\def\FAUX{1}\input{minimal-juge3}"           deux vrais défauts : doivent tomber
%   lualatex -jobname=trace    "\def\TRACE{1}\input{minimal-juge3}"                      trace de l'appel
% Les lignes du correctif sont tirées du préambule commité par extraire_correctifs.py, pas retapées.
\documentclass{scrbook}
\ifdefined\SANSCROCHET \RemoveFromHook{package/amsthm/after}[firstaid/aliascounter]\fi
\usepackage{amsthm}
\usepackage{hyperref}
\usepackage[french, nameinlink, noabbrev]{cleveref}
\ifdefined\AMSSEUL
  \newtheorem{theoreme}{Théorème}[chapter]
  \newtheorem{proposition}[theoreme]{Proposition}
  \newtheorem{corollaire}[theoreme]{Corollaire}
  \newtheorem{remarque}{Remarque}[chapter]
\else
  \usepackage{thmtools}
  \ifdefined\CORR \input{correctif-a-pose}\fi
  \declaretheorem[numberwithin=chapter, name=Théorème]{theoreme}
  \ifdefined\TRACE \tracingmacros=2 \tracingonline=0 \fi
  \declaretheorem[sibling=theoreme, name=Proposition]{proposition}
  \ifdefined\TRACE \tracingmacros=0 \fi
  \declaretheorem[sibling=theoreme, name=Corollaire]{corollaire}
  \declaretheorem[numberwithin=chapter, name=Remarque]{remarque}
  \ifdefined\FAUX \newcounteralias{corollaire}{section}\newcounteralias{orphelin}{inexistant}\fi
  \ifdefined\CORR \input{correctif-a-remise}\fi
\fi
\crefname{theoreme}{théorème}{théorèmes}
\crefname{proposition}{proposition}{propositions}
\crefname{corollaire}{corollaire}{corollaires}
\crefname{remarque}{remarque}{remarques}
\begin{document}
\chapter{Premier}
\begin{proposition}\label{p1}a\end{proposition}
\begin{theoreme}\label{t1}b\end{theoreme}
\begin{remarque}\label{r1}c\end{remarque}
\begin{corollaire}\label{c1}d\end{corollaire}
\chapter{Second}
\begin{corollaire}\label{c2}e\end{corollaire}
\begin{theoreme}\label{t2}f\end{theoreme}
RENVOIS \ref{p1} \ref{t1} \ref{r1} \ref{c1} \ref{c2} \ref{t2} ; \cref{p1} ; \cref{t1} ; \cref{r1} ; \cref{c1} ; \cref{c2} ; \cref{t2} FIN
\end{document}
"                      trace de l'appel
% Les lignes du correctif sont tirées du préambule commité par extraire_correctifs.py, pas retapées.
\documentclass{scrbook}
\ifdefined\SANSCROCHET \RemoveFromHook{package/amsthm/after}[firstaid/aliascounter]\fi
\usepackage{amsthm}
\usepackage{hyperref}
\usepackage[french, nameinlink, noabbrev]{cleveref}
\ifdefined\AMSSEUL
  \newtheorem{theoreme}{Théorème}[chapter]
  \newtheorem{proposition}[theoreme]{Proposition}
  \newtheorem{corollaire}[theoreme]{Corollaire}
  \newtheorem{remarque}{Remarque}[chapter]
\else
  \usepackage{thmtools}
  \ifdefined\CORR \input{correctif-a-pose}\fi
  \declaretheorem[numberwithin=chapter, name=Théorème]{theoreme}
  \ifdefined\TRACE \tracingmacros=2 \tracingonline=0 \fi
  \declaretheorem[sibling=theoreme, name=Proposition]{proposition}
  \ifdefined\TRACE \tracingmacros=0 \fi
  \declaretheorem[sibling=theoreme, name=Corollaire]{corollaire}
  \declaretheorem[numberwithin=chapter, name=Remarque]{remarque}
  \ifdefined\FAUX \newcounteralias{corollaire}{section}\newcounteralias{orphelin}{inexistant}\fi
  \ifdefined\CORR \input{correctif-a-remise}\fi
\fi
\crefname{theoreme}{théorème}{théorèmes}
\crefname{proposition}{proposition}{propositions}
\crefname{corollaire}{corollaire}{corollaires}
\crefname{remarque}{remarque}{remarques}
\begin{document}
\chapter{Premier}
\begin{proposition}\label{p1}a\end{proposition}
\begin{theoreme}\label{t1}b\end{theoreme}
\begin{remarque}\label{r1}c\end{remarque}
\begin{corollaire}\label{c1}d\end{corollaire}
\chapter{Second}
\begin{corollaire}\label{c2}e\end{corollaire}
\begin{theoreme}\label{t2}f\end{theoreme}
RENVOIS \ref{p1} \ref{t1} \ref{r1} \ref{c1} \ref{c2} \ref{t2} ; \cref{p1} ; \cref{t1} ; \cref{r1} ; \cref{c1} ; \cref{c2} ; \cref{t2} FIN
\end{document}
"           deux vrais défauts : doivent tomber
%   lualatex -jobname=trace    "\def\TRACE{1}% Juge de la troisième passe : document minimal écrit de zéro pour la cause A (thmtools sous LaTeX 2026-06-01).
% Variantes, par une macro posée sur la ligne de commande (deux passes chacune, depuis memoire/essai/juge3/minimal/) :
%   lualatex -jobname=nu       "% Juge de la troisième passe : document minimal écrit de zéro pour la cause A (thmtools sous LaTeX 2026-06-01).
% Variantes, par une macro posée sur la ligne de commande (deux passes chacune, depuis memoire/essai/juge3/minimal/) :
%   lualatex -jobname=nu       "\input{minimal-juge3}"                                   sans correctif
%   lualatex -jobname=corr     "\def\CORR{1}\input{minimal-juge3}"                       correctif A du préambule
%   lualatex -jobname=crochet  "\def\SANSCROCHET{1}\input{minimal-juge3}"                sans correctif, crochet retiré
%   lualatex -jobname=amsseul  "\def\AMSSEUL{1}\input{minimal-juge3}"                    amsthm seul (témoin)
%   lualatex -jobname=faux     "\def\CORR{1}\def\FAUX{1}\input{minimal-juge3}"           deux vrais défauts : doivent tomber
%   lualatex -jobname=trace    "\def\TRACE{1}\input{minimal-juge3}"                      trace de l'appel
% Les lignes du correctif sont tirées du préambule commité par extraire_correctifs.py, pas retapées.
\documentclass{scrbook}
\ifdefined\SANSCROCHET \RemoveFromHook{package/amsthm/after}[firstaid/aliascounter]\fi
\usepackage{amsthm}
\usepackage{hyperref}
\usepackage[french, nameinlink, noabbrev]{cleveref}
\ifdefined\AMSSEUL
  \newtheorem{theoreme}{Théorème}[chapter]
  \newtheorem{proposition}[theoreme]{Proposition}
  \newtheorem{corollaire}[theoreme]{Corollaire}
  \newtheorem{remarque}{Remarque}[chapter]
\else
  \usepackage{thmtools}
  \ifdefined\CORR \input{correctif-a-pose}\fi
  \declaretheorem[numberwithin=chapter, name=Théorème]{theoreme}
  \ifdefined\TRACE \tracingmacros=2 \tracingonline=0 \fi
  \declaretheorem[sibling=theoreme, name=Proposition]{proposition}
  \ifdefined\TRACE \tracingmacros=0 \fi
  \declaretheorem[sibling=theoreme, name=Corollaire]{corollaire}
  \declaretheorem[numberwithin=chapter, name=Remarque]{remarque}
  \ifdefined\FAUX \newcounteralias{corollaire}{section}\newcounteralias{orphelin}{inexistant}\fi
  \ifdefined\CORR \input{correctif-a-remise}\fi
\fi
\crefname{theoreme}{théorème}{théorèmes}
\crefname{proposition}{proposition}{propositions}
\crefname{corollaire}{corollaire}{corollaires}
\crefname{remarque}{remarque}{remarques}
\begin{document}
\chapter{Premier}
\begin{proposition}\label{p1}a\end{proposition}
\begin{theoreme}\label{t1}b\end{theoreme}
\begin{remarque}\label{r1}c\end{remarque}
\begin{corollaire}\label{c1}d\end{corollaire}
\chapter{Second}
\begin{corollaire}\label{c2}e\end{corollaire}
\begin{theoreme}\label{t2}f\end{theoreme}
RENVOIS \ref{p1} \ref{t1} \ref{r1} \ref{c1} \ref{c2} \ref{t2} ; \cref{p1} ; \cref{t1} ; \cref{r1} ; \cref{c1} ; \cref{c2} ; \cref{t2} FIN
\end{document}
"                                   sans correctif
%   lualatex -jobname=corr     "\def\CORR{1}% Juge de la troisième passe : document minimal écrit de zéro pour la cause A (thmtools sous LaTeX 2026-06-01).
% Variantes, par une macro posée sur la ligne de commande (deux passes chacune, depuis memoire/essai/juge3/minimal/) :
%   lualatex -jobname=nu       "\input{minimal-juge3}"                                   sans correctif
%   lualatex -jobname=corr     "\def\CORR{1}\input{minimal-juge3}"                       correctif A du préambule
%   lualatex -jobname=crochet  "\def\SANSCROCHET{1}\input{minimal-juge3}"                sans correctif, crochet retiré
%   lualatex -jobname=amsseul  "\def\AMSSEUL{1}\input{minimal-juge3}"                    amsthm seul (témoin)
%   lualatex -jobname=faux     "\def\CORR{1}\def\FAUX{1}\input{minimal-juge3}"           deux vrais défauts : doivent tomber
%   lualatex -jobname=trace    "\def\TRACE{1}\input{minimal-juge3}"                      trace de l'appel
% Les lignes du correctif sont tirées du préambule commité par extraire_correctifs.py, pas retapées.
\documentclass{scrbook}
\ifdefined\SANSCROCHET \RemoveFromHook{package/amsthm/after}[firstaid/aliascounter]\fi
\usepackage{amsthm}
\usepackage{hyperref}
\usepackage[french, nameinlink, noabbrev]{cleveref}
\ifdefined\AMSSEUL
  \newtheorem{theoreme}{Théorème}[chapter]
  \newtheorem{proposition}[theoreme]{Proposition}
  \newtheorem{corollaire}[theoreme]{Corollaire}
  \newtheorem{remarque}{Remarque}[chapter]
\else
  \usepackage{thmtools}
  \ifdefined\CORR \input{correctif-a-pose}\fi
  \declaretheorem[numberwithin=chapter, name=Théorème]{theoreme}
  \ifdefined\TRACE \tracingmacros=2 \tracingonline=0 \fi
  \declaretheorem[sibling=theoreme, name=Proposition]{proposition}
  \ifdefined\TRACE \tracingmacros=0 \fi
  \declaretheorem[sibling=theoreme, name=Corollaire]{corollaire}
  \declaretheorem[numberwithin=chapter, name=Remarque]{remarque}
  \ifdefined\FAUX \newcounteralias{corollaire}{section}\newcounteralias{orphelin}{inexistant}\fi
  \ifdefined\CORR \input{correctif-a-remise}\fi
\fi
\crefname{theoreme}{théorème}{théorèmes}
\crefname{proposition}{proposition}{propositions}
\crefname{corollaire}{corollaire}{corollaires}
\crefname{remarque}{remarque}{remarques}
\begin{document}
\chapter{Premier}
\begin{proposition}\label{p1}a\end{proposition}
\begin{theoreme}\label{t1}b\end{theoreme}
\begin{remarque}\label{r1}c\end{remarque}
\begin{corollaire}\label{c1}d\end{corollaire}
\chapter{Second}
\begin{corollaire}\label{c2}e\end{corollaire}
\begin{theoreme}\label{t2}f\end{theoreme}
RENVOIS \ref{p1} \ref{t1} \ref{r1} \ref{c1} \ref{c2} \ref{t2} ; \cref{p1} ; \cref{t1} ; \cref{r1} ; \cref{c1} ; \cref{c2} ; \cref{t2} FIN
\end{document}
"                       correctif A du préambule
%   lualatex -jobname=crochet  "\def\SANSCROCHET{1}% Juge de la troisième passe : document minimal écrit de zéro pour la cause A (thmtools sous LaTeX 2026-06-01).
% Variantes, par une macro posée sur la ligne de commande (deux passes chacune, depuis memoire/essai/juge3/minimal/) :
%   lualatex -jobname=nu       "\input{minimal-juge3}"                                   sans correctif
%   lualatex -jobname=corr     "\def\CORR{1}\input{minimal-juge3}"                       correctif A du préambule
%   lualatex -jobname=crochet  "\def\SANSCROCHET{1}\input{minimal-juge3}"                sans correctif, crochet retiré
%   lualatex -jobname=amsseul  "\def\AMSSEUL{1}\input{minimal-juge3}"                    amsthm seul (témoin)
%   lualatex -jobname=faux     "\def\CORR{1}\def\FAUX{1}\input{minimal-juge3}"           deux vrais défauts : doivent tomber
%   lualatex -jobname=trace    "\def\TRACE{1}\input{minimal-juge3}"                      trace de l'appel
% Les lignes du correctif sont tirées du préambule commité par extraire_correctifs.py, pas retapées.
\documentclass{scrbook}
\ifdefined\SANSCROCHET \RemoveFromHook{package/amsthm/after}[firstaid/aliascounter]\fi
\usepackage{amsthm}
\usepackage{hyperref}
\usepackage[french, nameinlink, noabbrev]{cleveref}
\ifdefined\AMSSEUL
  \newtheorem{theoreme}{Théorème}[chapter]
  \newtheorem{proposition}[theoreme]{Proposition}
  \newtheorem{corollaire}[theoreme]{Corollaire}
  \newtheorem{remarque}{Remarque}[chapter]
\else
  \usepackage{thmtools}
  \ifdefined\CORR \input{correctif-a-pose}\fi
  \declaretheorem[numberwithin=chapter, name=Théorème]{theoreme}
  \ifdefined\TRACE \tracingmacros=2 \tracingonline=0 \fi
  \declaretheorem[sibling=theoreme, name=Proposition]{proposition}
  \ifdefined\TRACE \tracingmacros=0 \fi
  \declaretheorem[sibling=theoreme, name=Corollaire]{corollaire}
  \declaretheorem[numberwithin=chapter, name=Remarque]{remarque}
  \ifdefined\FAUX \newcounteralias{corollaire}{section}\newcounteralias{orphelin}{inexistant}\fi
  \ifdefined\CORR \input{correctif-a-remise}\fi
\fi
\crefname{theoreme}{théorème}{théorèmes}
\crefname{proposition}{proposition}{propositions}
\crefname{corollaire}{corollaire}{corollaires}
\crefname{remarque}{remarque}{remarques}
\begin{document}
\chapter{Premier}
\begin{proposition}\label{p1}a\end{proposition}
\begin{theoreme}\label{t1}b\end{theoreme}
\begin{remarque}\label{r1}c\end{remarque}
\begin{corollaire}\label{c1}d\end{corollaire}
\chapter{Second}
\begin{corollaire}\label{c2}e\end{corollaire}
\begin{theoreme}\label{t2}f\end{theoreme}
RENVOIS \ref{p1} \ref{t1} \ref{r1} \ref{c1} \ref{c2} \ref{t2} ; \cref{p1} ; \cref{t1} ; \cref{r1} ; \cref{c1} ; \cref{c2} ; \cref{t2} FIN
\end{document}
"                sans correctif, crochet retiré
%   lualatex -jobname=amsseul  "\def\AMSSEUL{1}% Juge de la troisième passe : document minimal écrit de zéro pour la cause A (thmtools sous LaTeX 2026-06-01).
% Variantes, par une macro posée sur la ligne de commande (deux passes chacune, depuis memoire/essai/juge3/minimal/) :
%   lualatex -jobname=nu       "\input{minimal-juge3}"                                   sans correctif
%   lualatex -jobname=corr     "\def\CORR{1}\input{minimal-juge3}"                       correctif A du préambule
%   lualatex -jobname=crochet  "\def\SANSCROCHET{1}\input{minimal-juge3}"                sans correctif, crochet retiré
%   lualatex -jobname=amsseul  "\def\AMSSEUL{1}\input{minimal-juge3}"                    amsthm seul (témoin)
%   lualatex -jobname=faux     "\def\CORR{1}\def\FAUX{1}\input{minimal-juge3}"           deux vrais défauts : doivent tomber
%   lualatex -jobname=trace    "\def\TRACE{1}\input{minimal-juge3}"                      trace de l'appel
% Les lignes du correctif sont tirées du préambule commité par extraire_correctifs.py, pas retapées.
\documentclass{scrbook}
\ifdefined\SANSCROCHET \RemoveFromHook{package/amsthm/after}[firstaid/aliascounter]\fi
\usepackage{amsthm}
\usepackage{hyperref}
\usepackage[french, nameinlink, noabbrev]{cleveref}
\ifdefined\AMSSEUL
  \newtheorem{theoreme}{Théorème}[chapter]
  \newtheorem{proposition}[theoreme]{Proposition}
  \newtheorem{corollaire}[theoreme]{Corollaire}
  \newtheorem{remarque}{Remarque}[chapter]
\else
  \usepackage{thmtools}
  \ifdefined\CORR \input{correctif-a-pose}\fi
  \declaretheorem[numberwithin=chapter, name=Théorème]{theoreme}
  \ifdefined\TRACE \tracingmacros=2 \tracingonline=0 \fi
  \declaretheorem[sibling=theoreme, name=Proposition]{proposition}
  \ifdefined\TRACE \tracingmacros=0 \fi
  \declaretheorem[sibling=theoreme, name=Corollaire]{corollaire}
  \declaretheorem[numberwithin=chapter, name=Remarque]{remarque}
  \ifdefined\FAUX \newcounteralias{corollaire}{section}\newcounteralias{orphelin}{inexistant}\fi
  \ifdefined\CORR \input{correctif-a-remise}\fi
\fi
\crefname{theoreme}{théorème}{théorèmes}
\crefname{proposition}{proposition}{propositions}
\crefname{corollaire}{corollaire}{corollaires}
\crefname{remarque}{remarque}{remarques}
\begin{document}
\chapter{Premier}
\begin{proposition}\label{p1}a\end{proposition}
\begin{theoreme}\label{t1}b\end{theoreme}
\begin{remarque}\label{r1}c\end{remarque}
\begin{corollaire}\label{c1}d\end{corollaire}
\chapter{Second}
\begin{corollaire}\label{c2}e\end{corollaire}
\begin{theoreme}\label{t2}f\end{theoreme}
RENVOIS \ref{p1} \ref{t1} \ref{r1} \ref{c1} \ref{c2} \ref{t2} ; \cref{p1} ; \cref{t1} ; \cref{r1} ; \cref{c1} ; \cref{c2} ; \cref{t2} FIN
\end{document}
"                    amsthm seul (témoin)
%   lualatex -jobname=faux     "\def\CORR{1}\def\FAUX{1}% Juge de la troisième passe : document minimal écrit de zéro pour la cause A (thmtools sous LaTeX 2026-06-01).
% Variantes, par une macro posée sur la ligne de commande (deux passes chacune, depuis memoire/essai/juge3/minimal/) :
%   lualatex -jobname=nu       "\input{minimal-juge3}"                                   sans correctif
%   lualatex -jobname=corr     "\def\CORR{1}\input{minimal-juge3}"                       correctif A du préambule
%   lualatex -jobname=crochet  "\def\SANSCROCHET{1}\input{minimal-juge3}"                sans correctif, crochet retiré
%   lualatex -jobname=amsseul  "\def\AMSSEUL{1}\input{minimal-juge3}"                    amsthm seul (témoin)
%   lualatex -jobname=faux     "\def\CORR{1}\def\FAUX{1}\input{minimal-juge3}"           deux vrais défauts : doivent tomber
%   lualatex -jobname=trace    "\def\TRACE{1}\input{minimal-juge3}"                      trace de l'appel
% Les lignes du correctif sont tirées du préambule commité par extraire_correctifs.py, pas retapées.
\documentclass{scrbook}
\ifdefined\SANSCROCHET \RemoveFromHook{package/amsthm/after}[firstaid/aliascounter]\fi
\usepackage{amsthm}
\usepackage{hyperref}
\usepackage[french, nameinlink, noabbrev]{cleveref}
\ifdefined\AMSSEUL
  \newtheorem{theoreme}{Théorème}[chapter]
  \newtheorem{proposition}[theoreme]{Proposition}
  \newtheorem{corollaire}[theoreme]{Corollaire}
  \newtheorem{remarque}{Remarque}[chapter]
\else
  \usepackage{thmtools}
  \ifdefined\CORR \input{correctif-a-pose}\fi
  \declaretheorem[numberwithin=chapter, name=Théorème]{theoreme}
  \ifdefined\TRACE \tracingmacros=2 \tracingonline=0 \fi
  \declaretheorem[sibling=theoreme, name=Proposition]{proposition}
  \ifdefined\TRACE \tracingmacros=0 \fi
  \declaretheorem[sibling=theoreme, name=Corollaire]{corollaire}
  \declaretheorem[numberwithin=chapter, name=Remarque]{remarque}
  \ifdefined\FAUX \newcounteralias{corollaire}{section}\newcounteralias{orphelin}{inexistant}\fi
  \ifdefined\CORR \input{correctif-a-remise}\fi
\fi
\crefname{theoreme}{théorème}{théorèmes}
\crefname{proposition}{proposition}{propositions}
\crefname{corollaire}{corollaire}{corollaires}
\crefname{remarque}{remarque}{remarques}
\begin{document}
\chapter{Premier}
\begin{proposition}\label{p1}a\end{proposition}
\begin{theoreme}\label{t1}b\end{theoreme}
\begin{remarque}\label{r1}c\end{remarque}
\begin{corollaire}\label{c1}d\end{corollaire}
\chapter{Second}
\begin{corollaire}\label{c2}e\end{corollaire}
\begin{theoreme}\label{t2}f\end{theoreme}
RENVOIS \ref{p1} \ref{t1} \ref{r1} \ref{c1} \ref{c2} \ref{t2} ; \cref{p1} ; \cref{t1} ; \cref{r1} ; \cref{c1} ; \cref{c2} ; \cref{t2} FIN
\end{document}
"           deux vrais défauts : doivent tomber
%   lualatex -jobname=trace    "\def\TRACE{1}% Juge de la troisième passe : document minimal écrit de zéro pour la cause A (thmtools sous LaTeX 2026-06-01).
% Variantes, par une macro posée sur la ligne de commande (deux passes chacune, depuis memoire/essai/juge3/minimal/) :
%   lualatex -jobname=nu       "\input{minimal-juge3}"                                   sans correctif
%   lualatex -jobname=corr     "\def\CORR{1}\input{minimal-juge3}"                       correctif A du préambule
%   lualatex -jobname=crochet  "\def\SANSCROCHET{1}\input{minimal-juge3}"                sans correctif, crochet retiré
%   lualatex -jobname=amsseul  "\def\AMSSEUL{1}\input{minimal-juge3}"                    amsthm seul (témoin)
%   lualatex -jobname=faux     "\def\CORR{1}\def\FAUX{1}\input{minimal-juge3}"           deux vrais défauts : doivent tomber
%   lualatex -jobname=trace    "\def\TRACE{1}\input{minimal-juge3}"                      trace de l'appel
% Les lignes du correctif sont tirées du préambule commité par extraire_correctifs.py, pas retapées.
\documentclass{scrbook}
\ifdefined\SANSCROCHET \RemoveFromHook{package/amsthm/after}[firstaid/aliascounter]\fi
\usepackage{amsthm}
\usepackage{hyperref}
\usepackage[french, nameinlink, noabbrev]{cleveref}
\ifdefined\AMSSEUL
  \newtheorem{theoreme}{Théorème}[chapter]
  \newtheorem{proposition}[theoreme]{Proposition}
  \newtheorem{corollaire}[theoreme]{Corollaire}
  \newtheorem{remarque}{Remarque}[chapter]
\else
  \usepackage{thmtools}
  \ifdefined\CORR \input{correctif-a-pose}\fi
  \declaretheorem[numberwithin=chapter, name=Théorème]{theoreme}
  \ifdefined\TRACE \tracingmacros=2 \tracingonline=0 \fi
  \declaretheorem[sibling=theoreme, name=Proposition]{proposition}
  \ifdefined\TRACE \tracingmacros=0 \fi
  \declaretheorem[sibling=theoreme, name=Corollaire]{corollaire}
  \declaretheorem[numberwithin=chapter, name=Remarque]{remarque}
  \ifdefined\FAUX \newcounteralias{corollaire}{section}\newcounteralias{orphelin}{inexistant}\fi
  \ifdefined\CORR \input{correctif-a-remise}\fi
\fi
\crefname{theoreme}{théorème}{théorèmes}
\crefname{proposition}{proposition}{propositions}
\crefname{corollaire}{corollaire}{corollaires}
\crefname{remarque}{remarque}{remarques}
\begin{document}
\chapter{Premier}
\begin{proposition}\label{p1}a\end{proposition}
\begin{theoreme}\label{t1}b\end{theoreme}
\begin{remarque}\label{r1}c\end{remarque}
\begin{corollaire}\label{c1}d\end{corollaire}
\chapter{Second}
\begin{corollaire}\label{c2}e\end{corollaire}
\begin{theoreme}\label{t2}f\end{theoreme}
RENVOIS \ref{p1} \ref{t1} \ref{r1} \ref{c1} \ref{c2} \ref{t2} ; \cref{p1} ; \cref{t1} ; \cref{r1} ; \cref{c1} ; \cref{c2} ; \cref{t2} FIN
\end{document}
"                      trace de l'appel
% Les lignes du correctif sont tirées du préambule commité par extraire_correctifs.py, pas retapées.
\documentclass{scrbook}
\ifdefined\SANSCROCHET \RemoveFromHook{package/amsthm/after}[firstaid/aliascounter]\fi
\usepackage{amsthm}
\usepackage{hyperref}
\usepackage[french, nameinlink, noabbrev]{cleveref}
\ifdefined\AMSSEUL
  \newtheorem{theoreme}{Théorème}[chapter]
  \newtheorem{proposition}[theoreme]{Proposition}
  \newtheorem{corollaire}[theoreme]{Corollaire}
  \newtheorem{remarque}{Remarque}[chapter]
\else
  \usepackage{thmtools}
  \ifdefined\CORR \input{correctif-a-pose}\fi
  \declaretheorem[numberwithin=chapter, name=Théorème]{theoreme}
  \ifdefined\TRACE \tracingmacros=2 \tracingonline=0 \fi
  \declaretheorem[sibling=theoreme, name=Proposition]{proposition}
  \ifdefined\TRACE \tracingmacros=0 \fi
  \declaretheorem[sibling=theoreme, name=Corollaire]{corollaire}
  \declaretheorem[numberwithin=chapter, name=Remarque]{remarque}
  \ifdefined\FAUX \newcounteralias{corollaire}{section}\newcounteralias{orphelin}{inexistant}\fi
  \ifdefined\CORR \input{correctif-a-remise}\fi
\fi
\crefname{theoreme}{théorème}{théorèmes}
\crefname{proposition}{proposition}{propositions}
\crefname{corollaire}{corollaire}{corollaires}
\crefname{remarque}{remarque}{remarques}
\begin{document}
\chapter{Premier}
\begin{proposition}\label{p1}a\end{proposition}
\begin{theoreme}\label{t1}b\end{theoreme}
\begin{remarque}\label{r1}c\end{remarque}
\begin{corollaire}\label{c1}d\end{corollaire}
\chapter{Second}
\begin{corollaire}\label{c2}e\end{corollaire}
\begin{theoreme}\label{t2}f\end{theoreme}
RENVOIS \ref{p1} \ref{t1} \ref{r1} \ref{c1} \ref{c2} \ref{t2} ; \cref{p1} ; \cref{t1} ; \cref{r1} ; \cref{c1} ; \cref{c2} ; \cref{t2} FIN
\end{document}
"                      trace de l'appel
% Les lignes du correctif sont tirées du préambule commité par extraire_correctifs.py, pas retapées.
\documentclass{scrbook}
\ifdefined\SANSCROCHET \RemoveFromHook{package/amsthm/after}[firstaid/aliascounter]\fi
\usepackage{amsthm}
\usepackage{hyperref}
\usepackage[french, nameinlink, noabbrev]{cleveref}
\ifdefined\AMSSEUL
  \newtheorem{theoreme}{Théorème}[chapter]
  \newtheorem{proposition}[theoreme]{Proposition}
  \newtheorem{corollaire}[theoreme]{Corollaire}
  \newtheorem{remarque}{Remarque}[chapter]
\else
  \usepackage{thmtools}
  \ifdefined\CORR \input{correctif-a-pose}\fi
  \declaretheorem[numberwithin=chapter, name=Théorème]{theoreme}
  \ifdefined\TRACE \tracingmacros=2 \tracingonline=0 \fi
  \declaretheorem[sibling=theoreme, name=Proposition]{proposition}
  \ifdefined\TRACE \tracingmacros=0 \fi
  \declaretheorem[sibling=theoreme, name=Corollaire]{corollaire}
  \declaretheorem[numberwithin=chapter, name=Remarque]{remarque}
  \ifdefined\FAUX \newcounteralias{corollaire}{section}\newcounteralias{orphelin}{inexistant}\fi
  \ifdefined\CORR \input{correctif-a-remise}\fi
\fi
\crefname{theoreme}{théorème}{théorèmes}
\crefname{proposition}{proposition}{propositions}
\crefname{corollaire}{corollaire}{corollaires}
\crefname{remarque}{remarque}{remarques}
\begin{document}
\chapter{Premier}
\begin{proposition}\label{p1}a\end{proposition}
\begin{theoreme}\label{t1}b\end{theoreme}
\begin{remarque}\label{r1}c\end{remarque}
\begin{corollaire}\label{c1}d\end{corollaire}
\chapter{Second}
\begin{corollaire}\label{c2}e\end{corollaire}
\begin{theoreme}\label{t2}f\end{theoreme}
RENVOIS \ref{p1} \ref{t1} \ref{r1} \ref{c1} \ref{c2} \ref{t2} ; \cref{p1} ; \cref{t1} ; \cref{r1} ; \cref{c1} ; \cref{c2} ; \cref{t2} FIN
\end{document}
"                                   sans correctif
%   lualatex -jobname=corr     "\def\CORR{1}% Juge de la troisième passe : document minimal écrit de zéro pour la cause A (thmtools sous LaTeX 2026-06-01).
% Variantes, par une macro posée sur la ligne de commande (deux passes chacune, depuis memoire/essai/juge3/minimal/) :
%   lualatex -jobname=nu       "% Juge de la troisième passe : document minimal écrit de zéro pour la cause A (thmtools sous LaTeX 2026-06-01).
% Variantes, par une macro posée sur la ligne de commande (deux passes chacune, depuis memoire/essai/juge3/minimal/) :
%   lualatex -jobname=nu       "% Juge de la troisième passe : document minimal écrit de zéro pour la cause A (thmtools sous LaTeX 2026-06-01).
% Variantes, par une macro posée sur la ligne de commande (deux passes chacune, depuis memoire/essai/juge3/minimal/) :
%   lualatex -jobname=nu       "\input{minimal-juge3}"                                   sans correctif
%   lualatex -jobname=corr     "\def\CORR{1}\input{minimal-juge3}"                       correctif A du préambule
%   lualatex -jobname=crochet  "\def\SANSCROCHET{1}\input{minimal-juge3}"                sans correctif, crochet retiré
%   lualatex -jobname=amsseul  "\def\AMSSEUL{1}\input{minimal-juge3}"                    amsthm seul (témoin)
%   lualatex -jobname=faux     "\def\CORR{1}\def\FAUX{1}\input{minimal-juge3}"           deux vrais défauts : doivent tomber
%   lualatex -jobname=trace    "\def\TRACE{1}\input{minimal-juge3}"                      trace de l'appel
% Les lignes du correctif sont tirées du préambule commité par extraire_correctifs.py, pas retapées.
\documentclass{scrbook}
\ifdefined\SANSCROCHET \RemoveFromHook{package/amsthm/after}[firstaid/aliascounter]\fi
\usepackage{amsthm}
\usepackage{hyperref}
\usepackage[french, nameinlink, noabbrev]{cleveref}
\ifdefined\AMSSEUL
  \newtheorem{theoreme}{Théorème}[chapter]
  \newtheorem{proposition}[theoreme]{Proposition}
  \newtheorem{corollaire}[theoreme]{Corollaire}
  \newtheorem{remarque}{Remarque}[chapter]
\else
  \usepackage{thmtools}
  \ifdefined\CORR \input{correctif-a-pose}\fi
  \declaretheorem[numberwithin=chapter, name=Théorème]{theoreme}
  \ifdefined\TRACE \tracingmacros=2 \tracingonline=0 \fi
  \declaretheorem[sibling=theoreme, name=Proposition]{proposition}
  \ifdefined\TRACE \tracingmacros=0 \fi
  \declaretheorem[sibling=theoreme, name=Corollaire]{corollaire}
  \declaretheorem[numberwithin=chapter, name=Remarque]{remarque}
  \ifdefined\FAUX \newcounteralias{corollaire}{section}\newcounteralias{orphelin}{inexistant}\fi
  \ifdefined\CORR \input{correctif-a-remise}\fi
\fi
\crefname{theoreme}{théorème}{théorèmes}
\crefname{proposition}{proposition}{propositions}
\crefname{corollaire}{corollaire}{corollaires}
\crefname{remarque}{remarque}{remarques}
\begin{document}
\chapter{Premier}
\begin{proposition}\label{p1}a\end{proposition}
\begin{theoreme}\label{t1}b\end{theoreme}
\begin{remarque}\label{r1}c\end{remarque}
\begin{corollaire}\label{c1}d\end{corollaire}
\chapter{Second}
\begin{corollaire}\label{c2}e\end{corollaire}
\begin{theoreme}\label{t2}f\end{theoreme}
RENVOIS \ref{p1} \ref{t1} \ref{r1} \ref{c1} \ref{c2} \ref{t2} ; \cref{p1} ; \cref{t1} ; \cref{r1} ; \cref{c1} ; \cref{c2} ; \cref{t2} FIN
\end{document}
"                                   sans correctif
%   lualatex -jobname=corr     "\def\CORR{1}% Juge de la troisième passe : document minimal écrit de zéro pour la cause A (thmtools sous LaTeX 2026-06-01).
% Variantes, par une macro posée sur la ligne de commande (deux passes chacune, depuis memoire/essai/juge3/minimal/) :
%   lualatex -jobname=nu       "\input{minimal-juge3}"                                   sans correctif
%   lualatex -jobname=corr     "\def\CORR{1}\input{minimal-juge3}"                       correctif A du préambule
%   lualatex -jobname=crochet  "\def\SANSCROCHET{1}\input{minimal-juge3}"                sans correctif, crochet retiré
%   lualatex -jobname=amsseul  "\def\AMSSEUL{1}\input{minimal-juge3}"                    amsthm seul (témoin)
%   lualatex -jobname=faux     "\def\CORR{1}\def\FAUX{1}\input{minimal-juge3}"           deux vrais défauts : doivent tomber
%   lualatex -jobname=trace    "\def\TRACE{1}\input{minimal-juge3}"                      trace de l'appel
% Les lignes du correctif sont tirées du préambule commité par extraire_correctifs.py, pas retapées.
\documentclass{scrbook}
\ifdefined\SANSCROCHET \RemoveFromHook{package/amsthm/after}[firstaid/aliascounter]\fi
\usepackage{amsthm}
\usepackage{hyperref}
\usepackage[french, nameinlink, noabbrev]{cleveref}
\ifdefined\AMSSEUL
  \newtheorem{theoreme}{Théorème}[chapter]
  \newtheorem{proposition}[theoreme]{Proposition}
  \newtheorem{corollaire}[theoreme]{Corollaire}
  \newtheorem{remarque}{Remarque}[chapter]
\else
  \usepackage{thmtools}
  \ifdefined\CORR \input{correctif-a-pose}\fi
  \declaretheorem[numberwithin=chapter, name=Théorème]{theoreme}
  \ifdefined\TRACE \tracingmacros=2 \tracingonline=0 \fi
  \declaretheorem[sibling=theoreme, name=Proposition]{proposition}
  \ifdefined\TRACE \tracingmacros=0 \fi
  \declaretheorem[sibling=theoreme, name=Corollaire]{corollaire}
  \declaretheorem[numberwithin=chapter, name=Remarque]{remarque}
  \ifdefined\FAUX \newcounteralias{corollaire}{section}\newcounteralias{orphelin}{inexistant}\fi
  \ifdefined\CORR \input{correctif-a-remise}\fi
\fi
\crefname{theoreme}{théorème}{théorèmes}
\crefname{proposition}{proposition}{propositions}
\crefname{corollaire}{corollaire}{corollaires}
\crefname{remarque}{remarque}{remarques}
\begin{document}
\chapter{Premier}
\begin{proposition}\label{p1}a\end{proposition}
\begin{theoreme}\label{t1}b\end{theoreme}
\begin{remarque}\label{r1}c\end{remarque}
\begin{corollaire}\label{c1}d\end{corollaire}
\chapter{Second}
\begin{corollaire}\label{c2}e\end{corollaire}
\begin{theoreme}\label{t2}f\end{theoreme}
RENVOIS \ref{p1} \ref{t1} \ref{r1} \ref{c1} \ref{c2} \ref{t2} ; \cref{p1} ; \cref{t1} ; \cref{r1} ; \cref{c1} ; \cref{c2} ; \cref{t2} FIN
\end{document}
"                       correctif A du préambule
%   lualatex -jobname=crochet  "\def\SANSCROCHET{1}% Juge de la troisième passe : document minimal écrit de zéro pour la cause A (thmtools sous LaTeX 2026-06-01).
% Variantes, par une macro posée sur la ligne de commande (deux passes chacune, depuis memoire/essai/juge3/minimal/) :
%   lualatex -jobname=nu       "\input{minimal-juge3}"                                   sans correctif
%   lualatex -jobname=corr     "\def\CORR{1}\input{minimal-juge3}"                       correctif A du préambule
%   lualatex -jobname=crochet  "\def\SANSCROCHET{1}\input{minimal-juge3}"                sans correctif, crochet retiré
%   lualatex -jobname=amsseul  "\def\AMSSEUL{1}\input{minimal-juge3}"                    amsthm seul (témoin)
%   lualatex -jobname=faux     "\def\CORR{1}\def\FAUX{1}\input{minimal-juge3}"           deux vrais défauts : doivent tomber
%   lualatex -jobname=trace    "\def\TRACE{1}\input{minimal-juge3}"                      trace de l'appel
% Les lignes du correctif sont tirées du préambule commité par extraire_correctifs.py, pas retapées.
\documentclass{scrbook}
\ifdefined\SANSCROCHET \RemoveFromHook{package/amsthm/after}[firstaid/aliascounter]\fi
\usepackage{amsthm}
\usepackage{hyperref}
\usepackage[french, nameinlink, noabbrev]{cleveref}
\ifdefined\AMSSEUL
  \newtheorem{theoreme}{Théorème}[chapter]
  \newtheorem{proposition}[theoreme]{Proposition}
  \newtheorem{corollaire}[theoreme]{Corollaire}
  \newtheorem{remarque}{Remarque}[chapter]
\else
  \usepackage{thmtools}
  \ifdefined\CORR \input{correctif-a-pose}\fi
  \declaretheorem[numberwithin=chapter, name=Théorème]{theoreme}
  \ifdefined\TRACE \tracingmacros=2 \tracingonline=0 \fi
  \declaretheorem[sibling=theoreme, name=Proposition]{proposition}
  \ifdefined\TRACE \tracingmacros=0 \fi
  \declaretheorem[sibling=theoreme, name=Corollaire]{corollaire}
  \declaretheorem[numberwithin=chapter, name=Remarque]{remarque}
  \ifdefined\FAUX \newcounteralias{corollaire}{section}\newcounteralias{orphelin}{inexistant}\fi
  \ifdefined\CORR \input{correctif-a-remise}\fi
\fi
\crefname{theoreme}{théorème}{théorèmes}
\crefname{proposition}{proposition}{propositions}
\crefname{corollaire}{corollaire}{corollaires}
\crefname{remarque}{remarque}{remarques}
\begin{document}
\chapter{Premier}
\begin{proposition}\label{p1}a\end{proposition}
\begin{theoreme}\label{t1}b\end{theoreme}
\begin{remarque}\label{r1}c\end{remarque}
\begin{corollaire}\label{c1}d\end{corollaire}
\chapter{Second}
\begin{corollaire}\label{c2}e\end{corollaire}
\begin{theoreme}\label{t2}f\end{theoreme}
RENVOIS \ref{p1} \ref{t1} \ref{r1} \ref{c1} \ref{c2} \ref{t2} ; \cref{p1} ; \cref{t1} ; \cref{r1} ; \cref{c1} ; \cref{c2} ; \cref{t2} FIN
\end{document}
"                sans correctif, crochet retiré
%   lualatex -jobname=amsseul  "\def\AMSSEUL{1}% Juge de la troisième passe : document minimal écrit de zéro pour la cause A (thmtools sous LaTeX 2026-06-01).
% Variantes, par une macro posée sur la ligne de commande (deux passes chacune, depuis memoire/essai/juge3/minimal/) :
%   lualatex -jobname=nu       "\input{minimal-juge3}"                                   sans correctif
%   lualatex -jobname=corr     "\def\CORR{1}\input{minimal-juge3}"                       correctif A du préambule
%   lualatex -jobname=crochet  "\def\SANSCROCHET{1}\input{minimal-juge3}"                sans correctif, crochet retiré
%   lualatex -jobname=amsseul  "\def\AMSSEUL{1}\input{minimal-juge3}"                    amsthm seul (témoin)
%   lualatex -jobname=faux     "\def\CORR{1}\def\FAUX{1}\input{minimal-juge3}"           deux vrais défauts : doivent tomber
%   lualatex -jobname=trace    "\def\TRACE{1}\input{minimal-juge3}"                      trace de l'appel
% Les lignes du correctif sont tirées du préambule commité par extraire_correctifs.py, pas retapées.
\documentclass{scrbook}
\ifdefined\SANSCROCHET \RemoveFromHook{package/amsthm/after}[firstaid/aliascounter]\fi
\usepackage{amsthm}
\usepackage{hyperref}
\usepackage[french, nameinlink, noabbrev]{cleveref}
\ifdefined\AMSSEUL
  \newtheorem{theoreme}{Théorème}[chapter]
  \newtheorem{proposition}[theoreme]{Proposition}
  \newtheorem{corollaire}[theoreme]{Corollaire}
  \newtheorem{remarque}{Remarque}[chapter]
\else
  \usepackage{thmtools}
  \ifdefined\CORR \input{correctif-a-pose}\fi
  \declaretheorem[numberwithin=chapter, name=Théorème]{theoreme}
  \ifdefined\TRACE \tracingmacros=2 \tracingonline=0 \fi
  \declaretheorem[sibling=theoreme, name=Proposition]{proposition}
  \ifdefined\TRACE \tracingmacros=0 \fi
  \declaretheorem[sibling=theoreme, name=Corollaire]{corollaire}
  \declaretheorem[numberwithin=chapter, name=Remarque]{remarque}
  \ifdefined\FAUX \newcounteralias{corollaire}{section}\newcounteralias{orphelin}{inexistant}\fi
  \ifdefined\CORR \input{correctif-a-remise}\fi
\fi
\crefname{theoreme}{théorème}{théorèmes}
\crefname{proposition}{proposition}{propositions}
\crefname{corollaire}{corollaire}{corollaires}
\crefname{remarque}{remarque}{remarques}
\begin{document}
\chapter{Premier}
\begin{proposition}\label{p1}a\end{proposition}
\begin{theoreme}\label{t1}b\end{theoreme}
\begin{remarque}\label{r1}c\end{remarque}
\begin{corollaire}\label{c1}d\end{corollaire}
\chapter{Second}
\begin{corollaire}\label{c2}e\end{corollaire}
\begin{theoreme}\label{t2}f\end{theoreme}
RENVOIS \ref{p1} \ref{t1} \ref{r1} \ref{c1} \ref{c2} \ref{t2} ; \cref{p1} ; \cref{t1} ; \cref{r1} ; \cref{c1} ; \cref{c2} ; \cref{t2} FIN
\end{document}
"                    amsthm seul (témoin)
%   lualatex -jobname=faux     "\def\CORR{1}\def\FAUX{1}% Juge de la troisième passe : document minimal écrit de zéro pour la cause A (thmtools sous LaTeX 2026-06-01).
% Variantes, par une macro posée sur la ligne de commande (deux passes chacune, depuis memoire/essai/juge3/minimal/) :
%   lualatex -jobname=nu       "\input{minimal-juge3}"                                   sans correctif
%   lualatex -jobname=corr     "\def\CORR{1}\input{minimal-juge3}"                       correctif A du préambule
%   lualatex -jobname=crochet  "\def\SANSCROCHET{1}\input{minimal-juge3}"                sans correctif, crochet retiré
%   lualatex -jobname=amsseul  "\def\AMSSEUL{1}\input{minimal-juge3}"                    amsthm seul (témoin)
%   lualatex -jobname=faux     "\def\CORR{1}\def\FAUX{1}\input{minimal-juge3}"           deux vrais défauts : doivent tomber
%   lualatex -jobname=trace    "\def\TRACE{1}\input{minimal-juge3}"                      trace de l'appel
% Les lignes du correctif sont tirées du préambule commité par extraire_correctifs.py, pas retapées.
\documentclass{scrbook}
\ifdefined\SANSCROCHET \RemoveFromHook{package/amsthm/after}[firstaid/aliascounter]\fi
\usepackage{amsthm}
\usepackage{hyperref}
\usepackage[french, nameinlink, noabbrev]{cleveref}
\ifdefined\AMSSEUL
  \newtheorem{theoreme}{Théorème}[chapter]
  \newtheorem{proposition}[theoreme]{Proposition}
  \newtheorem{corollaire}[theoreme]{Corollaire}
  \newtheorem{remarque}{Remarque}[chapter]
\else
  \usepackage{thmtools}
  \ifdefined\CORR \input{correctif-a-pose}\fi
  \declaretheorem[numberwithin=chapter, name=Théorème]{theoreme}
  \ifdefined\TRACE \tracingmacros=2 \tracingonline=0 \fi
  \declaretheorem[sibling=theoreme, name=Proposition]{proposition}
  \ifdefined\TRACE \tracingmacros=0 \fi
  \declaretheorem[sibling=theoreme, name=Corollaire]{corollaire}
  \declaretheorem[numberwithin=chapter, name=Remarque]{remarque}
  \ifdefined\FAUX \newcounteralias{corollaire}{section}\newcounteralias{orphelin}{inexistant}\fi
  \ifdefined\CORR \input{correctif-a-remise}\fi
\fi
\crefname{theoreme}{théorème}{théorèmes}
\crefname{proposition}{proposition}{propositions}
\crefname{corollaire}{corollaire}{corollaires}
\crefname{remarque}{remarque}{remarques}
\begin{document}
\chapter{Premier}
\begin{proposition}\label{p1}a\end{proposition}
\begin{theoreme}\label{t1}b\end{theoreme}
\begin{remarque}\label{r1}c\end{remarque}
\begin{corollaire}\label{c1}d\end{corollaire}
\chapter{Second}
\begin{corollaire}\label{c2}e\end{corollaire}
\begin{theoreme}\label{t2}f\end{theoreme}
RENVOIS \ref{p1} \ref{t1} \ref{r1} \ref{c1} \ref{c2} \ref{t2} ; \cref{p1} ; \cref{t1} ; \cref{r1} ; \cref{c1} ; \cref{c2} ; \cref{t2} FIN
\end{document}
"           deux vrais défauts : doivent tomber
%   lualatex -jobname=trace    "\def\TRACE{1}% Juge de la troisième passe : document minimal écrit de zéro pour la cause A (thmtools sous LaTeX 2026-06-01).
% Variantes, par une macro posée sur la ligne de commande (deux passes chacune, depuis memoire/essai/juge3/minimal/) :
%   lualatex -jobname=nu       "\input{minimal-juge3}"                                   sans correctif
%   lualatex -jobname=corr     "\def\CORR{1}\input{minimal-juge3}"                       correctif A du préambule
%   lualatex -jobname=crochet  "\def\SANSCROCHET{1}\input{minimal-juge3}"                sans correctif, crochet retiré
%   lualatex -jobname=amsseul  "\def\AMSSEUL{1}\input{minimal-juge3}"                    amsthm seul (témoin)
%   lualatex -jobname=faux     "\def\CORR{1}\def\FAUX{1}\input{minimal-juge3}"           deux vrais défauts : doivent tomber
%   lualatex -jobname=trace    "\def\TRACE{1}\input{minimal-juge3}"                      trace de l'appel
% Les lignes du correctif sont tirées du préambule commité par extraire_correctifs.py, pas retapées.
\documentclass{scrbook}
\ifdefined\SANSCROCHET \RemoveFromHook{package/amsthm/after}[firstaid/aliascounter]\fi
\usepackage{amsthm}
\usepackage{hyperref}
\usepackage[french, nameinlink, noabbrev]{cleveref}
\ifdefined\AMSSEUL
  \newtheorem{theoreme}{Théorème}[chapter]
  \newtheorem{proposition}[theoreme]{Proposition}
  \newtheorem{corollaire}[theoreme]{Corollaire}
  \newtheorem{remarque}{Remarque}[chapter]
\else
  \usepackage{thmtools}
  \ifdefined\CORR \input{correctif-a-pose}\fi
  \declaretheorem[numberwithin=chapter, name=Théorème]{theoreme}
  \ifdefined\TRACE \tracingmacros=2 \tracingonline=0 \fi
  \declaretheorem[sibling=theoreme, name=Proposition]{proposition}
  \ifdefined\TRACE \tracingmacros=0 \fi
  \declaretheorem[sibling=theoreme, name=Corollaire]{corollaire}
  \declaretheorem[numberwithin=chapter, name=Remarque]{remarque}
  \ifdefined\FAUX \newcounteralias{corollaire}{section}\newcounteralias{orphelin}{inexistant}\fi
  \ifdefined\CORR \input{correctif-a-remise}\fi
\fi
\crefname{theoreme}{théorème}{théorèmes}
\crefname{proposition}{proposition}{propositions}
\crefname{corollaire}{corollaire}{corollaires}
\crefname{remarque}{remarque}{remarques}
\begin{document}
\chapter{Premier}
\begin{proposition}\label{p1}a\end{proposition}
\begin{theoreme}\label{t1}b\end{theoreme}
\begin{remarque}\label{r1}c\end{remarque}
\begin{corollaire}\label{c1}d\end{corollaire}
\chapter{Second}
\begin{corollaire}\label{c2}e\end{corollaire}
\begin{theoreme}\label{t2}f\end{theoreme}
RENVOIS \ref{p1} \ref{t1} \ref{r1} \ref{c1} \ref{c2} \ref{t2} ; \cref{p1} ; \cref{t1} ; \cref{r1} ; \cref{c1} ; \cref{c2} ; \cref{t2} FIN
\end{document}
"                      trace de l'appel
% Les lignes du correctif sont tirées du préambule commité par extraire_correctifs.py, pas retapées.
\documentclass{scrbook}
\ifdefined\SANSCROCHET \RemoveFromHook{package/amsthm/after}[firstaid/aliascounter]\fi
\usepackage{amsthm}
\usepackage{hyperref}
\usepackage[french, nameinlink, noabbrev]{cleveref}
\ifdefined\AMSSEUL
  \newtheorem{theoreme}{Théorème}[chapter]
  \newtheorem{proposition}[theoreme]{Proposition}
  \newtheorem{corollaire}[theoreme]{Corollaire}
  \newtheorem{remarque}{Remarque}[chapter]
\else
  \usepackage{thmtools}
  \ifdefined\CORR \input{correctif-a-pose}\fi
  \declaretheorem[numberwithin=chapter, name=Théorème]{theoreme}
  \ifdefined\TRACE \tracingmacros=2 \tracingonline=0 \fi
  \declaretheorem[sibling=theoreme, name=Proposition]{proposition}
  \ifdefined\TRACE \tracingmacros=0 \fi
  \declaretheorem[sibling=theoreme, name=Corollaire]{corollaire}
  \declaretheorem[numberwithin=chapter, name=Remarque]{remarque}
  \ifdefined\FAUX \newcounteralias{corollaire}{section}\newcounteralias{orphelin}{inexistant}\fi
  \ifdefined\CORR \input{correctif-a-remise}\fi
\fi
\crefname{theoreme}{théorème}{théorèmes}
\crefname{proposition}{proposition}{propositions}
\crefname{corollaire}{corollaire}{corollaires}
\crefname{remarque}{remarque}{remarques}
\begin{document}
\chapter{Premier}
\begin{proposition}\label{p1}a\end{proposition}
\begin{theoreme}\label{t1}b\end{theoreme}
\begin{remarque}\label{r1}c\end{remarque}
\begin{corollaire}\label{c1}d\end{corollaire}
\chapter{Second}
\begin{corollaire}\label{c2}e\end{corollaire}
\begin{theoreme}\label{t2}f\end{theoreme}
RENVOIS \ref{p1} \ref{t1} \ref{r1} \ref{c1} \ref{c2} \ref{t2} ; \cref{p1} ; \cref{t1} ; \cref{r1} ; \cref{c1} ; \cref{c2} ; \cref{t2} FIN
\end{document}
"                                   sans correctif
%   lualatex -jobname=corr     "\def\CORR{1}% Juge de la troisième passe : document minimal écrit de zéro pour la cause A (thmtools sous LaTeX 2026-06-01).
% Variantes, par une macro posée sur la ligne de commande (deux passes chacune, depuis memoire/essai/juge3/minimal/) :
%   lualatex -jobname=nu       "% Juge de la troisième passe : document minimal écrit de zéro pour la cause A (thmtools sous LaTeX 2026-06-01).
% Variantes, par une macro posée sur la ligne de commande (deux passes chacune, depuis memoire/essai/juge3/minimal/) :
%   lualatex -jobname=nu       "\input{minimal-juge3}"                                   sans correctif
%   lualatex -jobname=corr     "\def\CORR{1}\input{minimal-juge3}"                       correctif A du préambule
%   lualatex -jobname=crochet  "\def\SANSCROCHET{1}\input{minimal-juge3}"                sans correctif, crochet retiré
%   lualatex -jobname=amsseul  "\def\AMSSEUL{1}\input{minimal-juge3}"                    amsthm seul (témoin)
%   lualatex -jobname=faux     "\def\CORR{1}\def\FAUX{1}\input{minimal-juge3}"           deux vrais défauts : doivent tomber
%   lualatex -jobname=trace    "\def\TRACE{1}\input{minimal-juge3}"                      trace de l'appel
% Les lignes du correctif sont tirées du préambule commité par extraire_correctifs.py, pas retapées.
\documentclass{scrbook}
\ifdefined\SANSCROCHET \RemoveFromHook{package/amsthm/after}[firstaid/aliascounter]\fi
\usepackage{amsthm}
\usepackage{hyperref}
\usepackage[french, nameinlink, noabbrev]{cleveref}
\ifdefined\AMSSEUL
  \newtheorem{theoreme}{Théorème}[chapter]
  \newtheorem{proposition}[theoreme]{Proposition}
  \newtheorem{corollaire}[theoreme]{Corollaire}
  \newtheorem{remarque}{Remarque}[chapter]
\else
  \usepackage{thmtools}
  \ifdefined\CORR \input{correctif-a-pose}\fi
  \declaretheorem[numberwithin=chapter, name=Théorème]{theoreme}
  \ifdefined\TRACE \tracingmacros=2 \tracingonline=0 \fi
  \declaretheorem[sibling=theoreme, name=Proposition]{proposition}
  \ifdefined\TRACE \tracingmacros=0 \fi
  \declaretheorem[sibling=theoreme, name=Corollaire]{corollaire}
  \declaretheorem[numberwithin=chapter, name=Remarque]{remarque}
  \ifdefined\FAUX \newcounteralias{corollaire}{section}\newcounteralias{orphelin}{inexistant}\fi
  \ifdefined\CORR \input{correctif-a-remise}\fi
\fi
\crefname{theoreme}{théorème}{théorèmes}
\crefname{proposition}{proposition}{propositions}
\crefname{corollaire}{corollaire}{corollaires}
\crefname{remarque}{remarque}{remarques}
\begin{document}
\chapter{Premier}
\begin{proposition}\label{p1}a\end{proposition}
\begin{theoreme}\label{t1}b\end{theoreme}
\begin{remarque}\label{r1}c\end{remarque}
\begin{corollaire}\label{c1}d\end{corollaire}
\chapter{Second}
\begin{corollaire}\label{c2}e\end{corollaire}
\begin{theoreme}\label{t2}f\end{theoreme}
RENVOIS \ref{p1} \ref{t1} \ref{r1} \ref{c1} \ref{c2} \ref{t2} ; \cref{p1} ; \cref{t1} ; \cref{r1} ; \cref{c1} ; \cref{c2} ; \cref{t2} FIN
\end{document}
"                                   sans correctif
%   lualatex -jobname=corr     "\def\CORR{1}% Juge de la troisième passe : document minimal écrit de zéro pour la cause A (thmtools sous LaTeX 2026-06-01).
% Variantes, par une macro posée sur la ligne de commande (deux passes chacune, depuis memoire/essai/juge3/minimal/) :
%   lualatex -jobname=nu       "\input{minimal-juge3}"                                   sans correctif
%   lualatex -jobname=corr     "\def\CORR{1}\input{minimal-juge3}"                       correctif A du préambule
%   lualatex -jobname=crochet  "\def\SANSCROCHET{1}\input{minimal-juge3}"                sans correctif, crochet retiré
%   lualatex -jobname=amsseul  "\def\AMSSEUL{1}\input{minimal-juge3}"                    amsthm seul (témoin)
%   lualatex -jobname=faux     "\def\CORR{1}\def\FAUX{1}\input{minimal-juge3}"           deux vrais défauts : doivent tomber
%   lualatex -jobname=trace    "\def\TRACE{1}\input{minimal-juge3}"                      trace de l'appel
% Les lignes du correctif sont tirées du préambule commité par extraire_correctifs.py, pas retapées.
\documentclass{scrbook}
\ifdefined\SANSCROCHET \RemoveFromHook{package/amsthm/after}[firstaid/aliascounter]\fi
\usepackage{amsthm}
\usepackage{hyperref}
\usepackage[french, nameinlink, noabbrev]{cleveref}
\ifdefined\AMSSEUL
  \newtheorem{theoreme}{Théorème}[chapter]
  \newtheorem{proposition}[theoreme]{Proposition}
  \newtheorem{corollaire}[theoreme]{Corollaire}
  \newtheorem{remarque}{Remarque}[chapter]
\else
  \usepackage{thmtools}
  \ifdefined\CORR \input{correctif-a-pose}\fi
  \declaretheorem[numberwithin=chapter, name=Théorème]{theoreme}
  \ifdefined\TRACE \tracingmacros=2 \tracingonline=0 \fi
  \declaretheorem[sibling=theoreme, name=Proposition]{proposition}
  \ifdefined\TRACE \tracingmacros=0 \fi
  \declaretheorem[sibling=theoreme, name=Corollaire]{corollaire}
  \declaretheorem[numberwithin=chapter, name=Remarque]{remarque}
  \ifdefined\FAUX \newcounteralias{corollaire}{section}\newcounteralias{orphelin}{inexistant}\fi
  \ifdefined\CORR \input{correctif-a-remise}\fi
\fi
\crefname{theoreme}{théorème}{théorèmes}
\crefname{proposition}{proposition}{propositions}
\crefname{corollaire}{corollaire}{corollaires}
\crefname{remarque}{remarque}{remarques}
\begin{document}
\chapter{Premier}
\begin{proposition}\label{p1}a\end{proposition}
\begin{theoreme}\label{t1}b\end{theoreme}
\begin{remarque}\label{r1}c\end{remarque}
\begin{corollaire}\label{c1}d\end{corollaire}
\chapter{Second}
\begin{corollaire}\label{c2}e\end{corollaire}
\begin{theoreme}\label{t2}f\end{theoreme}
RENVOIS \ref{p1} \ref{t1} \ref{r1} \ref{c1} \ref{c2} \ref{t2} ; \cref{p1} ; \cref{t1} ; \cref{r1} ; \cref{c1} ; \cref{c2} ; \cref{t2} FIN
\end{document}
"                       correctif A du préambule
%   lualatex -jobname=crochet  "\def\SANSCROCHET{1}% Juge de la troisième passe : document minimal écrit de zéro pour la cause A (thmtools sous LaTeX 2026-06-01).
% Variantes, par une macro posée sur la ligne de commande (deux passes chacune, depuis memoire/essai/juge3/minimal/) :
%   lualatex -jobname=nu       "\input{minimal-juge3}"                                   sans correctif
%   lualatex -jobname=corr     "\def\CORR{1}\input{minimal-juge3}"                       correctif A du préambule
%   lualatex -jobname=crochet  "\def\SANSCROCHET{1}\input{minimal-juge3}"                sans correctif, crochet retiré
%   lualatex -jobname=amsseul  "\def\AMSSEUL{1}\input{minimal-juge3}"                    amsthm seul (témoin)
%   lualatex -jobname=faux     "\def\CORR{1}\def\FAUX{1}\input{minimal-juge3}"           deux vrais défauts : doivent tomber
%   lualatex -jobname=trace    "\def\TRACE{1}\input{minimal-juge3}"                      trace de l'appel
% Les lignes du correctif sont tirées du préambule commité par extraire_correctifs.py, pas retapées.
\documentclass{scrbook}
\ifdefined\SANSCROCHET \RemoveFromHook{package/amsthm/after}[firstaid/aliascounter]\fi
\usepackage{amsthm}
\usepackage{hyperref}
\usepackage[french, nameinlink, noabbrev]{cleveref}
\ifdefined\AMSSEUL
  \newtheorem{theoreme}{Théorème}[chapter]
  \newtheorem{proposition}[theoreme]{Proposition}
  \newtheorem{corollaire}[theoreme]{Corollaire}
  \newtheorem{remarque}{Remarque}[chapter]
\else
  \usepackage{thmtools}
  \ifdefined\CORR \input{correctif-a-pose}\fi
  \declaretheorem[numberwithin=chapter, name=Théorème]{theoreme}
  \ifdefined\TRACE \tracingmacros=2 \tracingonline=0 \fi
  \declaretheorem[sibling=theoreme, name=Proposition]{proposition}
  \ifdefined\TRACE \tracingmacros=0 \fi
  \declaretheorem[sibling=theoreme, name=Corollaire]{corollaire}
  \declaretheorem[numberwithin=chapter, name=Remarque]{remarque}
  \ifdefined\FAUX \newcounteralias{corollaire}{section}\newcounteralias{orphelin}{inexistant}\fi
  \ifdefined\CORR \input{correctif-a-remise}\fi
\fi
\crefname{theoreme}{théorème}{théorèmes}
\crefname{proposition}{proposition}{propositions}
\crefname{corollaire}{corollaire}{corollaires}
\crefname{remarque}{remarque}{remarques}
\begin{document}
\chapter{Premier}
\begin{proposition}\label{p1}a\end{proposition}
\begin{theoreme}\label{t1}b\end{theoreme}
\begin{remarque}\label{r1}c\end{remarque}
\begin{corollaire}\label{c1}d\end{corollaire}
\chapter{Second}
\begin{corollaire}\label{c2}e\end{corollaire}
\begin{theoreme}\label{t2}f\end{theoreme}
RENVOIS \ref{p1} \ref{t1} \ref{r1} \ref{c1} \ref{c2} \ref{t2} ; \cref{p1} ; \cref{t1} ; \cref{r1} ; \cref{c1} ; \cref{c2} ; \cref{t2} FIN
\end{document}
"                sans correctif, crochet retiré
%   lualatex -jobname=amsseul  "\def\AMSSEUL{1}% Juge de la troisième passe : document minimal écrit de zéro pour la cause A (thmtools sous LaTeX 2026-06-01).
% Variantes, par une macro posée sur la ligne de commande (deux passes chacune, depuis memoire/essai/juge3/minimal/) :
%   lualatex -jobname=nu       "\input{minimal-juge3}"                                   sans correctif
%   lualatex -jobname=corr     "\def\CORR{1}\input{minimal-juge3}"                       correctif A du préambule
%   lualatex -jobname=crochet  "\def\SANSCROCHET{1}\input{minimal-juge3}"                sans correctif, crochet retiré
%   lualatex -jobname=amsseul  "\def\AMSSEUL{1}\input{minimal-juge3}"                    amsthm seul (témoin)
%   lualatex -jobname=faux     "\def\CORR{1}\def\FAUX{1}\input{minimal-juge3}"           deux vrais défauts : doivent tomber
%   lualatex -jobname=trace    "\def\TRACE{1}\input{minimal-juge3}"                      trace de l'appel
% Les lignes du correctif sont tirées du préambule commité par extraire_correctifs.py, pas retapées.
\documentclass{scrbook}
\ifdefined\SANSCROCHET \RemoveFromHook{package/amsthm/after}[firstaid/aliascounter]\fi
\usepackage{amsthm}
\usepackage{hyperref}
\usepackage[french, nameinlink, noabbrev]{cleveref}
\ifdefined\AMSSEUL
  \newtheorem{theoreme}{Théorème}[chapter]
  \newtheorem{proposition}[theoreme]{Proposition}
  \newtheorem{corollaire}[theoreme]{Corollaire}
  \newtheorem{remarque}{Remarque}[chapter]
\else
  \usepackage{thmtools}
  \ifdefined\CORR \input{correctif-a-pose}\fi
  \declaretheorem[numberwithin=chapter, name=Théorème]{theoreme}
  \ifdefined\TRACE \tracingmacros=2 \tracingonline=0 \fi
  \declaretheorem[sibling=theoreme, name=Proposition]{proposition}
  \ifdefined\TRACE \tracingmacros=0 \fi
  \declaretheorem[sibling=theoreme, name=Corollaire]{corollaire}
  \declaretheorem[numberwithin=chapter, name=Remarque]{remarque}
  \ifdefined\FAUX \newcounteralias{corollaire}{section}\newcounteralias{orphelin}{inexistant}\fi
  \ifdefined\CORR \input{correctif-a-remise}\fi
\fi
\crefname{theoreme}{théorème}{théorèmes}
\crefname{proposition}{proposition}{propositions}
\crefname{corollaire}{corollaire}{corollaires}
\crefname{remarque}{remarque}{remarques}
\begin{document}
\chapter{Premier}
\begin{proposition}\label{p1}a\end{proposition}
\begin{theoreme}\label{t1}b\end{theoreme}
\begin{remarque}\label{r1}c\end{remarque}
\begin{corollaire}\label{c1}d\end{corollaire}
\chapter{Second}
\begin{corollaire}\label{c2}e\end{corollaire}
\begin{theoreme}\label{t2}f\end{theoreme}
RENVOIS \ref{p1} \ref{t1} \ref{r1} \ref{c1} \ref{c2} \ref{t2} ; \cref{p1} ; \cref{t1} ; \cref{r1} ; \cref{c1} ; \cref{c2} ; \cref{t2} FIN
\end{document}
"                    amsthm seul (témoin)
%   lualatex -jobname=faux     "\def\CORR{1}\def\FAUX{1}% Juge de la troisième passe : document minimal écrit de zéro pour la cause A (thmtools sous LaTeX 2026-06-01).
% Variantes, par une macro posée sur la ligne de commande (deux passes chacune, depuis memoire/essai/juge3/minimal/) :
%   lualatex -jobname=nu       "\input{minimal-juge3}"                                   sans correctif
%   lualatex -jobname=corr     "\def\CORR{1}\input{minimal-juge3}"                       correctif A du préambule
%   lualatex -jobname=crochet  "\def\SANSCROCHET{1}\input{minimal-juge3}"                sans correctif, crochet retiré
%   lualatex -jobname=amsseul  "\def\AMSSEUL{1}\input{minimal-juge3}"                    amsthm seul (témoin)
%   lualatex -jobname=faux     "\def\CORR{1}\def\FAUX{1}\input{minimal-juge3}"           deux vrais défauts : doivent tomber
%   lualatex -jobname=trace    "\def\TRACE{1}\input{minimal-juge3}"                      trace de l'appel
% Les lignes du correctif sont tirées du préambule commité par extraire_correctifs.py, pas retapées.
\documentclass{scrbook}
\ifdefined\SANSCROCHET \RemoveFromHook{package/amsthm/after}[firstaid/aliascounter]\fi
\usepackage{amsthm}
\usepackage{hyperref}
\usepackage[french, nameinlink, noabbrev]{cleveref}
\ifdefined\AMSSEUL
  \newtheorem{theoreme}{Théorème}[chapter]
  \newtheorem{proposition}[theoreme]{Proposition}
  \newtheorem{corollaire}[theoreme]{Corollaire}
  \newtheorem{remarque}{Remarque}[chapter]
\else
  \usepackage{thmtools}
  \ifdefined\CORR \input{correctif-a-pose}\fi
  \declaretheorem[numberwithin=chapter, name=Théorème]{theoreme}
  \ifdefined\TRACE \tracingmacros=2 \tracingonline=0 \fi
  \declaretheorem[sibling=theoreme, name=Proposition]{proposition}
  \ifdefined\TRACE \tracingmacros=0 \fi
  \declaretheorem[sibling=theoreme, name=Corollaire]{corollaire}
  \declaretheorem[numberwithin=chapter, name=Remarque]{remarque}
  \ifdefined\FAUX \newcounteralias{corollaire}{section}\newcounteralias{orphelin}{inexistant}\fi
  \ifdefined\CORR \input{correctif-a-remise}\fi
\fi
\crefname{theoreme}{théorème}{théorèmes}
\crefname{proposition}{proposition}{propositions}
\crefname{corollaire}{corollaire}{corollaires}
\crefname{remarque}{remarque}{remarques}
\begin{document}
\chapter{Premier}
\begin{proposition}\label{p1}a\end{proposition}
\begin{theoreme}\label{t1}b\end{theoreme}
\begin{remarque}\label{r1}c\end{remarque}
\begin{corollaire}\label{c1}d\end{corollaire}
\chapter{Second}
\begin{corollaire}\label{c2}e\end{corollaire}
\begin{theoreme}\label{t2}f\end{theoreme}
RENVOIS \ref{p1} \ref{t1} \ref{r1} \ref{c1} \ref{c2} \ref{t2} ; \cref{p1} ; \cref{t1} ; \cref{r1} ; \cref{c1} ; \cref{c2} ; \cref{t2} FIN
\end{document}
"           deux vrais défauts : doivent tomber
%   lualatex -jobname=trace    "\def\TRACE{1}% Juge de la troisième passe : document minimal écrit de zéro pour la cause A (thmtools sous LaTeX 2026-06-01).
% Variantes, par une macro posée sur la ligne de commande (deux passes chacune, depuis memoire/essai/juge3/minimal/) :
%   lualatex -jobname=nu       "\input{minimal-juge3}"                                   sans correctif
%   lualatex -jobname=corr     "\def\CORR{1}\input{minimal-juge3}"                       correctif A du préambule
%   lualatex -jobname=crochet  "\def\SANSCROCHET{1}\input{minimal-juge3}"                sans correctif, crochet retiré
%   lualatex -jobname=amsseul  "\def\AMSSEUL{1}\input{minimal-juge3}"                    amsthm seul (témoin)
%   lualatex -jobname=faux     "\def\CORR{1}\def\FAUX{1}\input{minimal-juge3}"           deux vrais défauts : doivent tomber
%   lualatex -jobname=trace    "\def\TRACE{1}\input{minimal-juge3}"                      trace de l'appel
% Les lignes du correctif sont tirées du préambule commité par extraire_correctifs.py, pas retapées.
\documentclass{scrbook}
\ifdefined\SANSCROCHET \RemoveFromHook{package/amsthm/after}[firstaid/aliascounter]\fi
\usepackage{amsthm}
\usepackage{hyperref}
\usepackage[french, nameinlink, noabbrev]{cleveref}
\ifdefined\AMSSEUL
  \newtheorem{theoreme}{Théorème}[chapter]
  \newtheorem{proposition}[theoreme]{Proposition}
  \newtheorem{corollaire}[theoreme]{Corollaire}
  \newtheorem{remarque}{Remarque}[chapter]
\else
  \usepackage{thmtools}
  \ifdefined\CORR \input{correctif-a-pose}\fi
  \declaretheorem[numberwithin=chapter, name=Théorème]{theoreme}
  \ifdefined\TRACE \tracingmacros=2 \tracingonline=0 \fi
  \declaretheorem[sibling=theoreme, name=Proposition]{proposition}
  \ifdefined\TRACE \tracingmacros=0 \fi
  \declaretheorem[sibling=theoreme, name=Corollaire]{corollaire}
  \declaretheorem[numberwithin=chapter, name=Remarque]{remarque}
  \ifdefined\FAUX \newcounteralias{corollaire}{section}\newcounteralias{orphelin}{inexistant}\fi
  \ifdefined\CORR \input{correctif-a-remise}\fi
\fi
\crefname{theoreme}{théorème}{théorèmes}
\crefname{proposition}{proposition}{propositions}
\crefname{corollaire}{corollaire}{corollaires}
\crefname{remarque}{remarque}{remarques}
\begin{document}
\chapter{Premier}
\begin{proposition}\label{p1}a\end{proposition}
\begin{theoreme}\label{t1}b\end{theoreme}
\begin{remarque}\label{r1}c\end{remarque}
\begin{corollaire}\label{c1}d\end{corollaire}
\chapter{Second}
\begin{corollaire}\label{c2}e\end{corollaire}
\begin{theoreme}\label{t2}f\end{theoreme}
RENVOIS \ref{p1} \ref{t1} \ref{r1} \ref{c1} \ref{c2} \ref{t2} ; \cref{p1} ; \cref{t1} ; \cref{r1} ; \cref{c1} ; \cref{c2} ; \cref{t2} FIN
\end{document}
"                      trace de l'appel
% Les lignes du correctif sont tirées du préambule commité par extraire_correctifs.py, pas retapées.
\documentclass{scrbook}
\ifdefined\SANSCROCHET \RemoveFromHook{package/amsthm/after}[firstaid/aliascounter]\fi
\usepackage{amsthm}
\usepackage{hyperref}
\usepackage[french, nameinlink, noabbrev]{cleveref}
\ifdefined\AMSSEUL
  \newtheorem{theoreme}{Théorème}[chapter]
  \newtheorem{proposition}[theoreme]{Proposition}
  \newtheorem{corollaire}[theoreme]{Corollaire}
  \newtheorem{remarque}{Remarque}[chapter]
\else
  \usepackage{thmtools}
  \ifdefined\CORR \input{correctif-a-pose}\fi
  \declaretheorem[numberwithin=chapter, name=Théorème]{theoreme}
  \ifdefined\TRACE \tracingmacros=2 \tracingonline=0 \fi
  \declaretheorem[sibling=theoreme, name=Proposition]{proposition}
  \ifdefined\TRACE \tracingmacros=0 \fi
  \declaretheorem[sibling=theoreme, name=Corollaire]{corollaire}
  \declaretheorem[numberwithin=chapter, name=Remarque]{remarque}
  \ifdefined\FAUX \newcounteralias{corollaire}{section}\newcounteralias{orphelin}{inexistant}\fi
  \ifdefined\CORR \input{correctif-a-remise}\fi
\fi
\crefname{theoreme}{théorème}{théorèmes}
\crefname{proposition}{proposition}{propositions}
\crefname{corollaire}{corollaire}{corollaires}
\crefname{remarque}{remarque}{remarques}
\begin{document}
\chapter{Premier}
\begin{proposition}\label{p1}a\end{proposition}
\begin{theoreme}\label{t1}b\end{theoreme}
\begin{remarque}\label{r1}c\end{remarque}
\begin{corollaire}\label{c1}d\end{corollaire}
\chapter{Second}
\begin{corollaire}\label{c2}e\end{corollaire}
\begin{theoreme}\label{t2}f\end{theoreme}
RENVOIS \ref{p1} \ref{t1} \ref{r1} \ref{c1} \ref{c2} \ref{t2} ; \cref{p1} ; \cref{t1} ; \cref{r1} ; \cref{c1} ; \cref{c2} ; \cref{t2} FIN
\end{document}
"                       correctif A du préambule
%   lualatex -jobname=crochet  "\def\SANSCROCHET{1}% Juge de la troisième passe : document minimal écrit de zéro pour la cause A (thmtools sous LaTeX 2026-06-01).
% Variantes, par une macro posée sur la ligne de commande (deux passes chacune, depuis memoire/essai/juge3/minimal/) :
%   lualatex -jobname=nu       "% Juge de la troisième passe : document minimal écrit de zéro pour la cause A (thmtools sous LaTeX 2026-06-01).
% Variantes, par une macro posée sur la ligne de commande (deux passes chacune, depuis memoire/essai/juge3/minimal/) :
%   lualatex -jobname=nu       "\input{minimal-juge3}"                                   sans correctif
%   lualatex -jobname=corr     "\def\CORR{1}\input{minimal-juge3}"                       correctif A du préambule
%   lualatex -jobname=crochet  "\def\SANSCROCHET{1}\input{minimal-juge3}"                sans correctif, crochet retiré
%   lualatex -jobname=amsseul  "\def\AMSSEUL{1}\input{minimal-juge3}"                    amsthm seul (témoin)
%   lualatex -jobname=faux     "\def\CORR{1}\def\FAUX{1}\input{minimal-juge3}"           deux vrais défauts : doivent tomber
%   lualatex -jobname=trace    "\def\TRACE{1}\input{minimal-juge3}"                      trace de l'appel
% Les lignes du correctif sont tirées du préambule commité par extraire_correctifs.py, pas retapées.
\documentclass{scrbook}
\ifdefined\SANSCROCHET \RemoveFromHook{package/amsthm/after}[firstaid/aliascounter]\fi
\usepackage{amsthm}
\usepackage{hyperref}
\usepackage[french, nameinlink, noabbrev]{cleveref}
\ifdefined\AMSSEUL
  \newtheorem{theoreme}{Théorème}[chapter]
  \newtheorem{proposition}[theoreme]{Proposition}
  \newtheorem{corollaire}[theoreme]{Corollaire}
  \newtheorem{remarque}{Remarque}[chapter]
\else
  \usepackage{thmtools}
  \ifdefined\CORR \input{correctif-a-pose}\fi
  \declaretheorem[numberwithin=chapter, name=Théorème]{theoreme}
  \ifdefined\TRACE \tracingmacros=2 \tracingonline=0 \fi
  \declaretheorem[sibling=theoreme, name=Proposition]{proposition}
  \ifdefined\TRACE \tracingmacros=0 \fi
  \declaretheorem[sibling=theoreme, name=Corollaire]{corollaire}
  \declaretheorem[numberwithin=chapter, name=Remarque]{remarque}
  \ifdefined\FAUX \newcounteralias{corollaire}{section}\newcounteralias{orphelin}{inexistant}\fi
  \ifdefined\CORR \input{correctif-a-remise}\fi
\fi
\crefname{theoreme}{théorème}{théorèmes}
\crefname{proposition}{proposition}{propositions}
\crefname{corollaire}{corollaire}{corollaires}
\crefname{remarque}{remarque}{remarques}
\begin{document}
\chapter{Premier}
\begin{proposition}\label{p1}a\end{proposition}
\begin{theoreme}\label{t1}b\end{theoreme}
\begin{remarque}\label{r1}c\end{remarque}
\begin{corollaire}\label{c1}d\end{corollaire}
\chapter{Second}
\begin{corollaire}\label{c2}e\end{corollaire}
\begin{theoreme}\label{t2}f\end{theoreme}
RENVOIS \ref{p1} \ref{t1} \ref{r1} \ref{c1} \ref{c2} \ref{t2} ; \cref{p1} ; \cref{t1} ; \cref{r1} ; \cref{c1} ; \cref{c2} ; \cref{t2} FIN
\end{document}
"                                   sans correctif
%   lualatex -jobname=corr     "\def\CORR{1}% Juge de la troisième passe : document minimal écrit de zéro pour la cause A (thmtools sous LaTeX 2026-06-01).
% Variantes, par une macro posée sur la ligne de commande (deux passes chacune, depuis memoire/essai/juge3/minimal/) :
%   lualatex -jobname=nu       "\input{minimal-juge3}"                                   sans correctif
%   lualatex -jobname=corr     "\def\CORR{1}\input{minimal-juge3}"                       correctif A du préambule
%   lualatex -jobname=crochet  "\def\SANSCROCHET{1}\input{minimal-juge3}"                sans correctif, crochet retiré
%   lualatex -jobname=amsseul  "\def\AMSSEUL{1}\input{minimal-juge3}"                    amsthm seul (témoin)
%   lualatex -jobname=faux     "\def\CORR{1}\def\FAUX{1}\input{minimal-juge3}"           deux vrais défauts : doivent tomber
%   lualatex -jobname=trace    "\def\TRACE{1}\input{minimal-juge3}"                      trace de l'appel
% Les lignes du correctif sont tirées du préambule commité par extraire_correctifs.py, pas retapées.
\documentclass{scrbook}
\ifdefined\SANSCROCHET \RemoveFromHook{package/amsthm/after}[firstaid/aliascounter]\fi
\usepackage{amsthm}
\usepackage{hyperref}
\usepackage[french, nameinlink, noabbrev]{cleveref}
\ifdefined\AMSSEUL
  \newtheorem{theoreme}{Théorème}[chapter]
  \newtheorem{proposition}[theoreme]{Proposition}
  \newtheorem{corollaire}[theoreme]{Corollaire}
  \newtheorem{remarque}{Remarque}[chapter]
\else
  \usepackage{thmtools}
  \ifdefined\CORR \input{correctif-a-pose}\fi
  \declaretheorem[numberwithin=chapter, name=Théorème]{theoreme}
  \ifdefined\TRACE \tracingmacros=2 \tracingonline=0 \fi
  \declaretheorem[sibling=theoreme, name=Proposition]{proposition}
  \ifdefined\TRACE \tracingmacros=0 \fi
  \declaretheorem[sibling=theoreme, name=Corollaire]{corollaire}
  \declaretheorem[numberwithin=chapter, name=Remarque]{remarque}
  \ifdefined\FAUX \newcounteralias{corollaire}{section}\newcounteralias{orphelin}{inexistant}\fi
  \ifdefined\CORR \input{correctif-a-remise}\fi
\fi
\crefname{theoreme}{théorème}{théorèmes}
\crefname{proposition}{proposition}{propositions}
\crefname{corollaire}{corollaire}{corollaires}
\crefname{remarque}{remarque}{remarques}
\begin{document}
\chapter{Premier}
\begin{proposition}\label{p1}a\end{proposition}
\begin{theoreme}\label{t1}b\end{theoreme}
\begin{remarque}\label{r1}c\end{remarque}
\begin{corollaire}\label{c1}d\end{corollaire}
\chapter{Second}
\begin{corollaire}\label{c2}e\end{corollaire}
\begin{theoreme}\label{t2}f\end{theoreme}
RENVOIS \ref{p1} \ref{t1} \ref{r1} \ref{c1} \ref{c2} \ref{t2} ; \cref{p1} ; \cref{t1} ; \cref{r1} ; \cref{c1} ; \cref{c2} ; \cref{t2} FIN
\end{document}
"                       correctif A du préambule
%   lualatex -jobname=crochet  "\def\SANSCROCHET{1}% Juge de la troisième passe : document minimal écrit de zéro pour la cause A (thmtools sous LaTeX 2026-06-01).
% Variantes, par une macro posée sur la ligne de commande (deux passes chacune, depuis memoire/essai/juge3/minimal/) :
%   lualatex -jobname=nu       "\input{minimal-juge3}"                                   sans correctif
%   lualatex -jobname=corr     "\def\CORR{1}\input{minimal-juge3}"                       correctif A du préambule
%   lualatex -jobname=crochet  "\def\SANSCROCHET{1}\input{minimal-juge3}"                sans correctif, crochet retiré
%   lualatex -jobname=amsseul  "\def\AMSSEUL{1}\input{minimal-juge3}"                    amsthm seul (témoin)
%   lualatex -jobname=faux     "\def\CORR{1}\def\FAUX{1}\input{minimal-juge3}"           deux vrais défauts : doivent tomber
%   lualatex -jobname=trace    "\def\TRACE{1}\input{minimal-juge3}"                      trace de l'appel
% Les lignes du correctif sont tirées du préambule commité par extraire_correctifs.py, pas retapées.
\documentclass{scrbook}
\ifdefined\SANSCROCHET \RemoveFromHook{package/amsthm/after}[firstaid/aliascounter]\fi
\usepackage{amsthm}
\usepackage{hyperref}
\usepackage[french, nameinlink, noabbrev]{cleveref}
\ifdefined\AMSSEUL
  \newtheorem{theoreme}{Théorème}[chapter]
  \newtheorem{proposition}[theoreme]{Proposition}
  \newtheorem{corollaire}[theoreme]{Corollaire}
  \newtheorem{remarque}{Remarque}[chapter]
\else
  \usepackage{thmtools}
  \ifdefined\CORR \input{correctif-a-pose}\fi
  \declaretheorem[numberwithin=chapter, name=Théorème]{theoreme}
  \ifdefined\TRACE \tracingmacros=2 \tracingonline=0 \fi
  \declaretheorem[sibling=theoreme, name=Proposition]{proposition}
  \ifdefined\TRACE \tracingmacros=0 \fi
  \declaretheorem[sibling=theoreme, name=Corollaire]{corollaire}
  \declaretheorem[numberwithin=chapter, name=Remarque]{remarque}
  \ifdefined\FAUX \newcounteralias{corollaire}{section}\newcounteralias{orphelin}{inexistant}\fi
  \ifdefined\CORR \input{correctif-a-remise}\fi
\fi
\crefname{theoreme}{théorème}{théorèmes}
\crefname{proposition}{proposition}{propositions}
\crefname{corollaire}{corollaire}{corollaires}
\crefname{remarque}{remarque}{remarques}
\begin{document}
\chapter{Premier}
\begin{proposition}\label{p1}a\end{proposition}
\begin{theoreme}\label{t1}b\end{theoreme}
\begin{remarque}\label{r1}c\end{remarque}
\begin{corollaire}\label{c1}d\end{corollaire}
\chapter{Second}
\begin{corollaire}\label{c2}e\end{corollaire}
\begin{theoreme}\label{t2}f\end{theoreme}
RENVOIS \ref{p1} \ref{t1} \ref{r1} \ref{c1} \ref{c2} \ref{t2} ; \cref{p1} ; \cref{t1} ; \cref{r1} ; \cref{c1} ; \cref{c2} ; \cref{t2} FIN
\end{document}
"                sans correctif, crochet retiré
%   lualatex -jobname=amsseul  "\def\AMSSEUL{1}% Juge de la troisième passe : document minimal écrit de zéro pour la cause A (thmtools sous LaTeX 2026-06-01).
% Variantes, par une macro posée sur la ligne de commande (deux passes chacune, depuis memoire/essai/juge3/minimal/) :
%   lualatex -jobname=nu       "\input{minimal-juge3}"                                   sans correctif
%   lualatex -jobname=corr     "\def\CORR{1}\input{minimal-juge3}"                       correctif A du préambule
%   lualatex -jobname=crochet  "\def\SANSCROCHET{1}\input{minimal-juge3}"                sans correctif, crochet retiré
%   lualatex -jobname=amsseul  "\def\AMSSEUL{1}\input{minimal-juge3}"                    amsthm seul (témoin)
%   lualatex -jobname=faux     "\def\CORR{1}\def\FAUX{1}\input{minimal-juge3}"           deux vrais défauts : doivent tomber
%   lualatex -jobname=trace    "\def\TRACE{1}\input{minimal-juge3}"                      trace de l'appel
% Les lignes du correctif sont tirées du préambule commité par extraire_correctifs.py, pas retapées.
\documentclass{scrbook}
\ifdefined\SANSCROCHET \RemoveFromHook{package/amsthm/after}[firstaid/aliascounter]\fi
\usepackage{amsthm}
\usepackage{hyperref}
\usepackage[french, nameinlink, noabbrev]{cleveref}
\ifdefined\AMSSEUL
  \newtheorem{theoreme}{Théorème}[chapter]
  \newtheorem{proposition}[theoreme]{Proposition}
  \newtheorem{corollaire}[theoreme]{Corollaire}
  \newtheorem{remarque}{Remarque}[chapter]
\else
  \usepackage{thmtools}
  \ifdefined\CORR \input{correctif-a-pose}\fi
  \declaretheorem[numberwithin=chapter, name=Théorème]{theoreme}
  \ifdefined\TRACE \tracingmacros=2 \tracingonline=0 \fi
  \declaretheorem[sibling=theoreme, name=Proposition]{proposition}
  \ifdefined\TRACE \tracingmacros=0 \fi
  \declaretheorem[sibling=theoreme, name=Corollaire]{corollaire}
  \declaretheorem[numberwithin=chapter, name=Remarque]{remarque}
  \ifdefined\FAUX \newcounteralias{corollaire}{section}\newcounteralias{orphelin}{inexistant}\fi
  \ifdefined\CORR \input{correctif-a-remise}\fi
\fi
\crefname{theoreme}{théorème}{théorèmes}
\crefname{proposition}{proposition}{propositions}
\crefname{corollaire}{corollaire}{corollaires}
\crefname{remarque}{remarque}{remarques}
\begin{document}
\chapter{Premier}
\begin{proposition}\label{p1}a\end{proposition}
\begin{theoreme}\label{t1}b\end{theoreme}
\begin{remarque}\label{r1}c\end{remarque}
\begin{corollaire}\label{c1}d\end{corollaire}
\chapter{Second}
\begin{corollaire}\label{c2}e\end{corollaire}
\begin{theoreme}\label{t2}f\end{theoreme}
RENVOIS \ref{p1} \ref{t1} \ref{r1} \ref{c1} \ref{c2} \ref{t2} ; \cref{p1} ; \cref{t1} ; \cref{r1} ; \cref{c1} ; \cref{c2} ; \cref{t2} FIN
\end{document}
"                    amsthm seul (témoin)
%   lualatex -jobname=faux     "\def\CORR{1}\def\FAUX{1}% Juge de la troisième passe : document minimal écrit de zéro pour la cause A (thmtools sous LaTeX 2026-06-01).
% Variantes, par une macro posée sur la ligne de commande (deux passes chacune, depuis memoire/essai/juge3/minimal/) :
%   lualatex -jobname=nu       "\input{minimal-juge3}"                                   sans correctif
%   lualatex -jobname=corr     "\def\CORR{1}\input{minimal-juge3}"                       correctif A du préambule
%   lualatex -jobname=crochet  "\def\SANSCROCHET{1}\input{minimal-juge3}"                sans correctif, crochet retiré
%   lualatex -jobname=amsseul  "\def\AMSSEUL{1}\input{minimal-juge3}"                    amsthm seul (témoin)
%   lualatex -jobname=faux     "\def\CORR{1}\def\FAUX{1}\input{minimal-juge3}"           deux vrais défauts : doivent tomber
%   lualatex -jobname=trace    "\def\TRACE{1}\input{minimal-juge3}"                      trace de l'appel
% Les lignes du correctif sont tirées du préambule commité par extraire_correctifs.py, pas retapées.
\documentclass{scrbook}
\ifdefined\SANSCROCHET \RemoveFromHook{package/amsthm/after}[firstaid/aliascounter]\fi
\usepackage{amsthm}
\usepackage{hyperref}
\usepackage[french, nameinlink, noabbrev]{cleveref}
\ifdefined\AMSSEUL
  \newtheorem{theoreme}{Théorème}[chapter]
  \newtheorem{proposition}[theoreme]{Proposition}
  \newtheorem{corollaire}[theoreme]{Corollaire}
  \newtheorem{remarque}{Remarque}[chapter]
\else
  \usepackage{thmtools}
  \ifdefined\CORR \input{correctif-a-pose}\fi
  \declaretheorem[numberwithin=chapter, name=Théorème]{theoreme}
  \ifdefined\TRACE \tracingmacros=2 \tracingonline=0 \fi
  \declaretheorem[sibling=theoreme, name=Proposition]{proposition}
  \ifdefined\TRACE \tracingmacros=0 \fi
  \declaretheorem[sibling=theoreme, name=Corollaire]{corollaire}
  \declaretheorem[numberwithin=chapter, name=Remarque]{remarque}
  \ifdefined\FAUX \newcounteralias{corollaire}{section}\newcounteralias{orphelin}{inexistant}\fi
  \ifdefined\CORR \input{correctif-a-remise}\fi
\fi
\crefname{theoreme}{théorème}{théorèmes}
\crefname{proposition}{proposition}{propositions}
\crefname{corollaire}{corollaire}{corollaires}
\crefname{remarque}{remarque}{remarques}
\begin{document}
\chapter{Premier}
\begin{proposition}\label{p1}a\end{proposition}
\begin{theoreme}\label{t1}b\end{theoreme}
\begin{remarque}\label{r1}c\end{remarque}
\begin{corollaire}\label{c1}d\end{corollaire}
\chapter{Second}
\begin{corollaire}\label{c2}e\end{corollaire}
\begin{theoreme}\label{t2}f\end{theoreme}
RENVOIS \ref{p1} \ref{t1} \ref{r1} \ref{c1} \ref{c2} \ref{t2} ; \cref{p1} ; \cref{t1} ; \cref{r1} ; \cref{c1} ; \cref{c2} ; \cref{t2} FIN
\end{document}
"           deux vrais défauts : doivent tomber
%   lualatex -jobname=trace    "\def\TRACE{1}% Juge de la troisième passe : document minimal écrit de zéro pour la cause A (thmtools sous LaTeX 2026-06-01).
% Variantes, par une macro posée sur la ligne de commande (deux passes chacune, depuis memoire/essai/juge3/minimal/) :
%   lualatex -jobname=nu       "\input{minimal-juge3}"                                   sans correctif
%   lualatex -jobname=corr     "\def\CORR{1}\input{minimal-juge3}"                       correctif A du préambule
%   lualatex -jobname=crochet  "\def\SANSCROCHET{1}\input{minimal-juge3}"                sans correctif, crochet retiré
%   lualatex -jobname=amsseul  "\def\AMSSEUL{1}\input{minimal-juge3}"                    amsthm seul (témoin)
%   lualatex -jobname=faux     "\def\CORR{1}\def\FAUX{1}\input{minimal-juge3}"           deux vrais défauts : doivent tomber
%   lualatex -jobname=trace    "\def\TRACE{1}\input{minimal-juge3}"                      trace de l'appel
% Les lignes du correctif sont tirées du préambule commité par extraire_correctifs.py, pas retapées.
\documentclass{scrbook}
\ifdefined\SANSCROCHET \RemoveFromHook{package/amsthm/after}[firstaid/aliascounter]\fi
\usepackage{amsthm}
\usepackage{hyperref}
\usepackage[french, nameinlink, noabbrev]{cleveref}
\ifdefined\AMSSEUL
  \newtheorem{theoreme}{Théorème}[chapter]
  \newtheorem{proposition}[theoreme]{Proposition}
  \newtheorem{corollaire}[theoreme]{Corollaire}
  \newtheorem{remarque}{Remarque}[chapter]
\else
  \usepackage{thmtools}
  \ifdefined\CORR \input{correctif-a-pose}\fi
  \declaretheorem[numberwithin=chapter, name=Théorème]{theoreme}
  \ifdefined\TRACE \tracingmacros=2 \tracingonline=0 \fi
  \declaretheorem[sibling=theoreme, name=Proposition]{proposition}
  \ifdefined\TRACE \tracingmacros=0 \fi
  \declaretheorem[sibling=theoreme, name=Corollaire]{corollaire}
  \declaretheorem[numberwithin=chapter, name=Remarque]{remarque}
  \ifdefined\FAUX \newcounteralias{corollaire}{section}\newcounteralias{orphelin}{inexistant}\fi
  \ifdefined\CORR \input{correctif-a-remise}\fi
\fi
\crefname{theoreme}{théorème}{théorèmes}
\crefname{proposition}{proposition}{propositions}
\crefname{corollaire}{corollaire}{corollaires}
\crefname{remarque}{remarque}{remarques}
\begin{document}
\chapter{Premier}
\begin{proposition}\label{p1}a\end{proposition}
\begin{theoreme}\label{t1}b\end{theoreme}
\begin{remarque}\label{r1}c\end{remarque}
\begin{corollaire}\label{c1}d\end{corollaire}
\chapter{Second}
\begin{corollaire}\label{c2}e\end{corollaire}
\begin{theoreme}\label{t2}f\end{theoreme}
RENVOIS \ref{p1} \ref{t1} \ref{r1} \ref{c1} \ref{c2} \ref{t2} ; \cref{p1} ; \cref{t1} ; \cref{r1} ; \cref{c1} ; \cref{c2} ; \cref{t2} FIN
\end{document}
"                      trace de l'appel
% Les lignes du correctif sont tirées du préambule commité par extraire_correctifs.py, pas retapées.
\documentclass{scrbook}
\ifdefined\SANSCROCHET \RemoveFromHook{package/amsthm/after}[firstaid/aliascounter]\fi
\usepackage{amsthm}
\usepackage{hyperref}
\usepackage[french, nameinlink, noabbrev]{cleveref}
\ifdefined\AMSSEUL
  \newtheorem{theoreme}{Théorème}[chapter]
  \newtheorem{proposition}[theoreme]{Proposition}
  \newtheorem{corollaire}[theoreme]{Corollaire}
  \newtheorem{remarque}{Remarque}[chapter]
\else
  \usepackage{thmtools}
  \ifdefined\CORR \input{correctif-a-pose}\fi
  \declaretheorem[numberwithin=chapter, name=Théorème]{theoreme}
  \ifdefined\TRACE \tracingmacros=2 \tracingonline=0 \fi
  \declaretheorem[sibling=theoreme, name=Proposition]{proposition}
  \ifdefined\TRACE \tracingmacros=0 \fi
  \declaretheorem[sibling=theoreme, name=Corollaire]{corollaire}
  \declaretheorem[numberwithin=chapter, name=Remarque]{remarque}
  \ifdefined\FAUX \newcounteralias{corollaire}{section}\newcounteralias{orphelin}{inexistant}\fi
  \ifdefined\CORR \input{correctif-a-remise}\fi
\fi
\crefname{theoreme}{théorème}{théorèmes}
\crefname{proposition}{proposition}{propositions}
\crefname{corollaire}{corollaire}{corollaires}
\crefname{remarque}{remarque}{remarques}
\begin{document}
\chapter{Premier}
\begin{proposition}\label{p1}a\end{proposition}
\begin{theoreme}\label{t1}b\end{theoreme}
\begin{remarque}\label{r1}c\end{remarque}
\begin{corollaire}\label{c1}d\end{corollaire}
\chapter{Second}
\begin{corollaire}\label{c2}e\end{corollaire}
\begin{theoreme}\label{t2}f\end{theoreme}
RENVOIS \ref{p1} \ref{t1} \ref{r1} \ref{c1} \ref{c2} \ref{t2} ; \cref{p1} ; \cref{t1} ; \cref{r1} ; \cref{c1} ; \cref{c2} ; \cref{t2} FIN
\end{document}
"                sans correctif, crochet retiré
%   lualatex -jobname=amsseul  "\def\AMSSEUL{1}% Juge de la troisième passe : document minimal écrit de zéro pour la cause A (thmtools sous LaTeX 2026-06-01).
% Variantes, par une macro posée sur la ligne de commande (deux passes chacune, depuis memoire/essai/juge3/minimal/) :
%   lualatex -jobname=nu       "% Juge de la troisième passe : document minimal écrit de zéro pour la cause A (thmtools sous LaTeX 2026-06-01).
% Variantes, par une macro posée sur la ligne de commande (deux passes chacune, depuis memoire/essai/juge3/minimal/) :
%   lualatex -jobname=nu       "\input{minimal-juge3}"                                   sans correctif
%   lualatex -jobname=corr     "\def\CORR{1}\input{minimal-juge3}"                       correctif A du préambule
%   lualatex -jobname=crochet  "\def\SANSCROCHET{1}\input{minimal-juge3}"                sans correctif, crochet retiré
%   lualatex -jobname=amsseul  "\def\AMSSEUL{1}\input{minimal-juge3}"                    amsthm seul (témoin)
%   lualatex -jobname=faux     "\def\CORR{1}\def\FAUX{1}\input{minimal-juge3}"           deux vrais défauts : doivent tomber
%   lualatex -jobname=trace    "\def\TRACE{1}\input{minimal-juge3}"                      trace de l'appel
% Les lignes du correctif sont tirées du préambule commité par extraire_correctifs.py, pas retapées.
\documentclass{scrbook}
\ifdefined\SANSCROCHET \RemoveFromHook{package/amsthm/after}[firstaid/aliascounter]\fi
\usepackage{amsthm}
\usepackage{hyperref}
\usepackage[french, nameinlink, noabbrev]{cleveref}
\ifdefined\AMSSEUL
  \newtheorem{theoreme}{Théorème}[chapter]
  \newtheorem{proposition}[theoreme]{Proposition}
  \newtheorem{corollaire}[theoreme]{Corollaire}
  \newtheorem{remarque}{Remarque}[chapter]
\else
  \usepackage{thmtools}
  \ifdefined\CORR \input{correctif-a-pose}\fi
  \declaretheorem[numberwithin=chapter, name=Théorème]{theoreme}
  \ifdefined\TRACE \tracingmacros=2 \tracingonline=0 \fi
  \declaretheorem[sibling=theoreme, name=Proposition]{proposition}
  \ifdefined\TRACE \tracingmacros=0 \fi
  \declaretheorem[sibling=theoreme, name=Corollaire]{corollaire}
  \declaretheorem[numberwithin=chapter, name=Remarque]{remarque}
  \ifdefined\FAUX \newcounteralias{corollaire}{section}\newcounteralias{orphelin}{inexistant}\fi
  \ifdefined\CORR \input{correctif-a-remise}\fi
\fi
\crefname{theoreme}{théorème}{théorèmes}
\crefname{proposition}{proposition}{propositions}
\crefname{corollaire}{corollaire}{corollaires}
\crefname{remarque}{remarque}{remarques}
\begin{document}
\chapter{Premier}
\begin{proposition}\label{p1}a\end{proposition}
\begin{theoreme}\label{t1}b\end{theoreme}
\begin{remarque}\label{r1}c\end{remarque}
\begin{corollaire}\label{c1}d\end{corollaire}
\chapter{Second}
\begin{corollaire}\label{c2}e\end{corollaire}
\begin{theoreme}\label{t2}f\end{theoreme}
RENVOIS \ref{p1} \ref{t1} \ref{r1} \ref{c1} \ref{c2} \ref{t2} ; \cref{p1} ; \cref{t1} ; \cref{r1} ; \cref{c1} ; \cref{c2} ; \cref{t2} FIN
\end{document}
"                                   sans correctif
%   lualatex -jobname=corr     "\def\CORR{1}% Juge de la troisième passe : document minimal écrit de zéro pour la cause A (thmtools sous LaTeX 2026-06-01).
% Variantes, par une macro posée sur la ligne de commande (deux passes chacune, depuis memoire/essai/juge3/minimal/) :
%   lualatex -jobname=nu       "\input{minimal-juge3}"                                   sans correctif
%   lualatex -jobname=corr     "\def\CORR{1}\input{minimal-juge3}"                       correctif A du préambule
%   lualatex -jobname=crochet  "\def\SANSCROCHET{1}\input{minimal-juge3}"                sans correctif, crochet retiré
%   lualatex -jobname=amsseul  "\def\AMSSEUL{1}\input{minimal-juge3}"                    amsthm seul (témoin)
%   lualatex -jobname=faux     "\def\CORR{1}\def\FAUX{1}\input{minimal-juge3}"           deux vrais défauts : doivent tomber
%   lualatex -jobname=trace    "\def\TRACE{1}\input{minimal-juge3}"                      trace de l'appel
% Les lignes du correctif sont tirées du préambule commité par extraire_correctifs.py, pas retapées.
\documentclass{scrbook}
\ifdefined\SANSCROCHET \RemoveFromHook{package/amsthm/after}[firstaid/aliascounter]\fi
\usepackage{amsthm}
\usepackage{hyperref}
\usepackage[french, nameinlink, noabbrev]{cleveref}
\ifdefined\AMSSEUL
  \newtheorem{theoreme}{Théorème}[chapter]
  \newtheorem{proposition}[theoreme]{Proposition}
  \newtheorem{corollaire}[theoreme]{Corollaire}
  \newtheorem{remarque}{Remarque}[chapter]
\else
  \usepackage{thmtools}
  \ifdefined\CORR \input{correctif-a-pose}\fi
  \declaretheorem[numberwithin=chapter, name=Théorème]{theoreme}
  \ifdefined\TRACE \tracingmacros=2 \tracingonline=0 \fi
  \declaretheorem[sibling=theoreme, name=Proposition]{proposition}
  \ifdefined\TRACE \tracingmacros=0 \fi
  \declaretheorem[sibling=theoreme, name=Corollaire]{corollaire}
  \declaretheorem[numberwithin=chapter, name=Remarque]{remarque}
  \ifdefined\FAUX \newcounteralias{corollaire}{section}\newcounteralias{orphelin}{inexistant}\fi
  \ifdefined\CORR \input{correctif-a-remise}\fi
\fi
\crefname{theoreme}{théorème}{théorèmes}
\crefname{proposition}{proposition}{propositions}
\crefname{corollaire}{corollaire}{corollaires}
\crefname{remarque}{remarque}{remarques}
\begin{document}
\chapter{Premier}
\begin{proposition}\label{p1}a\end{proposition}
\begin{theoreme}\label{t1}b\end{theoreme}
\begin{remarque}\label{r1}c\end{remarque}
\begin{corollaire}\label{c1}d\end{corollaire}
\chapter{Second}
\begin{corollaire}\label{c2}e\end{corollaire}
\begin{theoreme}\label{t2}f\end{theoreme}
RENVOIS \ref{p1} \ref{t1} \ref{r1} \ref{c1} \ref{c2} \ref{t2} ; \cref{p1} ; \cref{t1} ; \cref{r1} ; \cref{c1} ; \cref{c2} ; \cref{t2} FIN
\end{document}
"                       correctif A du préambule
%   lualatex -jobname=crochet  "\def\SANSCROCHET{1}% Juge de la troisième passe : document minimal écrit de zéro pour la cause A (thmtools sous LaTeX 2026-06-01).
% Variantes, par une macro posée sur la ligne de commande (deux passes chacune, depuis memoire/essai/juge3/minimal/) :
%   lualatex -jobname=nu       "\input{minimal-juge3}"                                   sans correctif
%   lualatex -jobname=corr     "\def\CORR{1}\input{minimal-juge3}"                       correctif A du préambule
%   lualatex -jobname=crochet  "\def\SANSCROCHET{1}\input{minimal-juge3}"                sans correctif, crochet retiré
%   lualatex -jobname=amsseul  "\def\AMSSEUL{1}\input{minimal-juge3}"                    amsthm seul (témoin)
%   lualatex -jobname=faux     "\def\CORR{1}\def\FAUX{1}\input{minimal-juge3}"           deux vrais défauts : doivent tomber
%   lualatex -jobname=trace    "\def\TRACE{1}\input{minimal-juge3}"                      trace de l'appel
% Les lignes du correctif sont tirées du préambule commité par extraire_correctifs.py, pas retapées.
\documentclass{scrbook}
\ifdefined\SANSCROCHET \RemoveFromHook{package/amsthm/after}[firstaid/aliascounter]\fi
\usepackage{amsthm}
\usepackage{hyperref}
\usepackage[french, nameinlink, noabbrev]{cleveref}
\ifdefined\AMSSEUL
  \newtheorem{theoreme}{Théorème}[chapter]
  \newtheorem{proposition}[theoreme]{Proposition}
  \newtheorem{corollaire}[theoreme]{Corollaire}
  \newtheorem{remarque}{Remarque}[chapter]
\else
  \usepackage{thmtools}
  \ifdefined\CORR \input{correctif-a-pose}\fi
  \declaretheorem[numberwithin=chapter, name=Théorème]{theoreme}
  \ifdefined\TRACE \tracingmacros=2 \tracingonline=0 \fi
  \declaretheorem[sibling=theoreme, name=Proposition]{proposition}
  \ifdefined\TRACE \tracingmacros=0 \fi
  \declaretheorem[sibling=theoreme, name=Corollaire]{corollaire}
  \declaretheorem[numberwithin=chapter, name=Remarque]{remarque}
  \ifdefined\FAUX \newcounteralias{corollaire}{section}\newcounteralias{orphelin}{inexistant}\fi
  \ifdefined\CORR \input{correctif-a-remise}\fi
\fi
\crefname{theoreme}{théorème}{théorèmes}
\crefname{proposition}{proposition}{propositions}
\crefname{corollaire}{corollaire}{corollaires}
\crefname{remarque}{remarque}{remarques}
\begin{document}
\chapter{Premier}
\begin{proposition}\label{p1}a\end{proposition}
\begin{theoreme}\label{t1}b\end{theoreme}
\begin{remarque}\label{r1}c\end{remarque}
\begin{corollaire}\label{c1}d\end{corollaire}
\chapter{Second}
\begin{corollaire}\label{c2}e\end{corollaire}
\begin{theoreme}\label{t2}f\end{theoreme}
RENVOIS \ref{p1} \ref{t1} \ref{r1} \ref{c1} \ref{c2} \ref{t2} ; \cref{p1} ; \cref{t1} ; \cref{r1} ; \cref{c1} ; \cref{c2} ; \cref{t2} FIN
\end{document}
"                sans correctif, crochet retiré
%   lualatex -jobname=amsseul  "\def\AMSSEUL{1}% Juge de la troisième passe : document minimal écrit de zéro pour la cause A (thmtools sous LaTeX 2026-06-01).
% Variantes, par une macro posée sur la ligne de commande (deux passes chacune, depuis memoire/essai/juge3/minimal/) :
%   lualatex -jobname=nu       "\input{minimal-juge3}"                                   sans correctif
%   lualatex -jobname=corr     "\def\CORR{1}\input{minimal-juge3}"                       correctif A du préambule
%   lualatex -jobname=crochet  "\def\SANSCROCHET{1}\input{minimal-juge3}"                sans correctif, crochet retiré
%   lualatex -jobname=amsseul  "\def\AMSSEUL{1}\input{minimal-juge3}"                    amsthm seul (témoin)
%   lualatex -jobname=faux     "\def\CORR{1}\def\FAUX{1}\input{minimal-juge3}"           deux vrais défauts : doivent tomber
%   lualatex -jobname=trace    "\def\TRACE{1}\input{minimal-juge3}"                      trace de l'appel
% Les lignes du correctif sont tirées du préambule commité par extraire_correctifs.py, pas retapées.
\documentclass{scrbook}
\ifdefined\SANSCROCHET \RemoveFromHook{package/amsthm/after}[firstaid/aliascounter]\fi
\usepackage{amsthm}
\usepackage{hyperref}
\usepackage[french, nameinlink, noabbrev]{cleveref}
\ifdefined\AMSSEUL
  \newtheorem{theoreme}{Théorème}[chapter]
  \newtheorem{proposition}[theoreme]{Proposition}
  \newtheorem{corollaire}[theoreme]{Corollaire}
  \newtheorem{remarque}{Remarque}[chapter]
\else
  \usepackage{thmtools}
  \ifdefined\CORR \input{correctif-a-pose}\fi
  \declaretheorem[numberwithin=chapter, name=Théorème]{theoreme}
  \ifdefined\TRACE \tracingmacros=2 \tracingonline=0 \fi
  \declaretheorem[sibling=theoreme, name=Proposition]{proposition}
  \ifdefined\TRACE \tracingmacros=0 \fi
  \declaretheorem[sibling=theoreme, name=Corollaire]{corollaire}
  \declaretheorem[numberwithin=chapter, name=Remarque]{remarque}
  \ifdefined\FAUX \newcounteralias{corollaire}{section}\newcounteralias{orphelin}{inexistant}\fi
  \ifdefined\CORR \input{correctif-a-remise}\fi
\fi
\crefname{theoreme}{théorème}{théorèmes}
\crefname{proposition}{proposition}{propositions}
\crefname{corollaire}{corollaire}{corollaires}
\crefname{remarque}{remarque}{remarques}
\begin{document}
\chapter{Premier}
\begin{proposition}\label{p1}a\end{proposition}
\begin{theoreme}\label{t1}b\end{theoreme}
\begin{remarque}\label{r1}c\end{remarque}
\begin{corollaire}\label{c1}d\end{corollaire}
\chapter{Second}
\begin{corollaire}\label{c2}e\end{corollaire}
\begin{theoreme}\label{t2}f\end{theoreme}
RENVOIS \ref{p1} \ref{t1} \ref{r1} \ref{c1} \ref{c2} \ref{t2} ; \cref{p1} ; \cref{t1} ; \cref{r1} ; \cref{c1} ; \cref{c2} ; \cref{t2} FIN
\end{document}
"                    amsthm seul (témoin)
%   lualatex -jobname=faux     "\def\CORR{1}\def\FAUX{1}% Juge de la troisième passe : document minimal écrit de zéro pour la cause A (thmtools sous LaTeX 2026-06-01).
% Variantes, par une macro posée sur la ligne de commande (deux passes chacune, depuis memoire/essai/juge3/minimal/) :
%   lualatex -jobname=nu       "\input{minimal-juge3}"                                   sans correctif
%   lualatex -jobname=corr     "\def\CORR{1}\input{minimal-juge3}"                       correctif A du préambule
%   lualatex -jobname=crochet  "\def\SANSCROCHET{1}\input{minimal-juge3}"                sans correctif, crochet retiré
%   lualatex -jobname=amsseul  "\def\AMSSEUL{1}\input{minimal-juge3}"                    amsthm seul (témoin)
%   lualatex -jobname=faux     "\def\CORR{1}\def\FAUX{1}\input{minimal-juge3}"           deux vrais défauts : doivent tomber
%   lualatex -jobname=trace    "\def\TRACE{1}\input{minimal-juge3}"                      trace de l'appel
% Les lignes du correctif sont tirées du préambule commité par extraire_correctifs.py, pas retapées.
\documentclass{scrbook}
\ifdefined\SANSCROCHET \RemoveFromHook{package/amsthm/after}[firstaid/aliascounter]\fi
\usepackage{amsthm}
\usepackage{hyperref}
\usepackage[french, nameinlink, noabbrev]{cleveref}
\ifdefined\AMSSEUL
  \newtheorem{theoreme}{Théorème}[chapter]
  \newtheorem{proposition}[theoreme]{Proposition}
  \newtheorem{corollaire}[theoreme]{Corollaire}
  \newtheorem{remarque}{Remarque}[chapter]
\else
  \usepackage{thmtools}
  \ifdefined\CORR \input{correctif-a-pose}\fi
  \declaretheorem[numberwithin=chapter, name=Théorème]{theoreme}
  \ifdefined\TRACE \tracingmacros=2 \tracingonline=0 \fi
  \declaretheorem[sibling=theoreme, name=Proposition]{proposition}
  \ifdefined\TRACE \tracingmacros=0 \fi
  \declaretheorem[sibling=theoreme, name=Corollaire]{corollaire}
  \declaretheorem[numberwithin=chapter, name=Remarque]{remarque}
  \ifdefined\FAUX \newcounteralias{corollaire}{section}\newcounteralias{orphelin}{inexistant}\fi
  \ifdefined\CORR \input{correctif-a-remise}\fi
\fi
\crefname{theoreme}{théorème}{théorèmes}
\crefname{proposition}{proposition}{propositions}
\crefname{corollaire}{corollaire}{corollaires}
\crefname{remarque}{remarque}{remarques}
\begin{document}
\chapter{Premier}
\begin{proposition}\label{p1}a\end{proposition}
\begin{theoreme}\label{t1}b\end{theoreme}
\begin{remarque}\label{r1}c\end{remarque}
\begin{corollaire}\label{c1}d\end{corollaire}
\chapter{Second}
\begin{corollaire}\label{c2}e\end{corollaire}
\begin{theoreme}\label{t2}f\end{theoreme}
RENVOIS \ref{p1} \ref{t1} \ref{r1} \ref{c1} \ref{c2} \ref{t2} ; \cref{p1} ; \cref{t1} ; \cref{r1} ; \cref{c1} ; \cref{c2} ; \cref{t2} FIN
\end{document}
"           deux vrais défauts : doivent tomber
%   lualatex -jobname=trace    "\def\TRACE{1}% Juge de la troisième passe : document minimal écrit de zéro pour la cause A (thmtools sous LaTeX 2026-06-01).
% Variantes, par une macro posée sur la ligne de commande (deux passes chacune, depuis memoire/essai/juge3/minimal/) :
%   lualatex -jobname=nu       "\input{minimal-juge3}"                                   sans correctif
%   lualatex -jobname=corr     "\def\CORR{1}\input{minimal-juge3}"                       correctif A du préambule
%   lualatex -jobname=crochet  "\def\SANSCROCHET{1}\input{minimal-juge3}"                sans correctif, crochet retiré
%   lualatex -jobname=amsseul  "\def\AMSSEUL{1}\input{minimal-juge3}"                    amsthm seul (témoin)
%   lualatex -jobname=faux     "\def\CORR{1}\def\FAUX{1}\input{minimal-juge3}"           deux vrais défauts : doivent tomber
%   lualatex -jobname=trace    "\def\TRACE{1}\input{minimal-juge3}"                      trace de l'appel
% Les lignes du correctif sont tirées du préambule commité par extraire_correctifs.py, pas retapées.
\documentclass{scrbook}
\ifdefined\SANSCROCHET \RemoveFromHook{package/amsthm/after}[firstaid/aliascounter]\fi
\usepackage{amsthm}
\usepackage{hyperref}
\usepackage[french, nameinlink, noabbrev]{cleveref}
\ifdefined\AMSSEUL
  \newtheorem{theoreme}{Théorème}[chapter]
  \newtheorem{proposition}[theoreme]{Proposition}
  \newtheorem{corollaire}[theoreme]{Corollaire}
  \newtheorem{remarque}{Remarque}[chapter]
\else
  \usepackage{thmtools}
  \ifdefined\CORR \input{correctif-a-pose}\fi
  \declaretheorem[numberwithin=chapter, name=Théorème]{theoreme}
  \ifdefined\TRACE \tracingmacros=2 \tracingonline=0 \fi
  \declaretheorem[sibling=theoreme, name=Proposition]{proposition}
  \ifdefined\TRACE \tracingmacros=0 \fi
  \declaretheorem[sibling=theoreme, name=Corollaire]{corollaire}
  \declaretheorem[numberwithin=chapter, name=Remarque]{remarque}
  \ifdefined\FAUX \newcounteralias{corollaire}{section}\newcounteralias{orphelin}{inexistant}\fi
  \ifdefined\CORR \input{correctif-a-remise}\fi
\fi
\crefname{theoreme}{théorème}{théorèmes}
\crefname{proposition}{proposition}{propositions}
\crefname{corollaire}{corollaire}{corollaires}
\crefname{remarque}{remarque}{remarques}
\begin{document}
\chapter{Premier}
\begin{proposition}\label{p1}a\end{proposition}
\begin{theoreme}\label{t1}b\end{theoreme}
\begin{remarque}\label{r1}c\end{remarque}
\begin{corollaire}\label{c1}d\end{corollaire}
\chapter{Second}
\begin{corollaire}\label{c2}e\end{corollaire}
\begin{theoreme}\label{t2}f\end{theoreme}
RENVOIS \ref{p1} \ref{t1} \ref{r1} \ref{c1} \ref{c2} \ref{t2} ; \cref{p1} ; \cref{t1} ; \cref{r1} ; \cref{c1} ; \cref{c2} ; \cref{t2} FIN
\end{document}
"                      trace de l'appel
% Les lignes du correctif sont tirées du préambule commité par extraire_correctifs.py, pas retapées.
\documentclass{scrbook}
\ifdefined\SANSCROCHET \RemoveFromHook{package/amsthm/after}[firstaid/aliascounter]\fi
\usepackage{amsthm}
\usepackage{hyperref}
\usepackage[french, nameinlink, noabbrev]{cleveref}
\ifdefined\AMSSEUL
  \newtheorem{theoreme}{Théorème}[chapter]
  \newtheorem{proposition}[theoreme]{Proposition}
  \newtheorem{corollaire}[theoreme]{Corollaire}
  \newtheorem{remarque}{Remarque}[chapter]
\else
  \usepackage{thmtools}
  \ifdefined\CORR \input{correctif-a-pose}\fi
  \declaretheorem[numberwithin=chapter, name=Théorème]{theoreme}
  \ifdefined\TRACE \tracingmacros=2 \tracingonline=0 \fi
  \declaretheorem[sibling=theoreme, name=Proposition]{proposition}
  \ifdefined\TRACE \tracingmacros=0 \fi
  \declaretheorem[sibling=theoreme, name=Corollaire]{corollaire}
  \declaretheorem[numberwithin=chapter, name=Remarque]{remarque}
  \ifdefined\FAUX \newcounteralias{corollaire}{section}\newcounteralias{orphelin}{inexistant}\fi
  \ifdefined\CORR \input{correctif-a-remise}\fi
\fi
\crefname{theoreme}{théorème}{théorèmes}
\crefname{proposition}{proposition}{propositions}
\crefname{corollaire}{corollaire}{corollaires}
\crefname{remarque}{remarque}{remarques}
\begin{document}
\chapter{Premier}
\begin{proposition}\label{p1}a\end{proposition}
\begin{theoreme}\label{t1}b\end{theoreme}
\begin{remarque}\label{r1}c\end{remarque}
\begin{corollaire}\label{c1}d\end{corollaire}
\chapter{Second}
\begin{corollaire}\label{c2}e\end{corollaire}
\begin{theoreme}\label{t2}f\end{theoreme}
RENVOIS \ref{p1} \ref{t1} \ref{r1} \ref{c1} \ref{c2} \ref{t2} ; \cref{p1} ; \cref{t1} ; \cref{r1} ; \cref{c1} ; \cref{c2} ; \cref{t2} FIN
\end{document}
"                    amsthm seul (témoin)
%   lualatex -jobname=faux     "\def\CORR{1}\def\FAUX{1}% Juge de la troisième passe : document minimal écrit de zéro pour la cause A (thmtools sous LaTeX 2026-06-01).
% Variantes, par une macro posée sur la ligne de commande (deux passes chacune, depuis memoire/essai/juge3/minimal/) :
%   lualatex -jobname=nu       "% Juge de la troisième passe : document minimal écrit de zéro pour la cause A (thmtools sous LaTeX 2026-06-01).
% Variantes, par une macro posée sur la ligne de commande (deux passes chacune, depuis memoire/essai/juge3/minimal/) :
%   lualatex -jobname=nu       "\input{minimal-juge3}"                                   sans correctif
%   lualatex -jobname=corr     "\def\CORR{1}\input{minimal-juge3}"                       correctif A du préambule
%   lualatex -jobname=crochet  "\def\SANSCROCHET{1}\input{minimal-juge3}"                sans correctif, crochet retiré
%   lualatex -jobname=amsseul  "\def\AMSSEUL{1}\input{minimal-juge3}"                    amsthm seul (témoin)
%   lualatex -jobname=faux     "\def\CORR{1}\def\FAUX{1}\input{minimal-juge3}"           deux vrais défauts : doivent tomber
%   lualatex -jobname=trace    "\def\TRACE{1}\input{minimal-juge3}"                      trace de l'appel
% Les lignes du correctif sont tirées du préambule commité par extraire_correctifs.py, pas retapées.
\documentclass{scrbook}
\ifdefined\SANSCROCHET \RemoveFromHook{package/amsthm/after}[firstaid/aliascounter]\fi
\usepackage{amsthm}
\usepackage{hyperref}
\usepackage[french, nameinlink, noabbrev]{cleveref}
\ifdefined\AMSSEUL
  \newtheorem{theoreme}{Théorème}[chapter]
  \newtheorem{proposition}[theoreme]{Proposition}
  \newtheorem{corollaire}[theoreme]{Corollaire}
  \newtheorem{remarque}{Remarque}[chapter]
\else
  \usepackage{thmtools}
  \ifdefined\CORR \input{correctif-a-pose}\fi
  \declaretheorem[numberwithin=chapter, name=Théorème]{theoreme}
  \ifdefined\TRACE \tracingmacros=2 \tracingonline=0 \fi
  \declaretheorem[sibling=theoreme, name=Proposition]{proposition}
  \ifdefined\TRACE \tracingmacros=0 \fi
  \declaretheorem[sibling=theoreme, name=Corollaire]{corollaire}
  \declaretheorem[numberwithin=chapter, name=Remarque]{remarque}
  \ifdefined\FAUX \newcounteralias{corollaire}{section}\newcounteralias{orphelin}{inexistant}\fi
  \ifdefined\CORR \input{correctif-a-remise}\fi
\fi
\crefname{theoreme}{théorème}{théorèmes}
\crefname{proposition}{proposition}{propositions}
\crefname{corollaire}{corollaire}{corollaires}
\crefname{remarque}{remarque}{remarques}
\begin{document}
\chapter{Premier}
\begin{proposition}\label{p1}a\end{proposition}
\begin{theoreme}\label{t1}b\end{theoreme}
\begin{remarque}\label{r1}c\end{remarque}
\begin{corollaire}\label{c1}d\end{corollaire}
\chapter{Second}
\begin{corollaire}\label{c2}e\end{corollaire}
\begin{theoreme}\label{t2}f\end{theoreme}
RENVOIS \ref{p1} \ref{t1} \ref{r1} \ref{c1} \ref{c2} \ref{t2} ; \cref{p1} ; \cref{t1} ; \cref{r1} ; \cref{c1} ; \cref{c2} ; \cref{t2} FIN
\end{document}
"                                   sans correctif
%   lualatex -jobname=corr     "\def\CORR{1}% Juge de la troisième passe : document minimal écrit de zéro pour la cause A (thmtools sous LaTeX 2026-06-01).
% Variantes, par une macro posée sur la ligne de commande (deux passes chacune, depuis memoire/essai/juge3/minimal/) :
%   lualatex -jobname=nu       "\input{minimal-juge3}"                                   sans correctif
%   lualatex -jobname=corr     "\def\CORR{1}\input{minimal-juge3}"                       correctif A du préambule
%   lualatex -jobname=crochet  "\def\SANSCROCHET{1}\input{minimal-juge3}"                sans correctif, crochet retiré
%   lualatex -jobname=amsseul  "\def\AMSSEUL{1}\input{minimal-juge3}"                    amsthm seul (témoin)
%   lualatex -jobname=faux     "\def\CORR{1}\def\FAUX{1}\input{minimal-juge3}"           deux vrais défauts : doivent tomber
%   lualatex -jobname=trace    "\def\TRACE{1}\input{minimal-juge3}"                      trace de l'appel
% Les lignes du correctif sont tirées du préambule commité par extraire_correctifs.py, pas retapées.
\documentclass{scrbook}
\ifdefined\SANSCROCHET \RemoveFromHook{package/amsthm/after}[firstaid/aliascounter]\fi
\usepackage{amsthm}
\usepackage{hyperref}
\usepackage[french, nameinlink, noabbrev]{cleveref}
\ifdefined\AMSSEUL
  \newtheorem{theoreme}{Théorème}[chapter]
  \newtheorem{proposition}[theoreme]{Proposition}
  \newtheorem{corollaire}[theoreme]{Corollaire}
  \newtheorem{remarque}{Remarque}[chapter]
\else
  \usepackage{thmtools}
  \ifdefined\CORR \input{correctif-a-pose}\fi
  \declaretheorem[numberwithin=chapter, name=Théorème]{theoreme}
  \ifdefined\TRACE \tracingmacros=2 \tracingonline=0 \fi
  \declaretheorem[sibling=theoreme, name=Proposition]{proposition}
  \ifdefined\TRACE \tracingmacros=0 \fi
  \declaretheorem[sibling=theoreme, name=Corollaire]{corollaire}
  \declaretheorem[numberwithin=chapter, name=Remarque]{remarque}
  \ifdefined\FAUX \newcounteralias{corollaire}{section}\newcounteralias{orphelin}{inexistant}\fi
  \ifdefined\CORR \input{correctif-a-remise}\fi
\fi
\crefname{theoreme}{théorème}{théorèmes}
\crefname{proposition}{proposition}{propositions}
\crefname{corollaire}{corollaire}{corollaires}
\crefname{remarque}{remarque}{remarques}
\begin{document}
\chapter{Premier}
\begin{proposition}\label{p1}a\end{proposition}
\begin{theoreme}\label{t1}b\end{theoreme}
\begin{remarque}\label{r1}c\end{remarque}
\begin{corollaire}\label{c1}d\end{corollaire}
\chapter{Second}
\begin{corollaire}\label{c2}e\end{corollaire}
\begin{theoreme}\label{t2}f\end{theoreme}
RENVOIS \ref{p1} \ref{t1} \ref{r1} \ref{c1} \ref{c2} \ref{t2} ; \cref{p1} ; \cref{t1} ; \cref{r1} ; \cref{c1} ; \cref{c2} ; \cref{t2} FIN
\end{document}
"                       correctif A du préambule
%   lualatex -jobname=crochet  "\def\SANSCROCHET{1}% Juge de la troisième passe : document minimal écrit de zéro pour la cause A (thmtools sous LaTeX 2026-06-01).
% Variantes, par une macro posée sur la ligne de commande (deux passes chacune, depuis memoire/essai/juge3/minimal/) :
%   lualatex -jobname=nu       "\input{minimal-juge3}"                                   sans correctif
%   lualatex -jobname=corr     "\def\CORR{1}\input{minimal-juge3}"                       correctif A du préambule
%   lualatex -jobname=crochet  "\def\SANSCROCHET{1}\input{minimal-juge3}"                sans correctif, crochet retiré
%   lualatex -jobname=amsseul  "\def\AMSSEUL{1}\input{minimal-juge3}"                    amsthm seul (témoin)
%   lualatex -jobname=faux     "\def\CORR{1}\def\FAUX{1}\input{minimal-juge3}"           deux vrais défauts : doivent tomber
%   lualatex -jobname=trace    "\def\TRACE{1}\input{minimal-juge3}"                      trace de l'appel
% Les lignes du correctif sont tirées du préambule commité par extraire_correctifs.py, pas retapées.
\documentclass{scrbook}
\ifdefined\SANSCROCHET \RemoveFromHook{package/amsthm/after}[firstaid/aliascounter]\fi
\usepackage{amsthm}
\usepackage{hyperref}
\usepackage[french, nameinlink, noabbrev]{cleveref}
\ifdefined\AMSSEUL
  \newtheorem{theoreme}{Théorème}[chapter]
  \newtheorem{proposition}[theoreme]{Proposition}
  \newtheorem{corollaire}[theoreme]{Corollaire}
  \newtheorem{remarque}{Remarque}[chapter]
\else
  \usepackage{thmtools}
  \ifdefined\CORR \input{correctif-a-pose}\fi
  \declaretheorem[numberwithin=chapter, name=Théorème]{theoreme}
  \ifdefined\TRACE \tracingmacros=2 \tracingonline=0 \fi
  \declaretheorem[sibling=theoreme, name=Proposition]{proposition}
  \ifdefined\TRACE \tracingmacros=0 \fi
  \declaretheorem[sibling=theoreme, name=Corollaire]{corollaire}
  \declaretheorem[numberwithin=chapter, name=Remarque]{remarque}
  \ifdefined\FAUX \newcounteralias{corollaire}{section}\newcounteralias{orphelin}{inexistant}\fi
  \ifdefined\CORR \input{correctif-a-remise}\fi
\fi
\crefname{theoreme}{théorème}{théorèmes}
\crefname{proposition}{proposition}{propositions}
\crefname{corollaire}{corollaire}{corollaires}
\crefname{remarque}{remarque}{remarques}
\begin{document}
\chapter{Premier}
\begin{proposition}\label{p1}a\end{proposition}
\begin{theoreme}\label{t1}b\end{theoreme}
\begin{remarque}\label{r1}c\end{remarque}
\begin{corollaire}\label{c1}d\end{corollaire}
\chapter{Second}
\begin{corollaire}\label{c2}e\end{corollaire}
\begin{theoreme}\label{t2}f\end{theoreme}
RENVOIS \ref{p1} \ref{t1} \ref{r1} \ref{c1} \ref{c2} \ref{t2} ; \cref{p1} ; \cref{t1} ; \cref{r1} ; \cref{c1} ; \cref{c2} ; \cref{t2} FIN
\end{document}
"                sans correctif, crochet retiré
%   lualatex -jobname=amsseul  "\def\AMSSEUL{1}% Juge de la troisième passe : document minimal écrit de zéro pour la cause A (thmtools sous LaTeX 2026-06-01).
% Variantes, par une macro posée sur la ligne de commande (deux passes chacune, depuis memoire/essai/juge3/minimal/) :
%   lualatex -jobname=nu       "\input{minimal-juge3}"                                   sans correctif
%   lualatex -jobname=corr     "\def\CORR{1}\input{minimal-juge3}"                       correctif A du préambule
%   lualatex -jobname=crochet  "\def\SANSCROCHET{1}\input{minimal-juge3}"                sans correctif, crochet retiré
%   lualatex -jobname=amsseul  "\def\AMSSEUL{1}\input{minimal-juge3}"                    amsthm seul (témoin)
%   lualatex -jobname=faux     "\def\CORR{1}\def\FAUX{1}\input{minimal-juge3}"           deux vrais défauts : doivent tomber
%   lualatex -jobname=trace    "\def\TRACE{1}\input{minimal-juge3}"                      trace de l'appel
% Les lignes du correctif sont tirées du préambule commité par extraire_correctifs.py, pas retapées.
\documentclass{scrbook}
\ifdefined\SANSCROCHET \RemoveFromHook{package/amsthm/after}[firstaid/aliascounter]\fi
\usepackage{amsthm}
\usepackage{hyperref}
\usepackage[french, nameinlink, noabbrev]{cleveref}
\ifdefined\AMSSEUL
  \newtheorem{theoreme}{Théorème}[chapter]
  \newtheorem{proposition}[theoreme]{Proposition}
  \newtheorem{corollaire}[theoreme]{Corollaire}
  \newtheorem{remarque}{Remarque}[chapter]
\else
  \usepackage{thmtools}
  \ifdefined\CORR \input{correctif-a-pose}\fi
  \declaretheorem[numberwithin=chapter, name=Théorème]{theoreme}
  \ifdefined\TRACE \tracingmacros=2 \tracingonline=0 \fi
  \declaretheorem[sibling=theoreme, name=Proposition]{proposition}
  \ifdefined\TRACE \tracingmacros=0 \fi
  \declaretheorem[sibling=theoreme, name=Corollaire]{corollaire}
  \declaretheorem[numberwithin=chapter, name=Remarque]{remarque}
  \ifdefined\FAUX \newcounteralias{corollaire}{section}\newcounteralias{orphelin}{inexistant}\fi
  \ifdefined\CORR \input{correctif-a-remise}\fi
\fi
\crefname{theoreme}{théorème}{théorèmes}
\crefname{proposition}{proposition}{propositions}
\crefname{corollaire}{corollaire}{corollaires}
\crefname{remarque}{remarque}{remarques}
\begin{document}
\chapter{Premier}
\begin{proposition}\label{p1}a\end{proposition}
\begin{theoreme}\label{t1}b\end{theoreme}
\begin{remarque}\label{r1}c\end{remarque}
\begin{corollaire}\label{c1}d\end{corollaire}
\chapter{Second}
\begin{corollaire}\label{c2}e\end{corollaire}
\begin{theoreme}\label{t2}f\end{theoreme}
RENVOIS \ref{p1} \ref{t1} \ref{r1} \ref{c1} \ref{c2} \ref{t2} ; \cref{p1} ; \cref{t1} ; \cref{r1} ; \cref{c1} ; \cref{c2} ; \cref{t2} FIN
\end{document}
"                    amsthm seul (témoin)
%   lualatex -jobname=faux     "\def\CORR{1}\def\FAUX{1}% Juge de la troisième passe : document minimal écrit de zéro pour la cause A (thmtools sous LaTeX 2026-06-01).
% Variantes, par une macro posée sur la ligne de commande (deux passes chacune, depuis memoire/essai/juge3/minimal/) :
%   lualatex -jobname=nu       "\input{minimal-juge3}"                                   sans correctif
%   lualatex -jobname=corr     "\def\CORR{1}\input{minimal-juge3}"                       correctif A du préambule
%   lualatex -jobname=crochet  "\def\SANSCROCHET{1}\input{minimal-juge3}"                sans correctif, crochet retiré
%   lualatex -jobname=amsseul  "\def\AMSSEUL{1}\input{minimal-juge3}"                    amsthm seul (témoin)
%   lualatex -jobname=faux     "\def\CORR{1}\def\FAUX{1}\input{minimal-juge3}"           deux vrais défauts : doivent tomber
%   lualatex -jobname=trace    "\def\TRACE{1}\input{minimal-juge3}"                      trace de l'appel
% Les lignes du correctif sont tirées du préambule commité par extraire_correctifs.py, pas retapées.
\documentclass{scrbook}
\ifdefined\SANSCROCHET \RemoveFromHook{package/amsthm/after}[firstaid/aliascounter]\fi
\usepackage{amsthm}
\usepackage{hyperref}
\usepackage[french, nameinlink, noabbrev]{cleveref}
\ifdefined\AMSSEUL
  \newtheorem{theoreme}{Théorème}[chapter]
  \newtheorem{proposition}[theoreme]{Proposition}
  \newtheorem{corollaire}[theoreme]{Corollaire}
  \newtheorem{remarque}{Remarque}[chapter]
\else
  \usepackage{thmtools}
  \ifdefined\CORR \input{correctif-a-pose}\fi
  \declaretheorem[numberwithin=chapter, name=Théorème]{theoreme}
  \ifdefined\TRACE \tracingmacros=2 \tracingonline=0 \fi
  \declaretheorem[sibling=theoreme, name=Proposition]{proposition}
  \ifdefined\TRACE \tracingmacros=0 \fi
  \declaretheorem[sibling=theoreme, name=Corollaire]{corollaire}
  \declaretheorem[numberwithin=chapter, name=Remarque]{remarque}
  \ifdefined\FAUX \newcounteralias{corollaire}{section}\newcounteralias{orphelin}{inexistant}\fi
  \ifdefined\CORR \input{correctif-a-remise}\fi
\fi
\crefname{theoreme}{théorème}{théorèmes}
\crefname{proposition}{proposition}{propositions}
\crefname{corollaire}{corollaire}{corollaires}
\crefname{remarque}{remarque}{remarques}
\begin{document}
\chapter{Premier}
\begin{proposition}\label{p1}a\end{proposition}
\begin{theoreme}\label{t1}b\end{theoreme}
\begin{remarque}\label{r1}c\end{remarque}
\begin{corollaire}\label{c1}d\end{corollaire}
\chapter{Second}
\begin{corollaire}\label{c2}e\end{corollaire}
\begin{theoreme}\label{t2}f\end{theoreme}
RENVOIS \ref{p1} \ref{t1} \ref{r1} \ref{c1} \ref{c2} \ref{t2} ; \cref{p1} ; \cref{t1} ; \cref{r1} ; \cref{c1} ; \cref{c2} ; \cref{t2} FIN
\end{document}
"           deux vrais défauts : doivent tomber
%   lualatex -jobname=trace    "\def\TRACE{1}% Juge de la troisième passe : document minimal écrit de zéro pour la cause A (thmtools sous LaTeX 2026-06-01).
% Variantes, par une macro posée sur la ligne de commande (deux passes chacune, depuis memoire/essai/juge3/minimal/) :
%   lualatex -jobname=nu       "\input{minimal-juge3}"                                   sans correctif
%   lualatex -jobname=corr     "\def\CORR{1}\input{minimal-juge3}"                       correctif A du préambule
%   lualatex -jobname=crochet  "\def\SANSCROCHET{1}\input{minimal-juge3}"                sans correctif, crochet retiré
%   lualatex -jobname=amsseul  "\def\AMSSEUL{1}\input{minimal-juge3}"                    amsthm seul (témoin)
%   lualatex -jobname=faux     "\def\CORR{1}\def\FAUX{1}\input{minimal-juge3}"           deux vrais défauts : doivent tomber
%   lualatex -jobname=trace    "\def\TRACE{1}\input{minimal-juge3}"                      trace de l'appel
% Les lignes du correctif sont tirées du préambule commité par extraire_correctifs.py, pas retapées.
\documentclass{scrbook}
\ifdefined\SANSCROCHET \RemoveFromHook{package/amsthm/after}[firstaid/aliascounter]\fi
\usepackage{amsthm}
\usepackage{hyperref}
\usepackage[french, nameinlink, noabbrev]{cleveref}
\ifdefined\AMSSEUL
  \newtheorem{theoreme}{Théorème}[chapter]
  \newtheorem{proposition}[theoreme]{Proposition}
  \newtheorem{corollaire}[theoreme]{Corollaire}
  \newtheorem{remarque}{Remarque}[chapter]
\else
  \usepackage{thmtools}
  \ifdefined\CORR \input{correctif-a-pose}\fi
  \declaretheorem[numberwithin=chapter, name=Théorème]{theoreme}
  \ifdefined\TRACE \tracingmacros=2 \tracingonline=0 \fi
  \declaretheorem[sibling=theoreme, name=Proposition]{proposition}
  \ifdefined\TRACE \tracingmacros=0 \fi
  \declaretheorem[sibling=theoreme, name=Corollaire]{corollaire}
  \declaretheorem[numberwithin=chapter, name=Remarque]{remarque}
  \ifdefined\FAUX \newcounteralias{corollaire}{section}\newcounteralias{orphelin}{inexistant}\fi
  \ifdefined\CORR \input{correctif-a-remise}\fi
\fi
\crefname{theoreme}{théorème}{théorèmes}
\crefname{proposition}{proposition}{propositions}
\crefname{corollaire}{corollaire}{corollaires}
\crefname{remarque}{remarque}{remarques}
\begin{document}
\chapter{Premier}
\begin{proposition}\label{p1}a\end{proposition}
\begin{theoreme}\label{t1}b\end{theoreme}
\begin{remarque}\label{r1}c\end{remarque}
\begin{corollaire}\label{c1}d\end{corollaire}
\chapter{Second}
\begin{corollaire}\label{c2}e\end{corollaire}
\begin{theoreme}\label{t2}f\end{theoreme}
RENVOIS \ref{p1} \ref{t1} \ref{r1} \ref{c1} \ref{c2} \ref{t2} ; \cref{p1} ; \cref{t1} ; \cref{r1} ; \cref{c1} ; \cref{c2} ; \cref{t2} FIN
\end{document}
"                      trace de l'appel
% Les lignes du correctif sont tirées du préambule commité par extraire_correctifs.py, pas retapées.
\documentclass{scrbook}
\ifdefined\SANSCROCHET \RemoveFromHook{package/amsthm/after}[firstaid/aliascounter]\fi
\usepackage{amsthm}
\usepackage{hyperref}
\usepackage[french, nameinlink, noabbrev]{cleveref}
\ifdefined\AMSSEUL
  \newtheorem{theoreme}{Théorème}[chapter]
  \newtheorem{proposition}[theoreme]{Proposition}
  \newtheorem{corollaire}[theoreme]{Corollaire}
  \newtheorem{remarque}{Remarque}[chapter]
\else
  \usepackage{thmtools}
  \ifdefined\CORR \input{correctif-a-pose}\fi
  \declaretheorem[numberwithin=chapter, name=Théorème]{theoreme}
  \ifdefined\TRACE \tracingmacros=2 \tracingonline=0 \fi
  \declaretheorem[sibling=theoreme, name=Proposition]{proposition}
  \ifdefined\TRACE \tracingmacros=0 \fi
  \declaretheorem[sibling=theoreme, name=Corollaire]{corollaire}
  \declaretheorem[numberwithin=chapter, name=Remarque]{remarque}
  \ifdefined\FAUX \newcounteralias{corollaire}{section}\newcounteralias{orphelin}{inexistant}\fi
  \ifdefined\CORR \input{correctif-a-remise}\fi
\fi
\crefname{theoreme}{théorème}{théorèmes}
\crefname{proposition}{proposition}{propositions}
\crefname{corollaire}{corollaire}{corollaires}
\crefname{remarque}{remarque}{remarques}
\begin{document}
\chapter{Premier}
\begin{proposition}\label{p1}a\end{proposition}
\begin{theoreme}\label{t1}b\end{theoreme}
\begin{remarque}\label{r1}c\end{remarque}
\begin{corollaire}\label{c1}d\end{corollaire}
\chapter{Second}
\begin{corollaire}\label{c2}e\end{corollaire}
\begin{theoreme}\label{t2}f\end{theoreme}
RENVOIS \ref{p1} \ref{t1} \ref{r1} \ref{c1} \ref{c2} \ref{t2} ; \cref{p1} ; \cref{t1} ; \cref{r1} ; \cref{c1} ; \cref{c2} ; \cref{t2} FIN
\end{document}
"           deux vrais défauts : doivent tomber
%   lualatex -jobname=trace    "\def\TRACE{1}% Juge de la troisième passe : document minimal écrit de zéro pour la cause A (thmtools sous LaTeX 2026-06-01).
% Variantes, par une macro posée sur la ligne de commande (deux passes chacune, depuis memoire/essai/juge3/minimal/) :
%   lualatex -jobname=nu       "% Juge de la troisième passe : document minimal écrit de zéro pour la cause A (thmtools sous LaTeX 2026-06-01).
% Variantes, par une macro posée sur la ligne de commande (deux passes chacune, depuis memoire/essai/juge3/minimal/) :
%   lualatex -jobname=nu       "\input{minimal-juge3}"                                   sans correctif
%   lualatex -jobname=corr     "\def\CORR{1}\input{minimal-juge3}"                       correctif A du préambule
%   lualatex -jobname=crochet  "\def\SANSCROCHET{1}\input{minimal-juge3}"                sans correctif, crochet retiré
%   lualatex -jobname=amsseul  "\def\AMSSEUL{1}\input{minimal-juge3}"                    amsthm seul (témoin)
%   lualatex -jobname=faux     "\def\CORR{1}\def\FAUX{1}\input{minimal-juge3}"           deux vrais défauts : doivent tomber
%   lualatex -jobname=trace    "\def\TRACE{1}\input{minimal-juge3}"                      trace de l'appel
% Les lignes du correctif sont tirées du préambule commité par extraire_correctifs.py, pas retapées.
\documentclass{scrbook}
\ifdefined\SANSCROCHET \RemoveFromHook{package/amsthm/after}[firstaid/aliascounter]\fi
\usepackage{amsthm}
\usepackage{hyperref}
\usepackage[french, nameinlink, noabbrev]{cleveref}
\ifdefined\AMSSEUL
  \newtheorem{theoreme}{Théorème}[chapter]
  \newtheorem{proposition}[theoreme]{Proposition}
  \newtheorem{corollaire}[theoreme]{Corollaire}
  \newtheorem{remarque}{Remarque}[chapter]
\else
  \usepackage{thmtools}
  \ifdefined\CORR \input{correctif-a-pose}\fi
  \declaretheorem[numberwithin=chapter, name=Théorème]{theoreme}
  \ifdefined\TRACE \tracingmacros=2 \tracingonline=0 \fi
  \declaretheorem[sibling=theoreme, name=Proposition]{proposition}
  \ifdefined\TRACE \tracingmacros=0 \fi
  \declaretheorem[sibling=theoreme, name=Corollaire]{corollaire}
  \declaretheorem[numberwithin=chapter, name=Remarque]{remarque}
  \ifdefined\FAUX \newcounteralias{corollaire}{section}\newcounteralias{orphelin}{inexistant}\fi
  \ifdefined\CORR \input{correctif-a-remise}\fi
\fi
\crefname{theoreme}{théorème}{théorèmes}
\crefname{proposition}{proposition}{propositions}
\crefname{corollaire}{corollaire}{corollaires}
\crefname{remarque}{remarque}{remarques}
\begin{document}
\chapter{Premier}
\begin{proposition}\label{p1}a\end{proposition}
\begin{theoreme}\label{t1}b\end{theoreme}
\begin{remarque}\label{r1}c\end{remarque}
\begin{corollaire}\label{c1}d\end{corollaire}
\chapter{Second}
\begin{corollaire}\label{c2}e\end{corollaire}
\begin{theoreme}\label{t2}f\end{theoreme}
RENVOIS \ref{p1} \ref{t1} \ref{r1} \ref{c1} \ref{c2} \ref{t2} ; \cref{p1} ; \cref{t1} ; \cref{r1} ; \cref{c1} ; \cref{c2} ; \cref{t2} FIN
\end{document}
"                                   sans correctif
%   lualatex -jobname=corr     "\def\CORR{1}% Juge de la troisième passe : document minimal écrit de zéro pour la cause A (thmtools sous LaTeX 2026-06-01).
% Variantes, par une macro posée sur la ligne de commande (deux passes chacune, depuis memoire/essai/juge3/minimal/) :
%   lualatex -jobname=nu       "\input{minimal-juge3}"                                   sans correctif
%   lualatex -jobname=corr     "\def\CORR{1}\input{minimal-juge3}"                       correctif A du préambule
%   lualatex -jobname=crochet  "\def\SANSCROCHET{1}\input{minimal-juge3}"                sans correctif, crochet retiré
%   lualatex -jobname=amsseul  "\def\AMSSEUL{1}\input{minimal-juge3}"                    amsthm seul (témoin)
%   lualatex -jobname=faux     "\def\CORR{1}\def\FAUX{1}\input{minimal-juge3}"           deux vrais défauts : doivent tomber
%   lualatex -jobname=trace    "\def\TRACE{1}\input{minimal-juge3}"                      trace de l'appel
% Les lignes du correctif sont tirées du préambule commité par extraire_correctifs.py, pas retapées.
\documentclass{scrbook}
\ifdefined\SANSCROCHET \RemoveFromHook{package/amsthm/after}[firstaid/aliascounter]\fi
\usepackage{amsthm}
\usepackage{hyperref}
\usepackage[french, nameinlink, noabbrev]{cleveref}
\ifdefined\AMSSEUL
  \newtheorem{theoreme}{Théorème}[chapter]
  \newtheorem{proposition}[theoreme]{Proposition}
  \newtheorem{corollaire}[theoreme]{Corollaire}
  \newtheorem{remarque}{Remarque}[chapter]
\else
  \usepackage{thmtools}
  \ifdefined\CORR \input{correctif-a-pose}\fi
  \declaretheorem[numberwithin=chapter, name=Théorème]{theoreme}
  \ifdefined\TRACE \tracingmacros=2 \tracingonline=0 \fi
  \declaretheorem[sibling=theoreme, name=Proposition]{proposition}
  \ifdefined\TRACE \tracingmacros=0 \fi
  \declaretheorem[sibling=theoreme, name=Corollaire]{corollaire}
  \declaretheorem[numberwithin=chapter, name=Remarque]{remarque}
  \ifdefined\FAUX \newcounteralias{corollaire}{section}\newcounteralias{orphelin}{inexistant}\fi
  \ifdefined\CORR \input{correctif-a-remise}\fi
\fi
\crefname{theoreme}{théorème}{théorèmes}
\crefname{proposition}{proposition}{propositions}
\crefname{corollaire}{corollaire}{corollaires}
\crefname{remarque}{remarque}{remarques}
\begin{document}
\chapter{Premier}
\begin{proposition}\label{p1}a\end{proposition}
\begin{theoreme}\label{t1}b\end{theoreme}
\begin{remarque}\label{r1}c\end{remarque}
\begin{corollaire}\label{c1}d\end{corollaire}
\chapter{Second}
\begin{corollaire}\label{c2}e\end{corollaire}
\begin{theoreme}\label{t2}f\end{theoreme}
RENVOIS \ref{p1} \ref{t1} \ref{r1} \ref{c1} \ref{c2} \ref{t2} ; \cref{p1} ; \cref{t1} ; \cref{r1} ; \cref{c1} ; \cref{c2} ; \cref{t2} FIN
\end{document}
"                       correctif A du préambule
%   lualatex -jobname=crochet  "\def\SANSCROCHET{1}% Juge de la troisième passe : document minimal écrit de zéro pour la cause A (thmtools sous LaTeX 2026-06-01).
% Variantes, par une macro posée sur la ligne de commande (deux passes chacune, depuis memoire/essai/juge3/minimal/) :
%   lualatex -jobname=nu       "\input{minimal-juge3}"                                   sans correctif
%   lualatex -jobname=corr     "\def\CORR{1}\input{minimal-juge3}"                       correctif A du préambule
%   lualatex -jobname=crochet  "\def\SANSCROCHET{1}\input{minimal-juge3}"                sans correctif, crochet retiré
%   lualatex -jobname=amsseul  "\def\AMSSEUL{1}\input{minimal-juge3}"                    amsthm seul (témoin)
%   lualatex -jobname=faux     "\def\CORR{1}\def\FAUX{1}\input{minimal-juge3}"           deux vrais défauts : doivent tomber
%   lualatex -jobname=trace    "\def\TRACE{1}\input{minimal-juge3}"                      trace de l'appel
% Les lignes du correctif sont tirées du préambule commité par extraire_correctifs.py, pas retapées.
\documentclass{scrbook}
\ifdefined\SANSCROCHET \RemoveFromHook{package/amsthm/after}[firstaid/aliascounter]\fi
\usepackage{amsthm}
\usepackage{hyperref}
\usepackage[french, nameinlink, noabbrev]{cleveref}
\ifdefined\AMSSEUL
  \newtheorem{theoreme}{Théorème}[chapter]
  \newtheorem{proposition}[theoreme]{Proposition}
  \newtheorem{corollaire}[theoreme]{Corollaire}
  \newtheorem{remarque}{Remarque}[chapter]
\else
  \usepackage{thmtools}
  \ifdefined\CORR \input{correctif-a-pose}\fi
  \declaretheorem[numberwithin=chapter, name=Théorème]{theoreme}
  \ifdefined\TRACE \tracingmacros=2 \tracingonline=0 \fi
  \declaretheorem[sibling=theoreme, name=Proposition]{proposition}
  \ifdefined\TRACE \tracingmacros=0 \fi
  \declaretheorem[sibling=theoreme, name=Corollaire]{corollaire}
  \declaretheorem[numberwithin=chapter, name=Remarque]{remarque}
  \ifdefined\FAUX \newcounteralias{corollaire}{section}\newcounteralias{orphelin}{inexistant}\fi
  \ifdefined\CORR \input{correctif-a-remise}\fi
\fi
\crefname{theoreme}{théorème}{théorèmes}
\crefname{proposition}{proposition}{propositions}
\crefname{corollaire}{corollaire}{corollaires}
\crefname{remarque}{remarque}{remarques}
\begin{document}
\chapter{Premier}
\begin{proposition}\label{p1}a\end{proposition}
\begin{theoreme}\label{t1}b\end{theoreme}
\begin{remarque}\label{r1}c\end{remarque}
\begin{corollaire}\label{c1}d\end{corollaire}
\chapter{Second}
\begin{corollaire}\label{c2}e\end{corollaire}
\begin{theoreme}\label{t2}f\end{theoreme}
RENVOIS \ref{p1} \ref{t1} \ref{r1} \ref{c1} \ref{c2} \ref{t2} ; \cref{p1} ; \cref{t1} ; \cref{r1} ; \cref{c1} ; \cref{c2} ; \cref{t2} FIN
\end{document}
"                sans correctif, crochet retiré
%   lualatex -jobname=amsseul  "\def\AMSSEUL{1}% Juge de la troisième passe : document minimal écrit de zéro pour la cause A (thmtools sous LaTeX 2026-06-01).
% Variantes, par une macro posée sur la ligne de commande (deux passes chacune, depuis memoire/essai/juge3/minimal/) :
%   lualatex -jobname=nu       "\input{minimal-juge3}"                                   sans correctif
%   lualatex -jobname=corr     "\def\CORR{1}\input{minimal-juge3}"                       correctif A du préambule
%   lualatex -jobname=crochet  "\def\SANSCROCHET{1}\input{minimal-juge3}"                sans correctif, crochet retiré
%   lualatex -jobname=amsseul  "\def\AMSSEUL{1}\input{minimal-juge3}"                    amsthm seul (témoin)
%   lualatex -jobname=faux     "\def\CORR{1}\def\FAUX{1}\input{minimal-juge3}"           deux vrais défauts : doivent tomber
%   lualatex -jobname=trace    "\def\TRACE{1}\input{minimal-juge3}"                      trace de l'appel
% Les lignes du correctif sont tirées du préambule commité par extraire_correctifs.py, pas retapées.
\documentclass{scrbook}
\ifdefined\SANSCROCHET \RemoveFromHook{package/amsthm/after}[firstaid/aliascounter]\fi
\usepackage{amsthm}
\usepackage{hyperref}
\usepackage[french, nameinlink, noabbrev]{cleveref}
\ifdefined\AMSSEUL
  \newtheorem{theoreme}{Théorème}[chapter]
  \newtheorem{proposition}[theoreme]{Proposition}
  \newtheorem{corollaire}[theoreme]{Corollaire}
  \newtheorem{remarque}{Remarque}[chapter]
\else
  \usepackage{thmtools}
  \ifdefined\CORR \input{correctif-a-pose}\fi
  \declaretheorem[numberwithin=chapter, name=Théorème]{theoreme}
  \ifdefined\TRACE \tracingmacros=2 \tracingonline=0 \fi
  \declaretheorem[sibling=theoreme, name=Proposition]{proposition}
  \ifdefined\TRACE \tracingmacros=0 \fi
  \declaretheorem[sibling=theoreme, name=Corollaire]{corollaire}
  \declaretheorem[numberwithin=chapter, name=Remarque]{remarque}
  \ifdefined\FAUX \newcounteralias{corollaire}{section}\newcounteralias{orphelin}{inexistant}\fi
  \ifdefined\CORR \input{correctif-a-remise}\fi
\fi
\crefname{theoreme}{théorème}{théorèmes}
\crefname{proposition}{proposition}{propositions}
\crefname{corollaire}{corollaire}{corollaires}
\crefname{remarque}{remarque}{remarques}
\begin{document}
\chapter{Premier}
\begin{proposition}\label{p1}a\end{proposition}
\begin{theoreme}\label{t1}b\end{theoreme}
\begin{remarque}\label{r1}c\end{remarque}
\begin{corollaire}\label{c1}d\end{corollaire}
\chapter{Second}
\begin{corollaire}\label{c2}e\end{corollaire}
\begin{theoreme}\label{t2}f\end{theoreme}
RENVOIS \ref{p1} \ref{t1} \ref{r1} \ref{c1} \ref{c2} \ref{t2} ; \cref{p1} ; \cref{t1} ; \cref{r1} ; \cref{c1} ; \cref{c2} ; \cref{t2} FIN
\end{document}
"                    amsthm seul (témoin)
%   lualatex -jobname=faux     "\def\CORR{1}\def\FAUX{1}% Juge de la troisième passe : document minimal écrit de zéro pour la cause A (thmtools sous LaTeX 2026-06-01).
% Variantes, par une macro posée sur la ligne de commande (deux passes chacune, depuis memoire/essai/juge3/minimal/) :
%   lualatex -jobname=nu       "\input{minimal-juge3}"                                   sans correctif
%   lualatex -jobname=corr     "\def\CORR{1}\input{minimal-juge3}"                       correctif A du préambule
%   lualatex -jobname=crochet  "\def\SANSCROCHET{1}\input{minimal-juge3}"                sans correctif, crochet retiré
%   lualatex -jobname=amsseul  "\def\AMSSEUL{1}\input{minimal-juge3}"                    amsthm seul (témoin)
%   lualatex -jobname=faux     "\def\CORR{1}\def\FAUX{1}\input{minimal-juge3}"           deux vrais défauts : doivent tomber
%   lualatex -jobname=trace    "\def\TRACE{1}\input{minimal-juge3}"                      trace de l'appel
% Les lignes du correctif sont tirées du préambule commité par extraire_correctifs.py, pas retapées.
\documentclass{scrbook}
\ifdefined\SANSCROCHET \RemoveFromHook{package/amsthm/after}[firstaid/aliascounter]\fi
\usepackage{amsthm}
\usepackage{hyperref}
\usepackage[french, nameinlink, noabbrev]{cleveref}
\ifdefined\AMSSEUL
  \newtheorem{theoreme}{Théorème}[chapter]
  \newtheorem{proposition}[theoreme]{Proposition}
  \newtheorem{corollaire}[theoreme]{Corollaire}
  \newtheorem{remarque}{Remarque}[chapter]
\else
  \usepackage{thmtools}
  \ifdefined\CORR \input{correctif-a-pose}\fi
  \declaretheorem[numberwithin=chapter, name=Théorème]{theoreme}
  \ifdefined\TRACE \tracingmacros=2 \tracingonline=0 \fi
  \declaretheorem[sibling=theoreme, name=Proposition]{proposition}
  \ifdefined\TRACE \tracingmacros=0 \fi
  \declaretheorem[sibling=theoreme, name=Corollaire]{corollaire}
  \declaretheorem[numberwithin=chapter, name=Remarque]{remarque}
  \ifdefined\FAUX \newcounteralias{corollaire}{section}\newcounteralias{orphelin}{inexistant}\fi
  \ifdefined\CORR \input{correctif-a-remise}\fi
\fi
\crefname{theoreme}{théorème}{théorèmes}
\crefname{proposition}{proposition}{propositions}
\crefname{corollaire}{corollaire}{corollaires}
\crefname{remarque}{remarque}{remarques}
\begin{document}
\chapter{Premier}
\begin{proposition}\label{p1}a\end{proposition}
\begin{theoreme}\label{t1}b\end{theoreme}
\begin{remarque}\label{r1}c\end{remarque}
\begin{corollaire}\label{c1}d\end{corollaire}
\chapter{Second}
\begin{corollaire}\label{c2}e\end{corollaire}
\begin{theoreme}\label{t2}f\end{theoreme}
RENVOIS \ref{p1} \ref{t1} \ref{r1} \ref{c1} \ref{c2} \ref{t2} ; \cref{p1} ; \cref{t1} ; \cref{r1} ; \cref{c1} ; \cref{c2} ; \cref{t2} FIN
\end{document}
"           deux vrais défauts : doivent tomber
%   lualatex -jobname=trace    "\def\TRACE{1}% Juge de la troisième passe : document minimal écrit de zéro pour la cause A (thmtools sous LaTeX 2026-06-01).
% Variantes, par une macro posée sur la ligne de commande (deux passes chacune, depuis memoire/essai/juge3/minimal/) :
%   lualatex -jobname=nu       "\input{minimal-juge3}"                                   sans correctif
%   lualatex -jobname=corr     "\def\CORR{1}\input{minimal-juge3}"                       correctif A du préambule
%   lualatex -jobname=crochet  "\def\SANSCROCHET{1}\input{minimal-juge3}"                sans correctif, crochet retiré
%   lualatex -jobname=amsseul  "\def\AMSSEUL{1}\input{minimal-juge3}"                    amsthm seul (témoin)
%   lualatex -jobname=faux     "\def\CORR{1}\def\FAUX{1}\input{minimal-juge3}"           deux vrais défauts : doivent tomber
%   lualatex -jobname=trace    "\def\TRACE{1}\input{minimal-juge3}"                      trace de l'appel
% Les lignes du correctif sont tirées du préambule commité par extraire_correctifs.py, pas retapées.
\documentclass{scrbook}
\ifdefined\SANSCROCHET \RemoveFromHook{package/amsthm/after}[firstaid/aliascounter]\fi
\usepackage{amsthm}
\usepackage{hyperref}
\usepackage[french, nameinlink, noabbrev]{cleveref}
\ifdefined\AMSSEUL
  \newtheorem{theoreme}{Théorème}[chapter]
  \newtheorem{proposition}[theoreme]{Proposition}
  \newtheorem{corollaire}[theoreme]{Corollaire}
  \newtheorem{remarque}{Remarque}[chapter]
\else
  \usepackage{thmtools}
  \ifdefined\CORR \input{correctif-a-pose}\fi
  \declaretheorem[numberwithin=chapter, name=Théorème]{theoreme}
  \ifdefined\TRACE \tracingmacros=2 \tracingonline=0 \fi
  \declaretheorem[sibling=theoreme, name=Proposition]{proposition}
  \ifdefined\TRACE \tracingmacros=0 \fi
  \declaretheorem[sibling=theoreme, name=Corollaire]{corollaire}
  \declaretheorem[numberwithin=chapter, name=Remarque]{remarque}
  \ifdefined\FAUX \newcounteralias{corollaire}{section}\newcounteralias{orphelin}{inexistant}\fi
  \ifdefined\CORR \input{correctif-a-remise}\fi
\fi
\crefname{theoreme}{théorème}{théorèmes}
\crefname{proposition}{proposition}{propositions}
\crefname{corollaire}{corollaire}{corollaires}
\crefname{remarque}{remarque}{remarques}
\begin{document}
\chapter{Premier}
\begin{proposition}\label{p1}a\end{proposition}
\begin{theoreme}\label{t1}b\end{theoreme}
\begin{remarque}\label{r1}c\end{remarque}
\begin{corollaire}\label{c1}d\end{corollaire}
\chapter{Second}
\begin{corollaire}\label{c2}e\end{corollaire}
\begin{theoreme}\label{t2}f\end{theoreme}
RENVOIS \ref{p1} \ref{t1} \ref{r1} \ref{c1} \ref{c2} \ref{t2} ; \cref{p1} ; \cref{t1} ; \cref{r1} ; \cref{c1} ; \cref{c2} ; \cref{t2} FIN
\end{document}
"                      trace de l'appel
% Les lignes du correctif sont tirées du préambule commité par extraire_correctifs.py, pas retapées.
\documentclass{scrbook}
\ifdefined\SANSCROCHET \RemoveFromHook{package/amsthm/after}[firstaid/aliascounter]\fi
\usepackage{amsthm}
\usepackage{hyperref}
\usepackage[french, nameinlink, noabbrev]{cleveref}
\ifdefined\AMSSEUL
  \newtheorem{theoreme}{Théorème}[chapter]
  \newtheorem{proposition}[theoreme]{Proposition}
  \newtheorem{corollaire}[theoreme]{Corollaire}
  \newtheorem{remarque}{Remarque}[chapter]
\else
  \usepackage{thmtools}
  \ifdefined\CORR \input{correctif-a-pose}\fi
  \declaretheorem[numberwithin=chapter, name=Théorème]{theoreme}
  \ifdefined\TRACE \tracingmacros=2 \tracingonline=0 \fi
  \declaretheorem[sibling=theoreme, name=Proposition]{proposition}
  \ifdefined\TRACE \tracingmacros=0 \fi
  \declaretheorem[sibling=theoreme, name=Corollaire]{corollaire}
  \declaretheorem[numberwithin=chapter, name=Remarque]{remarque}
  \ifdefined\FAUX \newcounteralias{corollaire}{section}\newcounteralias{orphelin}{inexistant}\fi
  \ifdefined\CORR \input{correctif-a-remise}\fi
\fi
\crefname{theoreme}{théorème}{théorèmes}
\crefname{proposition}{proposition}{propositions}
\crefname{corollaire}{corollaire}{corollaires}
\crefname{remarque}{remarque}{remarques}
\begin{document}
\chapter{Premier}
\begin{proposition}\label{p1}a\end{proposition}
\begin{theoreme}\label{t1}b\end{theoreme}
\begin{remarque}\label{r1}c\end{remarque}
\begin{corollaire}\label{c1}d\end{corollaire}
\chapter{Second}
\begin{corollaire}\label{c2}e\end{corollaire}
\begin{theoreme}\label{t2}f\end{theoreme}
RENVOIS \ref{p1} \ref{t1} \ref{r1} \ref{c1} \ref{c2} \ref{t2} ; \cref{p1} ; \cref{t1} ; \cref{r1} ; \cref{c1} ; \cref{c2} ; \cref{t2} FIN
\end{document}
"                      trace de l'appel
% Les lignes du correctif sont tirées du préambule commité par extraire_correctifs.py, pas retapées.
\documentclass{scrbook}
\ifdefined\SANSCROCHET \RemoveFromHook{package/amsthm/after}[firstaid/aliascounter]\fi
\usepackage{amsthm}
\usepackage{hyperref}
\usepackage[french, nameinlink, noabbrev]{cleveref}
\ifdefined\AMSSEUL
  \newtheorem{theoreme}{Théorème}[chapter]
  \newtheorem{proposition}[theoreme]{Proposition}
  \newtheorem{corollaire}[theoreme]{Corollaire}
  \newtheorem{remarque}{Remarque}[chapter]
\else
  \usepackage{thmtools}
  \ifdefined\CORR \input{correctif-a-pose}\fi
  \declaretheorem[numberwithin=chapter, name=Théorème]{theoreme}
  \ifdefined\TRACE \tracingmacros=2 \tracingonline=0 \fi
  \declaretheorem[sibling=theoreme, name=Proposition]{proposition}
  \ifdefined\TRACE \tracingmacros=0 \fi
  \declaretheorem[sibling=theoreme, name=Corollaire]{corollaire}
  \declaretheorem[numberwithin=chapter, name=Remarque]{remarque}
  \ifdefined\FAUX \newcounteralias{corollaire}{section}\newcounteralias{orphelin}{inexistant}\fi
  \ifdefined\CORR \input{correctif-a-remise}\fi
\fi
\crefname{theoreme}{théorème}{théorèmes}
\crefname{proposition}{proposition}{propositions}
\crefname{corollaire}{corollaire}{corollaires}
\crefname{remarque}{remarque}{remarques}
\begin{document}
\chapter{Premier}
\begin{proposition}\label{p1}a\end{proposition}
\begin{theoreme}\label{t1}b\end{theoreme}
\begin{remarque}\label{r1}c\end{remarque}
\begin{corollaire}\label{c1}d\end{corollaire}
\chapter{Second}
\begin{corollaire}\label{c2}e\end{corollaire}
\begin{theoreme}\label{t2}f\end{theoreme}
RENVOIS \ref{p1} \ref{t1} \ref{r1} \ref{c1} \ref{c2} \ref{t2} ; \cref{p1} ; \cref{t1} ; \cref{r1} ; \cref{c1} ; \cref{c2} ; \cref{t2} FIN
\end{document}
"                       correctif A du préambule
%   lualatex -jobname=crochet  "\def\SANSCROCHET{1}% Juge de la troisième passe : document minimal écrit de zéro pour la cause A (thmtools sous LaTeX 2026-06-01).
% Variantes, par une macro posée sur la ligne de commande (deux passes chacune, depuis memoire/essai/juge3/minimal/) :
%   lualatex -jobname=nu       "% Juge de la troisième passe : document minimal écrit de zéro pour la cause A (thmtools sous LaTeX 2026-06-01).
% Variantes, par une macro posée sur la ligne de commande (deux passes chacune, depuis memoire/essai/juge3/minimal/) :
%   lualatex -jobname=nu       "% Juge de la troisième passe : document minimal écrit de zéro pour la cause A (thmtools sous LaTeX 2026-06-01).
% Variantes, par une macro posée sur la ligne de commande (deux passes chacune, depuis memoire/essai/juge3/minimal/) :
%   lualatex -jobname=nu       "\input{minimal-juge3}"                                   sans correctif
%   lualatex -jobname=corr     "\def\CORR{1}\input{minimal-juge3}"                       correctif A du préambule
%   lualatex -jobname=crochet  "\def\SANSCROCHET{1}\input{minimal-juge3}"                sans correctif, crochet retiré
%   lualatex -jobname=amsseul  "\def\AMSSEUL{1}\input{minimal-juge3}"                    amsthm seul (témoin)
%   lualatex -jobname=faux     "\def\CORR{1}\def\FAUX{1}\input{minimal-juge3}"           deux vrais défauts : doivent tomber
%   lualatex -jobname=trace    "\def\TRACE{1}\input{minimal-juge3}"                      trace de l'appel
% Les lignes du correctif sont tirées du préambule commité par extraire_correctifs.py, pas retapées.
\documentclass{scrbook}
\ifdefined\SANSCROCHET \RemoveFromHook{package/amsthm/after}[firstaid/aliascounter]\fi
\usepackage{amsthm}
\usepackage{hyperref}
\usepackage[french, nameinlink, noabbrev]{cleveref}
\ifdefined\AMSSEUL
  \newtheorem{theoreme}{Théorème}[chapter]
  \newtheorem{proposition}[theoreme]{Proposition}
  \newtheorem{corollaire}[theoreme]{Corollaire}
  \newtheorem{remarque}{Remarque}[chapter]
\else
  \usepackage{thmtools}
  \ifdefined\CORR \input{correctif-a-pose}\fi
  \declaretheorem[numberwithin=chapter, name=Théorème]{theoreme}
  \ifdefined\TRACE \tracingmacros=2 \tracingonline=0 \fi
  \declaretheorem[sibling=theoreme, name=Proposition]{proposition}
  \ifdefined\TRACE \tracingmacros=0 \fi
  \declaretheorem[sibling=theoreme, name=Corollaire]{corollaire}
  \declaretheorem[numberwithin=chapter, name=Remarque]{remarque}
  \ifdefined\FAUX \newcounteralias{corollaire}{section}\newcounteralias{orphelin}{inexistant}\fi
  \ifdefined\CORR \input{correctif-a-remise}\fi
\fi
\crefname{theoreme}{théorème}{théorèmes}
\crefname{proposition}{proposition}{propositions}
\crefname{corollaire}{corollaire}{corollaires}
\crefname{remarque}{remarque}{remarques}
\begin{document}
\chapter{Premier}
\begin{proposition}\label{p1}a\end{proposition}
\begin{theoreme}\label{t1}b\end{theoreme}
\begin{remarque}\label{r1}c\end{remarque}
\begin{corollaire}\label{c1}d\end{corollaire}
\chapter{Second}
\begin{corollaire}\label{c2}e\end{corollaire}
\begin{theoreme}\label{t2}f\end{theoreme}
RENVOIS \ref{p1} \ref{t1} \ref{r1} \ref{c1} \ref{c2} \ref{t2} ; \cref{p1} ; \cref{t1} ; \cref{r1} ; \cref{c1} ; \cref{c2} ; \cref{t2} FIN
\end{document}
"                                   sans correctif
%   lualatex -jobname=corr     "\def\CORR{1}% Juge de la troisième passe : document minimal écrit de zéro pour la cause A (thmtools sous LaTeX 2026-06-01).
% Variantes, par une macro posée sur la ligne de commande (deux passes chacune, depuis memoire/essai/juge3/minimal/) :
%   lualatex -jobname=nu       "\input{minimal-juge3}"                                   sans correctif
%   lualatex -jobname=corr     "\def\CORR{1}\input{minimal-juge3}"                       correctif A du préambule
%   lualatex -jobname=crochet  "\def\SANSCROCHET{1}\input{minimal-juge3}"                sans correctif, crochet retiré
%   lualatex -jobname=amsseul  "\def\AMSSEUL{1}\input{minimal-juge3}"                    amsthm seul (témoin)
%   lualatex -jobname=faux     "\def\CORR{1}\def\FAUX{1}\input{minimal-juge3}"           deux vrais défauts : doivent tomber
%   lualatex -jobname=trace    "\def\TRACE{1}\input{minimal-juge3}"                      trace de l'appel
% Les lignes du correctif sont tirées du préambule commité par extraire_correctifs.py, pas retapées.
\documentclass{scrbook}
\ifdefined\SANSCROCHET \RemoveFromHook{package/amsthm/after}[firstaid/aliascounter]\fi
\usepackage{amsthm}
\usepackage{hyperref}
\usepackage[french, nameinlink, noabbrev]{cleveref}
\ifdefined\AMSSEUL
  \newtheorem{theoreme}{Théorème}[chapter]
  \newtheorem{proposition}[theoreme]{Proposition}
  \newtheorem{corollaire}[theoreme]{Corollaire}
  \newtheorem{remarque}{Remarque}[chapter]
\else
  \usepackage{thmtools}
  \ifdefined\CORR \input{correctif-a-pose}\fi
  \declaretheorem[numberwithin=chapter, name=Théorème]{theoreme}
  \ifdefined\TRACE \tracingmacros=2 \tracingonline=0 \fi
  \declaretheorem[sibling=theoreme, name=Proposition]{proposition}
  \ifdefined\TRACE \tracingmacros=0 \fi
  \declaretheorem[sibling=theoreme, name=Corollaire]{corollaire}
  \declaretheorem[numberwithin=chapter, name=Remarque]{remarque}
  \ifdefined\FAUX \newcounteralias{corollaire}{section}\newcounteralias{orphelin}{inexistant}\fi
  \ifdefined\CORR \input{correctif-a-remise}\fi
\fi
\crefname{theoreme}{théorème}{théorèmes}
\crefname{proposition}{proposition}{propositions}
\crefname{corollaire}{corollaire}{corollaires}
\crefname{remarque}{remarque}{remarques}
\begin{document}
\chapter{Premier}
\begin{proposition}\label{p1}a\end{proposition}
\begin{theoreme}\label{t1}b\end{theoreme}
\begin{remarque}\label{r1}c\end{remarque}
\begin{corollaire}\label{c1}d\end{corollaire}
\chapter{Second}
\begin{corollaire}\label{c2}e\end{corollaire}
\begin{theoreme}\label{t2}f\end{theoreme}
RENVOIS \ref{p1} \ref{t1} \ref{r1} \ref{c1} \ref{c2} \ref{t2} ; \cref{p1} ; \cref{t1} ; \cref{r1} ; \cref{c1} ; \cref{c2} ; \cref{t2} FIN
\end{document}
"                       correctif A du préambule
%   lualatex -jobname=crochet  "\def\SANSCROCHET{1}% Juge de la troisième passe : document minimal écrit de zéro pour la cause A (thmtools sous LaTeX 2026-06-01).
% Variantes, par une macro posée sur la ligne de commande (deux passes chacune, depuis memoire/essai/juge3/minimal/) :
%   lualatex -jobname=nu       "\input{minimal-juge3}"                                   sans correctif
%   lualatex -jobname=corr     "\def\CORR{1}\input{minimal-juge3}"                       correctif A du préambule
%   lualatex -jobname=crochet  "\def\SANSCROCHET{1}\input{minimal-juge3}"                sans correctif, crochet retiré
%   lualatex -jobname=amsseul  "\def\AMSSEUL{1}\input{minimal-juge3}"                    amsthm seul (témoin)
%   lualatex -jobname=faux     "\def\CORR{1}\def\FAUX{1}\input{minimal-juge3}"           deux vrais défauts : doivent tomber
%   lualatex -jobname=trace    "\def\TRACE{1}\input{minimal-juge3}"                      trace de l'appel
% Les lignes du correctif sont tirées du préambule commité par extraire_correctifs.py, pas retapées.
\documentclass{scrbook}
\ifdefined\SANSCROCHET \RemoveFromHook{package/amsthm/after}[firstaid/aliascounter]\fi
\usepackage{amsthm}
\usepackage{hyperref}
\usepackage[french, nameinlink, noabbrev]{cleveref}
\ifdefined\AMSSEUL
  \newtheorem{theoreme}{Théorème}[chapter]
  \newtheorem{proposition}[theoreme]{Proposition}
  \newtheorem{corollaire}[theoreme]{Corollaire}
  \newtheorem{remarque}{Remarque}[chapter]
\else
  \usepackage{thmtools}
  \ifdefined\CORR \input{correctif-a-pose}\fi
  \declaretheorem[numberwithin=chapter, name=Théorème]{theoreme}
  \ifdefined\TRACE \tracingmacros=2 \tracingonline=0 \fi
  \declaretheorem[sibling=theoreme, name=Proposition]{proposition}
  \ifdefined\TRACE \tracingmacros=0 \fi
  \declaretheorem[sibling=theoreme, name=Corollaire]{corollaire}
  \declaretheorem[numberwithin=chapter, name=Remarque]{remarque}
  \ifdefined\FAUX \newcounteralias{corollaire}{section}\newcounteralias{orphelin}{inexistant}\fi
  \ifdefined\CORR \input{correctif-a-remise}\fi
\fi
\crefname{theoreme}{théorème}{théorèmes}
\crefname{proposition}{proposition}{propositions}
\crefname{corollaire}{corollaire}{corollaires}
\crefname{remarque}{remarque}{remarques}
\begin{document}
\chapter{Premier}
\begin{proposition}\label{p1}a\end{proposition}
\begin{theoreme}\label{t1}b\end{theoreme}
\begin{remarque}\label{r1}c\end{remarque}
\begin{corollaire}\label{c1}d\end{corollaire}
\chapter{Second}
\begin{corollaire}\label{c2}e\end{corollaire}
\begin{theoreme}\label{t2}f\end{theoreme}
RENVOIS \ref{p1} \ref{t1} \ref{r1} \ref{c1} \ref{c2} \ref{t2} ; \cref{p1} ; \cref{t1} ; \cref{r1} ; \cref{c1} ; \cref{c2} ; \cref{t2} FIN
\end{document}
"                sans correctif, crochet retiré
%   lualatex -jobname=amsseul  "\def\AMSSEUL{1}% Juge de la troisième passe : document minimal écrit de zéro pour la cause A (thmtools sous LaTeX 2026-06-01).
% Variantes, par une macro posée sur la ligne de commande (deux passes chacune, depuis memoire/essai/juge3/minimal/) :
%   lualatex -jobname=nu       "\input{minimal-juge3}"                                   sans correctif
%   lualatex -jobname=corr     "\def\CORR{1}\input{minimal-juge3}"                       correctif A du préambule
%   lualatex -jobname=crochet  "\def\SANSCROCHET{1}\input{minimal-juge3}"                sans correctif, crochet retiré
%   lualatex -jobname=amsseul  "\def\AMSSEUL{1}\input{minimal-juge3}"                    amsthm seul (témoin)
%   lualatex -jobname=faux     "\def\CORR{1}\def\FAUX{1}\input{minimal-juge3}"           deux vrais défauts : doivent tomber
%   lualatex -jobname=trace    "\def\TRACE{1}\input{minimal-juge3}"                      trace de l'appel
% Les lignes du correctif sont tirées du préambule commité par extraire_correctifs.py, pas retapées.
\documentclass{scrbook}
\ifdefined\SANSCROCHET \RemoveFromHook{package/amsthm/after}[firstaid/aliascounter]\fi
\usepackage{amsthm}
\usepackage{hyperref}
\usepackage[french, nameinlink, noabbrev]{cleveref}
\ifdefined\AMSSEUL
  \newtheorem{theoreme}{Théorème}[chapter]
  \newtheorem{proposition}[theoreme]{Proposition}
  \newtheorem{corollaire}[theoreme]{Corollaire}
  \newtheorem{remarque}{Remarque}[chapter]
\else
  \usepackage{thmtools}
  \ifdefined\CORR \input{correctif-a-pose}\fi
  \declaretheorem[numberwithin=chapter, name=Théorème]{theoreme}
  \ifdefined\TRACE \tracingmacros=2 \tracingonline=0 \fi
  \declaretheorem[sibling=theoreme, name=Proposition]{proposition}
  \ifdefined\TRACE \tracingmacros=0 \fi
  \declaretheorem[sibling=theoreme, name=Corollaire]{corollaire}
  \declaretheorem[numberwithin=chapter, name=Remarque]{remarque}
  \ifdefined\FAUX \newcounteralias{corollaire}{section}\newcounteralias{orphelin}{inexistant}\fi
  \ifdefined\CORR \input{correctif-a-remise}\fi
\fi
\crefname{theoreme}{théorème}{théorèmes}
\crefname{proposition}{proposition}{propositions}
\crefname{corollaire}{corollaire}{corollaires}
\crefname{remarque}{remarque}{remarques}
\begin{document}
\chapter{Premier}
\begin{proposition}\label{p1}a\end{proposition}
\begin{theoreme}\label{t1}b\end{theoreme}
\begin{remarque}\label{r1}c\end{remarque}
\begin{corollaire}\label{c1}d\end{corollaire}
\chapter{Second}
\begin{corollaire}\label{c2}e\end{corollaire}
\begin{theoreme}\label{t2}f\end{theoreme}
RENVOIS \ref{p1} \ref{t1} \ref{r1} \ref{c1} \ref{c2} \ref{t2} ; \cref{p1} ; \cref{t1} ; \cref{r1} ; \cref{c1} ; \cref{c2} ; \cref{t2} FIN
\end{document}
"                    amsthm seul (témoin)
%   lualatex -jobname=faux     "\def\CORR{1}\def\FAUX{1}% Juge de la troisième passe : document minimal écrit de zéro pour la cause A (thmtools sous LaTeX 2026-06-01).
% Variantes, par une macro posée sur la ligne de commande (deux passes chacune, depuis memoire/essai/juge3/minimal/) :
%   lualatex -jobname=nu       "\input{minimal-juge3}"                                   sans correctif
%   lualatex -jobname=corr     "\def\CORR{1}\input{minimal-juge3}"                       correctif A du préambule
%   lualatex -jobname=crochet  "\def\SANSCROCHET{1}\input{minimal-juge3}"                sans correctif, crochet retiré
%   lualatex -jobname=amsseul  "\def\AMSSEUL{1}\input{minimal-juge3}"                    amsthm seul (témoin)
%   lualatex -jobname=faux     "\def\CORR{1}\def\FAUX{1}\input{minimal-juge3}"           deux vrais défauts : doivent tomber
%   lualatex -jobname=trace    "\def\TRACE{1}\input{minimal-juge3}"                      trace de l'appel
% Les lignes du correctif sont tirées du préambule commité par extraire_correctifs.py, pas retapées.
\documentclass{scrbook}
\ifdefined\SANSCROCHET \RemoveFromHook{package/amsthm/after}[firstaid/aliascounter]\fi
\usepackage{amsthm}
\usepackage{hyperref}
\usepackage[french, nameinlink, noabbrev]{cleveref}
\ifdefined\AMSSEUL
  \newtheorem{theoreme}{Théorème}[chapter]
  \newtheorem{proposition}[theoreme]{Proposition}
  \newtheorem{corollaire}[theoreme]{Corollaire}
  \newtheorem{remarque}{Remarque}[chapter]
\else
  \usepackage{thmtools}
  \ifdefined\CORR \input{correctif-a-pose}\fi
  \declaretheorem[numberwithin=chapter, name=Théorème]{theoreme}
  \ifdefined\TRACE \tracingmacros=2 \tracingonline=0 \fi
  \declaretheorem[sibling=theoreme, name=Proposition]{proposition}
  \ifdefined\TRACE \tracingmacros=0 \fi
  \declaretheorem[sibling=theoreme, name=Corollaire]{corollaire}
  \declaretheorem[numberwithin=chapter, name=Remarque]{remarque}
  \ifdefined\FAUX \newcounteralias{corollaire}{section}\newcounteralias{orphelin}{inexistant}\fi
  \ifdefined\CORR \input{correctif-a-remise}\fi
\fi
\crefname{theoreme}{théorème}{théorèmes}
\crefname{proposition}{proposition}{propositions}
\crefname{corollaire}{corollaire}{corollaires}
\crefname{remarque}{remarque}{remarques}
\begin{document}
\chapter{Premier}
\begin{proposition}\label{p1}a\end{proposition}
\begin{theoreme}\label{t1}b\end{theoreme}
\begin{remarque}\label{r1}c\end{remarque}
\begin{corollaire}\label{c1}d\end{corollaire}
\chapter{Second}
\begin{corollaire}\label{c2}e\end{corollaire}
\begin{theoreme}\label{t2}f\end{theoreme}
RENVOIS \ref{p1} \ref{t1} \ref{r1} \ref{c1} \ref{c2} \ref{t2} ; \cref{p1} ; \cref{t1} ; \cref{r1} ; \cref{c1} ; \cref{c2} ; \cref{t2} FIN
\end{document}
"           deux vrais défauts : doivent tomber
%   lualatex -jobname=trace    "\def\TRACE{1}% Juge de la troisième passe : document minimal écrit de zéro pour la cause A (thmtools sous LaTeX 2026-06-01).
% Variantes, par une macro posée sur la ligne de commande (deux passes chacune, depuis memoire/essai/juge3/minimal/) :
%   lualatex -jobname=nu       "\input{minimal-juge3}"                                   sans correctif
%   lualatex -jobname=corr     "\def\CORR{1}\input{minimal-juge3}"                       correctif A du préambule
%   lualatex -jobname=crochet  "\def\SANSCROCHET{1}\input{minimal-juge3}"                sans correctif, crochet retiré
%   lualatex -jobname=amsseul  "\def\AMSSEUL{1}\input{minimal-juge3}"                    amsthm seul (témoin)
%   lualatex -jobname=faux     "\def\CORR{1}\def\FAUX{1}\input{minimal-juge3}"           deux vrais défauts : doivent tomber
%   lualatex -jobname=trace    "\def\TRACE{1}\input{minimal-juge3}"                      trace de l'appel
% Les lignes du correctif sont tirées du préambule commité par extraire_correctifs.py, pas retapées.
\documentclass{scrbook}
\ifdefined\SANSCROCHET \RemoveFromHook{package/amsthm/after}[firstaid/aliascounter]\fi
\usepackage{amsthm}
\usepackage{hyperref}
\usepackage[french, nameinlink, noabbrev]{cleveref}
\ifdefined\AMSSEUL
  \newtheorem{theoreme}{Théorème}[chapter]
  \newtheorem{proposition}[theoreme]{Proposition}
  \newtheorem{corollaire}[theoreme]{Corollaire}
  \newtheorem{remarque}{Remarque}[chapter]
\else
  \usepackage{thmtools}
  \ifdefined\CORR \input{correctif-a-pose}\fi
  \declaretheorem[numberwithin=chapter, name=Théorème]{theoreme}
  \ifdefined\TRACE \tracingmacros=2 \tracingonline=0 \fi
  \declaretheorem[sibling=theoreme, name=Proposition]{proposition}
  \ifdefined\TRACE \tracingmacros=0 \fi
  \declaretheorem[sibling=theoreme, name=Corollaire]{corollaire}
  \declaretheorem[numberwithin=chapter, name=Remarque]{remarque}
  \ifdefined\FAUX \newcounteralias{corollaire}{section}\newcounteralias{orphelin}{inexistant}\fi
  \ifdefined\CORR \input{correctif-a-remise}\fi
\fi
\crefname{theoreme}{théorème}{théorèmes}
\crefname{proposition}{proposition}{propositions}
\crefname{corollaire}{corollaire}{corollaires}
\crefname{remarque}{remarque}{remarques}
\begin{document}
\chapter{Premier}
\begin{proposition}\label{p1}a\end{proposition}
\begin{theoreme}\label{t1}b\end{theoreme}
\begin{remarque}\label{r1}c\end{remarque}
\begin{corollaire}\label{c1}d\end{corollaire}
\chapter{Second}
\begin{corollaire}\label{c2}e\end{corollaire}
\begin{theoreme}\label{t2}f\end{theoreme}
RENVOIS \ref{p1} \ref{t1} \ref{r1} \ref{c1} \ref{c2} \ref{t2} ; \cref{p1} ; \cref{t1} ; \cref{r1} ; \cref{c1} ; \cref{c2} ; \cref{t2} FIN
\end{document}
"                      trace de l'appel
% Les lignes du correctif sont tirées du préambule commité par extraire_correctifs.py, pas retapées.
\documentclass{scrbook}
\ifdefined\SANSCROCHET \RemoveFromHook{package/amsthm/after}[firstaid/aliascounter]\fi
\usepackage{amsthm}
\usepackage{hyperref}
\usepackage[french, nameinlink, noabbrev]{cleveref}
\ifdefined\AMSSEUL
  \newtheorem{theoreme}{Théorème}[chapter]
  \newtheorem{proposition}[theoreme]{Proposition}
  \newtheorem{corollaire}[theoreme]{Corollaire}
  \newtheorem{remarque}{Remarque}[chapter]
\else
  \usepackage{thmtools}
  \ifdefined\CORR \input{correctif-a-pose}\fi
  \declaretheorem[numberwithin=chapter, name=Théorème]{theoreme}
  \ifdefined\TRACE \tracingmacros=2 \tracingonline=0 \fi
  \declaretheorem[sibling=theoreme, name=Proposition]{proposition}
  \ifdefined\TRACE \tracingmacros=0 \fi
  \declaretheorem[sibling=theoreme, name=Corollaire]{corollaire}
  \declaretheorem[numberwithin=chapter, name=Remarque]{remarque}
  \ifdefined\FAUX \newcounteralias{corollaire}{section}\newcounteralias{orphelin}{inexistant}\fi
  \ifdefined\CORR \input{correctif-a-remise}\fi
\fi
\crefname{theoreme}{théorème}{théorèmes}
\crefname{proposition}{proposition}{propositions}
\crefname{corollaire}{corollaire}{corollaires}
\crefname{remarque}{remarque}{remarques}
\begin{document}
\chapter{Premier}
\begin{proposition}\label{p1}a\end{proposition}
\begin{theoreme}\label{t1}b\end{theoreme}
\begin{remarque}\label{r1}c\end{remarque}
\begin{corollaire}\label{c1}d\end{corollaire}
\chapter{Second}
\begin{corollaire}\label{c2}e\end{corollaire}
\begin{theoreme}\label{t2}f\end{theoreme}
RENVOIS \ref{p1} \ref{t1} \ref{r1} \ref{c1} \ref{c2} \ref{t2} ; \cref{p1} ; \cref{t1} ; \cref{r1} ; \cref{c1} ; \cref{c2} ; \cref{t2} FIN
\end{document}
"                                   sans correctif
%   lualatex -jobname=corr     "\def\CORR{1}% Juge de la troisième passe : document minimal écrit de zéro pour la cause A (thmtools sous LaTeX 2026-06-01).
% Variantes, par une macro posée sur la ligne de commande (deux passes chacune, depuis memoire/essai/juge3/minimal/) :
%   lualatex -jobname=nu       "% Juge de la troisième passe : document minimal écrit de zéro pour la cause A (thmtools sous LaTeX 2026-06-01).
% Variantes, par une macro posée sur la ligne de commande (deux passes chacune, depuis memoire/essai/juge3/minimal/) :
%   lualatex -jobname=nu       "\input{minimal-juge3}"                                   sans correctif
%   lualatex -jobname=corr     "\def\CORR{1}\input{minimal-juge3}"                       correctif A du préambule
%   lualatex -jobname=crochet  "\def\SANSCROCHET{1}\input{minimal-juge3}"                sans correctif, crochet retiré
%   lualatex -jobname=amsseul  "\def\AMSSEUL{1}\input{minimal-juge3}"                    amsthm seul (témoin)
%   lualatex -jobname=faux     "\def\CORR{1}\def\FAUX{1}\input{minimal-juge3}"           deux vrais défauts : doivent tomber
%   lualatex -jobname=trace    "\def\TRACE{1}\input{minimal-juge3}"                      trace de l'appel
% Les lignes du correctif sont tirées du préambule commité par extraire_correctifs.py, pas retapées.
\documentclass{scrbook}
\ifdefined\SANSCROCHET \RemoveFromHook{package/amsthm/after}[firstaid/aliascounter]\fi
\usepackage{amsthm}
\usepackage{hyperref}
\usepackage[french, nameinlink, noabbrev]{cleveref}
\ifdefined\AMSSEUL
  \newtheorem{theoreme}{Théorème}[chapter]
  \newtheorem{proposition}[theoreme]{Proposition}
  \newtheorem{corollaire}[theoreme]{Corollaire}
  \newtheorem{remarque}{Remarque}[chapter]
\else
  \usepackage{thmtools}
  \ifdefined\CORR \input{correctif-a-pose}\fi
  \declaretheorem[numberwithin=chapter, name=Théorème]{theoreme}
  \ifdefined\TRACE \tracingmacros=2 \tracingonline=0 \fi
  \declaretheorem[sibling=theoreme, name=Proposition]{proposition}
  \ifdefined\TRACE \tracingmacros=0 \fi
  \declaretheorem[sibling=theoreme, name=Corollaire]{corollaire}
  \declaretheorem[numberwithin=chapter, name=Remarque]{remarque}
  \ifdefined\FAUX \newcounteralias{corollaire}{section}\newcounteralias{orphelin}{inexistant}\fi
  \ifdefined\CORR \input{correctif-a-remise}\fi
\fi
\crefname{theoreme}{théorème}{théorèmes}
\crefname{proposition}{proposition}{propositions}
\crefname{corollaire}{corollaire}{corollaires}
\crefname{remarque}{remarque}{remarques}
\begin{document}
\chapter{Premier}
\begin{proposition}\label{p1}a\end{proposition}
\begin{theoreme}\label{t1}b\end{theoreme}
\begin{remarque}\label{r1}c\end{remarque}
\begin{corollaire}\label{c1}d\end{corollaire}
\chapter{Second}
\begin{corollaire}\label{c2}e\end{corollaire}
\begin{theoreme}\label{t2}f\end{theoreme}
RENVOIS \ref{p1} \ref{t1} \ref{r1} \ref{c1} \ref{c2} \ref{t2} ; \cref{p1} ; \cref{t1} ; \cref{r1} ; \cref{c1} ; \cref{c2} ; \cref{t2} FIN
\end{document}
"                                   sans correctif
%   lualatex -jobname=corr     "\def\CORR{1}% Juge de la troisième passe : document minimal écrit de zéro pour la cause A (thmtools sous LaTeX 2026-06-01).
% Variantes, par une macro posée sur la ligne de commande (deux passes chacune, depuis memoire/essai/juge3/minimal/) :
%   lualatex -jobname=nu       "\input{minimal-juge3}"                                   sans correctif
%   lualatex -jobname=corr     "\def\CORR{1}\input{minimal-juge3}"                       correctif A du préambule
%   lualatex -jobname=crochet  "\def\SANSCROCHET{1}\input{minimal-juge3}"                sans correctif, crochet retiré
%   lualatex -jobname=amsseul  "\def\AMSSEUL{1}\input{minimal-juge3}"                    amsthm seul (témoin)
%   lualatex -jobname=faux     "\def\CORR{1}\def\FAUX{1}\input{minimal-juge3}"           deux vrais défauts : doivent tomber
%   lualatex -jobname=trace    "\def\TRACE{1}\input{minimal-juge3}"                      trace de l'appel
% Les lignes du correctif sont tirées du préambule commité par extraire_correctifs.py, pas retapées.
\documentclass{scrbook}
\ifdefined\SANSCROCHET \RemoveFromHook{package/amsthm/after}[firstaid/aliascounter]\fi
\usepackage{amsthm}
\usepackage{hyperref}
\usepackage[french, nameinlink, noabbrev]{cleveref}
\ifdefined\AMSSEUL
  \newtheorem{theoreme}{Théorème}[chapter]
  \newtheorem{proposition}[theoreme]{Proposition}
  \newtheorem{corollaire}[theoreme]{Corollaire}
  \newtheorem{remarque}{Remarque}[chapter]
\else
  \usepackage{thmtools}
  \ifdefined\CORR \input{correctif-a-pose}\fi
  \declaretheorem[numberwithin=chapter, name=Théorème]{theoreme}
  \ifdefined\TRACE \tracingmacros=2 \tracingonline=0 \fi
  \declaretheorem[sibling=theoreme, name=Proposition]{proposition}
  \ifdefined\TRACE \tracingmacros=0 \fi
  \declaretheorem[sibling=theoreme, name=Corollaire]{corollaire}
  \declaretheorem[numberwithin=chapter, name=Remarque]{remarque}
  \ifdefined\FAUX \newcounteralias{corollaire}{section}\newcounteralias{orphelin}{inexistant}\fi
  \ifdefined\CORR \input{correctif-a-remise}\fi
\fi
\crefname{theoreme}{théorème}{théorèmes}
\crefname{proposition}{proposition}{propositions}
\crefname{corollaire}{corollaire}{corollaires}
\crefname{remarque}{remarque}{remarques}
\begin{document}
\chapter{Premier}
\begin{proposition}\label{p1}a\end{proposition}
\begin{theoreme}\label{t1}b\end{theoreme}
\begin{remarque}\label{r1}c\end{remarque}
\begin{corollaire}\label{c1}d\end{corollaire}
\chapter{Second}
\begin{corollaire}\label{c2}e\end{corollaire}
\begin{theoreme}\label{t2}f\end{theoreme}
RENVOIS \ref{p1} \ref{t1} \ref{r1} \ref{c1} \ref{c2} \ref{t2} ; \cref{p1} ; \cref{t1} ; \cref{r1} ; \cref{c1} ; \cref{c2} ; \cref{t2} FIN
\end{document}
"                       correctif A du préambule
%   lualatex -jobname=crochet  "\def\SANSCROCHET{1}% Juge de la troisième passe : document minimal écrit de zéro pour la cause A (thmtools sous LaTeX 2026-06-01).
% Variantes, par une macro posée sur la ligne de commande (deux passes chacune, depuis memoire/essai/juge3/minimal/) :
%   lualatex -jobname=nu       "\input{minimal-juge3}"                                   sans correctif
%   lualatex -jobname=corr     "\def\CORR{1}\input{minimal-juge3}"                       correctif A du préambule
%   lualatex -jobname=crochet  "\def\SANSCROCHET{1}\input{minimal-juge3}"                sans correctif, crochet retiré
%   lualatex -jobname=amsseul  "\def\AMSSEUL{1}\input{minimal-juge3}"                    amsthm seul (témoin)
%   lualatex -jobname=faux     "\def\CORR{1}\def\FAUX{1}\input{minimal-juge3}"           deux vrais défauts : doivent tomber
%   lualatex -jobname=trace    "\def\TRACE{1}\input{minimal-juge3}"                      trace de l'appel
% Les lignes du correctif sont tirées du préambule commité par extraire_correctifs.py, pas retapées.
\documentclass{scrbook}
\ifdefined\SANSCROCHET \RemoveFromHook{package/amsthm/after}[firstaid/aliascounter]\fi
\usepackage{amsthm}
\usepackage{hyperref}
\usepackage[french, nameinlink, noabbrev]{cleveref}
\ifdefined\AMSSEUL
  \newtheorem{theoreme}{Théorème}[chapter]
  \newtheorem{proposition}[theoreme]{Proposition}
  \newtheorem{corollaire}[theoreme]{Corollaire}
  \newtheorem{remarque}{Remarque}[chapter]
\else
  \usepackage{thmtools}
  \ifdefined\CORR \input{correctif-a-pose}\fi
  \declaretheorem[numberwithin=chapter, name=Théorème]{theoreme}
  \ifdefined\TRACE \tracingmacros=2 \tracingonline=0 \fi
  \declaretheorem[sibling=theoreme, name=Proposition]{proposition}
  \ifdefined\TRACE \tracingmacros=0 \fi
  \declaretheorem[sibling=theoreme, name=Corollaire]{corollaire}
  \declaretheorem[numberwithin=chapter, name=Remarque]{remarque}
  \ifdefined\FAUX \newcounteralias{corollaire}{section}\newcounteralias{orphelin}{inexistant}\fi
  \ifdefined\CORR \input{correctif-a-remise}\fi
\fi
\crefname{theoreme}{théorème}{théorèmes}
\crefname{proposition}{proposition}{propositions}
\crefname{corollaire}{corollaire}{corollaires}
\crefname{remarque}{remarque}{remarques}
\begin{document}
\chapter{Premier}
\begin{proposition}\label{p1}a\end{proposition}
\begin{theoreme}\label{t1}b\end{theoreme}
\begin{remarque}\label{r1}c\end{remarque}
\begin{corollaire}\label{c1}d\end{corollaire}
\chapter{Second}
\begin{corollaire}\label{c2}e\end{corollaire}
\begin{theoreme}\label{t2}f\end{theoreme}
RENVOIS \ref{p1} \ref{t1} \ref{r1} \ref{c1} \ref{c2} \ref{t2} ; \cref{p1} ; \cref{t1} ; \cref{r1} ; \cref{c1} ; \cref{c2} ; \cref{t2} FIN
\end{document}
"                sans correctif, crochet retiré
%   lualatex -jobname=amsseul  "\def\AMSSEUL{1}% Juge de la troisième passe : document minimal écrit de zéro pour la cause A (thmtools sous LaTeX 2026-06-01).
% Variantes, par une macro posée sur la ligne de commande (deux passes chacune, depuis memoire/essai/juge3/minimal/) :
%   lualatex -jobname=nu       "\input{minimal-juge3}"                                   sans correctif
%   lualatex -jobname=corr     "\def\CORR{1}\input{minimal-juge3}"                       correctif A du préambule
%   lualatex -jobname=crochet  "\def\SANSCROCHET{1}\input{minimal-juge3}"                sans correctif, crochet retiré
%   lualatex -jobname=amsseul  "\def\AMSSEUL{1}\input{minimal-juge3}"                    amsthm seul (témoin)
%   lualatex -jobname=faux     "\def\CORR{1}\def\FAUX{1}\input{minimal-juge3}"           deux vrais défauts : doivent tomber
%   lualatex -jobname=trace    "\def\TRACE{1}\input{minimal-juge3}"                      trace de l'appel
% Les lignes du correctif sont tirées du préambule commité par extraire_correctifs.py, pas retapées.
\documentclass{scrbook}
\ifdefined\SANSCROCHET \RemoveFromHook{package/amsthm/after}[firstaid/aliascounter]\fi
\usepackage{amsthm}
\usepackage{hyperref}
\usepackage[french, nameinlink, noabbrev]{cleveref}
\ifdefined\AMSSEUL
  \newtheorem{theoreme}{Théorème}[chapter]
  \newtheorem{proposition}[theoreme]{Proposition}
  \newtheorem{corollaire}[theoreme]{Corollaire}
  \newtheorem{remarque}{Remarque}[chapter]
\else
  \usepackage{thmtools}
  \ifdefined\CORR \input{correctif-a-pose}\fi
  \declaretheorem[numberwithin=chapter, name=Théorème]{theoreme}
  \ifdefined\TRACE \tracingmacros=2 \tracingonline=0 \fi
  \declaretheorem[sibling=theoreme, name=Proposition]{proposition}
  \ifdefined\TRACE \tracingmacros=0 \fi
  \declaretheorem[sibling=theoreme, name=Corollaire]{corollaire}
  \declaretheorem[numberwithin=chapter, name=Remarque]{remarque}
  \ifdefined\FAUX \newcounteralias{corollaire}{section}\newcounteralias{orphelin}{inexistant}\fi
  \ifdefined\CORR \input{correctif-a-remise}\fi
\fi
\crefname{theoreme}{théorème}{théorèmes}
\crefname{proposition}{proposition}{propositions}
\crefname{corollaire}{corollaire}{corollaires}
\crefname{remarque}{remarque}{remarques}
\begin{document}
\chapter{Premier}
\begin{proposition}\label{p1}a\end{proposition}
\begin{theoreme}\label{t1}b\end{theoreme}
\begin{remarque}\label{r1}c\end{remarque}
\begin{corollaire}\label{c1}d\end{corollaire}
\chapter{Second}
\begin{corollaire}\label{c2}e\end{corollaire}
\begin{theoreme}\label{t2}f\end{theoreme}
RENVOIS \ref{p1} \ref{t1} \ref{r1} \ref{c1} \ref{c2} \ref{t2} ; \cref{p1} ; \cref{t1} ; \cref{r1} ; \cref{c1} ; \cref{c2} ; \cref{t2} FIN
\end{document}
"                    amsthm seul (témoin)
%   lualatex -jobname=faux     "\def\CORR{1}\def\FAUX{1}% Juge de la troisième passe : document minimal écrit de zéro pour la cause A (thmtools sous LaTeX 2026-06-01).
% Variantes, par une macro posée sur la ligne de commande (deux passes chacune, depuis memoire/essai/juge3/minimal/) :
%   lualatex -jobname=nu       "\input{minimal-juge3}"                                   sans correctif
%   lualatex -jobname=corr     "\def\CORR{1}\input{minimal-juge3}"                       correctif A du préambule
%   lualatex -jobname=crochet  "\def\SANSCROCHET{1}\input{minimal-juge3}"                sans correctif, crochet retiré
%   lualatex -jobname=amsseul  "\def\AMSSEUL{1}\input{minimal-juge3}"                    amsthm seul (témoin)
%   lualatex -jobname=faux     "\def\CORR{1}\def\FAUX{1}\input{minimal-juge3}"           deux vrais défauts : doivent tomber
%   lualatex -jobname=trace    "\def\TRACE{1}\input{minimal-juge3}"                      trace de l'appel
% Les lignes du correctif sont tirées du préambule commité par extraire_correctifs.py, pas retapées.
\documentclass{scrbook}
\ifdefined\SANSCROCHET \RemoveFromHook{package/amsthm/after}[firstaid/aliascounter]\fi
\usepackage{amsthm}
\usepackage{hyperref}
\usepackage[french, nameinlink, noabbrev]{cleveref}
\ifdefined\AMSSEUL
  \newtheorem{theoreme}{Théorème}[chapter]
  \newtheorem{proposition}[theoreme]{Proposition}
  \newtheorem{corollaire}[theoreme]{Corollaire}
  \newtheorem{remarque}{Remarque}[chapter]
\else
  \usepackage{thmtools}
  \ifdefined\CORR \input{correctif-a-pose}\fi
  \declaretheorem[numberwithin=chapter, name=Théorème]{theoreme}
  \ifdefined\TRACE \tracingmacros=2 \tracingonline=0 \fi
  \declaretheorem[sibling=theoreme, name=Proposition]{proposition}
  \ifdefined\TRACE \tracingmacros=0 \fi
  \declaretheorem[sibling=theoreme, name=Corollaire]{corollaire}
  \declaretheorem[numberwithin=chapter, name=Remarque]{remarque}
  \ifdefined\FAUX \newcounteralias{corollaire}{section}\newcounteralias{orphelin}{inexistant}\fi
  \ifdefined\CORR \input{correctif-a-remise}\fi
\fi
\crefname{theoreme}{théorème}{théorèmes}
\crefname{proposition}{proposition}{propositions}
\crefname{corollaire}{corollaire}{corollaires}
\crefname{remarque}{remarque}{remarques}
\begin{document}
\chapter{Premier}
\begin{proposition}\label{p1}a\end{proposition}
\begin{theoreme}\label{t1}b\end{theoreme}
\begin{remarque}\label{r1}c\end{remarque}
\begin{corollaire}\label{c1}d\end{corollaire}
\chapter{Second}
\begin{corollaire}\label{c2}e\end{corollaire}
\begin{theoreme}\label{t2}f\end{theoreme}
RENVOIS \ref{p1} \ref{t1} \ref{r1} \ref{c1} \ref{c2} \ref{t2} ; \cref{p1} ; \cref{t1} ; \cref{r1} ; \cref{c1} ; \cref{c2} ; \cref{t2} FIN
\end{document}
"           deux vrais défauts : doivent tomber
%   lualatex -jobname=trace    "\def\TRACE{1}% Juge de la troisième passe : document minimal écrit de zéro pour la cause A (thmtools sous LaTeX 2026-06-01).
% Variantes, par une macro posée sur la ligne de commande (deux passes chacune, depuis memoire/essai/juge3/minimal/) :
%   lualatex -jobname=nu       "\input{minimal-juge3}"                                   sans correctif
%   lualatex -jobname=corr     "\def\CORR{1}\input{minimal-juge3}"                       correctif A du préambule
%   lualatex -jobname=crochet  "\def\SANSCROCHET{1}\input{minimal-juge3}"                sans correctif, crochet retiré
%   lualatex -jobname=amsseul  "\def\AMSSEUL{1}\input{minimal-juge3}"                    amsthm seul (témoin)
%   lualatex -jobname=faux     "\def\CORR{1}\def\FAUX{1}\input{minimal-juge3}"           deux vrais défauts : doivent tomber
%   lualatex -jobname=trace    "\def\TRACE{1}\input{minimal-juge3}"                      trace de l'appel
% Les lignes du correctif sont tirées du préambule commité par extraire_correctifs.py, pas retapées.
\documentclass{scrbook}
\ifdefined\SANSCROCHET \RemoveFromHook{package/amsthm/after}[firstaid/aliascounter]\fi
\usepackage{amsthm}
\usepackage{hyperref}
\usepackage[french, nameinlink, noabbrev]{cleveref}
\ifdefined\AMSSEUL
  \newtheorem{theoreme}{Théorème}[chapter]
  \newtheorem{proposition}[theoreme]{Proposition}
  \newtheorem{corollaire}[theoreme]{Corollaire}
  \newtheorem{remarque}{Remarque}[chapter]
\else
  \usepackage{thmtools}
  \ifdefined\CORR \input{correctif-a-pose}\fi
  \declaretheorem[numberwithin=chapter, name=Théorème]{theoreme}
  \ifdefined\TRACE \tracingmacros=2 \tracingonline=0 \fi
  \declaretheorem[sibling=theoreme, name=Proposition]{proposition}
  \ifdefined\TRACE \tracingmacros=0 \fi
  \declaretheorem[sibling=theoreme, name=Corollaire]{corollaire}
  \declaretheorem[numberwithin=chapter, name=Remarque]{remarque}
  \ifdefined\FAUX \newcounteralias{corollaire}{section}\newcounteralias{orphelin}{inexistant}\fi
  \ifdefined\CORR \input{correctif-a-remise}\fi
\fi
\crefname{theoreme}{théorème}{théorèmes}
\crefname{proposition}{proposition}{propositions}
\crefname{corollaire}{corollaire}{corollaires}
\crefname{remarque}{remarque}{remarques}
\begin{document}
\chapter{Premier}
\begin{proposition}\label{p1}a\end{proposition}
\begin{theoreme}\label{t1}b\end{theoreme}
\begin{remarque}\label{r1}c\end{remarque}
\begin{corollaire}\label{c1}d\end{corollaire}
\chapter{Second}
\begin{corollaire}\label{c2}e\end{corollaire}
\begin{theoreme}\label{t2}f\end{theoreme}
RENVOIS \ref{p1} \ref{t1} \ref{r1} \ref{c1} \ref{c2} \ref{t2} ; \cref{p1} ; \cref{t1} ; \cref{r1} ; \cref{c1} ; \cref{c2} ; \cref{t2} FIN
\end{document}
"                      trace de l'appel
% Les lignes du correctif sont tirées du préambule commité par extraire_correctifs.py, pas retapées.
\documentclass{scrbook}
\ifdefined\SANSCROCHET \RemoveFromHook{package/amsthm/after}[firstaid/aliascounter]\fi
\usepackage{amsthm}
\usepackage{hyperref}
\usepackage[french, nameinlink, noabbrev]{cleveref}
\ifdefined\AMSSEUL
  \newtheorem{theoreme}{Théorème}[chapter]
  \newtheorem{proposition}[theoreme]{Proposition}
  \newtheorem{corollaire}[theoreme]{Corollaire}
  \newtheorem{remarque}{Remarque}[chapter]
\else
  \usepackage{thmtools}
  \ifdefined\CORR \input{correctif-a-pose}\fi
  \declaretheorem[numberwithin=chapter, name=Théorème]{theoreme}
  \ifdefined\TRACE \tracingmacros=2 \tracingonline=0 \fi
  \declaretheorem[sibling=theoreme, name=Proposition]{proposition}
  \ifdefined\TRACE \tracingmacros=0 \fi
  \declaretheorem[sibling=theoreme, name=Corollaire]{corollaire}
  \declaretheorem[numberwithin=chapter, name=Remarque]{remarque}
  \ifdefined\FAUX \newcounteralias{corollaire}{section}\newcounteralias{orphelin}{inexistant}\fi
  \ifdefined\CORR \input{correctif-a-remise}\fi
\fi
\crefname{theoreme}{théorème}{théorèmes}
\crefname{proposition}{proposition}{propositions}
\crefname{corollaire}{corollaire}{corollaires}
\crefname{remarque}{remarque}{remarques}
\begin{document}
\chapter{Premier}
\begin{proposition}\label{p1}a\end{proposition}
\begin{theoreme}\label{t1}b\end{theoreme}
\begin{remarque}\label{r1}c\end{remarque}
\begin{corollaire}\label{c1}d\end{corollaire}
\chapter{Second}
\begin{corollaire}\label{c2}e\end{corollaire}
\begin{theoreme}\label{t2}f\end{theoreme}
RENVOIS \ref{p1} \ref{t1} \ref{r1} \ref{c1} \ref{c2} \ref{t2} ; \cref{p1} ; \cref{t1} ; \cref{r1} ; \cref{c1} ; \cref{c2} ; \cref{t2} FIN
\end{document}
"                       correctif A du préambule
%   lualatex -jobname=crochet  "\def\SANSCROCHET{1}% Juge de la troisième passe : document minimal écrit de zéro pour la cause A (thmtools sous LaTeX 2026-06-01).
% Variantes, par une macro posée sur la ligne de commande (deux passes chacune, depuis memoire/essai/juge3/minimal/) :
%   lualatex -jobname=nu       "% Juge de la troisième passe : document minimal écrit de zéro pour la cause A (thmtools sous LaTeX 2026-06-01).
% Variantes, par une macro posée sur la ligne de commande (deux passes chacune, depuis memoire/essai/juge3/minimal/) :
%   lualatex -jobname=nu       "\input{minimal-juge3}"                                   sans correctif
%   lualatex -jobname=corr     "\def\CORR{1}\input{minimal-juge3}"                       correctif A du préambule
%   lualatex -jobname=crochet  "\def\SANSCROCHET{1}\input{minimal-juge3}"                sans correctif, crochet retiré
%   lualatex -jobname=amsseul  "\def\AMSSEUL{1}\input{minimal-juge3}"                    amsthm seul (témoin)
%   lualatex -jobname=faux     "\def\CORR{1}\def\FAUX{1}\input{minimal-juge3}"           deux vrais défauts : doivent tomber
%   lualatex -jobname=trace    "\def\TRACE{1}\input{minimal-juge3}"                      trace de l'appel
% Les lignes du correctif sont tirées du préambule commité par extraire_correctifs.py, pas retapées.
\documentclass{scrbook}
\ifdefined\SANSCROCHET \RemoveFromHook{package/amsthm/after}[firstaid/aliascounter]\fi
\usepackage{amsthm}
\usepackage{hyperref}
\usepackage[french, nameinlink, noabbrev]{cleveref}
\ifdefined\AMSSEUL
  \newtheorem{theoreme}{Théorème}[chapter]
  \newtheorem{proposition}[theoreme]{Proposition}
  \newtheorem{corollaire}[theoreme]{Corollaire}
  \newtheorem{remarque}{Remarque}[chapter]
\else
  \usepackage{thmtools}
  \ifdefined\CORR \input{correctif-a-pose}\fi
  \declaretheorem[numberwithin=chapter, name=Théorème]{theoreme}
  \ifdefined\TRACE \tracingmacros=2 \tracingonline=0 \fi
  \declaretheorem[sibling=theoreme, name=Proposition]{proposition}
  \ifdefined\TRACE \tracingmacros=0 \fi
  \declaretheorem[sibling=theoreme, name=Corollaire]{corollaire}
  \declaretheorem[numberwithin=chapter, name=Remarque]{remarque}
  \ifdefined\FAUX \newcounteralias{corollaire}{section}\newcounteralias{orphelin}{inexistant}\fi
  \ifdefined\CORR \input{correctif-a-remise}\fi
\fi
\crefname{theoreme}{théorème}{théorèmes}
\crefname{proposition}{proposition}{propositions}
\crefname{corollaire}{corollaire}{corollaires}
\crefname{remarque}{remarque}{remarques}
\begin{document}
\chapter{Premier}
\begin{proposition}\label{p1}a\end{proposition}
\begin{theoreme}\label{t1}b\end{theoreme}
\begin{remarque}\label{r1}c\end{remarque}
\begin{corollaire}\label{c1}d\end{corollaire}
\chapter{Second}
\begin{corollaire}\label{c2}e\end{corollaire}
\begin{theoreme}\label{t2}f\end{theoreme}
RENVOIS \ref{p1} \ref{t1} \ref{r1} \ref{c1} \ref{c2} \ref{t2} ; \cref{p1} ; \cref{t1} ; \cref{r1} ; \cref{c1} ; \cref{c2} ; \cref{t2} FIN
\end{document}
"                                   sans correctif
%   lualatex -jobname=corr     "\def\CORR{1}% Juge de la troisième passe : document minimal écrit de zéro pour la cause A (thmtools sous LaTeX 2026-06-01).
% Variantes, par une macro posée sur la ligne de commande (deux passes chacune, depuis memoire/essai/juge3/minimal/) :
%   lualatex -jobname=nu       "\input{minimal-juge3}"                                   sans correctif
%   lualatex -jobname=corr     "\def\CORR{1}\input{minimal-juge3}"                       correctif A du préambule
%   lualatex -jobname=crochet  "\def\SANSCROCHET{1}\input{minimal-juge3}"                sans correctif, crochet retiré
%   lualatex -jobname=amsseul  "\def\AMSSEUL{1}\input{minimal-juge3}"                    amsthm seul (témoin)
%   lualatex -jobname=faux     "\def\CORR{1}\def\FAUX{1}\input{minimal-juge3}"           deux vrais défauts : doivent tomber
%   lualatex -jobname=trace    "\def\TRACE{1}\input{minimal-juge3}"                      trace de l'appel
% Les lignes du correctif sont tirées du préambule commité par extraire_correctifs.py, pas retapées.
\documentclass{scrbook}
\ifdefined\SANSCROCHET \RemoveFromHook{package/amsthm/after}[firstaid/aliascounter]\fi
\usepackage{amsthm}
\usepackage{hyperref}
\usepackage[french, nameinlink, noabbrev]{cleveref}
\ifdefined\AMSSEUL
  \newtheorem{theoreme}{Théorème}[chapter]
  \newtheorem{proposition}[theoreme]{Proposition}
  \newtheorem{corollaire}[theoreme]{Corollaire}
  \newtheorem{remarque}{Remarque}[chapter]
\else
  \usepackage{thmtools}
  \ifdefined\CORR \input{correctif-a-pose}\fi
  \declaretheorem[numberwithin=chapter, name=Théorème]{theoreme}
  \ifdefined\TRACE \tracingmacros=2 \tracingonline=0 \fi
  \declaretheorem[sibling=theoreme, name=Proposition]{proposition}
  \ifdefined\TRACE \tracingmacros=0 \fi
  \declaretheorem[sibling=theoreme, name=Corollaire]{corollaire}
  \declaretheorem[numberwithin=chapter, name=Remarque]{remarque}
  \ifdefined\FAUX \newcounteralias{corollaire}{section}\newcounteralias{orphelin}{inexistant}\fi
  \ifdefined\CORR \input{correctif-a-remise}\fi
\fi
\crefname{theoreme}{théorème}{théorèmes}
\crefname{proposition}{proposition}{propositions}
\crefname{corollaire}{corollaire}{corollaires}
\crefname{remarque}{remarque}{remarques}
\begin{document}
\chapter{Premier}
\begin{proposition}\label{p1}a\end{proposition}
\begin{theoreme}\label{t1}b\end{theoreme}
\begin{remarque}\label{r1}c\end{remarque}
\begin{corollaire}\label{c1}d\end{corollaire}
\chapter{Second}
\begin{corollaire}\label{c2}e\end{corollaire}
\begin{theoreme}\label{t2}f\end{theoreme}
RENVOIS \ref{p1} \ref{t1} \ref{r1} \ref{c1} \ref{c2} \ref{t2} ; \cref{p1} ; \cref{t1} ; \cref{r1} ; \cref{c1} ; \cref{c2} ; \cref{t2} FIN
\end{document}
"                       correctif A du préambule
%   lualatex -jobname=crochet  "\def\SANSCROCHET{1}% Juge de la troisième passe : document minimal écrit de zéro pour la cause A (thmtools sous LaTeX 2026-06-01).
% Variantes, par une macro posée sur la ligne de commande (deux passes chacune, depuis memoire/essai/juge3/minimal/) :
%   lualatex -jobname=nu       "\input{minimal-juge3}"                                   sans correctif
%   lualatex -jobname=corr     "\def\CORR{1}\input{minimal-juge3}"                       correctif A du préambule
%   lualatex -jobname=crochet  "\def\SANSCROCHET{1}\input{minimal-juge3}"                sans correctif, crochet retiré
%   lualatex -jobname=amsseul  "\def\AMSSEUL{1}\input{minimal-juge3}"                    amsthm seul (témoin)
%   lualatex -jobname=faux     "\def\CORR{1}\def\FAUX{1}\input{minimal-juge3}"           deux vrais défauts : doivent tomber
%   lualatex -jobname=trace    "\def\TRACE{1}\input{minimal-juge3}"                      trace de l'appel
% Les lignes du correctif sont tirées du préambule commité par extraire_correctifs.py, pas retapées.
\documentclass{scrbook}
\ifdefined\SANSCROCHET \RemoveFromHook{package/amsthm/after}[firstaid/aliascounter]\fi
\usepackage{amsthm}
\usepackage{hyperref}
\usepackage[french, nameinlink, noabbrev]{cleveref}
\ifdefined\AMSSEUL
  \newtheorem{theoreme}{Théorème}[chapter]
  \newtheorem{proposition}[theoreme]{Proposition}
  \newtheorem{corollaire}[theoreme]{Corollaire}
  \newtheorem{remarque}{Remarque}[chapter]
\else
  \usepackage{thmtools}
  \ifdefined\CORR \input{correctif-a-pose}\fi
  \declaretheorem[numberwithin=chapter, name=Théorème]{theoreme}
  \ifdefined\TRACE \tracingmacros=2 \tracingonline=0 \fi
  \declaretheorem[sibling=theoreme, name=Proposition]{proposition}
  \ifdefined\TRACE \tracingmacros=0 \fi
  \declaretheorem[sibling=theoreme, name=Corollaire]{corollaire}
  \declaretheorem[numberwithin=chapter, name=Remarque]{remarque}
  \ifdefined\FAUX \newcounteralias{corollaire}{section}\newcounteralias{orphelin}{inexistant}\fi
  \ifdefined\CORR \input{correctif-a-remise}\fi
\fi
\crefname{theoreme}{théorème}{théorèmes}
\crefname{proposition}{proposition}{propositions}
\crefname{corollaire}{corollaire}{corollaires}
\crefname{remarque}{remarque}{remarques}
\begin{document}
\chapter{Premier}
\begin{proposition}\label{p1}a\end{proposition}
\begin{theoreme}\label{t1}b\end{theoreme}
\begin{remarque}\label{r1}c\end{remarque}
\begin{corollaire}\label{c1}d\end{corollaire}
\chapter{Second}
\begin{corollaire}\label{c2}e\end{corollaire}
\begin{theoreme}\label{t2}f\end{theoreme}
RENVOIS \ref{p1} \ref{t1} \ref{r1} \ref{c1} \ref{c2} \ref{t2} ; \cref{p1} ; \cref{t1} ; \cref{r1} ; \cref{c1} ; \cref{c2} ; \cref{t2} FIN
\end{document}
"                sans correctif, crochet retiré
%   lualatex -jobname=amsseul  "\def\AMSSEUL{1}% Juge de la troisième passe : document minimal écrit de zéro pour la cause A (thmtools sous LaTeX 2026-06-01).
% Variantes, par une macro posée sur la ligne de commande (deux passes chacune, depuis memoire/essai/juge3/minimal/) :
%   lualatex -jobname=nu       "\input{minimal-juge3}"                                   sans correctif
%   lualatex -jobname=corr     "\def\CORR{1}\input{minimal-juge3}"                       correctif A du préambule
%   lualatex -jobname=crochet  "\def\SANSCROCHET{1}\input{minimal-juge3}"                sans correctif, crochet retiré
%   lualatex -jobname=amsseul  "\def\AMSSEUL{1}\input{minimal-juge3}"                    amsthm seul (témoin)
%   lualatex -jobname=faux     "\def\CORR{1}\def\FAUX{1}\input{minimal-juge3}"           deux vrais défauts : doivent tomber
%   lualatex -jobname=trace    "\def\TRACE{1}\input{minimal-juge3}"                      trace de l'appel
% Les lignes du correctif sont tirées du préambule commité par extraire_correctifs.py, pas retapées.
\documentclass{scrbook}
\ifdefined\SANSCROCHET \RemoveFromHook{package/amsthm/after}[firstaid/aliascounter]\fi
\usepackage{amsthm}
\usepackage{hyperref}
\usepackage[french, nameinlink, noabbrev]{cleveref}
\ifdefined\AMSSEUL
  \newtheorem{theoreme}{Théorème}[chapter]
  \newtheorem{proposition}[theoreme]{Proposition}
  \newtheorem{corollaire}[theoreme]{Corollaire}
  \newtheorem{remarque}{Remarque}[chapter]
\else
  \usepackage{thmtools}
  \ifdefined\CORR \input{correctif-a-pose}\fi
  \declaretheorem[numberwithin=chapter, name=Théorème]{theoreme}
  \ifdefined\TRACE \tracingmacros=2 \tracingonline=0 \fi
  \declaretheorem[sibling=theoreme, name=Proposition]{proposition}
  \ifdefined\TRACE \tracingmacros=0 \fi
  \declaretheorem[sibling=theoreme, name=Corollaire]{corollaire}
  \declaretheorem[numberwithin=chapter, name=Remarque]{remarque}
  \ifdefined\FAUX \newcounteralias{corollaire}{section}\newcounteralias{orphelin}{inexistant}\fi
  \ifdefined\CORR \input{correctif-a-remise}\fi
\fi
\crefname{theoreme}{théorème}{théorèmes}
\crefname{proposition}{proposition}{propositions}
\crefname{corollaire}{corollaire}{corollaires}
\crefname{remarque}{remarque}{remarques}
\begin{document}
\chapter{Premier}
\begin{proposition}\label{p1}a\end{proposition}
\begin{theoreme}\label{t1}b\end{theoreme}
\begin{remarque}\label{r1}c\end{remarque}
\begin{corollaire}\label{c1}d\end{corollaire}
\chapter{Second}
\begin{corollaire}\label{c2}e\end{corollaire}
\begin{theoreme}\label{t2}f\end{theoreme}
RENVOIS \ref{p1} \ref{t1} \ref{r1} \ref{c1} \ref{c2} \ref{t2} ; \cref{p1} ; \cref{t1} ; \cref{r1} ; \cref{c1} ; \cref{c2} ; \cref{t2} FIN
\end{document}
"                    amsthm seul (témoin)
%   lualatex -jobname=faux     "\def\CORR{1}\def\FAUX{1}% Juge de la troisième passe : document minimal écrit de zéro pour la cause A (thmtools sous LaTeX 2026-06-01).
% Variantes, par une macro posée sur la ligne de commande (deux passes chacune, depuis memoire/essai/juge3/minimal/) :
%   lualatex -jobname=nu       "\input{minimal-juge3}"                                   sans correctif
%   lualatex -jobname=corr     "\def\CORR{1}\input{minimal-juge3}"                       correctif A du préambule
%   lualatex -jobname=crochet  "\def\SANSCROCHET{1}\input{minimal-juge3}"                sans correctif, crochet retiré
%   lualatex -jobname=amsseul  "\def\AMSSEUL{1}\input{minimal-juge3}"                    amsthm seul (témoin)
%   lualatex -jobname=faux     "\def\CORR{1}\def\FAUX{1}\input{minimal-juge3}"           deux vrais défauts : doivent tomber
%   lualatex -jobname=trace    "\def\TRACE{1}\input{minimal-juge3}"                      trace de l'appel
% Les lignes du correctif sont tirées du préambule commité par extraire_correctifs.py, pas retapées.
\documentclass{scrbook}
\ifdefined\SANSCROCHET \RemoveFromHook{package/amsthm/after}[firstaid/aliascounter]\fi
\usepackage{amsthm}
\usepackage{hyperref}
\usepackage[french, nameinlink, noabbrev]{cleveref}
\ifdefined\AMSSEUL
  \newtheorem{theoreme}{Théorème}[chapter]
  \newtheorem{proposition}[theoreme]{Proposition}
  \newtheorem{corollaire}[theoreme]{Corollaire}
  \newtheorem{remarque}{Remarque}[chapter]
\else
  \usepackage{thmtools}
  \ifdefined\CORR \input{correctif-a-pose}\fi
  \declaretheorem[numberwithin=chapter, name=Théorème]{theoreme}
  \ifdefined\TRACE \tracingmacros=2 \tracingonline=0 \fi
  \declaretheorem[sibling=theoreme, name=Proposition]{proposition}
  \ifdefined\TRACE \tracingmacros=0 \fi
  \declaretheorem[sibling=theoreme, name=Corollaire]{corollaire}
  \declaretheorem[numberwithin=chapter, name=Remarque]{remarque}
  \ifdefined\FAUX \newcounteralias{corollaire}{section}\newcounteralias{orphelin}{inexistant}\fi
  \ifdefined\CORR \input{correctif-a-remise}\fi
\fi
\crefname{theoreme}{théorème}{théorèmes}
\crefname{proposition}{proposition}{propositions}
\crefname{corollaire}{corollaire}{corollaires}
\crefname{remarque}{remarque}{remarques}
\begin{document}
\chapter{Premier}
\begin{proposition}\label{p1}a\end{proposition}
\begin{theoreme}\label{t1}b\end{theoreme}
\begin{remarque}\label{r1}c\end{remarque}
\begin{corollaire}\label{c1}d\end{corollaire}
\chapter{Second}
\begin{corollaire}\label{c2}e\end{corollaire}
\begin{theoreme}\label{t2}f\end{theoreme}
RENVOIS \ref{p1} \ref{t1} \ref{r1} \ref{c1} \ref{c2} \ref{t2} ; \cref{p1} ; \cref{t1} ; \cref{r1} ; \cref{c1} ; \cref{c2} ; \cref{t2} FIN
\end{document}
"           deux vrais défauts : doivent tomber
%   lualatex -jobname=trace    "\def\TRACE{1}% Juge de la troisième passe : document minimal écrit de zéro pour la cause A (thmtools sous LaTeX 2026-06-01).
% Variantes, par une macro posée sur la ligne de commande (deux passes chacune, depuis memoire/essai/juge3/minimal/) :
%   lualatex -jobname=nu       "\input{minimal-juge3}"                                   sans correctif
%   lualatex -jobname=corr     "\def\CORR{1}\input{minimal-juge3}"                       correctif A du préambule
%   lualatex -jobname=crochet  "\def\SANSCROCHET{1}\input{minimal-juge3}"                sans correctif, crochet retiré
%   lualatex -jobname=amsseul  "\def\AMSSEUL{1}\input{minimal-juge3}"                    amsthm seul (témoin)
%   lualatex -jobname=faux     "\def\CORR{1}\def\FAUX{1}\input{minimal-juge3}"           deux vrais défauts : doivent tomber
%   lualatex -jobname=trace    "\def\TRACE{1}\input{minimal-juge3}"                      trace de l'appel
% Les lignes du correctif sont tirées du préambule commité par extraire_correctifs.py, pas retapées.
\documentclass{scrbook}
\ifdefined\SANSCROCHET \RemoveFromHook{package/amsthm/after}[firstaid/aliascounter]\fi
\usepackage{amsthm}
\usepackage{hyperref}
\usepackage[french, nameinlink, noabbrev]{cleveref}
\ifdefined\AMSSEUL
  \newtheorem{theoreme}{Théorème}[chapter]
  \newtheorem{proposition}[theoreme]{Proposition}
  \newtheorem{corollaire}[theoreme]{Corollaire}
  \newtheorem{remarque}{Remarque}[chapter]
\else
  \usepackage{thmtools}
  \ifdefined\CORR \input{correctif-a-pose}\fi
  \declaretheorem[numberwithin=chapter, name=Théorème]{theoreme}
  \ifdefined\TRACE \tracingmacros=2 \tracingonline=0 \fi
  \declaretheorem[sibling=theoreme, name=Proposition]{proposition}
  \ifdefined\TRACE \tracingmacros=0 \fi
  \declaretheorem[sibling=theoreme, name=Corollaire]{corollaire}
  \declaretheorem[numberwithin=chapter, name=Remarque]{remarque}
  \ifdefined\FAUX \newcounteralias{corollaire}{section}\newcounteralias{orphelin}{inexistant}\fi
  \ifdefined\CORR \input{correctif-a-remise}\fi
\fi
\crefname{theoreme}{théorème}{théorèmes}
\crefname{proposition}{proposition}{propositions}
\crefname{corollaire}{corollaire}{corollaires}
\crefname{remarque}{remarque}{remarques}
\begin{document}
\chapter{Premier}
\begin{proposition}\label{p1}a\end{proposition}
\begin{theoreme}\label{t1}b\end{theoreme}
\begin{remarque}\label{r1}c\end{remarque}
\begin{corollaire}\label{c1}d\end{corollaire}
\chapter{Second}
\begin{corollaire}\label{c2}e\end{corollaire}
\begin{theoreme}\label{t2}f\end{theoreme}
RENVOIS \ref{p1} \ref{t1} \ref{r1} \ref{c1} \ref{c2} \ref{t2} ; \cref{p1} ; \cref{t1} ; \cref{r1} ; \cref{c1} ; \cref{c2} ; \cref{t2} FIN
\end{document}
"                      trace de l'appel
% Les lignes du correctif sont tirées du préambule commité par extraire_correctifs.py, pas retapées.
\documentclass{scrbook}
\ifdefined\SANSCROCHET \RemoveFromHook{package/amsthm/after}[firstaid/aliascounter]\fi
\usepackage{amsthm}
\usepackage{hyperref}
\usepackage[french, nameinlink, noabbrev]{cleveref}
\ifdefined\AMSSEUL
  \newtheorem{theoreme}{Théorème}[chapter]
  \newtheorem{proposition}[theoreme]{Proposition}
  \newtheorem{corollaire}[theoreme]{Corollaire}
  \newtheorem{remarque}{Remarque}[chapter]
\else
  \usepackage{thmtools}
  \ifdefined\CORR \input{correctif-a-pose}\fi
  \declaretheorem[numberwithin=chapter, name=Théorème]{theoreme}
  \ifdefined\TRACE \tracingmacros=2 \tracingonline=0 \fi
  \declaretheorem[sibling=theoreme, name=Proposition]{proposition}
  \ifdefined\TRACE \tracingmacros=0 \fi
  \declaretheorem[sibling=theoreme, name=Corollaire]{corollaire}
  \declaretheorem[numberwithin=chapter, name=Remarque]{remarque}
  \ifdefined\FAUX \newcounteralias{corollaire}{section}\newcounteralias{orphelin}{inexistant}\fi
  \ifdefined\CORR \input{correctif-a-remise}\fi
\fi
\crefname{theoreme}{théorème}{théorèmes}
\crefname{proposition}{proposition}{propositions}
\crefname{corollaire}{corollaire}{corollaires}
\crefname{remarque}{remarque}{remarques}
\begin{document}
\chapter{Premier}
\begin{proposition}\label{p1}a\end{proposition}
\begin{theoreme}\label{t1}b\end{theoreme}
\begin{remarque}\label{r1}c\end{remarque}
\begin{corollaire}\label{c1}d\end{corollaire}
\chapter{Second}
\begin{corollaire}\label{c2}e\end{corollaire}
\begin{theoreme}\label{t2}f\end{theoreme}
RENVOIS \ref{p1} \ref{t1} \ref{r1} \ref{c1} \ref{c2} \ref{t2} ; \cref{p1} ; \cref{t1} ; \cref{r1} ; \cref{c1} ; \cref{c2} ; \cref{t2} FIN
\end{document}
"                sans correctif, crochet retiré
%   lualatex -jobname=amsseul  "\def\AMSSEUL{1}% Juge de la troisième passe : document minimal écrit de zéro pour la cause A (thmtools sous LaTeX 2026-06-01).
% Variantes, par une macro posée sur la ligne de commande (deux passes chacune, depuis memoire/essai/juge3/minimal/) :
%   lualatex -jobname=nu       "% Juge de la troisième passe : document minimal écrit de zéro pour la cause A (thmtools sous LaTeX 2026-06-01).
% Variantes, par une macro posée sur la ligne de commande (deux passes chacune, depuis memoire/essai/juge3/minimal/) :
%   lualatex -jobname=nu       "\input{minimal-juge3}"                                   sans correctif
%   lualatex -jobname=corr     "\def\CORR{1}\input{minimal-juge3}"                       correctif A du préambule
%   lualatex -jobname=crochet  "\def\SANSCROCHET{1}\input{minimal-juge3}"                sans correctif, crochet retiré
%   lualatex -jobname=amsseul  "\def\AMSSEUL{1}\input{minimal-juge3}"                    amsthm seul (témoin)
%   lualatex -jobname=faux     "\def\CORR{1}\def\FAUX{1}\input{minimal-juge3}"           deux vrais défauts : doivent tomber
%   lualatex -jobname=trace    "\def\TRACE{1}\input{minimal-juge3}"                      trace de l'appel
% Les lignes du correctif sont tirées du préambule commité par extraire_correctifs.py, pas retapées.
\documentclass{scrbook}
\ifdefined\SANSCROCHET \RemoveFromHook{package/amsthm/after}[firstaid/aliascounter]\fi
\usepackage{amsthm}
\usepackage{hyperref}
\usepackage[french, nameinlink, noabbrev]{cleveref}
\ifdefined\AMSSEUL
  \newtheorem{theoreme}{Théorème}[chapter]
  \newtheorem{proposition}[theoreme]{Proposition}
  \newtheorem{corollaire}[theoreme]{Corollaire}
  \newtheorem{remarque}{Remarque}[chapter]
\else
  \usepackage{thmtools}
  \ifdefined\CORR \input{correctif-a-pose}\fi
  \declaretheorem[numberwithin=chapter, name=Théorème]{theoreme}
  \ifdefined\TRACE \tracingmacros=2 \tracingonline=0 \fi
  \declaretheorem[sibling=theoreme, name=Proposition]{proposition}
  \ifdefined\TRACE \tracingmacros=0 \fi
  \declaretheorem[sibling=theoreme, name=Corollaire]{corollaire}
  \declaretheorem[numberwithin=chapter, name=Remarque]{remarque}
  \ifdefined\FAUX \newcounteralias{corollaire}{section}\newcounteralias{orphelin}{inexistant}\fi
  \ifdefined\CORR \input{correctif-a-remise}\fi
\fi
\crefname{theoreme}{théorème}{théorèmes}
\crefname{proposition}{proposition}{propositions}
\crefname{corollaire}{corollaire}{corollaires}
\crefname{remarque}{remarque}{remarques}
\begin{document}
\chapter{Premier}
\begin{proposition}\label{p1}a\end{proposition}
\begin{theoreme}\label{t1}b\end{theoreme}
\begin{remarque}\label{r1}c\end{remarque}
\begin{corollaire}\label{c1}d\end{corollaire}
\chapter{Second}
\begin{corollaire}\label{c2}e\end{corollaire}
\begin{theoreme}\label{t2}f\end{theoreme}
RENVOIS \ref{p1} \ref{t1} \ref{r1} \ref{c1} \ref{c2} \ref{t2} ; \cref{p1} ; \cref{t1} ; \cref{r1} ; \cref{c1} ; \cref{c2} ; \cref{t2} FIN
\end{document}
"                                   sans correctif
%   lualatex -jobname=corr     "\def\CORR{1}% Juge de la troisième passe : document minimal écrit de zéro pour la cause A (thmtools sous LaTeX 2026-06-01).
% Variantes, par une macro posée sur la ligne de commande (deux passes chacune, depuis memoire/essai/juge3/minimal/) :
%   lualatex -jobname=nu       "\input{minimal-juge3}"                                   sans correctif
%   lualatex -jobname=corr     "\def\CORR{1}\input{minimal-juge3}"                       correctif A du préambule
%   lualatex -jobname=crochet  "\def\SANSCROCHET{1}\input{minimal-juge3}"                sans correctif, crochet retiré
%   lualatex -jobname=amsseul  "\def\AMSSEUL{1}\input{minimal-juge3}"                    amsthm seul (témoin)
%   lualatex -jobname=faux     "\def\CORR{1}\def\FAUX{1}\input{minimal-juge3}"           deux vrais défauts : doivent tomber
%   lualatex -jobname=trace    "\def\TRACE{1}\input{minimal-juge3}"                      trace de l'appel
% Les lignes du correctif sont tirées du préambule commité par extraire_correctifs.py, pas retapées.
\documentclass{scrbook}
\ifdefined\SANSCROCHET \RemoveFromHook{package/amsthm/after}[firstaid/aliascounter]\fi
\usepackage{amsthm}
\usepackage{hyperref}
\usepackage[french, nameinlink, noabbrev]{cleveref}
\ifdefined\AMSSEUL
  \newtheorem{theoreme}{Théorème}[chapter]
  \newtheorem{proposition}[theoreme]{Proposition}
  \newtheorem{corollaire}[theoreme]{Corollaire}
  \newtheorem{remarque}{Remarque}[chapter]
\else
  \usepackage{thmtools}
  \ifdefined\CORR \input{correctif-a-pose}\fi
  \declaretheorem[numberwithin=chapter, name=Théorème]{theoreme}
  \ifdefined\TRACE \tracingmacros=2 \tracingonline=0 \fi
  \declaretheorem[sibling=theoreme, name=Proposition]{proposition}
  \ifdefined\TRACE \tracingmacros=0 \fi
  \declaretheorem[sibling=theoreme, name=Corollaire]{corollaire}
  \declaretheorem[numberwithin=chapter, name=Remarque]{remarque}
  \ifdefined\FAUX \newcounteralias{corollaire}{section}\newcounteralias{orphelin}{inexistant}\fi
  \ifdefined\CORR \input{correctif-a-remise}\fi
\fi
\crefname{theoreme}{théorème}{théorèmes}
\crefname{proposition}{proposition}{propositions}
\crefname{corollaire}{corollaire}{corollaires}
\crefname{remarque}{remarque}{remarques}
\begin{document}
\chapter{Premier}
\begin{proposition}\label{p1}a\end{proposition}
\begin{theoreme}\label{t1}b\end{theoreme}
\begin{remarque}\label{r1}c\end{remarque}
\begin{corollaire}\label{c1}d\end{corollaire}
\chapter{Second}
\begin{corollaire}\label{c2}e\end{corollaire}
\begin{theoreme}\label{t2}f\end{theoreme}
RENVOIS \ref{p1} \ref{t1} \ref{r1} \ref{c1} \ref{c2} \ref{t2} ; \cref{p1} ; \cref{t1} ; \cref{r1} ; \cref{c1} ; \cref{c2} ; \cref{t2} FIN
\end{document}
"                       correctif A du préambule
%   lualatex -jobname=crochet  "\def\SANSCROCHET{1}% Juge de la troisième passe : document minimal écrit de zéro pour la cause A (thmtools sous LaTeX 2026-06-01).
% Variantes, par une macro posée sur la ligne de commande (deux passes chacune, depuis memoire/essai/juge3/minimal/) :
%   lualatex -jobname=nu       "\input{minimal-juge3}"                                   sans correctif
%   lualatex -jobname=corr     "\def\CORR{1}\input{minimal-juge3}"                       correctif A du préambule
%   lualatex -jobname=crochet  "\def\SANSCROCHET{1}\input{minimal-juge3}"                sans correctif, crochet retiré
%   lualatex -jobname=amsseul  "\def\AMSSEUL{1}\input{minimal-juge3}"                    amsthm seul (témoin)
%   lualatex -jobname=faux     "\def\CORR{1}\def\FAUX{1}\input{minimal-juge3}"           deux vrais défauts : doivent tomber
%   lualatex -jobname=trace    "\def\TRACE{1}\input{minimal-juge3}"                      trace de l'appel
% Les lignes du correctif sont tirées du préambule commité par extraire_correctifs.py, pas retapées.
\documentclass{scrbook}
\ifdefined\SANSCROCHET \RemoveFromHook{package/amsthm/after}[firstaid/aliascounter]\fi
\usepackage{amsthm}
\usepackage{hyperref}
\usepackage[french, nameinlink, noabbrev]{cleveref}
\ifdefined\AMSSEUL
  \newtheorem{theoreme}{Théorème}[chapter]
  \newtheorem{proposition}[theoreme]{Proposition}
  \newtheorem{corollaire}[theoreme]{Corollaire}
  \newtheorem{remarque}{Remarque}[chapter]
\else
  \usepackage{thmtools}
  \ifdefined\CORR \input{correctif-a-pose}\fi
  \declaretheorem[numberwithin=chapter, name=Théorème]{theoreme}
  \ifdefined\TRACE \tracingmacros=2 \tracingonline=0 \fi
  \declaretheorem[sibling=theoreme, name=Proposition]{proposition}
  \ifdefined\TRACE \tracingmacros=0 \fi
  \declaretheorem[sibling=theoreme, name=Corollaire]{corollaire}
  \declaretheorem[numberwithin=chapter, name=Remarque]{remarque}
  \ifdefined\FAUX \newcounteralias{corollaire}{section}\newcounteralias{orphelin}{inexistant}\fi
  \ifdefined\CORR \input{correctif-a-remise}\fi
\fi
\crefname{theoreme}{théorème}{théorèmes}
\crefname{proposition}{proposition}{propositions}
\crefname{corollaire}{corollaire}{corollaires}
\crefname{remarque}{remarque}{remarques}
\begin{document}
\chapter{Premier}
\begin{proposition}\label{p1}a\end{proposition}
\begin{theoreme}\label{t1}b\end{theoreme}
\begin{remarque}\label{r1}c\end{remarque}
\begin{corollaire}\label{c1}d\end{corollaire}
\chapter{Second}
\begin{corollaire}\label{c2}e\end{corollaire}
\begin{theoreme}\label{t2}f\end{theoreme}
RENVOIS \ref{p1} \ref{t1} \ref{r1} \ref{c1} \ref{c2} \ref{t2} ; \cref{p1} ; \cref{t1} ; \cref{r1} ; \cref{c1} ; \cref{c2} ; \cref{t2} FIN
\end{document}
"                sans correctif, crochet retiré
%   lualatex -jobname=amsseul  "\def\AMSSEUL{1}% Juge de la troisième passe : document minimal écrit de zéro pour la cause A (thmtools sous LaTeX 2026-06-01).
% Variantes, par une macro posée sur la ligne de commande (deux passes chacune, depuis memoire/essai/juge3/minimal/) :
%   lualatex -jobname=nu       "\input{minimal-juge3}"                                   sans correctif
%   lualatex -jobname=corr     "\def\CORR{1}\input{minimal-juge3}"                       correctif A du préambule
%   lualatex -jobname=crochet  "\def\SANSCROCHET{1}\input{minimal-juge3}"                sans correctif, crochet retiré
%   lualatex -jobname=amsseul  "\def\AMSSEUL{1}\input{minimal-juge3}"                    amsthm seul (témoin)
%   lualatex -jobname=faux     "\def\CORR{1}\def\FAUX{1}\input{minimal-juge3}"           deux vrais défauts : doivent tomber
%   lualatex -jobname=trace    "\def\TRACE{1}\input{minimal-juge3}"                      trace de l'appel
% Les lignes du correctif sont tirées du préambule commité par extraire_correctifs.py, pas retapées.
\documentclass{scrbook}
\ifdefined\SANSCROCHET \RemoveFromHook{package/amsthm/after}[firstaid/aliascounter]\fi
\usepackage{amsthm}
\usepackage{hyperref}
\usepackage[french, nameinlink, noabbrev]{cleveref}
\ifdefined\AMSSEUL
  \newtheorem{theoreme}{Théorème}[chapter]
  \newtheorem{proposition}[theoreme]{Proposition}
  \newtheorem{corollaire}[theoreme]{Corollaire}
  \newtheorem{remarque}{Remarque}[chapter]
\else
  \usepackage{thmtools}
  \ifdefined\CORR \input{correctif-a-pose}\fi
  \declaretheorem[numberwithin=chapter, name=Théorème]{theoreme}
  \ifdefined\TRACE \tracingmacros=2 \tracingonline=0 \fi
  \declaretheorem[sibling=theoreme, name=Proposition]{proposition}
  \ifdefined\TRACE \tracingmacros=0 \fi
  \declaretheorem[sibling=theoreme, name=Corollaire]{corollaire}
  \declaretheorem[numberwithin=chapter, name=Remarque]{remarque}
  \ifdefined\FAUX \newcounteralias{corollaire}{section}\newcounteralias{orphelin}{inexistant}\fi
  \ifdefined\CORR \input{correctif-a-remise}\fi
\fi
\crefname{theoreme}{théorème}{théorèmes}
\crefname{proposition}{proposition}{propositions}
\crefname{corollaire}{corollaire}{corollaires}
\crefname{remarque}{remarque}{remarques}
\begin{document}
\chapter{Premier}
\begin{proposition}\label{p1}a\end{proposition}
\begin{theoreme}\label{t1}b\end{theoreme}
\begin{remarque}\label{r1}c\end{remarque}
\begin{corollaire}\label{c1}d\end{corollaire}
\chapter{Second}
\begin{corollaire}\label{c2}e\end{corollaire}
\begin{theoreme}\label{t2}f\end{theoreme}
RENVOIS \ref{p1} \ref{t1} \ref{r1} \ref{c1} \ref{c2} \ref{t2} ; \cref{p1} ; \cref{t1} ; \cref{r1} ; \cref{c1} ; \cref{c2} ; \cref{t2} FIN
\end{document}
"                    amsthm seul (témoin)
%   lualatex -jobname=faux     "\def\CORR{1}\def\FAUX{1}% Juge de la troisième passe : document minimal écrit de zéro pour la cause A (thmtools sous LaTeX 2026-06-01).
% Variantes, par une macro posée sur la ligne de commande (deux passes chacune, depuis memoire/essai/juge3/minimal/) :
%   lualatex -jobname=nu       "\input{minimal-juge3}"                                   sans correctif
%   lualatex -jobname=corr     "\def\CORR{1}\input{minimal-juge3}"                       correctif A du préambule
%   lualatex -jobname=crochet  "\def\SANSCROCHET{1}\input{minimal-juge3}"                sans correctif, crochet retiré
%   lualatex -jobname=amsseul  "\def\AMSSEUL{1}\input{minimal-juge3}"                    amsthm seul (témoin)
%   lualatex -jobname=faux     "\def\CORR{1}\def\FAUX{1}\input{minimal-juge3}"           deux vrais défauts : doivent tomber
%   lualatex -jobname=trace    "\def\TRACE{1}\input{minimal-juge3}"                      trace de l'appel
% Les lignes du correctif sont tirées du préambule commité par extraire_correctifs.py, pas retapées.
\documentclass{scrbook}
\ifdefined\SANSCROCHET \RemoveFromHook{package/amsthm/after}[firstaid/aliascounter]\fi
\usepackage{amsthm}
\usepackage{hyperref}
\usepackage[french, nameinlink, noabbrev]{cleveref}
\ifdefined\AMSSEUL
  \newtheorem{theoreme}{Théorème}[chapter]
  \newtheorem{proposition}[theoreme]{Proposition}
  \newtheorem{corollaire}[theoreme]{Corollaire}
  \newtheorem{remarque}{Remarque}[chapter]
\else
  \usepackage{thmtools}
  \ifdefined\CORR \input{correctif-a-pose}\fi
  \declaretheorem[numberwithin=chapter, name=Théorème]{theoreme}
  \ifdefined\TRACE \tracingmacros=2 \tracingonline=0 \fi
  \declaretheorem[sibling=theoreme, name=Proposition]{proposition}
  \ifdefined\TRACE \tracingmacros=0 \fi
  \declaretheorem[sibling=theoreme, name=Corollaire]{corollaire}
  \declaretheorem[numberwithin=chapter, name=Remarque]{remarque}
  \ifdefined\FAUX \newcounteralias{corollaire}{section}\newcounteralias{orphelin}{inexistant}\fi
  \ifdefined\CORR \input{correctif-a-remise}\fi
\fi
\crefname{theoreme}{théorème}{théorèmes}
\crefname{proposition}{proposition}{propositions}
\crefname{corollaire}{corollaire}{corollaires}
\crefname{remarque}{remarque}{remarques}
\begin{document}
\chapter{Premier}
\begin{proposition}\label{p1}a\end{proposition}
\begin{theoreme}\label{t1}b\end{theoreme}
\begin{remarque}\label{r1}c\end{remarque}
\begin{corollaire}\label{c1}d\end{corollaire}
\chapter{Second}
\begin{corollaire}\label{c2}e\end{corollaire}
\begin{theoreme}\label{t2}f\end{theoreme}
RENVOIS \ref{p1} \ref{t1} \ref{r1} \ref{c1} \ref{c2} \ref{t2} ; \cref{p1} ; \cref{t1} ; \cref{r1} ; \cref{c1} ; \cref{c2} ; \cref{t2} FIN
\end{document}
"           deux vrais défauts : doivent tomber
%   lualatex -jobname=trace    "\def\TRACE{1}% Juge de la troisième passe : document minimal écrit de zéro pour la cause A (thmtools sous LaTeX 2026-06-01).
% Variantes, par une macro posée sur la ligne de commande (deux passes chacune, depuis memoire/essai/juge3/minimal/) :
%   lualatex -jobname=nu       "\input{minimal-juge3}"                                   sans correctif
%   lualatex -jobname=corr     "\def\CORR{1}\input{minimal-juge3}"                       correctif A du préambule
%   lualatex -jobname=crochet  "\def\SANSCROCHET{1}\input{minimal-juge3}"                sans correctif, crochet retiré
%   lualatex -jobname=amsseul  "\def\AMSSEUL{1}\input{minimal-juge3}"                    amsthm seul (témoin)
%   lualatex -jobname=faux     "\def\CORR{1}\def\FAUX{1}\input{minimal-juge3}"           deux vrais défauts : doivent tomber
%   lualatex -jobname=trace    "\def\TRACE{1}\input{minimal-juge3}"                      trace de l'appel
% Les lignes du correctif sont tirées du préambule commité par extraire_correctifs.py, pas retapées.
\documentclass{scrbook}
\ifdefined\SANSCROCHET \RemoveFromHook{package/amsthm/after}[firstaid/aliascounter]\fi
\usepackage{amsthm}
\usepackage{hyperref}
\usepackage[french, nameinlink, noabbrev]{cleveref}
\ifdefined\AMSSEUL
  \newtheorem{theoreme}{Théorème}[chapter]
  \newtheorem{proposition}[theoreme]{Proposition}
  \newtheorem{corollaire}[theoreme]{Corollaire}
  \newtheorem{remarque}{Remarque}[chapter]
\else
  \usepackage{thmtools}
  \ifdefined\CORR \input{correctif-a-pose}\fi
  \declaretheorem[numberwithin=chapter, name=Théorème]{theoreme}
  \ifdefined\TRACE \tracingmacros=2 \tracingonline=0 \fi
  \declaretheorem[sibling=theoreme, name=Proposition]{proposition}
  \ifdefined\TRACE \tracingmacros=0 \fi
  \declaretheorem[sibling=theoreme, name=Corollaire]{corollaire}
  \declaretheorem[numberwithin=chapter, name=Remarque]{remarque}
  \ifdefined\FAUX \newcounteralias{corollaire}{section}\newcounteralias{orphelin}{inexistant}\fi
  \ifdefined\CORR \input{correctif-a-remise}\fi
\fi
\crefname{theoreme}{théorème}{théorèmes}
\crefname{proposition}{proposition}{propositions}
\crefname{corollaire}{corollaire}{corollaires}
\crefname{remarque}{remarque}{remarques}
\begin{document}
\chapter{Premier}
\begin{proposition}\label{p1}a\end{proposition}
\begin{theoreme}\label{t1}b\end{theoreme}
\begin{remarque}\label{r1}c\end{remarque}
\begin{corollaire}\label{c1}d\end{corollaire}
\chapter{Second}
\begin{corollaire}\label{c2}e\end{corollaire}
\begin{theoreme}\label{t2}f\end{theoreme}
RENVOIS \ref{p1} \ref{t1} \ref{r1} \ref{c1} \ref{c2} \ref{t2} ; \cref{p1} ; \cref{t1} ; \cref{r1} ; \cref{c1} ; \cref{c2} ; \cref{t2} FIN
\end{document}
"                      trace de l'appel
% Les lignes du correctif sont tirées du préambule commité par extraire_correctifs.py, pas retapées.
\documentclass{scrbook}
\ifdefined\SANSCROCHET \RemoveFromHook{package/amsthm/after}[firstaid/aliascounter]\fi
\usepackage{amsthm}
\usepackage{hyperref}
\usepackage[french, nameinlink, noabbrev]{cleveref}
\ifdefined\AMSSEUL
  \newtheorem{theoreme}{Théorème}[chapter]
  \newtheorem{proposition}[theoreme]{Proposition}
  \newtheorem{corollaire}[theoreme]{Corollaire}
  \newtheorem{remarque}{Remarque}[chapter]
\else
  \usepackage{thmtools}
  \ifdefined\CORR \input{correctif-a-pose}\fi
  \declaretheorem[numberwithin=chapter, name=Théorème]{theoreme}
  \ifdefined\TRACE \tracingmacros=2 \tracingonline=0 \fi
  \declaretheorem[sibling=theoreme, name=Proposition]{proposition}
  \ifdefined\TRACE \tracingmacros=0 \fi
  \declaretheorem[sibling=theoreme, name=Corollaire]{corollaire}
  \declaretheorem[numberwithin=chapter, name=Remarque]{remarque}
  \ifdefined\FAUX \newcounteralias{corollaire}{section}\newcounteralias{orphelin}{inexistant}\fi
  \ifdefined\CORR \input{correctif-a-remise}\fi
\fi
\crefname{theoreme}{théorème}{théorèmes}
\crefname{proposition}{proposition}{propositions}
\crefname{corollaire}{corollaire}{corollaires}
\crefname{remarque}{remarque}{remarques}
\begin{document}
\chapter{Premier}
\begin{proposition}\label{p1}a\end{proposition}
\begin{theoreme}\label{t1}b\end{theoreme}
\begin{remarque}\label{r1}c\end{remarque}
\begin{corollaire}\label{c1}d\end{corollaire}
\chapter{Second}
\begin{corollaire}\label{c2}e\end{corollaire}
\begin{theoreme}\label{t2}f\end{theoreme}
RENVOIS \ref{p1} \ref{t1} \ref{r1} \ref{c1} \ref{c2} \ref{t2} ; \cref{p1} ; \cref{t1} ; \cref{r1} ; \cref{c1} ; \cref{c2} ; \cref{t2} FIN
\end{document}
"                    amsthm seul (témoin)
%   lualatex -jobname=faux     "\def\CORR{1}\def\FAUX{1}% Juge de la troisième passe : document minimal écrit de zéro pour la cause A (thmtools sous LaTeX 2026-06-01).
% Variantes, par une macro posée sur la ligne de commande (deux passes chacune, depuis memoire/essai/juge3/minimal/) :
%   lualatex -jobname=nu       "% Juge de la troisième passe : document minimal écrit de zéro pour la cause A (thmtools sous LaTeX 2026-06-01).
% Variantes, par une macro posée sur la ligne de commande (deux passes chacune, depuis memoire/essai/juge3/minimal/) :
%   lualatex -jobname=nu       "\input{minimal-juge3}"                                   sans correctif
%   lualatex -jobname=corr     "\def\CORR{1}\input{minimal-juge3}"                       correctif A du préambule
%   lualatex -jobname=crochet  "\def\SANSCROCHET{1}\input{minimal-juge3}"                sans correctif, crochet retiré
%   lualatex -jobname=amsseul  "\def\AMSSEUL{1}\input{minimal-juge3}"                    amsthm seul (témoin)
%   lualatex -jobname=faux     "\def\CORR{1}\def\FAUX{1}\input{minimal-juge3}"           deux vrais défauts : doivent tomber
%   lualatex -jobname=trace    "\def\TRACE{1}\input{minimal-juge3}"                      trace de l'appel
% Les lignes du correctif sont tirées du préambule commité par extraire_correctifs.py, pas retapées.
\documentclass{scrbook}
\ifdefined\SANSCROCHET \RemoveFromHook{package/amsthm/after}[firstaid/aliascounter]\fi
\usepackage{amsthm}
\usepackage{hyperref}
\usepackage[french, nameinlink, noabbrev]{cleveref}
\ifdefined\AMSSEUL
  \newtheorem{theoreme}{Théorème}[chapter]
  \newtheorem{proposition}[theoreme]{Proposition}
  \newtheorem{corollaire}[theoreme]{Corollaire}
  \newtheorem{remarque}{Remarque}[chapter]
\else
  \usepackage{thmtools}
  \ifdefined\CORR \input{correctif-a-pose}\fi
  \declaretheorem[numberwithin=chapter, name=Théorème]{theoreme}
  \ifdefined\TRACE \tracingmacros=2 \tracingonline=0 \fi
  \declaretheorem[sibling=theoreme, name=Proposition]{proposition}
  \ifdefined\TRACE \tracingmacros=0 \fi
  \declaretheorem[sibling=theoreme, name=Corollaire]{corollaire}
  \declaretheorem[numberwithin=chapter, name=Remarque]{remarque}
  \ifdefined\FAUX \newcounteralias{corollaire}{section}\newcounteralias{orphelin}{inexistant}\fi
  \ifdefined\CORR \input{correctif-a-remise}\fi
\fi
\crefname{theoreme}{théorème}{théorèmes}
\crefname{proposition}{proposition}{propositions}
\crefname{corollaire}{corollaire}{corollaires}
\crefname{remarque}{remarque}{remarques}
\begin{document}
\chapter{Premier}
\begin{proposition}\label{p1}a\end{proposition}
\begin{theoreme}\label{t1}b\end{theoreme}
\begin{remarque}\label{r1}c\end{remarque}
\begin{corollaire}\label{c1}d\end{corollaire}
\chapter{Second}
\begin{corollaire}\label{c2}e\end{corollaire}
\begin{theoreme}\label{t2}f\end{theoreme}
RENVOIS \ref{p1} \ref{t1} \ref{r1} \ref{c1} \ref{c2} \ref{t2} ; \cref{p1} ; \cref{t1} ; \cref{r1} ; \cref{c1} ; \cref{c2} ; \cref{t2} FIN
\end{document}
"                                   sans correctif
%   lualatex -jobname=corr     "\def\CORR{1}% Juge de la troisième passe : document minimal écrit de zéro pour la cause A (thmtools sous LaTeX 2026-06-01).
% Variantes, par une macro posée sur la ligne de commande (deux passes chacune, depuis memoire/essai/juge3/minimal/) :
%   lualatex -jobname=nu       "\input{minimal-juge3}"                                   sans correctif
%   lualatex -jobname=corr     "\def\CORR{1}\input{minimal-juge3}"                       correctif A du préambule
%   lualatex -jobname=crochet  "\def\SANSCROCHET{1}\input{minimal-juge3}"                sans correctif, crochet retiré
%   lualatex -jobname=amsseul  "\def\AMSSEUL{1}\input{minimal-juge3}"                    amsthm seul (témoin)
%   lualatex -jobname=faux     "\def\CORR{1}\def\FAUX{1}\input{minimal-juge3}"           deux vrais défauts : doivent tomber
%   lualatex -jobname=trace    "\def\TRACE{1}\input{minimal-juge3}"                      trace de l'appel
% Les lignes du correctif sont tirées du préambule commité par extraire_correctifs.py, pas retapées.
\documentclass{scrbook}
\ifdefined\SANSCROCHET \RemoveFromHook{package/amsthm/after}[firstaid/aliascounter]\fi
\usepackage{amsthm}
\usepackage{hyperref}
\usepackage[french, nameinlink, noabbrev]{cleveref}
\ifdefined\AMSSEUL
  \newtheorem{theoreme}{Théorème}[chapter]
  \newtheorem{proposition}[theoreme]{Proposition}
  \newtheorem{corollaire}[theoreme]{Corollaire}
  \newtheorem{remarque}{Remarque}[chapter]
\else
  \usepackage{thmtools}
  \ifdefined\CORR \input{correctif-a-pose}\fi
  \declaretheorem[numberwithin=chapter, name=Théorème]{theoreme}
  \ifdefined\TRACE \tracingmacros=2 \tracingonline=0 \fi
  \declaretheorem[sibling=theoreme, name=Proposition]{proposition}
  \ifdefined\TRACE \tracingmacros=0 \fi
  \declaretheorem[sibling=theoreme, name=Corollaire]{corollaire}
  \declaretheorem[numberwithin=chapter, name=Remarque]{remarque}
  \ifdefined\FAUX \newcounteralias{corollaire}{section}\newcounteralias{orphelin}{inexistant}\fi
  \ifdefined\CORR \input{correctif-a-remise}\fi
\fi
\crefname{theoreme}{théorème}{théorèmes}
\crefname{proposition}{proposition}{propositions}
\crefname{corollaire}{corollaire}{corollaires}
\crefname{remarque}{remarque}{remarques}
\begin{document}
\chapter{Premier}
\begin{proposition}\label{p1}a\end{proposition}
\begin{theoreme}\label{t1}b\end{theoreme}
\begin{remarque}\label{r1}c\end{remarque}
\begin{corollaire}\label{c1}d\end{corollaire}
\chapter{Second}
\begin{corollaire}\label{c2}e\end{corollaire}
\begin{theoreme}\label{t2}f\end{theoreme}
RENVOIS \ref{p1} \ref{t1} \ref{r1} \ref{c1} \ref{c2} \ref{t2} ; \cref{p1} ; \cref{t1} ; \cref{r1} ; \cref{c1} ; \cref{c2} ; \cref{t2} FIN
\end{document}
"                       correctif A du préambule
%   lualatex -jobname=crochet  "\def\SANSCROCHET{1}% Juge de la troisième passe : document minimal écrit de zéro pour la cause A (thmtools sous LaTeX 2026-06-01).
% Variantes, par une macro posée sur la ligne de commande (deux passes chacune, depuis memoire/essai/juge3/minimal/) :
%   lualatex -jobname=nu       "\input{minimal-juge3}"                                   sans correctif
%   lualatex -jobname=corr     "\def\CORR{1}\input{minimal-juge3}"                       correctif A du préambule
%   lualatex -jobname=crochet  "\def\SANSCROCHET{1}\input{minimal-juge3}"                sans correctif, crochet retiré
%   lualatex -jobname=amsseul  "\def\AMSSEUL{1}\input{minimal-juge3}"                    amsthm seul (témoin)
%   lualatex -jobname=faux     "\def\CORR{1}\def\FAUX{1}\input{minimal-juge3}"           deux vrais défauts : doivent tomber
%   lualatex -jobname=trace    "\def\TRACE{1}\input{minimal-juge3}"                      trace de l'appel
% Les lignes du correctif sont tirées du préambule commité par extraire_correctifs.py, pas retapées.
\documentclass{scrbook}
\ifdefined\SANSCROCHET \RemoveFromHook{package/amsthm/after}[firstaid/aliascounter]\fi
\usepackage{amsthm}
\usepackage{hyperref}
\usepackage[french, nameinlink, noabbrev]{cleveref}
\ifdefined\AMSSEUL
  \newtheorem{theoreme}{Théorème}[chapter]
  \newtheorem{proposition}[theoreme]{Proposition}
  \newtheorem{corollaire}[theoreme]{Corollaire}
  \newtheorem{remarque}{Remarque}[chapter]
\else
  \usepackage{thmtools}
  \ifdefined\CORR \input{correctif-a-pose}\fi
  \declaretheorem[numberwithin=chapter, name=Théorème]{theoreme}
  \ifdefined\TRACE \tracingmacros=2 \tracingonline=0 \fi
  \declaretheorem[sibling=theoreme, name=Proposition]{proposition}
  \ifdefined\TRACE \tracingmacros=0 \fi
  \declaretheorem[sibling=theoreme, name=Corollaire]{corollaire}
  \declaretheorem[numberwithin=chapter, name=Remarque]{remarque}
  \ifdefined\FAUX \newcounteralias{corollaire}{section}\newcounteralias{orphelin}{inexistant}\fi
  \ifdefined\CORR \input{correctif-a-remise}\fi
\fi
\crefname{theoreme}{théorème}{théorèmes}
\crefname{proposition}{proposition}{propositions}
\crefname{corollaire}{corollaire}{corollaires}
\crefname{remarque}{remarque}{remarques}
\begin{document}
\chapter{Premier}
\begin{proposition}\label{p1}a\end{proposition}
\begin{theoreme}\label{t1}b\end{theoreme}
\begin{remarque}\label{r1}c\end{remarque}
\begin{corollaire}\label{c1}d\end{corollaire}
\chapter{Second}
\begin{corollaire}\label{c2}e\end{corollaire}
\begin{theoreme}\label{t2}f\end{theoreme}
RENVOIS \ref{p1} \ref{t1} \ref{r1} \ref{c1} \ref{c2} \ref{t2} ; \cref{p1} ; \cref{t1} ; \cref{r1} ; \cref{c1} ; \cref{c2} ; \cref{t2} FIN
\end{document}
"                sans correctif, crochet retiré
%   lualatex -jobname=amsseul  "\def\AMSSEUL{1}% Juge de la troisième passe : document minimal écrit de zéro pour la cause A (thmtools sous LaTeX 2026-06-01).
% Variantes, par une macro posée sur la ligne de commande (deux passes chacune, depuis memoire/essai/juge3/minimal/) :
%   lualatex -jobname=nu       "\input{minimal-juge3}"                                   sans correctif
%   lualatex -jobname=corr     "\def\CORR{1}\input{minimal-juge3}"                       correctif A du préambule
%   lualatex -jobname=crochet  "\def\SANSCROCHET{1}\input{minimal-juge3}"                sans correctif, crochet retiré
%   lualatex -jobname=amsseul  "\def\AMSSEUL{1}\input{minimal-juge3}"                    amsthm seul (témoin)
%   lualatex -jobname=faux     "\def\CORR{1}\def\FAUX{1}\input{minimal-juge3}"           deux vrais défauts : doivent tomber
%   lualatex -jobname=trace    "\def\TRACE{1}\input{minimal-juge3}"                      trace de l'appel
% Les lignes du correctif sont tirées du préambule commité par extraire_correctifs.py, pas retapées.
\documentclass{scrbook}
\ifdefined\SANSCROCHET \RemoveFromHook{package/amsthm/after}[firstaid/aliascounter]\fi
\usepackage{amsthm}
\usepackage{hyperref}
\usepackage[french, nameinlink, noabbrev]{cleveref}
\ifdefined\AMSSEUL
  \newtheorem{theoreme}{Théorème}[chapter]
  \newtheorem{proposition}[theoreme]{Proposition}
  \newtheorem{corollaire}[theoreme]{Corollaire}
  \newtheorem{remarque}{Remarque}[chapter]
\else
  \usepackage{thmtools}
  \ifdefined\CORR \input{correctif-a-pose}\fi
  \declaretheorem[numberwithin=chapter, name=Théorème]{theoreme}
  \ifdefined\TRACE \tracingmacros=2 \tracingonline=0 \fi
  \declaretheorem[sibling=theoreme, name=Proposition]{proposition}
  \ifdefined\TRACE \tracingmacros=0 \fi
  \declaretheorem[sibling=theoreme, name=Corollaire]{corollaire}
  \declaretheorem[numberwithin=chapter, name=Remarque]{remarque}
  \ifdefined\FAUX \newcounteralias{corollaire}{section}\newcounteralias{orphelin}{inexistant}\fi
  \ifdefined\CORR \input{correctif-a-remise}\fi
\fi
\crefname{theoreme}{théorème}{théorèmes}
\crefname{proposition}{proposition}{propositions}
\crefname{corollaire}{corollaire}{corollaires}
\crefname{remarque}{remarque}{remarques}
\begin{document}
\chapter{Premier}
\begin{proposition}\label{p1}a\end{proposition}
\begin{theoreme}\label{t1}b\end{theoreme}
\begin{remarque}\label{r1}c\end{remarque}
\begin{corollaire}\label{c1}d\end{corollaire}
\chapter{Second}
\begin{corollaire}\label{c2}e\end{corollaire}
\begin{theoreme}\label{t2}f\end{theoreme}
RENVOIS \ref{p1} \ref{t1} \ref{r1} \ref{c1} \ref{c2} \ref{t2} ; \cref{p1} ; \cref{t1} ; \cref{r1} ; \cref{c1} ; \cref{c2} ; \cref{t2} FIN
\end{document}
"                    amsthm seul (témoin)
%   lualatex -jobname=faux     "\def\CORR{1}\def\FAUX{1}% Juge de la troisième passe : document minimal écrit de zéro pour la cause A (thmtools sous LaTeX 2026-06-01).
% Variantes, par une macro posée sur la ligne de commande (deux passes chacune, depuis memoire/essai/juge3/minimal/) :
%   lualatex -jobname=nu       "\input{minimal-juge3}"                                   sans correctif
%   lualatex -jobname=corr     "\def\CORR{1}\input{minimal-juge3}"                       correctif A du préambule
%   lualatex -jobname=crochet  "\def\SANSCROCHET{1}\input{minimal-juge3}"                sans correctif, crochet retiré
%   lualatex -jobname=amsseul  "\def\AMSSEUL{1}\input{minimal-juge3}"                    amsthm seul (témoin)
%   lualatex -jobname=faux     "\def\CORR{1}\def\FAUX{1}\input{minimal-juge3}"           deux vrais défauts : doivent tomber
%   lualatex -jobname=trace    "\def\TRACE{1}\input{minimal-juge3}"                      trace de l'appel
% Les lignes du correctif sont tirées du préambule commité par extraire_correctifs.py, pas retapées.
\documentclass{scrbook}
\ifdefined\SANSCROCHET \RemoveFromHook{package/amsthm/after}[firstaid/aliascounter]\fi
\usepackage{amsthm}
\usepackage{hyperref}
\usepackage[french, nameinlink, noabbrev]{cleveref}
\ifdefined\AMSSEUL
  \newtheorem{theoreme}{Théorème}[chapter]
  \newtheorem{proposition}[theoreme]{Proposition}
  \newtheorem{corollaire}[theoreme]{Corollaire}
  \newtheorem{remarque}{Remarque}[chapter]
\else
  \usepackage{thmtools}
  \ifdefined\CORR \input{correctif-a-pose}\fi
  \declaretheorem[numberwithin=chapter, name=Théorème]{theoreme}
  \ifdefined\TRACE \tracingmacros=2 \tracingonline=0 \fi
  \declaretheorem[sibling=theoreme, name=Proposition]{proposition}
  \ifdefined\TRACE \tracingmacros=0 \fi
  \declaretheorem[sibling=theoreme, name=Corollaire]{corollaire}
  \declaretheorem[numberwithin=chapter, name=Remarque]{remarque}
  \ifdefined\FAUX \newcounteralias{corollaire}{section}\newcounteralias{orphelin}{inexistant}\fi
  \ifdefined\CORR \input{correctif-a-remise}\fi
\fi
\crefname{theoreme}{théorème}{théorèmes}
\crefname{proposition}{proposition}{propositions}
\crefname{corollaire}{corollaire}{corollaires}
\crefname{remarque}{remarque}{remarques}
\begin{document}
\chapter{Premier}
\begin{proposition}\label{p1}a\end{proposition}
\begin{theoreme}\label{t1}b\end{theoreme}
\begin{remarque}\label{r1}c\end{remarque}
\begin{corollaire}\label{c1}d\end{corollaire}
\chapter{Second}
\begin{corollaire}\label{c2}e\end{corollaire}
\begin{theoreme}\label{t2}f\end{theoreme}
RENVOIS \ref{p1} \ref{t1} \ref{r1} \ref{c1} \ref{c2} \ref{t2} ; \cref{p1} ; \cref{t1} ; \cref{r1} ; \cref{c1} ; \cref{c2} ; \cref{t2} FIN
\end{document}
"           deux vrais défauts : doivent tomber
%   lualatex -jobname=trace    "\def\TRACE{1}% Juge de la troisième passe : document minimal écrit de zéro pour la cause A (thmtools sous LaTeX 2026-06-01).
% Variantes, par une macro posée sur la ligne de commande (deux passes chacune, depuis memoire/essai/juge3/minimal/) :
%   lualatex -jobname=nu       "\input{minimal-juge3}"                                   sans correctif
%   lualatex -jobname=corr     "\def\CORR{1}\input{minimal-juge3}"                       correctif A du préambule
%   lualatex -jobname=crochet  "\def\SANSCROCHET{1}\input{minimal-juge3}"                sans correctif, crochet retiré
%   lualatex -jobname=amsseul  "\def\AMSSEUL{1}\input{minimal-juge3}"                    amsthm seul (témoin)
%   lualatex -jobname=faux     "\def\CORR{1}\def\FAUX{1}\input{minimal-juge3}"           deux vrais défauts : doivent tomber
%   lualatex -jobname=trace    "\def\TRACE{1}\input{minimal-juge3}"                      trace de l'appel
% Les lignes du correctif sont tirées du préambule commité par extraire_correctifs.py, pas retapées.
\documentclass{scrbook}
\ifdefined\SANSCROCHET \RemoveFromHook{package/amsthm/after}[firstaid/aliascounter]\fi
\usepackage{amsthm}
\usepackage{hyperref}
\usepackage[french, nameinlink, noabbrev]{cleveref}
\ifdefined\AMSSEUL
  \newtheorem{theoreme}{Théorème}[chapter]
  \newtheorem{proposition}[theoreme]{Proposition}
  \newtheorem{corollaire}[theoreme]{Corollaire}
  \newtheorem{remarque}{Remarque}[chapter]
\else
  \usepackage{thmtools}
  \ifdefined\CORR \input{correctif-a-pose}\fi
  \declaretheorem[numberwithin=chapter, name=Théorème]{theoreme}
  \ifdefined\TRACE \tracingmacros=2 \tracingonline=0 \fi
  \declaretheorem[sibling=theoreme, name=Proposition]{proposition}
  \ifdefined\TRACE \tracingmacros=0 \fi
  \declaretheorem[sibling=theoreme, name=Corollaire]{corollaire}
  \declaretheorem[numberwithin=chapter, name=Remarque]{remarque}
  \ifdefined\FAUX \newcounteralias{corollaire}{section}\newcounteralias{orphelin}{inexistant}\fi
  \ifdefined\CORR \input{correctif-a-remise}\fi
\fi
\crefname{theoreme}{théorème}{théorèmes}
\crefname{proposition}{proposition}{propositions}
\crefname{corollaire}{corollaire}{corollaires}
\crefname{remarque}{remarque}{remarques}
\begin{document}
\chapter{Premier}
\begin{proposition}\label{p1}a\end{proposition}
\begin{theoreme}\label{t1}b\end{theoreme}
\begin{remarque}\label{r1}c\end{remarque}
\begin{corollaire}\label{c1}d\end{corollaire}
\chapter{Second}
\begin{corollaire}\label{c2}e\end{corollaire}
\begin{theoreme}\label{t2}f\end{theoreme}
RENVOIS \ref{p1} \ref{t1} \ref{r1} \ref{c1} \ref{c2} \ref{t2} ; \cref{p1} ; \cref{t1} ; \cref{r1} ; \cref{c1} ; \cref{c2} ; \cref{t2} FIN
\end{document}
"                      trace de l'appel
% Les lignes du correctif sont tirées du préambule commité par extraire_correctifs.py, pas retapées.
\documentclass{scrbook}
\ifdefined\SANSCROCHET \RemoveFromHook{package/amsthm/after}[firstaid/aliascounter]\fi
\usepackage{amsthm}
\usepackage{hyperref}
\usepackage[french, nameinlink, noabbrev]{cleveref}
\ifdefined\AMSSEUL
  \newtheorem{theoreme}{Théorème}[chapter]
  \newtheorem{proposition}[theoreme]{Proposition}
  \newtheorem{corollaire}[theoreme]{Corollaire}
  \newtheorem{remarque}{Remarque}[chapter]
\else
  \usepackage{thmtools}
  \ifdefined\CORR \input{correctif-a-pose}\fi
  \declaretheorem[numberwithin=chapter, name=Théorème]{theoreme}
  \ifdefined\TRACE \tracingmacros=2 \tracingonline=0 \fi
  \declaretheorem[sibling=theoreme, name=Proposition]{proposition}
  \ifdefined\TRACE \tracingmacros=0 \fi
  \declaretheorem[sibling=theoreme, name=Corollaire]{corollaire}
  \declaretheorem[numberwithin=chapter, name=Remarque]{remarque}
  \ifdefined\FAUX \newcounteralias{corollaire}{section}\newcounteralias{orphelin}{inexistant}\fi
  \ifdefined\CORR \input{correctif-a-remise}\fi
\fi
\crefname{theoreme}{théorème}{théorèmes}
\crefname{proposition}{proposition}{propositions}
\crefname{corollaire}{corollaire}{corollaires}
\crefname{remarque}{remarque}{remarques}
\begin{document}
\chapter{Premier}
\begin{proposition}\label{p1}a\end{proposition}
\begin{theoreme}\label{t1}b\end{theoreme}
\begin{remarque}\label{r1}c\end{remarque}
\begin{corollaire}\label{c1}d\end{corollaire}
\chapter{Second}
\begin{corollaire}\label{c2}e\end{corollaire}
\begin{theoreme}\label{t2}f\end{theoreme}
RENVOIS \ref{p1} \ref{t1} \ref{r1} \ref{c1} \ref{c2} \ref{t2} ; \cref{p1} ; \cref{t1} ; \cref{r1} ; \cref{c1} ; \cref{c2} ; \cref{t2} FIN
\end{document}
"           deux vrais défauts : doivent tomber
%   lualatex -jobname=trace    "\def\TRACE{1}% Juge de la troisième passe : document minimal écrit de zéro pour la cause A (thmtools sous LaTeX 2026-06-01).
% Variantes, par une macro posée sur la ligne de commande (deux passes chacune, depuis memoire/essai/juge3/minimal/) :
%   lualatex -jobname=nu       "% Juge de la troisième passe : document minimal écrit de zéro pour la cause A (thmtools sous LaTeX 2026-06-01).
% Variantes, par une macro posée sur la ligne de commande (deux passes chacune, depuis memoire/essai/juge3/minimal/) :
%   lualatex -jobname=nu       "\input{minimal-juge3}"                                   sans correctif
%   lualatex -jobname=corr     "\def\CORR{1}\input{minimal-juge3}"                       correctif A du préambule
%   lualatex -jobname=crochet  "\def\SANSCROCHET{1}\input{minimal-juge3}"                sans correctif, crochet retiré
%   lualatex -jobname=amsseul  "\def\AMSSEUL{1}\input{minimal-juge3}"                    amsthm seul (témoin)
%   lualatex -jobname=faux     "\def\CORR{1}\def\FAUX{1}\input{minimal-juge3}"           deux vrais défauts : doivent tomber
%   lualatex -jobname=trace    "\def\TRACE{1}\input{minimal-juge3}"                      trace de l'appel
% Les lignes du correctif sont tirées du préambule commité par extraire_correctifs.py, pas retapées.
\documentclass{scrbook}
\ifdefined\SANSCROCHET \RemoveFromHook{package/amsthm/after}[firstaid/aliascounter]\fi
\usepackage{amsthm}
\usepackage{hyperref}
\usepackage[french, nameinlink, noabbrev]{cleveref}
\ifdefined\AMSSEUL
  \newtheorem{theoreme}{Théorème}[chapter]
  \newtheorem{proposition}[theoreme]{Proposition}
  \newtheorem{corollaire}[theoreme]{Corollaire}
  \newtheorem{remarque}{Remarque}[chapter]
\else
  \usepackage{thmtools}
  \ifdefined\CORR \input{correctif-a-pose}\fi
  \declaretheorem[numberwithin=chapter, name=Théorème]{theoreme}
  \ifdefined\TRACE \tracingmacros=2 \tracingonline=0 \fi
  \declaretheorem[sibling=theoreme, name=Proposition]{proposition}
  \ifdefined\TRACE \tracingmacros=0 \fi
  \declaretheorem[sibling=theoreme, name=Corollaire]{corollaire}
  \declaretheorem[numberwithin=chapter, name=Remarque]{remarque}
  \ifdefined\FAUX \newcounteralias{corollaire}{section}\newcounteralias{orphelin}{inexistant}\fi
  \ifdefined\CORR \input{correctif-a-remise}\fi
\fi
\crefname{theoreme}{théorème}{théorèmes}
\crefname{proposition}{proposition}{propositions}
\crefname{corollaire}{corollaire}{corollaires}
\crefname{remarque}{remarque}{remarques}
\begin{document}
\chapter{Premier}
\begin{proposition}\label{p1}a\end{proposition}
\begin{theoreme}\label{t1}b\end{theoreme}
\begin{remarque}\label{r1}c\end{remarque}
\begin{corollaire}\label{c1}d\end{corollaire}
\chapter{Second}
\begin{corollaire}\label{c2}e\end{corollaire}
\begin{theoreme}\label{t2}f\end{theoreme}
RENVOIS \ref{p1} \ref{t1} \ref{r1} \ref{c1} \ref{c2} \ref{t2} ; \cref{p1} ; \cref{t1} ; \cref{r1} ; \cref{c1} ; \cref{c2} ; \cref{t2} FIN
\end{document}
"                                   sans correctif
%   lualatex -jobname=corr     "\def\CORR{1}% Juge de la troisième passe : document minimal écrit de zéro pour la cause A (thmtools sous LaTeX 2026-06-01).
% Variantes, par une macro posée sur la ligne de commande (deux passes chacune, depuis memoire/essai/juge3/minimal/) :
%   lualatex -jobname=nu       "\input{minimal-juge3}"                                   sans correctif
%   lualatex -jobname=corr     "\def\CORR{1}\input{minimal-juge3}"                       correctif A du préambule
%   lualatex -jobname=crochet  "\def\SANSCROCHET{1}\input{minimal-juge3}"                sans correctif, crochet retiré
%   lualatex -jobname=amsseul  "\def\AMSSEUL{1}\input{minimal-juge3}"                    amsthm seul (témoin)
%   lualatex -jobname=faux     "\def\CORR{1}\def\FAUX{1}\input{minimal-juge3}"           deux vrais défauts : doivent tomber
%   lualatex -jobname=trace    "\def\TRACE{1}\input{minimal-juge3}"                      trace de l'appel
% Les lignes du correctif sont tirées du préambule commité par extraire_correctifs.py, pas retapées.
\documentclass{scrbook}
\ifdefined\SANSCROCHET \RemoveFromHook{package/amsthm/after}[firstaid/aliascounter]\fi
\usepackage{amsthm}
\usepackage{hyperref}
\usepackage[french, nameinlink, noabbrev]{cleveref}
\ifdefined\AMSSEUL
  \newtheorem{theoreme}{Théorème}[chapter]
  \newtheorem{proposition}[theoreme]{Proposition}
  \newtheorem{corollaire}[theoreme]{Corollaire}
  \newtheorem{remarque}{Remarque}[chapter]
\else
  \usepackage{thmtools}
  \ifdefined\CORR \input{correctif-a-pose}\fi
  \declaretheorem[numberwithin=chapter, name=Théorème]{theoreme}
  \ifdefined\TRACE \tracingmacros=2 \tracingonline=0 \fi
  \declaretheorem[sibling=theoreme, name=Proposition]{proposition}
  \ifdefined\TRACE \tracingmacros=0 \fi
  \declaretheorem[sibling=theoreme, name=Corollaire]{corollaire}
  \declaretheorem[numberwithin=chapter, name=Remarque]{remarque}
  \ifdefined\FAUX \newcounteralias{corollaire}{section}\newcounteralias{orphelin}{inexistant}\fi
  \ifdefined\CORR \input{correctif-a-remise}\fi
\fi
\crefname{theoreme}{théorème}{théorèmes}
\crefname{proposition}{proposition}{propositions}
\crefname{corollaire}{corollaire}{corollaires}
\crefname{remarque}{remarque}{remarques}
\begin{document}
\chapter{Premier}
\begin{proposition}\label{p1}a\end{proposition}
\begin{theoreme}\label{t1}b\end{theoreme}
\begin{remarque}\label{r1}c\end{remarque}
\begin{corollaire}\label{c1}d\end{corollaire}
\chapter{Second}
\begin{corollaire}\label{c2}e\end{corollaire}
\begin{theoreme}\label{t2}f\end{theoreme}
RENVOIS \ref{p1} \ref{t1} \ref{r1} \ref{c1} \ref{c2} \ref{t2} ; \cref{p1} ; \cref{t1} ; \cref{r1} ; \cref{c1} ; \cref{c2} ; \cref{t2} FIN
\end{document}
"                       correctif A du préambule
%   lualatex -jobname=crochet  "\def\SANSCROCHET{1}% Juge de la troisième passe : document minimal écrit de zéro pour la cause A (thmtools sous LaTeX 2026-06-01).
% Variantes, par une macro posée sur la ligne de commande (deux passes chacune, depuis memoire/essai/juge3/minimal/) :
%   lualatex -jobname=nu       "\input{minimal-juge3}"                                   sans correctif
%   lualatex -jobname=corr     "\def\CORR{1}\input{minimal-juge3}"                       correctif A du préambule
%   lualatex -jobname=crochet  "\def\SANSCROCHET{1}\input{minimal-juge3}"                sans correctif, crochet retiré
%   lualatex -jobname=amsseul  "\def\AMSSEUL{1}\input{minimal-juge3}"                    amsthm seul (témoin)
%   lualatex -jobname=faux     "\def\CORR{1}\def\FAUX{1}\input{minimal-juge3}"           deux vrais défauts : doivent tomber
%   lualatex -jobname=trace    "\def\TRACE{1}\input{minimal-juge3}"                      trace de l'appel
% Les lignes du correctif sont tirées du préambule commité par extraire_correctifs.py, pas retapées.
\documentclass{scrbook}
\ifdefined\SANSCROCHET \RemoveFromHook{package/amsthm/after}[firstaid/aliascounter]\fi
\usepackage{amsthm}
\usepackage{hyperref}
\usepackage[french, nameinlink, noabbrev]{cleveref}
\ifdefined\AMSSEUL
  \newtheorem{theoreme}{Théorème}[chapter]
  \newtheorem{proposition}[theoreme]{Proposition}
  \newtheorem{corollaire}[theoreme]{Corollaire}
  \newtheorem{remarque}{Remarque}[chapter]
\else
  \usepackage{thmtools}
  \ifdefined\CORR \input{correctif-a-pose}\fi
  \declaretheorem[numberwithin=chapter, name=Théorème]{theoreme}
  \ifdefined\TRACE \tracingmacros=2 \tracingonline=0 \fi
  \declaretheorem[sibling=theoreme, name=Proposition]{proposition}
  \ifdefined\TRACE \tracingmacros=0 \fi
  \declaretheorem[sibling=theoreme, name=Corollaire]{corollaire}
  \declaretheorem[numberwithin=chapter, name=Remarque]{remarque}
  \ifdefined\FAUX \newcounteralias{corollaire}{section}\newcounteralias{orphelin}{inexistant}\fi
  \ifdefined\CORR \input{correctif-a-remise}\fi
\fi
\crefname{theoreme}{théorème}{théorèmes}
\crefname{proposition}{proposition}{propositions}
\crefname{corollaire}{corollaire}{corollaires}
\crefname{remarque}{remarque}{remarques}
\begin{document}
\chapter{Premier}
\begin{proposition}\label{p1}a\end{proposition}
\begin{theoreme}\label{t1}b\end{theoreme}
\begin{remarque}\label{r1}c\end{remarque}
\begin{corollaire}\label{c1}d\end{corollaire}
\chapter{Second}
\begin{corollaire}\label{c2}e\end{corollaire}
\begin{theoreme}\label{t2}f\end{theoreme}
RENVOIS \ref{p1} \ref{t1} \ref{r1} \ref{c1} \ref{c2} \ref{t2} ; \cref{p1} ; \cref{t1} ; \cref{r1} ; \cref{c1} ; \cref{c2} ; \cref{t2} FIN
\end{document}
"                sans correctif, crochet retiré
%   lualatex -jobname=amsseul  "\def\AMSSEUL{1}% Juge de la troisième passe : document minimal écrit de zéro pour la cause A (thmtools sous LaTeX 2026-06-01).
% Variantes, par une macro posée sur la ligne de commande (deux passes chacune, depuis memoire/essai/juge3/minimal/) :
%   lualatex -jobname=nu       "\input{minimal-juge3}"                                   sans correctif
%   lualatex -jobname=corr     "\def\CORR{1}\input{minimal-juge3}"                       correctif A du préambule
%   lualatex -jobname=crochet  "\def\SANSCROCHET{1}\input{minimal-juge3}"                sans correctif, crochet retiré
%   lualatex -jobname=amsseul  "\def\AMSSEUL{1}\input{minimal-juge3}"                    amsthm seul (témoin)
%   lualatex -jobname=faux     "\def\CORR{1}\def\FAUX{1}\input{minimal-juge3}"           deux vrais défauts : doivent tomber
%   lualatex -jobname=trace    "\def\TRACE{1}\input{minimal-juge3}"                      trace de l'appel
% Les lignes du correctif sont tirées du préambule commité par extraire_correctifs.py, pas retapées.
\documentclass{scrbook}
\ifdefined\SANSCROCHET \RemoveFromHook{package/amsthm/after}[firstaid/aliascounter]\fi
\usepackage{amsthm}
\usepackage{hyperref}
\usepackage[french, nameinlink, noabbrev]{cleveref}
\ifdefined\AMSSEUL
  \newtheorem{theoreme}{Théorème}[chapter]
  \newtheorem{proposition}[theoreme]{Proposition}
  \newtheorem{corollaire}[theoreme]{Corollaire}
  \newtheorem{remarque}{Remarque}[chapter]
\else
  \usepackage{thmtools}
  \ifdefined\CORR \input{correctif-a-pose}\fi
  \declaretheorem[numberwithin=chapter, name=Théorème]{theoreme}
  \ifdefined\TRACE \tracingmacros=2 \tracingonline=0 \fi
  \declaretheorem[sibling=theoreme, name=Proposition]{proposition}
  \ifdefined\TRACE \tracingmacros=0 \fi
  \declaretheorem[sibling=theoreme, name=Corollaire]{corollaire}
  \declaretheorem[numberwithin=chapter, name=Remarque]{remarque}
  \ifdefined\FAUX \newcounteralias{corollaire}{section}\newcounteralias{orphelin}{inexistant}\fi
  \ifdefined\CORR \input{correctif-a-remise}\fi
\fi
\crefname{theoreme}{théorème}{théorèmes}
\crefname{proposition}{proposition}{propositions}
\crefname{corollaire}{corollaire}{corollaires}
\crefname{remarque}{remarque}{remarques}
\begin{document}
\chapter{Premier}
\begin{proposition}\label{p1}a\end{proposition}
\begin{theoreme}\label{t1}b\end{theoreme}
\begin{remarque}\label{r1}c\end{remarque}
\begin{corollaire}\label{c1}d\end{corollaire}
\chapter{Second}
\begin{corollaire}\label{c2}e\end{corollaire}
\begin{theoreme}\label{t2}f\end{theoreme}
RENVOIS \ref{p1} \ref{t1} \ref{r1} \ref{c1} \ref{c2} \ref{t2} ; \cref{p1} ; \cref{t1} ; \cref{r1} ; \cref{c1} ; \cref{c2} ; \cref{t2} FIN
\end{document}
"                    amsthm seul (témoin)
%   lualatex -jobname=faux     "\def\CORR{1}\def\FAUX{1}% Juge de la troisième passe : document minimal écrit de zéro pour la cause A (thmtools sous LaTeX 2026-06-01).
% Variantes, par une macro posée sur la ligne de commande (deux passes chacune, depuis memoire/essai/juge3/minimal/) :
%   lualatex -jobname=nu       "\input{minimal-juge3}"                                   sans correctif
%   lualatex -jobname=corr     "\def\CORR{1}\input{minimal-juge3}"                       correctif A du préambule
%   lualatex -jobname=crochet  "\def\SANSCROCHET{1}\input{minimal-juge3}"                sans correctif, crochet retiré
%   lualatex -jobname=amsseul  "\def\AMSSEUL{1}\input{minimal-juge3}"                    amsthm seul (témoin)
%   lualatex -jobname=faux     "\def\CORR{1}\def\FAUX{1}\input{minimal-juge3}"           deux vrais défauts : doivent tomber
%   lualatex -jobname=trace    "\def\TRACE{1}\input{minimal-juge3}"                      trace de l'appel
% Les lignes du correctif sont tirées du préambule commité par extraire_correctifs.py, pas retapées.
\documentclass{scrbook}
\ifdefined\SANSCROCHET \RemoveFromHook{package/amsthm/after}[firstaid/aliascounter]\fi
\usepackage{amsthm}
\usepackage{hyperref}
\usepackage[french, nameinlink, noabbrev]{cleveref}
\ifdefined\AMSSEUL
  \newtheorem{theoreme}{Théorème}[chapter]
  \newtheorem{proposition}[theoreme]{Proposition}
  \newtheorem{corollaire}[theoreme]{Corollaire}
  \newtheorem{remarque}{Remarque}[chapter]
\else
  \usepackage{thmtools}
  \ifdefined\CORR \input{correctif-a-pose}\fi
  \declaretheorem[numberwithin=chapter, name=Théorème]{theoreme}
  \ifdefined\TRACE \tracingmacros=2 \tracingonline=0 \fi
  \declaretheorem[sibling=theoreme, name=Proposition]{proposition}
  \ifdefined\TRACE \tracingmacros=0 \fi
  \declaretheorem[sibling=theoreme, name=Corollaire]{corollaire}
  \declaretheorem[numberwithin=chapter, name=Remarque]{remarque}
  \ifdefined\FAUX \newcounteralias{corollaire}{section}\newcounteralias{orphelin}{inexistant}\fi
  \ifdefined\CORR \input{correctif-a-remise}\fi
\fi
\crefname{theoreme}{théorème}{théorèmes}
\crefname{proposition}{proposition}{propositions}
\crefname{corollaire}{corollaire}{corollaires}
\crefname{remarque}{remarque}{remarques}
\begin{document}
\chapter{Premier}
\begin{proposition}\label{p1}a\end{proposition}
\begin{theoreme}\label{t1}b\end{theoreme}
\begin{remarque}\label{r1}c\end{remarque}
\begin{corollaire}\label{c1}d\end{corollaire}
\chapter{Second}
\begin{corollaire}\label{c2}e\end{corollaire}
\begin{theoreme}\label{t2}f\end{theoreme}
RENVOIS \ref{p1} \ref{t1} \ref{r1} \ref{c1} \ref{c2} \ref{t2} ; \cref{p1} ; \cref{t1} ; \cref{r1} ; \cref{c1} ; \cref{c2} ; \cref{t2} FIN
\end{document}
"           deux vrais défauts : doivent tomber
%   lualatex -jobname=trace    "\def\TRACE{1}% Juge de la troisième passe : document minimal écrit de zéro pour la cause A (thmtools sous LaTeX 2026-06-01).
% Variantes, par une macro posée sur la ligne de commande (deux passes chacune, depuis memoire/essai/juge3/minimal/) :
%   lualatex -jobname=nu       "\input{minimal-juge3}"                                   sans correctif
%   lualatex -jobname=corr     "\def\CORR{1}\input{minimal-juge3}"                       correctif A du préambule
%   lualatex -jobname=crochet  "\def\SANSCROCHET{1}\input{minimal-juge3}"                sans correctif, crochet retiré
%   lualatex -jobname=amsseul  "\def\AMSSEUL{1}\input{minimal-juge3}"                    amsthm seul (témoin)
%   lualatex -jobname=faux     "\def\CORR{1}\def\FAUX{1}\input{minimal-juge3}"           deux vrais défauts : doivent tomber
%   lualatex -jobname=trace    "\def\TRACE{1}\input{minimal-juge3}"                      trace de l'appel
% Les lignes du correctif sont tirées du préambule commité par extraire_correctifs.py, pas retapées.
\documentclass{scrbook}
\ifdefined\SANSCROCHET \RemoveFromHook{package/amsthm/after}[firstaid/aliascounter]\fi
\usepackage{amsthm}
\usepackage{hyperref}
\usepackage[french, nameinlink, noabbrev]{cleveref}
\ifdefined\AMSSEUL
  \newtheorem{theoreme}{Théorème}[chapter]
  \newtheorem{proposition}[theoreme]{Proposition}
  \newtheorem{corollaire}[theoreme]{Corollaire}
  \newtheorem{remarque}{Remarque}[chapter]
\else
  \usepackage{thmtools}
  \ifdefined\CORR \input{correctif-a-pose}\fi
  \declaretheorem[numberwithin=chapter, name=Théorème]{theoreme}
  \ifdefined\TRACE \tracingmacros=2 \tracingonline=0 \fi
  \declaretheorem[sibling=theoreme, name=Proposition]{proposition}
  \ifdefined\TRACE \tracingmacros=0 \fi
  \declaretheorem[sibling=theoreme, name=Corollaire]{corollaire}
  \declaretheorem[numberwithin=chapter, name=Remarque]{remarque}
  \ifdefined\FAUX \newcounteralias{corollaire}{section}\newcounteralias{orphelin}{inexistant}\fi
  \ifdefined\CORR \input{correctif-a-remise}\fi
\fi
\crefname{theoreme}{théorème}{théorèmes}
\crefname{proposition}{proposition}{propositions}
\crefname{corollaire}{corollaire}{corollaires}
\crefname{remarque}{remarque}{remarques}
\begin{document}
\chapter{Premier}
\begin{proposition}\label{p1}a\end{proposition}
\begin{theoreme}\label{t1}b\end{theoreme}
\begin{remarque}\label{r1}c\end{remarque}
\begin{corollaire}\label{c1}d\end{corollaire}
\chapter{Second}
\begin{corollaire}\label{c2}e\end{corollaire}
\begin{theoreme}\label{t2}f\end{theoreme}
RENVOIS \ref{p1} \ref{t1} \ref{r1} \ref{c1} \ref{c2} \ref{t2} ; \cref{p1} ; \cref{t1} ; \cref{r1} ; \cref{c1} ; \cref{c2} ; \cref{t2} FIN
\end{document}
"                      trace de l'appel
% Les lignes du correctif sont tirées du préambule commité par extraire_correctifs.py, pas retapées.
\documentclass{scrbook}
\ifdefined\SANSCROCHET \RemoveFromHook{package/amsthm/after}[firstaid/aliascounter]\fi
\usepackage{amsthm}
\usepackage{hyperref}
\usepackage[french, nameinlink, noabbrev]{cleveref}
\ifdefined\AMSSEUL
  \newtheorem{theoreme}{Théorème}[chapter]
  \newtheorem{proposition}[theoreme]{Proposition}
  \newtheorem{corollaire}[theoreme]{Corollaire}
  \newtheorem{remarque}{Remarque}[chapter]
\else
  \usepackage{thmtools}
  \ifdefined\CORR \input{correctif-a-pose}\fi
  \declaretheorem[numberwithin=chapter, name=Théorème]{theoreme}
  \ifdefined\TRACE \tracingmacros=2 \tracingonline=0 \fi
  \declaretheorem[sibling=theoreme, name=Proposition]{proposition}
  \ifdefined\TRACE \tracingmacros=0 \fi
  \declaretheorem[sibling=theoreme, name=Corollaire]{corollaire}
  \declaretheorem[numberwithin=chapter, name=Remarque]{remarque}
  \ifdefined\FAUX \newcounteralias{corollaire}{section}\newcounteralias{orphelin}{inexistant}\fi
  \ifdefined\CORR \input{correctif-a-remise}\fi
\fi
\crefname{theoreme}{théorème}{théorèmes}
\crefname{proposition}{proposition}{propositions}
\crefname{corollaire}{corollaire}{corollaires}
\crefname{remarque}{remarque}{remarques}
\begin{document}
\chapter{Premier}
\begin{proposition}\label{p1}a\end{proposition}
\begin{theoreme}\label{t1}b\end{theoreme}
\begin{remarque}\label{r1}c\end{remarque}
\begin{corollaire}\label{c1}d\end{corollaire}
\chapter{Second}
\begin{corollaire}\label{c2}e\end{corollaire}
\begin{theoreme}\label{t2}f\end{theoreme}
RENVOIS \ref{p1} \ref{t1} \ref{r1} \ref{c1} \ref{c2} \ref{t2} ; \cref{p1} ; \cref{t1} ; \cref{r1} ; \cref{c1} ; \cref{c2} ; \cref{t2} FIN
\end{document}
"                      trace de l'appel
% Les lignes du correctif sont tirées du préambule commité par extraire_correctifs.py, pas retapées.
\documentclass{scrbook}
\ifdefined\SANSCROCHET \RemoveFromHook{package/amsthm/after}[firstaid/aliascounter]\fi
\usepackage{amsthm}
\usepackage{hyperref}
\usepackage[french, nameinlink, noabbrev]{cleveref}
\ifdefined\AMSSEUL
  \newtheorem{theoreme}{Théorème}[chapter]
  \newtheorem{proposition}[theoreme]{Proposition}
  \newtheorem{corollaire}[theoreme]{Corollaire}
  \newtheorem{remarque}{Remarque}[chapter]
\else
  \usepackage{thmtools}
  \ifdefined\CORR \input{correctif-a-pose}\fi
  \declaretheorem[numberwithin=chapter, name=Théorème]{theoreme}
  \ifdefined\TRACE \tracingmacros=2 \tracingonline=0 \fi
  \declaretheorem[sibling=theoreme, name=Proposition]{proposition}
  \ifdefined\TRACE \tracingmacros=0 \fi
  \declaretheorem[sibling=theoreme, name=Corollaire]{corollaire}
  \declaretheorem[numberwithin=chapter, name=Remarque]{remarque}
  \ifdefined\FAUX \newcounteralias{corollaire}{section}\newcounteralias{orphelin}{inexistant}\fi
  \ifdefined\CORR \input{correctif-a-remise}\fi
\fi
\crefname{theoreme}{théorème}{théorèmes}
\crefname{proposition}{proposition}{propositions}
\crefname{corollaire}{corollaire}{corollaires}
\crefname{remarque}{remarque}{remarques}
\begin{document}
\chapter{Premier}
\begin{proposition}\label{p1}a\end{proposition}
\begin{theoreme}\label{t1}b\end{theoreme}
\begin{remarque}\label{r1}c\end{remarque}
\begin{corollaire}\label{c1}d\end{corollaire}
\chapter{Second}
\begin{corollaire}\label{c2}e\end{corollaire}
\begin{theoreme}\label{t2}f\end{theoreme}
RENVOIS \ref{p1} \ref{t1} \ref{r1} \ref{c1} \ref{c2} \ref{t2} ; \cref{p1} ; \cref{t1} ; \cref{r1} ; \cref{c1} ; \cref{c2} ; \cref{t2} FIN
\end{document}
"                sans correctif, crochet retiré
%   lualatex -jobname=amsseul  "\def\AMSSEUL{1}% Juge de la troisième passe : document minimal écrit de zéro pour la cause A (thmtools sous LaTeX 2026-06-01).
% Variantes, par une macro posée sur la ligne de commande (deux passes chacune, depuis memoire/essai/juge3/minimal/) :
%   lualatex -jobname=nu       "% Juge de la troisième passe : document minimal écrit de zéro pour la cause A (thmtools sous LaTeX 2026-06-01).
% Variantes, par une macro posée sur la ligne de commande (deux passes chacune, depuis memoire/essai/juge3/minimal/) :
%   lualatex -jobname=nu       "% Juge de la troisième passe : document minimal écrit de zéro pour la cause A (thmtools sous LaTeX 2026-06-01).
% Variantes, par une macro posée sur la ligne de commande (deux passes chacune, depuis memoire/essai/juge3/minimal/) :
%   lualatex -jobname=nu       "\input{minimal-juge3}"                                   sans correctif
%   lualatex -jobname=corr     "\def\CORR{1}\input{minimal-juge3}"                       correctif A du préambule
%   lualatex -jobname=crochet  "\def\SANSCROCHET{1}\input{minimal-juge3}"                sans correctif, crochet retiré
%   lualatex -jobname=amsseul  "\def\AMSSEUL{1}\input{minimal-juge3}"                    amsthm seul (témoin)
%   lualatex -jobname=faux     "\def\CORR{1}\def\FAUX{1}\input{minimal-juge3}"           deux vrais défauts : doivent tomber
%   lualatex -jobname=trace    "\def\TRACE{1}\input{minimal-juge3}"                      trace de l'appel
% Les lignes du correctif sont tirées du préambule commité par extraire_correctifs.py, pas retapées.
\documentclass{scrbook}
\ifdefined\SANSCROCHET \RemoveFromHook{package/amsthm/after}[firstaid/aliascounter]\fi
\usepackage{amsthm}
\usepackage{hyperref}
\usepackage[french, nameinlink, noabbrev]{cleveref}
\ifdefined\AMSSEUL
  \newtheorem{theoreme}{Théorème}[chapter]
  \newtheorem{proposition}[theoreme]{Proposition}
  \newtheorem{corollaire}[theoreme]{Corollaire}
  \newtheorem{remarque}{Remarque}[chapter]
\else
  \usepackage{thmtools}
  \ifdefined\CORR \input{correctif-a-pose}\fi
  \declaretheorem[numberwithin=chapter, name=Théorème]{theoreme}
  \ifdefined\TRACE \tracingmacros=2 \tracingonline=0 \fi
  \declaretheorem[sibling=theoreme, name=Proposition]{proposition}
  \ifdefined\TRACE \tracingmacros=0 \fi
  \declaretheorem[sibling=theoreme, name=Corollaire]{corollaire}
  \declaretheorem[numberwithin=chapter, name=Remarque]{remarque}
  \ifdefined\FAUX \newcounteralias{corollaire}{section}\newcounteralias{orphelin}{inexistant}\fi
  \ifdefined\CORR \input{correctif-a-remise}\fi
\fi
\crefname{theoreme}{théorème}{théorèmes}
\crefname{proposition}{proposition}{propositions}
\crefname{corollaire}{corollaire}{corollaires}
\crefname{remarque}{remarque}{remarques}
\begin{document}
\chapter{Premier}
\begin{proposition}\label{p1}a\end{proposition}
\begin{theoreme}\label{t1}b\end{theoreme}
\begin{remarque}\label{r1}c\end{remarque}
\begin{corollaire}\label{c1}d\end{corollaire}
\chapter{Second}
\begin{corollaire}\label{c2}e\end{corollaire}
\begin{theoreme}\label{t2}f\end{theoreme}
RENVOIS \ref{p1} \ref{t1} \ref{r1} \ref{c1} \ref{c2} \ref{t2} ; \cref{p1} ; \cref{t1} ; \cref{r1} ; \cref{c1} ; \cref{c2} ; \cref{t2} FIN
\end{document}
"                                   sans correctif
%   lualatex -jobname=corr     "\def\CORR{1}% Juge de la troisième passe : document minimal écrit de zéro pour la cause A (thmtools sous LaTeX 2026-06-01).
% Variantes, par une macro posée sur la ligne de commande (deux passes chacune, depuis memoire/essai/juge3/minimal/) :
%   lualatex -jobname=nu       "\input{minimal-juge3}"                                   sans correctif
%   lualatex -jobname=corr     "\def\CORR{1}\input{minimal-juge3}"                       correctif A du préambule
%   lualatex -jobname=crochet  "\def\SANSCROCHET{1}\input{minimal-juge3}"                sans correctif, crochet retiré
%   lualatex -jobname=amsseul  "\def\AMSSEUL{1}\input{minimal-juge3}"                    amsthm seul (témoin)
%   lualatex -jobname=faux     "\def\CORR{1}\def\FAUX{1}\input{minimal-juge3}"           deux vrais défauts : doivent tomber
%   lualatex -jobname=trace    "\def\TRACE{1}\input{minimal-juge3}"                      trace de l'appel
% Les lignes du correctif sont tirées du préambule commité par extraire_correctifs.py, pas retapées.
\documentclass{scrbook}
\ifdefined\SANSCROCHET \RemoveFromHook{package/amsthm/after}[firstaid/aliascounter]\fi
\usepackage{amsthm}
\usepackage{hyperref}
\usepackage[french, nameinlink, noabbrev]{cleveref}
\ifdefined\AMSSEUL
  \newtheorem{theoreme}{Théorème}[chapter]
  \newtheorem{proposition}[theoreme]{Proposition}
  \newtheorem{corollaire}[theoreme]{Corollaire}
  \newtheorem{remarque}{Remarque}[chapter]
\else
  \usepackage{thmtools}
  \ifdefined\CORR \input{correctif-a-pose}\fi
  \declaretheorem[numberwithin=chapter, name=Théorème]{theoreme}
  \ifdefined\TRACE \tracingmacros=2 \tracingonline=0 \fi
  \declaretheorem[sibling=theoreme, name=Proposition]{proposition}
  \ifdefined\TRACE \tracingmacros=0 \fi
  \declaretheorem[sibling=theoreme, name=Corollaire]{corollaire}
  \declaretheorem[numberwithin=chapter, name=Remarque]{remarque}
  \ifdefined\FAUX \newcounteralias{corollaire}{section}\newcounteralias{orphelin}{inexistant}\fi
  \ifdefined\CORR \input{correctif-a-remise}\fi
\fi
\crefname{theoreme}{théorème}{théorèmes}
\crefname{proposition}{proposition}{propositions}
\crefname{corollaire}{corollaire}{corollaires}
\crefname{remarque}{remarque}{remarques}
\begin{document}
\chapter{Premier}
\begin{proposition}\label{p1}a\end{proposition}
\begin{theoreme}\label{t1}b\end{theoreme}
\begin{remarque}\label{r1}c\end{remarque}
\begin{corollaire}\label{c1}d\end{corollaire}
\chapter{Second}
\begin{corollaire}\label{c2}e\end{corollaire}
\begin{theoreme}\label{t2}f\end{theoreme}
RENVOIS \ref{p1} \ref{t1} \ref{r1} \ref{c1} \ref{c2} \ref{t2} ; \cref{p1} ; \cref{t1} ; \cref{r1} ; \cref{c1} ; \cref{c2} ; \cref{t2} FIN
\end{document}
"                       correctif A du préambule
%   lualatex -jobname=crochet  "\def\SANSCROCHET{1}% Juge de la troisième passe : document minimal écrit de zéro pour la cause A (thmtools sous LaTeX 2026-06-01).
% Variantes, par une macro posée sur la ligne de commande (deux passes chacune, depuis memoire/essai/juge3/minimal/) :
%   lualatex -jobname=nu       "\input{minimal-juge3}"                                   sans correctif
%   lualatex -jobname=corr     "\def\CORR{1}\input{minimal-juge3}"                       correctif A du préambule
%   lualatex -jobname=crochet  "\def\SANSCROCHET{1}\input{minimal-juge3}"                sans correctif, crochet retiré
%   lualatex -jobname=amsseul  "\def\AMSSEUL{1}\input{minimal-juge3}"                    amsthm seul (témoin)
%   lualatex -jobname=faux     "\def\CORR{1}\def\FAUX{1}\input{minimal-juge3}"           deux vrais défauts : doivent tomber
%   lualatex -jobname=trace    "\def\TRACE{1}\input{minimal-juge3}"                      trace de l'appel
% Les lignes du correctif sont tirées du préambule commité par extraire_correctifs.py, pas retapées.
\documentclass{scrbook}
\ifdefined\SANSCROCHET \RemoveFromHook{package/amsthm/after}[firstaid/aliascounter]\fi
\usepackage{amsthm}
\usepackage{hyperref}
\usepackage[french, nameinlink, noabbrev]{cleveref}
\ifdefined\AMSSEUL
  \newtheorem{theoreme}{Théorème}[chapter]
  \newtheorem{proposition}[theoreme]{Proposition}
  \newtheorem{corollaire}[theoreme]{Corollaire}
  \newtheorem{remarque}{Remarque}[chapter]
\else
  \usepackage{thmtools}
  \ifdefined\CORR \input{correctif-a-pose}\fi
  \declaretheorem[numberwithin=chapter, name=Théorème]{theoreme}
  \ifdefined\TRACE \tracingmacros=2 \tracingonline=0 \fi
  \declaretheorem[sibling=theoreme, name=Proposition]{proposition}
  \ifdefined\TRACE \tracingmacros=0 \fi
  \declaretheorem[sibling=theoreme, name=Corollaire]{corollaire}
  \declaretheorem[numberwithin=chapter, name=Remarque]{remarque}
  \ifdefined\FAUX \newcounteralias{corollaire}{section}\newcounteralias{orphelin}{inexistant}\fi
  \ifdefined\CORR \input{correctif-a-remise}\fi
\fi
\crefname{theoreme}{théorème}{théorèmes}
\crefname{proposition}{proposition}{propositions}
\crefname{corollaire}{corollaire}{corollaires}
\crefname{remarque}{remarque}{remarques}
\begin{document}
\chapter{Premier}
\begin{proposition}\label{p1}a\end{proposition}
\begin{theoreme}\label{t1}b\end{theoreme}
\begin{remarque}\label{r1}c\end{remarque}
\begin{corollaire}\label{c1}d\end{corollaire}
\chapter{Second}
\begin{corollaire}\label{c2}e\end{corollaire}
\begin{theoreme}\label{t2}f\end{theoreme}
RENVOIS \ref{p1} \ref{t1} \ref{r1} \ref{c1} \ref{c2} \ref{t2} ; \cref{p1} ; \cref{t1} ; \cref{r1} ; \cref{c1} ; \cref{c2} ; \cref{t2} FIN
\end{document}
"                sans correctif, crochet retiré
%   lualatex -jobname=amsseul  "\def\AMSSEUL{1}% Juge de la troisième passe : document minimal écrit de zéro pour la cause A (thmtools sous LaTeX 2026-06-01).
% Variantes, par une macro posée sur la ligne de commande (deux passes chacune, depuis memoire/essai/juge3/minimal/) :
%   lualatex -jobname=nu       "\input{minimal-juge3}"                                   sans correctif
%   lualatex -jobname=corr     "\def\CORR{1}\input{minimal-juge3}"                       correctif A du préambule
%   lualatex -jobname=crochet  "\def\SANSCROCHET{1}\input{minimal-juge3}"                sans correctif, crochet retiré
%   lualatex -jobname=amsseul  "\def\AMSSEUL{1}\input{minimal-juge3}"                    amsthm seul (témoin)
%   lualatex -jobname=faux     "\def\CORR{1}\def\FAUX{1}\input{minimal-juge3}"           deux vrais défauts : doivent tomber
%   lualatex -jobname=trace    "\def\TRACE{1}\input{minimal-juge3}"                      trace de l'appel
% Les lignes du correctif sont tirées du préambule commité par extraire_correctifs.py, pas retapées.
\documentclass{scrbook}
\ifdefined\SANSCROCHET \RemoveFromHook{package/amsthm/after}[firstaid/aliascounter]\fi
\usepackage{amsthm}
\usepackage{hyperref}
\usepackage[french, nameinlink, noabbrev]{cleveref}
\ifdefined\AMSSEUL
  \newtheorem{theoreme}{Théorème}[chapter]
  \newtheorem{proposition}[theoreme]{Proposition}
  \newtheorem{corollaire}[theoreme]{Corollaire}
  \newtheorem{remarque}{Remarque}[chapter]
\else
  \usepackage{thmtools}
  \ifdefined\CORR \input{correctif-a-pose}\fi
  \declaretheorem[numberwithin=chapter, name=Théorème]{theoreme}
  \ifdefined\TRACE \tracingmacros=2 \tracingonline=0 \fi
  \declaretheorem[sibling=theoreme, name=Proposition]{proposition}
  \ifdefined\TRACE \tracingmacros=0 \fi
  \declaretheorem[sibling=theoreme, name=Corollaire]{corollaire}
  \declaretheorem[numberwithin=chapter, name=Remarque]{remarque}
  \ifdefined\FAUX \newcounteralias{corollaire}{section}\newcounteralias{orphelin}{inexistant}\fi
  \ifdefined\CORR \input{correctif-a-remise}\fi
\fi
\crefname{theoreme}{théorème}{théorèmes}
\crefname{proposition}{proposition}{propositions}
\crefname{corollaire}{corollaire}{corollaires}
\crefname{remarque}{remarque}{remarques}
\begin{document}
\chapter{Premier}
\begin{proposition}\label{p1}a\end{proposition}
\begin{theoreme}\label{t1}b\end{theoreme}
\begin{remarque}\label{r1}c\end{remarque}
\begin{corollaire}\label{c1}d\end{corollaire}
\chapter{Second}
\begin{corollaire}\label{c2}e\end{corollaire}
\begin{theoreme}\label{t2}f\end{theoreme}
RENVOIS \ref{p1} \ref{t1} \ref{r1} \ref{c1} \ref{c2} \ref{t2} ; \cref{p1} ; \cref{t1} ; \cref{r1} ; \cref{c1} ; \cref{c2} ; \cref{t2} FIN
\end{document}
"                    amsthm seul (témoin)
%   lualatex -jobname=faux     "\def\CORR{1}\def\FAUX{1}% Juge de la troisième passe : document minimal écrit de zéro pour la cause A (thmtools sous LaTeX 2026-06-01).
% Variantes, par une macro posée sur la ligne de commande (deux passes chacune, depuis memoire/essai/juge3/minimal/) :
%   lualatex -jobname=nu       "\input{minimal-juge3}"                                   sans correctif
%   lualatex -jobname=corr     "\def\CORR{1}\input{minimal-juge3}"                       correctif A du préambule
%   lualatex -jobname=crochet  "\def\SANSCROCHET{1}\input{minimal-juge3}"                sans correctif, crochet retiré
%   lualatex -jobname=amsseul  "\def\AMSSEUL{1}\input{minimal-juge3}"                    amsthm seul (témoin)
%   lualatex -jobname=faux     "\def\CORR{1}\def\FAUX{1}\input{minimal-juge3}"           deux vrais défauts : doivent tomber
%   lualatex -jobname=trace    "\def\TRACE{1}\input{minimal-juge3}"                      trace de l'appel
% Les lignes du correctif sont tirées du préambule commité par extraire_correctifs.py, pas retapées.
\documentclass{scrbook}
\ifdefined\SANSCROCHET \RemoveFromHook{package/amsthm/after}[firstaid/aliascounter]\fi
\usepackage{amsthm}
\usepackage{hyperref}
\usepackage[french, nameinlink, noabbrev]{cleveref}
\ifdefined\AMSSEUL
  \newtheorem{theoreme}{Théorème}[chapter]
  \newtheorem{proposition}[theoreme]{Proposition}
  \newtheorem{corollaire}[theoreme]{Corollaire}
  \newtheorem{remarque}{Remarque}[chapter]
\else
  \usepackage{thmtools}
  \ifdefined\CORR \input{correctif-a-pose}\fi
  \declaretheorem[numberwithin=chapter, name=Théorème]{theoreme}
  \ifdefined\TRACE \tracingmacros=2 \tracingonline=0 \fi
  \declaretheorem[sibling=theoreme, name=Proposition]{proposition}
  \ifdefined\TRACE \tracingmacros=0 \fi
  \declaretheorem[sibling=theoreme, name=Corollaire]{corollaire}
  \declaretheorem[numberwithin=chapter, name=Remarque]{remarque}
  \ifdefined\FAUX \newcounteralias{corollaire}{section}\newcounteralias{orphelin}{inexistant}\fi
  \ifdefined\CORR \input{correctif-a-remise}\fi
\fi
\crefname{theoreme}{théorème}{théorèmes}
\crefname{proposition}{proposition}{propositions}
\crefname{corollaire}{corollaire}{corollaires}
\crefname{remarque}{remarque}{remarques}
\begin{document}
\chapter{Premier}
\begin{proposition}\label{p1}a\end{proposition}
\begin{theoreme}\label{t1}b\end{theoreme}
\begin{remarque}\label{r1}c\end{remarque}
\begin{corollaire}\label{c1}d\end{corollaire}
\chapter{Second}
\begin{corollaire}\label{c2}e\end{corollaire}
\begin{theoreme}\label{t2}f\end{theoreme}
RENVOIS \ref{p1} \ref{t1} \ref{r1} \ref{c1} \ref{c2} \ref{t2} ; \cref{p1} ; \cref{t1} ; \cref{r1} ; \cref{c1} ; \cref{c2} ; \cref{t2} FIN
\end{document}
"           deux vrais défauts : doivent tomber
%   lualatex -jobname=trace    "\def\TRACE{1}% Juge de la troisième passe : document minimal écrit de zéro pour la cause A (thmtools sous LaTeX 2026-06-01).
% Variantes, par une macro posée sur la ligne de commande (deux passes chacune, depuis memoire/essai/juge3/minimal/) :
%   lualatex -jobname=nu       "\input{minimal-juge3}"                                   sans correctif
%   lualatex -jobname=corr     "\def\CORR{1}\input{minimal-juge3}"                       correctif A du préambule
%   lualatex -jobname=crochet  "\def\SANSCROCHET{1}\input{minimal-juge3}"                sans correctif, crochet retiré
%   lualatex -jobname=amsseul  "\def\AMSSEUL{1}\input{minimal-juge3}"                    amsthm seul (témoin)
%   lualatex -jobname=faux     "\def\CORR{1}\def\FAUX{1}\input{minimal-juge3}"           deux vrais défauts : doivent tomber
%   lualatex -jobname=trace    "\def\TRACE{1}\input{minimal-juge3}"                      trace de l'appel
% Les lignes du correctif sont tirées du préambule commité par extraire_correctifs.py, pas retapées.
\documentclass{scrbook}
\ifdefined\SANSCROCHET \RemoveFromHook{package/amsthm/after}[firstaid/aliascounter]\fi
\usepackage{amsthm}
\usepackage{hyperref}
\usepackage[french, nameinlink, noabbrev]{cleveref}
\ifdefined\AMSSEUL
  \newtheorem{theoreme}{Théorème}[chapter]
  \newtheorem{proposition}[theoreme]{Proposition}
  \newtheorem{corollaire}[theoreme]{Corollaire}
  \newtheorem{remarque}{Remarque}[chapter]
\else
  \usepackage{thmtools}
  \ifdefined\CORR \input{correctif-a-pose}\fi
  \declaretheorem[numberwithin=chapter, name=Théorème]{theoreme}
  \ifdefined\TRACE \tracingmacros=2 \tracingonline=0 \fi
  \declaretheorem[sibling=theoreme, name=Proposition]{proposition}
  \ifdefined\TRACE \tracingmacros=0 \fi
  \declaretheorem[sibling=theoreme, name=Corollaire]{corollaire}
  \declaretheorem[numberwithin=chapter, name=Remarque]{remarque}
  \ifdefined\FAUX \newcounteralias{corollaire}{section}\newcounteralias{orphelin}{inexistant}\fi
  \ifdefined\CORR \input{correctif-a-remise}\fi
\fi
\crefname{theoreme}{théorème}{théorèmes}
\crefname{proposition}{proposition}{propositions}
\crefname{corollaire}{corollaire}{corollaires}
\crefname{remarque}{remarque}{remarques}
\begin{document}
\chapter{Premier}
\begin{proposition}\label{p1}a\end{proposition}
\begin{theoreme}\label{t1}b\end{theoreme}
\begin{remarque}\label{r1}c\end{remarque}
\begin{corollaire}\label{c1}d\end{corollaire}
\chapter{Second}
\begin{corollaire}\label{c2}e\end{corollaire}
\begin{theoreme}\label{t2}f\end{theoreme}
RENVOIS \ref{p1} \ref{t1} \ref{r1} \ref{c1} \ref{c2} \ref{t2} ; \cref{p1} ; \cref{t1} ; \cref{r1} ; \cref{c1} ; \cref{c2} ; \cref{t2} FIN
\end{document}
"                      trace de l'appel
% Les lignes du correctif sont tirées du préambule commité par extraire_correctifs.py, pas retapées.
\documentclass{scrbook}
\ifdefined\SANSCROCHET \RemoveFromHook{package/amsthm/after}[firstaid/aliascounter]\fi
\usepackage{amsthm}
\usepackage{hyperref}
\usepackage[french, nameinlink, noabbrev]{cleveref}
\ifdefined\AMSSEUL
  \newtheorem{theoreme}{Théorème}[chapter]
  \newtheorem{proposition}[theoreme]{Proposition}
  \newtheorem{corollaire}[theoreme]{Corollaire}
  \newtheorem{remarque}{Remarque}[chapter]
\else
  \usepackage{thmtools}
  \ifdefined\CORR \input{correctif-a-pose}\fi
  \declaretheorem[numberwithin=chapter, name=Théorème]{theoreme}
  \ifdefined\TRACE \tracingmacros=2 \tracingonline=0 \fi
  \declaretheorem[sibling=theoreme, name=Proposition]{proposition}
  \ifdefined\TRACE \tracingmacros=0 \fi
  \declaretheorem[sibling=theoreme, name=Corollaire]{corollaire}
  \declaretheorem[numberwithin=chapter, name=Remarque]{remarque}
  \ifdefined\FAUX \newcounteralias{corollaire}{section}\newcounteralias{orphelin}{inexistant}\fi
  \ifdefined\CORR \input{correctif-a-remise}\fi
\fi
\crefname{theoreme}{théorème}{théorèmes}
\crefname{proposition}{proposition}{propositions}
\crefname{corollaire}{corollaire}{corollaires}
\crefname{remarque}{remarque}{remarques}
\begin{document}
\chapter{Premier}
\begin{proposition}\label{p1}a\end{proposition}
\begin{theoreme}\label{t1}b\end{theoreme}
\begin{remarque}\label{r1}c\end{remarque}
\begin{corollaire}\label{c1}d\end{corollaire}
\chapter{Second}
\begin{corollaire}\label{c2}e\end{corollaire}
\begin{theoreme}\label{t2}f\end{theoreme}
RENVOIS \ref{p1} \ref{t1} \ref{r1} \ref{c1} \ref{c2} \ref{t2} ; \cref{p1} ; \cref{t1} ; \cref{r1} ; \cref{c1} ; \cref{c2} ; \cref{t2} FIN
\end{document}
"                                   sans correctif
%   lualatex -jobname=corr     "\def\CORR{1}% Juge de la troisième passe : document minimal écrit de zéro pour la cause A (thmtools sous LaTeX 2026-06-01).
% Variantes, par une macro posée sur la ligne de commande (deux passes chacune, depuis memoire/essai/juge3/minimal/) :
%   lualatex -jobname=nu       "% Juge de la troisième passe : document minimal écrit de zéro pour la cause A (thmtools sous LaTeX 2026-06-01).
% Variantes, par une macro posée sur la ligne de commande (deux passes chacune, depuis memoire/essai/juge3/minimal/) :
%   lualatex -jobname=nu       "\input{minimal-juge3}"                                   sans correctif
%   lualatex -jobname=corr     "\def\CORR{1}\input{minimal-juge3}"                       correctif A du préambule
%   lualatex -jobname=crochet  "\def\SANSCROCHET{1}\input{minimal-juge3}"                sans correctif, crochet retiré
%   lualatex -jobname=amsseul  "\def\AMSSEUL{1}\input{minimal-juge3}"                    amsthm seul (témoin)
%   lualatex -jobname=faux     "\def\CORR{1}\def\FAUX{1}\input{minimal-juge3}"           deux vrais défauts : doivent tomber
%   lualatex -jobname=trace    "\def\TRACE{1}\input{minimal-juge3}"                      trace de l'appel
% Les lignes du correctif sont tirées du préambule commité par extraire_correctifs.py, pas retapées.
\documentclass{scrbook}
\ifdefined\SANSCROCHET \RemoveFromHook{package/amsthm/after}[firstaid/aliascounter]\fi
\usepackage{amsthm}
\usepackage{hyperref}
\usepackage[french, nameinlink, noabbrev]{cleveref}
\ifdefined\AMSSEUL
  \newtheorem{theoreme}{Théorème}[chapter]
  \newtheorem{proposition}[theoreme]{Proposition}
  \newtheorem{corollaire}[theoreme]{Corollaire}
  \newtheorem{remarque}{Remarque}[chapter]
\else
  \usepackage{thmtools}
  \ifdefined\CORR \input{correctif-a-pose}\fi
  \declaretheorem[numberwithin=chapter, name=Théorème]{theoreme}
  \ifdefined\TRACE \tracingmacros=2 \tracingonline=0 \fi
  \declaretheorem[sibling=theoreme, name=Proposition]{proposition}
  \ifdefined\TRACE \tracingmacros=0 \fi
  \declaretheorem[sibling=theoreme, name=Corollaire]{corollaire}
  \declaretheorem[numberwithin=chapter, name=Remarque]{remarque}
  \ifdefined\FAUX \newcounteralias{corollaire}{section}\newcounteralias{orphelin}{inexistant}\fi
  \ifdefined\CORR \input{correctif-a-remise}\fi
\fi
\crefname{theoreme}{théorème}{théorèmes}
\crefname{proposition}{proposition}{propositions}
\crefname{corollaire}{corollaire}{corollaires}
\crefname{remarque}{remarque}{remarques}
\begin{document}
\chapter{Premier}
\begin{proposition}\label{p1}a\end{proposition}
\begin{theoreme}\label{t1}b\end{theoreme}
\begin{remarque}\label{r1}c\end{remarque}
\begin{corollaire}\label{c1}d\end{corollaire}
\chapter{Second}
\begin{corollaire}\label{c2}e\end{corollaire}
\begin{theoreme}\label{t2}f\end{theoreme}
RENVOIS \ref{p1} \ref{t1} \ref{r1} \ref{c1} \ref{c2} \ref{t2} ; \cref{p1} ; \cref{t1} ; \cref{r1} ; \cref{c1} ; \cref{c2} ; \cref{t2} FIN
\end{document}
"                                   sans correctif
%   lualatex -jobname=corr     "\def\CORR{1}% Juge de la troisième passe : document minimal écrit de zéro pour la cause A (thmtools sous LaTeX 2026-06-01).
% Variantes, par une macro posée sur la ligne de commande (deux passes chacune, depuis memoire/essai/juge3/minimal/) :
%   lualatex -jobname=nu       "\input{minimal-juge3}"                                   sans correctif
%   lualatex -jobname=corr     "\def\CORR{1}\input{minimal-juge3}"                       correctif A du préambule
%   lualatex -jobname=crochet  "\def\SANSCROCHET{1}\input{minimal-juge3}"                sans correctif, crochet retiré
%   lualatex -jobname=amsseul  "\def\AMSSEUL{1}\input{minimal-juge3}"                    amsthm seul (témoin)
%   lualatex -jobname=faux     "\def\CORR{1}\def\FAUX{1}\input{minimal-juge3}"           deux vrais défauts : doivent tomber
%   lualatex -jobname=trace    "\def\TRACE{1}\input{minimal-juge3}"                      trace de l'appel
% Les lignes du correctif sont tirées du préambule commité par extraire_correctifs.py, pas retapées.
\documentclass{scrbook}
\ifdefined\SANSCROCHET \RemoveFromHook{package/amsthm/after}[firstaid/aliascounter]\fi
\usepackage{amsthm}
\usepackage{hyperref}
\usepackage[french, nameinlink, noabbrev]{cleveref}
\ifdefined\AMSSEUL
  \newtheorem{theoreme}{Théorème}[chapter]
  \newtheorem{proposition}[theoreme]{Proposition}
  \newtheorem{corollaire}[theoreme]{Corollaire}
  \newtheorem{remarque}{Remarque}[chapter]
\else
  \usepackage{thmtools}
  \ifdefined\CORR \input{correctif-a-pose}\fi
  \declaretheorem[numberwithin=chapter, name=Théorème]{theoreme}
  \ifdefined\TRACE \tracingmacros=2 \tracingonline=0 \fi
  \declaretheorem[sibling=theoreme, name=Proposition]{proposition}
  \ifdefined\TRACE \tracingmacros=0 \fi
  \declaretheorem[sibling=theoreme, name=Corollaire]{corollaire}
  \declaretheorem[numberwithin=chapter, name=Remarque]{remarque}
  \ifdefined\FAUX \newcounteralias{corollaire}{section}\newcounteralias{orphelin}{inexistant}\fi
  \ifdefined\CORR \input{correctif-a-remise}\fi
\fi
\crefname{theoreme}{théorème}{théorèmes}
\crefname{proposition}{proposition}{propositions}
\crefname{corollaire}{corollaire}{corollaires}
\crefname{remarque}{remarque}{remarques}
\begin{document}
\chapter{Premier}
\begin{proposition}\label{p1}a\end{proposition}
\begin{theoreme}\label{t1}b\end{theoreme}
\begin{remarque}\label{r1}c\end{remarque}
\begin{corollaire}\label{c1}d\end{corollaire}
\chapter{Second}
\begin{corollaire}\label{c2}e\end{corollaire}
\begin{theoreme}\label{t2}f\end{theoreme}
RENVOIS \ref{p1} \ref{t1} \ref{r1} \ref{c1} \ref{c2} \ref{t2} ; \cref{p1} ; \cref{t1} ; \cref{r1} ; \cref{c1} ; \cref{c2} ; \cref{t2} FIN
\end{document}
"                       correctif A du préambule
%   lualatex -jobname=crochet  "\def\SANSCROCHET{1}% Juge de la troisième passe : document minimal écrit de zéro pour la cause A (thmtools sous LaTeX 2026-06-01).
% Variantes, par une macro posée sur la ligne de commande (deux passes chacune, depuis memoire/essai/juge3/minimal/) :
%   lualatex -jobname=nu       "\input{minimal-juge3}"                                   sans correctif
%   lualatex -jobname=corr     "\def\CORR{1}\input{minimal-juge3}"                       correctif A du préambule
%   lualatex -jobname=crochet  "\def\SANSCROCHET{1}\input{minimal-juge3}"                sans correctif, crochet retiré
%   lualatex -jobname=amsseul  "\def\AMSSEUL{1}\input{minimal-juge3}"                    amsthm seul (témoin)
%   lualatex -jobname=faux     "\def\CORR{1}\def\FAUX{1}\input{minimal-juge3}"           deux vrais défauts : doivent tomber
%   lualatex -jobname=trace    "\def\TRACE{1}\input{minimal-juge3}"                      trace de l'appel
% Les lignes du correctif sont tirées du préambule commité par extraire_correctifs.py, pas retapées.
\documentclass{scrbook}
\ifdefined\SANSCROCHET \RemoveFromHook{package/amsthm/after}[firstaid/aliascounter]\fi
\usepackage{amsthm}
\usepackage{hyperref}
\usepackage[french, nameinlink, noabbrev]{cleveref}
\ifdefined\AMSSEUL
  \newtheorem{theoreme}{Théorème}[chapter]
  \newtheorem{proposition}[theoreme]{Proposition}
  \newtheorem{corollaire}[theoreme]{Corollaire}
  \newtheorem{remarque}{Remarque}[chapter]
\else
  \usepackage{thmtools}
  \ifdefined\CORR \input{correctif-a-pose}\fi
  \declaretheorem[numberwithin=chapter, name=Théorème]{theoreme}
  \ifdefined\TRACE \tracingmacros=2 \tracingonline=0 \fi
  \declaretheorem[sibling=theoreme, name=Proposition]{proposition}
  \ifdefined\TRACE \tracingmacros=0 \fi
  \declaretheorem[sibling=theoreme, name=Corollaire]{corollaire}
  \declaretheorem[numberwithin=chapter, name=Remarque]{remarque}
  \ifdefined\FAUX \newcounteralias{corollaire}{section}\newcounteralias{orphelin}{inexistant}\fi
  \ifdefined\CORR \input{correctif-a-remise}\fi
\fi
\crefname{theoreme}{théorème}{théorèmes}
\crefname{proposition}{proposition}{propositions}
\crefname{corollaire}{corollaire}{corollaires}
\crefname{remarque}{remarque}{remarques}
\begin{document}
\chapter{Premier}
\begin{proposition}\label{p1}a\end{proposition}
\begin{theoreme}\label{t1}b\end{theoreme}
\begin{remarque}\label{r1}c\end{remarque}
\begin{corollaire}\label{c1}d\end{corollaire}
\chapter{Second}
\begin{corollaire}\label{c2}e\end{corollaire}
\begin{theoreme}\label{t2}f\end{theoreme}
RENVOIS \ref{p1} \ref{t1} \ref{r1} \ref{c1} \ref{c2} \ref{t2} ; \cref{p1} ; \cref{t1} ; \cref{r1} ; \cref{c1} ; \cref{c2} ; \cref{t2} FIN
\end{document}
"                sans correctif, crochet retiré
%   lualatex -jobname=amsseul  "\def\AMSSEUL{1}% Juge de la troisième passe : document minimal écrit de zéro pour la cause A (thmtools sous LaTeX 2026-06-01).
% Variantes, par une macro posée sur la ligne de commande (deux passes chacune, depuis memoire/essai/juge3/minimal/) :
%   lualatex -jobname=nu       "\input{minimal-juge3}"                                   sans correctif
%   lualatex -jobname=corr     "\def\CORR{1}\input{minimal-juge3}"                       correctif A du préambule
%   lualatex -jobname=crochet  "\def\SANSCROCHET{1}\input{minimal-juge3}"                sans correctif, crochet retiré
%   lualatex -jobname=amsseul  "\def\AMSSEUL{1}\input{minimal-juge3}"                    amsthm seul (témoin)
%   lualatex -jobname=faux     "\def\CORR{1}\def\FAUX{1}\input{minimal-juge3}"           deux vrais défauts : doivent tomber
%   lualatex -jobname=trace    "\def\TRACE{1}\input{minimal-juge3}"                      trace de l'appel
% Les lignes du correctif sont tirées du préambule commité par extraire_correctifs.py, pas retapées.
\documentclass{scrbook}
\ifdefined\SANSCROCHET \RemoveFromHook{package/amsthm/after}[firstaid/aliascounter]\fi
\usepackage{amsthm}
\usepackage{hyperref}
\usepackage[french, nameinlink, noabbrev]{cleveref}
\ifdefined\AMSSEUL
  \newtheorem{theoreme}{Théorème}[chapter]
  \newtheorem{proposition}[theoreme]{Proposition}
  \newtheorem{corollaire}[theoreme]{Corollaire}
  \newtheorem{remarque}{Remarque}[chapter]
\else
  \usepackage{thmtools}
  \ifdefined\CORR \input{correctif-a-pose}\fi
  \declaretheorem[numberwithin=chapter, name=Théorème]{theoreme}
  \ifdefined\TRACE \tracingmacros=2 \tracingonline=0 \fi
  \declaretheorem[sibling=theoreme, name=Proposition]{proposition}
  \ifdefined\TRACE \tracingmacros=0 \fi
  \declaretheorem[sibling=theoreme, name=Corollaire]{corollaire}
  \declaretheorem[numberwithin=chapter, name=Remarque]{remarque}
  \ifdefined\FAUX \newcounteralias{corollaire}{section}\newcounteralias{orphelin}{inexistant}\fi
  \ifdefined\CORR \input{correctif-a-remise}\fi
\fi
\crefname{theoreme}{théorème}{théorèmes}
\crefname{proposition}{proposition}{propositions}
\crefname{corollaire}{corollaire}{corollaires}
\crefname{remarque}{remarque}{remarques}
\begin{document}
\chapter{Premier}
\begin{proposition}\label{p1}a\end{proposition}
\begin{theoreme}\label{t1}b\end{theoreme}
\begin{remarque}\label{r1}c\end{remarque}
\begin{corollaire}\label{c1}d\end{corollaire}
\chapter{Second}
\begin{corollaire}\label{c2}e\end{corollaire}
\begin{theoreme}\label{t2}f\end{theoreme}
RENVOIS \ref{p1} \ref{t1} \ref{r1} \ref{c1} \ref{c2} \ref{t2} ; \cref{p1} ; \cref{t1} ; \cref{r1} ; \cref{c1} ; \cref{c2} ; \cref{t2} FIN
\end{document}
"                    amsthm seul (témoin)
%   lualatex -jobname=faux     "\def\CORR{1}\def\FAUX{1}% Juge de la troisième passe : document minimal écrit de zéro pour la cause A (thmtools sous LaTeX 2026-06-01).
% Variantes, par une macro posée sur la ligne de commande (deux passes chacune, depuis memoire/essai/juge3/minimal/) :
%   lualatex -jobname=nu       "\input{minimal-juge3}"                                   sans correctif
%   lualatex -jobname=corr     "\def\CORR{1}\input{minimal-juge3}"                       correctif A du préambule
%   lualatex -jobname=crochet  "\def\SANSCROCHET{1}\input{minimal-juge3}"                sans correctif, crochet retiré
%   lualatex -jobname=amsseul  "\def\AMSSEUL{1}\input{minimal-juge3}"                    amsthm seul (témoin)
%   lualatex -jobname=faux     "\def\CORR{1}\def\FAUX{1}\input{minimal-juge3}"           deux vrais défauts : doivent tomber
%   lualatex -jobname=trace    "\def\TRACE{1}\input{minimal-juge3}"                      trace de l'appel
% Les lignes du correctif sont tirées du préambule commité par extraire_correctifs.py, pas retapées.
\documentclass{scrbook}
\ifdefined\SANSCROCHET \RemoveFromHook{package/amsthm/after}[firstaid/aliascounter]\fi
\usepackage{amsthm}
\usepackage{hyperref}
\usepackage[french, nameinlink, noabbrev]{cleveref}
\ifdefined\AMSSEUL
  \newtheorem{theoreme}{Théorème}[chapter]
  \newtheorem{proposition}[theoreme]{Proposition}
  \newtheorem{corollaire}[theoreme]{Corollaire}
  \newtheorem{remarque}{Remarque}[chapter]
\else
  \usepackage{thmtools}
  \ifdefined\CORR \input{correctif-a-pose}\fi
  \declaretheorem[numberwithin=chapter, name=Théorème]{theoreme}
  \ifdefined\TRACE \tracingmacros=2 \tracingonline=0 \fi
  \declaretheorem[sibling=theoreme, name=Proposition]{proposition}
  \ifdefined\TRACE \tracingmacros=0 \fi
  \declaretheorem[sibling=theoreme, name=Corollaire]{corollaire}
  \declaretheorem[numberwithin=chapter, name=Remarque]{remarque}
  \ifdefined\FAUX \newcounteralias{corollaire}{section}\newcounteralias{orphelin}{inexistant}\fi
  \ifdefined\CORR \input{correctif-a-remise}\fi
\fi
\crefname{theoreme}{théorème}{théorèmes}
\crefname{proposition}{proposition}{propositions}
\crefname{corollaire}{corollaire}{corollaires}
\crefname{remarque}{remarque}{remarques}
\begin{document}
\chapter{Premier}
\begin{proposition}\label{p1}a\end{proposition}
\begin{theoreme}\label{t1}b\end{theoreme}
\begin{remarque}\label{r1}c\end{remarque}
\begin{corollaire}\label{c1}d\end{corollaire}
\chapter{Second}
\begin{corollaire}\label{c2}e\end{corollaire}
\begin{theoreme}\label{t2}f\end{theoreme}
RENVOIS \ref{p1} \ref{t1} \ref{r1} \ref{c1} \ref{c2} \ref{t2} ; \cref{p1} ; \cref{t1} ; \cref{r1} ; \cref{c1} ; \cref{c2} ; \cref{t2} FIN
\end{document}
"           deux vrais défauts : doivent tomber
%   lualatex -jobname=trace    "\def\TRACE{1}% Juge de la troisième passe : document minimal écrit de zéro pour la cause A (thmtools sous LaTeX 2026-06-01).
% Variantes, par une macro posée sur la ligne de commande (deux passes chacune, depuis memoire/essai/juge3/minimal/) :
%   lualatex -jobname=nu       "\input{minimal-juge3}"                                   sans correctif
%   lualatex -jobname=corr     "\def\CORR{1}\input{minimal-juge3}"                       correctif A du préambule
%   lualatex -jobname=crochet  "\def\SANSCROCHET{1}\input{minimal-juge3}"                sans correctif, crochet retiré
%   lualatex -jobname=amsseul  "\def\AMSSEUL{1}\input{minimal-juge3}"                    amsthm seul (témoin)
%   lualatex -jobname=faux     "\def\CORR{1}\def\FAUX{1}\input{minimal-juge3}"           deux vrais défauts : doivent tomber
%   lualatex -jobname=trace    "\def\TRACE{1}\input{minimal-juge3}"                      trace de l'appel
% Les lignes du correctif sont tirées du préambule commité par extraire_correctifs.py, pas retapées.
\documentclass{scrbook}
\ifdefined\SANSCROCHET \RemoveFromHook{package/amsthm/after}[firstaid/aliascounter]\fi
\usepackage{amsthm}
\usepackage{hyperref}
\usepackage[french, nameinlink, noabbrev]{cleveref}
\ifdefined\AMSSEUL
  \newtheorem{theoreme}{Théorème}[chapter]
  \newtheorem{proposition}[theoreme]{Proposition}
  \newtheorem{corollaire}[theoreme]{Corollaire}
  \newtheorem{remarque}{Remarque}[chapter]
\else
  \usepackage{thmtools}
  \ifdefined\CORR \input{correctif-a-pose}\fi
  \declaretheorem[numberwithin=chapter, name=Théorème]{theoreme}
  \ifdefined\TRACE \tracingmacros=2 \tracingonline=0 \fi
  \declaretheorem[sibling=theoreme, name=Proposition]{proposition}
  \ifdefined\TRACE \tracingmacros=0 \fi
  \declaretheorem[sibling=theoreme, name=Corollaire]{corollaire}
  \declaretheorem[numberwithin=chapter, name=Remarque]{remarque}
  \ifdefined\FAUX \newcounteralias{corollaire}{section}\newcounteralias{orphelin}{inexistant}\fi
  \ifdefined\CORR \input{correctif-a-remise}\fi
\fi
\crefname{theoreme}{théorème}{théorèmes}
\crefname{proposition}{proposition}{propositions}
\crefname{corollaire}{corollaire}{corollaires}
\crefname{remarque}{remarque}{remarques}
\begin{document}
\chapter{Premier}
\begin{proposition}\label{p1}a\end{proposition}
\begin{theoreme}\label{t1}b\end{theoreme}
\begin{remarque}\label{r1}c\end{remarque}
\begin{corollaire}\label{c1}d\end{corollaire}
\chapter{Second}
\begin{corollaire}\label{c2}e\end{corollaire}
\begin{theoreme}\label{t2}f\end{theoreme}
RENVOIS \ref{p1} \ref{t1} \ref{r1} \ref{c1} \ref{c2} \ref{t2} ; \cref{p1} ; \cref{t1} ; \cref{r1} ; \cref{c1} ; \cref{c2} ; \cref{t2} FIN
\end{document}
"                      trace de l'appel
% Les lignes du correctif sont tirées du préambule commité par extraire_correctifs.py, pas retapées.
\documentclass{scrbook}
\ifdefined\SANSCROCHET \RemoveFromHook{package/amsthm/after}[firstaid/aliascounter]\fi
\usepackage{amsthm}
\usepackage{hyperref}
\usepackage[french, nameinlink, noabbrev]{cleveref}
\ifdefined\AMSSEUL
  \newtheorem{theoreme}{Théorème}[chapter]
  \newtheorem{proposition}[theoreme]{Proposition}
  \newtheorem{corollaire}[theoreme]{Corollaire}
  \newtheorem{remarque}{Remarque}[chapter]
\else
  \usepackage{thmtools}
  \ifdefined\CORR \input{correctif-a-pose}\fi
  \declaretheorem[numberwithin=chapter, name=Théorème]{theoreme}
  \ifdefined\TRACE \tracingmacros=2 \tracingonline=0 \fi
  \declaretheorem[sibling=theoreme, name=Proposition]{proposition}
  \ifdefined\TRACE \tracingmacros=0 \fi
  \declaretheorem[sibling=theoreme, name=Corollaire]{corollaire}
  \declaretheorem[numberwithin=chapter, name=Remarque]{remarque}
  \ifdefined\FAUX \newcounteralias{corollaire}{section}\newcounteralias{orphelin}{inexistant}\fi
  \ifdefined\CORR \input{correctif-a-remise}\fi
\fi
\crefname{theoreme}{théorème}{théorèmes}
\crefname{proposition}{proposition}{propositions}
\crefname{corollaire}{corollaire}{corollaires}
\crefname{remarque}{remarque}{remarques}
\begin{document}
\chapter{Premier}
\begin{proposition}\label{p1}a\end{proposition}
\begin{theoreme}\label{t1}b\end{theoreme}
\begin{remarque}\label{r1}c\end{remarque}
\begin{corollaire}\label{c1}d\end{corollaire}
\chapter{Second}
\begin{corollaire}\label{c2}e\end{corollaire}
\begin{theoreme}\label{t2}f\end{theoreme}
RENVOIS \ref{p1} \ref{t1} \ref{r1} \ref{c1} \ref{c2} \ref{t2} ; \cref{p1} ; \cref{t1} ; \cref{r1} ; \cref{c1} ; \cref{c2} ; \cref{t2} FIN
\end{document}
"                       correctif A du préambule
%   lualatex -jobname=crochet  "\def\SANSCROCHET{1}% Juge de la troisième passe : document minimal écrit de zéro pour la cause A (thmtools sous LaTeX 2026-06-01).
% Variantes, par une macro posée sur la ligne de commande (deux passes chacune, depuis memoire/essai/juge3/minimal/) :
%   lualatex -jobname=nu       "% Juge de la troisième passe : document minimal écrit de zéro pour la cause A (thmtools sous LaTeX 2026-06-01).
% Variantes, par une macro posée sur la ligne de commande (deux passes chacune, depuis memoire/essai/juge3/minimal/) :
%   lualatex -jobname=nu       "\input{minimal-juge3}"                                   sans correctif
%   lualatex -jobname=corr     "\def\CORR{1}\input{minimal-juge3}"                       correctif A du préambule
%   lualatex -jobname=crochet  "\def\SANSCROCHET{1}\input{minimal-juge3}"                sans correctif, crochet retiré
%   lualatex -jobname=amsseul  "\def\AMSSEUL{1}\input{minimal-juge3}"                    amsthm seul (témoin)
%   lualatex -jobname=faux     "\def\CORR{1}\def\FAUX{1}\input{minimal-juge3}"           deux vrais défauts : doivent tomber
%   lualatex -jobname=trace    "\def\TRACE{1}\input{minimal-juge3}"                      trace de l'appel
% Les lignes du correctif sont tirées du préambule commité par extraire_correctifs.py, pas retapées.
\documentclass{scrbook}
\ifdefined\SANSCROCHET \RemoveFromHook{package/amsthm/after}[firstaid/aliascounter]\fi
\usepackage{amsthm}
\usepackage{hyperref}
\usepackage[french, nameinlink, noabbrev]{cleveref}
\ifdefined\AMSSEUL
  \newtheorem{theoreme}{Théorème}[chapter]
  \newtheorem{proposition}[theoreme]{Proposition}
  \newtheorem{corollaire}[theoreme]{Corollaire}
  \newtheorem{remarque}{Remarque}[chapter]
\else
  \usepackage{thmtools}
  \ifdefined\CORR \input{correctif-a-pose}\fi
  \declaretheorem[numberwithin=chapter, name=Théorème]{theoreme}
  \ifdefined\TRACE \tracingmacros=2 \tracingonline=0 \fi
  \declaretheorem[sibling=theoreme, name=Proposition]{proposition}
  \ifdefined\TRACE \tracingmacros=0 \fi
  \declaretheorem[sibling=theoreme, name=Corollaire]{corollaire}
  \declaretheorem[numberwithin=chapter, name=Remarque]{remarque}
  \ifdefined\FAUX \newcounteralias{corollaire}{section}\newcounteralias{orphelin}{inexistant}\fi
  \ifdefined\CORR \input{correctif-a-remise}\fi
\fi
\crefname{theoreme}{théorème}{théorèmes}
\crefname{proposition}{proposition}{propositions}
\crefname{corollaire}{corollaire}{corollaires}
\crefname{remarque}{remarque}{remarques}
\begin{document}
\chapter{Premier}
\begin{proposition}\label{p1}a\end{proposition}
\begin{theoreme}\label{t1}b\end{theoreme}
\begin{remarque}\label{r1}c\end{remarque}
\begin{corollaire}\label{c1}d\end{corollaire}
\chapter{Second}
\begin{corollaire}\label{c2}e\end{corollaire}
\begin{theoreme}\label{t2}f\end{theoreme}
RENVOIS \ref{p1} \ref{t1} \ref{r1} \ref{c1} \ref{c2} \ref{t2} ; \cref{p1} ; \cref{t1} ; \cref{r1} ; \cref{c1} ; \cref{c2} ; \cref{t2} FIN
\end{document}
"                                   sans correctif
%   lualatex -jobname=corr     "\def\CORR{1}% Juge de la troisième passe : document minimal écrit de zéro pour la cause A (thmtools sous LaTeX 2026-06-01).
% Variantes, par une macro posée sur la ligne de commande (deux passes chacune, depuis memoire/essai/juge3/minimal/) :
%   lualatex -jobname=nu       "\input{minimal-juge3}"                                   sans correctif
%   lualatex -jobname=corr     "\def\CORR{1}\input{minimal-juge3}"                       correctif A du préambule
%   lualatex -jobname=crochet  "\def\SANSCROCHET{1}\input{minimal-juge3}"                sans correctif, crochet retiré
%   lualatex -jobname=amsseul  "\def\AMSSEUL{1}\input{minimal-juge3}"                    amsthm seul (témoin)
%   lualatex -jobname=faux     "\def\CORR{1}\def\FAUX{1}\input{minimal-juge3}"           deux vrais défauts : doivent tomber
%   lualatex -jobname=trace    "\def\TRACE{1}\input{minimal-juge3}"                      trace de l'appel
% Les lignes du correctif sont tirées du préambule commité par extraire_correctifs.py, pas retapées.
\documentclass{scrbook}
\ifdefined\SANSCROCHET \RemoveFromHook{package/amsthm/after}[firstaid/aliascounter]\fi
\usepackage{amsthm}
\usepackage{hyperref}
\usepackage[french, nameinlink, noabbrev]{cleveref}
\ifdefined\AMSSEUL
  \newtheorem{theoreme}{Théorème}[chapter]
  \newtheorem{proposition}[theoreme]{Proposition}
  \newtheorem{corollaire}[theoreme]{Corollaire}
  \newtheorem{remarque}{Remarque}[chapter]
\else
  \usepackage{thmtools}
  \ifdefined\CORR \input{correctif-a-pose}\fi
  \declaretheorem[numberwithin=chapter, name=Théorème]{theoreme}
  \ifdefined\TRACE \tracingmacros=2 \tracingonline=0 \fi
  \declaretheorem[sibling=theoreme, name=Proposition]{proposition}
  \ifdefined\TRACE \tracingmacros=0 \fi
  \declaretheorem[sibling=theoreme, name=Corollaire]{corollaire}
  \declaretheorem[numberwithin=chapter, name=Remarque]{remarque}
  \ifdefined\FAUX \newcounteralias{corollaire}{section}\newcounteralias{orphelin}{inexistant}\fi
  \ifdefined\CORR \input{correctif-a-remise}\fi
\fi
\crefname{theoreme}{théorème}{théorèmes}
\crefname{proposition}{proposition}{propositions}
\crefname{corollaire}{corollaire}{corollaires}
\crefname{remarque}{remarque}{remarques}
\begin{document}
\chapter{Premier}
\begin{proposition}\label{p1}a\end{proposition}
\begin{theoreme}\label{t1}b\end{theoreme}
\begin{remarque}\label{r1}c\end{remarque}
\begin{corollaire}\label{c1}d\end{corollaire}
\chapter{Second}
\begin{corollaire}\label{c2}e\end{corollaire}
\begin{theoreme}\label{t2}f\end{theoreme}
RENVOIS \ref{p1} \ref{t1} \ref{r1} \ref{c1} \ref{c2} \ref{t2} ; \cref{p1} ; \cref{t1} ; \cref{r1} ; \cref{c1} ; \cref{c2} ; \cref{t2} FIN
\end{document}
"                       correctif A du préambule
%   lualatex -jobname=crochet  "\def\SANSCROCHET{1}% Juge de la troisième passe : document minimal écrit de zéro pour la cause A (thmtools sous LaTeX 2026-06-01).
% Variantes, par une macro posée sur la ligne de commande (deux passes chacune, depuis memoire/essai/juge3/minimal/) :
%   lualatex -jobname=nu       "\input{minimal-juge3}"                                   sans correctif
%   lualatex -jobname=corr     "\def\CORR{1}\input{minimal-juge3}"                       correctif A du préambule
%   lualatex -jobname=crochet  "\def\SANSCROCHET{1}\input{minimal-juge3}"                sans correctif, crochet retiré
%   lualatex -jobname=amsseul  "\def\AMSSEUL{1}\input{minimal-juge3}"                    amsthm seul (témoin)
%   lualatex -jobname=faux     "\def\CORR{1}\def\FAUX{1}\input{minimal-juge3}"           deux vrais défauts : doivent tomber
%   lualatex -jobname=trace    "\def\TRACE{1}\input{minimal-juge3}"                      trace de l'appel
% Les lignes du correctif sont tirées du préambule commité par extraire_correctifs.py, pas retapées.
\documentclass{scrbook}
\ifdefined\SANSCROCHET \RemoveFromHook{package/amsthm/after}[firstaid/aliascounter]\fi
\usepackage{amsthm}
\usepackage{hyperref}
\usepackage[french, nameinlink, noabbrev]{cleveref}
\ifdefined\AMSSEUL
  \newtheorem{theoreme}{Théorème}[chapter]
  \newtheorem{proposition}[theoreme]{Proposition}
  \newtheorem{corollaire}[theoreme]{Corollaire}
  \newtheorem{remarque}{Remarque}[chapter]
\else
  \usepackage{thmtools}
  \ifdefined\CORR \input{correctif-a-pose}\fi
  \declaretheorem[numberwithin=chapter, name=Théorème]{theoreme}
  \ifdefined\TRACE \tracingmacros=2 \tracingonline=0 \fi
  \declaretheorem[sibling=theoreme, name=Proposition]{proposition}
  \ifdefined\TRACE \tracingmacros=0 \fi
  \declaretheorem[sibling=theoreme, name=Corollaire]{corollaire}
  \declaretheorem[numberwithin=chapter, name=Remarque]{remarque}
  \ifdefined\FAUX \newcounteralias{corollaire}{section}\newcounteralias{orphelin}{inexistant}\fi
  \ifdefined\CORR \input{correctif-a-remise}\fi
\fi
\crefname{theoreme}{théorème}{théorèmes}
\crefname{proposition}{proposition}{propositions}
\crefname{corollaire}{corollaire}{corollaires}
\crefname{remarque}{remarque}{remarques}
\begin{document}
\chapter{Premier}
\begin{proposition}\label{p1}a\end{proposition}
\begin{theoreme}\label{t1}b\end{theoreme}
\begin{remarque}\label{r1}c\end{remarque}
\begin{corollaire}\label{c1}d\end{corollaire}
\chapter{Second}
\begin{corollaire}\label{c2}e\end{corollaire}
\begin{theoreme}\label{t2}f\end{theoreme}
RENVOIS \ref{p1} \ref{t1} \ref{r1} \ref{c1} \ref{c2} \ref{t2} ; \cref{p1} ; \cref{t1} ; \cref{r1} ; \cref{c1} ; \cref{c2} ; \cref{t2} FIN
\end{document}
"                sans correctif, crochet retiré
%   lualatex -jobname=amsseul  "\def\AMSSEUL{1}% Juge de la troisième passe : document minimal écrit de zéro pour la cause A (thmtools sous LaTeX 2026-06-01).
% Variantes, par une macro posée sur la ligne de commande (deux passes chacune, depuis memoire/essai/juge3/minimal/) :
%   lualatex -jobname=nu       "\input{minimal-juge3}"                                   sans correctif
%   lualatex -jobname=corr     "\def\CORR{1}\input{minimal-juge3}"                       correctif A du préambule
%   lualatex -jobname=crochet  "\def\SANSCROCHET{1}\input{minimal-juge3}"                sans correctif, crochet retiré
%   lualatex -jobname=amsseul  "\def\AMSSEUL{1}\input{minimal-juge3}"                    amsthm seul (témoin)
%   lualatex -jobname=faux     "\def\CORR{1}\def\FAUX{1}\input{minimal-juge3}"           deux vrais défauts : doivent tomber
%   lualatex -jobname=trace    "\def\TRACE{1}\input{minimal-juge3}"                      trace de l'appel
% Les lignes du correctif sont tirées du préambule commité par extraire_correctifs.py, pas retapées.
\documentclass{scrbook}
\ifdefined\SANSCROCHET \RemoveFromHook{package/amsthm/after}[firstaid/aliascounter]\fi
\usepackage{amsthm}
\usepackage{hyperref}
\usepackage[french, nameinlink, noabbrev]{cleveref}
\ifdefined\AMSSEUL
  \newtheorem{theoreme}{Théorème}[chapter]
  \newtheorem{proposition}[theoreme]{Proposition}
  \newtheorem{corollaire}[theoreme]{Corollaire}
  \newtheorem{remarque}{Remarque}[chapter]
\else
  \usepackage{thmtools}
  \ifdefined\CORR \input{correctif-a-pose}\fi
  \declaretheorem[numberwithin=chapter, name=Théorème]{theoreme}
  \ifdefined\TRACE \tracingmacros=2 \tracingonline=0 \fi
  \declaretheorem[sibling=theoreme, name=Proposition]{proposition}
  \ifdefined\TRACE \tracingmacros=0 \fi
  \declaretheorem[sibling=theoreme, name=Corollaire]{corollaire}
  \declaretheorem[numberwithin=chapter, name=Remarque]{remarque}
  \ifdefined\FAUX \newcounteralias{corollaire}{section}\newcounteralias{orphelin}{inexistant}\fi
  \ifdefined\CORR \input{correctif-a-remise}\fi
\fi
\crefname{theoreme}{théorème}{théorèmes}
\crefname{proposition}{proposition}{propositions}
\crefname{corollaire}{corollaire}{corollaires}
\crefname{remarque}{remarque}{remarques}
\begin{document}
\chapter{Premier}
\begin{proposition}\label{p1}a\end{proposition}
\begin{theoreme}\label{t1}b\end{theoreme}
\begin{remarque}\label{r1}c\end{remarque}
\begin{corollaire}\label{c1}d\end{corollaire}
\chapter{Second}
\begin{corollaire}\label{c2}e\end{corollaire}
\begin{theoreme}\label{t2}f\end{theoreme}
RENVOIS \ref{p1} \ref{t1} \ref{r1} \ref{c1} \ref{c2} \ref{t2} ; \cref{p1} ; \cref{t1} ; \cref{r1} ; \cref{c1} ; \cref{c2} ; \cref{t2} FIN
\end{document}
"                    amsthm seul (témoin)
%   lualatex -jobname=faux     "\def\CORR{1}\def\FAUX{1}% Juge de la troisième passe : document minimal écrit de zéro pour la cause A (thmtools sous LaTeX 2026-06-01).
% Variantes, par une macro posée sur la ligne de commande (deux passes chacune, depuis memoire/essai/juge3/minimal/) :
%   lualatex -jobname=nu       "\input{minimal-juge3}"                                   sans correctif
%   lualatex -jobname=corr     "\def\CORR{1}\input{minimal-juge3}"                       correctif A du préambule
%   lualatex -jobname=crochet  "\def\SANSCROCHET{1}\input{minimal-juge3}"                sans correctif, crochet retiré
%   lualatex -jobname=amsseul  "\def\AMSSEUL{1}\input{minimal-juge3}"                    amsthm seul (témoin)
%   lualatex -jobname=faux     "\def\CORR{1}\def\FAUX{1}\input{minimal-juge3}"           deux vrais défauts : doivent tomber
%   lualatex -jobname=trace    "\def\TRACE{1}\input{minimal-juge3}"                      trace de l'appel
% Les lignes du correctif sont tirées du préambule commité par extraire_correctifs.py, pas retapées.
\documentclass{scrbook}
\ifdefined\SANSCROCHET \RemoveFromHook{package/amsthm/after}[firstaid/aliascounter]\fi
\usepackage{amsthm}
\usepackage{hyperref}
\usepackage[french, nameinlink, noabbrev]{cleveref}
\ifdefined\AMSSEUL
  \newtheorem{theoreme}{Théorème}[chapter]
  \newtheorem{proposition}[theoreme]{Proposition}
  \newtheorem{corollaire}[theoreme]{Corollaire}
  \newtheorem{remarque}{Remarque}[chapter]
\else
  \usepackage{thmtools}
  \ifdefined\CORR \input{correctif-a-pose}\fi
  \declaretheorem[numberwithin=chapter, name=Théorème]{theoreme}
  \ifdefined\TRACE \tracingmacros=2 \tracingonline=0 \fi
  \declaretheorem[sibling=theoreme, name=Proposition]{proposition}
  \ifdefined\TRACE \tracingmacros=0 \fi
  \declaretheorem[sibling=theoreme, name=Corollaire]{corollaire}
  \declaretheorem[numberwithin=chapter, name=Remarque]{remarque}
  \ifdefined\FAUX \newcounteralias{corollaire}{section}\newcounteralias{orphelin}{inexistant}\fi
  \ifdefined\CORR \input{correctif-a-remise}\fi
\fi
\crefname{theoreme}{théorème}{théorèmes}
\crefname{proposition}{proposition}{propositions}
\crefname{corollaire}{corollaire}{corollaires}
\crefname{remarque}{remarque}{remarques}
\begin{document}
\chapter{Premier}
\begin{proposition}\label{p1}a\end{proposition}
\begin{theoreme}\label{t1}b\end{theoreme}
\begin{remarque}\label{r1}c\end{remarque}
\begin{corollaire}\label{c1}d\end{corollaire}
\chapter{Second}
\begin{corollaire}\label{c2}e\end{corollaire}
\begin{theoreme}\label{t2}f\end{theoreme}
RENVOIS \ref{p1} \ref{t1} \ref{r1} \ref{c1} \ref{c2} \ref{t2} ; \cref{p1} ; \cref{t1} ; \cref{r1} ; \cref{c1} ; \cref{c2} ; \cref{t2} FIN
\end{document}
"           deux vrais défauts : doivent tomber
%   lualatex -jobname=trace    "\def\TRACE{1}% Juge de la troisième passe : document minimal écrit de zéro pour la cause A (thmtools sous LaTeX 2026-06-01).
% Variantes, par une macro posée sur la ligne de commande (deux passes chacune, depuis memoire/essai/juge3/minimal/) :
%   lualatex -jobname=nu       "\input{minimal-juge3}"                                   sans correctif
%   lualatex -jobname=corr     "\def\CORR{1}\input{minimal-juge3}"                       correctif A du préambule
%   lualatex -jobname=crochet  "\def\SANSCROCHET{1}\input{minimal-juge3}"                sans correctif, crochet retiré
%   lualatex -jobname=amsseul  "\def\AMSSEUL{1}\input{minimal-juge3}"                    amsthm seul (témoin)
%   lualatex -jobname=faux     "\def\CORR{1}\def\FAUX{1}\input{minimal-juge3}"           deux vrais défauts : doivent tomber
%   lualatex -jobname=trace    "\def\TRACE{1}\input{minimal-juge3}"                      trace de l'appel
% Les lignes du correctif sont tirées du préambule commité par extraire_correctifs.py, pas retapées.
\documentclass{scrbook}
\ifdefined\SANSCROCHET \RemoveFromHook{package/amsthm/after}[firstaid/aliascounter]\fi
\usepackage{amsthm}
\usepackage{hyperref}
\usepackage[french, nameinlink, noabbrev]{cleveref}
\ifdefined\AMSSEUL
  \newtheorem{theoreme}{Théorème}[chapter]
  \newtheorem{proposition}[theoreme]{Proposition}
  \newtheorem{corollaire}[theoreme]{Corollaire}
  \newtheorem{remarque}{Remarque}[chapter]
\else
  \usepackage{thmtools}
  \ifdefined\CORR \input{correctif-a-pose}\fi
  \declaretheorem[numberwithin=chapter, name=Théorème]{theoreme}
  \ifdefined\TRACE \tracingmacros=2 \tracingonline=0 \fi
  \declaretheorem[sibling=theoreme, name=Proposition]{proposition}
  \ifdefined\TRACE \tracingmacros=0 \fi
  \declaretheorem[sibling=theoreme, name=Corollaire]{corollaire}
  \declaretheorem[numberwithin=chapter, name=Remarque]{remarque}
  \ifdefined\FAUX \newcounteralias{corollaire}{section}\newcounteralias{orphelin}{inexistant}\fi
  \ifdefined\CORR \input{correctif-a-remise}\fi
\fi
\crefname{theoreme}{théorème}{théorèmes}
\crefname{proposition}{proposition}{propositions}
\crefname{corollaire}{corollaire}{corollaires}
\crefname{remarque}{remarque}{remarques}
\begin{document}
\chapter{Premier}
\begin{proposition}\label{p1}a\end{proposition}
\begin{theoreme}\label{t1}b\end{theoreme}
\begin{remarque}\label{r1}c\end{remarque}
\begin{corollaire}\label{c1}d\end{corollaire}
\chapter{Second}
\begin{corollaire}\label{c2}e\end{corollaire}
\begin{theoreme}\label{t2}f\end{theoreme}
RENVOIS \ref{p1} \ref{t1} \ref{r1} \ref{c1} \ref{c2} \ref{t2} ; \cref{p1} ; \cref{t1} ; \cref{r1} ; \cref{c1} ; \cref{c2} ; \cref{t2} FIN
\end{document}
"                      trace de l'appel
% Les lignes du correctif sont tirées du préambule commité par extraire_correctifs.py, pas retapées.
\documentclass{scrbook}
\ifdefined\SANSCROCHET \RemoveFromHook{package/amsthm/after}[firstaid/aliascounter]\fi
\usepackage{amsthm}
\usepackage{hyperref}
\usepackage[french, nameinlink, noabbrev]{cleveref}
\ifdefined\AMSSEUL
  \newtheorem{theoreme}{Théorème}[chapter]
  \newtheorem{proposition}[theoreme]{Proposition}
  \newtheorem{corollaire}[theoreme]{Corollaire}
  \newtheorem{remarque}{Remarque}[chapter]
\else
  \usepackage{thmtools}
  \ifdefined\CORR \input{correctif-a-pose}\fi
  \declaretheorem[numberwithin=chapter, name=Théorème]{theoreme}
  \ifdefined\TRACE \tracingmacros=2 \tracingonline=0 \fi
  \declaretheorem[sibling=theoreme, name=Proposition]{proposition}
  \ifdefined\TRACE \tracingmacros=0 \fi
  \declaretheorem[sibling=theoreme, name=Corollaire]{corollaire}
  \declaretheorem[numberwithin=chapter, name=Remarque]{remarque}
  \ifdefined\FAUX \newcounteralias{corollaire}{section}\newcounteralias{orphelin}{inexistant}\fi
  \ifdefined\CORR \input{correctif-a-remise}\fi
\fi
\crefname{theoreme}{théorème}{théorèmes}
\crefname{proposition}{proposition}{propositions}
\crefname{corollaire}{corollaire}{corollaires}
\crefname{remarque}{remarque}{remarques}
\begin{document}
\chapter{Premier}
\begin{proposition}\label{p1}a\end{proposition}
\begin{theoreme}\label{t1}b\end{theoreme}
\begin{remarque}\label{r1}c\end{remarque}
\begin{corollaire}\label{c1}d\end{corollaire}
\chapter{Second}
\begin{corollaire}\label{c2}e\end{corollaire}
\begin{theoreme}\label{t2}f\end{theoreme}
RENVOIS \ref{p1} \ref{t1} \ref{r1} \ref{c1} \ref{c2} \ref{t2} ; \cref{p1} ; \cref{t1} ; \cref{r1} ; \cref{c1} ; \cref{c2} ; \cref{t2} FIN
\end{document}
"                sans correctif, crochet retiré
%   lualatex -jobname=amsseul  "\def\AMSSEUL{1}% Juge de la troisième passe : document minimal écrit de zéro pour la cause A (thmtools sous LaTeX 2026-06-01).
% Variantes, par une macro posée sur la ligne de commande (deux passes chacune, depuis memoire/essai/juge3/minimal/) :
%   lualatex -jobname=nu       "% Juge de la troisième passe : document minimal écrit de zéro pour la cause A (thmtools sous LaTeX 2026-06-01).
% Variantes, par une macro posée sur la ligne de commande (deux passes chacune, depuis memoire/essai/juge3/minimal/) :
%   lualatex -jobname=nu       "\input{minimal-juge3}"                                   sans correctif
%   lualatex -jobname=corr     "\def\CORR{1}\input{minimal-juge3}"                       correctif A du préambule
%   lualatex -jobname=crochet  "\def\SANSCROCHET{1}\input{minimal-juge3}"                sans correctif, crochet retiré
%   lualatex -jobname=amsseul  "\def\AMSSEUL{1}\input{minimal-juge3}"                    amsthm seul (témoin)
%   lualatex -jobname=faux     "\def\CORR{1}\def\FAUX{1}\input{minimal-juge3}"           deux vrais défauts : doivent tomber
%   lualatex -jobname=trace    "\def\TRACE{1}\input{minimal-juge3}"                      trace de l'appel
% Les lignes du correctif sont tirées du préambule commité par extraire_correctifs.py, pas retapées.
\documentclass{scrbook}
\ifdefined\SANSCROCHET \RemoveFromHook{package/amsthm/after}[firstaid/aliascounter]\fi
\usepackage{amsthm}
\usepackage{hyperref}
\usepackage[french, nameinlink, noabbrev]{cleveref}
\ifdefined\AMSSEUL
  \newtheorem{theoreme}{Théorème}[chapter]
  \newtheorem{proposition}[theoreme]{Proposition}
  \newtheorem{corollaire}[theoreme]{Corollaire}
  \newtheorem{remarque}{Remarque}[chapter]
\else
  \usepackage{thmtools}
  \ifdefined\CORR \input{correctif-a-pose}\fi
  \declaretheorem[numberwithin=chapter, name=Théorème]{theoreme}
  \ifdefined\TRACE \tracingmacros=2 \tracingonline=0 \fi
  \declaretheorem[sibling=theoreme, name=Proposition]{proposition}
  \ifdefined\TRACE \tracingmacros=0 \fi
  \declaretheorem[sibling=theoreme, name=Corollaire]{corollaire}
  \declaretheorem[numberwithin=chapter, name=Remarque]{remarque}
  \ifdefined\FAUX \newcounteralias{corollaire}{section}\newcounteralias{orphelin}{inexistant}\fi
  \ifdefined\CORR \input{correctif-a-remise}\fi
\fi
\crefname{theoreme}{théorème}{théorèmes}
\crefname{proposition}{proposition}{propositions}
\crefname{corollaire}{corollaire}{corollaires}
\crefname{remarque}{remarque}{remarques}
\begin{document}
\chapter{Premier}
\begin{proposition}\label{p1}a\end{proposition}
\begin{theoreme}\label{t1}b\end{theoreme}
\begin{remarque}\label{r1}c\end{remarque}
\begin{corollaire}\label{c1}d\end{corollaire}
\chapter{Second}
\begin{corollaire}\label{c2}e\end{corollaire}
\begin{theoreme}\label{t2}f\end{theoreme}
RENVOIS \ref{p1} \ref{t1} \ref{r1} \ref{c1} \ref{c2} \ref{t2} ; \cref{p1} ; \cref{t1} ; \cref{r1} ; \cref{c1} ; \cref{c2} ; \cref{t2} FIN
\end{document}
"                                   sans correctif
%   lualatex -jobname=corr     "\def\CORR{1}% Juge de la troisième passe : document minimal écrit de zéro pour la cause A (thmtools sous LaTeX 2026-06-01).
% Variantes, par une macro posée sur la ligne de commande (deux passes chacune, depuis memoire/essai/juge3/minimal/) :
%   lualatex -jobname=nu       "\input{minimal-juge3}"                                   sans correctif
%   lualatex -jobname=corr     "\def\CORR{1}\input{minimal-juge3}"                       correctif A du préambule
%   lualatex -jobname=crochet  "\def\SANSCROCHET{1}\input{minimal-juge3}"                sans correctif, crochet retiré
%   lualatex -jobname=amsseul  "\def\AMSSEUL{1}\input{minimal-juge3}"                    amsthm seul (témoin)
%   lualatex -jobname=faux     "\def\CORR{1}\def\FAUX{1}\input{minimal-juge3}"           deux vrais défauts : doivent tomber
%   lualatex -jobname=trace    "\def\TRACE{1}\input{minimal-juge3}"                      trace de l'appel
% Les lignes du correctif sont tirées du préambule commité par extraire_correctifs.py, pas retapées.
\documentclass{scrbook}
\ifdefined\SANSCROCHET \RemoveFromHook{package/amsthm/after}[firstaid/aliascounter]\fi
\usepackage{amsthm}
\usepackage{hyperref}
\usepackage[french, nameinlink, noabbrev]{cleveref}
\ifdefined\AMSSEUL
  \newtheorem{theoreme}{Théorème}[chapter]
  \newtheorem{proposition}[theoreme]{Proposition}
  \newtheorem{corollaire}[theoreme]{Corollaire}
  \newtheorem{remarque}{Remarque}[chapter]
\else
  \usepackage{thmtools}
  \ifdefined\CORR \input{correctif-a-pose}\fi
  \declaretheorem[numberwithin=chapter, name=Théorème]{theoreme}
  \ifdefined\TRACE \tracingmacros=2 \tracingonline=0 \fi
  \declaretheorem[sibling=theoreme, name=Proposition]{proposition}
  \ifdefined\TRACE \tracingmacros=0 \fi
  \declaretheorem[sibling=theoreme, name=Corollaire]{corollaire}
  \declaretheorem[numberwithin=chapter, name=Remarque]{remarque}
  \ifdefined\FAUX \newcounteralias{corollaire}{section}\newcounteralias{orphelin}{inexistant}\fi
  \ifdefined\CORR \input{correctif-a-remise}\fi
\fi
\crefname{theoreme}{théorème}{théorèmes}
\crefname{proposition}{proposition}{propositions}
\crefname{corollaire}{corollaire}{corollaires}
\crefname{remarque}{remarque}{remarques}
\begin{document}
\chapter{Premier}
\begin{proposition}\label{p1}a\end{proposition}
\begin{theoreme}\label{t1}b\end{theoreme}
\begin{remarque}\label{r1}c\end{remarque}
\begin{corollaire}\label{c1}d\end{corollaire}
\chapter{Second}
\begin{corollaire}\label{c2}e\end{corollaire}
\begin{theoreme}\label{t2}f\end{theoreme}
RENVOIS \ref{p1} \ref{t1} \ref{r1} \ref{c1} \ref{c2} \ref{t2} ; \cref{p1} ; \cref{t1} ; \cref{r1} ; \cref{c1} ; \cref{c2} ; \cref{t2} FIN
\end{document}
"                       correctif A du préambule
%   lualatex -jobname=crochet  "\def\SANSCROCHET{1}% Juge de la troisième passe : document minimal écrit de zéro pour la cause A (thmtools sous LaTeX 2026-06-01).
% Variantes, par une macro posée sur la ligne de commande (deux passes chacune, depuis memoire/essai/juge3/minimal/) :
%   lualatex -jobname=nu       "\input{minimal-juge3}"                                   sans correctif
%   lualatex -jobname=corr     "\def\CORR{1}\input{minimal-juge3}"                       correctif A du préambule
%   lualatex -jobname=crochet  "\def\SANSCROCHET{1}\input{minimal-juge3}"                sans correctif, crochet retiré
%   lualatex -jobname=amsseul  "\def\AMSSEUL{1}\input{minimal-juge3}"                    amsthm seul (témoin)
%   lualatex -jobname=faux     "\def\CORR{1}\def\FAUX{1}\input{minimal-juge3}"           deux vrais défauts : doivent tomber
%   lualatex -jobname=trace    "\def\TRACE{1}\input{minimal-juge3}"                      trace de l'appel
% Les lignes du correctif sont tirées du préambule commité par extraire_correctifs.py, pas retapées.
\documentclass{scrbook}
\ifdefined\SANSCROCHET \RemoveFromHook{package/amsthm/after}[firstaid/aliascounter]\fi
\usepackage{amsthm}
\usepackage{hyperref}
\usepackage[french, nameinlink, noabbrev]{cleveref}
\ifdefined\AMSSEUL
  \newtheorem{theoreme}{Théorème}[chapter]
  \newtheorem{proposition}[theoreme]{Proposition}
  \newtheorem{corollaire}[theoreme]{Corollaire}
  \newtheorem{remarque}{Remarque}[chapter]
\else
  \usepackage{thmtools}
  \ifdefined\CORR \input{correctif-a-pose}\fi
  \declaretheorem[numberwithin=chapter, name=Théorème]{theoreme}
  \ifdefined\TRACE \tracingmacros=2 \tracingonline=0 \fi
  \declaretheorem[sibling=theoreme, name=Proposition]{proposition}
  \ifdefined\TRACE \tracingmacros=0 \fi
  \declaretheorem[sibling=theoreme, name=Corollaire]{corollaire}
  \declaretheorem[numberwithin=chapter, name=Remarque]{remarque}
  \ifdefined\FAUX \newcounteralias{corollaire}{section}\newcounteralias{orphelin}{inexistant}\fi
  \ifdefined\CORR \input{correctif-a-remise}\fi
\fi
\crefname{theoreme}{théorème}{théorèmes}
\crefname{proposition}{proposition}{propositions}
\crefname{corollaire}{corollaire}{corollaires}
\crefname{remarque}{remarque}{remarques}
\begin{document}
\chapter{Premier}
\begin{proposition}\label{p1}a\end{proposition}
\begin{theoreme}\label{t1}b\end{theoreme}
\begin{remarque}\label{r1}c\end{remarque}
\begin{corollaire}\label{c1}d\end{corollaire}
\chapter{Second}
\begin{corollaire}\label{c2}e\end{corollaire}
\begin{theoreme}\label{t2}f\end{theoreme}
RENVOIS \ref{p1} \ref{t1} \ref{r1} \ref{c1} \ref{c2} \ref{t2} ; \cref{p1} ; \cref{t1} ; \cref{r1} ; \cref{c1} ; \cref{c2} ; \cref{t2} FIN
\end{document}
"                sans correctif, crochet retiré
%   lualatex -jobname=amsseul  "\def\AMSSEUL{1}% Juge de la troisième passe : document minimal écrit de zéro pour la cause A (thmtools sous LaTeX 2026-06-01).
% Variantes, par une macro posée sur la ligne de commande (deux passes chacune, depuis memoire/essai/juge3/minimal/) :
%   lualatex -jobname=nu       "\input{minimal-juge3}"                                   sans correctif
%   lualatex -jobname=corr     "\def\CORR{1}\input{minimal-juge3}"                       correctif A du préambule
%   lualatex -jobname=crochet  "\def\SANSCROCHET{1}\input{minimal-juge3}"                sans correctif, crochet retiré
%   lualatex -jobname=amsseul  "\def\AMSSEUL{1}\input{minimal-juge3}"                    amsthm seul (témoin)
%   lualatex -jobname=faux     "\def\CORR{1}\def\FAUX{1}\input{minimal-juge3}"           deux vrais défauts : doivent tomber
%   lualatex -jobname=trace    "\def\TRACE{1}\input{minimal-juge3}"                      trace de l'appel
% Les lignes du correctif sont tirées du préambule commité par extraire_correctifs.py, pas retapées.
\documentclass{scrbook}
\ifdefined\SANSCROCHET \RemoveFromHook{package/amsthm/after}[firstaid/aliascounter]\fi
\usepackage{amsthm}
\usepackage{hyperref}
\usepackage[french, nameinlink, noabbrev]{cleveref}
\ifdefined\AMSSEUL
  \newtheorem{theoreme}{Théorème}[chapter]
  \newtheorem{proposition}[theoreme]{Proposition}
  \newtheorem{corollaire}[theoreme]{Corollaire}
  \newtheorem{remarque}{Remarque}[chapter]
\else
  \usepackage{thmtools}
  \ifdefined\CORR \input{correctif-a-pose}\fi
  \declaretheorem[numberwithin=chapter, name=Théorème]{theoreme}
  \ifdefined\TRACE \tracingmacros=2 \tracingonline=0 \fi
  \declaretheorem[sibling=theoreme, name=Proposition]{proposition}
  \ifdefined\TRACE \tracingmacros=0 \fi
  \declaretheorem[sibling=theoreme, name=Corollaire]{corollaire}
  \declaretheorem[numberwithin=chapter, name=Remarque]{remarque}
  \ifdefined\FAUX \newcounteralias{corollaire}{section}\newcounteralias{orphelin}{inexistant}\fi
  \ifdefined\CORR \input{correctif-a-remise}\fi
\fi
\crefname{theoreme}{théorème}{théorèmes}
\crefname{proposition}{proposition}{propositions}
\crefname{corollaire}{corollaire}{corollaires}
\crefname{remarque}{remarque}{remarques}
\begin{document}
\chapter{Premier}
\begin{proposition}\label{p1}a\end{proposition}
\begin{theoreme}\label{t1}b\end{theoreme}
\begin{remarque}\label{r1}c\end{remarque}
\begin{corollaire}\label{c1}d\end{corollaire}
\chapter{Second}
\begin{corollaire}\label{c2}e\end{corollaire}
\begin{theoreme}\label{t2}f\end{theoreme}
RENVOIS \ref{p1} \ref{t1} \ref{r1} \ref{c1} \ref{c2} \ref{t2} ; \cref{p1} ; \cref{t1} ; \cref{r1} ; \cref{c1} ; \cref{c2} ; \cref{t2} FIN
\end{document}
"                    amsthm seul (témoin)
%   lualatex -jobname=faux     "\def\CORR{1}\def\FAUX{1}% Juge de la troisième passe : document minimal écrit de zéro pour la cause A (thmtools sous LaTeX 2026-06-01).
% Variantes, par une macro posée sur la ligne de commande (deux passes chacune, depuis memoire/essai/juge3/minimal/) :
%   lualatex -jobname=nu       "\input{minimal-juge3}"                                   sans correctif
%   lualatex -jobname=corr     "\def\CORR{1}\input{minimal-juge3}"                       correctif A du préambule
%   lualatex -jobname=crochet  "\def\SANSCROCHET{1}\input{minimal-juge3}"                sans correctif, crochet retiré
%   lualatex -jobname=amsseul  "\def\AMSSEUL{1}\input{minimal-juge3}"                    amsthm seul (témoin)
%   lualatex -jobname=faux     "\def\CORR{1}\def\FAUX{1}\input{minimal-juge3}"           deux vrais défauts : doivent tomber
%   lualatex -jobname=trace    "\def\TRACE{1}\input{minimal-juge3}"                      trace de l'appel
% Les lignes du correctif sont tirées du préambule commité par extraire_correctifs.py, pas retapées.
\documentclass{scrbook}
\ifdefined\SANSCROCHET \RemoveFromHook{package/amsthm/after}[firstaid/aliascounter]\fi
\usepackage{amsthm}
\usepackage{hyperref}
\usepackage[french, nameinlink, noabbrev]{cleveref}
\ifdefined\AMSSEUL
  \newtheorem{theoreme}{Théorème}[chapter]
  \newtheorem{proposition}[theoreme]{Proposition}
  \newtheorem{corollaire}[theoreme]{Corollaire}
  \newtheorem{remarque}{Remarque}[chapter]
\else
  \usepackage{thmtools}
  \ifdefined\CORR \input{correctif-a-pose}\fi
  \declaretheorem[numberwithin=chapter, name=Théorème]{theoreme}
  \ifdefined\TRACE \tracingmacros=2 \tracingonline=0 \fi
  \declaretheorem[sibling=theoreme, name=Proposition]{proposition}
  \ifdefined\TRACE \tracingmacros=0 \fi
  \declaretheorem[sibling=theoreme, name=Corollaire]{corollaire}
  \declaretheorem[numberwithin=chapter, name=Remarque]{remarque}
  \ifdefined\FAUX \newcounteralias{corollaire}{section}\newcounteralias{orphelin}{inexistant}\fi
  \ifdefined\CORR \input{correctif-a-remise}\fi
\fi
\crefname{theoreme}{théorème}{théorèmes}
\crefname{proposition}{proposition}{propositions}
\crefname{corollaire}{corollaire}{corollaires}
\crefname{remarque}{remarque}{remarques}
\begin{document}
\chapter{Premier}
\begin{proposition}\label{p1}a\end{proposition}
\begin{theoreme}\label{t1}b\end{theoreme}
\begin{remarque}\label{r1}c\end{remarque}
\begin{corollaire}\label{c1}d\end{corollaire}
\chapter{Second}
\begin{corollaire}\label{c2}e\end{corollaire}
\begin{theoreme}\label{t2}f\end{theoreme}
RENVOIS \ref{p1} \ref{t1} \ref{r1} \ref{c1} \ref{c2} \ref{t2} ; \cref{p1} ; \cref{t1} ; \cref{r1} ; \cref{c1} ; \cref{c2} ; \cref{t2} FIN
\end{document}
"           deux vrais défauts : doivent tomber
%   lualatex -jobname=trace    "\def\TRACE{1}% Juge de la troisième passe : document minimal écrit de zéro pour la cause A (thmtools sous LaTeX 2026-06-01).
% Variantes, par une macro posée sur la ligne de commande (deux passes chacune, depuis memoire/essai/juge3/minimal/) :
%   lualatex -jobname=nu       "\input{minimal-juge3}"                                   sans correctif
%   lualatex -jobname=corr     "\def\CORR{1}\input{minimal-juge3}"                       correctif A du préambule
%   lualatex -jobname=crochet  "\def\SANSCROCHET{1}\input{minimal-juge3}"                sans correctif, crochet retiré
%   lualatex -jobname=amsseul  "\def\AMSSEUL{1}\input{minimal-juge3}"                    amsthm seul (témoin)
%   lualatex -jobname=faux     "\def\CORR{1}\def\FAUX{1}\input{minimal-juge3}"           deux vrais défauts : doivent tomber
%   lualatex -jobname=trace    "\def\TRACE{1}\input{minimal-juge3}"                      trace de l'appel
% Les lignes du correctif sont tirées du préambule commité par extraire_correctifs.py, pas retapées.
\documentclass{scrbook}
\ifdefined\SANSCROCHET \RemoveFromHook{package/amsthm/after}[firstaid/aliascounter]\fi
\usepackage{amsthm}
\usepackage{hyperref}
\usepackage[french, nameinlink, noabbrev]{cleveref}
\ifdefined\AMSSEUL
  \newtheorem{theoreme}{Théorème}[chapter]
  \newtheorem{proposition}[theoreme]{Proposition}
  \newtheorem{corollaire}[theoreme]{Corollaire}
  \newtheorem{remarque}{Remarque}[chapter]
\else
  \usepackage{thmtools}
  \ifdefined\CORR \input{correctif-a-pose}\fi
  \declaretheorem[numberwithin=chapter, name=Théorème]{theoreme}
  \ifdefined\TRACE \tracingmacros=2 \tracingonline=0 \fi
  \declaretheorem[sibling=theoreme, name=Proposition]{proposition}
  \ifdefined\TRACE \tracingmacros=0 \fi
  \declaretheorem[sibling=theoreme, name=Corollaire]{corollaire}
  \declaretheorem[numberwithin=chapter, name=Remarque]{remarque}
  \ifdefined\FAUX \newcounteralias{corollaire}{section}\newcounteralias{orphelin}{inexistant}\fi
  \ifdefined\CORR \input{correctif-a-remise}\fi
\fi
\crefname{theoreme}{théorème}{théorèmes}
\crefname{proposition}{proposition}{propositions}
\crefname{corollaire}{corollaire}{corollaires}
\crefname{remarque}{remarque}{remarques}
\begin{document}
\chapter{Premier}
\begin{proposition}\label{p1}a\end{proposition}
\begin{theoreme}\label{t1}b\end{theoreme}
\begin{remarque}\label{r1}c\end{remarque}
\begin{corollaire}\label{c1}d\end{corollaire}
\chapter{Second}
\begin{corollaire}\label{c2}e\end{corollaire}
\begin{theoreme}\label{t2}f\end{theoreme}
RENVOIS \ref{p1} \ref{t1} \ref{r1} \ref{c1} \ref{c2} \ref{t2} ; \cref{p1} ; \cref{t1} ; \cref{r1} ; \cref{c1} ; \cref{c2} ; \cref{t2} FIN
\end{document}
"                      trace de l'appel
% Les lignes du correctif sont tirées du préambule commité par extraire_correctifs.py, pas retapées.
\documentclass{scrbook}
\ifdefined\SANSCROCHET \RemoveFromHook{package/amsthm/after}[firstaid/aliascounter]\fi
\usepackage{amsthm}
\usepackage{hyperref}
\usepackage[french, nameinlink, noabbrev]{cleveref}
\ifdefined\AMSSEUL
  \newtheorem{theoreme}{Théorème}[chapter]
  \newtheorem{proposition}[theoreme]{Proposition}
  \newtheorem{corollaire}[theoreme]{Corollaire}
  \newtheorem{remarque}{Remarque}[chapter]
\else
  \usepackage{thmtools}
  \ifdefined\CORR \input{correctif-a-pose}\fi
  \declaretheorem[numberwithin=chapter, name=Théorème]{theoreme}
  \ifdefined\TRACE \tracingmacros=2 \tracingonline=0 \fi
  \declaretheorem[sibling=theoreme, name=Proposition]{proposition}
  \ifdefined\TRACE \tracingmacros=0 \fi
  \declaretheorem[sibling=theoreme, name=Corollaire]{corollaire}
  \declaretheorem[numberwithin=chapter, name=Remarque]{remarque}
  \ifdefined\FAUX \newcounteralias{corollaire}{section}\newcounteralias{orphelin}{inexistant}\fi
  \ifdefined\CORR \input{correctif-a-remise}\fi
\fi
\crefname{theoreme}{théorème}{théorèmes}
\crefname{proposition}{proposition}{propositions}
\crefname{corollaire}{corollaire}{corollaires}
\crefname{remarque}{remarque}{remarques}
\begin{document}
\chapter{Premier}
\begin{proposition}\label{p1}a\end{proposition}
\begin{theoreme}\label{t1}b\end{theoreme}
\begin{remarque}\label{r1}c\end{remarque}
\begin{corollaire}\label{c1}d\end{corollaire}
\chapter{Second}
\begin{corollaire}\label{c2}e\end{corollaire}
\begin{theoreme}\label{t2}f\end{theoreme}
RENVOIS \ref{p1} \ref{t1} \ref{r1} \ref{c1} \ref{c2} \ref{t2} ; \cref{p1} ; \cref{t1} ; \cref{r1} ; \cref{c1} ; \cref{c2} ; \cref{t2} FIN
\end{document}
"                    amsthm seul (témoin)
%   lualatex -jobname=faux     "\def\CORR{1}\def\FAUX{1}% Juge de la troisième passe : document minimal écrit de zéro pour la cause A (thmtools sous LaTeX 2026-06-01).
% Variantes, par une macro posée sur la ligne de commande (deux passes chacune, depuis memoire/essai/juge3/minimal/) :
%   lualatex -jobname=nu       "% Juge de la troisième passe : document minimal écrit de zéro pour la cause A (thmtools sous LaTeX 2026-06-01).
% Variantes, par une macro posée sur la ligne de commande (deux passes chacune, depuis memoire/essai/juge3/minimal/) :
%   lualatex -jobname=nu       "\input{minimal-juge3}"                                   sans correctif
%   lualatex -jobname=corr     "\def\CORR{1}\input{minimal-juge3}"                       correctif A du préambule
%   lualatex -jobname=crochet  "\def\SANSCROCHET{1}\input{minimal-juge3}"                sans correctif, crochet retiré
%   lualatex -jobname=amsseul  "\def\AMSSEUL{1}\input{minimal-juge3}"                    amsthm seul (témoin)
%   lualatex -jobname=faux     "\def\CORR{1}\def\FAUX{1}\input{minimal-juge3}"           deux vrais défauts : doivent tomber
%   lualatex -jobname=trace    "\def\TRACE{1}\input{minimal-juge3}"                      trace de l'appel
% Les lignes du correctif sont tirées du préambule commité par extraire_correctifs.py, pas retapées.
\documentclass{scrbook}
\ifdefined\SANSCROCHET \RemoveFromHook{package/amsthm/after}[firstaid/aliascounter]\fi
\usepackage{amsthm}
\usepackage{hyperref}
\usepackage[french, nameinlink, noabbrev]{cleveref}
\ifdefined\AMSSEUL
  \newtheorem{theoreme}{Théorème}[chapter]
  \newtheorem{proposition}[theoreme]{Proposition}
  \newtheorem{corollaire}[theoreme]{Corollaire}
  \newtheorem{remarque}{Remarque}[chapter]
\else
  \usepackage{thmtools}
  \ifdefined\CORR \input{correctif-a-pose}\fi
  \declaretheorem[numberwithin=chapter, name=Théorème]{theoreme}
  \ifdefined\TRACE \tracingmacros=2 \tracingonline=0 \fi
  \declaretheorem[sibling=theoreme, name=Proposition]{proposition}
  \ifdefined\TRACE \tracingmacros=0 \fi
  \declaretheorem[sibling=theoreme, name=Corollaire]{corollaire}
  \declaretheorem[numberwithin=chapter, name=Remarque]{remarque}
  \ifdefined\FAUX \newcounteralias{corollaire}{section}\newcounteralias{orphelin}{inexistant}\fi
  \ifdefined\CORR \input{correctif-a-remise}\fi
\fi
\crefname{theoreme}{théorème}{théorèmes}
\crefname{proposition}{proposition}{propositions}
\crefname{corollaire}{corollaire}{corollaires}
\crefname{remarque}{remarque}{remarques}
\begin{document}
\chapter{Premier}
\begin{proposition}\label{p1}a\end{proposition}
\begin{theoreme}\label{t1}b\end{theoreme}
\begin{remarque}\label{r1}c\end{remarque}
\begin{corollaire}\label{c1}d\end{corollaire}
\chapter{Second}
\begin{corollaire}\label{c2}e\end{corollaire}
\begin{theoreme}\label{t2}f\end{theoreme}
RENVOIS \ref{p1} \ref{t1} \ref{r1} \ref{c1} \ref{c2} \ref{t2} ; \cref{p1} ; \cref{t1} ; \cref{r1} ; \cref{c1} ; \cref{c2} ; \cref{t2} FIN
\end{document}
"                                   sans correctif
%   lualatex -jobname=corr     "\def\CORR{1}% Juge de la troisième passe : document minimal écrit de zéro pour la cause A (thmtools sous LaTeX 2026-06-01).
% Variantes, par une macro posée sur la ligne de commande (deux passes chacune, depuis memoire/essai/juge3/minimal/) :
%   lualatex -jobname=nu       "\input{minimal-juge3}"                                   sans correctif
%   lualatex -jobname=corr     "\def\CORR{1}\input{minimal-juge3}"                       correctif A du préambule
%   lualatex -jobname=crochet  "\def\SANSCROCHET{1}\input{minimal-juge3}"                sans correctif, crochet retiré
%   lualatex -jobname=amsseul  "\def\AMSSEUL{1}\input{minimal-juge3}"                    amsthm seul (témoin)
%   lualatex -jobname=faux     "\def\CORR{1}\def\FAUX{1}\input{minimal-juge3}"           deux vrais défauts : doivent tomber
%   lualatex -jobname=trace    "\def\TRACE{1}\input{minimal-juge3}"                      trace de l'appel
% Les lignes du correctif sont tirées du préambule commité par extraire_correctifs.py, pas retapées.
\documentclass{scrbook}
\ifdefined\SANSCROCHET \RemoveFromHook{package/amsthm/after}[firstaid/aliascounter]\fi
\usepackage{amsthm}
\usepackage{hyperref}
\usepackage[french, nameinlink, noabbrev]{cleveref}
\ifdefined\AMSSEUL
  \newtheorem{theoreme}{Théorème}[chapter]
  \newtheorem{proposition}[theoreme]{Proposition}
  \newtheorem{corollaire}[theoreme]{Corollaire}
  \newtheorem{remarque}{Remarque}[chapter]
\else
  \usepackage{thmtools}
  \ifdefined\CORR \input{correctif-a-pose}\fi
  \declaretheorem[numberwithin=chapter, name=Théorème]{theoreme}
  \ifdefined\TRACE \tracingmacros=2 \tracingonline=0 \fi
  \declaretheorem[sibling=theoreme, name=Proposition]{proposition}
  \ifdefined\TRACE \tracingmacros=0 \fi
  \declaretheorem[sibling=theoreme, name=Corollaire]{corollaire}
  \declaretheorem[numberwithin=chapter, name=Remarque]{remarque}
  \ifdefined\FAUX \newcounteralias{corollaire}{section}\newcounteralias{orphelin}{inexistant}\fi
  \ifdefined\CORR \input{correctif-a-remise}\fi
\fi
\crefname{theoreme}{théorème}{théorèmes}
\crefname{proposition}{proposition}{propositions}
\crefname{corollaire}{corollaire}{corollaires}
\crefname{remarque}{remarque}{remarques}
\begin{document}
\chapter{Premier}
\begin{proposition}\label{p1}a\end{proposition}
\begin{theoreme}\label{t1}b\end{theoreme}
\begin{remarque}\label{r1}c\end{remarque}
\begin{corollaire}\label{c1}d\end{corollaire}
\chapter{Second}
\begin{corollaire}\label{c2}e\end{corollaire}
\begin{theoreme}\label{t2}f\end{theoreme}
RENVOIS \ref{p1} \ref{t1} \ref{r1} \ref{c1} \ref{c2} \ref{t2} ; \cref{p1} ; \cref{t1} ; \cref{r1} ; \cref{c1} ; \cref{c2} ; \cref{t2} FIN
\end{document}
"                       correctif A du préambule
%   lualatex -jobname=crochet  "\def\SANSCROCHET{1}% Juge de la troisième passe : document minimal écrit de zéro pour la cause A (thmtools sous LaTeX 2026-06-01).
% Variantes, par une macro posée sur la ligne de commande (deux passes chacune, depuis memoire/essai/juge3/minimal/) :
%   lualatex -jobname=nu       "\input{minimal-juge3}"                                   sans correctif
%   lualatex -jobname=corr     "\def\CORR{1}\input{minimal-juge3}"                       correctif A du préambule
%   lualatex -jobname=crochet  "\def\SANSCROCHET{1}\input{minimal-juge3}"                sans correctif, crochet retiré
%   lualatex -jobname=amsseul  "\def\AMSSEUL{1}\input{minimal-juge3}"                    amsthm seul (témoin)
%   lualatex -jobname=faux     "\def\CORR{1}\def\FAUX{1}\input{minimal-juge3}"           deux vrais défauts : doivent tomber
%   lualatex -jobname=trace    "\def\TRACE{1}\input{minimal-juge3}"                      trace de l'appel
% Les lignes du correctif sont tirées du préambule commité par extraire_correctifs.py, pas retapées.
\documentclass{scrbook}
\ifdefined\SANSCROCHET \RemoveFromHook{package/amsthm/after}[firstaid/aliascounter]\fi
\usepackage{amsthm}
\usepackage{hyperref}
\usepackage[french, nameinlink, noabbrev]{cleveref}
\ifdefined\AMSSEUL
  \newtheorem{theoreme}{Théorème}[chapter]
  \newtheorem{proposition}[theoreme]{Proposition}
  \newtheorem{corollaire}[theoreme]{Corollaire}
  \newtheorem{remarque}{Remarque}[chapter]
\else
  \usepackage{thmtools}
  \ifdefined\CORR \input{correctif-a-pose}\fi
  \declaretheorem[numberwithin=chapter, name=Théorème]{theoreme}
  \ifdefined\TRACE \tracingmacros=2 \tracingonline=0 \fi
  \declaretheorem[sibling=theoreme, name=Proposition]{proposition}
  \ifdefined\TRACE \tracingmacros=0 \fi
  \declaretheorem[sibling=theoreme, name=Corollaire]{corollaire}
  \declaretheorem[numberwithin=chapter, name=Remarque]{remarque}
  \ifdefined\FAUX \newcounteralias{corollaire}{section}\newcounteralias{orphelin}{inexistant}\fi
  \ifdefined\CORR \input{correctif-a-remise}\fi
\fi
\crefname{theoreme}{théorème}{théorèmes}
\crefname{proposition}{proposition}{propositions}
\crefname{corollaire}{corollaire}{corollaires}
\crefname{remarque}{remarque}{remarques}
\begin{document}
\chapter{Premier}
\begin{proposition}\label{p1}a\end{proposition}
\begin{theoreme}\label{t1}b\end{theoreme}
\begin{remarque}\label{r1}c\end{remarque}
\begin{corollaire}\label{c1}d\end{corollaire}
\chapter{Second}
\begin{corollaire}\label{c2}e\end{corollaire}
\begin{theoreme}\label{t2}f\end{theoreme}
RENVOIS \ref{p1} \ref{t1} \ref{r1} \ref{c1} \ref{c2} \ref{t2} ; \cref{p1} ; \cref{t1} ; \cref{r1} ; \cref{c1} ; \cref{c2} ; \cref{t2} FIN
\end{document}
"                sans correctif, crochet retiré
%   lualatex -jobname=amsseul  "\def\AMSSEUL{1}% Juge de la troisième passe : document minimal écrit de zéro pour la cause A (thmtools sous LaTeX 2026-06-01).
% Variantes, par une macro posée sur la ligne de commande (deux passes chacune, depuis memoire/essai/juge3/minimal/) :
%   lualatex -jobname=nu       "\input{minimal-juge3}"                                   sans correctif
%   lualatex -jobname=corr     "\def\CORR{1}\input{minimal-juge3}"                       correctif A du préambule
%   lualatex -jobname=crochet  "\def\SANSCROCHET{1}\input{minimal-juge3}"                sans correctif, crochet retiré
%   lualatex -jobname=amsseul  "\def\AMSSEUL{1}\input{minimal-juge3}"                    amsthm seul (témoin)
%   lualatex -jobname=faux     "\def\CORR{1}\def\FAUX{1}\input{minimal-juge3}"           deux vrais défauts : doivent tomber
%   lualatex -jobname=trace    "\def\TRACE{1}\input{minimal-juge3}"                      trace de l'appel
% Les lignes du correctif sont tirées du préambule commité par extraire_correctifs.py, pas retapées.
\documentclass{scrbook}
\ifdefined\SANSCROCHET \RemoveFromHook{package/amsthm/after}[firstaid/aliascounter]\fi
\usepackage{amsthm}
\usepackage{hyperref}
\usepackage[french, nameinlink, noabbrev]{cleveref}
\ifdefined\AMSSEUL
  \newtheorem{theoreme}{Théorème}[chapter]
  \newtheorem{proposition}[theoreme]{Proposition}
  \newtheorem{corollaire}[theoreme]{Corollaire}
  \newtheorem{remarque}{Remarque}[chapter]
\else
  \usepackage{thmtools}
  \ifdefined\CORR \input{correctif-a-pose}\fi
  \declaretheorem[numberwithin=chapter, name=Théorème]{theoreme}
  \ifdefined\TRACE \tracingmacros=2 \tracingonline=0 \fi
  \declaretheorem[sibling=theoreme, name=Proposition]{proposition}
  \ifdefined\TRACE \tracingmacros=0 \fi
  \declaretheorem[sibling=theoreme, name=Corollaire]{corollaire}
  \declaretheorem[numberwithin=chapter, name=Remarque]{remarque}
  \ifdefined\FAUX \newcounteralias{corollaire}{section}\newcounteralias{orphelin}{inexistant}\fi
  \ifdefined\CORR \input{correctif-a-remise}\fi
\fi
\crefname{theoreme}{théorème}{théorèmes}
\crefname{proposition}{proposition}{propositions}
\crefname{corollaire}{corollaire}{corollaires}
\crefname{remarque}{remarque}{remarques}
\begin{document}
\chapter{Premier}
\begin{proposition}\label{p1}a\end{proposition}
\begin{theoreme}\label{t1}b\end{theoreme}
\begin{remarque}\label{r1}c\end{remarque}
\begin{corollaire}\label{c1}d\end{corollaire}
\chapter{Second}
\begin{corollaire}\label{c2}e\end{corollaire}
\begin{theoreme}\label{t2}f\end{theoreme}
RENVOIS \ref{p1} \ref{t1} \ref{r1} \ref{c1} \ref{c2} \ref{t2} ; \cref{p1} ; \cref{t1} ; \cref{r1} ; \cref{c1} ; \cref{c2} ; \cref{t2} FIN
\end{document}
"                    amsthm seul (témoin)
%   lualatex -jobname=faux     "\def\CORR{1}\def\FAUX{1}% Juge de la troisième passe : document minimal écrit de zéro pour la cause A (thmtools sous LaTeX 2026-06-01).
% Variantes, par une macro posée sur la ligne de commande (deux passes chacune, depuis memoire/essai/juge3/minimal/) :
%   lualatex -jobname=nu       "\input{minimal-juge3}"                                   sans correctif
%   lualatex -jobname=corr     "\def\CORR{1}\input{minimal-juge3}"                       correctif A du préambule
%   lualatex -jobname=crochet  "\def\SANSCROCHET{1}\input{minimal-juge3}"                sans correctif, crochet retiré
%   lualatex -jobname=amsseul  "\def\AMSSEUL{1}\input{minimal-juge3}"                    amsthm seul (témoin)
%   lualatex -jobname=faux     "\def\CORR{1}\def\FAUX{1}\input{minimal-juge3}"           deux vrais défauts : doivent tomber
%   lualatex -jobname=trace    "\def\TRACE{1}\input{minimal-juge3}"                      trace de l'appel
% Les lignes du correctif sont tirées du préambule commité par extraire_correctifs.py, pas retapées.
\documentclass{scrbook}
\ifdefined\SANSCROCHET \RemoveFromHook{package/amsthm/after}[firstaid/aliascounter]\fi
\usepackage{amsthm}
\usepackage{hyperref}
\usepackage[french, nameinlink, noabbrev]{cleveref}
\ifdefined\AMSSEUL
  \newtheorem{theoreme}{Théorème}[chapter]
  \newtheorem{proposition}[theoreme]{Proposition}
  \newtheorem{corollaire}[theoreme]{Corollaire}
  \newtheorem{remarque}{Remarque}[chapter]
\else
  \usepackage{thmtools}
  \ifdefined\CORR \input{correctif-a-pose}\fi
  \declaretheorem[numberwithin=chapter, name=Théorème]{theoreme}
  \ifdefined\TRACE \tracingmacros=2 \tracingonline=0 \fi
  \declaretheorem[sibling=theoreme, name=Proposition]{proposition}
  \ifdefined\TRACE \tracingmacros=0 \fi
  \declaretheorem[sibling=theoreme, name=Corollaire]{corollaire}
  \declaretheorem[numberwithin=chapter, name=Remarque]{remarque}
  \ifdefined\FAUX \newcounteralias{corollaire}{section}\newcounteralias{orphelin}{inexistant}\fi
  \ifdefined\CORR \input{correctif-a-remise}\fi
\fi
\crefname{theoreme}{théorème}{théorèmes}
\crefname{proposition}{proposition}{propositions}
\crefname{corollaire}{corollaire}{corollaires}
\crefname{remarque}{remarque}{remarques}
\begin{document}
\chapter{Premier}
\begin{proposition}\label{p1}a\end{proposition}
\begin{theoreme}\label{t1}b\end{theoreme}
\begin{remarque}\label{r1}c\end{remarque}
\begin{corollaire}\label{c1}d\end{corollaire}
\chapter{Second}
\begin{corollaire}\label{c2}e\end{corollaire}
\begin{theoreme}\label{t2}f\end{theoreme}
RENVOIS \ref{p1} \ref{t1} \ref{r1} \ref{c1} \ref{c2} \ref{t2} ; \cref{p1} ; \cref{t1} ; \cref{r1} ; \cref{c1} ; \cref{c2} ; \cref{t2} FIN
\end{document}
"           deux vrais défauts : doivent tomber
%   lualatex -jobname=trace    "\def\TRACE{1}% Juge de la troisième passe : document minimal écrit de zéro pour la cause A (thmtools sous LaTeX 2026-06-01).
% Variantes, par une macro posée sur la ligne de commande (deux passes chacune, depuis memoire/essai/juge3/minimal/) :
%   lualatex -jobname=nu       "\input{minimal-juge3}"                                   sans correctif
%   lualatex -jobname=corr     "\def\CORR{1}\input{minimal-juge3}"                       correctif A du préambule
%   lualatex -jobname=crochet  "\def\SANSCROCHET{1}\input{minimal-juge3}"                sans correctif, crochet retiré
%   lualatex -jobname=amsseul  "\def\AMSSEUL{1}\input{minimal-juge3}"                    amsthm seul (témoin)
%   lualatex -jobname=faux     "\def\CORR{1}\def\FAUX{1}\input{minimal-juge3}"           deux vrais défauts : doivent tomber
%   lualatex -jobname=trace    "\def\TRACE{1}\input{minimal-juge3}"                      trace de l'appel
% Les lignes du correctif sont tirées du préambule commité par extraire_correctifs.py, pas retapées.
\documentclass{scrbook}
\ifdefined\SANSCROCHET \RemoveFromHook{package/amsthm/after}[firstaid/aliascounter]\fi
\usepackage{amsthm}
\usepackage{hyperref}
\usepackage[french, nameinlink, noabbrev]{cleveref}
\ifdefined\AMSSEUL
  \newtheorem{theoreme}{Théorème}[chapter]
  \newtheorem{proposition}[theoreme]{Proposition}
  \newtheorem{corollaire}[theoreme]{Corollaire}
  \newtheorem{remarque}{Remarque}[chapter]
\else
  \usepackage{thmtools}
  \ifdefined\CORR \input{correctif-a-pose}\fi
  \declaretheorem[numberwithin=chapter, name=Théorème]{theoreme}
  \ifdefined\TRACE \tracingmacros=2 \tracingonline=0 \fi
  \declaretheorem[sibling=theoreme, name=Proposition]{proposition}
  \ifdefined\TRACE \tracingmacros=0 \fi
  \declaretheorem[sibling=theoreme, name=Corollaire]{corollaire}
  \declaretheorem[numberwithin=chapter, name=Remarque]{remarque}
  \ifdefined\FAUX \newcounteralias{corollaire}{section}\newcounteralias{orphelin}{inexistant}\fi
  \ifdefined\CORR \input{correctif-a-remise}\fi
\fi
\crefname{theoreme}{théorème}{théorèmes}
\crefname{proposition}{proposition}{propositions}
\crefname{corollaire}{corollaire}{corollaires}
\crefname{remarque}{remarque}{remarques}
\begin{document}
\chapter{Premier}
\begin{proposition}\label{p1}a\end{proposition}
\begin{theoreme}\label{t1}b\end{theoreme}
\begin{remarque}\label{r1}c\end{remarque}
\begin{corollaire}\label{c1}d\end{corollaire}
\chapter{Second}
\begin{corollaire}\label{c2}e\end{corollaire}
\begin{theoreme}\label{t2}f\end{theoreme}
RENVOIS \ref{p1} \ref{t1} \ref{r1} \ref{c1} \ref{c2} \ref{t2} ; \cref{p1} ; \cref{t1} ; \cref{r1} ; \cref{c1} ; \cref{c2} ; \cref{t2} FIN
\end{document}
"                      trace de l'appel
% Les lignes du correctif sont tirées du préambule commité par extraire_correctifs.py, pas retapées.
\documentclass{scrbook}
\ifdefined\SANSCROCHET \RemoveFromHook{package/amsthm/after}[firstaid/aliascounter]\fi
\usepackage{amsthm}
\usepackage{hyperref}
\usepackage[french, nameinlink, noabbrev]{cleveref}
\ifdefined\AMSSEUL
  \newtheorem{theoreme}{Théorème}[chapter]
  \newtheorem{proposition}[theoreme]{Proposition}
  \newtheorem{corollaire}[theoreme]{Corollaire}
  \newtheorem{remarque}{Remarque}[chapter]
\else
  \usepackage{thmtools}
  \ifdefined\CORR \input{correctif-a-pose}\fi
  \declaretheorem[numberwithin=chapter, name=Théorème]{theoreme}
  \ifdefined\TRACE \tracingmacros=2 \tracingonline=0 \fi
  \declaretheorem[sibling=theoreme, name=Proposition]{proposition}
  \ifdefined\TRACE \tracingmacros=0 \fi
  \declaretheorem[sibling=theoreme, name=Corollaire]{corollaire}
  \declaretheorem[numberwithin=chapter, name=Remarque]{remarque}
  \ifdefined\FAUX \newcounteralias{corollaire}{section}\newcounteralias{orphelin}{inexistant}\fi
  \ifdefined\CORR \input{correctif-a-remise}\fi
\fi
\crefname{theoreme}{théorème}{théorèmes}
\crefname{proposition}{proposition}{propositions}
\crefname{corollaire}{corollaire}{corollaires}
\crefname{remarque}{remarque}{remarques}
\begin{document}
\chapter{Premier}
\begin{proposition}\label{p1}a\end{proposition}
\begin{theoreme}\label{t1}b\end{theoreme}
\begin{remarque}\label{r1}c\end{remarque}
\begin{corollaire}\label{c1}d\end{corollaire}
\chapter{Second}
\begin{corollaire}\label{c2}e\end{corollaire}
\begin{theoreme}\label{t2}f\end{theoreme}
RENVOIS \ref{p1} \ref{t1} \ref{r1} \ref{c1} \ref{c2} \ref{t2} ; \cref{p1} ; \cref{t1} ; \cref{r1} ; \cref{c1} ; \cref{c2} ; \cref{t2} FIN
\end{document}
"           deux vrais défauts : doivent tomber
%   lualatex -jobname=trace    "\def\TRACE{1}% Juge de la troisième passe : document minimal écrit de zéro pour la cause A (thmtools sous LaTeX 2026-06-01).
% Variantes, par une macro posée sur la ligne de commande (deux passes chacune, depuis memoire/essai/juge3/minimal/) :
%   lualatex -jobname=nu       "% Juge de la troisième passe : document minimal écrit de zéro pour la cause A (thmtools sous LaTeX 2026-06-01).
% Variantes, par une macro posée sur la ligne de commande (deux passes chacune, depuis memoire/essai/juge3/minimal/) :
%   lualatex -jobname=nu       "\input{minimal-juge3}"                                   sans correctif
%   lualatex -jobname=corr     "\def\CORR{1}\input{minimal-juge3}"                       correctif A du préambule
%   lualatex -jobname=crochet  "\def\SANSCROCHET{1}\input{minimal-juge3}"                sans correctif, crochet retiré
%   lualatex -jobname=amsseul  "\def\AMSSEUL{1}\input{minimal-juge3}"                    amsthm seul (témoin)
%   lualatex -jobname=faux     "\def\CORR{1}\def\FAUX{1}\input{minimal-juge3}"           deux vrais défauts : doivent tomber
%   lualatex -jobname=trace    "\def\TRACE{1}\input{minimal-juge3}"                      trace de l'appel
% Les lignes du correctif sont tirées du préambule commité par extraire_correctifs.py, pas retapées.
\documentclass{scrbook}
\ifdefined\SANSCROCHET \RemoveFromHook{package/amsthm/after}[firstaid/aliascounter]\fi
\usepackage{amsthm}
\usepackage{hyperref}
\usepackage[french, nameinlink, noabbrev]{cleveref}
\ifdefined\AMSSEUL
  \newtheorem{theoreme}{Théorème}[chapter]
  \newtheorem{proposition}[theoreme]{Proposition}
  \newtheorem{corollaire}[theoreme]{Corollaire}
  \newtheorem{remarque}{Remarque}[chapter]
\else
  \usepackage{thmtools}
  \ifdefined\CORR \input{correctif-a-pose}\fi
  \declaretheorem[numberwithin=chapter, name=Théorème]{theoreme}
  \ifdefined\TRACE \tracingmacros=2 \tracingonline=0 \fi
  \declaretheorem[sibling=theoreme, name=Proposition]{proposition}
  \ifdefined\TRACE \tracingmacros=0 \fi
  \declaretheorem[sibling=theoreme, name=Corollaire]{corollaire}
  \declaretheorem[numberwithin=chapter, name=Remarque]{remarque}
  \ifdefined\FAUX \newcounteralias{corollaire}{section}\newcounteralias{orphelin}{inexistant}\fi
  \ifdefined\CORR \input{correctif-a-remise}\fi
\fi
\crefname{theoreme}{théorème}{théorèmes}
\crefname{proposition}{proposition}{propositions}
\crefname{corollaire}{corollaire}{corollaires}
\crefname{remarque}{remarque}{remarques}
\begin{document}
\chapter{Premier}
\begin{proposition}\label{p1}a\end{proposition}
\begin{theoreme}\label{t1}b\end{theoreme}
\begin{remarque}\label{r1}c\end{remarque}
\begin{corollaire}\label{c1}d\end{corollaire}
\chapter{Second}
\begin{corollaire}\label{c2}e\end{corollaire}
\begin{theoreme}\label{t2}f\end{theoreme}
RENVOIS \ref{p1} \ref{t1} \ref{r1} \ref{c1} \ref{c2} \ref{t2} ; \cref{p1} ; \cref{t1} ; \cref{r1} ; \cref{c1} ; \cref{c2} ; \cref{t2} FIN
\end{document}
"                                   sans correctif
%   lualatex -jobname=corr     "\def\CORR{1}% Juge de la troisième passe : document minimal écrit de zéro pour la cause A (thmtools sous LaTeX 2026-06-01).
% Variantes, par une macro posée sur la ligne de commande (deux passes chacune, depuis memoire/essai/juge3/minimal/) :
%   lualatex -jobname=nu       "\input{minimal-juge3}"                                   sans correctif
%   lualatex -jobname=corr     "\def\CORR{1}\input{minimal-juge3}"                       correctif A du préambule
%   lualatex -jobname=crochet  "\def\SANSCROCHET{1}\input{minimal-juge3}"                sans correctif, crochet retiré
%   lualatex -jobname=amsseul  "\def\AMSSEUL{1}\input{minimal-juge3}"                    amsthm seul (témoin)
%   lualatex -jobname=faux     "\def\CORR{1}\def\FAUX{1}\input{minimal-juge3}"           deux vrais défauts : doivent tomber
%   lualatex -jobname=trace    "\def\TRACE{1}\input{minimal-juge3}"                      trace de l'appel
% Les lignes du correctif sont tirées du préambule commité par extraire_correctifs.py, pas retapées.
\documentclass{scrbook}
\ifdefined\SANSCROCHET \RemoveFromHook{package/amsthm/after}[firstaid/aliascounter]\fi
\usepackage{amsthm}
\usepackage{hyperref}
\usepackage[french, nameinlink, noabbrev]{cleveref}
\ifdefined\AMSSEUL
  \newtheorem{theoreme}{Théorème}[chapter]
  \newtheorem{proposition}[theoreme]{Proposition}
  \newtheorem{corollaire}[theoreme]{Corollaire}
  \newtheorem{remarque}{Remarque}[chapter]
\else
  \usepackage{thmtools}
  \ifdefined\CORR \input{correctif-a-pose}\fi
  \declaretheorem[numberwithin=chapter, name=Théorème]{theoreme}
  \ifdefined\TRACE \tracingmacros=2 \tracingonline=0 \fi
  \declaretheorem[sibling=theoreme, name=Proposition]{proposition}
  \ifdefined\TRACE \tracingmacros=0 \fi
  \declaretheorem[sibling=theoreme, name=Corollaire]{corollaire}
  \declaretheorem[numberwithin=chapter, name=Remarque]{remarque}
  \ifdefined\FAUX \newcounteralias{corollaire}{section}\newcounteralias{orphelin}{inexistant}\fi
  \ifdefined\CORR \input{correctif-a-remise}\fi
\fi
\crefname{theoreme}{théorème}{théorèmes}
\crefname{proposition}{proposition}{propositions}
\crefname{corollaire}{corollaire}{corollaires}
\crefname{remarque}{remarque}{remarques}
\begin{document}
\chapter{Premier}
\begin{proposition}\label{p1}a\end{proposition}
\begin{theoreme}\label{t1}b\end{theoreme}
\begin{remarque}\label{r1}c\end{remarque}
\begin{corollaire}\label{c1}d\end{corollaire}
\chapter{Second}
\begin{corollaire}\label{c2}e\end{corollaire}
\begin{theoreme}\label{t2}f\end{theoreme}
RENVOIS \ref{p1} \ref{t1} \ref{r1} \ref{c1} \ref{c2} \ref{t2} ; \cref{p1} ; \cref{t1} ; \cref{r1} ; \cref{c1} ; \cref{c2} ; \cref{t2} FIN
\end{document}
"                       correctif A du préambule
%   lualatex -jobname=crochet  "\def\SANSCROCHET{1}% Juge de la troisième passe : document minimal écrit de zéro pour la cause A (thmtools sous LaTeX 2026-06-01).
% Variantes, par une macro posée sur la ligne de commande (deux passes chacune, depuis memoire/essai/juge3/minimal/) :
%   lualatex -jobname=nu       "\input{minimal-juge3}"                                   sans correctif
%   lualatex -jobname=corr     "\def\CORR{1}\input{minimal-juge3}"                       correctif A du préambule
%   lualatex -jobname=crochet  "\def\SANSCROCHET{1}\input{minimal-juge3}"                sans correctif, crochet retiré
%   lualatex -jobname=amsseul  "\def\AMSSEUL{1}\input{minimal-juge3}"                    amsthm seul (témoin)
%   lualatex -jobname=faux     "\def\CORR{1}\def\FAUX{1}\input{minimal-juge3}"           deux vrais défauts : doivent tomber
%   lualatex -jobname=trace    "\def\TRACE{1}\input{minimal-juge3}"                      trace de l'appel
% Les lignes du correctif sont tirées du préambule commité par extraire_correctifs.py, pas retapées.
\documentclass{scrbook}
\ifdefined\SANSCROCHET \RemoveFromHook{package/amsthm/after}[firstaid/aliascounter]\fi
\usepackage{amsthm}
\usepackage{hyperref}
\usepackage[french, nameinlink, noabbrev]{cleveref}
\ifdefined\AMSSEUL
  \newtheorem{theoreme}{Théorème}[chapter]
  \newtheorem{proposition}[theoreme]{Proposition}
  \newtheorem{corollaire}[theoreme]{Corollaire}
  \newtheorem{remarque}{Remarque}[chapter]
\else
  \usepackage{thmtools}
  \ifdefined\CORR \input{correctif-a-pose}\fi
  \declaretheorem[numberwithin=chapter, name=Théorème]{theoreme}
  \ifdefined\TRACE \tracingmacros=2 \tracingonline=0 \fi
  \declaretheorem[sibling=theoreme, name=Proposition]{proposition}
  \ifdefined\TRACE \tracingmacros=0 \fi
  \declaretheorem[sibling=theoreme, name=Corollaire]{corollaire}
  \declaretheorem[numberwithin=chapter, name=Remarque]{remarque}
  \ifdefined\FAUX \newcounteralias{corollaire}{section}\newcounteralias{orphelin}{inexistant}\fi
  \ifdefined\CORR \input{correctif-a-remise}\fi
\fi
\crefname{theoreme}{théorème}{théorèmes}
\crefname{proposition}{proposition}{propositions}
\crefname{corollaire}{corollaire}{corollaires}
\crefname{remarque}{remarque}{remarques}
\begin{document}
\chapter{Premier}
\begin{proposition}\label{p1}a\end{proposition}
\begin{theoreme}\label{t1}b\end{theoreme}
\begin{remarque}\label{r1}c\end{remarque}
\begin{corollaire}\label{c1}d\end{corollaire}
\chapter{Second}
\begin{corollaire}\label{c2}e\end{corollaire}
\begin{theoreme}\label{t2}f\end{theoreme}
RENVOIS \ref{p1} \ref{t1} \ref{r1} \ref{c1} \ref{c2} \ref{t2} ; \cref{p1} ; \cref{t1} ; \cref{r1} ; \cref{c1} ; \cref{c2} ; \cref{t2} FIN
\end{document}
"                sans correctif, crochet retiré
%   lualatex -jobname=amsseul  "\def\AMSSEUL{1}% Juge de la troisième passe : document minimal écrit de zéro pour la cause A (thmtools sous LaTeX 2026-06-01).
% Variantes, par une macro posée sur la ligne de commande (deux passes chacune, depuis memoire/essai/juge3/minimal/) :
%   lualatex -jobname=nu       "\input{minimal-juge3}"                                   sans correctif
%   lualatex -jobname=corr     "\def\CORR{1}\input{minimal-juge3}"                       correctif A du préambule
%   lualatex -jobname=crochet  "\def\SANSCROCHET{1}\input{minimal-juge3}"                sans correctif, crochet retiré
%   lualatex -jobname=amsseul  "\def\AMSSEUL{1}\input{minimal-juge3}"                    amsthm seul (témoin)
%   lualatex -jobname=faux     "\def\CORR{1}\def\FAUX{1}\input{minimal-juge3}"           deux vrais défauts : doivent tomber
%   lualatex -jobname=trace    "\def\TRACE{1}\input{minimal-juge3}"                      trace de l'appel
% Les lignes du correctif sont tirées du préambule commité par extraire_correctifs.py, pas retapées.
\documentclass{scrbook}
\ifdefined\SANSCROCHET \RemoveFromHook{package/amsthm/after}[firstaid/aliascounter]\fi
\usepackage{amsthm}
\usepackage{hyperref}
\usepackage[french, nameinlink, noabbrev]{cleveref}
\ifdefined\AMSSEUL
  \newtheorem{theoreme}{Théorème}[chapter]
  \newtheorem{proposition}[theoreme]{Proposition}
  \newtheorem{corollaire}[theoreme]{Corollaire}
  \newtheorem{remarque}{Remarque}[chapter]
\else
  \usepackage{thmtools}
  \ifdefined\CORR \input{correctif-a-pose}\fi
  \declaretheorem[numberwithin=chapter, name=Théorème]{theoreme}
  \ifdefined\TRACE \tracingmacros=2 \tracingonline=0 \fi
  \declaretheorem[sibling=theoreme, name=Proposition]{proposition}
  \ifdefined\TRACE \tracingmacros=0 \fi
  \declaretheorem[sibling=theoreme, name=Corollaire]{corollaire}
  \declaretheorem[numberwithin=chapter, name=Remarque]{remarque}
  \ifdefined\FAUX \newcounteralias{corollaire}{section}\newcounteralias{orphelin}{inexistant}\fi
  \ifdefined\CORR \input{correctif-a-remise}\fi
\fi
\crefname{theoreme}{théorème}{théorèmes}
\crefname{proposition}{proposition}{propositions}
\crefname{corollaire}{corollaire}{corollaires}
\crefname{remarque}{remarque}{remarques}
\begin{document}
\chapter{Premier}
\begin{proposition}\label{p1}a\end{proposition}
\begin{theoreme}\label{t1}b\end{theoreme}
\begin{remarque}\label{r1}c\end{remarque}
\begin{corollaire}\label{c1}d\end{corollaire}
\chapter{Second}
\begin{corollaire}\label{c2}e\end{corollaire}
\begin{theoreme}\label{t2}f\end{theoreme}
RENVOIS \ref{p1} \ref{t1} \ref{r1} \ref{c1} \ref{c2} \ref{t2} ; \cref{p1} ; \cref{t1} ; \cref{r1} ; \cref{c1} ; \cref{c2} ; \cref{t2} FIN
\end{document}
"                    amsthm seul (témoin)
%   lualatex -jobname=faux     "\def\CORR{1}\def\FAUX{1}% Juge de la troisième passe : document minimal écrit de zéro pour la cause A (thmtools sous LaTeX 2026-06-01).
% Variantes, par une macro posée sur la ligne de commande (deux passes chacune, depuis memoire/essai/juge3/minimal/) :
%   lualatex -jobname=nu       "\input{minimal-juge3}"                                   sans correctif
%   lualatex -jobname=corr     "\def\CORR{1}\input{minimal-juge3}"                       correctif A du préambule
%   lualatex -jobname=crochet  "\def\SANSCROCHET{1}\input{minimal-juge3}"                sans correctif, crochet retiré
%   lualatex -jobname=amsseul  "\def\AMSSEUL{1}\input{minimal-juge3}"                    amsthm seul (témoin)
%   lualatex -jobname=faux     "\def\CORR{1}\def\FAUX{1}\input{minimal-juge3}"           deux vrais défauts : doivent tomber
%   lualatex -jobname=trace    "\def\TRACE{1}\input{minimal-juge3}"                      trace de l'appel
% Les lignes du correctif sont tirées du préambule commité par extraire_correctifs.py, pas retapées.
\documentclass{scrbook}
\ifdefined\SANSCROCHET \RemoveFromHook{package/amsthm/after}[firstaid/aliascounter]\fi
\usepackage{amsthm}
\usepackage{hyperref}
\usepackage[french, nameinlink, noabbrev]{cleveref}
\ifdefined\AMSSEUL
  \newtheorem{theoreme}{Théorème}[chapter]
  \newtheorem{proposition}[theoreme]{Proposition}
  \newtheorem{corollaire}[theoreme]{Corollaire}
  \newtheorem{remarque}{Remarque}[chapter]
\else
  \usepackage{thmtools}
  \ifdefined\CORR \input{correctif-a-pose}\fi
  \declaretheorem[numberwithin=chapter, name=Théorème]{theoreme}
  \ifdefined\TRACE \tracingmacros=2 \tracingonline=0 \fi
  \declaretheorem[sibling=theoreme, name=Proposition]{proposition}
  \ifdefined\TRACE \tracingmacros=0 \fi
  \declaretheorem[sibling=theoreme, name=Corollaire]{corollaire}
  \declaretheorem[numberwithin=chapter, name=Remarque]{remarque}
  \ifdefined\FAUX \newcounteralias{corollaire}{section}\newcounteralias{orphelin}{inexistant}\fi
  \ifdefined\CORR \input{correctif-a-remise}\fi
\fi
\crefname{theoreme}{théorème}{théorèmes}
\crefname{proposition}{proposition}{propositions}
\crefname{corollaire}{corollaire}{corollaires}
\crefname{remarque}{remarque}{remarques}
\begin{document}
\chapter{Premier}
\begin{proposition}\label{p1}a\end{proposition}
\begin{theoreme}\label{t1}b\end{theoreme}
\begin{remarque}\label{r1}c\end{remarque}
\begin{corollaire}\label{c1}d\end{corollaire}
\chapter{Second}
\begin{corollaire}\label{c2}e\end{corollaire}
\begin{theoreme}\label{t2}f\end{theoreme}
RENVOIS \ref{p1} \ref{t1} \ref{r1} \ref{c1} \ref{c2} \ref{t2} ; \cref{p1} ; \cref{t1} ; \cref{r1} ; \cref{c1} ; \cref{c2} ; \cref{t2} FIN
\end{document}
"           deux vrais défauts : doivent tomber
%   lualatex -jobname=trace    "\def\TRACE{1}% Juge de la troisième passe : document minimal écrit de zéro pour la cause A (thmtools sous LaTeX 2026-06-01).
% Variantes, par une macro posée sur la ligne de commande (deux passes chacune, depuis memoire/essai/juge3/minimal/) :
%   lualatex -jobname=nu       "\input{minimal-juge3}"                                   sans correctif
%   lualatex -jobname=corr     "\def\CORR{1}\input{minimal-juge3}"                       correctif A du préambule
%   lualatex -jobname=crochet  "\def\SANSCROCHET{1}\input{minimal-juge3}"                sans correctif, crochet retiré
%   lualatex -jobname=amsseul  "\def\AMSSEUL{1}\input{minimal-juge3}"                    amsthm seul (témoin)
%   lualatex -jobname=faux     "\def\CORR{1}\def\FAUX{1}\input{minimal-juge3}"           deux vrais défauts : doivent tomber
%   lualatex -jobname=trace    "\def\TRACE{1}\input{minimal-juge3}"                      trace de l'appel
% Les lignes du correctif sont tirées du préambule commité par extraire_correctifs.py, pas retapées.
\documentclass{scrbook}
\ifdefined\SANSCROCHET \RemoveFromHook{package/amsthm/after}[firstaid/aliascounter]\fi
\usepackage{amsthm}
\usepackage{hyperref}
\usepackage[french, nameinlink, noabbrev]{cleveref}
\ifdefined\AMSSEUL
  \newtheorem{theoreme}{Théorème}[chapter]
  \newtheorem{proposition}[theoreme]{Proposition}
  \newtheorem{corollaire}[theoreme]{Corollaire}
  \newtheorem{remarque}{Remarque}[chapter]
\else
  \usepackage{thmtools}
  \ifdefined\CORR \input{correctif-a-pose}\fi
  \declaretheorem[numberwithin=chapter, name=Théorème]{theoreme}
  \ifdefined\TRACE \tracingmacros=2 \tracingonline=0 \fi
  \declaretheorem[sibling=theoreme, name=Proposition]{proposition}
  \ifdefined\TRACE \tracingmacros=0 \fi
  \declaretheorem[sibling=theoreme, name=Corollaire]{corollaire}
  \declaretheorem[numberwithin=chapter, name=Remarque]{remarque}
  \ifdefined\FAUX \newcounteralias{corollaire}{section}\newcounteralias{orphelin}{inexistant}\fi
  \ifdefined\CORR \input{correctif-a-remise}\fi
\fi
\crefname{theoreme}{théorème}{théorèmes}
\crefname{proposition}{proposition}{propositions}
\crefname{corollaire}{corollaire}{corollaires}
\crefname{remarque}{remarque}{remarques}
\begin{document}
\chapter{Premier}
\begin{proposition}\label{p1}a\end{proposition}
\begin{theoreme}\label{t1}b\end{theoreme}
\begin{remarque}\label{r1}c\end{remarque}
\begin{corollaire}\label{c1}d\end{corollaire}
\chapter{Second}
\begin{corollaire}\label{c2}e\end{corollaire}
\begin{theoreme}\label{t2}f\end{theoreme}
RENVOIS \ref{p1} \ref{t1} \ref{r1} \ref{c1} \ref{c2} \ref{t2} ; \cref{p1} ; \cref{t1} ; \cref{r1} ; \cref{c1} ; \cref{c2} ; \cref{t2} FIN
\end{document}
"                      trace de l'appel
% Les lignes du correctif sont tirées du préambule commité par extraire_correctifs.py, pas retapées.
\documentclass{scrbook}
\ifdefined\SANSCROCHET \RemoveFromHook{package/amsthm/after}[firstaid/aliascounter]\fi
\usepackage{amsthm}
\usepackage{hyperref}
\usepackage[french, nameinlink, noabbrev]{cleveref}
\ifdefined\AMSSEUL
  \newtheorem{theoreme}{Théorème}[chapter]
  \newtheorem{proposition}[theoreme]{Proposition}
  \newtheorem{corollaire}[theoreme]{Corollaire}
  \newtheorem{remarque}{Remarque}[chapter]
\else
  \usepackage{thmtools}
  \ifdefined\CORR \input{correctif-a-pose}\fi
  \declaretheorem[numberwithin=chapter, name=Théorème]{theoreme}
  \ifdefined\TRACE \tracingmacros=2 \tracingonline=0 \fi
  \declaretheorem[sibling=theoreme, name=Proposition]{proposition}
  \ifdefined\TRACE \tracingmacros=0 \fi
  \declaretheorem[sibling=theoreme, name=Corollaire]{corollaire}
  \declaretheorem[numberwithin=chapter, name=Remarque]{remarque}
  \ifdefined\FAUX \newcounteralias{corollaire}{section}\newcounteralias{orphelin}{inexistant}\fi
  \ifdefined\CORR \input{correctif-a-remise}\fi
\fi
\crefname{theoreme}{théorème}{théorèmes}
\crefname{proposition}{proposition}{propositions}
\crefname{corollaire}{corollaire}{corollaires}
\crefname{remarque}{remarque}{remarques}
\begin{document}
\chapter{Premier}
\begin{proposition}\label{p1}a\end{proposition}
\begin{theoreme}\label{t1}b\end{theoreme}
\begin{remarque}\label{r1}c\end{remarque}
\begin{corollaire}\label{c1}d\end{corollaire}
\chapter{Second}
\begin{corollaire}\label{c2}e\end{corollaire}
\begin{theoreme}\label{t2}f\end{theoreme}
RENVOIS \ref{p1} \ref{t1} \ref{r1} \ref{c1} \ref{c2} \ref{t2} ; \cref{p1} ; \cref{t1} ; \cref{r1} ; \cref{c1} ; \cref{c2} ; \cref{t2} FIN
\end{document}
"                      trace de l'appel
% Les lignes du correctif sont tirées du préambule commité par extraire_correctifs.py, pas retapées.
\documentclass{scrbook}
\ifdefined\SANSCROCHET \RemoveFromHook{package/amsthm/after}[firstaid/aliascounter]\fi
\usepackage{amsthm}
\usepackage{hyperref}
\usepackage[french, nameinlink, noabbrev]{cleveref}
\ifdefined\AMSSEUL
  \newtheorem{theoreme}{Théorème}[chapter]
  \newtheorem{proposition}[theoreme]{Proposition}
  \newtheorem{corollaire}[theoreme]{Corollaire}
  \newtheorem{remarque}{Remarque}[chapter]
\else
  \usepackage{thmtools}
  \ifdefined\CORR \input{correctif-a-pose}\fi
  \declaretheorem[numberwithin=chapter, name=Théorème]{theoreme}
  \ifdefined\TRACE \tracingmacros=2 \tracingonline=0 \fi
  \declaretheorem[sibling=theoreme, name=Proposition]{proposition}
  \ifdefined\TRACE \tracingmacros=0 \fi
  \declaretheorem[sibling=theoreme, name=Corollaire]{corollaire}
  \declaretheorem[numberwithin=chapter, name=Remarque]{remarque}
  \ifdefined\FAUX \newcounteralias{corollaire}{section}\newcounteralias{orphelin}{inexistant}\fi
  \ifdefined\CORR \input{correctif-a-remise}\fi
\fi
\crefname{theoreme}{théorème}{théorèmes}
\crefname{proposition}{proposition}{propositions}
\crefname{corollaire}{corollaire}{corollaires}
\crefname{remarque}{remarque}{remarques}
\begin{document}
\chapter{Premier}
\begin{proposition}\label{p1}a\end{proposition}
\begin{theoreme}\label{t1}b\end{theoreme}
\begin{remarque}\label{r1}c\end{remarque}
\begin{corollaire}\label{c1}d\end{corollaire}
\chapter{Second}
\begin{corollaire}\label{c2}e\end{corollaire}
\begin{theoreme}\label{t2}f\end{theoreme}
RENVOIS \ref{p1} \ref{t1} \ref{r1} \ref{c1} \ref{c2} \ref{t2} ; \cref{p1} ; \cref{t1} ; \cref{r1} ; \cref{c1} ; \cref{c2} ; \cref{t2} FIN
\end{document}
"                    amsthm seul (témoin)
%   lualatex -jobname=faux     "\def\CORR{1}\def\FAUX{1}% Juge de la troisième passe : document minimal écrit de zéro pour la cause A (thmtools sous LaTeX 2026-06-01).
% Variantes, par une macro posée sur la ligne de commande (deux passes chacune, depuis memoire/essai/juge3/minimal/) :
%   lualatex -jobname=nu       "% Juge de la troisième passe : document minimal écrit de zéro pour la cause A (thmtools sous LaTeX 2026-06-01).
% Variantes, par une macro posée sur la ligne de commande (deux passes chacune, depuis memoire/essai/juge3/minimal/) :
%   lualatex -jobname=nu       "% Juge de la troisième passe : document minimal écrit de zéro pour la cause A (thmtools sous LaTeX 2026-06-01).
% Variantes, par une macro posée sur la ligne de commande (deux passes chacune, depuis memoire/essai/juge3/minimal/) :
%   lualatex -jobname=nu       "\input{minimal-juge3}"                                   sans correctif
%   lualatex -jobname=corr     "\def\CORR{1}\input{minimal-juge3}"                       correctif A du préambule
%   lualatex -jobname=crochet  "\def\SANSCROCHET{1}\input{minimal-juge3}"                sans correctif, crochet retiré
%   lualatex -jobname=amsseul  "\def\AMSSEUL{1}\input{minimal-juge3}"                    amsthm seul (témoin)
%   lualatex -jobname=faux     "\def\CORR{1}\def\FAUX{1}\input{minimal-juge3}"           deux vrais défauts : doivent tomber
%   lualatex -jobname=trace    "\def\TRACE{1}\input{minimal-juge3}"                      trace de l'appel
% Les lignes du correctif sont tirées du préambule commité par extraire_correctifs.py, pas retapées.
\documentclass{scrbook}
\ifdefined\SANSCROCHET \RemoveFromHook{package/amsthm/after}[firstaid/aliascounter]\fi
\usepackage{amsthm}
\usepackage{hyperref}
\usepackage[french, nameinlink, noabbrev]{cleveref}
\ifdefined\AMSSEUL
  \newtheorem{theoreme}{Théorème}[chapter]
  \newtheorem{proposition}[theoreme]{Proposition}
  \newtheorem{corollaire}[theoreme]{Corollaire}
  \newtheorem{remarque}{Remarque}[chapter]
\else
  \usepackage{thmtools}
  \ifdefined\CORR \input{correctif-a-pose}\fi
  \declaretheorem[numberwithin=chapter, name=Théorème]{theoreme}
  \ifdefined\TRACE \tracingmacros=2 \tracingonline=0 \fi
  \declaretheorem[sibling=theoreme, name=Proposition]{proposition}
  \ifdefined\TRACE \tracingmacros=0 \fi
  \declaretheorem[sibling=theoreme, name=Corollaire]{corollaire}
  \declaretheorem[numberwithin=chapter, name=Remarque]{remarque}
  \ifdefined\FAUX \newcounteralias{corollaire}{section}\newcounteralias{orphelin}{inexistant}\fi
  \ifdefined\CORR \input{correctif-a-remise}\fi
\fi
\crefname{theoreme}{théorème}{théorèmes}
\crefname{proposition}{proposition}{propositions}
\crefname{corollaire}{corollaire}{corollaires}
\crefname{remarque}{remarque}{remarques}
\begin{document}
\chapter{Premier}
\begin{proposition}\label{p1}a\end{proposition}
\begin{theoreme}\label{t1}b\end{theoreme}
\begin{remarque}\label{r1}c\end{remarque}
\begin{corollaire}\label{c1}d\end{corollaire}
\chapter{Second}
\begin{corollaire}\label{c2}e\end{corollaire}
\begin{theoreme}\label{t2}f\end{theoreme}
RENVOIS \ref{p1} \ref{t1} \ref{r1} \ref{c1} \ref{c2} \ref{t2} ; \cref{p1} ; \cref{t1} ; \cref{r1} ; \cref{c1} ; \cref{c2} ; \cref{t2} FIN
\end{document}
"                                   sans correctif
%   lualatex -jobname=corr     "\def\CORR{1}% Juge de la troisième passe : document minimal écrit de zéro pour la cause A (thmtools sous LaTeX 2026-06-01).
% Variantes, par une macro posée sur la ligne de commande (deux passes chacune, depuis memoire/essai/juge3/minimal/) :
%   lualatex -jobname=nu       "\input{minimal-juge3}"                                   sans correctif
%   lualatex -jobname=corr     "\def\CORR{1}\input{minimal-juge3}"                       correctif A du préambule
%   lualatex -jobname=crochet  "\def\SANSCROCHET{1}\input{minimal-juge3}"                sans correctif, crochet retiré
%   lualatex -jobname=amsseul  "\def\AMSSEUL{1}\input{minimal-juge3}"                    amsthm seul (témoin)
%   lualatex -jobname=faux     "\def\CORR{1}\def\FAUX{1}\input{minimal-juge3}"           deux vrais défauts : doivent tomber
%   lualatex -jobname=trace    "\def\TRACE{1}\input{minimal-juge3}"                      trace de l'appel
% Les lignes du correctif sont tirées du préambule commité par extraire_correctifs.py, pas retapées.
\documentclass{scrbook}
\ifdefined\SANSCROCHET \RemoveFromHook{package/amsthm/after}[firstaid/aliascounter]\fi
\usepackage{amsthm}
\usepackage{hyperref}
\usepackage[french, nameinlink, noabbrev]{cleveref}
\ifdefined\AMSSEUL
  \newtheorem{theoreme}{Théorème}[chapter]
  \newtheorem{proposition}[theoreme]{Proposition}
  \newtheorem{corollaire}[theoreme]{Corollaire}
  \newtheorem{remarque}{Remarque}[chapter]
\else
  \usepackage{thmtools}
  \ifdefined\CORR \input{correctif-a-pose}\fi
  \declaretheorem[numberwithin=chapter, name=Théorème]{theoreme}
  \ifdefined\TRACE \tracingmacros=2 \tracingonline=0 \fi
  \declaretheorem[sibling=theoreme, name=Proposition]{proposition}
  \ifdefined\TRACE \tracingmacros=0 \fi
  \declaretheorem[sibling=theoreme, name=Corollaire]{corollaire}
  \declaretheorem[numberwithin=chapter, name=Remarque]{remarque}
  \ifdefined\FAUX \newcounteralias{corollaire}{section}\newcounteralias{orphelin}{inexistant}\fi
  \ifdefined\CORR \input{correctif-a-remise}\fi
\fi
\crefname{theoreme}{théorème}{théorèmes}
\crefname{proposition}{proposition}{propositions}
\crefname{corollaire}{corollaire}{corollaires}
\crefname{remarque}{remarque}{remarques}
\begin{document}
\chapter{Premier}
\begin{proposition}\label{p1}a\end{proposition}
\begin{theoreme}\label{t1}b\end{theoreme}
\begin{remarque}\label{r1}c\end{remarque}
\begin{corollaire}\label{c1}d\end{corollaire}
\chapter{Second}
\begin{corollaire}\label{c2}e\end{corollaire}
\begin{theoreme}\label{t2}f\end{theoreme}
RENVOIS \ref{p1} \ref{t1} \ref{r1} \ref{c1} \ref{c2} \ref{t2} ; \cref{p1} ; \cref{t1} ; \cref{r1} ; \cref{c1} ; \cref{c2} ; \cref{t2} FIN
\end{document}
"                       correctif A du préambule
%   lualatex -jobname=crochet  "\def\SANSCROCHET{1}% Juge de la troisième passe : document minimal écrit de zéro pour la cause A (thmtools sous LaTeX 2026-06-01).
% Variantes, par une macro posée sur la ligne de commande (deux passes chacune, depuis memoire/essai/juge3/minimal/) :
%   lualatex -jobname=nu       "\input{minimal-juge3}"                                   sans correctif
%   lualatex -jobname=corr     "\def\CORR{1}\input{minimal-juge3}"                       correctif A du préambule
%   lualatex -jobname=crochet  "\def\SANSCROCHET{1}\input{minimal-juge3}"                sans correctif, crochet retiré
%   lualatex -jobname=amsseul  "\def\AMSSEUL{1}\input{minimal-juge3}"                    amsthm seul (témoin)
%   lualatex -jobname=faux     "\def\CORR{1}\def\FAUX{1}\input{minimal-juge3}"           deux vrais défauts : doivent tomber
%   lualatex -jobname=trace    "\def\TRACE{1}\input{minimal-juge3}"                      trace de l'appel
% Les lignes du correctif sont tirées du préambule commité par extraire_correctifs.py, pas retapées.
\documentclass{scrbook}
\ifdefined\SANSCROCHET \RemoveFromHook{package/amsthm/after}[firstaid/aliascounter]\fi
\usepackage{amsthm}
\usepackage{hyperref}
\usepackage[french, nameinlink, noabbrev]{cleveref}
\ifdefined\AMSSEUL
  \newtheorem{theoreme}{Théorème}[chapter]
  \newtheorem{proposition}[theoreme]{Proposition}
  \newtheorem{corollaire}[theoreme]{Corollaire}
  \newtheorem{remarque}{Remarque}[chapter]
\else
  \usepackage{thmtools}
  \ifdefined\CORR \input{correctif-a-pose}\fi
  \declaretheorem[numberwithin=chapter, name=Théorème]{theoreme}
  \ifdefined\TRACE \tracingmacros=2 \tracingonline=0 \fi
  \declaretheorem[sibling=theoreme, name=Proposition]{proposition}
  \ifdefined\TRACE \tracingmacros=0 \fi
  \declaretheorem[sibling=theoreme, name=Corollaire]{corollaire}
  \declaretheorem[numberwithin=chapter, name=Remarque]{remarque}
  \ifdefined\FAUX \newcounteralias{corollaire}{section}\newcounteralias{orphelin}{inexistant}\fi
  \ifdefined\CORR \input{correctif-a-remise}\fi
\fi
\crefname{theoreme}{théorème}{théorèmes}
\crefname{proposition}{proposition}{propositions}
\crefname{corollaire}{corollaire}{corollaires}
\crefname{remarque}{remarque}{remarques}
\begin{document}
\chapter{Premier}
\begin{proposition}\label{p1}a\end{proposition}
\begin{theoreme}\label{t1}b\end{theoreme}
\begin{remarque}\label{r1}c\end{remarque}
\begin{corollaire}\label{c1}d\end{corollaire}
\chapter{Second}
\begin{corollaire}\label{c2}e\end{corollaire}
\begin{theoreme}\label{t2}f\end{theoreme}
RENVOIS \ref{p1} \ref{t1} \ref{r1} \ref{c1} \ref{c2} \ref{t2} ; \cref{p1} ; \cref{t1} ; \cref{r1} ; \cref{c1} ; \cref{c2} ; \cref{t2} FIN
\end{document}
"                sans correctif, crochet retiré
%   lualatex -jobname=amsseul  "\def\AMSSEUL{1}% Juge de la troisième passe : document minimal écrit de zéro pour la cause A (thmtools sous LaTeX 2026-06-01).
% Variantes, par une macro posée sur la ligne de commande (deux passes chacune, depuis memoire/essai/juge3/minimal/) :
%   lualatex -jobname=nu       "\input{minimal-juge3}"                                   sans correctif
%   lualatex -jobname=corr     "\def\CORR{1}\input{minimal-juge3}"                       correctif A du préambule
%   lualatex -jobname=crochet  "\def\SANSCROCHET{1}\input{minimal-juge3}"                sans correctif, crochet retiré
%   lualatex -jobname=amsseul  "\def\AMSSEUL{1}\input{minimal-juge3}"                    amsthm seul (témoin)
%   lualatex -jobname=faux     "\def\CORR{1}\def\FAUX{1}\input{minimal-juge3}"           deux vrais défauts : doivent tomber
%   lualatex -jobname=trace    "\def\TRACE{1}\input{minimal-juge3}"                      trace de l'appel
% Les lignes du correctif sont tirées du préambule commité par extraire_correctifs.py, pas retapées.
\documentclass{scrbook}
\ifdefined\SANSCROCHET \RemoveFromHook{package/amsthm/after}[firstaid/aliascounter]\fi
\usepackage{amsthm}
\usepackage{hyperref}
\usepackage[french, nameinlink, noabbrev]{cleveref}
\ifdefined\AMSSEUL
  \newtheorem{theoreme}{Théorème}[chapter]
  \newtheorem{proposition}[theoreme]{Proposition}
  \newtheorem{corollaire}[theoreme]{Corollaire}
  \newtheorem{remarque}{Remarque}[chapter]
\else
  \usepackage{thmtools}
  \ifdefined\CORR \input{correctif-a-pose}\fi
  \declaretheorem[numberwithin=chapter, name=Théorème]{theoreme}
  \ifdefined\TRACE \tracingmacros=2 \tracingonline=0 \fi
  \declaretheorem[sibling=theoreme, name=Proposition]{proposition}
  \ifdefined\TRACE \tracingmacros=0 \fi
  \declaretheorem[sibling=theoreme, name=Corollaire]{corollaire}
  \declaretheorem[numberwithin=chapter, name=Remarque]{remarque}
  \ifdefined\FAUX \newcounteralias{corollaire}{section}\newcounteralias{orphelin}{inexistant}\fi
  \ifdefined\CORR \input{correctif-a-remise}\fi
\fi
\crefname{theoreme}{théorème}{théorèmes}
\crefname{proposition}{proposition}{propositions}
\crefname{corollaire}{corollaire}{corollaires}
\crefname{remarque}{remarque}{remarques}
\begin{document}
\chapter{Premier}
\begin{proposition}\label{p1}a\end{proposition}
\begin{theoreme}\label{t1}b\end{theoreme}
\begin{remarque}\label{r1}c\end{remarque}
\begin{corollaire}\label{c1}d\end{corollaire}
\chapter{Second}
\begin{corollaire}\label{c2}e\end{corollaire}
\begin{theoreme}\label{t2}f\end{theoreme}
RENVOIS \ref{p1} \ref{t1} \ref{r1} \ref{c1} \ref{c2} \ref{t2} ; \cref{p1} ; \cref{t1} ; \cref{r1} ; \cref{c1} ; \cref{c2} ; \cref{t2} FIN
\end{document}
"                    amsthm seul (témoin)
%   lualatex -jobname=faux     "\def\CORR{1}\def\FAUX{1}% Juge de la troisième passe : document minimal écrit de zéro pour la cause A (thmtools sous LaTeX 2026-06-01).
% Variantes, par une macro posée sur la ligne de commande (deux passes chacune, depuis memoire/essai/juge3/minimal/) :
%   lualatex -jobname=nu       "\input{minimal-juge3}"                                   sans correctif
%   lualatex -jobname=corr     "\def\CORR{1}\input{minimal-juge3}"                       correctif A du préambule
%   lualatex -jobname=crochet  "\def\SANSCROCHET{1}\input{minimal-juge3}"                sans correctif, crochet retiré
%   lualatex -jobname=amsseul  "\def\AMSSEUL{1}\input{minimal-juge3}"                    amsthm seul (témoin)
%   lualatex -jobname=faux     "\def\CORR{1}\def\FAUX{1}\input{minimal-juge3}"           deux vrais défauts : doivent tomber
%   lualatex -jobname=trace    "\def\TRACE{1}\input{minimal-juge3}"                      trace de l'appel
% Les lignes du correctif sont tirées du préambule commité par extraire_correctifs.py, pas retapées.
\documentclass{scrbook}
\ifdefined\SANSCROCHET \RemoveFromHook{package/amsthm/after}[firstaid/aliascounter]\fi
\usepackage{amsthm}
\usepackage{hyperref}
\usepackage[french, nameinlink, noabbrev]{cleveref}
\ifdefined\AMSSEUL
  \newtheorem{theoreme}{Théorème}[chapter]
  \newtheorem{proposition}[theoreme]{Proposition}
  \newtheorem{corollaire}[theoreme]{Corollaire}
  \newtheorem{remarque}{Remarque}[chapter]
\else
  \usepackage{thmtools}
  \ifdefined\CORR \input{correctif-a-pose}\fi
  \declaretheorem[numberwithin=chapter, name=Théorème]{theoreme}
  \ifdefined\TRACE \tracingmacros=2 \tracingonline=0 \fi
  \declaretheorem[sibling=theoreme, name=Proposition]{proposition}
  \ifdefined\TRACE \tracingmacros=0 \fi
  \declaretheorem[sibling=theoreme, name=Corollaire]{corollaire}
  \declaretheorem[numberwithin=chapter, name=Remarque]{remarque}
  \ifdefined\FAUX \newcounteralias{corollaire}{section}\newcounteralias{orphelin}{inexistant}\fi
  \ifdefined\CORR \input{correctif-a-remise}\fi
\fi
\crefname{theoreme}{théorème}{théorèmes}
\crefname{proposition}{proposition}{propositions}
\crefname{corollaire}{corollaire}{corollaires}
\crefname{remarque}{remarque}{remarques}
\begin{document}
\chapter{Premier}
\begin{proposition}\label{p1}a\end{proposition}
\begin{theoreme}\label{t1}b\end{theoreme}
\begin{remarque}\label{r1}c\end{remarque}
\begin{corollaire}\label{c1}d\end{corollaire}
\chapter{Second}
\begin{corollaire}\label{c2}e\end{corollaire}
\begin{theoreme}\label{t2}f\end{theoreme}
RENVOIS \ref{p1} \ref{t1} \ref{r1} \ref{c1} \ref{c2} \ref{t2} ; \cref{p1} ; \cref{t1} ; \cref{r1} ; \cref{c1} ; \cref{c2} ; \cref{t2} FIN
\end{document}
"           deux vrais défauts : doivent tomber
%   lualatex -jobname=trace    "\def\TRACE{1}% Juge de la troisième passe : document minimal écrit de zéro pour la cause A (thmtools sous LaTeX 2026-06-01).
% Variantes, par une macro posée sur la ligne de commande (deux passes chacune, depuis memoire/essai/juge3/minimal/) :
%   lualatex -jobname=nu       "\input{minimal-juge3}"                                   sans correctif
%   lualatex -jobname=corr     "\def\CORR{1}\input{minimal-juge3}"                       correctif A du préambule
%   lualatex -jobname=crochet  "\def\SANSCROCHET{1}\input{minimal-juge3}"                sans correctif, crochet retiré
%   lualatex -jobname=amsseul  "\def\AMSSEUL{1}\input{minimal-juge3}"                    amsthm seul (témoin)
%   lualatex -jobname=faux     "\def\CORR{1}\def\FAUX{1}\input{minimal-juge3}"           deux vrais défauts : doivent tomber
%   lualatex -jobname=trace    "\def\TRACE{1}\input{minimal-juge3}"                      trace de l'appel
% Les lignes du correctif sont tirées du préambule commité par extraire_correctifs.py, pas retapées.
\documentclass{scrbook}
\ifdefined\SANSCROCHET \RemoveFromHook{package/amsthm/after}[firstaid/aliascounter]\fi
\usepackage{amsthm}
\usepackage{hyperref}
\usepackage[french, nameinlink, noabbrev]{cleveref}
\ifdefined\AMSSEUL
  \newtheorem{theoreme}{Théorème}[chapter]
  \newtheorem{proposition}[theoreme]{Proposition}
  \newtheorem{corollaire}[theoreme]{Corollaire}
  \newtheorem{remarque}{Remarque}[chapter]
\else
  \usepackage{thmtools}
  \ifdefined\CORR \input{correctif-a-pose}\fi
  \declaretheorem[numberwithin=chapter, name=Théorème]{theoreme}
  \ifdefined\TRACE \tracingmacros=2 \tracingonline=0 \fi
  \declaretheorem[sibling=theoreme, name=Proposition]{proposition}
  \ifdefined\TRACE \tracingmacros=0 \fi
  \declaretheorem[sibling=theoreme, name=Corollaire]{corollaire}
  \declaretheorem[numberwithin=chapter, name=Remarque]{remarque}
  \ifdefined\FAUX \newcounteralias{corollaire}{section}\newcounteralias{orphelin}{inexistant}\fi
  \ifdefined\CORR \input{correctif-a-remise}\fi
\fi
\crefname{theoreme}{théorème}{théorèmes}
\crefname{proposition}{proposition}{propositions}
\crefname{corollaire}{corollaire}{corollaires}
\crefname{remarque}{remarque}{remarques}
\begin{document}
\chapter{Premier}
\begin{proposition}\label{p1}a\end{proposition}
\begin{theoreme}\label{t1}b\end{theoreme}
\begin{remarque}\label{r1}c\end{remarque}
\begin{corollaire}\label{c1}d\end{corollaire}
\chapter{Second}
\begin{corollaire}\label{c2}e\end{corollaire}
\begin{theoreme}\label{t2}f\end{theoreme}
RENVOIS \ref{p1} \ref{t1} \ref{r1} \ref{c1} \ref{c2} \ref{t2} ; \cref{p1} ; \cref{t1} ; \cref{r1} ; \cref{c1} ; \cref{c2} ; \cref{t2} FIN
\end{document}
"                      trace de l'appel
% Les lignes du correctif sont tirées du préambule commité par extraire_correctifs.py, pas retapées.
\documentclass{scrbook}
\ifdefined\SANSCROCHET \RemoveFromHook{package/amsthm/after}[firstaid/aliascounter]\fi
\usepackage{amsthm}
\usepackage{hyperref}
\usepackage[french, nameinlink, noabbrev]{cleveref}
\ifdefined\AMSSEUL
  \newtheorem{theoreme}{Théorème}[chapter]
  \newtheorem{proposition}[theoreme]{Proposition}
  \newtheorem{corollaire}[theoreme]{Corollaire}
  \newtheorem{remarque}{Remarque}[chapter]
\else
  \usepackage{thmtools}
  \ifdefined\CORR \input{correctif-a-pose}\fi
  \declaretheorem[numberwithin=chapter, name=Théorème]{theoreme}
  \ifdefined\TRACE \tracingmacros=2 \tracingonline=0 \fi
  \declaretheorem[sibling=theoreme, name=Proposition]{proposition}
  \ifdefined\TRACE \tracingmacros=0 \fi
  \declaretheorem[sibling=theoreme, name=Corollaire]{corollaire}
  \declaretheorem[numberwithin=chapter, name=Remarque]{remarque}
  \ifdefined\FAUX \newcounteralias{corollaire}{section}\newcounteralias{orphelin}{inexistant}\fi
  \ifdefined\CORR \input{correctif-a-remise}\fi
\fi
\crefname{theoreme}{théorème}{théorèmes}
\crefname{proposition}{proposition}{propositions}
\crefname{corollaire}{corollaire}{corollaires}
\crefname{remarque}{remarque}{remarques}
\begin{document}
\chapter{Premier}
\begin{proposition}\label{p1}a\end{proposition}
\begin{theoreme}\label{t1}b\end{theoreme}
\begin{remarque}\label{r1}c\end{remarque}
\begin{corollaire}\label{c1}d\end{corollaire}
\chapter{Second}
\begin{corollaire}\label{c2}e\end{corollaire}
\begin{theoreme}\label{t2}f\end{theoreme}
RENVOIS \ref{p1} \ref{t1} \ref{r1} \ref{c1} \ref{c2} \ref{t2} ; \cref{p1} ; \cref{t1} ; \cref{r1} ; \cref{c1} ; \cref{c2} ; \cref{t2} FIN
\end{document}
"                                   sans correctif
%   lualatex -jobname=corr     "\def\CORR{1}% Juge de la troisième passe : document minimal écrit de zéro pour la cause A (thmtools sous LaTeX 2026-06-01).
% Variantes, par une macro posée sur la ligne de commande (deux passes chacune, depuis memoire/essai/juge3/minimal/) :
%   lualatex -jobname=nu       "% Juge de la troisième passe : document minimal écrit de zéro pour la cause A (thmtools sous LaTeX 2026-06-01).
% Variantes, par une macro posée sur la ligne de commande (deux passes chacune, depuis memoire/essai/juge3/minimal/) :
%   lualatex -jobname=nu       "\input{minimal-juge3}"                                   sans correctif
%   lualatex -jobname=corr     "\def\CORR{1}\input{minimal-juge3}"                       correctif A du préambule
%   lualatex -jobname=crochet  "\def\SANSCROCHET{1}\input{minimal-juge3}"                sans correctif, crochet retiré
%   lualatex -jobname=amsseul  "\def\AMSSEUL{1}\input{minimal-juge3}"                    amsthm seul (témoin)
%   lualatex -jobname=faux     "\def\CORR{1}\def\FAUX{1}\input{minimal-juge3}"           deux vrais défauts : doivent tomber
%   lualatex -jobname=trace    "\def\TRACE{1}\input{minimal-juge3}"                      trace de l'appel
% Les lignes du correctif sont tirées du préambule commité par extraire_correctifs.py, pas retapées.
\documentclass{scrbook}
\ifdefined\SANSCROCHET \RemoveFromHook{package/amsthm/after}[firstaid/aliascounter]\fi
\usepackage{amsthm}
\usepackage{hyperref}
\usepackage[french, nameinlink, noabbrev]{cleveref}
\ifdefined\AMSSEUL
  \newtheorem{theoreme}{Théorème}[chapter]
  \newtheorem{proposition}[theoreme]{Proposition}
  \newtheorem{corollaire}[theoreme]{Corollaire}
  \newtheorem{remarque}{Remarque}[chapter]
\else
  \usepackage{thmtools}
  \ifdefined\CORR \input{correctif-a-pose}\fi
  \declaretheorem[numberwithin=chapter, name=Théorème]{theoreme}
  \ifdefined\TRACE \tracingmacros=2 \tracingonline=0 \fi
  \declaretheorem[sibling=theoreme, name=Proposition]{proposition}
  \ifdefined\TRACE \tracingmacros=0 \fi
  \declaretheorem[sibling=theoreme, name=Corollaire]{corollaire}
  \declaretheorem[numberwithin=chapter, name=Remarque]{remarque}
  \ifdefined\FAUX \newcounteralias{corollaire}{section}\newcounteralias{orphelin}{inexistant}\fi
  \ifdefined\CORR \input{correctif-a-remise}\fi
\fi
\crefname{theoreme}{théorème}{théorèmes}
\crefname{proposition}{proposition}{propositions}
\crefname{corollaire}{corollaire}{corollaires}
\crefname{remarque}{remarque}{remarques}
\begin{document}
\chapter{Premier}
\begin{proposition}\label{p1}a\end{proposition}
\begin{theoreme}\label{t1}b\end{theoreme}
\begin{remarque}\label{r1}c\end{remarque}
\begin{corollaire}\label{c1}d\end{corollaire}
\chapter{Second}
\begin{corollaire}\label{c2}e\end{corollaire}
\begin{theoreme}\label{t2}f\end{theoreme}
RENVOIS \ref{p1} \ref{t1} \ref{r1} \ref{c1} \ref{c2} \ref{t2} ; \cref{p1} ; \cref{t1} ; \cref{r1} ; \cref{c1} ; \cref{c2} ; \cref{t2} FIN
\end{document}
"                                   sans correctif
%   lualatex -jobname=corr     "\def\CORR{1}% Juge de la troisième passe : document minimal écrit de zéro pour la cause A (thmtools sous LaTeX 2026-06-01).
% Variantes, par une macro posée sur la ligne de commande (deux passes chacune, depuis memoire/essai/juge3/minimal/) :
%   lualatex -jobname=nu       "\input{minimal-juge3}"                                   sans correctif
%   lualatex -jobname=corr     "\def\CORR{1}\input{minimal-juge3}"                       correctif A du préambule
%   lualatex -jobname=crochet  "\def\SANSCROCHET{1}\input{minimal-juge3}"                sans correctif, crochet retiré
%   lualatex -jobname=amsseul  "\def\AMSSEUL{1}\input{minimal-juge3}"                    amsthm seul (témoin)
%   lualatex -jobname=faux     "\def\CORR{1}\def\FAUX{1}\input{minimal-juge3}"           deux vrais défauts : doivent tomber
%   lualatex -jobname=trace    "\def\TRACE{1}\input{minimal-juge3}"                      trace de l'appel
% Les lignes du correctif sont tirées du préambule commité par extraire_correctifs.py, pas retapées.
\documentclass{scrbook}
\ifdefined\SANSCROCHET \RemoveFromHook{package/amsthm/after}[firstaid/aliascounter]\fi
\usepackage{amsthm}
\usepackage{hyperref}
\usepackage[french, nameinlink, noabbrev]{cleveref}
\ifdefined\AMSSEUL
  \newtheorem{theoreme}{Théorème}[chapter]
  \newtheorem{proposition}[theoreme]{Proposition}
  \newtheorem{corollaire}[theoreme]{Corollaire}
  \newtheorem{remarque}{Remarque}[chapter]
\else
  \usepackage{thmtools}
  \ifdefined\CORR \input{correctif-a-pose}\fi
  \declaretheorem[numberwithin=chapter, name=Théorème]{theoreme}
  \ifdefined\TRACE \tracingmacros=2 \tracingonline=0 \fi
  \declaretheorem[sibling=theoreme, name=Proposition]{proposition}
  \ifdefined\TRACE \tracingmacros=0 \fi
  \declaretheorem[sibling=theoreme, name=Corollaire]{corollaire}
  \declaretheorem[numberwithin=chapter, name=Remarque]{remarque}
  \ifdefined\FAUX \newcounteralias{corollaire}{section}\newcounteralias{orphelin}{inexistant}\fi
  \ifdefined\CORR \input{correctif-a-remise}\fi
\fi
\crefname{theoreme}{théorème}{théorèmes}
\crefname{proposition}{proposition}{propositions}
\crefname{corollaire}{corollaire}{corollaires}
\crefname{remarque}{remarque}{remarques}
\begin{document}
\chapter{Premier}
\begin{proposition}\label{p1}a\end{proposition}
\begin{theoreme}\label{t1}b\end{theoreme}
\begin{remarque}\label{r1}c\end{remarque}
\begin{corollaire}\label{c1}d\end{corollaire}
\chapter{Second}
\begin{corollaire}\label{c2}e\end{corollaire}
\begin{theoreme}\label{t2}f\end{theoreme}
RENVOIS \ref{p1} \ref{t1} \ref{r1} \ref{c1} \ref{c2} \ref{t2} ; \cref{p1} ; \cref{t1} ; \cref{r1} ; \cref{c1} ; \cref{c2} ; \cref{t2} FIN
\end{document}
"                       correctif A du préambule
%   lualatex -jobname=crochet  "\def\SANSCROCHET{1}% Juge de la troisième passe : document minimal écrit de zéro pour la cause A (thmtools sous LaTeX 2026-06-01).
% Variantes, par une macro posée sur la ligne de commande (deux passes chacune, depuis memoire/essai/juge3/minimal/) :
%   lualatex -jobname=nu       "\input{minimal-juge3}"                                   sans correctif
%   lualatex -jobname=corr     "\def\CORR{1}\input{minimal-juge3}"                       correctif A du préambule
%   lualatex -jobname=crochet  "\def\SANSCROCHET{1}\input{minimal-juge3}"                sans correctif, crochet retiré
%   lualatex -jobname=amsseul  "\def\AMSSEUL{1}\input{minimal-juge3}"                    amsthm seul (témoin)
%   lualatex -jobname=faux     "\def\CORR{1}\def\FAUX{1}\input{minimal-juge3}"           deux vrais défauts : doivent tomber
%   lualatex -jobname=trace    "\def\TRACE{1}\input{minimal-juge3}"                      trace de l'appel
% Les lignes du correctif sont tirées du préambule commité par extraire_correctifs.py, pas retapées.
\documentclass{scrbook}
\ifdefined\SANSCROCHET \RemoveFromHook{package/amsthm/after}[firstaid/aliascounter]\fi
\usepackage{amsthm}
\usepackage{hyperref}
\usepackage[french, nameinlink, noabbrev]{cleveref}
\ifdefined\AMSSEUL
  \newtheorem{theoreme}{Théorème}[chapter]
  \newtheorem{proposition}[theoreme]{Proposition}
  \newtheorem{corollaire}[theoreme]{Corollaire}
  \newtheorem{remarque}{Remarque}[chapter]
\else
  \usepackage{thmtools}
  \ifdefined\CORR \input{correctif-a-pose}\fi
  \declaretheorem[numberwithin=chapter, name=Théorème]{theoreme}
  \ifdefined\TRACE \tracingmacros=2 \tracingonline=0 \fi
  \declaretheorem[sibling=theoreme, name=Proposition]{proposition}
  \ifdefined\TRACE \tracingmacros=0 \fi
  \declaretheorem[sibling=theoreme, name=Corollaire]{corollaire}
  \declaretheorem[numberwithin=chapter, name=Remarque]{remarque}
  \ifdefined\FAUX \newcounteralias{corollaire}{section}\newcounteralias{orphelin}{inexistant}\fi
  \ifdefined\CORR \input{correctif-a-remise}\fi
\fi
\crefname{theoreme}{théorème}{théorèmes}
\crefname{proposition}{proposition}{propositions}
\crefname{corollaire}{corollaire}{corollaires}
\crefname{remarque}{remarque}{remarques}
\begin{document}
\chapter{Premier}
\begin{proposition}\label{p1}a\end{proposition}
\begin{theoreme}\label{t1}b\end{theoreme}
\begin{remarque}\label{r1}c\end{remarque}
\begin{corollaire}\label{c1}d\end{corollaire}
\chapter{Second}
\begin{corollaire}\label{c2}e\end{corollaire}
\begin{theoreme}\label{t2}f\end{theoreme}
RENVOIS \ref{p1} \ref{t1} \ref{r1} \ref{c1} \ref{c2} \ref{t2} ; \cref{p1} ; \cref{t1} ; \cref{r1} ; \cref{c1} ; \cref{c2} ; \cref{t2} FIN
\end{document}
"                sans correctif, crochet retiré
%   lualatex -jobname=amsseul  "\def\AMSSEUL{1}% Juge de la troisième passe : document minimal écrit de zéro pour la cause A (thmtools sous LaTeX 2026-06-01).
% Variantes, par une macro posée sur la ligne de commande (deux passes chacune, depuis memoire/essai/juge3/minimal/) :
%   lualatex -jobname=nu       "\input{minimal-juge3}"                                   sans correctif
%   lualatex -jobname=corr     "\def\CORR{1}\input{minimal-juge3}"                       correctif A du préambule
%   lualatex -jobname=crochet  "\def\SANSCROCHET{1}\input{minimal-juge3}"                sans correctif, crochet retiré
%   lualatex -jobname=amsseul  "\def\AMSSEUL{1}\input{minimal-juge3}"                    amsthm seul (témoin)
%   lualatex -jobname=faux     "\def\CORR{1}\def\FAUX{1}\input{minimal-juge3}"           deux vrais défauts : doivent tomber
%   lualatex -jobname=trace    "\def\TRACE{1}\input{minimal-juge3}"                      trace de l'appel
% Les lignes du correctif sont tirées du préambule commité par extraire_correctifs.py, pas retapées.
\documentclass{scrbook}
\ifdefined\SANSCROCHET \RemoveFromHook{package/amsthm/after}[firstaid/aliascounter]\fi
\usepackage{amsthm}
\usepackage{hyperref}
\usepackage[french, nameinlink, noabbrev]{cleveref}
\ifdefined\AMSSEUL
  \newtheorem{theoreme}{Théorème}[chapter]
  \newtheorem{proposition}[theoreme]{Proposition}
  \newtheorem{corollaire}[theoreme]{Corollaire}
  \newtheorem{remarque}{Remarque}[chapter]
\else
  \usepackage{thmtools}
  \ifdefined\CORR \input{correctif-a-pose}\fi
  \declaretheorem[numberwithin=chapter, name=Théorème]{theoreme}
  \ifdefined\TRACE \tracingmacros=2 \tracingonline=0 \fi
  \declaretheorem[sibling=theoreme, name=Proposition]{proposition}
  \ifdefined\TRACE \tracingmacros=0 \fi
  \declaretheorem[sibling=theoreme, name=Corollaire]{corollaire}
  \declaretheorem[numberwithin=chapter, name=Remarque]{remarque}
  \ifdefined\FAUX \newcounteralias{corollaire}{section}\newcounteralias{orphelin}{inexistant}\fi
  \ifdefined\CORR \input{correctif-a-remise}\fi
\fi
\crefname{theoreme}{théorème}{théorèmes}
\crefname{proposition}{proposition}{propositions}
\crefname{corollaire}{corollaire}{corollaires}
\crefname{remarque}{remarque}{remarques}
\begin{document}
\chapter{Premier}
\begin{proposition}\label{p1}a\end{proposition}
\begin{theoreme}\label{t1}b\end{theoreme}
\begin{remarque}\label{r1}c\end{remarque}
\begin{corollaire}\label{c1}d\end{corollaire}
\chapter{Second}
\begin{corollaire}\label{c2}e\end{corollaire}
\begin{theoreme}\label{t2}f\end{theoreme}
RENVOIS \ref{p1} \ref{t1} \ref{r1} \ref{c1} \ref{c2} \ref{t2} ; \cref{p1} ; \cref{t1} ; \cref{r1} ; \cref{c1} ; \cref{c2} ; \cref{t2} FIN
\end{document}
"                    amsthm seul (témoin)
%   lualatex -jobname=faux     "\def\CORR{1}\def\FAUX{1}% Juge de la troisième passe : document minimal écrit de zéro pour la cause A (thmtools sous LaTeX 2026-06-01).
% Variantes, par une macro posée sur la ligne de commande (deux passes chacune, depuis memoire/essai/juge3/minimal/) :
%   lualatex -jobname=nu       "\input{minimal-juge3}"                                   sans correctif
%   lualatex -jobname=corr     "\def\CORR{1}\input{minimal-juge3}"                       correctif A du préambule
%   lualatex -jobname=crochet  "\def\SANSCROCHET{1}\input{minimal-juge3}"                sans correctif, crochet retiré
%   lualatex -jobname=amsseul  "\def\AMSSEUL{1}\input{minimal-juge3}"                    amsthm seul (témoin)
%   lualatex -jobname=faux     "\def\CORR{1}\def\FAUX{1}\input{minimal-juge3}"           deux vrais défauts : doivent tomber
%   lualatex -jobname=trace    "\def\TRACE{1}\input{minimal-juge3}"                      trace de l'appel
% Les lignes du correctif sont tirées du préambule commité par extraire_correctifs.py, pas retapées.
\documentclass{scrbook}
\ifdefined\SANSCROCHET \RemoveFromHook{package/amsthm/after}[firstaid/aliascounter]\fi
\usepackage{amsthm}
\usepackage{hyperref}
\usepackage[french, nameinlink, noabbrev]{cleveref}
\ifdefined\AMSSEUL
  \newtheorem{theoreme}{Théorème}[chapter]
  \newtheorem{proposition}[theoreme]{Proposition}
  \newtheorem{corollaire}[theoreme]{Corollaire}
  \newtheorem{remarque}{Remarque}[chapter]
\else
  \usepackage{thmtools}
  \ifdefined\CORR \input{correctif-a-pose}\fi
  \declaretheorem[numberwithin=chapter, name=Théorème]{theoreme}
  \ifdefined\TRACE \tracingmacros=2 \tracingonline=0 \fi
  \declaretheorem[sibling=theoreme, name=Proposition]{proposition}
  \ifdefined\TRACE \tracingmacros=0 \fi
  \declaretheorem[sibling=theoreme, name=Corollaire]{corollaire}
  \declaretheorem[numberwithin=chapter, name=Remarque]{remarque}
  \ifdefined\FAUX \newcounteralias{corollaire}{section}\newcounteralias{orphelin}{inexistant}\fi
  \ifdefined\CORR \input{correctif-a-remise}\fi
\fi
\crefname{theoreme}{théorème}{théorèmes}
\crefname{proposition}{proposition}{propositions}
\crefname{corollaire}{corollaire}{corollaires}
\crefname{remarque}{remarque}{remarques}
\begin{document}
\chapter{Premier}
\begin{proposition}\label{p1}a\end{proposition}
\begin{theoreme}\label{t1}b\end{theoreme}
\begin{remarque}\label{r1}c\end{remarque}
\begin{corollaire}\label{c1}d\end{corollaire}
\chapter{Second}
\begin{corollaire}\label{c2}e\end{corollaire}
\begin{theoreme}\label{t2}f\end{theoreme}
RENVOIS \ref{p1} \ref{t1} \ref{r1} \ref{c1} \ref{c2} \ref{t2} ; \cref{p1} ; \cref{t1} ; \cref{r1} ; \cref{c1} ; \cref{c2} ; \cref{t2} FIN
\end{document}
"           deux vrais défauts : doivent tomber
%   lualatex -jobname=trace    "\def\TRACE{1}% Juge de la troisième passe : document minimal écrit de zéro pour la cause A (thmtools sous LaTeX 2026-06-01).
% Variantes, par une macro posée sur la ligne de commande (deux passes chacune, depuis memoire/essai/juge3/minimal/) :
%   lualatex -jobname=nu       "\input{minimal-juge3}"                                   sans correctif
%   lualatex -jobname=corr     "\def\CORR{1}\input{minimal-juge3}"                       correctif A du préambule
%   lualatex -jobname=crochet  "\def\SANSCROCHET{1}\input{minimal-juge3}"                sans correctif, crochet retiré
%   lualatex -jobname=amsseul  "\def\AMSSEUL{1}\input{minimal-juge3}"                    amsthm seul (témoin)
%   lualatex -jobname=faux     "\def\CORR{1}\def\FAUX{1}\input{minimal-juge3}"           deux vrais défauts : doivent tomber
%   lualatex -jobname=trace    "\def\TRACE{1}\input{minimal-juge3}"                      trace de l'appel
% Les lignes du correctif sont tirées du préambule commité par extraire_correctifs.py, pas retapées.
\documentclass{scrbook}
\ifdefined\SANSCROCHET \RemoveFromHook{package/amsthm/after}[firstaid/aliascounter]\fi
\usepackage{amsthm}
\usepackage{hyperref}
\usepackage[french, nameinlink, noabbrev]{cleveref}
\ifdefined\AMSSEUL
  \newtheorem{theoreme}{Théorème}[chapter]
  \newtheorem{proposition}[theoreme]{Proposition}
  \newtheorem{corollaire}[theoreme]{Corollaire}
  \newtheorem{remarque}{Remarque}[chapter]
\else
  \usepackage{thmtools}
  \ifdefined\CORR \input{correctif-a-pose}\fi
  \declaretheorem[numberwithin=chapter, name=Théorème]{theoreme}
  \ifdefined\TRACE \tracingmacros=2 \tracingonline=0 \fi
  \declaretheorem[sibling=theoreme, name=Proposition]{proposition}
  \ifdefined\TRACE \tracingmacros=0 \fi
  \declaretheorem[sibling=theoreme, name=Corollaire]{corollaire}
  \declaretheorem[numberwithin=chapter, name=Remarque]{remarque}
  \ifdefined\FAUX \newcounteralias{corollaire}{section}\newcounteralias{orphelin}{inexistant}\fi
  \ifdefined\CORR \input{correctif-a-remise}\fi
\fi
\crefname{theoreme}{théorème}{théorèmes}
\crefname{proposition}{proposition}{propositions}
\crefname{corollaire}{corollaire}{corollaires}
\crefname{remarque}{remarque}{remarques}
\begin{document}
\chapter{Premier}
\begin{proposition}\label{p1}a\end{proposition}
\begin{theoreme}\label{t1}b\end{theoreme}
\begin{remarque}\label{r1}c\end{remarque}
\begin{corollaire}\label{c1}d\end{corollaire}
\chapter{Second}
\begin{corollaire}\label{c2}e\end{corollaire}
\begin{theoreme}\label{t2}f\end{theoreme}
RENVOIS \ref{p1} \ref{t1} \ref{r1} \ref{c1} \ref{c2} \ref{t2} ; \cref{p1} ; \cref{t1} ; \cref{r1} ; \cref{c1} ; \cref{c2} ; \cref{t2} FIN
\end{document}
"                      trace de l'appel
% Les lignes du correctif sont tirées du préambule commité par extraire_correctifs.py, pas retapées.
\documentclass{scrbook}
\ifdefined\SANSCROCHET \RemoveFromHook{package/amsthm/after}[firstaid/aliascounter]\fi
\usepackage{amsthm}
\usepackage{hyperref}
\usepackage[french, nameinlink, noabbrev]{cleveref}
\ifdefined\AMSSEUL
  \newtheorem{theoreme}{Théorème}[chapter]
  \newtheorem{proposition}[theoreme]{Proposition}
  \newtheorem{corollaire}[theoreme]{Corollaire}
  \newtheorem{remarque}{Remarque}[chapter]
\else
  \usepackage{thmtools}
  \ifdefined\CORR \input{correctif-a-pose}\fi
  \declaretheorem[numberwithin=chapter, name=Théorème]{theoreme}
  \ifdefined\TRACE \tracingmacros=2 \tracingonline=0 \fi
  \declaretheorem[sibling=theoreme, name=Proposition]{proposition}
  \ifdefined\TRACE \tracingmacros=0 \fi
  \declaretheorem[sibling=theoreme, name=Corollaire]{corollaire}
  \declaretheorem[numberwithin=chapter, name=Remarque]{remarque}
  \ifdefined\FAUX \newcounteralias{corollaire}{section}\newcounteralias{orphelin}{inexistant}\fi
  \ifdefined\CORR \input{correctif-a-remise}\fi
\fi
\crefname{theoreme}{théorème}{théorèmes}
\crefname{proposition}{proposition}{propositions}
\crefname{corollaire}{corollaire}{corollaires}
\crefname{remarque}{remarque}{remarques}
\begin{document}
\chapter{Premier}
\begin{proposition}\label{p1}a\end{proposition}
\begin{theoreme}\label{t1}b\end{theoreme}
\begin{remarque}\label{r1}c\end{remarque}
\begin{corollaire}\label{c1}d\end{corollaire}
\chapter{Second}
\begin{corollaire}\label{c2}e\end{corollaire}
\begin{theoreme}\label{t2}f\end{theoreme}
RENVOIS \ref{p1} \ref{t1} \ref{r1} \ref{c1} \ref{c2} \ref{t2} ; \cref{p1} ; \cref{t1} ; \cref{r1} ; \cref{c1} ; \cref{c2} ; \cref{t2} FIN
\end{document}
"                       correctif A du préambule
%   lualatex -jobname=crochet  "\def\SANSCROCHET{1}% Juge de la troisième passe : document minimal écrit de zéro pour la cause A (thmtools sous LaTeX 2026-06-01).
% Variantes, par une macro posée sur la ligne de commande (deux passes chacune, depuis memoire/essai/juge3/minimal/) :
%   lualatex -jobname=nu       "% Juge de la troisième passe : document minimal écrit de zéro pour la cause A (thmtools sous LaTeX 2026-06-01).
% Variantes, par une macro posée sur la ligne de commande (deux passes chacune, depuis memoire/essai/juge3/minimal/) :
%   lualatex -jobname=nu       "\input{minimal-juge3}"                                   sans correctif
%   lualatex -jobname=corr     "\def\CORR{1}\input{minimal-juge3}"                       correctif A du préambule
%   lualatex -jobname=crochet  "\def\SANSCROCHET{1}\input{minimal-juge3}"                sans correctif, crochet retiré
%   lualatex -jobname=amsseul  "\def\AMSSEUL{1}\input{minimal-juge3}"                    amsthm seul (témoin)
%   lualatex -jobname=faux     "\def\CORR{1}\def\FAUX{1}\input{minimal-juge3}"           deux vrais défauts : doivent tomber
%   lualatex -jobname=trace    "\def\TRACE{1}\input{minimal-juge3}"                      trace de l'appel
% Les lignes du correctif sont tirées du préambule commité par extraire_correctifs.py, pas retapées.
\documentclass{scrbook}
\ifdefined\SANSCROCHET \RemoveFromHook{package/amsthm/after}[firstaid/aliascounter]\fi
\usepackage{amsthm}
\usepackage{hyperref}
\usepackage[french, nameinlink, noabbrev]{cleveref}
\ifdefined\AMSSEUL
  \newtheorem{theoreme}{Théorème}[chapter]
  \newtheorem{proposition}[theoreme]{Proposition}
  \newtheorem{corollaire}[theoreme]{Corollaire}
  \newtheorem{remarque}{Remarque}[chapter]
\else
  \usepackage{thmtools}
  \ifdefined\CORR \input{correctif-a-pose}\fi
  \declaretheorem[numberwithin=chapter, name=Théorème]{theoreme}
  \ifdefined\TRACE \tracingmacros=2 \tracingonline=0 \fi
  \declaretheorem[sibling=theoreme, name=Proposition]{proposition}
  \ifdefined\TRACE \tracingmacros=0 \fi
  \declaretheorem[sibling=theoreme, name=Corollaire]{corollaire}
  \declaretheorem[numberwithin=chapter, name=Remarque]{remarque}
  \ifdefined\FAUX \newcounteralias{corollaire}{section}\newcounteralias{orphelin}{inexistant}\fi
  \ifdefined\CORR \input{correctif-a-remise}\fi
\fi
\crefname{theoreme}{théorème}{théorèmes}
\crefname{proposition}{proposition}{propositions}
\crefname{corollaire}{corollaire}{corollaires}
\crefname{remarque}{remarque}{remarques}
\begin{document}
\chapter{Premier}
\begin{proposition}\label{p1}a\end{proposition}
\begin{theoreme}\label{t1}b\end{theoreme}
\begin{remarque}\label{r1}c\end{remarque}
\begin{corollaire}\label{c1}d\end{corollaire}
\chapter{Second}
\begin{corollaire}\label{c2}e\end{corollaire}
\begin{theoreme}\label{t2}f\end{theoreme}
RENVOIS \ref{p1} \ref{t1} \ref{r1} \ref{c1} \ref{c2} \ref{t2} ; \cref{p1} ; \cref{t1} ; \cref{r1} ; \cref{c1} ; \cref{c2} ; \cref{t2} FIN
\end{document}
"                                   sans correctif
%   lualatex -jobname=corr     "\def\CORR{1}% Juge de la troisième passe : document minimal écrit de zéro pour la cause A (thmtools sous LaTeX 2026-06-01).
% Variantes, par une macro posée sur la ligne de commande (deux passes chacune, depuis memoire/essai/juge3/minimal/) :
%   lualatex -jobname=nu       "\input{minimal-juge3}"                                   sans correctif
%   lualatex -jobname=corr     "\def\CORR{1}\input{minimal-juge3}"                       correctif A du préambule
%   lualatex -jobname=crochet  "\def\SANSCROCHET{1}\input{minimal-juge3}"                sans correctif, crochet retiré
%   lualatex -jobname=amsseul  "\def\AMSSEUL{1}\input{minimal-juge3}"                    amsthm seul (témoin)
%   lualatex -jobname=faux     "\def\CORR{1}\def\FAUX{1}\input{minimal-juge3}"           deux vrais défauts : doivent tomber
%   lualatex -jobname=trace    "\def\TRACE{1}\input{minimal-juge3}"                      trace de l'appel
% Les lignes du correctif sont tirées du préambule commité par extraire_correctifs.py, pas retapées.
\documentclass{scrbook}
\ifdefined\SANSCROCHET \RemoveFromHook{package/amsthm/after}[firstaid/aliascounter]\fi
\usepackage{amsthm}
\usepackage{hyperref}
\usepackage[french, nameinlink, noabbrev]{cleveref}
\ifdefined\AMSSEUL
  \newtheorem{theoreme}{Théorème}[chapter]
  \newtheorem{proposition}[theoreme]{Proposition}
  \newtheorem{corollaire}[theoreme]{Corollaire}
  \newtheorem{remarque}{Remarque}[chapter]
\else
  \usepackage{thmtools}
  \ifdefined\CORR \input{correctif-a-pose}\fi
  \declaretheorem[numberwithin=chapter, name=Théorème]{theoreme}
  \ifdefined\TRACE \tracingmacros=2 \tracingonline=0 \fi
  \declaretheorem[sibling=theoreme, name=Proposition]{proposition}
  \ifdefined\TRACE \tracingmacros=0 \fi
  \declaretheorem[sibling=theoreme, name=Corollaire]{corollaire}
  \declaretheorem[numberwithin=chapter, name=Remarque]{remarque}
  \ifdefined\FAUX \newcounteralias{corollaire}{section}\newcounteralias{orphelin}{inexistant}\fi
  \ifdefined\CORR \input{correctif-a-remise}\fi
\fi
\crefname{theoreme}{théorème}{théorèmes}
\crefname{proposition}{proposition}{propositions}
\crefname{corollaire}{corollaire}{corollaires}
\crefname{remarque}{remarque}{remarques}
\begin{document}
\chapter{Premier}
\begin{proposition}\label{p1}a\end{proposition}
\begin{theoreme}\label{t1}b\end{theoreme}
\begin{remarque}\label{r1}c\end{remarque}
\begin{corollaire}\label{c1}d\end{corollaire}
\chapter{Second}
\begin{corollaire}\label{c2}e\end{corollaire}
\begin{theoreme}\label{t2}f\end{theoreme}
RENVOIS \ref{p1} \ref{t1} \ref{r1} \ref{c1} \ref{c2} \ref{t2} ; \cref{p1} ; \cref{t1} ; \cref{r1} ; \cref{c1} ; \cref{c2} ; \cref{t2} FIN
\end{document}
"                       correctif A du préambule
%   lualatex -jobname=crochet  "\def\SANSCROCHET{1}% Juge de la troisième passe : document minimal écrit de zéro pour la cause A (thmtools sous LaTeX 2026-06-01).
% Variantes, par une macro posée sur la ligne de commande (deux passes chacune, depuis memoire/essai/juge3/minimal/) :
%   lualatex -jobname=nu       "\input{minimal-juge3}"                                   sans correctif
%   lualatex -jobname=corr     "\def\CORR{1}\input{minimal-juge3}"                       correctif A du préambule
%   lualatex -jobname=crochet  "\def\SANSCROCHET{1}\input{minimal-juge3}"                sans correctif, crochet retiré
%   lualatex -jobname=amsseul  "\def\AMSSEUL{1}\input{minimal-juge3}"                    amsthm seul (témoin)
%   lualatex -jobname=faux     "\def\CORR{1}\def\FAUX{1}\input{minimal-juge3}"           deux vrais défauts : doivent tomber
%   lualatex -jobname=trace    "\def\TRACE{1}\input{minimal-juge3}"                      trace de l'appel
% Les lignes du correctif sont tirées du préambule commité par extraire_correctifs.py, pas retapées.
\documentclass{scrbook}
\ifdefined\SANSCROCHET \RemoveFromHook{package/amsthm/after}[firstaid/aliascounter]\fi
\usepackage{amsthm}
\usepackage{hyperref}
\usepackage[french, nameinlink, noabbrev]{cleveref}
\ifdefined\AMSSEUL
  \newtheorem{theoreme}{Théorème}[chapter]
  \newtheorem{proposition}[theoreme]{Proposition}
  \newtheorem{corollaire}[theoreme]{Corollaire}
  \newtheorem{remarque}{Remarque}[chapter]
\else
  \usepackage{thmtools}
  \ifdefined\CORR \input{correctif-a-pose}\fi
  \declaretheorem[numberwithin=chapter, name=Théorème]{theoreme}
  \ifdefined\TRACE \tracingmacros=2 \tracingonline=0 \fi
  \declaretheorem[sibling=theoreme, name=Proposition]{proposition}
  \ifdefined\TRACE \tracingmacros=0 \fi
  \declaretheorem[sibling=theoreme, name=Corollaire]{corollaire}
  \declaretheorem[numberwithin=chapter, name=Remarque]{remarque}
  \ifdefined\FAUX \newcounteralias{corollaire}{section}\newcounteralias{orphelin}{inexistant}\fi
  \ifdefined\CORR \input{correctif-a-remise}\fi
\fi
\crefname{theoreme}{théorème}{théorèmes}
\crefname{proposition}{proposition}{propositions}
\crefname{corollaire}{corollaire}{corollaires}
\crefname{remarque}{remarque}{remarques}
\begin{document}
\chapter{Premier}
\begin{proposition}\label{p1}a\end{proposition}
\begin{theoreme}\label{t1}b\end{theoreme}
\begin{remarque}\label{r1}c\end{remarque}
\begin{corollaire}\label{c1}d\end{corollaire}
\chapter{Second}
\begin{corollaire}\label{c2}e\end{corollaire}
\begin{theoreme}\label{t2}f\end{theoreme}
RENVOIS \ref{p1} \ref{t1} \ref{r1} \ref{c1} \ref{c2} \ref{t2} ; \cref{p1} ; \cref{t1} ; \cref{r1} ; \cref{c1} ; \cref{c2} ; \cref{t2} FIN
\end{document}
"                sans correctif, crochet retiré
%   lualatex -jobname=amsseul  "\def\AMSSEUL{1}% Juge de la troisième passe : document minimal écrit de zéro pour la cause A (thmtools sous LaTeX 2026-06-01).
% Variantes, par une macro posée sur la ligne de commande (deux passes chacune, depuis memoire/essai/juge3/minimal/) :
%   lualatex -jobname=nu       "\input{minimal-juge3}"                                   sans correctif
%   lualatex -jobname=corr     "\def\CORR{1}\input{minimal-juge3}"                       correctif A du préambule
%   lualatex -jobname=crochet  "\def\SANSCROCHET{1}\input{minimal-juge3}"                sans correctif, crochet retiré
%   lualatex -jobname=amsseul  "\def\AMSSEUL{1}\input{minimal-juge3}"                    amsthm seul (témoin)
%   lualatex -jobname=faux     "\def\CORR{1}\def\FAUX{1}\input{minimal-juge3}"           deux vrais défauts : doivent tomber
%   lualatex -jobname=trace    "\def\TRACE{1}\input{minimal-juge3}"                      trace de l'appel
% Les lignes du correctif sont tirées du préambule commité par extraire_correctifs.py, pas retapées.
\documentclass{scrbook}
\ifdefined\SANSCROCHET \RemoveFromHook{package/amsthm/after}[firstaid/aliascounter]\fi
\usepackage{amsthm}
\usepackage{hyperref}
\usepackage[french, nameinlink, noabbrev]{cleveref}
\ifdefined\AMSSEUL
  \newtheorem{theoreme}{Théorème}[chapter]
  \newtheorem{proposition}[theoreme]{Proposition}
  \newtheorem{corollaire}[theoreme]{Corollaire}
  \newtheorem{remarque}{Remarque}[chapter]
\else
  \usepackage{thmtools}
  \ifdefined\CORR \input{correctif-a-pose}\fi
  \declaretheorem[numberwithin=chapter, name=Théorème]{theoreme}
  \ifdefined\TRACE \tracingmacros=2 \tracingonline=0 \fi
  \declaretheorem[sibling=theoreme, name=Proposition]{proposition}
  \ifdefined\TRACE \tracingmacros=0 \fi
  \declaretheorem[sibling=theoreme, name=Corollaire]{corollaire}
  \declaretheorem[numberwithin=chapter, name=Remarque]{remarque}
  \ifdefined\FAUX \newcounteralias{corollaire}{section}\newcounteralias{orphelin}{inexistant}\fi
  \ifdefined\CORR \input{correctif-a-remise}\fi
\fi
\crefname{theoreme}{théorème}{théorèmes}
\crefname{proposition}{proposition}{propositions}
\crefname{corollaire}{corollaire}{corollaires}
\crefname{remarque}{remarque}{remarques}
\begin{document}
\chapter{Premier}
\begin{proposition}\label{p1}a\end{proposition}
\begin{theoreme}\label{t1}b\end{theoreme}
\begin{remarque}\label{r1}c\end{remarque}
\begin{corollaire}\label{c1}d\end{corollaire}
\chapter{Second}
\begin{corollaire}\label{c2}e\end{corollaire}
\begin{theoreme}\label{t2}f\end{theoreme}
RENVOIS \ref{p1} \ref{t1} \ref{r1} \ref{c1} \ref{c2} \ref{t2} ; \cref{p1} ; \cref{t1} ; \cref{r1} ; \cref{c1} ; \cref{c2} ; \cref{t2} FIN
\end{document}
"                    amsthm seul (témoin)
%   lualatex -jobname=faux     "\def\CORR{1}\def\FAUX{1}% Juge de la troisième passe : document minimal écrit de zéro pour la cause A (thmtools sous LaTeX 2026-06-01).
% Variantes, par une macro posée sur la ligne de commande (deux passes chacune, depuis memoire/essai/juge3/minimal/) :
%   lualatex -jobname=nu       "\input{minimal-juge3}"                                   sans correctif
%   lualatex -jobname=corr     "\def\CORR{1}\input{minimal-juge3}"                       correctif A du préambule
%   lualatex -jobname=crochet  "\def\SANSCROCHET{1}\input{minimal-juge3}"                sans correctif, crochet retiré
%   lualatex -jobname=amsseul  "\def\AMSSEUL{1}\input{minimal-juge3}"                    amsthm seul (témoin)
%   lualatex -jobname=faux     "\def\CORR{1}\def\FAUX{1}\input{minimal-juge3}"           deux vrais défauts : doivent tomber
%   lualatex -jobname=trace    "\def\TRACE{1}\input{minimal-juge3}"                      trace de l'appel
% Les lignes du correctif sont tirées du préambule commité par extraire_correctifs.py, pas retapées.
\documentclass{scrbook}
\ifdefined\SANSCROCHET \RemoveFromHook{package/amsthm/after}[firstaid/aliascounter]\fi
\usepackage{amsthm}
\usepackage{hyperref}
\usepackage[french, nameinlink, noabbrev]{cleveref}
\ifdefined\AMSSEUL
  \newtheorem{theoreme}{Théorème}[chapter]
  \newtheorem{proposition}[theoreme]{Proposition}
  \newtheorem{corollaire}[theoreme]{Corollaire}
  \newtheorem{remarque}{Remarque}[chapter]
\else
  \usepackage{thmtools}
  \ifdefined\CORR \input{correctif-a-pose}\fi
  \declaretheorem[numberwithin=chapter, name=Théorème]{theoreme}
  \ifdefined\TRACE \tracingmacros=2 \tracingonline=0 \fi
  \declaretheorem[sibling=theoreme, name=Proposition]{proposition}
  \ifdefined\TRACE \tracingmacros=0 \fi
  \declaretheorem[sibling=theoreme, name=Corollaire]{corollaire}
  \declaretheorem[numberwithin=chapter, name=Remarque]{remarque}
  \ifdefined\FAUX \newcounteralias{corollaire}{section}\newcounteralias{orphelin}{inexistant}\fi
  \ifdefined\CORR \input{correctif-a-remise}\fi
\fi
\crefname{theoreme}{théorème}{théorèmes}
\crefname{proposition}{proposition}{propositions}
\crefname{corollaire}{corollaire}{corollaires}
\crefname{remarque}{remarque}{remarques}
\begin{document}
\chapter{Premier}
\begin{proposition}\label{p1}a\end{proposition}
\begin{theoreme}\label{t1}b\end{theoreme}
\begin{remarque}\label{r1}c\end{remarque}
\begin{corollaire}\label{c1}d\end{corollaire}
\chapter{Second}
\begin{corollaire}\label{c2}e\end{corollaire}
\begin{theoreme}\label{t2}f\end{theoreme}
RENVOIS \ref{p1} \ref{t1} \ref{r1} \ref{c1} \ref{c2} \ref{t2} ; \cref{p1} ; \cref{t1} ; \cref{r1} ; \cref{c1} ; \cref{c2} ; \cref{t2} FIN
\end{document}
"           deux vrais défauts : doivent tomber
%   lualatex -jobname=trace    "\def\TRACE{1}% Juge de la troisième passe : document minimal écrit de zéro pour la cause A (thmtools sous LaTeX 2026-06-01).
% Variantes, par une macro posée sur la ligne de commande (deux passes chacune, depuis memoire/essai/juge3/minimal/) :
%   lualatex -jobname=nu       "\input{minimal-juge3}"                                   sans correctif
%   lualatex -jobname=corr     "\def\CORR{1}\input{minimal-juge3}"                       correctif A du préambule
%   lualatex -jobname=crochet  "\def\SANSCROCHET{1}\input{minimal-juge3}"                sans correctif, crochet retiré
%   lualatex -jobname=amsseul  "\def\AMSSEUL{1}\input{minimal-juge3}"                    amsthm seul (témoin)
%   lualatex -jobname=faux     "\def\CORR{1}\def\FAUX{1}\input{minimal-juge3}"           deux vrais défauts : doivent tomber
%   lualatex -jobname=trace    "\def\TRACE{1}\input{minimal-juge3}"                      trace de l'appel
% Les lignes du correctif sont tirées du préambule commité par extraire_correctifs.py, pas retapées.
\documentclass{scrbook}
\ifdefined\SANSCROCHET \RemoveFromHook{package/amsthm/after}[firstaid/aliascounter]\fi
\usepackage{amsthm}
\usepackage{hyperref}
\usepackage[french, nameinlink, noabbrev]{cleveref}
\ifdefined\AMSSEUL
  \newtheorem{theoreme}{Théorème}[chapter]
  \newtheorem{proposition}[theoreme]{Proposition}
  \newtheorem{corollaire}[theoreme]{Corollaire}
  \newtheorem{remarque}{Remarque}[chapter]
\else
  \usepackage{thmtools}
  \ifdefined\CORR \input{correctif-a-pose}\fi
  \declaretheorem[numberwithin=chapter, name=Théorème]{theoreme}
  \ifdefined\TRACE \tracingmacros=2 \tracingonline=0 \fi
  \declaretheorem[sibling=theoreme, name=Proposition]{proposition}
  \ifdefined\TRACE \tracingmacros=0 \fi
  \declaretheorem[sibling=theoreme, name=Corollaire]{corollaire}
  \declaretheorem[numberwithin=chapter, name=Remarque]{remarque}
  \ifdefined\FAUX \newcounteralias{corollaire}{section}\newcounteralias{orphelin}{inexistant}\fi
  \ifdefined\CORR \input{correctif-a-remise}\fi
\fi
\crefname{theoreme}{théorème}{théorèmes}
\crefname{proposition}{proposition}{propositions}
\crefname{corollaire}{corollaire}{corollaires}
\crefname{remarque}{remarque}{remarques}
\begin{document}
\chapter{Premier}
\begin{proposition}\label{p1}a\end{proposition}
\begin{theoreme}\label{t1}b\end{theoreme}
\begin{remarque}\label{r1}c\end{remarque}
\begin{corollaire}\label{c1}d\end{corollaire}
\chapter{Second}
\begin{corollaire}\label{c2}e\end{corollaire}
\begin{theoreme}\label{t2}f\end{theoreme}
RENVOIS \ref{p1} \ref{t1} \ref{r1} \ref{c1} \ref{c2} \ref{t2} ; \cref{p1} ; \cref{t1} ; \cref{r1} ; \cref{c1} ; \cref{c2} ; \cref{t2} FIN
\end{document}
"                      trace de l'appel
% Les lignes du correctif sont tirées du préambule commité par extraire_correctifs.py, pas retapées.
\documentclass{scrbook}
\ifdefined\SANSCROCHET \RemoveFromHook{package/amsthm/after}[firstaid/aliascounter]\fi
\usepackage{amsthm}
\usepackage{hyperref}
\usepackage[french, nameinlink, noabbrev]{cleveref}
\ifdefined\AMSSEUL
  \newtheorem{theoreme}{Théorème}[chapter]
  \newtheorem{proposition}[theoreme]{Proposition}
  \newtheorem{corollaire}[theoreme]{Corollaire}
  \newtheorem{remarque}{Remarque}[chapter]
\else
  \usepackage{thmtools}
  \ifdefined\CORR \input{correctif-a-pose}\fi
  \declaretheorem[numberwithin=chapter, name=Théorème]{theoreme}
  \ifdefined\TRACE \tracingmacros=2 \tracingonline=0 \fi
  \declaretheorem[sibling=theoreme, name=Proposition]{proposition}
  \ifdefined\TRACE \tracingmacros=0 \fi
  \declaretheorem[sibling=theoreme, name=Corollaire]{corollaire}
  \declaretheorem[numberwithin=chapter, name=Remarque]{remarque}
  \ifdefined\FAUX \newcounteralias{corollaire}{section}\newcounteralias{orphelin}{inexistant}\fi
  \ifdefined\CORR \input{correctif-a-remise}\fi
\fi
\crefname{theoreme}{théorème}{théorèmes}
\crefname{proposition}{proposition}{propositions}
\crefname{corollaire}{corollaire}{corollaires}
\crefname{remarque}{remarque}{remarques}
\begin{document}
\chapter{Premier}
\begin{proposition}\label{p1}a\end{proposition}
\begin{theoreme}\label{t1}b\end{theoreme}
\begin{remarque}\label{r1}c\end{remarque}
\begin{corollaire}\label{c1}d\end{corollaire}
\chapter{Second}
\begin{corollaire}\label{c2}e\end{corollaire}
\begin{theoreme}\label{t2}f\end{theoreme}
RENVOIS \ref{p1} \ref{t1} \ref{r1} \ref{c1} \ref{c2} \ref{t2} ; \cref{p1} ; \cref{t1} ; \cref{r1} ; \cref{c1} ; \cref{c2} ; \cref{t2} FIN
\end{document}
"                sans correctif, crochet retiré
%   lualatex -jobname=amsseul  "\def\AMSSEUL{1}% Juge de la troisième passe : document minimal écrit de zéro pour la cause A (thmtools sous LaTeX 2026-06-01).
% Variantes, par une macro posée sur la ligne de commande (deux passes chacune, depuis memoire/essai/juge3/minimal/) :
%   lualatex -jobname=nu       "% Juge de la troisième passe : document minimal écrit de zéro pour la cause A (thmtools sous LaTeX 2026-06-01).
% Variantes, par une macro posée sur la ligne de commande (deux passes chacune, depuis memoire/essai/juge3/minimal/) :
%   lualatex -jobname=nu       "\input{minimal-juge3}"                                   sans correctif
%   lualatex -jobname=corr     "\def\CORR{1}\input{minimal-juge3}"                       correctif A du préambule
%   lualatex -jobname=crochet  "\def\SANSCROCHET{1}\input{minimal-juge3}"                sans correctif, crochet retiré
%   lualatex -jobname=amsseul  "\def\AMSSEUL{1}\input{minimal-juge3}"                    amsthm seul (témoin)
%   lualatex -jobname=faux     "\def\CORR{1}\def\FAUX{1}\input{minimal-juge3}"           deux vrais défauts : doivent tomber
%   lualatex -jobname=trace    "\def\TRACE{1}\input{minimal-juge3}"                      trace de l'appel
% Les lignes du correctif sont tirées du préambule commité par extraire_correctifs.py, pas retapées.
\documentclass{scrbook}
\ifdefined\SANSCROCHET \RemoveFromHook{package/amsthm/after}[firstaid/aliascounter]\fi
\usepackage{amsthm}
\usepackage{hyperref}
\usepackage[french, nameinlink, noabbrev]{cleveref}
\ifdefined\AMSSEUL
  \newtheorem{theoreme}{Théorème}[chapter]
  \newtheorem{proposition}[theoreme]{Proposition}
  \newtheorem{corollaire}[theoreme]{Corollaire}
  \newtheorem{remarque}{Remarque}[chapter]
\else
  \usepackage{thmtools}
  \ifdefined\CORR \input{correctif-a-pose}\fi
  \declaretheorem[numberwithin=chapter, name=Théorème]{theoreme}
  \ifdefined\TRACE \tracingmacros=2 \tracingonline=0 \fi
  \declaretheorem[sibling=theoreme, name=Proposition]{proposition}
  \ifdefined\TRACE \tracingmacros=0 \fi
  \declaretheorem[sibling=theoreme, name=Corollaire]{corollaire}
  \declaretheorem[numberwithin=chapter, name=Remarque]{remarque}
  \ifdefined\FAUX \newcounteralias{corollaire}{section}\newcounteralias{orphelin}{inexistant}\fi
  \ifdefined\CORR \input{correctif-a-remise}\fi
\fi
\crefname{theoreme}{théorème}{théorèmes}
\crefname{proposition}{proposition}{propositions}
\crefname{corollaire}{corollaire}{corollaires}
\crefname{remarque}{remarque}{remarques}
\begin{document}
\chapter{Premier}
\begin{proposition}\label{p1}a\end{proposition}
\begin{theoreme}\label{t1}b\end{theoreme}
\begin{remarque}\label{r1}c\end{remarque}
\begin{corollaire}\label{c1}d\end{corollaire}
\chapter{Second}
\begin{corollaire}\label{c2}e\end{corollaire}
\begin{theoreme}\label{t2}f\end{theoreme}
RENVOIS \ref{p1} \ref{t1} \ref{r1} \ref{c1} \ref{c2} \ref{t2} ; \cref{p1} ; \cref{t1} ; \cref{r1} ; \cref{c1} ; \cref{c2} ; \cref{t2} FIN
\end{document}
"                                   sans correctif
%   lualatex -jobname=corr     "\def\CORR{1}% Juge de la troisième passe : document minimal écrit de zéro pour la cause A (thmtools sous LaTeX 2026-06-01).
% Variantes, par une macro posée sur la ligne de commande (deux passes chacune, depuis memoire/essai/juge3/minimal/) :
%   lualatex -jobname=nu       "\input{minimal-juge3}"                                   sans correctif
%   lualatex -jobname=corr     "\def\CORR{1}\input{minimal-juge3}"                       correctif A du préambule
%   lualatex -jobname=crochet  "\def\SANSCROCHET{1}\input{minimal-juge3}"                sans correctif, crochet retiré
%   lualatex -jobname=amsseul  "\def\AMSSEUL{1}\input{minimal-juge3}"                    amsthm seul (témoin)
%   lualatex -jobname=faux     "\def\CORR{1}\def\FAUX{1}\input{minimal-juge3}"           deux vrais défauts : doivent tomber
%   lualatex -jobname=trace    "\def\TRACE{1}\input{minimal-juge3}"                      trace de l'appel
% Les lignes du correctif sont tirées du préambule commité par extraire_correctifs.py, pas retapées.
\documentclass{scrbook}
\ifdefined\SANSCROCHET \RemoveFromHook{package/amsthm/after}[firstaid/aliascounter]\fi
\usepackage{amsthm}
\usepackage{hyperref}
\usepackage[french, nameinlink, noabbrev]{cleveref}
\ifdefined\AMSSEUL
  \newtheorem{theoreme}{Théorème}[chapter]
  \newtheorem{proposition}[theoreme]{Proposition}
  \newtheorem{corollaire}[theoreme]{Corollaire}
  \newtheorem{remarque}{Remarque}[chapter]
\else
  \usepackage{thmtools}
  \ifdefined\CORR \input{correctif-a-pose}\fi
  \declaretheorem[numberwithin=chapter, name=Théorème]{theoreme}
  \ifdefined\TRACE \tracingmacros=2 \tracingonline=0 \fi
  \declaretheorem[sibling=theoreme, name=Proposition]{proposition}
  \ifdefined\TRACE \tracingmacros=0 \fi
  \declaretheorem[sibling=theoreme, name=Corollaire]{corollaire}
  \declaretheorem[numberwithin=chapter, name=Remarque]{remarque}
  \ifdefined\FAUX \newcounteralias{corollaire}{section}\newcounteralias{orphelin}{inexistant}\fi
  \ifdefined\CORR \input{correctif-a-remise}\fi
\fi
\crefname{theoreme}{théorème}{théorèmes}
\crefname{proposition}{proposition}{propositions}
\crefname{corollaire}{corollaire}{corollaires}
\crefname{remarque}{remarque}{remarques}
\begin{document}
\chapter{Premier}
\begin{proposition}\label{p1}a\end{proposition}
\begin{theoreme}\label{t1}b\end{theoreme}
\begin{remarque}\label{r1}c\end{remarque}
\begin{corollaire}\label{c1}d\end{corollaire}
\chapter{Second}
\begin{corollaire}\label{c2}e\end{corollaire}
\begin{theoreme}\label{t2}f\end{theoreme}
RENVOIS \ref{p1} \ref{t1} \ref{r1} \ref{c1} \ref{c2} \ref{t2} ; \cref{p1} ; \cref{t1} ; \cref{r1} ; \cref{c1} ; \cref{c2} ; \cref{t2} FIN
\end{document}
"                       correctif A du préambule
%   lualatex -jobname=crochet  "\def\SANSCROCHET{1}% Juge de la troisième passe : document minimal écrit de zéro pour la cause A (thmtools sous LaTeX 2026-06-01).
% Variantes, par une macro posée sur la ligne de commande (deux passes chacune, depuis memoire/essai/juge3/minimal/) :
%   lualatex -jobname=nu       "\input{minimal-juge3}"                                   sans correctif
%   lualatex -jobname=corr     "\def\CORR{1}\input{minimal-juge3}"                       correctif A du préambule
%   lualatex -jobname=crochet  "\def\SANSCROCHET{1}\input{minimal-juge3}"                sans correctif, crochet retiré
%   lualatex -jobname=amsseul  "\def\AMSSEUL{1}\input{minimal-juge3}"                    amsthm seul (témoin)
%   lualatex -jobname=faux     "\def\CORR{1}\def\FAUX{1}\input{minimal-juge3}"           deux vrais défauts : doivent tomber
%   lualatex -jobname=trace    "\def\TRACE{1}\input{minimal-juge3}"                      trace de l'appel
% Les lignes du correctif sont tirées du préambule commité par extraire_correctifs.py, pas retapées.
\documentclass{scrbook}
\ifdefined\SANSCROCHET \RemoveFromHook{package/amsthm/after}[firstaid/aliascounter]\fi
\usepackage{amsthm}
\usepackage{hyperref}
\usepackage[french, nameinlink, noabbrev]{cleveref}
\ifdefined\AMSSEUL
  \newtheorem{theoreme}{Théorème}[chapter]
  \newtheorem{proposition}[theoreme]{Proposition}
  \newtheorem{corollaire}[theoreme]{Corollaire}
  \newtheorem{remarque}{Remarque}[chapter]
\else
  \usepackage{thmtools}
  \ifdefined\CORR \input{correctif-a-pose}\fi
  \declaretheorem[numberwithin=chapter, name=Théorème]{theoreme}
  \ifdefined\TRACE \tracingmacros=2 \tracingonline=0 \fi
  \declaretheorem[sibling=theoreme, name=Proposition]{proposition}
  \ifdefined\TRACE \tracingmacros=0 \fi
  \declaretheorem[sibling=theoreme, name=Corollaire]{corollaire}
  \declaretheorem[numberwithin=chapter, name=Remarque]{remarque}
  \ifdefined\FAUX \newcounteralias{corollaire}{section}\newcounteralias{orphelin}{inexistant}\fi
  \ifdefined\CORR \input{correctif-a-remise}\fi
\fi
\crefname{theoreme}{théorème}{théorèmes}
\crefname{proposition}{proposition}{propositions}
\crefname{corollaire}{corollaire}{corollaires}
\crefname{remarque}{remarque}{remarques}
\begin{document}
\chapter{Premier}
\begin{proposition}\label{p1}a\end{proposition}
\begin{theoreme}\label{t1}b\end{theoreme}
\begin{remarque}\label{r1}c\end{remarque}
\begin{corollaire}\label{c1}d\end{corollaire}
\chapter{Second}
\begin{corollaire}\label{c2}e\end{corollaire}
\begin{theoreme}\label{t2}f\end{theoreme}
RENVOIS \ref{p1} \ref{t1} \ref{r1} \ref{c1} \ref{c2} \ref{t2} ; \cref{p1} ; \cref{t1} ; \cref{r1} ; \cref{c1} ; \cref{c2} ; \cref{t2} FIN
\end{document}
"                sans correctif, crochet retiré
%   lualatex -jobname=amsseul  "\def\AMSSEUL{1}% Juge de la troisième passe : document minimal écrit de zéro pour la cause A (thmtools sous LaTeX 2026-06-01).
% Variantes, par une macro posée sur la ligne de commande (deux passes chacune, depuis memoire/essai/juge3/minimal/) :
%   lualatex -jobname=nu       "\input{minimal-juge3}"                                   sans correctif
%   lualatex -jobname=corr     "\def\CORR{1}\input{minimal-juge3}"                       correctif A du préambule
%   lualatex -jobname=crochet  "\def\SANSCROCHET{1}\input{minimal-juge3}"                sans correctif, crochet retiré
%   lualatex -jobname=amsseul  "\def\AMSSEUL{1}\input{minimal-juge3}"                    amsthm seul (témoin)
%   lualatex -jobname=faux     "\def\CORR{1}\def\FAUX{1}\input{minimal-juge3}"           deux vrais défauts : doivent tomber
%   lualatex -jobname=trace    "\def\TRACE{1}\input{minimal-juge3}"                      trace de l'appel
% Les lignes du correctif sont tirées du préambule commité par extraire_correctifs.py, pas retapées.
\documentclass{scrbook}
\ifdefined\SANSCROCHET \RemoveFromHook{package/amsthm/after}[firstaid/aliascounter]\fi
\usepackage{amsthm}
\usepackage{hyperref}
\usepackage[french, nameinlink, noabbrev]{cleveref}
\ifdefined\AMSSEUL
  \newtheorem{theoreme}{Théorème}[chapter]
  \newtheorem{proposition}[theoreme]{Proposition}
  \newtheorem{corollaire}[theoreme]{Corollaire}
  \newtheorem{remarque}{Remarque}[chapter]
\else
  \usepackage{thmtools}
  \ifdefined\CORR \input{correctif-a-pose}\fi
  \declaretheorem[numberwithin=chapter, name=Théorème]{theoreme}
  \ifdefined\TRACE \tracingmacros=2 \tracingonline=0 \fi
  \declaretheorem[sibling=theoreme, name=Proposition]{proposition}
  \ifdefined\TRACE \tracingmacros=0 \fi
  \declaretheorem[sibling=theoreme, name=Corollaire]{corollaire}
  \declaretheorem[numberwithin=chapter, name=Remarque]{remarque}
  \ifdefined\FAUX \newcounteralias{corollaire}{section}\newcounteralias{orphelin}{inexistant}\fi
  \ifdefined\CORR \input{correctif-a-remise}\fi
\fi
\crefname{theoreme}{théorème}{théorèmes}
\crefname{proposition}{proposition}{propositions}
\crefname{corollaire}{corollaire}{corollaires}
\crefname{remarque}{remarque}{remarques}
\begin{document}
\chapter{Premier}
\begin{proposition}\label{p1}a\end{proposition}
\begin{theoreme}\label{t1}b\end{theoreme}
\begin{remarque}\label{r1}c\end{remarque}
\begin{corollaire}\label{c1}d\end{corollaire}
\chapter{Second}
\begin{corollaire}\label{c2}e\end{corollaire}
\begin{theoreme}\label{t2}f\end{theoreme}
RENVOIS \ref{p1} \ref{t1} \ref{r1} \ref{c1} \ref{c2} \ref{t2} ; \cref{p1} ; \cref{t1} ; \cref{r1} ; \cref{c1} ; \cref{c2} ; \cref{t2} FIN
\end{document}
"                    amsthm seul (témoin)
%   lualatex -jobname=faux     "\def\CORR{1}\def\FAUX{1}% Juge de la troisième passe : document minimal écrit de zéro pour la cause A (thmtools sous LaTeX 2026-06-01).
% Variantes, par une macro posée sur la ligne de commande (deux passes chacune, depuis memoire/essai/juge3/minimal/) :
%   lualatex -jobname=nu       "\input{minimal-juge3}"                                   sans correctif
%   lualatex -jobname=corr     "\def\CORR{1}\input{minimal-juge3}"                       correctif A du préambule
%   lualatex -jobname=crochet  "\def\SANSCROCHET{1}\input{minimal-juge3}"                sans correctif, crochet retiré
%   lualatex -jobname=amsseul  "\def\AMSSEUL{1}\input{minimal-juge3}"                    amsthm seul (témoin)
%   lualatex -jobname=faux     "\def\CORR{1}\def\FAUX{1}\input{minimal-juge3}"           deux vrais défauts : doivent tomber
%   lualatex -jobname=trace    "\def\TRACE{1}\input{minimal-juge3}"                      trace de l'appel
% Les lignes du correctif sont tirées du préambule commité par extraire_correctifs.py, pas retapées.
\documentclass{scrbook}
\ifdefined\SANSCROCHET \RemoveFromHook{package/amsthm/after}[firstaid/aliascounter]\fi
\usepackage{amsthm}
\usepackage{hyperref}
\usepackage[french, nameinlink, noabbrev]{cleveref}
\ifdefined\AMSSEUL
  \newtheorem{theoreme}{Théorème}[chapter]
  \newtheorem{proposition}[theoreme]{Proposition}
  \newtheorem{corollaire}[theoreme]{Corollaire}
  \newtheorem{remarque}{Remarque}[chapter]
\else
  \usepackage{thmtools}
  \ifdefined\CORR \input{correctif-a-pose}\fi
  \declaretheorem[numberwithin=chapter, name=Théorème]{theoreme}
  \ifdefined\TRACE \tracingmacros=2 \tracingonline=0 \fi
  \declaretheorem[sibling=theoreme, name=Proposition]{proposition}
  \ifdefined\TRACE \tracingmacros=0 \fi
  \declaretheorem[sibling=theoreme, name=Corollaire]{corollaire}
  \declaretheorem[numberwithin=chapter, name=Remarque]{remarque}
  \ifdefined\FAUX \newcounteralias{corollaire}{section}\newcounteralias{orphelin}{inexistant}\fi
  \ifdefined\CORR \input{correctif-a-remise}\fi
\fi
\crefname{theoreme}{théorème}{théorèmes}
\crefname{proposition}{proposition}{propositions}
\crefname{corollaire}{corollaire}{corollaires}
\crefname{remarque}{remarque}{remarques}
\begin{document}
\chapter{Premier}
\begin{proposition}\label{p1}a\end{proposition}
\begin{theoreme}\label{t1}b\end{theoreme}
\begin{remarque}\label{r1}c\end{remarque}
\begin{corollaire}\label{c1}d\end{corollaire}
\chapter{Second}
\begin{corollaire}\label{c2}e\end{corollaire}
\begin{theoreme}\label{t2}f\end{theoreme}
RENVOIS \ref{p1} \ref{t1} \ref{r1} \ref{c1} \ref{c2} \ref{t2} ; \cref{p1} ; \cref{t1} ; \cref{r1} ; \cref{c1} ; \cref{c2} ; \cref{t2} FIN
\end{document}
"           deux vrais défauts : doivent tomber
%   lualatex -jobname=trace    "\def\TRACE{1}% Juge de la troisième passe : document minimal écrit de zéro pour la cause A (thmtools sous LaTeX 2026-06-01).
% Variantes, par une macro posée sur la ligne de commande (deux passes chacune, depuis memoire/essai/juge3/minimal/) :
%   lualatex -jobname=nu       "\input{minimal-juge3}"                                   sans correctif
%   lualatex -jobname=corr     "\def\CORR{1}\input{minimal-juge3}"                       correctif A du préambule
%   lualatex -jobname=crochet  "\def\SANSCROCHET{1}\input{minimal-juge3}"                sans correctif, crochet retiré
%   lualatex -jobname=amsseul  "\def\AMSSEUL{1}\input{minimal-juge3}"                    amsthm seul (témoin)
%   lualatex -jobname=faux     "\def\CORR{1}\def\FAUX{1}\input{minimal-juge3}"           deux vrais défauts : doivent tomber
%   lualatex -jobname=trace    "\def\TRACE{1}\input{minimal-juge3}"                      trace de l'appel
% Les lignes du correctif sont tirées du préambule commité par extraire_correctifs.py, pas retapées.
\documentclass{scrbook}
\ifdefined\SANSCROCHET \RemoveFromHook{package/amsthm/after}[firstaid/aliascounter]\fi
\usepackage{amsthm}
\usepackage{hyperref}
\usepackage[french, nameinlink, noabbrev]{cleveref}
\ifdefined\AMSSEUL
  \newtheorem{theoreme}{Théorème}[chapter]
  \newtheorem{proposition}[theoreme]{Proposition}
  \newtheorem{corollaire}[theoreme]{Corollaire}
  \newtheorem{remarque}{Remarque}[chapter]
\else
  \usepackage{thmtools}
  \ifdefined\CORR \input{correctif-a-pose}\fi
  \declaretheorem[numberwithin=chapter, name=Théorème]{theoreme}
  \ifdefined\TRACE \tracingmacros=2 \tracingonline=0 \fi
  \declaretheorem[sibling=theoreme, name=Proposition]{proposition}
  \ifdefined\TRACE \tracingmacros=0 \fi
  \declaretheorem[sibling=theoreme, name=Corollaire]{corollaire}
  \declaretheorem[numberwithin=chapter, name=Remarque]{remarque}
  \ifdefined\FAUX \newcounteralias{corollaire}{section}\newcounteralias{orphelin}{inexistant}\fi
  \ifdefined\CORR \input{correctif-a-remise}\fi
\fi
\crefname{theoreme}{théorème}{théorèmes}
\crefname{proposition}{proposition}{propositions}
\crefname{corollaire}{corollaire}{corollaires}
\crefname{remarque}{remarque}{remarques}
\begin{document}
\chapter{Premier}
\begin{proposition}\label{p1}a\end{proposition}
\begin{theoreme}\label{t1}b\end{theoreme}
\begin{remarque}\label{r1}c\end{remarque}
\begin{corollaire}\label{c1}d\end{corollaire}
\chapter{Second}
\begin{corollaire}\label{c2}e\end{corollaire}
\begin{theoreme}\label{t2}f\end{theoreme}
RENVOIS \ref{p1} \ref{t1} \ref{r1} \ref{c1} \ref{c2} \ref{t2} ; \cref{p1} ; \cref{t1} ; \cref{r1} ; \cref{c1} ; \cref{c2} ; \cref{t2} FIN
\end{document}
"                      trace de l'appel
% Les lignes du correctif sont tirées du préambule commité par extraire_correctifs.py, pas retapées.
\documentclass{scrbook}
\ifdefined\SANSCROCHET \RemoveFromHook{package/amsthm/after}[firstaid/aliascounter]\fi
\usepackage{amsthm}
\usepackage{hyperref}
\usepackage[french, nameinlink, noabbrev]{cleveref}
\ifdefined\AMSSEUL
  \newtheorem{theoreme}{Théorème}[chapter]
  \newtheorem{proposition}[theoreme]{Proposition}
  \newtheorem{corollaire}[theoreme]{Corollaire}
  \newtheorem{remarque}{Remarque}[chapter]
\else
  \usepackage{thmtools}
  \ifdefined\CORR \input{correctif-a-pose}\fi
  \declaretheorem[numberwithin=chapter, name=Théorème]{theoreme}
  \ifdefined\TRACE \tracingmacros=2 \tracingonline=0 \fi
  \declaretheorem[sibling=theoreme, name=Proposition]{proposition}
  \ifdefined\TRACE \tracingmacros=0 \fi
  \declaretheorem[sibling=theoreme, name=Corollaire]{corollaire}
  \declaretheorem[numberwithin=chapter, name=Remarque]{remarque}
  \ifdefined\FAUX \newcounteralias{corollaire}{section}\newcounteralias{orphelin}{inexistant}\fi
  \ifdefined\CORR \input{correctif-a-remise}\fi
\fi
\crefname{theoreme}{théorème}{théorèmes}
\crefname{proposition}{proposition}{propositions}
\crefname{corollaire}{corollaire}{corollaires}
\crefname{remarque}{remarque}{remarques}
\begin{document}
\chapter{Premier}
\begin{proposition}\label{p1}a\end{proposition}
\begin{theoreme}\label{t1}b\end{theoreme}
\begin{remarque}\label{r1}c\end{remarque}
\begin{corollaire}\label{c1}d\end{corollaire}
\chapter{Second}
\begin{corollaire}\label{c2}e\end{corollaire}
\begin{theoreme}\label{t2}f\end{theoreme}
RENVOIS \ref{p1} \ref{t1} \ref{r1} \ref{c1} \ref{c2} \ref{t2} ; \cref{p1} ; \cref{t1} ; \cref{r1} ; \cref{c1} ; \cref{c2} ; \cref{t2} FIN
\end{document}
"                    amsthm seul (témoin)
%   lualatex -jobname=faux     "\def\CORR{1}\def\FAUX{1}% Juge de la troisième passe : document minimal écrit de zéro pour la cause A (thmtools sous LaTeX 2026-06-01).
% Variantes, par une macro posée sur la ligne de commande (deux passes chacune, depuis memoire/essai/juge3/minimal/) :
%   lualatex -jobname=nu       "% Juge de la troisième passe : document minimal écrit de zéro pour la cause A (thmtools sous LaTeX 2026-06-01).
% Variantes, par une macro posée sur la ligne de commande (deux passes chacune, depuis memoire/essai/juge3/minimal/) :
%   lualatex -jobname=nu       "\input{minimal-juge3}"                                   sans correctif
%   lualatex -jobname=corr     "\def\CORR{1}\input{minimal-juge3}"                       correctif A du préambule
%   lualatex -jobname=crochet  "\def\SANSCROCHET{1}\input{minimal-juge3}"                sans correctif, crochet retiré
%   lualatex -jobname=amsseul  "\def\AMSSEUL{1}\input{minimal-juge3}"                    amsthm seul (témoin)
%   lualatex -jobname=faux     "\def\CORR{1}\def\FAUX{1}\input{minimal-juge3}"           deux vrais défauts : doivent tomber
%   lualatex -jobname=trace    "\def\TRACE{1}\input{minimal-juge3}"                      trace de l'appel
% Les lignes du correctif sont tirées du préambule commité par extraire_correctifs.py, pas retapées.
\documentclass{scrbook}
\ifdefined\SANSCROCHET \RemoveFromHook{package/amsthm/after}[firstaid/aliascounter]\fi
\usepackage{amsthm}
\usepackage{hyperref}
\usepackage[french, nameinlink, noabbrev]{cleveref}
\ifdefined\AMSSEUL
  \newtheorem{theoreme}{Théorème}[chapter]
  \newtheorem{proposition}[theoreme]{Proposition}
  \newtheorem{corollaire}[theoreme]{Corollaire}
  \newtheorem{remarque}{Remarque}[chapter]
\else
  \usepackage{thmtools}
  \ifdefined\CORR \input{correctif-a-pose}\fi
  \declaretheorem[numberwithin=chapter, name=Théorème]{theoreme}
  \ifdefined\TRACE \tracingmacros=2 \tracingonline=0 \fi
  \declaretheorem[sibling=theoreme, name=Proposition]{proposition}
  \ifdefined\TRACE \tracingmacros=0 \fi
  \declaretheorem[sibling=theoreme, name=Corollaire]{corollaire}
  \declaretheorem[numberwithin=chapter, name=Remarque]{remarque}
  \ifdefined\FAUX \newcounteralias{corollaire}{section}\newcounteralias{orphelin}{inexistant}\fi
  \ifdefined\CORR \input{correctif-a-remise}\fi
\fi
\crefname{theoreme}{théorème}{théorèmes}
\crefname{proposition}{proposition}{propositions}
\crefname{corollaire}{corollaire}{corollaires}
\crefname{remarque}{remarque}{remarques}
\begin{document}
\chapter{Premier}
\begin{proposition}\label{p1}a\end{proposition}
\begin{theoreme}\label{t1}b\end{theoreme}
\begin{remarque}\label{r1}c\end{remarque}
\begin{corollaire}\label{c1}d\end{corollaire}
\chapter{Second}
\begin{corollaire}\label{c2}e\end{corollaire}
\begin{theoreme}\label{t2}f\end{theoreme}
RENVOIS \ref{p1} \ref{t1} \ref{r1} \ref{c1} \ref{c2} \ref{t2} ; \cref{p1} ; \cref{t1} ; \cref{r1} ; \cref{c1} ; \cref{c2} ; \cref{t2} FIN
\end{document}
"                                   sans correctif
%   lualatex -jobname=corr     "\def\CORR{1}% Juge de la troisième passe : document minimal écrit de zéro pour la cause A (thmtools sous LaTeX 2026-06-01).
% Variantes, par une macro posée sur la ligne de commande (deux passes chacune, depuis memoire/essai/juge3/minimal/) :
%   lualatex -jobname=nu       "\input{minimal-juge3}"                                   sans correctif
%   lualatex -jobname=corr     "\def\CORR{1}\input{minimal-juge3}"                       correctif A du préambule
%   lualatex -jobname=crochet  "\def\SANSCROCHET{1}\input{minimal-juge3}"                sans correctif, crochet retiré
%   lualatex -jobname=amsseul  "\def\AMSSEUL{1}\input{minimal-juge3}"                    amsthm seul (témoin)
%   lualatex -jobname=faux     "\def\CORR{1}\def\FAUX{1}\input{minimal-juge3}"           deux vrais défauts : doivent tomber
%   lualatex -jobname=trace    "\def\TRACE{1}\input{minimal-juge3}"                      trace de l'appel
% Les lignes du correctif sont tirées du préambule commité par extraire_correctifs.py, pas retapées.
\documentclass{scrbook}
\ifdefined\SANSCROCHET \RemoveFromHook{package/amsthm/after}[firstaid/aliascounter]\fi
\usepackage{amsthm}
\usepackage{hyperref}
\usepackage[french, nameinlink, noabbrev]{cleveref}
\ifdefined\AMSSEUL
  \newtheorem{theoreme}{Théorème}[chapter]
  \newtheorem{proposition}[theoreme]{Proposition}
  \newtheorem{corollaire}[theoreme]{Corollaire}
  \newtheorem{remarque}{Remarque}[chapter]
\else
  \usepackage{thmtools}
  \ifdefined\CORR \input{correctif-a-pose}\fi
  \declaretheorem[numberwithin=chapter, name=Théorème]{theoreme}
  \ifdefined\TRACE \tracingmacros=2 \tracingonline=0 \fi
  \declaretheorem[sibling=theoreme, name=Proposition]{proposition}
  \ifdefined\TRACE \tracingmacros=0 \fi
  \declaretheorem[sibling=theoreme, name=Corollaire]{corollaire}
  \declaretheorem[numberwithin=chapter, name=Remarque]{remarque}
  \ifdefined\FAUX \newcounteralias{corollaire}{section}\newcounteralias{orphelin}{inexistant}\fi
  \ifdefined\CORR \input{correctif-a-remise}\fi
\fi
\crefname{theoreme}{théorème}{théorèmes}
\crefname{proposition}{proposition}{propositions}
\crefname{corollaire}{corollaire}{corollaires}
\crefname{remarque}{remarque}{remarques}
\begin{document}
\chapter{Premier}
\begin{proposition}\label{p1}a\end{proposition}
\begin{theoreme}\label{t1}b\end{theoreme}
\begin{remarque}\label{r1}c\end{remarque}
\begin{corollaire}\label{c1}d\end{corollaire}
\chapter{Second}
\begin{corollaire}\label{c2}e\end{corollaire}
\begin{theoreme}\label{t2}f\end{theoreme}
RENVOIS \ref{p1} \ref{t1} \ref{r1} \ref{c1} \ref{c2} \ref{t2} ; \cref{p1} ; \cref{t1} ; \cref{r1} ; \cref{c1} ; \cref{c2} ; \cref{t2} FIN
\end{document}
"                       correctif A du préambule
%   lualatex -jobname=crochet  "\def\SANSCROCHET{1}% Juge de la troisième passe : document minimal écrit de zéro pour la cause A (thmtools sous LaTeX 2026-06-01).
% Variantes, par une macro posée sur la ligne de commande (deux passes chacune, depuis memoire/essai/juge3/minimal/) :
%   lualatex -jobname=nu       "\input{minimal-juge3}"                                   sans correctif
%   lualatex -jobname=corr     "\def\CORR{1}\input{minimal-juge3}"                       correctif A du préambule
%   lualatex -jobname=crochet  "\def\SANSCROCHET{1}\input{minimal-juge3}"                sans correctif, crochet retiré
%   lualatex -jobname=amsseul  "\def\AMSSEUL{1}\input{minimal-juge3}"                    amsthm seul (témoin)
%   lualatex -jobname=faux     "\def\CORR{1}\def\FAUX{1}\input{minimal-juge3}"           deux vrais défauts : doivent tomber
%   lualatex -jobname=trace    "\def\TRACE{1}\input{minimal-juge3}"                      trace de l'appel
% Les lignes du correctif sont tirées du préambule commité par extraire_correctifs.py, pas retapées.
\documentclass{scrbook}
\ifdefined\SANSCROCHET \RemoveFromHook{package/amsthm/after}[firstaid/aliascounter]\fi
\usepackage{amsthm}
\usepackage{hyperref}
\usepackage[french, nameinlink, noabbrev]{cleveref}
\ifdefined\AMSSEUL
  \newtheorem{theoreme}{Théorème}[chapter]
  \newtheorem{proposition}[theoreme]{Proposition}
  \newtheorem{corollaire}[theoreme]{Corollaire}
  \newtheorem{remarque}{Remarque}[chapter]
\else
  \usepackage{thmtools}
  \ifdefined\CORR \input{correctif-a-pose}\fi
  \declaretheorem[numberwithin=chapter, name=Théorème]{theoreme}
  \ifdefined\TRACE \tracingmacros=2 \tracingonline=0 \fi
  \declaretheorem[sibling=theoreme, name=Proposition]{proposition}
  \ifdefined\TRACE \tracingmacros=0 \fi
  \declaretheorem[sibling=theoreme, name=Corollaire]{corollaire}
  \declaretheorem[numberwithin=chapter, name=Remarque]{remarque}
  \ifdefined\FAUX \newcounteralias{corollaire}{section}\newcounteralias{orphelin}{inexistant}\fi
  \ifdefined\CORR \input{correctif-a-remise}\fi
\fi
\crefname{theoreme}{théorème}{théorèmes}
\crefname{proposition}{proposition}{propositions}
\crefname{corollaire}{corollaire}{corollaires}
\crefname{remarque}{remarque}{remarques}
\begin{document}
\chapter{Premier}
\begin{proposition}\label{p1}a\end{proposition}
\begin{theoreme}\label{t1}b\end{theoreme}
\begin{remarque}\label{r1}c\end{remarque}
\begin{corollaire}\label{c1}d\end{corollaire}
\chapter{Second}
\begin{corollaire}\label{c2}e\end{corollaire}
\begin{theoreme}\label{t2}f\end{theoreme}
RENVOIS \ref{p1} \ref{t1} \ref{r1} \ref{c1} \ref{c2} \ref{t2} ; \cref{p1} ; \cref{t1} ; \cref{r1} ; \cref{c1} ; \cref{c2} ; \cref{t2} FIN
\end{document}
"                sans correctif, crochet retiré
%   lualatex -jobname=amsseul  "\def\AMSSEUL{1}% Juge de la troisième passe : document minimal écrit de zéro pour la cause A (thmtools sous LaTeX 2026-06-01).
% Variantes, par une macro posée sur la ligne de commande (deux passes chacune, depuis memoire/essai/juge3/minimal/) :
%   lualatex -jobname=nu       "\input{minimal-juge3}"                                   sans correctif
%   lualatex -jobname=corr     "\def\CORR{1}\input{minimal-juge3}"                       correctif A du préambule
%   lualatex -jobname=crochet  "\def\SANSCROCHET{1}\input{minimal-juge3}"                sans correctif, crochet retiré
%   lualatex -jobname=amsseul  "\def\AMSSEUL{1}\input{minimal-juge3}"                    amsthm seul (témoin)
%   lualatex -jobname=faux     "\def\CORR{1}\def\FAUX{1}\input{minimal-juge3}"           deux vrais défauts : doivent tomber
%   lualatex -jobname=trace    "\def\TRACE{1}\input{minimal-juge3}"                      trace de l'appel
% Les lignes du correctif sont tirées du préambule commité par extraire_correctifs.py, pas retapées.
\documentclass{scrbook}
\ifdefined\SANSCROCHET \RemoveFromHook{package/amsthm/after}[firstaid/aliascounter]\fi
\usepackage{amsthm}
\usepackage{hyperref}
\usepackage[french, nameinlink, noabbrev]{cleveref}
\ifdefined\AMSSEUL
  \newtheorem{theoreme}{Théorème}[chapter]
  \newtheorem{proposition}[theoreme]{Proposition}
  \newtheorem{corollaire}[theoreme]{Corollaire}
  \newtheorem{remarque}{Remarque}[chapter]
\else
  \usepackage{thmtools}
  \ifdefined\CORR \input{correctif-a-pose}\fi
  \declaretheorem[numberwithin=chapter, name=Théorème]{theoreme}
  \ifdefined\TRACE \tracingmacros=2 \tracingonline=0 \fi
  \declaretheorem[sibling=theoreme, name=Proposition]{proposition}
  \ifdefined\TRACE \tracingmacros=0 \fi
  \declaretheorem[sibling=theoreme, name=Corollaire]{corollaire}
  \declaretheorem[numberwithin=chapter, name=Remarque]{remarque}
  \ifdefined\FAUX \newcounteralias{corollaire}{section}\newcounteralias{orphelin}{inexistant}\fi
  \ifdefined\CORR \input{correctif-a-remise}\fi
\fi
\crefname{theoreme}{théorème}{théorèmes}
\crefname{proposition}{proposition}{propositions}
\crefname{corollaire}{corollaire}{corollaires}
\crefname{remarque}{remarque}{remarques}
\begin{document}
\chapter{Premier}
\begin{proposition}\label{p1}a\end{proposition}
\begin{theoreme}\label{t1}b\end{theoreme}
\begin{remarque}\label{r1}c\end{remarque}
\begin{corollaire}\label{c1}d\end{corollaire}
\chapter{Second}
\begin{corollaire}\label{c2}e\end{corollaire}
\begin{theoreme}\label{t2}f\end{theoreme}
RENVOIS \ref{p1} \ref{t1} \ref{r1} \ref{c1} \ref{c2} \ref{t2} ; \cref{p1} ; \cref{t1} ; \cref{r1} ; \cref{c1} ; \cref{c2} ; \cref{t2} FIN
\end{document}
"                    amsthm seul (témoin)
%   lualatex -jobname=faux     "\def\CORR{1}\def\FAUX{1}% Juge de la troisième passe : document minimal écrit de zéro pour la cause A (thmtools sous LaTeX 2026-06-01).
% Variantes, par une macro posée sur la ligne de commande (deux passes chacune, depuis memoire/essai/juge3/minimal/) :
%   lualatex -jobname=nu       "\input{minimal-juge3}"                                   sans correctif
%   lualatex -jobname=corr     "\def\CORR{1}\input{minimal-juge3}"                       correctif A du préambule
%   lualatex -jobname=crochet  "\def\SANSCROCHET{1}\input{minimal-juge3}"                sans correctif, crochet retiré
%   lualatex -jobname=amsseul  "\def\AMSSEUL{1}\input{minimal-juge3}"                    amsthm seul (témoin)
%   lualatex -jobname=faux     "\def\CORR{1}\def\FAUX{1}\input{minimal-juge3}"           deux vrais défauts : doivent tomber
%   lualatex -jobname=trace    "\def\TRACE{1}\input{minimal-juge3}"                      trace de l'appel
% Les lignes du correctif sont tirées du préambule commité par extraire_correctifs.py, pas retapées.
\documentclass{scrbook}
\ifdefined\SANSCROCHET \RemoveFromHook{package/amsthm/after}[firstaid/aliascounter]\fi
\usepackage{amsthm}
\usepackage{hyperref}
\usepackage[french, nameinlink, noabbrev]{cleveref}
\ifdefined\AMSSEUL
  \newtheorem{theoreme}{Théorème}[chapter]
  \newtheorem{proposition}[theoreme]{Proposition}
  \newtheorem{corollaire}[theoreme]{Corollaire}
  \newtheorem{remarque}{Remarque}[chapter]
\else
  \usepackage{thmtools}
  \ifdefined\CORR \input{correctif-a-pose}\fi
  \declaretheorem[numberwithin=chapter, name=Théorème]{theoreme}
  \ifdefined\TRACE \tracingmacros=2 \tracingonline=0 \fi
  \declaretheorem[sibling=theoreme, name=Proposition]{proposition}
  \ifdefined\TRACE \tracingmacros=0 \fi
  \declaretheorem[sibling=theoreme, name=Corollaire]{corollaire}
  \declaretheorem[numberwithin=chapter, name=Remarque]{remarque}
  \ifdefined\FAUX \newcounteralias{corollaire}{section}\newcounteralias{orphelin}{inexistant}\fi
  \ifdefined\CORR \input{correctif-a-remise}\fi
\fi
\crefname{theoreme}{théorème}{théorèmes}
\crefname{proposition}{proposition}{propositions}
\crefname{corollaire}{corollaire}{corollaires}
\crefname{remarque}{remarque}{remarques}
\begin{document}
\chapter{Premier}
\begin{proposition}\label{p1}a\end{proposition}
\begin{theoreme}\label{t1}b\end{theoreme}
\begin{remarque}\label{r1}c\end{remarque}
\begin{corollaire}\label{c1}d\end{corollaire}
\chapter{Second}
\begin{corollaire}\label{c2}e\end{corollaire}
\begin{theoreme}\label{t2}f\end{theoreme}
RENVOIS \ref{p1} \ref{t1} \ref{r1} \ref{c1} \ref{c2} \ref{t2} ; \cref{p1} ; \cref{t1} ; \cref{r1} ; \cref{c1} ; \cref{c2} ; \cref{t2} FIN
\end{document}
"           deux vrais défauts : doivent tomber
%   lualatex -jobname=trace    "\def\TRACE{1}% Juge de la troisième passe : document minimal écrit de zéro pour la cause A (thmtools sous LaTeX 2026-06-01).
% Variantes, par une macro posée sur la ligne de commande (deux passes chacune, depuis memoire/essai/juge3/minimal/) :
%   lualatex -jobname=nu       "\input{minimal-juge3}"                                   sans correctif
%   lualatex -jobname=corr     "\def\CORR{1}\input{minimal-juge3}"                       correctif A du préambule
%   lualatex -jobname=crochet  "\def\SANSCROCHET{1}\input{minimal-juge3}"                sans correctif, crochet retiré
%   lualatex -jobname=amsseul  "\def\AMSSEUL{1}\input{minimal-juge3}"                    amsthm seul (témoin)
%   lualatex -jobname=faux     "\def\CORR{1}\def\FAUX{1}\input{minimal-juge3}"           deux vrais défauts : doivent tomber
%   lualatex -jobname=trace    "\def\TRACE{1}\input{minimal-juge3}"                      trace de l'appel
% Les lignes du correctif sont tirées du préambule commité par extraire_correctifs.py, pas retapées.
\documentclass{scrbook}
\ifdefined\SANSCROCHET \RemoveFromHook{package/amsthm/after}[firstaid/aliascounter]\fi
\usepackage{amsthm}
\usepackage{hyperref}
\usepackage[french, nameinlink, noabbrev]{cleveref}
\ifdefined\AMSSEUL
  \newtheorem{theoreme}{Théorème}[chapter]
  \newtheorem{proposition}[theoreme]{Proposition}
  \newtheorem{corollaire}[theoreme]{Corollaire}
  \newtheorem{remarque}{Remarque}[chapter]
\else
  \usepackage{thmtools}
  \ifdefined\CORR \input{correctif-a-pose}\fi
  \declaretheorem[numberwithin=chapter, name=Théorème]{theoreme}
  \ifdefined\TRACE \tracingmacros=2 \tracingonline=0 \fi
  \declaretheorem[sibling=theoreme, name=Proposition]{proposition}
  \ifdefined\TRACE \tracingmacros=0 \fi
  \declaretheorem[sibling=theoreme, name=Corollaire]{corollaire}
  \declaretheorem[numberwithin=chapter, name=Remarque]{remarque}
  \ifdefined\FAUX \newcounteralias{corollaire}{section}\newcounteralias{orphelin}{inexistant}\fi
  \ifdefined\CORR \input{correctif-a-remise}\fi
\fi
\crefname{theoreme}{théorème}{théorèmes}
\crefname{proposition}{proposition}{propositions}
\crefname{corollaire}{corollaire}{corollaires}
\crefname{remarque}{remarque}{remarques}
\begin{document}
\chapter{Premier}
\begin{proposition}\label{p1}a\end{proposition}
\begin{theoreme}\label{t1}b\end{theoreme}
\begin{remarque}\label{r1}c\end{remarque}
\begin{corollaire}\label{c1}d\end{corollaire}
\chapter{Second}
\begin{corollaire}\label{c2}e\end{corollaire}
\begin{theoreme}\label{t2}f\end{theoreme}
RENVOIS \ref{p1} \ref{t1} \ref{r1} \ref{c1} \ref{c2} \ref{t2} ; \cref{p1} ; \cref{t1} ; \cref{r1} ; \cref{c1} ; \cref{c2} ; \cref{t2} FIN
\end{document}
"                      trace de l'appel
% Les lignes du correctif sont tirées du préambule commité par extraire_correctifs.py, pas retapées.
\documentclass{scrbook}
\ifdefined\SANSCROCHET \RemoveFromHook{package/amsthm/after}[firstaid/aliascounter]\fi
\usepackage{amsthm}
\usepackage{hyperref}
\usepackage[french, nameinlink, noabbrev]{cleveref}
\ifdefined\AMSSEUL
  \newtheorem{theoreme}{Théorème}[chapter]
  \newtheorem{proposition}[theoreme]{Proposition}
  \newtheorem{corollaire}[theoreme]{Corollaire}
  \newtheorem{remarque}{Remarque}[chapter]
\else
  \usepackage{thmtools}
  \ifdefined\CORR \input{correctif-a-pose}\fi
  \declaretheorem[numberwithin=chapter, name=Théorème]{theoreme}
  \ifdefined\TRACE \tracingmacros=2 \tracingonline=0 \fi
  \declaretheorem[sibling=theoreme, name=Proposition]{proposition}
  \ifdefined\TRACE \tracingmacros=0 \fi
  \declaretheorem[sibling=theoreme, name=Corollaire]{corollaire}
  \declaretheorem[numberwithin=chapter, name=Remarque]{remarque}
  \ifdefined\FAUX \newcounteralias{corollaire}{section}\newcounteralias{orphelin}{inexistant}\fi
  \ifdefined\CORR \input{correctif-a-remise}\fi
\fi
\crefname{theoreme}{théorème}{théorèmes}
\crefname{proposition}{proposition}{propositions}
\crefname{corollaire}{corollaire}{corollaires}
\crefname{remarque}{remarque}{remarques}
\begin{document}
\chapter{Premier}
\begin{proposition}\label{p1}a\end{proposition}
\begin{theoreme}\label{t1}b\end{theoreme}
\begin{remarque}\label{r1}c\end{remarque}
\begin{corollaire}\label{c1}d\end{corollaire}
\chapter{Second}
\begin{corollaire}\label{c2}e\end{corollaire}
\begin{theoreme}\label{t2}f\end{theoreme}
RENVOIS \ref{p1} \ref{t1} \ref{r1} \ref{c1} \ref{c2} \ref{t2} ; \cref{p1} ; \cref{t1} ; \cref{r1} ; \cref{c1} ; \cref{c2} ; \cref{t2} FIN
\end{document}
"           deux vrais défauts : doivent tomber
%   lualatex -jobname=trace    "\def\TRACE{1}% Juge de la troisième passe : document minimal écrit de zéro pour la cause A (thmtools sous LaTeX 2026-06-01).
% Variantes, par une macro posée sur la ligne de commande (deux passes chacune, depuis memoire/essai/juge3/minimal/) :
%   lualatex -jobname=nu       "% Juge de la troisième passe : document minimal écrit de zéro pour la cause A (thmtools sous LaTeX 2026-06-01).
% Variantes, par une macro posée sur la ligne de commande (deux passes chacune, depuis memoire/essai/juge3/minimal/) :
%   lualatex -jobname=nu       "\input{minimal-juge3}"                                   sans correctif
%   lualatex -jobname=corr     "\def\CORR{1}\input{minimal-juge3}"                       correctif A du préambule
%   lualatex -jobname=crochet  "\def\SANSCROCHET{1}\input{minimal-juge3}"                sans correctif, crochet retiré
%   lualatex -jobname=amsseul  "\def\AMSSEUL{1}\input{minimal-juge3}"                    amsthm seul (témoin)
%   lualatex -jobname=faux     "\def\CORR{1}\def\FAUX{1}\input{minimal-juge3}"           deux vrais défauts : doivent tomber
%   lualatex -jobname=trace    "\def\TRACE{1}\input{minimal-juge3}"                      trace de l'appel
% Les lignes du correctif sont tirées du préambule commité par extraire_correctifs.py, pas retapées.
\documentclass{scrbook}
\ifdefined\SANSCROCHET \RemoveFromHook{package/amsthm/after}[firstaid/aliascounter]\fi
\usepackage{amsthm}
\usepackage{hyperref}
\usepackage[french, nameinlink, noabbrev]{cleveref}
\ifdefined\AMSSEUL
  \newtheorem{theoreme}{Théorème}[chapter]
  \newtheorem{proposition}[theoreme]{Proposition}
  \newtheorem{corollaire}[theoreme]{Corollaire}
  \newtheorem{remarque}{Remarque}[chapter]
\else
  \usepackage{thmtools}
  \ifdefined\CORR \input{correctif-a-pose}\fi
  \declaretheorem[numberwithin=chapter, name=Théorème]{theoreme}
  \ifdefined\TRACE \tracingmacros=2 \tracingonline=0 \fi
  \declaretheorem[sibling=theoreme, name=Proposition]{proposition}
  \ifdefined\TRACE \tracingmacros=0 \fi
  \declaretheorem[sibling=theoreme, name=Corollaire]{corollaire}
  \declaretheorem[numberwithin=chapter, name=Remarque]{remarque}
  \ifdefined\FAUX \newcounteralias{corollaire}{section}\newcounteralias{orphelin}{inexistant}\fi
  \ifdefined\CORR \input{correctif-a-remise}\fi
\fi
\crefname{theoreme}{théorème}{théorèmes}
\crefname{proposition}{proposition}{propositions}
\crefname{corollaire}{corollaire}{corollaires}
\crefname{remarque}{remarque}{remarques}
\begin{document}
\chapter{Premier}
\begin{proposition}\label{p1}a\end{proposition}
\begin{theoreme}\label{t1}b\end{theoreme}
\begin{remarque}\label{r1}c\end{remarque}
\begin{corollaire}\label{c1}d\end{corollaire}
\chapter{Second}
\begin{corollaire}\label{c2}e\end{corollaire}
\begin{theoreme}\label{t2}f\end{theoreme}
RENVOIS \ref{p1} \ref{t1} \ref{r1} \ref{c1} \ref{c2} \ref{t2} ; \cref{p1} ; \cref{t1} ; \cref{r1} ; \cref{c1} ; \cref{c2} ; \cref{t2} FIN
\end{document}
"                                   sans correctif
%   lualatex -jobname=corr     "\def\CORR{1}% Juge de la troisième passe : document minimal écrit de zéro pour la cause A (thmtools sous LaTeX 2026-06-01).
% Variantes, par une macro posée sur la ligne de commande (deux passes chacune, depuis memoire/essai/juge3/minimal/) :
%   lualatex -jobname=nu       "\input{minimal-juge3}"                                   sans correctif
%   lualatex -jobname=corr     "\def\CORR{1}\input{minimal-juge3}"                       correctif A du préambule
%   lualatex -jobname=crochet  "\def\SANSCROCHET{1}\input{minimal-juge3}"                sans correctif, crochet retiré
%   lualatex -jobname=amsseul  "\def\AMSSEUL{1}\input{minimal-juge3}"                    amsthm seul (témoin)
%   lualatex -jobname=faux     "\def\CORR{1}\def\FAUX{1}\input{minimal-juge3}"           deux vrais défauts : doivent tomber
%   lualatex -jobname=trace    "\def\TRACE{1}\input{minimal-juge3}"                      trace de l'appel
% Les lignes du correctif sont tirées du préambule commité par extraire_correctifs.py, pas retapées.
\documentclass{scrbook}
\ifdefined\SANSCROCHET \RemoveFromHook{package/amsthm/after}[firstaid/aliascounter]\fi
\usepackage{amsthm}
\usepackage{hyperref}
\usepackage[french, nameinlink, noabbrev]{cleveref}
\ifdefined\AMSSEUL
  \newtheorem{theoreme}{Théorème}[chapter]
  \newtheorem{proposition}[theoreme]{Proposition}
  \newtheorem{corollaire}[theoreme]{Corollaire}
  \newtheorem{remarque}{Remarque}[chapter]
\else
  \usepackage{thmtools}
  \ifdefined\CORR \input{correctif-a-pose}\fi
  \declaretheorem[numberwithin=chapter, name=Théorème]{theoreme}
  \ifdefined\TRACE \tracingmacros=2 \tracingonline=0 \fi
  \declaretheorem[sibling=theoreme, name=Proposition]{proposition}
  \ifdefined\TRACE \tracingmacros=0 \fi
  \declaretheorem[sibling=theoreme, name=Corollaire]{corollaire}
  \declaretheorem[numberwithin=chapter, name=Remarque]{remarque}
  \ifdefined\FAUX \newcounteralias{corollaire}{section}\newcounteralias{orphelin}{inexistant}\fi
  \ifdefined\CORR \input{correctif-a-remise}\fi
\fi
\crefname{theoreme}{théorème}{théorèmes}
\crefname{proposition}{proposition}{propositions}
\crefname{corollaire}{corollaire}{corollaires}
\crefname{remarque}{remarque}{remarques}
\begin{document}
\chapter{Premier}
\begin{proposition}\label{p1}a\end{proposition}
\begin{theoreme}\label{t1}b\end{theoreme}
\begin{remarque}\label{r1}c\end{remarque}
\begin{corollaire}\label{c1}d\end{corollaire}
\chapter{Second}
\begin{corollaire}\label{c2}e\end{corollaire}
\begin{theoreme}\label{t2}f\end{theoreme}
RENVOIS \ref{p1} \ref{t1} \ref{r1} \ref{c1} \ref{c2} \ref{t2} ; \cref{p1} ; \cref{t1} ; \cref{r1} ; \cref{c1} ; \cref{c2} ; \cref{t2} FIN
\end{document}
"                       correctif A du préambule
%   lualatex -jobname=crochet  "\def\SANSCROCHET{1}% Juge de la troisième passe : document minimal écrit de zéro pour la cause A (thmtools sous LaTeX 2026-06-01).
% Variantes, par une macro posée sur la ligne de commande (deux passes chacune, depuis memoire/essai/juge3/minimal/) :
%   lualatex -jobname=nu       "\input{minimal-juge3}"                                   sans correctif
%   lualatex -jobname=corr     "\def\CORR{1}\input{minimal-juge3}"                       correctif A du préambule
%   lualatex -jobname=crochet  "\def\SANSCROCHET{1}\input{minimal-juge3}"                sans correctif, crochet retiré
%   lualatex -jobname=amsseul  "\def\AMSSEUL{1}\input{minimal-juge3}"                    amsthm seul (témoin)
%   lualatex -jobname=faux     "\def\CORR{1}\def\FAUX{1}\input{minimal-juge3}"           deux vrais défauts : doivent tomber
%   lualatex -jobname=trace    "\def\TRACE{1}\input{minimal-juge3}"                      trace de l'appel
% Les lignes du correctif sont tirées du préambule commité par extraire_correctifs.py, pas retapées.
\documentclass{scrbook}
\ifdefined\SANSCROCHET \RemoveFromHook{package/amsthm/after}[firstaid/aliascounter]\fi
\usepackage{amsthm}
\usepackage{hyperref}
\usepackage[french, nameinlink, noabbrev]{cleveref}
\ifdefined\AMSSEUL
  \newtheorem{theoreme}{Théorème}[chapter]
  \newtheorem{proposition}[theoreme]{Proposition}
  \newtheorem{corollaire}[theoreme]{Corollaire}
  \newtheorem{remarque}{Remarque}[chapter]
\else
  \usepackage{thmtools}
  \ifdefined\CORR \input{correctif-a-pose}\fi
  \declaretheorem[numberwithin=chapter, name=Théorème]{theoreme}
  \ifdefined\TRACE \tracingmacros=2 \tracingonline=0 \fi
  \declaretheorem[sibling=theoreme, name=Proposition]{proposition}
  \ifdefined\TRACE \tracingmacros=0 \fi
  \declaretheorem[sibling=theoreme, name=Corollaire]{corollaire}
  \declaretheorem[numberwithin=chapter, name=Remarque]{remarque}
  \ifdefined\FAUX \newcounteralias{corollaire}{section}\newcounteralias{orphelin}{inexistant}\fi
  \ifdefined\CORR \input{correctif-a-remise}\fi
\fi
\crefname{theoreme}{théorème}{théorèmes}
\crefname{proposition}{proposition}{propositions}
\crefname{corollaire}{corollaire}{corollaires}
\crefname{remarque}{remarque}{remarques}
\begin{document}
\chapter{Premier}
\begin{proposition}\label{p1}a\end{proposition}
\begin{theoreme}\label{t1}b\end{theoreme}
\begin{remarque}\label{r1}c\end{remarque}
\begin{corollaire}\label{c1}d\end{corollaire}
\chapter{Second}
\begin{corollaire}\label{c2}e\end{corollaire}
\begin{theoreme}\label{t2}f\end{theoreme}
RENVOIS \ref{p1} \ref{t1} \ref{r1} \ref{c1} \ref{c2} \ref{t2} ; \cref{p1} ; \cref{t1} ; \cref{r1} ; \cref{c1} ; \cref{c2} ; \cref{t2} FIN
\end{document}
"                sans correctif, crochet retiré
%   lualatex -jobname=amsseul  "\def\AMSSEUL{1}% Juge de la troisième passe : document minimal écrit de zéro pour la cause A (thmtools sous LaTeX 2026-06-01).
% Variantes, par une macro posée sur la ligne de commande (deux passes chacune, depuis memoire/essai/juge3/minimal/) :
%   lualatex -jobname=nu       "\input{minimal-juge3}"                                   sans correctif
%   lualatex -jobname=corr     "\def\CORR{1}\input{minimal-juge3}"                       correctif A du préambule
%   lualatex -jobname=crochet  "\def\SANSCROCHET{1}\input{minimal-juge3}"                sans correctif, crochet retiré
%   lualatex -jobname=amsseul  "\def\AMSSEUL{1}\input{minimal-juge3}"                    amsthm seul (témoin)
%   lualatex -jobname=faux     "\def\CORR{1}\def\FAUX{1}\input{minimal-juge3}"           deux vrais défauts : doivent tomber
%   lualatex -jobname=trace    "\def\TRACE{1}\input{minimal-juge3}"                      trace de l'appel
% Les lignes du correctif sont tirées du préambule commité par extraire_correctifs.py, pas retapées.
\documentclass{scrbook}
\ifdefined\SANSCROCHET \RemoveFromHook{package/amsthm/after}[firstaid/aliascounter]\fi
\usepackage{amsthm}
\usepackage{hyperref}
\usepackage[french, nameinlink, noabbrev]{cleveref}
\ifdefined\AMSSEUL
  \newtheorem{theoreme}{Théorème}[chapter]
  \newtheorem{proposition}[theoreme]{Proposition}
  \newtheorem{corollaire}[theoreme]{Corollaire}
  \newtheorem{remarque}{Remarque}[chapter]
\else
  \usepackage{thmtools}
  \ifdefined\CORR \input{correctif-a-pose}\fi
  \declaretheorem[numberwithin=chapter, name=Théorème]{theoreme}
  \ifdefined\TRACE \tracingmacros=2 \tracingonline=0 \fi
  \declaretheorem[sibling=theoreme, name=Proposition]{proposition}
  \ifdefined\TRACE \tracingmacros=0 \fi
  \declaretheorem[sibling=theoreme, name=Corollaire]{corollaire}
  \declaretheorem[numberwithin=chapter, name=Remarque]{remarque}
  \ifdefined\FAUX \newcounteralias{corollaire}{section}\newcounteralias{orphelin}{inexistant}\fi
  \ifdefined\CORR \input{correctif-a-remise}\fi
\fi
\crefname{theoreme}{théorème}{théorèmes}
\crefname{proposition}{proposition}{propositions}
\crefname{corollaire}{corollaire}{corollaires}
\crefname{remarque}{remarque}{remarques}
\begin{document}
\chapter{Premier}
\begin{proposition}\label{p1}a\end{proposition}
\begin{theoreme}\label{t1}b\end{theoreme}
\begin{remarque}\label{r1}c\end{remarque}
\begin{corollaire}\label{c1}d\end{corollaire}
\chapter{Second}
\begin{corollaire}\label{c2}e\end{corollaire}
\begin{theoreme}\label{t2}f\end{theoreme}
RENVOIS \ref{p1} \ref{t1} \ref{r1} \ref{c1} \ref{c2} \ref{t2} ; \cref{p1} ; \cref{t1} ; \cref{r1} ; \cref{c1} ; \cref{c2} ; \cref{t2} FIN
\end{document}
"                    amsthm seul (témoin)
%   lualatex -jobname=faux     "\def\CORR{1}\def\FAUX{1}% Juge de la troisième passe : document minimal écrit de zéro pour la cause A (thmtools sous LaTeX 2026-06-01).
% Variantes, par une macro posée sur la ligne de commande (deux passes chacune, depuis memoire/essai/juge3/minimal/) :
%   lualatex -jobname=nu       "\input{minimal-juge3}"                                   sans correctif
%   lualatex -jobname=corr     "\def\CORR{1}\input{minimal-juge3}"                       correctif A du préambule
%   lualatex -jobname=crochet  "\def\SANSCROCHET{1}\input{minimal-juge3}"                sans correctif, crochet retiré
%   lualatex -jobname=amsseul  "\def\AMSSEUL{1}\input{minimal-juge3}"                    amsthm seul (témoin)
%   lualatex -jobname=faux     "\def\CORR{1}\def\FAUX{1}\input{minimal-juge3}"           deux vrais défauts : doivent tomber
%   lualatex -jobname=trace    "\def\TRACE{1}\input{minimal-juge3}"                      trace de l'appel
% Les lignes du correctif sont tirées du préambule commité par extraire_correctifs.py, pas retapées.
\documentclass{scrbook}
\ifdefined\SANSCROCHET \RemoveFromHook{package/amsthm/after}[firstaid/aliascounter]\fi
\usepackage{amsthm}
\usepackage{hyperref}
\usepackage[french, nameinlink, noabbrev]{cleveref}
\ifdefined\AMSSEUL
  \newtheorem{theoreme}{Théorème}[chapter]
  \newtheorem{proposition}[theoreme]{Proposition}
  \newtheorem{corollaire}[theoreme]{Corollaire}
  \newtheorem{remarque}{Remarque}[chapter]
\else
  \usepackage{thmtools}
  \ifdefined\CORR \input{correctif-a-pose}\fi
  \declaretheorem[numberwithin=chapter, name=Théorème]{theoreme}
  \ifdefined\TRACE \tracingmacros=2 \tracingonline=0 \fi
  \declaretheorem[sibling=theoreme, name=Proposition]{proposition}
  \ifdefined\TRACE \tracingmacros=0 \fi
  \declaretheorem[sibling=theoreme, name=Corollaire]{corollaire}
  \declaretheorem[numberwithin=chapter, name=Remarque]{remarque}
  \ifdefined\FAUX \newcounteralias{corollaire}{section}\newcounteralias{orphelin}{inexistant}\fi
  \ifdefined\CORR \input{correctif-a-remise}\fi
\fi
\crefname{theoreme}{théorème}{théorèmes}
\crefname{proposition}{proposition}{propositions}
\crefname{corollaire}{corollaire}{corollaires}
\crefname{remarque}{remarque}{remarques}
\begin{document}
\chapter{Premier}
\begin{proposition}\label{p1}a\end{proposition}
\begin{theoreme}\label{t1}b\end{theoreme}
\begin{remarque}\label{r1}c\end{remarque}
\begin{corollaire}\label{c1}d\end{corollaire}
\chapter{Second}
\begin{corollaire}\label{c2}e\end{corollaire}
\begin{theoreme}\label{t2}f\end{theoreme}
RENVOIS \ref{p1} \ref{t1} \ref{r1} \ref{c1} \ref{c2} \ref{t2} ; \cref{p1} ; \cref{t1} ; \cref{r1} ; \cref{c1} ; \cref{c2} ; \cref{t2} FIN
\end{document}
"           deux vrais défauts : doivent tomber
%   lualatex -jobname=trace    "\def\TRACE{1}% Juge de la troisième passe : document minimal écrit de zéro pour la cause A (thmtools sous LaTeX 2026-06-01).
% Variantes, par une macro posée sur la ligne de commande (deux passes chacune, depuis memoire/essai/juge3/minimal/) :
%   lualatex -jobname=nu       "\input{minimal-juge3}"                                   sans correctif
%   lualatex -jobname=corr     "\def\CORR{1}\input{minimal-juge3}"                       correctif A du préambule
%   lualatex -jobname=crochet  "\def\SANSCROCHET{1}\input{minimal-juge3}"                sans correctif, crochet retiré
%   lualatex -jobname=amsseul  "\def\AMSSEUL{1}\input{minimal-juge3}"                    amsthm seul (témoin)
%   lualatex -jobname=faux     "\def\CORR{1}\def\FAUX{1}\input{minimal-juge3}"           deux vrais défauts : doivent tomber
%   lualatex -jobname=trace    "\def\TRACE{1}\input{minimal-juge3}"                      trace de l'appel
% Les lignes du correctif sont tirées du préambule commité par extraire_correctifs.py, pas retapées.
\documentclass{scrbook}
\ifdefined\SANSCROCHET \RemoveFromHook{package/amsthm/after}[firstaid/aliascounter]\fi
\usepackage{amsthm}
\usepackage{hyperref}
\usepackage[french, nameinlink, noabbrev]{cleveref}
\ifdefined\AMSSEUL
  \newtheorem{theoreme}{Théorème}[chapter]
  \newtheorem{proposition}[theoreme]{Proposition}
  \newtheorem{corollaire}[theoreme]{Corollaire}
  \newtheorem{remarque}{Remarque}[chapter]
\else
  \usepackage{thmtools}
  \ifdefined\CORR \input{correctif-a-pose}\fi
  \declaretheorem[numberwithin=chapter, name=Théorème]{theoreme}
  \ifdefined\TRACE \tracingmacros=2 \tracingonline=0 \fi
  \declaretheorem[sibling=theoreme, name=Proposition]{proposition}
  \ifdefined\TRACE \tracingmacros=0 \fi
  \declaretheorem[sibling=theoreme, name=Corollaire]{corollaire}
  \declaretheorem[numberwithin=chapter, name=Remarque]{remarque}
  \ifdefined\FAUX \newcounteralias{corollaire}{section}\newcounteralias{orphelin}{inexistant}\fi
  \ifdefined\CORR \input{correctif-a-remise}\fi
\fi
\crefname{theoreme}{théorème}{théorèmes}
\crefname{proposition}{proposition}{propositions}
\crefname{corollaire}{corollaire}{corollaires}
\crefname{remarque}{remarque}{remarques}
\begin{document}
\chapter{Premier}
\begin{proposition}\label{p1}a\end{proposition}
\begin{theoreme}\label{t1}b\end{theoreme}
\begin{remarque}\label{r1}c\end{remarque}
\begin{corollaire}\label{c1}d\end{corollaire}
\chapter{Second}
\begin{corollaire}\label{c2}e\end{corollaire}
\begin{theoreme}\label{t2}f\end{theoreme}
RENVOIS \ref{p1} \ref{t1} \ref{r1} \ref{c1} \ref{c2} \ref{t2} ; \cref{p1} ; \cref{t1} ; \cref{r1} ; \cref{c1} ; \cref{c2} ; \cref{t2} FIN
\end{document}
"                      trace de l'appel
% Les lignes du correctif sont tirées du préambule commité par extraire_correctifs.py, pas retapées.
\documentclass{scrbook}
\ifdefined\SANSCROCHET \RemoveFromHook{package/amsthm/after}[firstaid/aliascounter]\fi
\usepackage{amsthm}
\usepackage{hyperref}
\usepackage[french, nameinlink, noabbrev]{cleveref}
\ifdefined\AMSSEUL
  \newtheorem{theoreme}{Théorème}[chapter]
  \newtheorem{proposition}[theoreme]{Proposition}
  \newtheorem{corollaire}[theoreme]{Corollaire}
  \newtheorem{remarque}{Remarque}[chapter]
\else
  \usepackage{thmtools}
  \ifdefined\CORR \input{correctif-a-pose}\fi
  \declaretheorem[numberwithin=chapter, name=Théorème]{theoreme}
  \ifdefined\TRACE \tracingmacros=2 \tracingonline=0 \fi
  \declaretheorem[sibling=theoreme, name=Proposition]{proposition}
  \ifdefined\TRACE \tracingmacros=0 \fi
  \declaretheorem[sibling=theoreme, name=Corollaire]{corollaire}
  \declaretheorem[numberwithin=chapter, name=Remarque]{remarque}
  \ifdefined\FAUX \newcounteralias{corollaire}{section}\newcounteralias{orphelin}{inexistant}\fi
  \ifdefined\CORR \input{correctif-a-remise}\fi
\fi
\crefname{theoreme}{théorème}{théorèmes}
\crefname{proposition}{proposition}{propositions}
\crefname{corollaire}{corollaire}{corollaires}
\crefname{remarque}{remarque}{remarques}
\begin{document}
\chapter{Premier}
\begin{proposition}\label{p1}a\end{proposition}
\begin{theoreme}\label{t1}b\end{theoreme}
\begin{remarque}\label{r1}c\end{remarque}
\begin{corollaire}\label{c1}d\end{corollaire}
\chapter{Second}
\begin{corollaire}\label{c2}e\end{corollaire}
\begin{theoreme}\label{t2}f\end{theoreme}
RENVOIS \ref{p1} \ref{t1} \ref{r1} \ref{c1} \ref{c2} \ref{t2} ; \cref{p1} ; \cref{t1} ; \cref{r1} ; \cref{c1} ; \cref{c2} ; \cref{t2} FIN
\end{document}
"                      trace de l'appel
% Les lignes du correctif sont tirées du préambule commité par extraire_correctifs.py, pas retapées.
\documentclass{scrbook}
\ifdefined\SANSCROCHET \RemoveFromHook{package/amsthm/after}[firstaid/aliascounter]\fi
\usepackage{amsthm}
\usepackage{hyperref}
\usepackage[french, nameinlink, noabbrev]{cleveref}
\ifdefined\AMSSEUL
  \newtheorem{theoreme}{Théorème}[chapter]
  \newtheorem{proposition}[theoreme]{Proposition}
  \newtheorem{corollaire}[theoreme]{Corollaire}
  \newtheorem{remarque}{Remarque}[chapter]
\else
  \usepackage{thmtools}
  \ifdefined\CORR \input{correctif-a-pose}\fi
  \declaretheorem[numberwithin=chapter, name=Théorème]{theoreme}
  \ifdefined\TRACE \tracingmacros=2 \tracingonline=0 \fi
  \declaretheorem[sibling=theoreme, name=Proposition]{proposition}
  \ifdefined\TRACE \tracingmacros=0 \fi
  \declaretheorem[sibling=theoreme, name=Corollaire]{corollaire}
  \declaretheorem[numberwithin=chapter, name=Remarque]{remarque}
  \ifdefined\FAUX \newcounteralias{corollaire}{section}\newcounteralias{orphelin}{inexistant}\fi
  \ifdefined\CORR \input{correctif-a-remise}\fi
\fi
\crefname{theoreme}{théorème}{théorèmes}
\crefname{proposition}{proposition}{propositions}
\crefname{corollaire}{corollaire}{corollaires}
\crefname{remarque}{remarque}{remarques}
\begin{document}
\chapter{Premier}
\begin{proposition}\label{p1}a\end{proposition}
\begin{theoreme}\label{t1}b\end{theoreme}
\begin{remarque}\label{r1}c\end{remarque}
\begin{corollaire}\label{c1}d\end{corollaire}
\chapter{Second}
\begin{corollaire}\label{c2}e\end{corollaire}
\begin{theoreme}\label{t2}f\end{theoreme}
RENVOIS \ref{p1} \ref{t1} \ref{r1} \ref{c1} \ref{c2} \ref{t2} ; \cref{p1} ; \cref{t1} ; \cref{r1} ; \cref{c1} ; \cref{c2} ; \cref{t2} FIN
\end{document}
"           deux vrais défauts : doivent tomber
%   lualatex -jobname=trace    "\def\TRACE{1}% Juge de la troisième passe : document minimal écrit de zéro pour la cause A (thmtools sous LaTeX 2026-06-01).
% Variantes, par une macro posée sur la ligne de commande (deux passes chacune, depuis memoire/essai/juge3/minimal/) :
%   lualatex -jobname=nu       "% Juge de la troisième passe : document minimal écrit de zéro pour la cause A (thmtools sous LaTeX 2026-06-01).
% Variantes, par une macro posée sur la ligne de commande (deux passes chacune, depuis memoire/essai/juge3/minimal/) :
%   lualatex -jobname=nu       "% Juge de la troisième passe : document minimal écrit de zéro pour la cause A (thmtools sous LaTeX 2026-06-01).
% Variantes, par une macro posée sur la ligne de commande (deux passes chacune, depuis memoire/essai/juge3/minimal/) :
%   lualatex -jobname=nu       "\input{minimal-juge3}"                                   sans correctif
%   lualatex -jobname=corr     "\def\CORR{1}\input{minimal-juge3}"                       correctif A du préambule
%   lualatex -jobname=crochet  "\def\SANSCROCHET{1}\input{minimal-juge3}"                sans correctif, crochet retiré
%   lualatex -jobname=amsseul  "\def\AMSSEUL{1}\input{minimal-juge3}"                    amsthm seul (témoin)
%   lualatex -jobname=faux     "\def\CORR{1}\def\FAUX{1}\input{minimal-juge3}"           deux vrais défauts : doivent tomber
%   lualatex -jobname=trace    "\def\TRACE{1}\input{minimal-juge3}"                      trace de l'appel
% Les lignes du correctif sont tirées du préambule commité par extraire_correctifs.py, pas retapées.
\documentclass{scrbook}
\ifdefined\SANSCROCHET \RemoveFromHook{package/amsthm/after}[firstaid/aliascounter]\fi
\usepackage{amsthm}
\usepackage{hyperref}
\usepackage[french, nameinlink, noabbrev]{cleveref}
\ifdefined\AMSSEUL
  \newtheorem{theoreme}{Théorème}[chapter]
  \newtheorem{proposition}[theoreme]{Proposition}
  \newtheorem{corollaire}[theoreme]{Corollaire}
  \newtheorem{remarque}{Remarque}[chapter]
\else
  \usepackage{thmtools}
  \ifdefined\CORR \input{correctif-a-pose}\fi
  \declaretheorem[numberwithin=chapter, name=Théorème]{theoreme}
  \ifdefined\TRACE \tracingmacros=2 \tracingonline=0 \fi
  \declaretheorem[sibling=theoreme, name=Proposition]{proposition}
  \ifdefined\TRACE \tracingmacros=0 \fi
  \declaretheorem[sibling=theoreme, name=Corollaire]{corollaire}
  \declaretheorem[numberwithin=chapter, name=Remarque]{remarque}
  \ifdefined\FAUX \newcounteralias{corollaire}{section}\newcounteralias{orphelin}{inexistant}\fi
  \ifdefined\CORR \input{correctif-a-remise}\fi
\fi
\crefname{theoreme}{théorème}{théorèmes}
\crefname{proposition}{proposition}{propositions}
\crefname{corollaire}{corollaire}{corollaires}
\crefname{remarque}{remarque}{remarques}
\begin{document}
\chapter{Premier}
\begin{proposition}\label{p1}a\end{proposition}
\begin{theoreme}\label{t1}b\end{theoreme}
\begin{remarque}\label{r1}c\end{remarque}
\begin{corollaire}\label{c1}d\end{corollaire}
\chapter{Second}
\begin{corollaire}\label{c2}e\end{corollaire}
\begin{theoreme}\label{t2}f\end{theoreme}
RENVOIS \ref{p1} \ref{t1} \ref{r1} \ref{c1} \ref{c2} \ref{t2} ; \cref{p1} ; \cref{t1} ; \cref{r1} ; \cref{c1} ; \cref{c2} ; \cref{t2} FIN
\end{document}
"                                   sans correctif
%   lualatex -jobname=corr     "\def\CORR{1}% Juge de la troisième passe : document minimal écrit de zéro pour la cause A (thmtools sous LaTeX 2026-06-01).
% Variantes, par une macro posée sur la ligne de commande (deux passes chacune, depuis memoire/essai/juge3/minimal/) :
%   lualatex -jobname=nu       "\input{minimal-juge3}"                                   sans correctif
%   lualatex -jobname=corr     "\def\CORR{1}\input{minimal-juge3}"                       correctif A du préambule
%   lualatex -jobname=crochet  "\def\SANSCROCHET{1}\input{minimal-juge3}"                sans correctif, crochet retiré
%   lualatex -jobname=amsseul  "\def\AMSSEUL{1}\input{minimal-juge3}"                    amsthm seul (témoin)
%   lualatex -jobname=faux     "\def\CORR{1}\def\FAUX{1}\input{minimal-juge3}"           deux vrais défauts : doivent tomber
%   lualatex -jobname=trace    "\def\TRACE{1}\input{minimal-juge3}"                      trace de l'appel
% Les lignes du correctif sont tirées du préambule commité par extraire_correctifs.py, pas retapées.
\documentclass{scrbook}
\ifdefined\SANSCROCHET \RemoveFromHook{package/amsthm/after}[firstaid/aliascounter]\fi
\usepackage{amsthm}
\usepackage{hyperref}
\usepackage[french, nameinlink, noabbrev]{cleveref}
\ifdefined\AMSSEUL
  \newtheorem{theoreme}{Théorème}[chapter]
  \newtheorem{proposition}[theoreme]{Proposition}
  \newtheorem{corollaire}[theoreme]{Corollaire}
  \newtheorem{remarque}{Remarque}[chapter]
\else
  \usepackage{thmtools}
  \ifdefined\CORR \input{correctif-a-pose}\fi
  \declaretheorem[numberwithin=chapter, name=Théorème]{theoreme}
  \ifdefined\TRACE \tracingmacros=2 \tracingonline=0 \fi
  \declaretheorem[sibling=theoreme, name=Proposition]{proposition}
  \ifdefined\TRACE \tracingmacros=0 \fi
  \declaretheorem[sibling=theoreme, name=Corollaire]{corollaire}
  \declaretheorem[numberwithin=chapter, name=Remarque]{remarque}
  \ifdefined\FAUX \newcounteralias{corollaire}{section}\newcounteralias{orphelin}{inexistant}\fi
  \ifdefined\CORR \input{correctif-a-remise}\fi
\fi
\crefname{theoreme}{théorème}{théorèmes}
\crefname{proposition}{proposition}{propositions}
\crefname{corollaire}{corollaire}{corollaires}
\crefname{remarque}{remarque}{remarques}
\begin{document}
\chapter{Premier}
\begin{proposition}\label{p1}a\end{proposition}
\begin{theoreme}\label{t1}b\end{theoreme}
\begin{remarque}\label{r1}c\end{remarque}
\begin{corollaire}\label{c1}d\end{corollaire}
\chapter{Second}
\begin{corollaire}\label{c2}e\end{corollaire}
\begin{theoreme}\label{t2}f\end{theoreme}
RENVOIS \ref{p1} \ref{t1} \ref{r1} \ref{c1} \ref{c2} \ref{t2} ; \cref{p1} ; \cref{t1} ; \cref{r1} ; \cref{c1} ; \cref{c2} ; \cref{t2} FIN
\end{document}
"                       correctif A du préambule
%   lualatex -jobname=crochet  "\def\SANSCROCHET{1}% Juge de la troisième passe : document minimal écrit de zéro pour la cause A (thmtools sous LaTeX 2026-06-01).
% Variantes, par une macro posée sur la ligne de commande (deux passes chacune, depuis memoire/essai/juge3/minimal/) :
%   lualatex -jobname=nu       "\input{minimal-juge3}"                                   sans correctif
%   lualatex -jobname=corr     "\def\CORR{1}\input{minimal-juge3}"                       correctif A du préambule
%   lualatex -jobname=crochet  "\def\SANSCROCHET{1}\input{minimal-juge3}"                sans correctif, crochet retiré
%   lualatex -jobname=amsseul  "\def\AMSSEUL{1}\input{minimal-juge3}"                    amsthm seul (témoin)
%   lualatex -jobname=faux     "\def\CORR{1}\def\FAUX{1}\input{minimal-juge3}"           deux vrais défauts : doivent tomber
%   lualatex -jobname=trace    "\def\TRACE{1}\input{minimal-juge3}"                      trace de l'appel
% Les lignes du correctif sont tirées du préambule commité par extraire_correctifs.py, pas retapées.
\documentclass{scrbook}
\ifdefined\SANSCROCHET \RemoveFromHook{package/amsthm/after}[firstaid/aliascounter]\fi
\usepackage{amsthm}
\usepackage{hyperref}
\usepackage[french, nameinlink, noabbrev]{cleveref}
\ifdefined\AMSSEUL
  \newtheorem{theoreme}{Théorème}[chapter]
  \newtheorem{proposition}[theoreme]{Proposition}
  \newtheorem{corollaire}[theoreme]{Corollaire}
  \newtheorem{remarque}{Remarque}[chapter]
\else
  \usepackage{thmtools}
  \ifdefined\CORR \input{correctif-a-pose}\fi
  \declaretheorem[numberwithin=chapter, name=Théorème]{theoreme}
  \ifdefined\TRACE \tracingmacros=2 \tracingonline=0 \fi
  \declaretheorem[sibling=theoreme, name=Proposition]{proposition}
  \ifdefined\TRACE \tracingmacros=0 \fi
  \declaretheorem[sibling=theoreme, name=Corollaire]{corollaire}
  \declaretheorem[numberwithin=chapter, name=Remarque]{remarque}
  \ifdefined\FAUX \newcounteralias{corollaire}{section}\newcounteralias{orphelin}{inexistant}\fi
  \ifdefined\CORR \input{correctif-a-remise}\fi
\fi
\crefname{theoreme}{théorème}{théorèmes}
\crefname{proposition}{proposition}{propositions}
\crefname{corollaire}{corollaire}{corollaires}
\crefname{remarque}{remarque}{remarques}
\begin{document}
\chapter{Premier}
\begin{proposition}\label{p1}a\end{proposition}
\begin{theoreme}\label{t1}b\end{theoreme}
\begin{remarque}\label{r1}c\end{remarque}
\begin{corollaire}\label{c1}d\end{corollaire}
\chapter{Second}
\begin{corollaire}\label{c2}e\end{corollaire}
\begin{theoreme}\label{t2}f\end{theoreme}
RENVOIS \ref{p1} \ref{t1} \ref{r1} \ref{c1} \ref{c2} \ref{t2} ; \cref{p1} ; \cref{t1} ; \cref{r1} ; \cref{c1} ; \cref{c2} ; \cref{t2} FIN
\end{document}
"                sans correctif, crochet retiré
%   lualatex -jobname=amsseul  "\def\AMSSEUL{1}% Juge de la troisième passe : document minimal écrit de zéro pour la cause A (thmtools sous LaTeX 2026-06-01).
% Variantes, par une macro posée sur la ligne de commande (deux passes chacune, depuis memoire/essai/juge3/minimal/) :
%   lualatex -jobname=nu       "\input{minimal-juge3}"                                   sans correctif
%   lualatex -jobname=corr     "\def\CORR{1}\input{minimal-juge3}"                       correctif A du préambule
%   lualatex -jobname=crochet  "\def\SANSCROCHET{1}\input{minimal-juge3}"                sans correctif, crochet retiré
%   lualatex -jobname=amsseul  "\def\AMSSEUL{1}\input{minimal-juge3}"                    amsthm seul (témoin)
%   lualatex -jobname=faux     "\def\CORR{1}\def\FAUX{1}\input{minimal-juge3}"           deux vrais défauts : doivent tomber
%   lualatex -jobname=trace    "\def\TRACE{1}\input{minimal-juge3}"                      trace de l'appel
% Les lignes du correctif sont tirées du préambule commité par extraire_correctifs.py, pas retapées.
\documentclass{scrbook}
\ifdefined\SANSCROCHET \RemoveFromHook{package/amsthm/after}[firstaid/aliascounter]\fi
\usepackage{amsthm}
\usepackage{hyperref}
\usepackage[french, nameinlink, noabbrev]{cleveref}
\ifdefined\AMSSEUL
  \newtheorem{theoreme}{Théorème}[chapter]
  \newtheorem{proposition}[theoreme]{Proposition}
  \newtheorem{corollaire}[theoreme]{Corollaire}
  \newtheorem{remarque}{Remarque}[chapter]
\else
  \usepackage{thmtools}
  \ifdefined\CORR \input{correctif-a-pose}\fi
  \declaretheorem[numberwithin=chapter, name=Théorème]{theoreme}
  \ifdefined\TRACE \tracingmacros=2 \tracingonline=0 \fi
  \declaretheorem[sibling=theoreme, name=Proposition]{proposition}
  \ifdefined\TRACE \tracingmacros=0 \fi
  \declaretheorem[sibling=theoreme, name=Corollaire]{corollaire}
  \declaretheorem[numberwithin=chapter, name=Remarque]{remarque}
  \ifdefined\FAUX \newcounteralias{corollaire}{section}\newcounteralias{orphelin}{inexistant}\fi
  \ifdefined\CORR \input{correctif-a-remise}\fi
\fi
\crefname{theoreme}{théorème}{théorèmes}
\crefname{proposition}{proposition}{propositions}
\crefname{corollaire}{corollaire}{corollaires}
\crefname{remarque}{remarque}{remarques}
\begin{document}
\chapter{Premier}
\begin{proposition}\label{p1}a\end{proposition}
\begin{theoreme}\label{t1}b\end{theoreme}
\begin{remarque}\label{r1}c\end{remarque}
\begin{corollaire}\label{c1}d\end{corollaire}
\chapter{Second}
\begin{corollaire}\label{c2}e\end{corollaire}
\begin{theoreme}\label{t2}f\end{theoreme}
RENVOIS \ref{p1} \ref{t1} \ref{r1} \ref{c1} \ref{c2} \ref{t2} ; \cref{p1} ; \cref{t1} ; \cref{r1} ; \cref{c1} ; \cref{c2} ; \cref{t2} FIN
\end{document}
"                    amsthm seul (témoin)
%   lualatex -jobname=faux     "\def\CORR{1}\def\FAUX{1}% Juge de la troisième passe : document minimal écrit de zéro pour la cause A (thmtools sous LaTeX 2026-06-01).
% Variantes, par une macro posée sur la ligne de commande (deux passes chacune, depuis memoire/essai/juge3/minimal/) :
%   lualatex -jobname=nu       "\input{minimal-juge3}"                                   sans correctif
%   lualatex -jobname=corr     "\def\CORR{1}\input{minimal-juge3}"                       correctif A du préambule
%   lualatex -jobname=crochet  "\def\SANSCROCHET{1}\input{minimal-juge3}"                sans correctif, crochet retiré
%   lualatex -jobname=amsseul  "\def\AMSSEUL{1}\input{minimal-juge3}"                    amsthm seul (témoin)
%   lualatex -jobname=faux     "\def\CORR{1}\def\FAUX{1}\input{minimal-juge3}"           deux vrais défauts : doivent tomber
%   lualatex -jobname=trace    "\def\TRACE{1}\input{minimal-juge3}"                      trace de l'appel
% Les lignes du correctif sont tirées du préambule commité par extraire_correctifs.py, pas retapées.
\documentclass{scrbook}
\ifdefined\SANSCROCHET \RemoveFromHook{package/amsthm/after}[firstaid/aliascounter]\fi
\usepackage{amsthm}
\usepackage{hyperref}
\usepackage[french, nameinlink, noabbrev]{cleveref}
\ifdefined\AMSSEUL
  \newtheorem{theoreme}{Théorème}[chapter]
  \newtheorem{proposition}[theoreme]{Proposition}
  \newtheorem{corollaire}[theoreme]{Corollaire}
  \newtheorem{remarque}{Remarque}[chapter]
\else
  \usepackage{thmtools}
  \ifdefined\CORR \input{correctif-a-pose}\fi
  \declaretheorem[numberwithin=chapter, name=Théorème]{theoreme}
  \ifdefined\TRACE \tracingmacros=2 \tracingonline=0 \fi
  \declaretheorem[sibling=theoreme, name=Proposition]{proposition}
  \ifdefined\TRACE \tracingmacros=0 \fi
  \declaretheorem[sibling=theoreme, name=Corollaire]{corollaire}
  \declaretheorem[numberwithin=chapter, name=Remarque]{remarque}
  \ifdefined\FAUX \newcounteralias{corollaire}{section}\newcounteralias{orphelin}{inexistant}\fi
  \ifdefined\CORR \input{correctif-a-remise}\fi
\fi
\crefname{theoreme}{théorème}{théorèmes}
\crefname{proposition}{proposition}{propositions}
\crefname{corollaire}{corollaire}{corollaires}
\crefname{remarque}{remarque}{remarques}
\begin{document}
\chapter{Premier}
\begin{proposition}\label{p1}a\end{proposition}
\begin{theoreme}\label{t1}b\end{theoreme}
\begin{remarque}\label{r1}c\end{remarque}
\begin{corollaire}\label{c1}d\end{corollaire}
\chapter{Second}
\begin{corollaire}\label{c2}e\end{corollaire}
\begin{theoreme}\label{t2}f\end{theoreme}
RENVOIS \ref{p1} \ref{t1} \ref{r1} \ref{c1} \ref{c2} \ref{t2} ; \cref{p1} ; \cref{t1} ; \cref{r1} ; \cref{c1} ; \cref{c2} ; \cref{t2} FIN
\end{document}
"           deux vrais défauts : doivent tomber
%   lualatex -jobname=trace    "\def\TRACE{1}% Juge de la troisième passe : document minimal écrit de zéro pour la cause A (thmtools sous LaTeX 2026-06-01).
% Variantes, par une macro posée sur la ligne de commande (deux passes chacune, depuis memoire/essai/juge3/minimal/) :
%   lualatex -jobname=nu       "\input{minimal-juge3}"                                   sans correctif
%   lualatex -jobname=corr     "\def\CORR{1}\input{minimal-juge3}"                       correctif A du préambule
%   lualatex -jobname=crochet  "\def\SANSCROCHET{1}\input{minimal-juge3}"                sans correctif, crochet retiré
%   lualatex -jobname=amsseul  "\def\AMSSEUL{1}\input{minimal-juge3}"                    amsthm seul (témoin)
%   lualatex -jobname=faux     "\def\CORR{1}\def\FAUX{1}\input{minimal-juge3}"           deux vrais défauts : doivent tomber
%   lualatex -jobname=trace    "\def\TRACE{1}\input{minimal-juge3}"                      trace de l'appel
% Les lignes du correctif sont tirées du préambule commité par extraire_correctifs.py, pas retapées.
\documentclass{scrbook}
\ifdefined\SANSCROCHET \RemoveFromHook{package/amsthm/after}[firstaid/aliascounter]\fi
\usepackage{amsthm}
\usepackage{hyperref}
\usepackage[french, nameinlink, noabbrev]{cleveref}
\ifdefined\AMSSEUL
  \newtheorem{theoreme}{Théorème}[chapter]
  \newtheorem{proposition}[theoreme]{Proposition}
  \newtheorem{corollaire}[theoreme]{Corollaire}
  \newtheorem{remarque}{Remarque}[chapter]
\else
  \usepackage{thmtools}
  \ifdefined\CORR \input{correctif-a-pose}\fi
  \declaretheorem[numberwithin=chapter, name=Théorème]{theoreme}
  \ifdefined\TRACE \tracingmacros=2 \tracingonline=0 \fi
  \declaretheorem[sibling=theoreme, name=Proposition]{proposition}
  \ifdefined\TRACE \tracingmacros=0 \fi
  \declaretheorem[sibling=theoreme, name=Corollaire]{corollaire}
  \declaretheorem[numberwithin=chapter, name=Remarque]{remarque}
  \ifdefined\FAUX \newcounteralias{corollaire}{section}\newcounteralias{orphelin}{inexistant}\fi
  \ifdefined\CORR \input{correctif-a-remise}\fi
\fi
\crefname{theoreme}{théorème}{théorèmes}
\crefname{proposition}{proposition}{propositions}
\crefname{corollaire}{corollaire}{corollaires}
\crefname{remarque}{remarque}{remarques}
\begin{document}
\chapter{Premier}
\begin{proposition}\label{p1}a\end{proposition}
\begin{theoreme}\label{t1}b\end{theoreme}
\begin{remarque}\label{r1}c\end{remarque}
\begin{corollaire}\label{c1}d\end{corollaire}
\chapter{Second}
\begin{corollaire}\label{c2}e\end{corollaire}
\begin{theoreme}\label{t2}f\end{theoreme}
RENVOIS \ref{p1} \ref{t1} \ref{r1} \ref{c1} \ref{c2} \ref{t2} ; \cref{p1} ; \cref{t1} ; \cref{r1} ; \cref{c1} ; \cref{c2} ; \cref{t2} FIN
\end{document}
"                      trace de l'appel
% Les lignes du correctif sont tirées du préambule commité par extraire_correctifs.py, pas retapées.
\documentclass{scrbook}
\ifdefined\SANSCROCHET \RemoveFromHook{package/amsthm/after}[firstaid/aliascounter]\fi
\usepackage{amsthm}
\usepackage{hyperref}
\usepackage[french, nameinlink, noabbrev]{cleveref}
\ifdefined\AMSSEUL
  \newtheorem{theoreme}{Théorème}[chapter]
  \newtheorem{proposition}[theoreme]{Proposition}
  \newtheorem{corollaire}[theoreme]{Corollaire}
  \newtheorem{remarque}{Remarque}[chapter]
\else
  \usepackage{thmtools}
  \ifdefined\CORR \input{correctif-a-pose}\fi
  \declaretheorem[numberwithin=chapter, name=Théorème]{theoreme}
  \ifdefined\TRACE \tracingmacros=2 \tracingonline=0 \fi
  \declaretheorem[sibling=theoreme, name=Proposition]{proposition}
  \ifdefined\TRACE \tracingmacros=0 \fi
  \declaretheorem[sibling=theoreme, name=Corollaire]{corollaire}
  \declaretheorem[numberwithin=chapter, name=Remarque]{remarque}
  \ifdefined\FAUX \newcounteralias{corollaire}{section}\newcounteralias{orphelin}{inexistant}\fi
  \ifdefined\CORR \input{correctif-a-remise}\fi
\fi
\crefname{theoreme}{théorème}{théorèmes}
\crefname{proposition}{proposition}{propositions}
\crefname{corollaire}{corollaire}{corollaires}
\crefname{remarque}{remarque}{remarques}
\begin{document}
\chapter{Premier}
\begin{proposition}\label{p1}a\end{proposition}
\begin{theoreme}\label{t1}b\end{theoreme}
\begin{remarque}\label{r1}c\end{remarque}
\begin{corollaire}\label{c1}d\end{corollaire}
\chapter{Second}
\begin{corollaire}\label{c2}e\end{corollaire}
\begin{theoreme}\label{t2}f\end{theoreme}
RENVOIS \ref{p1} \ref{t1} \ref{r1} \ref{c1} \ref{c2} \ref{t2} ; \cref{p1} ; \cref{t1} ; \cref{r1} ; \cref{c1} ; \cref{c2} ; \cref{t2} FIN
\end{document}
"                                   sans correctif
%   lualatex -jobname=corr     "\def\CORR{1}% Juge de la troisième passe : document minimal écrit de zéro pour la cause A (thmtools sous LaTeX 2026-06-01).
% Variantes, par une macro posée sur la ligne de commande (deux passes chacune, depuis memoire/essai/juge3/minimal/) :
%   lualatex -jobname=nu       "% Juge de la troisième passe : document minimal écrit de zéro pour la cause A (thmtools sous LaTeX 2026-06-01).
% Variantes, par une macro posée sur la ligne de commande (deux passes chacune, depuis memoire/essai/juge3/minimal/) :
%   lualatex -jobname=nu       "\input{minimal-juge3}"                                   sans correctif
%   lualatex -jobname=corr     "\def\CORR{1}\input{minimal-juge3}"                       correctif A du préambule
%   lualatex -jobname=crochet  "\def\SANSCROCHET{1}\input{minimal-juge3}"                sans correctif, crochet retiré
%   lualatex -jobname=amsseul  "\def\AMSSEUL{1}\input{minimal-juge3}"                    amsthm seul (témoin)
%   lualatex -jobname=faux     "\def\CORR{1}\def\FAUX{1}\input{minimal-juge3}"           deux vrais défauts : doivent tomber
%   lualatex -jobname=trace    "\def\TRACE{1}\input{minimal-juge3}"                      trace de l'appel
% Les lignes du correctif sont tirées du préambule commité par extraire_correctifs.py, pas retapées.
\documentclass{scrbook}
\ifdefined\SANSCROCHET \RemoveFromHook{package/amsthm/after}[firstaid/aliascounter]\fi
\usepackage{amsthm}
\usepackage{hyperref}
\usepackage[french, nameinlink, noabbrev]{cleveref}
\ifdefined\AMSSEUL
  \newtheorem{theoreme}{Théorème}[chapter]
  \newtheorem{proposition}[theoreme]{Proposition}
  \newtheorem{corollaire}[theoreme]{Corollaire}
  \newtheorem{remarque}{Remarque}[chapter]
\else
  \usepackage{thmtools}
  \ifdefined\CORR \input{correctif-a-pose}\fi
  \declaretheorem[numberwithin=chapter, name=Théorème]{theoreme}
  \ifdefined\TRACE \tracingmacros=2 \tracingonline=0 \fi
  \declaretheorem[sibling=theoreme, name=Proposition]{proposition}
  \ifdefined\TRACE \tracingmacros=0 \fi
  \declaretheorem[sibling=theoreme, name=Corollaire]{corollaire}
  \declaretheorem[numberwithin=chapter, name=Remarque]{remarque}
  \ifdefined\FAUX \newcounteralias{corollaire}{section}\newcounteralias{orphelin}{inexistant}\fi
  \ifdefined\CORR \input{correctif-a-remise}\fi
\fi
\crefname{theoreme}{théorème}{théorèmes}
\crefname{proposition}{proposition}{propositions}
\crefname{corollaire}{corollaire}{corollaires}
\crefname{remarque}{remarque}{remarques}
\begin{document}
\chapter{Premier}
\begin{proposition}\label{p1}a\end{proposition}
\begin{theoreme}\label{t1}b\end{theoreme}
\begin{remarque}\label{r1}c\end{remarque}
\begin{corollaire}\label{c1}d\end{corollaire}
\chapter{Second}
\begin{corollaire}\label{c2}e\end{corollaire}
\begin{theoreme}\label{t2}f\end{theoreme}
RENVOIS \ref{p1} \ref{t1} \ref{r1} \ref{c1} \ref{c2} \ref{t2} ; \cref{p1} ; \cref{t1} ; \cref{r1} ; \cref{c1} ; \cref{c2} ; \cref{t2} FIN
\end{document}
"                                   sans correctif
%   lualatex -jobname=corr     "\def\CORR{1}% Juge de la troisième passe : document minimal écrit de zéro pour la cause A (thmtools sous LaTeX 2026-06-01).
% Variantes, par une macro posée sur la ligne de commande (deux passes chacune, depuis memoire/essai/juge3/minimal/) :
%   lualatex -jobname=nu       "\input{minimal-juge3}"                                   sans correctif
%   lualatex -jobname=corr     "\def\CORR{1}\input{minimal-juge3}"                       correctif A du préambule
%   lualatex -jobname=crochet  "\def\SANSCROCHET{1}\input{minimal-juge3}"                sans correctif, crochet retiré
%   lualatex -jobname=amsseul  "\def\AMSSEUL{1}\input{minimal-juge3}"                    amsthm seul (témoin)
%   lualatex -jobname=faux     "\def\CORR{1}\def\FAUX{1}\input{minimal-juge3}"           deux vrais défauts : doivent tomber
%   lualatex -jobname=trace    "\def\TRACE{1}\input{minimal-juge3}"                      trace de l'appel
% Les lignes du correctif sont tirées du préambule commité par extraire_correctifs.py, pas retapées.
\documentclass{scrbook}
\ifdefined\SANSCROCHET \RemoveFromHook{package/amsthm/after}[firstaid/aliascounter]\fi
\usepackage{amsthm}
\usepackage{hyperref}
\usepackage[french, nameinlink, noabbrev]{cleveref}
\ifdefined\AMSSEUL
  \newtheorem{theoreme}{Théorème}[chapter]
  \newtheorem{proposition}[theoreme]{Proposition}
  \newtheorem{corollaire}[theoreme]{Corollaire}
  \newtheorem{remarque}{Remarque}[chapter]
\else
  \usepackage{thmtools}
  \ifdefined\CORR \input{correctif-a-pose}\fi
  \declaretheorem[numberwithin=chapter, name=Théorème]{theoreme}
  \ifdefined\TRACE \tracingmacros=2 \tracingonline=0 \fi
  \declaretheorem[sibling=theoreme, name=Proposition]{proposition}
  \ifdefined\TRACE \tracingmacros=0 \fi
  \declaretheorem[sibling=theoreme, name=Corollaire]{corollaire}
  \declaretheorem[numberwithin=chapter, name=Remarque]{remarque}
  \ifdefined\FAUX \newcounteralias{corollaire}{section}\newcounteralias{orphelin}{inexistant}\fi
  \ifdefined\CORR \input{correctif-a-remise}\fi
\fi
\crefname{theoreme}{théorème}{théorèmes}
\crefname{proposition}{proposition}{propositions}
\crefname{corollaire}{corollaire}{corollaires}
\crefname{remarque}{remarque}{remarques}
\begin{document}
\chapter{Premier}
\begin{proposition}\label{p1}a\end{proposition}
\begin{theoreme}\label{t1}b\end{theoreme}
\begin{remarque}\label{r1}c\end{remarque}
\begin{corollaire}\label{c1}d\end{corollaire}
\chapter{Second}
\begin{corollaire}\label{c2}e\end{corollaire}
\begin{theoreme}\label{t2}f\end{theoreme}
RENVOIS \ref{p1} \ref{t1} \ref{r1} \ref{c1} \ref{c2} \ref{t2} ; \cref{p1} ; \cref{t1} ; \cref{r1} ; \cref{c1} ; \cref{c2} ; \cref{t2} FIN
\end{document}
"                       correctif A du préambule
%   lualatex -jobname=crochet  "\def\SANSCROCHET{1}% Juge de la troisième passe : document minimal écrit de zéro pour la cause A (thmtools sous LaTeX 2026-06-01).
% Variantes, par une macro posée sur la ligne de commande (deux passes chacune, depuis memoire/essai/juge3/minimal/) :
%   lualatex -jobname=nu       "\input{minimal-juge3}"                                   sans correctif
%   lualatex -jobname=corr     "\def\CORR{1}\input{minimal-juge3}"                       correctif A du préambule
%   lualatex -jobname=crochet  "\def\SANSCROCHET{1}\input{minimal-juge3}"                sans correctif, crochet retiré
%   lualatex -jobname=amsseul  "\def\AMSSEUL{1}\input{minimal-juge3}"                    amsthm seul (témoin)
%   lualatex -jobname=faux     "\def\CORR{1}\def\FAUX{1}\input{minimal-juge3}"           deux vrais défauts : doivent tomber
%   lualatex -jobname=trace    "\def\TRACE{1}\input{minimal-juge3}"                      trace de l'appel
% Les lignes du correctif sont tirées du préambule commité par extraire_correctifs.py, pas retapées.
\documentclass{scrbook}
\ifdefined\SANSCROCHET \RemoveFromHook{package/amsthm/after}[firstaid/aliascounter]\fi
\usepackage{amsthm}
\usepackage{hyperref}
\usepackage[french, nameinlink, noabbrev]{cleveref}
\ifdefined\AMSSEUL
  \newtheorem{theoreme}{Théorème}[chapter]
  \newtheorem{proposition}[theoreme]{Proposition}
  \newtheorem{corollaire}[theoreme]{Corollaire}
  \newtheorem{remarque}{Remarque}[chapter]
\else
  \usepackage{thmtools}
  \ifdefined\CORR \input{correctif-a-pose}\fi
  \declaretheorem[numberwithin=chapter, name=Théorème]{theoreme}
  \ifdefined\TRACE \tracingmacros=2 \tracingonline=0 \fi
  \declaretheorem[sibling=theoreme, name=Proposition]{proposition}
  \ifdefined\TRACE \tracingmacros=0 \fi
  \declaretheorem[sibling=theoreme, name=Corollaire]{corollaire}
  \declaretheorem[numberwithin=chapter, name=Remarque]{remarque}
  \ifdefined\FAUX \newcounteralias{corollaire}{section}\newcounteralias{orphelin}{inexistant}\fi
  \ifdefined\CORR \input{correctif-a-remise}\fi
\fi
\crefname{theoreme}{théorème}{théorèmes}
\crefname{proposition}{proposition}{propositions}
\crefname{corollaire}{corollaire}{corollaires}
\crefname{remarque}{remarque}{remarques}
\begin{document}
\chapter{Premier}
\begin{proposition}\label{p1}a\end{proposition}
\begin{theoreme}\label{t1}b\end{theoreme}
\begin{remarque}\label{r1}c\end{remarque}
\begin{corollaire}\label{c1}d\end{corollaire}
\chapter{Second}
\begin{corollaire}\label{c2}e\end{corollaire}
\begin{theoreme}\label{t2}f\end{theoreme}
RENVOIS \ref{p1} \ref{t1} \ref{r1} \ref{c1} \ref{c2} \ref{t2} ; \cref{p1} ; \cref{t1} ; \cref{r1} ; \cref{c1} ; \cref{c2} ; \cref{t2} FIN
\end{document}
"                sans correctif, crochet retiré
%   lualatex -jobname=amsseul  "\def\AMSSEUL{1}% Juge de la troisième passe : document minimal écrit de zéro pour la cause A (thmtools sous LaTeX 2026-06-01).
% Variantes, par une macro posée sur la ligne de commande (deux passes chacune, depuis memoire/essai/juge3/minimal/) :
%   lualatex -jobname=nu       "\input{minimal-juge3}"                                   sans correctif
%   lualatex -jobname=corr     "\def\CORR{1}\input{minimal-juge3}"                       correctif A du préambule
%   lualatex -jobname=crochet  "\def\SANSCROCHET{1}\input{minimal-juge3}"                sans correctif, crochet retiré
%   lualatex -jobname=amsseul  "\def\AMSSEUL{1}\input{minimal-juge3}"                    amsthm seul (témoin)
%   lualatex -jobname=faux     "\def\CORR{1}\def\FAUX{1}\input{minimal-juge3}"           deux vrais défauts : doivent tomber
%   lualatex -jobname=trace    "\def\TRACE{1}\input{minimal-juge3}"                      trace de l'appel
% Les lignes du correctif sont tirées du préambule commité par extraire_correctifs.py, pas retapées.
\documentclass{scrbook}
\ifdefined\SANSCROCHET \RemoveFromHook{package/amsthm/after}[firstaid/aliascounter]\fi
\usepackage{amsthm}
\usepackage{hyperref}
\usepackage[french, nameinlink, noabbrev]{cleveref}
\ifdefined\AMSSEUL
  \newtheorem{theoreme}{Théorème}[chapter]
  \newtheorem{proposition}[theoreme]{Proposition}
  \newtheorem{corollaire}[theoreme]{Corollaire}
  \newtheorem{remarque}{Remarque}[chapter]
\else
  \usepackage{thmtools}
  \ifdefined\CORR \input{correctif-a-pose}\fi
  \declaretheorem[numberwithin=chapter, name=Théorème]{theoreme}
  \ifdefined\TRACE \tracingmacros=2 \tracingonline=0 \fi
  \declaretheorem[sibling=theoreme, name=Proposition]{proposition}
  \ifdefined\TRACE \tracingmacros=0 \fi
  \declaretheorem[sibling=theoreme, name=Corollaire]{corollaire}
  \declaretheorem[numberwithin=chapter, name=Remarque]{remarque}
  \ifdefined\FAUX \newcounteralias{corollaire}{section}\newcounteralias{orphelin}{inexistant}\fi
  \ifdefined\CORR \input{correctif-a-remise}\fi
\fi
\crefname{theoreme}{théorème}{théorèmes}
\crefname{proposition}{proposition}{propositions}
\crefname{corollaire}{corollaire}{corollaires}
\crefname{remarque}{remarque}{remarques}
\begin{document}
\chapter{Premier}
\begin{proposition}\label{p1}a\end{proposition}
\begin{theoreme}\label{t1}b\end{theoreme}
\begin{remarque}\label{r1}c\end{remarque}
\begin{corollaire}\label{c1}d\end{corollaire}
\chapter{Second}
\begin{corollaire}\label{c2}e\end{corollaire}
\begin{theoreme}\label{t2}f\end{theoreme}
RENVOIS \ref{p1} \ref{t1} \ref{r1} \ref{c1} \ref{c2} \ref{t2} ; \cref{p1} ; \cref{t1} ; \cref{r1} ; \cref{c1} ; \cref{c2} ; \cref{t2} FIN
\end{document}
"                    amsthm seul (témoin)
%   lualatex -jobname=faux     "\def\CORR{1}\def\FAUX{1}% Juge de la troisième passe : document minimal écrit de zéro pour la cause A (thmtools sous LaTeX 2026-06-01).
% Variantes, par une macro posée sur la ligne de commande (deux passes chacune, depuis memoire/essai/juge3/minimal/) :
%   lualatex -jobname=nu       "\input{minimal-juge3}"                                   sans correctif
%   lualatex -jobname=corr     "\def\CORR{1}\input{minimal-juge3}"                       correctif A du préambule
%   lualatex -jobname=crochet  "\def\SANSCROCHET{1}\input{minimal-juge3}"                sans correctif, crochet retiré
%   lualatex -jobname=amsseul  "\def\AMSSEUL{1}\input{minimal-juge3}"                    amsthm seul (témoin)
%   lualatex -jobname=faux     "\def\CORR{1}\def\FAUX{1}\input{minimal-juge3}"           deux vrais défauts : doivent tomber
%   lualatex -jobname=trace    "\def\TRACE{1}\input{minimal-juge3}"                      trace de l'appel
% Les lignes du correctif sont tirées du préambule commité par extraire_correctifs.py, pas retapées.
\documentclass{scrbook}
\ifdefined\SANSCROCHET \RemoveFromHook{package/amsthm/after}[firstaid/aliascounter]\fi
\usepackage{amsthm}
\usepackage{hyperref}
\usepackage[french, nameinlink, noabbrev]{cleveref}
\ifdefined\AMSSEUL
  \newtheorem{theoreme}{Théorème}[chapter]
  \newtheorem{proposition}[theoreme]{Proposition}
  \newtheorem{corollaire}[theoreme]{Corollaire}
  \newtheorem{remarque}{Remarque}[chapter]
\else
  \usepackage{thmtools}
  \ifdefined\CORR \input{correctif-a-pose}\fi
  \declaretheorem[numberwithin=chapter, name=Théorème]{theoreme}
  \ifdefined\TRACE \tracingmacros=2 \tracingonline=0 \fi
  \declaretheorem[sibling=theoreme, name=Proposition]{proposition}
  \ifdefined\TRACE \tracingmacros=0 \fi
  \declaretheorem[sibling=theoreme, name=Corollaire]{corollaire}
  \declaretheorem[numberwithin=chapter, name=Remarque]{remarque}
  \ifdefined\FAUX \newcounteralias{corollaire}{section}\newcounteralias{orphelin}{inexistant}\fi
  \ifdefined\CORR \input{correctif-a-remise}\fi
\fi
\crefname{theoreme}{théorème}{théorèmes}
\crefname{proposition}{proposition}{propositions}
\crefname{corollaire}{corollaire}{corollaires}
\crefname{remarque}{remarque}{remarques}
\begin{document}
\chapter{Premier}
\begin{proposition}\label{p1}a\end{proposition}
\begin{theoreme}\label{t1}b\end{theoreme}
\begin{remarque}\label{r1}c\end{remarque}
\begin{corollaire}\label{c1}d\end{corollaire}
\chapter{Second}
\begin{corollaire}\label{c2}e\end{corollaire}
\begin{theoreme}\label{t2}f\end{theoreme}
RENVOIS \ref{p1} \ref{t1} \ref{r1} \ref{c1} \ref{c2} \ref{t2} ; \cref{p1} ; \cref{t1} ; \cref{r1} ; \cref{c1} ; \cref{c2} ; \cref{t2} FIN
\end{document}
"           deux vrais défauts : doivent tomber
%   lualatex -jobname=trace    "\def\TRACE{1}% Juge de la troisième passe : document minimal écrit de zéro pour la cause A (thmtools sous LaTeX 2026-06-01).
% Variantes, par une macro posée sur la ligne de commande (deux passes chacune, depuis memoire/essai/juge3/minimal/) :
%   lualatex -jobname=nu       "\input{minimal-juge3}"                                   sans correctif
%   lualatex -jobname=corr     "\def\CORR{1}\input{minimal-juge3}"                       correctif A du préambule
%   lualatex -jobname=crochet  "\def\SANSCROCHET{1}\input{minimal-juge3}"                sans correctif, crochet retiré
%   lualatex -jobname=amsseul  "\def\AMSSEUL{1}\input{minimal-juge3}"                    amsthm seul (témoin)
%   lualatex -jobname=faux     "\def\CORR{1}\def\FAUX{1}\input{minimal-juge3}"           deux vrais défauts : doivent tomber
%   lualatex -jobname=trace    "\def\TRACE{1}\input{minimal-juge3}"                      trace de l'appel
% Les lignes du correctif sont tirées du préambule commité par extraire_correctifs.py, pas retapées.
\documentclass{scrbook}
\ifdefined\SANSCROCHET \RemoveFromHook{package/amsthm/after}[firstaid/aliascounter]\fi
\usepackage{amsthm}
\usepackage{hyperref}
\usepackage[french, nameinlink, noabbrev]{cleveref}
\ifdefined\AMSSEUL
  \newtheorem{theoreme}{Théorème}[chapter]
  \newtheorem{proposition}[theoreme]{Proposition}
  \newtheorem{corollaire}[theoreme]{Corollaire}
  \newtheorem{remarque}{Remarque}[chapter]
\else
  \usepackage{thmtools}
  \ifdefined\CORR \input{correctif-a-pose}\fi
  \declaretheorem[numberwithin=chapter, name=Théorème]{theoreme}
  \ifdefined\TRACE \tracingmacros=2 \tracingonline=0 \fi
  \declaretheorem[sibling=theoreme, name=Proposition]{proposition}
  \ifdefined\TRACE \tracingmacros=0 \fi
  \declaretheorem[sibling=theoreme, name=Corollaire]{corollaire}
  \declaretheorem[numberwithin=chapter, name=Remarque]{remarque}
  \ifdefined\FAUX \newcounteralias{corollaire}{section}\newcounteralias{orphelin}{inexistant}\fi
  \ifdefined\CORR \input{correctif-a-remise}\fi
\fi
\crefname{theoreme}{théorème}{théorèmes}
\crefname{proposition}{proposition}{propositions}
\crefname{corollaire}{corollaire}{corollaires}
\crefname{remarque}{remarque}{remarques}
\begin{document}
\chapter{Premier}
\begin{proposition}\label{p1}a\end{proposition}
\begin{theoreme}\label{t1}b\end{theoreme}
\begin{remarque}\label{r1}c\end{remarque}
\begin{corollaire}\label{c1}d\end{corollaire}
\chapter{Second}
\begin{corollaire}\label{c2}e\end{corollaire}
\begin{theoreme}\label{t2}f\end{theoreme}
RENVOIS \ref{p1} \ref{t1} \ref{r1} \ref{c1} \ref{c2} \ref{t2} ; \cref{p1} ; \cref{t1} ; \cref{r1} ; \cref{c1} ; \cref{c2} ; \cref{t2} FIN
\end{document}
"                      trace de l'appel
% Les lignes du correctif sont tirées du préambule commité par extraire_correctifs.py, pas retapées.
\documentclass{scrbook}
\ifdefined\SANSCROCHET \RemoveFromHook{package/amsthm/after}[firstaid/aliascounter]\fi
\usepackage{amsthm}
\usepackage{hyperref}
\usepackage[french, nameinlink, noabbrev]{cleveref}
\ifdefined\AMSSEUL
  \newtheorem{theoreme}{Théorème}[chapter]
  \newtheorem{proposition}[theoreme]{Proposition}
  \newtheorem{corollaire}[theoreme]{Corollaire}
  \newtheorem{remarque}{Remarque}[chapter]
\else
  \usepackage{thmtools}
  \ifdefined\CORR \input{correctif-a-pose}\fi
  \declaretheorem[numberwithin=chapter, name=Théorème]{theoreme}
  \ifdefined\TRACE \tracingmacros=2 \tracingonline=0 \fi
  \declaretheorem[sibling=theoreme, name=Proposition]{proposition}
  \ifdefined\TRACE \tracingmacros=0 \fi
  \declaretheorem[sibling=theoreme, name=Corollaire]{corollaire}
  \declaretheorem[numberwithin=chapter, name=Remarque]{remarque}
  \ifdefined\FAUX \newcounteralias{corollaire}{section}\newcounteralias{orphelin}{inexistant}\fi
  \ifdefined\CORR \input{correctif-a-remise}\fi
\fi
\crefname{theoreme}{théorème}{théorèmes}
\crefname{proposition}{proposition}{propositions}
\crefname{corollaire}{corollaire}{corollaires}
\crefname{remarque}{remarque}{remarques}
\begin{document}
\chapter{Premier}
\begin{proposition}\label{p1}a\end{proposition}
\begin{theoreme}\label{t1}b\end{theoreme}
\begin{remarque}\label{r1}c\end{remarque}
\begin{corollaire}\label{c1}d\end{corollaire}
\chapter{Second}
\begin{corollaire}\label{c2}e\end{corollaire}
\begin{theoreme}\label{t2}f\end{theoreme}
RENVOIS \ref{p1} \ref{t1} \ref{r1} \ref{c1} \ref{c2} \ref{t2} ; \cref{p1} ; \cref{t1} ; \cref{r1} ; \cref{c1} ; \cref{c2} ; \cref{t2} FIN
\end{document}
"                       correctif A du préambule
%   lualatex -jobname=crochet  "\def\SANSCROCHET{1}% Juge de la troisième passe : document minimal écrit de zéro pour la cause A (thmtools sous LaTeX 2026-06-01).
% Variantes, par une macro posée sur la ligne de commande (deux passes chacune, depuis memoire/essai/juge3/minimal/) :
%   lualatex -jobname=nu       "% Juge de la troisième passe : document minimal écrit de zéro pour la cause A (thmtools sous LaTeX 2026-06-01).
% Variantes, par une macro posée sur la ligne de commande (deux passes chacune, depuis memoire/essai/juge3/minimal/) :
%   lualatex -jobname=nu       "\input{minimal-juge3}"                                   sans correctif
%   lualatex -jobname=corr     "\def\CORR{1}\input{minimal-juge3}"                       correctif A du préambule
%   lualatex -jobname=crochet  "\def\SANSCROCHET{1}\input{minimal-juge3}"                sans correctif, crochet retiré
%   lualatex -jobname=amsseul  "\def\AMSSEUL{1}\input{minimal-juge3}"                    amsthm seul (témoin)
%   lualatex -jobname=faux     "\def\CORR{1}\def\FAUX{1}\input{minimal-juge3}"           deux vrais défauts : doivent tomber
%   lualatex -jobname=trace    "\def\TRACE{1}\input{minimal-juge3}"                      trace de l'appel
% Les lignes du correctif sont tirées du préambule commité par extraire_correctifs.py, pas retapées.
\documentclass{scrbook}
\ifdefined\SANSCROCHET \RemoveFromHook{package/amsthm/after}[firstaid/aliascounter]\fi
\usepackage{amsthm}
\usepackage{hyperref}
\usepackage[french, nameinlink, noabbrev]{cleveref}
\ifdefined\AMSSEUL
  \newtheorem{theoreme}{Théorème}[chapter]
  \newtheorem{proposition}[theoreme]{Proposition}
  \newtheorem{corollaire}[theoreme]{Corollaire}
  \newtheorem{remarque}{Remarque}[chapter]
\else
  \usepackage{thmtools}
  \ifdefined\CORR \input{correctif-a-pose}\fi
  \declaretheorem[numberwithin=chapter, name=Théorème]{theoreme}
  \ifdefined\TRACE \tracingmacros=2 \tracingonline=0 \fi
  \declaretheorem[sibling=theoreme, name=Proposition]{proposition}
  \ifdefined\TRACE \tracingmacros=0 \fi
  \declaretheorem[sibling=theoreme, name=Corollaire]{corollaire}
  \declaretheorem[numberwithin=chapter, name=Remarque]{remarque}
  \ifdefined\FAUX \newcounteralias{corollaire}{section}\newcounteralias{orphelin}{inexistant}\fi
  \ifdefined\CORR \input{correctif-a-remise}\fi
\fi
\crefname{theoreme}{théorème}{théorèmes}
\crefname{proposition}{proposition}{propositions}
\crefname{corollaire}{corollaire}{corollaires}
\crefname{remarque}{remarque}{remarques}
\begin{document}
\chapter{Premier}
\begin{proposition}\label{p1}a\end{proposition}
\begin{theoreme}\label{t1}b\end{theoreme}
\begin{remarque}\label{r1}c\end{remarque}
\begin{corollaire}\label{c1}d\end{corollaire}
\chapter{Second}
\begin{corollaire}\label{c2}e\end{corollaire}
\begin{theoreme}\label{t2}f\end{theoreme}
RENVOIS \ref{p1} \ref{t1} \ref{r1} \ref{c1} \ref{c2} \ref{t2} ; \cref{p1} ; \cref{t1} ; \cref{r1} ; \cref{c1} ; \cref{c2} ; \cref{t2} FIN
\end{document}
"                                   sans correctif
%   lualatex -jobname=corr     "\def\CORR{1}% Juge de la troisième passe : document minimal écrit de zéro pour la cause A (thmtools sous LaTeX 2026-06-01).
% Variantes, par une macro posée sur la ligne de commande (deux passes chacune, depuis memoire/essai/juge3/minimal/) :
%   lualatex -jobname=nu       "\input{minimal-juge3}"                                   sans correctif
%   lualatex -jobname=corr     "\def\CORR{1}\input{minimal-juge3}"                       correctif A du préambule
%   lualatex -jobname=crochet  "\def\SANSCROCHET{1}\input{minimal-juge3}"                sans correctif, crochet retiré
%   lualatex -jobname=amsseul  "\def\AMSSEUL{1}\input{minimal-juge3}"                    amsthm seul (témoin)
%   lualatex -jobname=faux     "\def\CORR{1}\def\FAUX{1}\input{minimal-juge3}"           deux vrais défauts : doivent tomber
%   lualatex -jobname=trace    "\def\TRACE{1}\input{minimal-juge3}"                      trace de l'appel
% Les lignes du correctif sont tirées du préambule commité par extraire_correctifs.py, pas retapées.
\documentclass{scrbook}
\ifdefined\SANSCROCHET \RemoveFromHook{package/amsthm/after}[firstaid/aliascounter]\fi
\usepackage{amsthm}
\usepackage{hyperref}
\usepackage[french, nameinlink, noabbrev]{cleveref}
\ifdefined\AMSSEUL
  \newtheorem{theoreme}{Théorème}[chapter]
  \newtheorem{proposition}[theoreme]{Proposition}
  \newtheorem{corollaire}[theoreme]{Corollaire}
  \newtheorem{remarque}{Remarque}[chapter]
\else
  \usepackage{thmtools}
  \ifdefined\CORR \input{correctif-a-pose}\fi
  \declaretheorem[numberwithin=chapter, name=Théorème]{theoreme}
  \ifdefined\TRACE \tracingmacros=2 \tracingonline=0 \fi
  \declaretheorem[sibling=theoreme, name=Proposition]{proposition}
  \ifdefined\TRACE \tracingmacros=0 \fi
  \declaretheorem[sibling=theoreme, name=Corollaire]{corollaire}
  \declaretheorem[numberwithin=chapter, name=Remarque]{remarque}
  \ifdefined\FAUX \newcounteralias{corollaire}{section}\newcounteralias{orphelin}{inexistant}\fi
  \ifdefined\CORR \input{correctif-a-remise}\fi
\fi
\crefname{theoreme}{théorème}{théorèmes}
\crefname{proposition}{proposition}{propositions}
\crefname{corollaire}{corollaire}{corollaires}
\crefname{remarque}{remarque}{remarques}
\begin{document}
\chapter{Premier}
\begin{proposition}\label{p1}a\end{proposition}
\begin{theoreme}\label{t1}b\end{theoreme}
\begin{remarque}\label{r1}c\end{remarque}
\begin{corollaire}\label{c1}d\end{corollaire}
\chapter{Second}
\begin{corollaire}\label{c2}e\end{corollaire}
\begin{theoreme}\label{t2}f\end{theoreme}
RENVOIS \ref{p1} \ref{t1} \ref{r1} \ref{c1} \ref{c2} \ref{t2} ; \cref{p1} ; \cref{t1} ; \cref{r1} ; \cref{c1} ; \cref{c2} ; \cref{t2} FIN
\end{document}
"                       correctif A du préambule
%   lualatex -jobname=crochet  "\def\SANSCROCHET{1}% Juge de la troisième passe : document minimal écrit de zéro pour la cause A (thmtools sous LaTeX 2026-06-01).
% Variantes, par une macro posée sur la ligne de commande (deux passes chacune, depuis memoire/essai/juge3/minimal/) :
%   lualatex -jobname=nu       "\input{minimal-juge3}"                                   sans correctif
%   lualatex -jobname=corr     "\def\CORR{1}\input{minimal-juge3}"                       correctif A du préambule
%   lualatex -jobname=crochet  "\def\SANSCROCHET{1}\input{minimal-juge3}"                sans correctif, crochet retiré
%   lualatex -jobname=amsseul  "\def\AMSSEUL{1}\input{minimal-juge3}"                    amsthm seul (témoin)
%   lualatex -jobname=faux     "\def\CORR{1}\def\FAUX{1}\input{minimal-juge3}"           deux vrais défauts : doivent tomber
%   lualatex -jobname=trace    "\def\TRACE{1}\input{minimal-juge3}"                      trace de l'appel
% Les lignes du correctif sont tirées du préambule commité par extraire_correctifs.py, pas retapées.
\documentclass{scrbook}
\ifdefined\SANSCROCHET \RemoveFromHook{package/amsthm/after}[firstaid/aliascounter]\fi
\usepackage{amsthm}
\usepackage{hyperref}
\usepackage[french, nameinlink, noabbrev]{cleveref}
\ifdefined\AMSSEUL
  \newtheorem{theoreme}{Théorème}[chapter]
  \newtheorem{proposition}[theoreme]{Proposition}
  \newtheorem{corollaire}[theoreme]{Corollaire}
  \newtheorem{remarque}{Remarque}[chapter]
\else
  \usepackage{thmtools}
  \ifdefined\CORR \input{correctif-a-pose}\fi
  \declaretheorem[numberwithin=chapter, name=Théorème]{theoreme}
  \ifdefined\TRACE \tracingmacros=2 \tracingonline=0 \fi
  \declaretheorem[sibling=theoreme, name=Proposition]{proposition}
  \ifdefined\TRACE \tracingmacros=0 \fi
  \declaretheorem[sibling=theoreme, name=Corollaire]{corollaire}
  \declaretheorem[numberwithin=chapter, name=Remarque]{remarque}
  \ifdefined\FAUX \newcounteralias{corollaire}{section}\newcounteralias{orphelin}{inexistant}\fi
  \ifdefined\CORR \input{correctif-a-remise}\fi
\fi
\crefname{theoreme}{théorème}{théorèmes}
\crefname{proposition}{proposition}{propositions}
\crefname{corollaire}{corollaire}{corollaires}
\crefname{remarque}{remarque}{remarques}
\begin{document}
\chapter{Premier}
\begin{proposition}\label{p1}a\end{proposition}
\begin{theoreme}\label{t1}b\end{theoreme}
\begin{remarque}\label{r1}c\end{remarque}
\begin{corollaire}\label{c1}d\end{corollaire}
\chapter{Second}
\begin{corollaire}\label{c2}e\end{corollaire}
\begin{theoreme}\label{t2}f\end{theoreme}
RENVOIS \ref{p1} \ref{t1} \ref{r1} \ref{c1} \ref{c2} \ref{t2} ; \cref{p1} ; \cref{t1} ; \cref{r1} ; \cref{c1} ; \cref{c2} ; \cref{t2} FIN
\end{document}
"                sans correctif, crochet retiré
%   lualatex -jobname=amsseul  "\def\AMSSEUL{1}% Juge de la troisième passe : document minimal écrit de zéro pour la cause A (thmtools sous LaTeX 2026-06-01).
% Variantes, par une macro posée sur la ligne de commande (deux passes chacune, depuis memoire/essai/juge3/minimal/) :
%   lualatex -jobname=nu       "\input{minimal-juge3}"                                   sans correctif
%   lualatex -jobname=corr     "\def\CORR{1}\input{minimal-juge3}"                       correctif A du préambule
%   lualatex -jobname=crochet  "\def\SANSCROCHET{1}\input{minimal-juge3}"                sans correctif, crochet retiré
%   lualatex -jobname=amsseul  "\def\AMSSEUL{1}\input{minimal-juge3}"                    amsthm seul (témoin)
%   lualatex -jobname=faux     "\def\CORR{1}\def\FAUX{1}\input{minimal-juge3}"           deux vrais défauts : doivent tomber
%   lualatex -jobname=trace    "\def\TRACE{1}\input{minimal-juge3}"                      trace de l'appel
% Les lignes du correctif sont tirées du préambule commité par extraire_correctifs.py, pas retapées.
\documentclass{scrbook}
\ifdefined\SANSCROCHET \RemoveFromHook{package/amsthm/after}[firstaid/aliascounter]\fi
\usepackage{amsthm}
\usepackage{hyperref}
\usepackage[french, nameinlink, noabbrev]{cleveref}
\ifdefined\AMSSEUL
  \newtheorem{theoreme}{Théorème}[chapter]
  \newtheorem{proposition}[theoreme]{Proposition}
  \newtheorem{corollaire}[theoreme]{Corollaire}
  \newtheorem{remarque}{Remarque}[chapter]
\else
  \usepackage{thmtools}
  \ifdefined\CORR \input{correctif-a-pose}\fi
  \declaretheorem[numberwithin=chapter, name=Théorème]{theoreme}
  \ifdefined\TRACE \tracingmacros=2 \tracingonline=0 \fi
  \declaretheorem[sibling=theoreme, name=Proposition]{proposition}
  \ifdefined\TRACE \tracingmacros=0 \fi
  \declaretheorem[sibling=theoreme, name=Corollaire]{corollaire}
  \declaretheorem[numberwithin=chapter, name=Remarque]{remarque}
  \ifdefined\FAUX \newcounteralias{corollaire}{section}\newcounteralias{orphelin}{inexistant}\fi
  \ifdefined\CORR \input{correctif-a-remise}\fi
\fi
\crefname{theoreme}{théorème}{théorèmes}
\crefname{proposition}{proposition}{propositions}
\crefname{corollaire}{corollaire}{corollaires}
\crefname{remarque}{remarque}{remarques}
\begin{document}
\chapter{Premier}
\begin{proposition}\label{p1}a\end{proposition}
\begin{theoreme}\label{t1}b\end{theoreme}
\begin{remarque}\label{r1}c\end{remarque}
\begin{corollaire}\label{c1}d\end{corollaire}
\chapter{Second}
\begin{corollaire}\label{c2}e\end{corollaire}
\begin{theoreme}\label{t2}f\end{theoreme}
RENVOIS \ref{p1} \ref{t1} \ref{r1} \ref{c1} \ref{c2} \ref{t2} ; \cref{p1} ; \cref{t1} ; \cref{r1} ; \cref{c1} ; \cref{c2} ; \cref{t2} FIN
\end{document}
"                    amsthm seul (témoin)
%   lualatex -jobname=faux     "\def\CORR{1}\def\FAUX{1}% Juge de la troisième passe : document minimal écrit de zéro pour la cause A (thmtools sous LaTeX 2026-06-01).
% Variantes, par une macro posée sur la ligne de commande (deux passes chacune, depuis memoire/essai/juge3/minimal/) :
%   lualatex -jobname=nu       "\input{minimal-juge3}"                                   sans correctif
%   lualatex -jobname=corr     "\def\CORR{1}\input{minimal-juge3}"                       correctif A du préambule
%   lualatex -jobname=crochet  "\def\SANSCROCHET{1}\input{minimal-juge3}"                sans correctif, crochet retiré
%   lualatex -jobname=amsseul  "\def\AMSSEUL{1}\input{minimal-juge3}"                    amsthm seul (témoin)
%   lualatex -jobname=faux     "\def\CORR{1}\def\FAUX{1}\input{minimal-juge3}"           deux vrais défauts : doivent tomber
%   lualatex -jobname=trace    "\def\TRACE{1}\input{minimal-juge3}"                      trace de l'appel
% Les lignes du correctif sont tirées du préambule commité par extraire_correctifs.py, pas retapées.
\documentclass{scrbook}
\ifdefined\SANSCROCHET \RemoveFromHook{package/amsthm/after}[firstaid/aliascounter]\fi
\usepackage{amsthm}
\usepackage{hyperref}
\usepackage[french, nameinlink, noabbrev]{cleveref}
\ifdefined\AMSSEUL
  \newtheorem{theoreme}{Théorème}[chapter]
  \newtheorem{proposition}[theoreme]{Proposition}
  \newtheorem{corollaire}[theoreme]{Corollaire}
  \newtheorem{remarque}{Remarque}[chapter]
\else
  \usepackage{thmtools}
  \ifdefined\CORR \input{correctif-a-pose}\fi
  \declaretheorem[numberwithin=chapter, name=Théorème]{theoreme}
  \ifdefined\TRACE \tracingmacros=2 \tracingonline=0 \fi
  \declaretheorem[sibling=theoreme, name=Proposition]{proposition}
  \ifdefined\TRACE \tracingmacros=0 \fi
  \declaretheorem[sibling=theoreme, name=Corollaire]{corollaire}
  \declaretheorem[numberwithin=chapter, name=Remarque]{remarque}
  \ifdefined\FAUX \newcounteralias{corollaire}{section}\newcounteralias{orphelin}{inexistant}\fi
  \ifdefined\CORR \input{correctif-a-remise}\fi
\fi
\crefname{theoreme}{théorème}{théorèmes}
\crefname{proposition}{proposition}{propositions}
\crefname{corollaire}{corollaire}{corollaires}
\crefname{remarque}{remarque}{remarques}
\begin{document}
\chapter{Premier}
\begin{proposition}\label{p1}a\end{proposition}
\begin{theoreme}\label{t1}b\end{theoreme}
\begin{remarque}\label{r1}c\end{remarque}
\begin{corollaire}\label{c1}d\end{corollaire}
\chapter{Second}
\begin{corollaire}\label{c2}e\end{corollaire}
\begin{theoreme}\label{t2}f\end{theoreme}
RENVOIS \ref{p1} \ref{t1} \ref{r1} \ref{c1} \ref{c2} \ref{t2} ; \cref{p1} ; \cref{t1} ; \cref{r1} ; \cref{c1} ; \cref{c2} ; \cref{t2} FIN
\end{document}
"           deux vrais défauts : doivent tomber
%   lualatex -jobname=trace    "\def\TRACE{1}% Juge de la troisième passe : document minimal écrit de zéro pour la cause A (thmtools sous LaTeX 2026-06-01).
% Variantes, par une macro posée sur la ligne de commande (deux passes chacune, depuis memoire/essai/juge3/minimal/) :
%   lualatex -jobname=nu       "\input{minimal-juge3}"                                   sans correctif
%   lualatex -jobname=corr     "\def\CORR{1}\input{minimal-juge3}"                       correctif A du préambule
%   lualatex -jobname=crochet  "\def\SANSCROCHET{1}\input{minimal-juge3}"                sans correctif, crochet retiré
%   lualatex -jobname=amsseul  "\def\AMSSEUL{1}\input{minimal-juge3}"                    amsthm seul (témoin)
%   lualatex -jobname=faux     "\def\CORR{1}\def\FAUX{1}\input{minimal-juge3}"           deux vrais défauts : doivent tomber
%   lualatex -jobname=trace    "\def\TRACE{1}\input{minimal-juge3}"                      trace de l'appel
% Les lignes du correctif sont tirées du préambule commité par extraire_correctifs.py, pas retapées.
\documentclass{scrbook}
\ifdefined\SANSCROCHET \RemoveFromHook{package/amsthm/after}[firstaid/aliascounter]\fi
\usepackage{amsthm}
\usepackage{hyperref}
\usepackage[french, nameinlink, noabbrev]{cleveref}
\ifdefined\AMSSEUL
  \newtheorem{theoreme}{Théorème}[chapter]
  \newtheorem{proposition}[theoreme]{Proposition}
  \newtheorem{corollaire}[theoreme]{Corollaire}
  \newtheorem{remarque}{Remarque}[chapter]
\else
  \usepackage{thmtools}
  \ifdefined\CORR \input{correctif-a-pose}\fi
  \declaretheorem[numberwithin=chapter, name=Théorème]{theoreme}
  \ifdefined\TRACE \tracingmacros=2 \tracingonline=0 \fi
  \declaretheorem[sibling=theoreme, name=Proposition]{proposition}
  \ifdefined\TRACE \tracingmacros=0 \fi
  \declaretheorem[sibling=theoreme, name=Corollaire]{corollaire}
  \declaretheorem[numberwithin=chapter, name=Remarque]{remarque}
  \ifdefined\FAUX \newcounteralias{corollaire}{section}\newcounteralias{orphelin}{inexistant}\fi
  \ifdefined\CORR \input{correctif-a-remise}\fi
\fi
\crefname{theoreme}{théorème}{théorèmes}
\crefname{proposition}{proposition}{propositions}
\crefname{corollaire}{corollaire}{corollaires}
\crefname{remarque}{remarque}{remarques}
\begin{document}
\chapter{Premier}
\begin{proposition}\label{p1}a\end{proposition}
\begin{theoreme}\label{t1}b\end{theoreme}
\begin{remarque}\label{r1}c\end{remarque}
\begin{corollaire}\label{c1}d\end{corollaire}
\chapter{Second}
\begin{corollaire}\label{c2}e\end{corollaire}
\begin{theoreme}\label{t2}f\end{theoreme}
RENVOIS \ref{p1} \ref{t1} \ref{r1} \ref{c1} \ref{c2} \ref{t2} ; \cref{p1} ; \cref{t1} ; \cref{r1} ; \cref{c1} ; \cref{c2} ; \cref{t2} FIN
\end{document}
"                      trace de l'appel
% Les lignes du correctif sont tirées du préambule commité par extraire_correctifs.py, pas retapées.
\documentclass{scrbook}
\ifdefined\SANSCROCHET \RemoveFromHook{package/amsthm/after}[firstaid/aliascounter]\fi
\usepackage{amsthm}
\usepackage{hyperref}
\usepackage[french, nameinlink, noabbrev]{cleveref}
\ifdefined\AMSSEUL
  \newtheorem{theoreme}{Théorème}[chapter]
  \newtheorem{proposition}[theoreme]{Proposition}
  \newtheorem{corollaire}[theoreme]{Corollaire}
  \newtheorem{remarque}{Remarque}[chapter]
\else
  \usepackage{thmtools}
  \ifdefined\CORR \input{correctif-a-pose}\fi
  \declaretheorem[numberwithin=chapter, name=Théorème]{theoreme}
  \ifdefined\TRACE \tracingmacros=2 \tracingonline=0 \fi
  \declaretheorem[sibling=theoreme, name=Proposition]{proposition}
  \ifdefined\TRACE \tracingmacros=0 \fi
  \declaretheorem[sibling=theoreme, name=Corollaire]{corollaire}
  \declaretheorem[numberwithin=chapter, name=Remarque]{remarque}
  \ifdefined\FAUX \newcounteralias{corollaire}{section}\newcounteralias{orphelin}{inexistant}\fi
  \ifdefined\CORR \input{correctif-a-remise}\fi
\fi
\crefname{theoreme}{théorème}{théorèmes}
\crefname{proposition}{proposition}{propositions}
\crefname{corollaire}{corollaire}{corollaires}
\crefname{remarque}{remarque}{remarques}
\begin{document}
\chapter{Premier}
\begin{proposition}\label{p1}a\end{proposition}
\begin{theoreme}\label{t1}b\end{theoreme}
\begin{remarque}\label{r1}c\end{remarque}
\begin{corollaire}\label{c1}d\end{corollaire}
\chapter{Second}
\begin{corollaire}\label{c2}e\end{corollaire}
\begin{theoreme}\label{t2}f\end{theoreme}
RENVOIS \ref{p1} \ref{t1} \ref{r1} \ref{c1} \ref{c2} \ref{t2} ; \cref{p1} ; \cref{t1} ; \cref{r1} ; \cref{c1} ; \cref{c2} ; \cref{t2} FIN
\end{document}
"                sans correctif, crochet retiré
%   lualatex -jobname=amsseul  "\def\AMSSEUL{1}% Juge de la troisième passe : document minimal écrit de zéro pour la cause A (thmtools sous LaTeX 2026-06-01).
% Variantes, par une macro posée sur la ligne de commande (deux passes chacune, depuis memoire/essai/juge3/minimal/) :
%   lualatex -jobname=nu       "% Juge de la troisième passe : document minimal écrit de zéro pour la cause A (thmtools sous LaTeX 2026-06-01).
% Variantes, par une macro posée sur la ligne de commande (deux passes chacune, depuis memoire/essai/juge3/minimal/) :
%   lualatex -jobname=nu       "\input{minimal-juge3}"                                   sans correctif
%   lualatex -jobname=corr     "\def\CORR{1}\input{minimal-juge3}"                       correctif A du préambule
%   lualatex -jobname=crochet  "\def\SANSCROCHET{1}\input{minimal-juge3}"                sans correctif, crochet retiré
%   lualatex -jobname=amsseul  "\def\AMSSEUL{1}\input{minimal-juge3}"                    amsthm seul (témoin)
%   lualatex -jobname=faux     "\def\CORR{1}\def\FAUX{1}\input{minimal-juge3}"           deux vrais défauts : doivent tomber
%   lualatex -jobname=trace    "\def\TRACE{1}\input{minimal-juge3}"                      trace de l'appel
% Les lignes du correctif sont tirées du préambule commité par extraire_correctifs.py, pas retapées.
\documentclass{scrbook}
\ifdefined\SANSCROCHET \RemoveFromHook{package/amsthm/after}[firstaid/aliascounter]\fi
\usepackage{amsthm}
\usepackage{hyperref}
\usepackage[french, nameinlink, noabbrev]{cleveref}
\ifdefined\AMSSEUL
  \newtheorem{theoreme}{Théorème}[chapter]
  \newtheorem{proposition}[theoreme]{Proposition}
  \newtheorem{corollaire}[theoreme]{Corollaire}
  \newtheorem{remarque}{Remarque}[chapter]
\else
  \usepackage{thmtools}
  \ifdefined\CORR \input{correctif-a-pose}\fi
  \declaretheorem[numberwithin=chapter, name=Théorème]{theoreme}
  \ifdefined\TRACE \tracingmacros=2 \tracingonline=0 \fi
  \declaretheorem[sibling=theoreme, name=Proposition]{proposition}
  \ifdefined\TRACE \tracingmacros=0 \fi
  \declaretheorem[sibling=theoreme, name=Corollaire]{corollaire}
  \declaretheorem[numberwithin=chapter, name=Remarque]{remarque}
  \ifdefined\FAUX \newcounteralias{corollaire}{section}\newcounteralias{orphelin}{inexistant}\fi
  \ifdefined\CORR \input{correctif-a-remise}\fi
\fi
\crefname{theoreme}{théorème}{théorèmes}
\crefname{proposition}{proposition}{propositions}
\crefname{corollaire}{corollaire}{corollaires}
\crefname{remarque}{remarque}{remarques}
\begin{document}
\chapter{Premier}
\begin{proposition}\label{p1}a\end{proposition}
\begin{theoreme}\label{t1}b\end{theoreme}
\begin{remarque}\label{r1}c\end{remarque}
\begin{corollaire}\label{c1}d\end{corollaire}
\chapter{Second}
\begin{corollaire}\label{c2}e\end{corollaire}
\begin{theoreme}\label{t2}f\end{theoreme}
RENVOIS \ref{p1} \ref{t1} \ref{r1} \ref{c1} \ref{c2} \ref{t2} ; \cref{p1} ; \cref{t1} ; \cref{r1} ; \cref{c1} ; \cref{c2} ; \cref{t2} FIN
\end{document}
"                                   sans correctif
%   lualatex -jobname=corr     "\def\CORR{1}% Juge de la troisième passe : document minimal écrit de zéro pour la cause A (thmtools sous LaTeX 2026-06-01).
% Variantes, par une macro posée sur la ligne de commande (deux passes chacune, depuis memoire/essai/juge3/minimal/) :
%   lualatex -jobname=nu       "\input{minimal-juge3}"                                   sans correctif
%   lualatex -jobname=corr     "\def\CORR{1}\input{minimal-juge3}"                       correctif A du préambule
%   lualatex -jobname=crochet  "\def\SANSCROCHET{1}\input{minimal-juge3}"                sans correctif, crochet retiré
%   lualatex -jobname=amsseul  "\def\AMSSEUL{1}\input{minimal-juge3}"                    amsthm seul (témoin)
%   lualatex -jobname=faux     "\def\CORR{1}\def\FAUX{1}\input{minimal-juge3}"           deux vrais défauts : doivent tomber
%   lualatex -jobname=trace    "\def\TRACE{1}\input{minimal-juge3}"                      trace de l'appel
% Les lignes du correctif sont tirées du préambule commité par extraire_correctifs.py, pas retapées.
\documentclass{scrbook}
\ifdefined\SANSCROCHET \RemoveFromHook{package/amsthm/after}[firstaid/aliascounter]\fi
\usepackage{amsthm}
\usepackage{hyperref}
\usepackage[french, nameinlink, noabbrev]{cleveref}
\ifdefined\AMSSEUL
  \newtheorem{theoreme}{Théorème}[chapter]
  \newtheorem{proposition}[theoreme]{Proposition}
  \newtheorem{corollaire}[theoreme]{Corollaire}
  \newtheorem{remarque}{Remarque}[chapter]
\else
  \usepackage{thmtools}
  \ifdefined\CORR \input{correctif-a-pose}\fi
  \declaretheorem[numberwithin=chapter, name=Théorème]{theoreme}
  \ifdefined\TRACE \tracingmacros=2 \tracingonline=0 \fi
  \declaretheorem[sibling=theoreme, name=Proposition]{proposition}
  \ifdefined\TRACE \tracingmacros=0 \fi
  \declaretheorem[sibling=theoreme, name=Corollaire]{corollaire}
  \declaretheorem[numberwithin=chapter, name=Remarque]{remarque}
  \ifdefined\FAUX \newcounteralias{corollaire}{section}\newcounteralias{orphelin}{inexistant}\fi
  \ifdefined\CORR \input{correctif-a-remise}\fi
\fi
\crefname{theoreme}{théorème}{théorèmes}
\crefname{proposition}{proposition}{propositions}
\crefname{corollaire}{corollaire}{corollaires}
\crefname{remarque}{remarque}{remarques}
\begin{document}
\chapter{Premier}
\begin{proposition}\label{p1}a\end{proposition}
\begin{theoreme}\label{t1}b\end{theoreme}
\begin{remarque}\label{r1}c\end{remarque}
\begin{corollaire}\label{c1}d\end{corollaire}
\chapter{Second}
\begin{corollaire}\label{c2}e\end{corollaire}
\begin{theoreme}\label{t2}f\end{theoreme}
RENVOIS \ref{p1} \ref{t1} \ref{r1} \ref{c1} \ref{c2} \ref{t2} ; \cref{p1} ; \cref{t1} ; \cref{r1} ; \cref{c1} ; \cref{c2} ; \cref{t2} FIN
\end{document}
"                       correctif A du préambule
%   lualatex -jobname=crochet  "\def\SANSCROCHET{1}% Juge de la troisième passe : document minimal écrit de zéro pour la cause A (thmtools sous LaTeX 2026-06-01).
% Variantes, par une macro posée sur la ligne de commande (deux passes chacune, depuis memoire/essai/juge3/minimal/) :
%   lualatex -jobname=nu       "\input{minimal-juge3}"                                   sans correctif
%   lualatex -jobname=corr     "\def\CORR{1}\input{minimal-juge3}"                       correctif A du préambule
%   lualatex -jobname=crochet  "\def\SANSCROCHET{1}\input{minimal-juge3}"                sans correctif, crochet retiré
%   lualatex -jobname=amsseul  "\def\AMSSEUL{1}\input{minimal-juge3}"                    amsthm seul (témoin)
%   lualatex -jobname=faux     "\def\CORR{1}\def\FAUX{1}\input{minimal-juge3}"           deux vrais défauts : doivent tomber
%   lualatex -jobname=trace    "\def\TRACE{1}\input{minimal-juge3}"                      trace de l'appel
% Les lignes du correctif sont tirées du préambule commité par extraire_correctifs.py, pas retapées.
\documentclass{scrbook}
\ifdefined\SANSCROCHET \RemoveFromHook{package/amsthm/after}[firstaid/aliascounter]\fi
\usepackage{amsthm}
\usepackage{hyperref}
\usepackage[french, nameinlink, noabbrev]{cleveref}
\ifdefined\AMSSEUL
  \newtheorem{theoreme}{Théorème}[chapter]
  \newtheorem{proposition}[theoreme]{Proposition}
  \newtheorem{corollaire}[theoreme]{Corollaire}
  \newtheorem{remarque}{Remarque}[chapter]
\else
  \usepackage{thmtools}
  \ifdefined\CORR \input{correctif-a-pose}\fi
  \declaretheorem[numberwithin=chapter, name=Théorème]{theoreme}
  \ifdefined\TRACE \tracingmacros=2 \tracingonline=0 \fi
  \declaretheorem[sibling=theoreme, name=Proposition]{proposition}
  \ifdefined\TRACE \tracingmacros=0 \fi
  \declaretheorem[sibling=theoreme, name=Corollaire]{corollaire}
  \declaretheorem[numberwithin=chapter, name=Remarque]{remarque}
  \ifdefined\FAUX \newcounteralias{corollaire}{section}\newcounteralias{orphelin}{inexistant}\fi
  \ifdefined\CORR \input{correctif-a-remise}\fi
\fi
\crefname{theoreme}{théorème}{théorèmes}
\crefname{proposition}{proposition}{propositions}
\crefname{corollaire}{corollaire}{corollaires}
\crefname{remarque}{remarque}{remarques}
\begin{document}
\chapter{Premier}
\begin{proposition}\label{p1}a\end{proposition}
\begin{theoreme}\label{t1}b\end{theoreme}
\begin{remarque}\label{r1}c\end{remarque}
\begin{corollaire}\label{c1}d\end{corollaire}
\chapter{Second}
\begin{corollaire}\label{c2}e\end{corollaire}
\begin{theoreme}\label{t2}f\end{theoreme}
RENVOIS \ref{p1} \ref{t1} \ref{r1} \ref{c1} \ref{c2} \ref{t2} ; \cref{p1} ; \cref{t1} ; \cref{r1} ; \cref{c1} ; \cref{c2} ; \cref{t2} FIN
\end{document}
"                sans correctif, crochet retiré
%   lualatex -jobname=amsseul  "\def\AMSSEUL{1}% Juge de la troisième passe : document minimal écrit de zéro pour la cause A (thmtools sous LaTeX 2026-06-01).
% Variantes, par une macro posée sur la ligne de commande (deux passes chacune, depuis memoire/essai/juge3/minimal/) :
%   lualatex -jobname=nu       "\input{minimal-juge3}"                                   sans correctif
%   lualatex -jobname=corr     "\def\CORR{1}\input{minimal-juge3}"                       correctif A du préambule
%   lualatex -jobname=crochet  "\def\SANSCROCHET{1}\input{minimal-juge3}"                sans correctif, crochet retiré
%   lualatex -jobname=amsseul  "\def\AMSSEUL{1}\input{minimal-juge3}"                    amsthm seul (témoin)
%   lualatex -jobname=faux     "\def\CORR{1}\def\FAUX{1}\input{minimal-juge3}"           deux vrais défauts : doivent tomber
%   lualatex -jobname=trace    "\def\TRACE{1}\input{minimal-juge3}"                      trace de l'appel
% Les lignes du correctif sont tirées du préambule commité par extraire_correctifs.py, pas retapées.
\documentclass{scrbook}
\ifdefined\SANSCROCHET \RemoveFromHook{package/amsthm/after}[firstaid/aliascounter]\fi
\usepackage{amsthm}
\usepackage{hyperref}
\usepackage[french, nameinlink, noabbrev]{cleveref}
\ifdefined\AMSSEUL
  \newtheorem{theoreme}{Théorème}[chapter]
  \newtheorem{proposition}[theoreme]{Proposition}
  \newtheorem{corollaire}[theoreme]{Corollaire}
  \newtheorem{remarque}{Remarque}[chapter]
\else
  \usepackage{thmtools}
  \ifdefined\CORR \input{correctif-a-pose}\fi
  \declaretheorem[numberwithin=chapter, name=Théorème]{theoreme}
  \ifdefined\TRACE \tracingmacros=2 \tracingonline=0 \fi
  \declaretheorem[sibling=theoreme, name=Proposition]{proposition}
  \ifdefined\TRACE \tracingmacros=0 \fi
  \declaretheorem[sibling=theoreme, name=Corollaire]{corollaire}
  \declaretheorem[numberwithin=chapter, name=Remarque]{remarque}
  \ifdefined\FAUX \newcounteralias{corollaire}{section}\newcounteralias{orphelin}{inexistant}\fi
  \ifdefined\CORR \input{correctif-a-remise}\fi
\fi
\crefname{theoreme}{théorème}{théorèmes}
\crefname{proposition}{proposition}{propositions}
\crefname{corollaire}{corollaire}{corollaires}
\crefname{remarque}{remarque}{remarques}
\begin{document}
\chapter{Premier}
\begin{proposition}\label{p1}a\end{proposition}
\begin{theoreme}\label{t1}b\end{theoreme}
\begin{remarque}\label{r1}c\end{remarque}
\begin{corollaire}\label{c1}d\end{corollaire}
\chapter{Second}
\begin{corollaire}\label{c2}e\end{corollaire}
\begin{theoreme}\label{t2}f\end{theoreme}
RENVOIS \ref{p1} \ref{t1} \ref{r1} \ref{c1} \ref{c2} \ref{t2} ; \cref{p1} ; \cref{t1} ; \cref{r1} ; \cref{c1} ; \cref{c2} ; \cref{t2} FIN
\end{document}
"                    amsthm seul (témoin)
%   lualatex -jobname=faux     "\def\CORR{1}\def\FAUX{1}% Juge de la troisième passe : document minimal écrit de zéro pour la cause A (thmtools sous LaTeX 2026-06-01).
% Variantes, par une macro posée sur la ligne de commande (deux passes chacune, depuis memoire/essai/juge3/minimal/) :
%   lualatex -jobname=nu       "\input{minimal-juge3}"                                   sans correctif
%   lualatex -jobname=corr     "\def\CORR{1}\input{minimal-juge3}"                       correctif A du préambule
%   lualatex -jobname=crochet  "\def\SANSCROCHET{1}\input{minimal-juge3}"                sans correctif, crochet retiré
%   lualatex -jobname=amsseul  "\def\AMSSEUL{1}\input{minimal-juge3}"                    amsthm seul (témoin)
%   lualatex -jobname=faux     "\def\CORR{1}\def\FAUX{1}\input{minimal-juge3}"           deux vrais défauts : doivent tomber
%   lualatex -jobname=trace    "\def\TRACE{1}\input{minimal-juge3}"                      trace de l'appel
% Les lignes du correctif sont tirées du préambule commité par extraire_correctifs.py, pas retapées.
\documentclass{scrbook}
\ifdefined\SANSCROCHET \RemoveFromHook{package/amsthm/after}[firstaid/aliascounter]\fi
\usepackage{amsthm}
\usepackage{hyperref}
\usepackage[french, nameinlink, noabbrev]{cleveref}
\ifdefined\AMSSEUL
  \newtheorem{theoreme}{Théorème}[chapter]
  \newtheorem{proposition}[theoreme]{Proposition}
  \newtheorem{corollaire}[theoreme]{Corollaire}
  \newtheorem{remarque}{Remarque}[chapter]
\else
  \usepackage{thmtools}
  \ifdefined\CORR \input{correctif-a-pose}\fi
  \declaretheorem[numberwithin=chapter, name=Théorème]{theoreme}
  \ifdefined\TRACE \tracingmacros=2 \tracingonline=0 \fi
  \declaretheorem[sibling=theoreme, name=Proposition]{proposition}
  \ifdefined\TRACE \tracingmacros=0 \fi
  \declaretheorem[sibling=theoreme, name=Corollaire]{corollaire}
  \declaretheorem[numberwithin=chapter, name=Remarque]{remarque}
  \ifdefined\FAUX \newcounteralias{corollaire}{section}\newcounteralias{orphelin}{inexistant}\fi
  \ifdefined\CORR \input{correctif-a-remise}\fi
\fi
\crefname{theoreme}{théorème}{théorèmes}
\crefname{proposition}{proposition}{propositions}
\crefname{corollaire}{corollaire}{corollaires}
\crefname{remarque}{remarque}{remarques}
\begin{document}
\chapter{Premier}
\begin{proposition}\label{p1}a\end{proposition}
\begin{theoreme}\label{t1}b\end{theoreme}
\begin{remarque}\label{r1}c\end{remarque}
\begin{corollaire}\label{c1}d\end{corollaire}
\chapter{Second}
\begin{corollaire}\label{c2}e\end{corollaire}
\begin{theoreme}\label{t2}f\end{theoreme}
RENVOIS \ref{p1} \ref{t1} \ref{r1} \ref{c1} \ref{c2} \ref{t2} ; \cref{p1} ; \cref{t1} ; \cref{r1} ; \cref{c1} ; \cref{c2} ; \cref{t2} FIN
\end{document}
"           deux vrais défauts : doivent tomber
%   lualatex -jobname=trace    "\def\TRACE{1}% Juge de la troisième passe : document minimal écrit de zéro pour la cause A (thmtools sous LaTeX 2026-06-01).
% Variantes, par une macro posée sur la ligne de commande (deux passes chacune, depuis memoire/essai/juge3/minimal/) :
%   lualatex -jobname=nu       "\input{minimal-juge3}"                                   sans correctif
%   lualatex -jobname=corr     "\def\CORR{1}\input{minimal-juge3}"                       correctif A du préambule
%   lualatex -jobname=crochet  "\def\SANSCROCHET{1}\input{minimal-juge3}"                sans correctif, crochet retiré
%   lualatex -jobname=amsseul  "\def\AMSSEUL{1}\input{minimal-juge3}"                    amsthm seul (témoin)
%   lualatex -jobname=faux     "\def\CORR{1}\def\FAUX{1}\input{minimal-juge3}"           deux vrais défauts : doivent tomber
%   lualatex -jobname=trace    "\def\TRACE{1}\input{minimal-juge3}"                      trace de l'appel
% Les lignes du correctif sont tirées du préambule commité par extraire_correctifs.py, pas retapées.
\documentclass{scrbook}
\ifdefined\SANSCROCHET \RemoveFromHook{package/amsthm/after}[firstaid/aliascounter]\fi
\usepackage{amsthm}
\usepackage{hyperref}
\usepackage[french, nameinlink, noabbrev]{cleveref}
\ifdefined\AMSSEUL
  \newtheorem{theoreme}{Théorème}[chapter]
  \newtheorem{proposition}[theoreme]{Proposition}
  \newtheorem{corollaire}[theoreme]{Corollaire}
  \newtheorem{remarque}{Remarque}[chapter]
\else
  \usepackage{thmtools}
  \ifdefined\CORR \input{correctif-a-pose}\fi
  \declaretheorem[numberwithin=chapter, name=Théorème]{theoreme}
  \ifdefined\TRACE \tracingmacros=2 \tracingonline=0 \fi
  \declaretheorem[sibling=theoreme, name=Proposition]{proposition}
  \ifdefined\TRACE \tracingmacros=0 \fi
  \declaretheorem[sibling=theoreme, name=Corollaire]{corollaire}
  \declaretheorem[numberwithin=chapter, name=Remarque]{remarque}
  \ifdefined\FAUX \newcounteralias{corollaire}{section}\newcounteralias{orphelin}{inexistant}\fi
  \ifdefined\CORR \input{correctif-a-remise}\fi
\fi
\crefname{theoreme}{théorème}{théorèmes}
\crefname{proposition}{proposition}{propositions}
\crefname{corollaire}{corollaire}{corollaires}
\crefname{remarque}{remarque}{remarques}
\begin{document}
\chapter{Premier}
\begin{proposition}\label{p1}a\end{proposition}
\begin{theoreme}\label{t1}b\end{theoreme}
\begin{remarque}\label{r1}c\end{remarque}
\begin{corollaire}\label{c1}d\end{corollaire}
\chapter{Second}
\begin{corollaire}\label{c2}e\end{corollaire}
\begin{theoreme}\label{t2}f\end{theoreme}
RENVOIS \ref{p1} \ref{t1} \ref{r1} \ref{c1} \ref{c2} \ref{t2} ; \cref{p1} ; \cref{t1} ; \cref{r1} ; \cref{c1} ; \cref{c2} ; \cref{t2} FIN
\end{document}
"                      trace de l'appel
% Les lignes du correctif sont tirées du préambule commité par extraire_correctifs.py, pas retapées.
\documentclass{scrbook}
\ifdefined\SANSCROCHET \RemoveFromHook{package/amsthm/after}[firstaid/aliascounter]\fi
\usepackage{amsthm}
\usepackage{hyperref}
\usepackage[french, nameinlink, noabbrev]{cleveref}
\ifdefined\AMSSEUL
  \newtheorem{theoreme}{Théorème}[chapter]
  \newtheorem{proposition}[theoreme]{Proposition}
  \newtheorem{corollaire}[theoreme]{Corollaire}
  \newtheorem{remarque}{Remarque}[chapter]
\else
  \usepackage{thmtools}
  \ifdefined\CORR \input{correctif-a-pose}\fi
  \declaretheorem[numberwithin=chapter, name=Théorème]{theoreme}
  \ifdefined\TRACE \tracingmacros=2 \tracingonline=0 \fi
  \declaretheorem[sibling=theoreme, name=Proposition]{proposition}
  \ifdefined\TRACE \tracingmacros=0 \fi
  \declaretheorem[sibling=theoreme, name=Corollaire]{corollaire}
  \declaretheorem[numberwithin=chapter, name=Remarque]{remarque}
  \ifdefined\FAUX \newcounteralias{corollaire}{section}\newcounteralias{orphelin}{inexistant}\fi
  \ifdefined\CORR \input{correctif-a-remise}\fi
\fi
\crefname{theoreme}{théorème}{théorèmes}
\crefname{proposition}{proposition}{propositions}
\crefname{corollaire}{corollaire}{corollaires}
\crefname{remarque}{remarque}{remarques}
\begin{document}
\chapter{Premier}
\begin{proposition}\label{p1}a\end{proposition}
\begin{theoreme}\label{t1}b\end{theoreme}
\begin{remarque}\label{r1}c\end{remarque}
\begin{corollaire}\label{c1}d\end{corollaire}
\chapter{Second}
\begin{corollaire}\label{c2}e\end{corollaire}
\begin{theoreme}\label{t2}f\end{theoreme}
RENVOIS \ref{p1} \ref{t1} \ref{r1} \ref{c1} \ref{c2} \ref{t2} ; \cref{p1} ; \cref{t1} ; \cref{r1} ; \cref{c1} ; \cref{c2} ; \cref{t2} FIN
\end{document}
"                    amsthm seul (témoin)
%   lualatex -jobname=faux     "\def\CORR{1}\def\FAUX{1}% Juge de la troisième passe : document minimal écrit de zéro pour la cause A (thmtools sous LaTeX 2026-06-01).
% Variantes, par une macro posée sur la ligne de commande (deux passes chacune, depuis memoire/essai/juge3/minimal/) :
%   lualatex -jobname=nu       "% Juge de la troisième passe : document minimal écrit de zéro pour la cause A (thmtools sous LaTeX 2026-06-01).
% Variantes, par une macro posée sur la ligne de commande (deux passes chacune, depuis memoire/essai/juge3/minimal/) :
%   lualatex -jobname=nu       "\input{minimal-juge3}"                                   sans correctif
%   lualatex -jobname=corr     "\def\CORR{1}\input{minimal-juge3}"                       correctif A du préambule
%   lualatex -jobname=crochet  "\def\SANSCROCHET{1}\input{minimal-juge3}"                sans correctif, crochet retiré
%   lualatex -jobname=amsseul  "\def\AMSSEUL{1}\input{minimal-juge3}"                    amsthm seul (témoin)
%   lualatex -jobname=faux     "\def\CORR{1}\def\FAUX{1}\input{minimal-juge3}"           deux vrais défauts : doivent tomber
%   lualatex -jobname=trace    "\def\TRACE{1}\input{minimal-juge3}"                      trace de l'appel
% Les lignes du correctif sont tirées du préambule commité par extraire_correctifs.py, pas retapées.
\documentclass{scrbook}
\ifdefined\SANSCROCHET \RemoveFromHook{package/amsthm/after}[firstaid/aliascounter]\fi
\usepackage{amsthm}
\usepackage{hyperref}
\usepackage[french, nameinlink, noabbrev]{cleveref}
\ifdefined\AMSSEUL
  \newtheorem{theoreme}{Théorème}[chapter]
  \newtheorem{proposition}[theoreme]{Proposition}
  \newtheorem{corollaire}[theoreme]{Corollaire}
  \newtheorem{remarque}{Remarque}[chapter]
\else
  \usepackage{thmtools}
  \ifdefined\CORR \input{correctif-a-pose}\fi
  \declaretheorem[numberwithin=chapter, name=Théorème]{theoreme}
  \ifdefined\TRACE \tracingmacros=2 \tracingonline=0 \fi
  \declaretheorem[sibling=theoreme, name=Proposition]{proposition}
  \ifdefined\TRACE \tracingmacros=0 \fi
  \declaretheorem[sibling=theoreme, name=Corollaire]{corollaire}
  \declaretheorem[numberwithin=chapter, name=Remarque]{remarque}
  \ifdefined\FAUX \newcounteralias{corollaire}{section}\newcounteralias{orphelin}{inexistant}\fi
  \ifdefined\CORR \input{correctif-a-remise}\fi
\fi
\crefname{theoreme}{théorème}{théorèmes}
\crefname{proposition}{proposition}{propositions}
\crefname{corollaire}{corollaire}{corollaires}
\crefname{remarque}{remarque}{remarques}
\begin{document}
\chapter{Premier}
\begin{proposition}\label{p1}a\end{proposition}
\begin{theoreme}\label{t1}b\end{theoreme}
\begin{remarque}\label{r1}c\end{remarque}
\begin{corollaire}\label{c1}d\end{corollaire}
\chapter{Second}
\begin{corollaire}\label{c2}e\end{corollaire}
\begin{theoreme}\label{t2}f\end{theoreme}
RENVOIS \ref{p1} \ref{t1} \ref{r1} \ref{c1} \ref{c2} \ref{t2} ; \cref{p1} ; \cref{t1} ; \cref{r1} ; \cref{c1} ; \cref{c2} ; \cref{t2} FIN
\end{document}
"                                   sans correctif
%   lualatex -jobname=corr     "\def\CORR{1}% Juge de la troisième passe : document minimal écrit de zéro pour la cause A (thmtools sous LaTeX 2026-06-01).
% Variantes, par une macro posée sur la ligne de commande (deux passes chacune, depuis memoire/essai/juge3/minimal/) :
%   lualatex -jobname=nu       "\input{minimal-juge3}"                                   sans correctif
%   lualatex -jobname=corr     "\def\CORR{1}\input{minimal-juge3}"                       correctif A du préambule
%   lualatex -jobname=crochet  "\def\SANSCROCHET{1}\input{minimal-juge3}"                sans correctif, crochet retiré
%   lualatex -jobname=amsseul  "\def\AMSSEUL{1}\input{minimal-juge3}"                    amsthm seul (témoin)
%   lualatex -jobname=faux     "\def\CORR{1}\def\FAUX{1}\input{minimal-juge3}"           deux vrais défauts : doivent tomber
%   lualatex -jobname=trace    "\def\TRACE{1}\input{minimal-juge3}"                      trace de l'appel
% Les lignes du correctif sont tirées du préambule commité par extraire_correctifs.py, pas retapées.
\documentclass{scrbook}
\ifdefined\SANSCROCHET \RemoveFromHook{package/amsthm/after}[firstaid/aliascounter]\fi
\usepackage{amsthm}
\usepackage{hyperref}
\usepackage[french, nameinlink, noabbrev]{cleveref}
\ifdefined\AMSSEUL
  \newtheorem{theoreme}{Théorème}[chapter]
  \newtheorem{proposition}[theoreme]{Proposition}
  \newtheorem{corollaire}[theoreme]{Corollaire}
  \newtheorem{remarque}{Remarque}[chapter]
\else
  \usepackage{thmtools}
  \ifdefined\CORR \input{correctif-a-pose}\fi
  \declaretheorem[numberwithin=chapter, name=Théorème]{theoreme}
  \ifdefined\TRACE \tracingmacros=2 \tracingonline=0 \fi
  \declaretheorem[sibling=theoreme, name=Proposition]{proposition}
  \ifdefined\TRACE \tracingmacros=0 \fi
  \declaretheorem[sibling=theoreme, name=Corollaire]{corollaire}
  \declaretheorem[numberwithin=chapter, name=Remarque]{remarque}
  \ifdefined\FAUX \newcounteralias{corollaire}{section}\newcounteralias{orphelin}{inexistant}\fi
  \ifdefined\CORR \input{correctif-a-remise}\fi
\fi
\crefname{theoreme}{théorème}{théorèmes}
\crefname{proposition}{proposition}{propositions}
\crefname{corollaire}{corollaire}{corollaires}
\crefname{remarque}{remarque}{remarques}
\begin{document}
\chapter{Premier}
\begin{proposition}\label{p1}a\end{proposition}
\begin{theoreme}\label{t1}b\end{theoreme}
\begin{remarque}\label{r1}c\end{remarque}
\begin{corollaire}\label{c1}d\end{corollaire}
\chapter{Second}
\begin{corollaire}\label{c2}e\end{corollaire}
\begin{theoreme}\label{t2}f\end{theoreme}
RENVOIS \ref{p1} \ref{t1} \ref{r1} \ref{c1} \ref{c2} \ref{t2} ; \cref{p1} ; \cref{t1} ; \cref{r1} ; \cref{c1} ; \cref{c2} ; \cref{t2} FIN
\end{document}
"                       correctif A du préambule
%   lualatex -jobname=crochet  "\def\SANSCROCHET{1}% Juge de la troisième passe : document minimal écrit de zéro pour la cause A (thmtools sous LaTeX 2026-06-01).
% Variantes, par une macro posée sur la ligne de commande (deux passes chacune, depuis memoire/essai/juge3/minimal/) :
%   lualatex -jobname=nu       "\input{minimal-juge3}"                                   sans correctif
%   lualatex -jobname=corr     "\def\CORR{1}\input{minimal-juge3}"                       correctif A du préambule
%   lualatex -jobname=crochet  "\def\SANSCROCHET{1}\input{minimal-juge3}"                sans correctif, crochet retiré
%   lualatex -jobname=amsseul  "\def\AMSSEUL{1}\input{minimal-juge3}"                    amsthm seul (témoin)
%   lualatex -jobname=faux     "\def\CORR{1}\def\FAUX{1}\input{minimal-juge3}"           deux vrais défauts : doivent tomber
%   lualatex -jobname=trace    "\def\TRACE{1}\input{minimal-juge3}"                      trace de l'appel
% Les lignes du correctif sont tirées du préambule commité par extraire_correctifs.py, pas retapées.
\documentclass{scrbook}
\ifdefined\SANSCROCHET \RemoveFromHook{package/amsthm/after}[firstaid/aliascounter]\fi
\usepackage{amsthm}
\usepackage{hyperref}
\usepackage[french, nameinlink, noabbrev]{cleveref}
\ifdefined\AMSSEUL
  \newtheorem{theoreme}{Théorème}[chapter]
  \newtheorem{proposition}[theoreme]{Proposition}
  \newtheorem{corollaire}[theoreme]{Corollaire}
  \newtheorem{remarque}{Remarque}[chapter]
\else
  \usepackage{thmtools}
  \ifdefined\CORR \input{correctif-a-pose}\fi
  \declaretheorem[numberwithin=chapter, name=Théorème]{theoreme}
  \ifdefined\TRACE \tracingmacros=2 \tracingonline=0 \fi
  \declaretheorem[sibling=theoreme, name=Proposition]{proposition}
  \ifdefined\TRACE \tracingmacros=0 \fi
  \declaretheorem[sibling=theoreme, name=Corollaire]{corollaire}
  \declaretheorem[numberwithin=chapter, name=Remarque]{remarque}
  \ifdefined\FAUX \newcounteralias{corollaire}{section}\newcounteralias{orphelin}{inexistant}\fi
  \ifdefined\CORR \input{correctif-a-remise}\fi
\fi
\crefname{theoreme}{théorème}{théorèmes}
\crefname{proposition}{proposition}{propositions}
\crefname{corollaire}{corollaire}{corollaires}
\crefname{remarque}{remarque}{remarques}
\begin{document}
\chapter{Premier}
\begin{proposition}\label{p1}a\end{proposition}
\begin{theoreme}\label{t1}b\end{theoreme}
\begin{remarque}\label{r1}c\end{remarque}
\begin{corollaire}\label{c1}d\end{corollaire}
\chapter{Second}
\begin{corollaire}\label{c2}e\end{corollaire}
\begin{theoreme}\label{t2}f\end{theoreme}
RENVOIS \ref{p1} \ref{t1} \ref{r1} \ref{c1} \ref{c2} \ref{t2} ; \cref{p1} ; \cref{t1} ; \cref{r1} ; \cref{c1} ; \cref{c2} ; \cref{t2} FIN
\end{document}
"                sans correctif, crochet retiré
%   lualatex -jobname=amsseul  "\def\AMSSEUL{1}% Juge de la troisième passe : document minimal écrit de zéro pour la cause A (thmtools sous LaTeX 2026-06-01).
% Variantes, par une macro posée sur la ligne de commande (deux passes chacune, depuis memoire/essai/juge3/minimal/) :
%   lualatex -jobname=nu       "\input{minimal-juge3}"                                   sans correctif
%   lualatex -jobname=corr     "\def\CORR{1}\input{minimal-juge3}"                       correctif A du préambule
%   lualatex -jobname=crochet  "\def\SANSCROCHET{1}\input{minimal-juge3}"                sans correctif, crochet retiré
%   lualatex -jobname=amsseul  "\def\AMSSEUL{1}\input{minimal-juge3}"                    amsthm seul (témoin)
%   lualatex -jobname=faux     "\def\CORR{1}\def\FAUX{1}\input{minimal-juge3}"           deux vrais défauts : doivent tomber
%   lualatex -jobname=trace    "\def\TRACE{1}\input{minimal-juge3}"                      trace de l'appel
% Les lignes du correctif sont tirées du préambule commité par extraire_correctifs.py, pas retapées.
\documentclass{scrbook}
\ifdefined\SANSCROCHET \RemoveFromHook{package/amsthm/after}[firstaid/aliascounter]\fi
\usepackage{amsthm}
\usepackage{hyperref}
\usepackage[french, nameinlink, noabbrev]{cleveref}
\ifdefined\AMSSEUL
  \newtheorem{theoreme}{Théorème}[chapter]
  \newtheorem{proposition}[theoreme]{Proposition}
  \newtheorem{corollaire}[theoreme]{Corollaire}
  \newtheorem{remarque}{Remarque}[chapter]
\else
  \usepackage{thmtools}
  \ifdefined\CORR \input{correctif-a-pose}\fi
  \declaretheorem[numberwithin=chapter, name=Théorème]{theoreme}
  \ifdefined\TRACE \tracingmacros=2 \tracingonline=0 \fi
  \declaretheorem[sibling=theoreme, name=Proposition]{proposition}
  \ifdefined\TRACE \tracingmacros=0 \fi
  \declaretheorem[sibling=theoreme, name=Corollaire]{corollaire}
  \declaretheorem[numberwithin=chapter, name=Remarque]{remarque}
  \ifdefined\FAUX \newcounteralias{corollaire}{section}\newcounteralias{orphelin}{inexistant}\fi
  \ifdefined\CORR \input{correctif-a-remise}\fi
\fi
\crefname{theoreme}{théorème}{théorèmes}
\crefname{proposition}{proposition}{propositions}
\crefname{corollaire}{corollaire}{corollaires}
\crefname{remarque}{remarque}{remarques}
\begin{document}
\chapter{Premier}
\begin{proposition}\label{p1}a\end{proposition}
\begin{theoreme}\label{t1}b\end{theoreme}
\begin{remarque}\label{r1}c\end{remarque}
\begin{corollaire}\label{c1}d\end{corollaire}
\chapter{Second}
\begin{corollaire}\label{c2}e\end{corollaire}
\begin{theoreme}\label{t2}f\end{theoreme}
RENVOIS \ref{p1} \ref{t1} \ref{r1} \ref{c1} \ref{c2} \ref{t2} ; \cref{p1} ; \cref{t1} ; \cref{r1} ; \cref{c1} ; \cref{c2} ; \cref{t2} FIN
\end{document}
"                    amsthm seul (témoin)
%   lualatex -jobname=faux     "\def\CORR{1}\def\FAUX{1}% Juge de la troisième passe : document minimal écrit de zéro pour la cause A (thmtools sous LaTeX 2026-06-01).
% Variantes, par une macro posée sur la ligne de commande (deux passes chacune, depuis memoire/essai/juge3/minimal/) :
%   lualatex -jobname=nu       "\input{minimal-juge3}"                                   sans correctif
%   lualatex -jobname=corr     "\def\CORR{1}\input{minimal-juge3}"                       correctif A du préambule
%   lualatex -jobname=crochet  "\def\SANSCROCHET{1}\input{minimal-juge3}"                sans correctif, crochet retiré
%   lualatex -jobname=amsseul  "\def\AMSSEUL{1}\input{minimal-juge3}"                    amsthm seul (témoin)
%   lualatex -jobname=faux     "\def\CORR{1}\def\FAUX{1}\input{minimal-juge3}"           deux vrais défauts : doivent tomber
%   lualatex -jobname=trace    "\def\TRACE{1}\input{minimal-juge3}"                      trace de l'appel
% Les lignes du correctif sont tirées du préambule commité par extraire_correctifs.py, pas retapées.
\documentclass{scrbook}
\ifdefined\SANSCROCHET \RemoveFromHook{package/amsthm/after}[firstaid/aliascounter]\fi
\usepackage{amsthm}
\usepackage{hyperref}
\usepackage[french, nameinlink, noabbrev]{cleveref}
\ifdefined\AMSSEUL
  \newtheorem{theoreme}{Théorème}[chapter]
  \newtheorem{proposition}[theoreme]{Proposition}
  \newtheorem{corollaire}[theoreme]{Corollaire}
  \newtheorem{remarque}{Remarque}[chapter]
\else
  \usepackage{thmtools}
  \ifdefined\CORR \input{correctif-a-pose}\fi
  \declaretheorem[numberwithin=chapter, name=Théorème]{theoreme}
  \ifdefined\TRACE \tracingmacros=2 \tracingonline=0 \fi
  \declaretheorem[sibling=theoreme, name=Proposition]{proposition}
  \ifdefined\TRACE \tracingmacros=0 \fi
  \declaretheorem[sibling=theoreme, name=Corollaire]{corollaire}
  \declaretheorem[numberwithin=chapter, name=Remarque]{remarque}
  \ifdefined\FAUX \newcounteralias{corollaire}{section}\newcounteralias{orphelin}{inexistant}\fi
  \ifdefined\CORR \input{correctif-a-remise}\fi
\fi
\crefname{theoreme}{théorème}{théorèmes}
\crefname{proposition}{proposition}{propositions}
\crefname{corollaire}{corollaire}{corollaires}
\crefname{remarque}{remarque}{remarques}
\begin{document}
\chapter{Premier}
\begin{proposition}\label{p1}a\end{proposition}
\begin{theoreme}\label{t1}b\end{theoreme}
\begin{remarque}\label{r1}c\end{remarque}
\begin{corollaire}\label{c1}d\end{corollaire}
\chapter{Second}
\begin{corollaire}\label{c2}e\end{corollaire}
\begin{theoreme}\label{t2}f\end{theoreme}
RENVOIS \ref{p1} \ref{t1} \ref{r1} \ref{c1} \ref{c2} \ref{t2} ; \cref{p1} ; \cref{t1} ; \cref{r1} ; \cref{c1} ; \cref{c2} ; \cref{t2} FIN
\end{document}
"           deux vrais défauts : doivent tomber
%   lualatex -jobname=trace    "\def\TRACE{1}% Juge de la troisième passe : document minimal écrit de zéro pour la cause A (thmtools sous LaTeX 2026-06-01).
% Variantes, par une macro posée sur la ligne de commande (deux passes chacune, depuis memoire/essai/juge3/minimal/) :
%   lualatex -jobname=nu       "\input{minimal-juge3}"                                   sans correctif
%   lualatex -jobname=corr     "\def\CORR{1}\input{minimal-juge3}"                       correctif A du préambule
%   lualatex -jobname=crochet  "\def\SANSCROCHET{1}\input{minimal-juge3}"                sans correctif, crochet retiré
%   lualatex -jobname=amsseul  "\def\AMSSEUL{1}\input{minimal-juge3}"                    amsthm seul (témoin)
%   lualatex -jobname=faux     "\def\CORR{1}\def\FAUX{1}\input{minimal-juge3}"           deux vrais défauts : doivent tomber
%   lualatex -jobname=trace    "\def\TRACE{1}\input{minimal-juge3}"                      trace de l'appel
% Les lignes du correctif sont tirées du préambule commité par extraire_correctifs.py, pas retapées.
\documentclass{scrbook}
\ifdefined\SANSCROCHET \RemoveFromHook{package/amsthm/after}[firstaid/aliascounter]\fi
\usepackage{amsthm}
\usepackage{hyperref}
\usepackage[french, nameinlink, noabbrev]{cleveref}
\ifdefined\AMSSEUL
  \newtheorem{theoreme}{Théorème}[chapter]
  \newtheorem{proposition}[theoreme]{Proposition}
  \newtheorem{corollaire}[theoreme]{Corollaire}
  \newtheorem{remarque}{Remarque}[chapter]
\else
  \usepackage{thmtools}
  \ifdefined\CORR \input{correctif-a-pose}\fi
  \declaretheorem[numberwithin=chapter, name=Théorème]{theoreme}
  \ifdefined\TRACE \tracingmacros=2 \tracingonline=0 \fi
  \declaretheorem[sibling=theoreme, name=Proposition]{proposition}
  \ifdefined\TRACE \tracingmacros=0 \fi
  \declaretheorem[sibling=theoreme, name=Corollaire]{corollaire}
  \declaretheorem[numberwithin=chapter, name=Remarque]{remarque}
  \ifdefined\FAUX \newcounteralias{corollaire}{section}\newcounteralias{orphelin}{inexistant}\fi
  \ifdefined\CORR \input{correctif-a-remise}\fi
\fi
\crefname{theoreme}{théorème}{théorèmes}
\crefname{proposition}{proposition}{propositions}
\crefname{corollaire}{corollaire}{corollaires}
\crefname{remarque}{remarque}{remarques}
\begin{document}
\chapter{Premier}
\begin{proposition}\label{p1}a\end{proposition}
\begin{theoreme}\label{t1}b\end{theoreme}
\begin{remarque}\label{r1}c\end{remarque}
\begin{corollaire}\label{c1}d\end{corollaire}
\chapter{Second}
\begin{corollaire}\label{c2}e\end{corollaire}
\begin{theoreme}\label{t2}f\end{theoreme}
RENVOIS \ref{p1} \ref{t1} \ref{r1} \ref{c1} \ref{c2} \ref{t2} ; \cref{p1} ; \cref{t1} ; \cref{r1} ; \cref{c1} ; \cref{c2} ; \cref{t2} FIN
\end{document}
"                      trace de l'appel
% Les lignes du correctif sont tirées du préambule commité par extraire_correctifs.py, pas retapées.
\documentclass{scrbook}
\ifdefined\SANSCROCHET \RemoveFromHook{package/amsthm/after}[firstaid/aliascounter]\fi
\usepackage{amsthm}
\usepackage{hyperref}
\usepackage[french, nameinlink, noabbrev]{cleveref}
\ifdefined\AMSSEUL
  \newtheorem{theoreme}{Théorème}[chapter]
  \newtheorem{proposition}[theoreme]{Proposition}
  \newtheorem{corollaire}[theoreme]{Corollaire}
  \newtheorem{remarque}{Remarque}[chapter]
\else
  \usepackage{thmtools}
  \ifdefined\CORR \input{correctif-a-pose}\fi
  \declaretheorem[numberwithin=chapter, name=Théorème]{theoreme}
  \ifdefined\TRACE \tracingmacros=2 \tracingonline=0 \fi
  \declaretheorem[sibling=theoreme, name=Proposition]{proposition}
  \ifdefined\TRACE \tracingmacros=0 \fi
  \declaretheorem[sibling=theoreme, name=Corollaire]{corollaire}
  \declaretheorem[numberwithin=chapter, name=Remarque]{remarque}
  \ifdefined\FAUX \newcounteralias{corollaire}{section}\newcounteralias{orphelin}{inexistant}\fi
  \ifdefined\CORR \input{correctif-a-remise}\fi
\fi
\crefname{theoreme}{théorème}{théorèmes}
\crefname{proposition}{proposition}{propositions}
\crefname{corollaire}{corollaire}{corollaires}
\crefname{remarque}{remarque}{remarques}
\begin{document}
\chapter{Premier}
\begin{proposition}\label{p1}a\end{proposition}
\begin{theoreme}\label{t1}b\end{theoreme}
\begin{remarque}\label{r1}c\end{remarque}
\begin{corollaire}\label{c1}d\end{corollaire}
\chapter{Second}
\begin{corollaire}\label{c2}e\end{corollaire}
\begin{theoreme}\label{t2}f\end{theoreme}
RENVOIS \ref{p1} \ref{t1} \ref{r1} \ref{c1} \ref{c2} \ref{t2} ; \cref{p1} ; \cref{t1} ; \cref{r1} ; \cref{c1} ; \cref{c2} ; \cref{t2} FIN
\end{document}
"           deux vrais défauts : doivent tomber
%   lualatex -jobname=trace    "\def\TRACE{1}% Juge de la troisième passe : document minimal écrit de zéro pour la cause A (thmtools sous LaTeX 2026-06-01).
% Variantes, par une macro posée sur la ligne de commande (deux passes chacune, depuis memoire/essai/juge3/minimal/) :
%   lualatex -jobname=nu       "% Juge de la troisième passe : document minimal écrit de zéro pour la cause A (thmtools sous LaTeX 2026-06-01).
% Variantes, par une macro posée sur la ligne de commande (deux passes chacune, depuis memoire/essai/juge3/minimal/) :
%   lualatex -jobname=nu       "\input{minimal-juge3}"                                   sans correctif
%   lualatex -jobname=corr     "\def\CORR{1}\input{minimal-juge3}"                       correctif A du préambule
%   lualatex -jobname=crochet  "\def\SANSCROCHET{1}\input{minimal-juge3}"                sans correctif, crochet retiré
%   lualatex -jobname=amsseul  "\def\AMSSEUL{1}\input{minimal-juge3}"                    amsthm seul (témoin)
%   lualatex -jobname=faux     "\def\CORR{1}\def\FAUX{1}\input{minimal-juge3}"           deux vrais défauts : doivent tomber
%   lualatex -jobname=trace    "\def\TRACE{1}\input{minimal-juge3}"                      trace de l'appel
% Les lignes du correctif sont tirées du préambule commité par extraire_correctifs.py, pas retapées.
\documentclass{scrbook}
\ifdefined\SANSCROCHET \RemoveFromHook{package/amsthm/after}[firstaid/aliascounter]\fi
\usepackage{amsthm}
\usepackage{hyperref}
\usepackage[french, nameinlink, noabbrev]{cleveref}
\ifdefined\AMSSEUL
  \newtheorem{theoreme}{Théorème}[chapter]
  \newtheorem{proposition}[theoreme]{Proposition}
  \newtheorem{corollaire}[theoreme]{Corollaire}
  \newtheorem{remarque}{Remarque}[chapter]
\else
  \usepackage{thmtools}
  \ifdefined\CORR \input{correctif-a-pose}\fi
  \declaretheorem[numberwithin=chapter, name=Théorème]{theoreme}
  \ifdefined\TRACE \tracingmacros=2 \tracingonline=0 \fi
  \declaretheorem[sibling=theoreme, name=Proposition]{proposition}
  \ifdefined\TRACE \tracingmacros=0 \fi
  \declaretheorem[sibling=theoreme, name=Corollaire]{corollaire}
  \declaretheorem[numberwithin=chapter, name=Remarque]{remarque}
  \ifdefined\FAUX \newcounteralias{corollaire}{section}\newcounteralias{orphelin}{inexistant}\fi
  \ifdefined\CORR \input{correctif-a-remise}\fi
\fi
\crefname{theoreme}{théorème}{théorèmes}
\crefname{proposition}{proposition}{propositions}
\crefname{corollaire}{corollaire}{corollaires}
\crefname{remarque}{remarque}{remarques}
\begin{document}
\chapter{Premier}
\begin{proposition}\label{p1}a\end{proposition}
\begin{theoreme}\label{t1}b\end{theoreme}
\begin{remarque}\label{r1}c\end{remarque}
\begin{corollaire}\label{c1}d\end{corollaire}
\chapter{Second}
\begin{corollaire}\label{c2}e\end{corollaire}
\begin{theoreme}\label{t2}f\end{theoreme}
RENVOIS \ref{p1} \ref{t1} \ref{r1} \ref{c1} \ref{c2} \ref{t2} ; \cref{p1} ; \cref{t1} ; \cref{r1} ; \cref{c1} ; \cref{c2} ; \cref{t2} FIN
\end{document}
"                                   sans correctif
%   lualatex -jobname=corr     "\def\CORR{1}% Juge de la troisième passe : document minimal écrit de zéro pour la cause A (thmtools sous LaTeX 2026-06-01).
% Variantes, par une macro posée sur la ligne de commande (deux passes chacune, depuis memoire/essai/juge3/minimal/) :
%   lualatex -jobname=nu       "\input{minimal-juge3}"                                   sans correctif
%   lualatex -jobname=corr     "\def\CORR{1}\input{minimal-juge3}"                       correctif A du préambule
%   lualatex -jobname=crochet  "\def\SANSCROCHET{1}\input{minimal-juge3}"                sans correctif, crochet retiré
%   lualatex -jobname=amsseul  "\def\AMSSEUL{1}\input{minimal-juge3}"                    amsthm seul (témoin)
%   lualatex -jobname=faux     "\def\CORR{1}\def\FAUX{1}\input{minimal-juge3}"           deux vrais défauts : doivent tomber
%   lualatex -jobname=trace    "\def\TRACE{1}\input{minimal-juge3}"                      trace de l'appel
% Les lignes du correctif sont tirées du préambule commité par extraire_correctifs.py, pas retapées.
\documentclass{scrbook}
\ifdefined\SANSCROCHET \RemoveFromHook{package/amsthm/after}[firstaid/aliascounter]\fi
\usepackage{amsthm}
\usepackage{hyperref}
\usepackage[french, nameinlink, noabbrev]{cleveref}
\ifdefined\AMSSEUL
  \newtheorem{theoreme}{Théorème}[chapter]
  \newtheorem{proposition}[theoreme]{Proposition}
  \newtheorem{corollaire}[theoreme]{Corollaire}
  \newtheorem{remarque}{Remarque}[chapter]
\else
  \usepackage{thmtools}
  \ifdefined\CORR \input{correctif-a-pose}\fi
  \declaretheorem[numberwithin=chapter, name=Théorème]{theoreme}
  \ifdefined\TRACE \tracingmacros=2 \tracingonline=0 \fi
  \declaretheorem[sibling=theoreme, name=Proposition]{proposition}
  \ifdefined\TRACE \tracingmacros=0 \fi
  \declaretheorem[sibling=theoreme, name=Corollaire]{corollaire}
  \declaretheorem[numberwithin=chapter, name=Remarque]{remarque}
  \ifdefined\FAUX \newcounteralias{corollaire}{section}\newcounteralias{orphelin}{inexistant}\fi
  \ifdefined\CORR \input{correctif-a-remise}\fi
\fi
\crefname{theoreme}{théorème}{théorèmes}
\crefname{proposition}{proposition}{propositions}
\crefname{corollaire}{corollaire}{corollaires}
\crefname{remarque}{remarque}{remarques}
\begin{document}
\chapter{Premier}
\begin{proposition}\label{p1}a\end{proposition}
\begin{theoreme}\label{t1}b\end{theoreme}
\begin{remarque}\label{r1}c\end{remarque}
\begin{corollaire}\label{c1}d\end{corollaire}
\chapter{Second}
\begin{corollaire}\label{c2}e\end{corollaire}
\begin{theoreme}\label{t2}f\end{theoreme}
RENVOIS \ref{p1} \ref{t1} \ref{r1} \ref{c1} \ref{c2} \ref{t2} ; \cref{p1} ; \cref{t1} ; \cref{r1} ; \cref{c1} ; \cref{c2} ; \cref{t2} FIN
\end{document}
"                       correctif A du préambule
%   lualatex -jobname=crochet  "\def\SANSCROCHET{1}% Juge de la troisième passe : document minimal écrit de zéro pour la cause A (thmtools sous LaTeX 2026-06-01).
% Variantes, par une macro posée sur la ligne de commande (deux passes chacune, depuis memoire/essai/juge3/minimal/) :
%   lualatex -jobname=nu       "\input{minimal-juge3}"                                   sans correctif
%   lualatex -jobname=corr     "\def\CORR{1}\input{minimal-juge3}"                       correctif A du préambule
%   lualatex -jobname=crochet  "\def\SANSCROCHET{1}\input{minimal-juge3}"                sans correctif, crochet retiré
%   lualatex -jobname=amsseul  "\def\AMSSEUL{1}\input{minimal-juge3}"                    amsthm seul (témoin)
%   lualatex -jobname=faux     "\def\CORR{1}\def\FAUX{1}\input{minimal-juge3}"           deux vrais défauts : doivent tomber
%   lualatex -jobname=trace    "\def\TRACE{1}\input{minimal-juge3}"                      trace de l'appel
% Les lignes du correctif sont tirées du préambule commité par extraire_correctifs.py, pas retapées.
\documentclass{scrbook}
\ifdefined\SANSCROCHET \RemoveFromHook{package/amsthm/after}[firstaid/aliascounter]\fi
\usepackage{amsthm}
\usepackage{hyperref}
\usepackage[french, nameinlink, noabbrev]{cleveref}
\ifdefined\AMSSEUL
  \newtheorem{theoreme}{Théorème}[chapter]
  \newtheorem{proposition}[theoreme]{Proposition}
  \newtheorem{corollaire}[theoreme]{Corollaire}
  \newtheorem{remarque}{Remarque}[chapter]
\else
  \usepackage{thmtools}
  \ifdefined\CORR \input{correctif-a-pose}\fi
  \declaretheorem[numberwithin=chapter, name=Théorème]{theoreme}
  \ifdefined\TRACE \tracingmacros=2 \tracingonline=0 \fi
  \declaretheorem[sibling=theoreme, name=Proposition]{proposition}
  \ifdefined\TRACE \tracingmacros=0 \fi
  \declaretheorem[sibling=theoreme, name=Corollaire]{corollaire}
  \declaretheorem[numberwithin=chapter, name=Remarque]{remarque}
  \ifdefined\FAUX \newcounteralias{corollaire}{section}\newcounteralias{orphelin}{inexistant}\fi
  \ifdefined\CORR \input{correctif-a-remise}\fi
\fi
\crefname{theoreme}{théorème}{théorèmes}
\crefname{proposition}{proposition}{propositions}
\crefname{corollaire}{corollaire}{corollaires}
\crefname{remarque}{remarque}{remarques}
\begin{document}
\chapter{Premier}
\begin{proposition}\label{p1}a\end{proposition}
\begin{theoreme}\label{t1}b\end{theoreme}
\begin{remarque}\label{r1}c\end{remarque}
\begin{corollaire}\label{c1}d\end{corollaire}
\chapter{Second}
\begin{corollaire}\label{c2}e\end{corollaire}
\begin{theoreme}\label{t2}f\end{theoreme}
RENVOIS \ref{p1} \ref{t1} \ref{r1} \ref{c1} \ref{c2} \ref{t2} ; \cref{p1} ; \cref{t1} ; \cref{r1} ; \cref{c1} ; \cref{c2} ; \cref{t2} FIN
\end{document}
"                sans correctif, crochet retiré
%   lualatex -jobname=amsseul  "\def\AMSSEUL{1}% Juge de la troisième passe : document minimal écrit de zéro pour la cause A (thmtools sous LaTeX 2026-06-01).
% Variantes, par une macro posée sur la ligne de commande (deux passes chacune, depuis memoire/essai/juge3/minimal/) :
%   lualatex -jobname=nu       "\input{minimal-juge3}"                                   sans correctif
%   lualatex -jobname=corr     "\def\CORR{1}\input{minimal-juge3}"                       correctif A du préambule
%   lualatex -jobname=crochet  "\def\SANSCROCHET{1}\input{minimal-juge3}"                sans correctif, crochet retiré
%   lualatex -jobname=amsseul  "\def\AMSSEUL{1}\input{minimal-juge3}"                    amsthm seul (témoin)
%   lualatex -jobname=faux     "\def\CORR{1}\def\FAUX{1}\input{minimal-juge3}"           deux vrais défauts : doivent tomber
%   lualatex -jobname=trace    "\def\TRACE{1}\input{minimal-juge3}"                      trace de l'appel
% Les lignes du correctif sont tirées du préambule commité par extraire_correctifs.py, pas retapées.
\documentclass{scrbook}
\ifdefined\SANSCROCHET \RemoveFromHook{package/amsthm/after}[firstaid/aliascounter]\fi
\usepackage{amsthm}
\usepackage{hyperref}
\usepackage[french, nameinlink, noabbrev]{cleveref}
\ifdefined\AMSSEUL
  \newtheorem{theoreme}{Théorème}[chapter]
  \newtheorem{proposition}[theoreme]{Proposition}
  \newtheorem{corollaire}[theoreme]{Corollaire}
  \newtheorem{remarque}{Remarque}[chapter]
\else
  \usepackage{thmtools}
  \ifdefined\CORR \input{correctif-a-pose}\fi
  \declaretheorem[numberwithin=chapter, name=Théorème]{theoreme}
  \ifdefined\TRACE \tracingmacros=2 \tracingonline=0 \fi
  \declaretheorem[sibling=theoreme, name=Proposition]{proposition}
  \ifdefined\TRACE \tracingmacros=0 \fi
  \declaretheorem[sibling=theoreme, name=Corollaire]{corollaire}
  \declaretheorem[numberwithin=chapter, name=Remarque]{remarque}
  \ifdefined\FAUX \newcounteralias{corollaire}{section}\newcounteralias{orphelin}{inexistant}\fi
  \ifdefined\CORR \input{correctif-a-remise}\fi
\fi
\crefname{theoreme}{théorème}{théorèmes}
\crefname{proposition}{proposition}{propositions}
\crefname{corollaire}{corollaire}{corollaires}
\crefname{remarque}{remarque}{remarques}
\begin{document}
\chapter{Premier}
\begin{proposition}\label{p1}a\end{proposition}
\begin{theoreme}\label{t1}b\end{theoreme}
\begin{remarque}\label{r1}c\end{remarque}
\begin{corollaire}\label{c1}d\end{corollaire}
\chapter{Second}
\begin{corollaire}\label{c2}e\end{corollaire}
\begin{theoreme}\label{t2}f\end{theoreme}
RENVOIS \ref{p1} \ref{t1} \ref{r1} \ref{c1} \ref{c2} \ref{t2} ; \cref{p1} ; \cref{t1} ; \cref{r1} ; \cref{c1} ; \cref{c2} ; \cref{t2} FIN
\end{document}
"                    amsthm seul (témoin)
%   lualatex -jobname=faux     "\def\CORR{1}\def\FAUX{1}% Juge de la troisième passe : document minimal écrit de zéro pour la cause A (thmtools sous LaTeX 2026-06-01).
% Variantes, par une macro posée sur la ligne de commande (deux passes chacune, depuis memoire/essai/juge3/minimal/) :
%   lualatex -jobname=nu       "\input{minimal-juge3}"                                   sans correctif
%   lualatex -jobname=corr     "\def\CORR{1}\input{minimal-juge3}"                       correctif A du préambule
%   lualatex -jobname=crochet  "\def\SANSCROCHET{1}\input{minimal-juge3}"                sans correctif, crochet retiré
%   lualatex -jobname=amsseul  "\def\AMSSEUL{1}\input{minimal-juge3}"                    amsthm seul (témoin)
%   lualatex -jobname=faux     "\def\CORR{1}\def\FAUX{1}\input{minimal-juge3}"           deux vrais défauts : doivent tomber
%   lualatex -jobname=trace    "\def\TRACE{1}\input{minimal-juge3}"                      trace de l'appel
% Les lignes du correctif sont tirées du préambule commité par extraire_correctifs.py, pas retapées.
\documentclass{scrbook}
\ifdefined\SANSCROCHET \RemoveFromHook{package/amsthm/after}[firstaid/aliascounter]\fi
\usepackage{amsthm}
\usepackage{hyperref}
\usepackage[french, nameinlink, noabbrev]{cleveref}
\ifdefined\AMSSEUL
  \newtheorem{theoreme}{Théorème}[chapter]
  \newtheorem{proposition}[theoreme]{Proposition}
  \newtheorem{corollaire}[theoreme]{Corollaire}
  \newtheorem{remarque}{Remarque}[chapter]
\else
  \usepackage{thmtools}
  \ifdefined\CORR \input{correctif-a-pose}\fi
  \declaretheorem[numberwithin=chapter, name=Théorème]{theoreme}
  \ifdefined\TRACE \tracingmacros=2 \tracingonline=0 \fi
  \declaretheorem[sibling=theoreme, name=Proposition]{proposition}
  \ifdefined\TRACE \tracingmacros=0 \fi
  \declaretheorem[sibling=theoreme, name=Corollaire]{corollaire}
  \declaretheorem[numberwithin=chapter, name=Remarque]{remarque}
  \ifdefined\FAUX \newcounteralias{corollaire}{section}\newcounteralias{orphelin}{inexistant}\fi
  \ifdefined\CORR \input{correctif-a-remise}\fi
\fi
\crefname{theoreme}{théorème}{théorèmes}
\crefname{proposition}{proposition}{propositions}
\crefname{corollaire}{corollaire}{corollaires}
\crefname{remarque}{remarque}{remarques}
\begin{document}
\chapter{Premier}
\begin{proposition}\label{p1}a\end{proposition}
\begin{theoreme}\label{t1}b\end{theoreme}
\begin{remarque}\label{r1}c\end{remarque}
\begin{corollaire}\label{c1}d\end{corollaire}
\chapter{Second}
\begin{corollaire}\label{c2}e\end{corollaire}
\begin{theoreme}\label{t2}f\end{theoreme}
RENVOIS \ref{p1} \ref{t1} \ref{r1} \ref{c1} \ref{c2} \ref{t2} ; \cref{p1} ; \cref{t1} ; \cref{r1} ; \cref{c1} ; \cref{c2} ; \cref{t2} FIN
\end{document}
"           deux vrais défauts : doivent tomber
%   lualatex -jobname=trace    "\def\TRACE{1}% Juge de la troisième passe : document minimal écrit de zéro pour la cause A (thmtools sous LaTeX 2026-06-01).
% Variantes, par une macro posée sur la ligne de commande (deux passes chacune, depuis memoire/essai/juge3/minimal/) :
%   lualatex -jobname=nu       "\input{minimal-juge3}"                                   sans correctif
%   lualatex -jobname=corr     "\def\CORR{1}\input{minimal-juge3}"                       correctif A du préambule
%   lualatex -jobname=crochet  "\def\SANSCROCHET{1}\input{minimal-juge3}"                sans correctif, crochet retiré
%   lualatex -jobname=amsseul  "\def\AMSSEUL{1}\input{minimal-juge3}"                    amsthm seul (témoin)
%   lualatex -jobname=faux     "\def\CORR{1}\def\FAUX{1}\input{minimal-juge3}"           deux vrais défauts : doivent tomber
%   lualatex -jobname=trace    "\def\TRACE{1}\input{minimal-juge3}"                      trace de l'appel
% Les lignes du correctif sont tirées du préambule commité par extraire_correctifs.py, pas retapées.
\documentclass{scrbook}
\ifdefined\SANSCROCHET \RemoveFromHook{package/amsthm/after}[firstaid/aliascounter]\fi
\usepackage{amsthm}
\usepackage{hyperref}
\usepackage[french, nameinlink, noabbrev]{cleveref}
\ifdefined\AMSSEUL
  \newtheorem{theoreme}{Théorème}[chapter]
  \newtheorem{proposition}[theoreme]{Proposition}
  \newtheorem{corollaire}[theoreme]{Corollaire}
  \newtheorem{remarque}{Remarque}[chapter]
\else
  \usepackage{thmtools}
  \ifdefined\CORR \input{correctif-a-pose}\fi
  \declaretheorem[numberwithin=chapter, name=Théorème]{theoreme}
  \ifdefined\TRACE \tracingmacros=2 \tracingonline=0 \fi
  \declaretheorem[sibling=theoreme, name=Proposition]{proposition}
  \ifdefined\TRACE \tracingmacros=0 \fi
  \declaretheorem[sibling=theoreme, name=Corollaire]{corollaire}
  \declaretheorem[numberwithin=chapter, name=Remarque]{remarque}
  \ifdefined\FAUX \newcounteralias{corollaire}{section}\newcounteralias{orphelin}{inexistant}\fi
  \ifdefined\CORR \input{correctif-a-remise}\fi
\fi
\crefname{theoreme}{théorème}{théorèmes}
\crefname{proposition}{proposition}{propositions}
\crefname{corollaire}{corollaire}{corollaires}
\crefname{remarque}{remarque}{remarques}
\begin{document}
\chapter{Premier}
\begin{proposition}\label{p1}a\end{proposition}
\begin{theoreme}\label{t1}b\end{theoreme}
\begin{remarque}\label{r1}c\end{remarque}
\begin{corollaire}\label{c1}d\end{corollaire}
\chapter{Second}
\begin{corollaire}\label{c2}e\end{corollaire}
\begin{theoreme}\label{t2}f\end{theoreme}
RENVOIS \ref{p1} \ref{t1} \ref{r1} \ref{c1} \ref{c2} \ref{t2} ; \cref{p1} ; \cref{t1} ; \cref{r1} ; \cref{c1} ; \cref{c2} ; \cref{t2} FIN
\end{document}
"                      trace de l'appel
% Les lignes du correctif sont tirées du préambule commité par extraire_correctifs.py, pas retapées.
\documentclass{scrbook}
\ifdefined\SANSCROCHET \RemoveFromHook{package/amsthm/after}[firstaid/aliascounter]\fi
\usepackage{amsthm}
\usepackage{hyperref}
\usepackage[french, nameinlink, noabbrev]{cleveref}
\ifdefined\AMSSEUL
  \newtheorem{theoreme}{Théorème}[chapter]
  \newtheorem{proposition}[theoreme]{Proposition}
  \newtheorem{corollaire}[theoreme]{Corollaire}
  \newtheorem{remarque}{Remarque}[chapter]
\else
  \usepackage{thmtools}
  \ifdefined\CORR \input{correctif-a-pose}\fi
  \declaretheorem[numberwithin=chapter, name=Théorème]{theoreme}
  \ifdefined\TRACE \tracingmacros=2 \tracingonline=0 \fi
  \declaretheorem[sibling=theoreme, name=Proposition]{proposition}
  \ifdefined\TRACE \tracingmacros=0 \fi
  \declaretheorem[sibling=theoreme, name=Corollaire]{corollaire}
  \declaretheorem[numberwithin=chapter, name=Remarque]{remarque}
  \ifdefined\FAUX \newcounteralias{corollaire}{section}\newcounteralias{orphelin}{inexistant}\fi
  \ifdefined\CORR \input{correctif-a-remise}\fi
\fi
\crefname{theoreme}{théorème}{théorèmes}
\crefname{proposition}{proposition}{propositions}
\crefname{corollaire}{corollaire}{corollaires}
\crefname{remarque}{remarque}{remarques}
\begin{document}
\chapter{Premier}
\begin{proposition}\label{p1}a\end{proposition}
\begin{theoreme}\label{t1}b\end{theoreme}
\begin{remarque}\label{r1}c\end{remarque}
\begin{corollaire}\label{c1}d\end{corollaire}
\chapter{Second}
\begin{corollaire}\label{c2}e\end{corollaire}
\begin{theoreme}\label{t2}f\end{theoreme}
RENVOIS \ref{p1} \ref{t1} \ref{r1} \ref{c1} \ref{c2} \ref{t2} ; \cref{p1} ; \cref{t1} ; \cref{r1} ; \cref{c1} ; \cref{c2} ; \cref{t2} FIN
\end{document}
"                      trace de l'appel
% Les lignes du correctif sont tirées du préambule commité par extraire_correctifs.py, pas retapées.
\documentclass{scrbook}
\ifdefined\SANSCROCHET \RemoveFromHook{package/amsthm/after}[firstaid/aliascounter]\fi
\usepackage{amsthm}
\usepackage{hyperref}
\usepackage[french, nameinlink, noabbrev]{cleveref}
\ifdefined\AMSSEUL
  \newtheorem{theoreme}{Théorème}[chapter]
  \newtheorem{proposition}[theoreme]{Proposition}
  \newtheorem{corollaire}[theoreme]{Corollaire}
  \newtheorem{remarque}{Remarque}[chapter]
\else
  \usepackage{thmtools}
  \ifdefined\CORR \input{correctif-a-pose}\fi
  \declaretheorem[numberwithin=chapter, name=Théorème]{theoreme}
  \ifdefined\TRACE \tracingmacros=2 \tracingonline=0 \fi
  \declaretheorem[sibling=theoreme, name=Proposition]{proposition}
  \ifdefined\TRACE \tracingmacros=0 \fi
  \declaretheorem[sibling=theoreme, name=Corollaire]{corollaire}
  \declaretheorem[numberwithin=chapter, name=Remarque]{remarque}
  \ifdefined\FAUX \newcounteralias{corollaire}{section}\newcounteralias{orphelin}{inexistant}\fi
  \ifdefined\CORR \input{correctif-a-remise}\fi
\fi
\crefname{theoreme}{théorème}{théorèmes}
\crefname{proposition}{proposition}{propositions}
\crefname{corollaire}{corollaire}{corollaires}
\crefname{remarque}{remarque}{remarques}
\begin{document}
\chapter{Premier}
\begin{proposition}\label{p1}a\end{proposition}
\begin{theoreme}\label{t1}b\end{theoreme}
\begin{remarque}\label{r1}c\end{remarque}
\begin{corollaire}\label{c1}d\end{corollaire}
\chapter{Second}
\begin{corollaire}\label{c2}e\end{corollaire}
\begin{theoreme}\label{t2}f\end{theoreme}
RENVOIS \ref{p1} \ref{t1} \ref{r1} \ref{c1} \ref{c2} \ref{t2} ; \cref{p1} ; \cref{t1} ; \cref{r1} ; \cref{c1} ; \cref{c2} ; \cref{t2} FIN
\end{document}
"                      trace de l'appel
% Les lignes du correctif sont tirées du préambule commité par extraire_correctifs.py, pas retapées.
\documentclass{scrbook}
\ifdefined\SANSCROCHET \RemoveFromHook{package/amsthm/after}[firstaid/aliascounter]\fi
\usepackage{amsthm}
\usepackage{hyperref}
\usepackage[french, nameinlink, noabbrev]{cleveref}
\ifdefined\AMSSEUL
  \newtheorem{theoreme}{Théorème}[chapter]
  \newtheorem{proposition}[theoreme]{Proposition}
  \newtheorem{corollaire}[theoreme]{Corollaire}
  \newtheorem{remarque}{Remarque}[chapter]
\else
  \usepackage{thmtools}
  \ifdefined\CORR \input{correctif-a-pose}\fi
  \declaretheorem[numberwithin=chapter, name=Théorème]{theoreme}
  \ifdefined\TRACE \tracingmacros=2 \tracingonline=0 \fi
  \declaretheorem[sibling=theoreme, name=Proposition]{proposition}
  \ifdefined\TRACE \tracingmacros=0 \fi
  \declaretheorem[sibling=theoreme, name=Corollaire]{corollaire}
  \declaretheorem[numberwithin=chapter, name=Remarque]{remarque}
  \ifdefined\FAUX \newcounteralias{corollaire}{section}\newcounteralias{orphelin}{inexistant}\fi
  \ifdefined\CORR \input{correctif-a-remise}\fi
\fi
\crefname{theoreme}{théorème}{théorèmes}
\crefname{proposition}{proposition}{propositions}
\crefname{corollaire}{corollaire}{corollaires}
\crefname{remarque}{remarque}{remarques}
\begin{document}
\chapter{Premier}
\begin{proposition}\label{p1}a\end{proposition}
\begin{theoreme}\label{t1}b\end{theoreme}
\begin{remarque}\label{r1}c\end{remarque}
\begin{corollaire}\label{c1}d\end{corollaire}
\chapter{Second}
\begin{corollaire}\label{c2}e\end{corollaire}
\begin{theoreme}\label{t2}f\end{theoreme}
RENVOIS \ref{p1} \ref{t1} \ref{r1} \ref{c1} \ref{c2} \ref{t2} ; \cref{p1} ; \cref{t1} ; \cref{r1} ; \cref{c1} ; \cref{c2} ; \cref{t2} FIN
\end{document}
